% Options for packages loaded elsewhere
\PassOptionsToPackage{unicode}{hyperref}
\PassOptionsToPackage{hyphens}{url}
\PassOptionsToPackage{dvipsnames,svgnames,x11names}{xcolor}
%
\documentclass[
  10pt,
]{article}
\usepackage{amsmath,amssymb}
\usepackage{iftex}
\ifPDFTeX
  \usepackage[T1]{fontenc}
  \usepackage[utf8]{inputenc}
  \usepackage{textcomp} % provide euro and other symbols
\else % if luatex or xetex
  \usepackage{unicode-math} % this also loads fontspec
  \defaultfontfeatures{Scale=MatchLowercase}
  \defaultfontfeatures[\rmfamily]{Ligatures=TeX,Scale=1}
\fi
\usepackage{lmodern}
\ifPDFTeX\else
  % xetex/luatex font selection
\fi
% Use upquote if available, for straight quotes in verbatim environments
\IfFileExists{upquote.sty}{\usepackage{upquote}}{}
\IfFileExists{microtype.sty}{% use microtype if available
  \usepackage[]{microtype}
  \UseMicrotypeSet[protrusion]{basicmath} % disable protrusion for tt fonts
}{}
\makeatletter
\@ifundefined{KOMAClassName}{% if non-KOMA class
  \IfFileExists{parskip.sty}{%
    \usepackage{parskip}
  }{% else
    \setlength{\parindent}{0pt}
    \setlength{\parskip}{6pt plus 2pt minus 1pt}}
}{% if KOMA class
  \KOMAoptions{parskip=half}}
\makeatother
\usepackage{xcolor}
\usepackage{longtable,booktabs,array}
\usepackage{calc} % for calculating minipage widths
% Correct order of tables after \paragraph or \subparagraph
\usepackage{etoolbox}
\makeatletter
\patchcmd\longtable{\par}{\if@noskipsec\mbox{}\fi\par}{}{}
\makeatother
% Allow footnotes in longtable head/foot
\IfFileExists{footnotehyper.sty}{\usepackage{footnotehyper}}{\usepackage{footnote}}
\makesavenoteenv{longtable}
\setlength{\emergencystretch}{3em} % prevent overfull lines
\providecommand{\tightlist}{%
  \setlength{\itemsep}{0pt}\setlength{\parskip}{0pt}}
\setcounter{secnumdepth}{5}

\usepackage{amsmath,amssymb}
\usepackage{graphicx}
\usepackage{longtable,booktabs,array}
\usepackage{xcolor}
\usepackage[margin=2.0cm]{geometry}
\usepackage{microtype}
\setlength{\emergencystretch}{3em}
\setcounter{MaxMatrixCols}{20}
\pdfinfoomitdate=1
\pdftrailerid{}
\pdfsuppressptexinfo=-1
\renewcommand{\arraystretch}{1.12}
\usepackage{fvextra}
\RecustomVerbatimEnvironment{verbatim}{Verbatim}{breaklines,breakanywhere,fontsize=\small}
\fvset{breaklines,breakanywhere}
\allowdisplaybreaks
\sloppy
\ifLuaTeX
  \usepackage{selnolig}  % disable illegal ligatures
\fi
\IfFileExists{bookmark.sty}{\usepackage{bookmark}}{\usepackage{hyperref}}
\IfFileExists{xurl.sty}{\usepackage{xurl}}{} % add URL line breaks if available
\urlstyle{same}
\hypersetup{
  pdftitle={UM-4 field equations},
  pdfauthor={Underlying theory: Patrick L. Nash; Derived and verified by Claude (Anthropic); Field equations of UM-4, 2026-10-25; Repository once-ere/Dirac\_claude, folder universe\_model\_4},
  colorlinks=true,
  linkcolor={blue},
  filecolor={Maroon},
  citecolor={Blue},
  urlcolor={blue},
  pdfcreator={LaTeX via pandoc}}

\title{UM-4 field equations}
\usepackage{etoolbox}
\makeatletter
\providecommand{\subtitle}[1]{% add subtitle to \maketitle
  \apptocmd{\@title}{\par {\large #1 \par}}{}{}
}
\makeatother
\subtitle{The complete set of coupled field equations of UM-4: Einstein
gravity in a bulk warped in all directions of the fold, dirac16complex
and the gauge fields of SU(3), SU(2) and U(1) in the bulk, the Higgs
doublet with the ad hoc vacuum and three generations of quarks and
leptons with three right-handed neutrinos on the fold}
\author{Underlying theory: Patrick L. Nash \and Derived and verified by
Claude (Anthropic) \and Field equations of UM-4,
2026-10-25 \and Repository once-ere/Dirac\_claude, folder
universe\_model\_4}
\date{}

\begin{document}
\maketitle

{
\hypersetup{linkcolor=}
\setcounter{tocdepth}{3}
\tableofcontents
}
\hypertarget{the-fold-definition-and-examples}{%
\section{The fold: definition and
examples}\label{the-fold-definition-and-examples}}

\textbf{Definition.} The \emph{fold} is the seven-dimensional surface of
the author's eight-dimensional spacetime on which the cosine of the
angle \(z = 6Hx_0\) of the hidden direction vanishes:

\[
\mathcal{F} = \big\{X = (x_0, x_1, \dots, x_7):\ \cos(6Hx_0) = 0\big\},\qquad
\text{in one period } 0 < z < \pi\text{ of the metric: } z = \frac{\pi}{2},\quad x_0 = \frac{\pi}{12H} \tag{FD-1}
\] \textbf{Symbols in (FD-1).}

\begin{itemize}
\tightlist
\item
  \(\mathcal{F}\): The fold: the seven-dimensional surface
  \(\cos(6Hx_0) = 0\) of the eight-dimensional spacetime (FD-1), in one
  period \(0 < z < \pi\) the surface \(z = \pi/2\), \(x_0 = \pi/(12H)\),
  where \(\sigma = 0\); \(\vert_{\mathcal{F}}\) means restricted to the
  fold. Status: DERIVED.
\item
  \(X = \{x_0, \dots, x_7\}\): The author's coordinates of the
  eight-dimensional spacetime, \(X = \{x_0, x_1, \dots, x_7\}\) with the
  hidden direction first (Dirac\_Einstein-Lovelock\_4+4\_Nash.nb,
  section DEFCHART); \(X^M\) is the coordinate with the index \(M\) and
  \(\mathrm{d}^8X = \mathrm{d}x_0\,\mathrm{d}x_1\cdots\mathrm{d}x_7\)
  the coordinate volume element. Unit: length (each coordinate). Status:
  AUTHOR.
\item
  \(x_0\): The hidden space-like coordinate (listed first); the metric
  depends on it through \(z = 6Hx_0\), periodically with the period
  \(\pi/(6H)\). Unit: length. Status: AUTHOR.
\item
  \(x_1\): The first coordinate of ordinary 3-space. Unit: length.
  Status: AUTHOR.
\item
  \(x_7\): The third of the three extra time-like coordinates. Unit:
  length. Status: AUTHOR.
\item
  \(H\): The author's constant (an inverse length) that fixes the period
  \(\pi/(6H)\) of \(x_0\), the warp factor
  \(\mathrm{e}^{-2H\lvert\sigma\rvert}\) of the metric, the cosmological
  constant \(\Lambda = -21H^2\) and the tension
  \(\kappa\lambda_{\mathrm{F}} = 12H\) of the fold; in UM-4 a free
  parameter, bounded from below by the propagation of the gauge fields
  in the bulk. Unit: GeV (1/length). Value: \(H > 2534\) GeV, about 2.5
  TeV (from the invisible width of the Z boson). Status: AUTHOR.
\item
  \(z\): The angle variable of the hidden direction, \(z = 6Hx_0\); one
  period of \(x_0\) is \(0 < z < \pi\), the fold is \(z = \pi/2\). Unit:
  dimensionless. Status: AUTHOR.
\item
  \(\pi\): The number pi, the ratio of the circumference of a circle to
  its diameter. Unit: dimensionless. Value: 3.141592654. Status:
  STANDARD.
\item
  \(\Lambda\): The eight-dimensional cosmological constant of UM-4,
  \(\Lambda = -21H^2\): the value for which the static block is a vacuum
  of Einstein gravity (its Einstein tensor is \(21H^2\) in every
  direction). The author writes the pure number \(\Lambda_{\mathrm{N}}\)
  in
  \(0 = -\Lambda_{\mathrm{N}} + \sum_p H^{-2p}w_{(p)}\mathcal{E}_{(p)}\),
  so that \(\Lambda = -\Lambda_{\mathrm{N}}H^2\). Unit: GeV\(^2\).
  Value: \(-21H^2\); the author's original \(-10H^2\); UM-3: \(-7H^2\).
  Status: AUTHOR (the term), DERIVED (the value).
\item
  \(\Lambda_{\mathrm{N}}\): The author's pure number in the vacuum
  equations of his notebook, \(\Lambda = -\Lambda_{\mathrm{N}}H^2\): his
  original value 10; 21 in UM-4 (Einstein gravity), 7 in UM-3. Unit:
  dimensionless. Value: 21 (the author's original 10). Status: AUTHOR.
\item
  \(w_{(p)}\): The Lovelock weights of the author's Einstein-Lovelock
  action, \(w_{(1)} = 1\) and \(w_{(2)}\), \(w_{(3)}\) of the
  Gauss-Bonnet and the cubic densities. UM-4 sets
  \(w_{(2)} = w_{(3)} = 0\) (Einstein gravity): the only weights for
  which gravitons propagate on the light cone on every background, not
  only around the static block. The author's original weights
  \(w_{(2)} = 1/32\), \(w_{(3)} = 1/1152\)
  (Dirac\_Einstein-Lovelock\_4+4\_Nash.nb, In{[}134{]}) keep the
  gravitons on the light cone only around the static block (with
  \(\Lambda = -161H^2/16\)); UM-3 used \((1, 1/20, 1/360)\). Unit:
  dimensionless. Value: \((1, 0, 0)\) in UM-4; the author's original
  \((1, 1/32, 1/1152)\). Status: AUTHOR (the form), DERIVED (the values
  that the order of 2026-10-25 selects).
\end{itemize}

\(X = \{x_0, \dots, x_7\}\) are the author's coordinates: \(x_0\) the
hidden space-like direction, \(x_1, x_2, x_3\) ordinary space, \(x_4\)
the time, \(x_5, x_6, x_7\) the three extra times; \(H\) is the author's
constant (an inverse length). The metric of UM-4 depends on \(x_0\) only
through \(\cot^2z\) and \((\sin^2z)^{1/6}\), so it repeats with the
period \(\pi/(6H)\) of \(x_0\), and every period \(0 < z < \pi\)
contains exactly one surface of \(\mathcal{F}\), at \(z = \pi/2\):
\(\mathcal{F}\) is the union of these copies, one per period (example
2), and ``the fold'' without a qualifier is the copy in the period that
contains our world. A point of the fold is labelled by the seven
coordinates \(x_1, \dots, x_7\), and the metric restricted to the fold
(\(\mathrm{d}x_0 = 0\), \(z = \pi/2\)) is

\[
g_{MN}\,\mathrm{d}X^M\mathrm{d}X^N\big\vert_{\mathcal{F}} = \mathrm{e}^{2a_4}\big(\mathrm{d}x_1^2 + \mathrm{d}x_2^2 + \mathrm{d}x_3^2\big)
 - \mathrm{d}x_4^2 - \mathrm{e}^{-2a_4}b_4^2\big(\mathrm{d}x_5^2 + \mathrm{d}x_6^2 + \mathrm{d}x_7^2\big) \tag{FD-2}
\] \textbf{Symbols in (FD-2).}

\begin{itemize}
\tightlist
\item
  \(g_{MN}\): The metric tensor of the eight-dimensional spacetime
  (signature \(+,+,+,+,-,-,-,-\) in the order \(x_0, \dots, x_7\)):
  \(g_{MN}\,\mathrm{d}X^M\mathrm{d}X^N\) is the squared interval,
  \(g^{MN}\) its inverse, \(\lvert g\rvert\) the absolute value of its
  determinant (\(\sqrt{\lvert g\rvert}\,\mathrm{d}^8X\) is the proper
  volume element). Unit: dimensionless. Status: AUTHOR.
\item
  \(M\): A coordinate index of eight dimensions (also with a numbered
  copy \(M_1, M_2, \dots\)): it runs over \(x_0, x_1, \dots, x_7\); a
  repeated index is summed. Status: STANDARD.
\item
  \(N\): A coordinate index of eight dimensions (also
  \(N_1, N_2, \dots\)), running over \(x_0, \dots, x_7\); a repeated
  index is summed. Status: STANDARD.
\item
  \(\mathrm{d}\): The differential: \(\mathrm{d}\) in front of a
  quantity is an infinitesimal change of it (as \(\mathrm{d}\sigma\), a
  small step of \(\sigma\)); the quotient of two such changes is a
  derivative; \(\mathrm{d}\) with an exponent, in front of several
  variables, is their volume element, the product of their
  differentials. Status: STANDARD.
\item
  \(X = \{x_0, \dots, x_7\}\): The author's coordinates of the
  eight-dimensional spacetime, \(X = \{x_0, x_1, \dots, x_7\}\) with the
  hidden direction first (Dirac\_Einstein-Lovelock\_4+4\_Nash.nb,
  section DEFCHART); \(X^M\) is the coordinate with the index \(M\) and
  \(\mathrm{d}^8X = \mathrm{d}x_0\,\mathrm{d}x_1\cdots\mathrm{d}x_7\)
  the coordinate volume element. Unit: length (each coordinate). Status:
  AUTHOR.
\item
  \(\mathcal{F}\): The fold: the seven-dimensional surface
  \(\cos(6Hx_0) = 0\) of the eight-dimensional spacetime (FD-1), in one
  period \(0 < z < \pi\) the surface \(z = \pi/2\), \(x_0 = \pi/(12H)\),
  where \(\sigma = 0\); \(\vert_{\mathcal{F}}\) means restricted to the
  fold. Status: DERIVED.
\item
  \(\mathrm{e}\): Euler's number, the base of the natural logarithm;
  \(\mathrm{e}^{X}\) is the exponential of \(X\). Unit: dimensionless.
  Value: 2.718281828. Status: STANDARD.
\item
  \(a_4\): The author's metric function of the time \(x_4\) that scales
  ordinary 3-space: on the fold \(\mathrm{e}^{a_4}\) is the scale factor
  of cosmology and \(a_4'\) the Hubble rate. Unit: dimensionless.
  Status: AUTHOR.
\item
  \(b_4\): The author's metric function of the time \(x_4\) in the
  extra-time components, \(b_4 > 0\), never identically 1 (the author's
  order of 2026-10-24); the scale factor (the factor of the lengths) of
  the extra times is \(\mathrm{e}^{-a_4}b_4\) on the fold and
  \(\mathrm{e}^{-a_4}b_4\,\mathrm{e}^{-H\lvert\sigma\rvert}\) at the
  distance \(\sigma\) from it (the square root of the warp factor times
  \(\mathrm{e}^{-a_4}b_4\)). Unit: dimensionless. Status: AUTHOR.
\item
  \(x_1\): The first coordinate of ordinary 3-space. Unit: length.
  Status: AUTHOR.
\item
  \(x_2\): The second coordinate of ordinary 3-space. Unit: length.
  Status: AUTHOR.
\item
  \(x_3\): The third coordinate of ordinary 3-space. Unit: length.
  Status: AUTHOR.
\item
  \(x_4\): The time; a prime on a function of \(x_4\) alone is
  \(\mathrm{d}/\mathrm{d}x_4\). Unit: length (= time, \(c = 1\)).
  Status: AUTHOR.
\item
  \(x_5\): The first of the three extra time-like coordinates. Unit:
  length. Status: AUTHOR.
\item
  \(x_6\): The second of the three extra time-like coordinates. Unit:
  length. Status: AUTHOR.
\item
  \(x_7\): The third of the three extra time-like coordinates. Unit:
  length. Status: AUTHOR.
\item
  \({}'\): The derivative of a function of the time \(x_4\) alone with
  respect to \(x_4\) (\(a_4' = \mathrm{d}a_4/\mathrm{d}x_4\); two primes
  the second derivative); on a function of \(\sigma\) alone (a profile)
  the derivative with respect to \(\sigma\). Status: STANDARD.
\end{itemize}

with three space-like and four time-like directions; \(a_4\) and \(b_4\)
are the author's two functions of the time \(x_4\). Notation: \(\ln\) is
the natural logarithm, \(\mathrm{e}\) Euler's number,
\(\exp(y) = \mathrm{e}^y\), \(\sin\), \(\cos\) and \(\cot = \cos/\sin\)
the trigonometric functions, \(\sinh\), \(\cosh\),
\(\tanh = \sinh/\cosh\) and \(\coth = \cosh/\sinh\) the hyperbolic
functions, \(\lvert y\rvert\) the absolute value, \(\max\) and \(\min\)
the larger and the smaller of their arguments; \(\det\), \(\mathrm{tr}\)
and \({}^T\) are the determinant, the trace and the transpose of a
square matrix, and \(\mathbf{1}_4\), \(\mathbf{1}_8\) and
\(\mathbf{1}_{16}\) the identity matrices of 4, 8 and 16 rows;
\(f(0^+)\) and \(f(0^-)\) are the limits of \(f(\sigma)\) as \(\sigma\)
approaches 0 from above and from below.

\textbf{Why this surface is called a fold.} Measured by the proper
distance along the hidden direction,

\[
\sigma = \mathrm{sgn}\Big(\frac{\pi}{2} - z\Big)\frac{\ln\sin z}{6H},\qquad \mathrm{d}\sigma = \lvert\cot z\rvert\,\mathrm{d}x_0,\qquad
\sin z = \mathrm{e}^{-6H\lvert\sigma\rvert} \tag{FD-3}
\] \textbf{Symbols in (FD-3).}

\begin{itemize}
\tightlist
\item
  \(\sigma\): The proper coordinate of the hidden direction:
  \(\mathrm{d}\sigma = \lvert\cot z\rvert\,\mathrm{d}x_0\),
  \(\sigma = \mathrm{sgn}(\pi/2 - z)\ln(\sin z)/(6H)\) on one period
  \(0 < z < \pi\); \(\sigma = 0\) is the fold \(z = \pi/2\),
  \(\sigma \to \mp\infty\) the tips \(z \to 0, \pi\); one period of
  \(x_0\) is the whole line \(-\infty < \sigma < \infty\). Unit: length.
  Status: DERIVED.
\item
  \(\mathrm{sgn}\): The sign function: \(+1\) for a positive argument
  and \(-1\) for a negative one. Status: STANDARD.
\item
  \(\pi\): The number pi, the ratio of the circumference of a circle to
  its diameter. Unit: dimensionless. Value: 3.141592654. Status:
  STANDARD.
\item
  \(z\): The angle variable of the hidden direction, \(z = 6Hx_0\); one
  period of \(x_0\) is \(0 < z < \pi\), the fold is \(z = \pi/2\). Unit:
  dimensionless. Status: AUTHOR.
\item
  \(H\): The author's constant (an inverse length) that fixes the period
  \(\pi/(6H)\) of \(x_0\), the warp factor
  \(\mathrm{e}^{-2H\lvert\sigma\rvert}\) of the metric, the cosmological
  constant \(\Lambda = -21H^2\) and the tension
  \(\kappa\lambda_{\mathrm{F}} = 12H\) of the fold; in UM-4 a free
  parameter, bounded from below by the propagation of the gauge fields
  in the bulk. Unit: GeV (1/length). Value: \(H > 2534\) GeV, about 2.5
  TeV (from the invisible width of the Z boson). Status: AUTHOR.
\item
  \(\mathrm{d}\): The differential: \(\mathrm{d}\) in front of a
  quantity is an infinitesimal change of it (as \(\mathrm{d}\sigma\), a
  small step of \(\sigma\)); the quotient of two such changes is a
  derivative; \(\mathrm{d}\) with an exponent, in front of several
  variables, is their volume element, the product of their
  differentials. Status: STANDARD.
\item
  \(x_0\): The hidden space-like coordinate (listed first); the metric
  depends on it through \(z = 6Hx_0\), periodically with the period
  \(\pi/(6H)\). Unit: length. Status: AUTHOR.
\item
  \(\mathrm{e}\): Euler's number, the base of the natural logarithm;
  \(\mathrm{e}^{X}\) is the exponential of \(X\). Unit: dimensionless.
  Value: 2.718281828. Status: STANDARD.
\item
  \(\infty\): Infinity: as a limit of an integral or a sum it means that
  the range has no bound. Status: STANDARD.
\end{itemize}

(the fold is \(\sigma = 0\); the half-period \(0 < z < \pi/2\) is
\(\sigma < 0\), the half-period \(\pi/2 < z < \pi\) is \(\sigma > 0\);
the two ends of the period, the tips \(z \to 0\) and \(z \to \pi\), lie
at \(\sigma \to -\infty\) and \(\sigma \to +\infty\)), the metric of
UM-4 is

\[
g_{MN}\,\mathrm{d}X^M\mathrm{d}X^N = \mathrm{d}\sigma^2 + \mathrm{e}^{-2H\lvert\sigma\rvert}\Big[\mathrm{e}^{2a_4}\big(\mathrm{d}x_1^2
 + \mathrm{d}x_2^2 + \mathrm{d}x_3^2\big) - \mathrm{d}x_4^2 - \mathrm{e}^{-2a_4}b_4^2\big(\mathrm{d}x_5^2 + \mathrm{d}x_6^2 + \mathrm{d}x_7^2\big)\Big] \tag{FD-4}
\] \textbf{Symbols in (FD-4).}

\begin{itemize}
\tightlist
\item
  \(g_{MN}\): The metric tensor of the eight-dimensional spacetime
  (signature \(+,+,+,+,-,-,-,-\) in the order \(x_0, \dots, x_7\)):
  \(g_{MN}\,\mathrm{d}X^M\mathrm{d}X^N\) is the squared interval,
  \(g^{MN}\) its inverse, \(\lvert g\rvert\) the absolute value of its
  determinant (\(\sqrt{\lvert g\rvert}\,\mathrm{d}^8X\) is the proper
  volume element). Unit: dimensionless. Status: AUTHOR.
\item
  \(M\): A coordinate index of eight dimensions (also with a numbered
  copy \(M_1, M_2, \dots\)): it runs over \(x_0, x_1, \dots, x_7\); a
  repeated index is summed. Status: STANDARD.
\item
  \(N\): A coordinate index of eight dimensions (also
  \(N_1, N_2, \dots\)), running over \(x_0, \dots, x_7\); a repeated
  index is summed. Status: STANDARD.
\item
  \(\mathrm{d}\): The differential: \(\mathrm{d}\) in front of a
  quantity is an infinitesimal change of it (as \(\mathrm{d}\sigma\), a
  small step of \(\sigma\)); the quotient of two such changes is a
  derivative; \(\mathrm{d}\) with an exponent, in front of several
  variables, is their volume element, the product of their
  differentials. Status: STANDARD.
\item
  \(X = \{x_0, \dots, x_7\}\): The author's coordinates of the
  eight-dimensional spacetime, \(X = \{x_0, x_1, \dots, x_7\}\) with the
  hidden direction first (Dirac\_Einstein-Lovelock\_4+4\_Nash.nb,
  section DEFCHART); \(X^M\) is the coordinate with the index \(M\) and
  \(\mathrm{d}^8X = \mathrm{d}x_0\,\mathrm{d}x_1\cdots\mathrm{d}x_7\)
  the coordinate volume element. Unit: length (each coordinate). Status:
  AUTHOR.
\item
  \(\sigma\): The proper coordinate of the hidden direction:
  \(\mathrm{d}\sigma = \lvert\cot z\rvert\,\mathrm{d}x_0\),
  \(\sigma = \mathrm{sgn}(\pi/2 - z)\ln(\sin z)/(6H)\) on one period
  \(0 < z < \pi\); \(\sigma = 0\) is the fold \(z = \pi/2\),
  \(\sigma \to \mp\infty\) the tips \(z \to 0, \pi\); one period of
  \(x_0\) is the whole line \(-\infty < \sigma < \infty\). Unit: length.
  Status: DERIVED.
\item
  \(\mathrm{e}\): Euler's number, the base of the natural logarithm;
  \(\mathrm{e}^{X}\) is the exponential of \(X\). Unit: dimensionless.
  Value: 2.718281828. Status: STANDARD.
\item
  \(H\): The author's constant (an inverse length) that fixes the period
  \(\pi/(6H)\) of \(x_0\), the warp factor
  \(\mathrm{e}^{-2H\lvert\sigma\rvert}\) of the metric, the cosmological
  constant \(\Lambda = -21H^2\) and the tension
  \(\kappa\lambda_{\mathrm{F}} = 12H\) of the fold; in UM-4 a free
  parameter, bounded from below by the propagation of the gauge fields
  in the bulk. Unit: GeV (1/length). Value: \(H > 2534\) GeV, about 2.5
  TeV (from the invisible width of the Z boson). Status: AUTHOR.
\item
  \(a_4\): The author's metric function of the time \(x_4\) that scales
  ordinary 3-space: on the fold \(\mathrm{e}^{a_4}\) is the scale factor
  of cosmology and \(a_4'\) the Hubble rate. Unit: dimensionless.
  Status: AUTHOR.
\item
  \(b_4\): The author's metric function of the time \(x_4\) in the
  extra-time components, \(b_4 > 0\), never identically 1 (the author's
  order of 2026-10-24); the scale factor (the factor of the lengths) of
  the extra times is \(\mathrm{e}^{-a_4}b_4\) on the fold and
  \(\mathrm{e}^{-a_4}b_4\,\mathrm{e}^{-H\lvert\sigma\rvert}\) at the
  distance \(\sigma\) from it (the square root of the warp factor times
  \(\mathrm{e}^{-a_4}b_4\)). Unit: dimensionless. Status: AUTHOR.
\item
  \(x_1\): The first coordinate of ordinary 3-space. Unit: length.
  Status: AUTHOR.
\item
  \(x_2\): The second coordinate of ordinary 3-space. Unit: length.
  Status: AUTHOR.
\item
  \(x_3\): The third coordinate of ordinary 3-space. Unit: length.
  Status: AUTHOR.
\item
  \(x_4\): The time; a prime on a function of \(x_4\) alone is
  \(\mathrm{d}/\mathrm{d}x_4\). Unit: length (= time, \(c = 1\)).
  Status: AUTHOR.
\item
  \(x_5\): The first of the three extra time-like coordinates. Unit:
  length. Status: AUTHOR.
\item
  \(x_6\): The second of the three extra time-like coordinates. Unit:
  length. Status: AUTHOR.
\item
  \(x_7\): The third of the three extra time-like coordinates. Unit:
  length. Status: AUTHOR.
\end{itemize}

It depends on \(\sigma\) only through \(\lvert\sigma\rvert\). The
function \(\mathcal{W}(\sigma) = \mathrm{e}^{-2H\lvert\sigma\rvert}\),
which multiplies the squared steps of all seven directions of the fold
(ordinary space, the time and the extra times), is the warp factor of
the metric (its square root \(\mathrm{e}^{-H\lvert\sigma\rvert}\)
multiplies the lengths). The time \(x_4\) is warped like the other
directions (the author's order of 2026-10-25, 13:15 PDT; in UM-3 it was
not), so that every surface \(\sigma = \text{const}\) has the light cone
of the fold. Four properties make the surface \(\sigma = 0\) a fold, and
together they single it out within the period:

\begin{enumerate}
\def\labelenumi{\arabic{enumi}.}
\tightlist
\item
  \emph{It is a mirror.} The reflection \(\sigma \to -\sigma\) (that is
  \(z \to \pi - z\), or \(x_0 \to \pi/(6H) - x_0\)) maps the metric to
  itself, and the fold is the only surface of the period that the
  reflection leaves in place: the two halves of the period are mirror
  images of each other, joined along the fold, as the two halves of a
  folded sheet of paper are joined along the crease.
\item
  \emph{The warp factor is largest on it.} The warp factor
  \((\sin^2z)^{1/6} = \mathrm{e}^{-2H\lvert\sigma\rvert}\) equals 1 on
  the fold and decreases on both sides.
\item
  \emph{The metric has a kink on it.} The metric is continuous across
  the fold, but its first derivative along \(\sigma\) jumps. The
  extrinsic curvature of the surfaces \(\sigma = \text{const}\) changes
  sign across the fold,
\end{enumerate}

\[
K^a{}_b = \frac12\,g^{ac}\,\partial_\sigma g_{cb} = -H\,\mathrm{sgn}(\sigma)\,\delta^a{}_b\quad(\text{all seven directions of the fold}) \tag{FD-5}
\] \textbf{Symbols in (FD-5).}

\begin{itemize}
\tightlist
\item
  \(K^a{}_b\): The extrinsic curvature of a surface of constant
  \(\sigma\) of the hidden direction,
  \(K^a{}_b = \frac12 g^{ac}\partial_\sigma g_{cb}\) for \(a, b\) along
  the surface (FD-5): \(-H\,\mathrm{sgn}(\sigma)\) along all seven
  directions of the fold for the author's form of the metric of UM-4;
  \(K = K^c{}_c\) its trace. Unit: GeV (1/length). Status: DERIVED.
\item
  \(a\): A frame (orthonormal) index running over \(0, 1, \dots, 7\)
  (frame direction \((0)\) along \(\sigma\), \((1), (2), (3)\) along
  ordinary space, \((4)\) along the time, \((5), (6), (7)\) along the
  extra times); raised and lowered with the frame metric
  \(\eta_{ab} = \mathrm{diag}(1, 1, 1, 1, -1, -1, -1, -1)\); on the fold
  (\(a, b \in \{1, \dots, 7\}\)) the frame of the fold; in the tensors
  of the fold written with its coordinates \(x^a\) (its metric
  \(\hat g_{ab}\) and the tensors built from it) a coordinate index of
  the fold, \(a \in \{1, \dots, 7\}\); a repeated index is summed.
  Status: STANDARD.
\item
  \(b\): A frame index (as \(a\)). Status: STANDARD.
\item
  \(c\): A frame index (as \(a\)). Status: STANDARD.
\item
  \(g_{MN}\): The metric tensor of the eight-dimensional spacetime
  (signature \(+,+,+,+,-,-,-,-\) in the order \(x_0, \dots, x_7\)):
  \(g_{MN}\,\mathrm{d}X^M\mathrm{d}X^N\) is the squared interval,
  \(g^{MN}\) its inverse, \(\lvert g\rvert\) the absolute value of its
  determinant (\(\sqrt{\lvert g\rvert}\,\mathrm{d}^8X\) is the proper
  volume element). Unit: dimensionless. Status: AUTHOR.
\item
  \(\partial\): A partial derivative: the derivative with respect to one
  variable at fixed values of the others (\(\partial_\sigma\) with
  respect to \(\sigma\); \(\partial_M\) and \(\partial_\mu\) with
  respect to the coordinate that carries the same index, of the eight
  dimensions or of the fold). Status: STANDARD.
\item
  \(\sigma\): The label \(\sigma\) in \(\partial_\sigma\): the
  derivative with respect to the proper hidden coordinate \(\sigma\).
  Status: STANDARD.
\item
  \(\sigma\): The proper coordinate of the hidden direction:
  \(\mathrm{d}\sigma = \lvert\cot z\rvert\,\mathrm{d}x_0\),
  \(\sigma = \mathrm{sgn}(\pi/2 - z)\ln(\sin z)/(6H)\) on one period
  \(0 < z < \pi\); \(\sigma = 0\) is the fold \(z = \pi/2\),
  \(\sigma \to \mp\infty\) the tips \(z \to 0, \pi\); one period of
  \(x_0\) is the whole line \(-\infty < \sigma < \infty\). Unit: length.
  Status: DERIVED.
\item
  \(H\): The author's constant (an inverse length) that fixes the period
  \(\pi/(6H)\) of \(x_0\), the warp factor
  \(\mathrm{e}^{-2H\lvert\sigma\rvert}\) of the metric, the cosmological
  constant \(\Lambda = -21H^2\) and the tension
  \(\kappa\lambda_{\mathrm{F}} = 12H\) of the fold; in UM-4 a free
  parameter, bounded from below by the propagation of the gauge fields
  in the bulk. Unit: GeV (1/length). Value: \(H > 2534\) GeV, about 2.5
  TeV (from the invisible width of the Z boson). Status: AUTHOR.
\item
  \(\mathrm{sgn}\): The sign function: \(+1\) for a positive argument
  and \(-1\) for a negative one. Status: STANDARD.
\item
  \(\delta\): The Kronecker delta: with two indices \(\delta^M_N = 1\)
  if \(M = N\) and 0 otherwise (likewise \(\delta_{X\mathrm{L}}\) is 1
  if the label \(X\) of a half of dirac16complex is L and 0 otherwise);
  with two index strings of equal length the generalised Kronecker delta
  \(\delta^{M_1M_2\cdots}_{N_1N_2\cdots}\), the determinant of the
  square matrix of the Kronecker deltas \(\delta^{M_a}_{N_b}\) (\(+1\)
  or \(-1\) if the upper indices are an even or odd permutation of the
  lower ones, 0 otherwise). Status: STANDARD.
\item
  \(\eta_{ab}\): The flat metric of the frame,
  \(\eta = \mathrm{diag}(1, 1, 1, 1, -1, -1, -1, -1)\) (eta4488 of
  4+4.nb); \(\eta^{ab}\) is its inverse (the same matrix). Unit:
  dimensionless. Status: AUTHOR.
\item
  \(\hat g_{ab}\): The metric of the seven-dimensional fold (FD-2) in
  its coordinates \(x_1, \dots, x_7\) (or in its frame), \(\hat g^{ab}\)
  its inverse; \(\hat h_{\mu\nu}\) is its part along the four observable
  directions. Unit: dimensionless. Status: AUTHOR.
\item
  \(\mathrm{diag}\): A diagonal matrix with the listed entries on its
  diagonal (in the order of the coordinates \(x_0, x_1, \dots, x_7\) for
  a tensor of eight dimensions, \(x_1, \dots, x_7\) for a tensor of the
  fold, \(x_1, \dots, x_4\) for a four-dimensional tensor). Status:
  STANDARD.
\item
  \(x\): The coordinates of the fold: \(x_1, x_2, x_3\) of ordinary
  space, \(x_4\) the time, \(x_5, x_6, x_7\) the extra times; a field of
  the fold is a function of the four observable coordinates
  \(x = (x_1, x_2, x_3, x_4)\), and
  \(\mathrm{d}^4x = \mathrm{d}x_1\,\mathrm{d}x_2\,\mathrm{d}x_3\,\mathrm{d}x_4\).
  Unit: length (each coordinate). Status: AUTHOR.
\end{itemize}

(\(a, b\) frame indices along the fold); it jumps by \(-2H\) in every
direction of the fold. The field equations of gravity of UM-4,
\(\mathcal{E}_{(1)}{}^M{}_N + \Lambda\delta^M_N = \kappa\mathcal{T}^M{}_N\)
(Einstein gravity: the author's Einstein-Lovelock equations with the
weights of the Gauss-Bonnet and cubic terms set to zero;
\(\mathcal{E}_{(1)}\) the Einstein tensor, \(\Lambda = -21H^2\) the
cosmological constant, \(\kappa = 8\pi G_8\) the gravitational coupling
of eight dimensions, \(\mathcal{T}^M{}_N\) the energy-momentum tensor,
\(\delta^M_N\) the Kronecker delta), are linear in the second
\(\sigma\)-derivatives of the metric, so the jump puts a Dirac delta
function \(\delta_{\mathrm{D}}(\sigma)\) into them: the metric solves
them only if the fold carries a thin source \(\mathcal{S}^a{}_b\), an
energy-momentum tensor concentrated on \(\sigma = 0\), given by the
integral of the equations across the fold (the junction condition of
Israel),

\[
\kappa\,\mathcal{S}^a{}_b = \int_{0^-}^{0^+}\mathrm{d}\sigma\,\mathcal{E}_{(1)}{}^a{}_b
 = -\Big(\big(K^a{}_b\big\vert_{0^+} - K^a{}_b\big\vert_{0^-}\big) - \delta^a{}_b\big(K\big\vert_{0^+} - K\big\vert_{0^-}\big)\Big) \tag{FD-11}
\] \textbf{Symbols in (FD-11).}

\begin{itemize}
\tightlist
\item
  \(\kappa\): The eight-dimensional gravitational coupling,
  \(\kappa = 8\pi G_8\). Unit: GeV\(^{-6}\). Status: AUTHOR.
\item
  \(\mathcal{S}^a{}_b\): The thin source of the fold: the
  energy-momentum tensor concentrated on the fold, integrated across it,
  \(\kappa\mathcal{S}^a{}_b = \int_{0^-}^{0^+}\mathrm{d}\sigma\,\mathcal{E}_{(1)}{}^a{}_b\)
  (the junction condition of Israel); \(\mathcal{S}^4{}_4\) is minus its
  energy density, \(\mathcal{S}^1{}_1\) and \(\mathcal{S}^5{}_5\) its
  pressures along ordinary space and along the extra times (frame
  components along the fold, eight-dimensional normalisation: an energy
  per unit volume of the six other directions of the fold). Unit:
  GeV\(^7\). Value: \(\kappa\mathcal{S}^a{}_b = -12H\delta^a{}_b\) for
  every \(a_4\), \(b_4\) of the author's form. Status: DERIVED.
\item
  \(a\): A frame (orthonormal) index running over \(0, 1, \dots, 7\)
  (frame direction \((0)\) along \(\sigma\), \((1), (2), (3)\) along
  ordinary space, \((4)\) along the time, \((5), (6), (7)\) along the
  extra times); raised and lowered with the frame metric
  \(\eta_{ab} = \mathrm{diag}(1, 1, 1, 1, -1, -1, -1, -1)\); on the fold
  (\(a, b \in \{1, \dots, 7\}\)) the frame of the fold; in the tensors
  of the fold written with its coordinates \(x^a\) (its metric
  \(\hat g_{ab}\) and the tensors built from it) a coordinate index of
  the fold, \(a \in \{1, \dots, 7\}\); a repeated index is summed.
  Status: STANDARD.
\item
  \(b\): A frame index (as \(a\)). Status: STANDARD.
\item
  \(\mathrm{d}\): The differential: \(\mathrm{d}\) in front of a
  quantity is an infinitesimal change of it (as \(\mathrm{d}\sigma\), a
  small step of \(\sigma\)); the quotient of two such changes is a
  derivative; \(\mathrm{d}\) with an exponent, in front of several
  variables, is their volume element, the product of their
  differentials. Status: STANDARD.
\item
  \(\sigma\): The proper coordinate of the hidden direction:
  \(\mathrm{d}\sigma = \lvert\cot z\rvert\,\mathrm{d}x_0\),
  \(\sigma = \mathrm{sgn}(\pi/2 - z)\ln(\sin z)/(6H)\) on one period
  \(0 < z < \pi\); \(\sigma = 0\) is the fold \(z = \pi/2\),
  \(\sigma \to \mp\infty\) the tips \(z \to 0, \pi\); one period of
  \(x_0\) is the whole line \(-\infty < \sigma < \infty\). Unit: length.
  Status: DERIVED.
\item
  \(\mathcal{E}_{(p)}{}^M{}_N\): The Lovelock tensor of order \(p\),
  \(\mathcal{E}_{(p)}{}^M{}_N = -2^{-(p+1)}\,\delta^{MP_1\cdots P_{2p}}_{NQ_1\cdots Q_{2p}}\,R^{Q_1Q_2}{}_{P_1P_2}\cdots R^{Q_{2p-1}Q_{2p}}{}_{P_{2p-1}P_{2p}}\)
  (with both indices down it is \(1/\sqrt{\lvert g\rvert}\) times the
  variation of \(\sqrt{\lvert g\rvert}\,\mathcal{L}_{(p)}\) with respect
  to \(g^{MN}\); \(\mathcal{E}_{(1)}\) the Einstein tensor,
  \(\mathcal{E}_{(2)}\) the Gauss-Bonnet tensor, \(\mathcal{E}_{(3)}\)
  the cubic Lovelock tensor); with frame indices its frame components.
  Unit: 1/length\(^{2p}\). Status: STANDARD.
\item
  \(p\): The Lovelock order \(p = 1, 2, 3\) (Einstein, Gauss-Bonnet,
  cubic), as an index and as a number (UM-4 keeps only \(p = 1\)); in
  the Kohn-Sham equations also a second label of a Kohn-Sham orbital,
  with the Kronecker delta of the two labels. Status: STANDARD.
\item
  \(K^a{}_b\): The extrinsic curvature of a surface of constant
  \(\sigma\) of the hidden direction,
  \(K^a{}_b = \frac12 g^{ac}\partial_\sigma g_{cb}\) for \(a, b\) along
  the surface (FD-5): \(-H\,\mathrm{sgn}(\sigma)\) along all seven
  directions of the fold for the author's form of the metric of UM-4;
  \(K = K^c{}_c\) its trace. Unit: GeV (1/length). Status: DERIVED.
\item
  \(\delta\): The Kronecker delta: with two indices \(\delta^M_N = 1\)
  if \(M = N\) and 0 otherwise (likewise \(\delta_{X\mathrm{L}}\) is 1
  if the label \(X\) of a half of dirac16complex is L and 0 otherwise);
  with two index strings of equal length the generalised Kronecker delta
  \(\delta^{M_1M_2\cdots}_{N_1N_2\cdots}\), the determinant of the
  square matrix of the Kronecker deltas \(\delta^{M_a}_{N_b}\) (\(+1\)
  or \(-1\) if the upper indices are an even or odd permutation of the
  lower ones, 0 otherwise). Status: STANDARD.
\item
  \(G_8\): The eight-dimensional Newton constant
  (\(\kappa = 8\pi G_8\)). Unit: GeV\(^{-6}\). Status: AUTHOR.
\item
  \(R^{MN}{}_{PQ},\ R_{MN},\ R\): The curvature of the metric
  \(g_{MN}\): with four indices the Riemann tensor \(R^M{}_{NPQ}\) (in
  the convention in which the commutator of two covariant derivatives,
  acting on a vector, is the Riemann tensor, with its last two indices
  those of the two derivatives, contracted with the vector), with two
  indices the Ricci tensor \(R_{MN} = R^P{}_{MPN}\), without indices the
  scalar curvature \(R = g^{MN}R_{MN}\); with frame indices
  \(R^{(a)(b)}{}_{(c)(d)}\) its frame components,
  \(R^{(a)(b)}{}_{(a)(b)}\) (no sum) the sectional curvature of the
  frame plane \((a, b)\). Unit: 1/length\(^2\). Status: STANDARD.
\item
  \(\mathcal{L}_{(p)}\): The Lovelock density of order \(p\),
  \(\mathcal{L}_{(p)} = 2^{-p}\,\delta^{P_1\cdots P_{2p}}_{Q_1\cdots Q_{2p}}\,R^{Q_1Q_2}{}_{P_1P_2}\cdots R^{Q_{2p-1}Q_{2p}}{}_{P_{2p-1}P_{2p}}\),
  a scalar built from \(p\) curvature tensors: \(\mathcal{L}_{(1)} = R\)
  (Einstein-Hilbert), \(\mathcal{L}_{(2)}\) the Gauss-Bonnet density,
  \(\mathcal{L}_{(3)}\) the cubic Lovelock density. Unit:
  1/length\(^{2p}\). Status: STANDARD.
\end{itemize}

(\(K = K^c{}_c\) the trace of the extrinsic curvature), with
\(\mathcal{S}^4{}_4 = -(\text{its energy density})\),
\(\mathcal{S}^1{}_1 = \mathcal{S}^2{}_2 = \mathcal{S}^3{}_3\) its
pressure along ordinary space and
\(\mathcal{S}^5{}_5 = \mathcal{S}^6{}_6 = \mathcal{S}^7{}_7\) its
pressure along the extra times; \(\mathcal{S}\) is \(\mathcal{T}\)
integrated across the fold, an energy per unit volume of the six other
directions of the fold (unit GeV\(^7\)). For the static block (\(a_4\)
and \(b_4\) constant) and, to the first order in the curvature of the
fold divided by \(H^2\), for every fold

\[
\begin{aligned}
&\text{static block:}\quad \kappa\mathcal{S}^a{}_b = -12H\,\delta^a{}_b\quad\text{(a pure tension)},\qquad \kappa\lambda_{\mathrm{F}} = 12H,\\
&\text{any fold:}\quad \kappa\mathcal{S}^a{}_b = -12H\,\delta^a{}_b + \frac{2}{5H}\,\hat G^a{}_b + \dots
\end{aligned} \tag{FD-12}
\] \textbf{Symbols in (FD-12).}

\begin{itemize}
\tightlist
\item
  \(\kappa\): The eight-dimensional gravitational coupling,
  \(\kappa = 8\pi G_8\). Unit: GeV\(^{-6}\). Status: AUTHOR.
\item
  \(\mathcal{S}^a{}_b\): The thin source of the fold: the
  energy-momentum tensor concentrated on the fold, integrated across it,
  \(\kappa\mathcal{S}^a{}_b = \int_{0^-}^{0^+}\mathrm{d}\sigma\,\mathcal{E}_{(1)}{}^a{}_b\)
  (the junction condition of Israel); \(\mathcal{S}^4{}_4\) is minus its
  energy density, \(\mathcal{S}^1{}_1\) and \(\mathcal{S}^5{}_5\) its
  pressures along ordinary space and along the extra times (frame
  components along the fold, eight-dimensional normalisation: an energy
  per unit volume of the six other directions of the fold). Unit:
  GeV\(^7\). Value: \(\kappa\mathcal{S}^a{}_b = -12H\delta^a{}_b\) for
  every \(a_4\), \(b_4\) of the author's form. Status: DERIVED.
\item
  \(a\): A frame (orthonormal) index running over \(0, 1, \dots, 7\)
  (frame direction \((0)\) along \(\sigma\), \((1), (2), (3)\) along
  ordinary space, \((4)\) along the time, \((5), (6), (7)\) along the
  extra times); raised and lowered with the frame metric
  \(\eta_{ab} = \mathrm{diag}(1, 1, 1, 1, -1, -1, -1, -1)\); on the fold
  (\(a, b \in \{1, \dots, 7\}\)) the frame of the fold; in the tensors
  of the fold written with its coordinates \(x^a\) (its metric
  \(\hat g_{ab}\) and the tensors built from it) a coordinate index of
  the fold, \(a \in \{1, \dots, 7\}\); a repeated index is summed.
  Status: STANDARD.
\item
  \(b\): A frame index (as \(a\)). Status: STANDARD.
\item
  \(H\): The author's constant (an inverse length) that fixes the period
  \(\pi/(6H)\) of \(x_0\), the warp factor
  \(\mathrm{e}^{-2H\lvert\sigma\rvert}\) of the metric, the cosmological
  constant \(\Lambda = -21H^2\) and the tension
  \(\kappa\lambda_{\mathrm{F}} = 12H\) of the fold; in UM-4 a free
  parameter, bounded from below by the propagation of the gauge fields
  in the bulk. Unit: GeV (1/length). Value: \(H > 2534\) GeV, about 2.5
  TeV (from the invisible width of the Z boson). Status: AUTHOR.
\item
  \(\delta\): The Kronecker delta: with two indices \(\delta^M_N = 1\)
  if \(M = N\) and 0 otherwise (likewise \(\delta_{X\mathrm{L}}\) is 1
  if the label \(X\) of a half of dirac16complex is L and 0 otherwise);
  with two index strings of equal length the generalised Kronecker delta
  \(\delta^{M_1M_2\cdots}_{N_1N_2\cdots}\), the determinant of the
  square matrix of the Kronecker deltas \(\delta^{M_a}_{N_b}\) (\(+1\)
  or \(-1\) if the upper indices are an even or odd permutation of the
  lower ones, 0 otherwise). Status: STANDARD.
\item
  \(\lambda_{\mathrm{F}}\): The tension that the static block requires
  of the fold, \(\kappa\lambda_{\mathrm{F}} = 12H\): the total vacuum
  energy of the seven-dimensional fold, made of the tension of the fold
  fluid and the vacuum energy of the fields of the fold (both have the
  pressure minus their energy density along ordinary space and along the
  extra times), per unit volume of the six other directions of the fold
  (eight-dimensional normalisation); in the four-dimensional
  normalisation it is \(\lambda_{\mathrm{F}}V_*\), and the excess of the
  total over \(\lambda_{\mathrm{F}}V_*\) is the dark energy
  \(\rho_\Lambda\). Unit: GeV\(^7\). Value: \(12H/\kappa\);
  \(\lambda_{\mathrm{F}}V_* = 30H^2M_{\mathrm{red}}^2 = 1.1\times10^{45}\)
  GeV\(^4\) for \(H = 2534\) GeV. Status: DERIVED.
\item
  \(\hat G^\mu{}_\nu\): The Einstein tensor of the metric
  \(\hat g_{ab}\) of the fold (FD-2), with frame indices along its seven
  directions; with four-dimensional indices, \(\hat G^\mu{}_\nu\), its
  components along the four observable directions, which for frozen
  extra times are those of the Einstein tensor of the four-dimensional
  metric \(\hat h_{\mu\nu}\). Unit: 1/length\(^2\). Status: STANDARD.
\item
  \(M_{\mathrm{red}}\): The reduced Planck mass,
  \(M_{\mathrm{red}} = (8\pi G_{\mathrm{N}})^{-1/2}\). Unit: GeV. Value:
  \(2.435\times10^{18}\) GeV. Status: STANDARD.
\item
  \(V_*\): The physical volume of the extra times \(x_5, x_6, x_7\) on
  the fold in the late universe,
  \(V_* = V_{\mathrm{t}}(\mathrm{e}^{-a_4}b_4)^3\) (with
  \(V_{\mathrm{t}} = \int\mathrm{d}x_5\,\mathrm{d}x_6\,\mathrm{d}x_7\)
  the coordinate volume of the extra times) at a time when
  \(\mathrm{e}^{-a_4}b_4\) is constant; it converts the
  eight-dimensional normalisation into the four-dimensional one. Unit:
  length\(^3\) (GeV\(^{-3}\)). Status: ASSUMED.
\item
  \(\hat h_{\mu\nu}\): The four-dimensional metric of the fold,
  \(\hat h_{\mu\nu} = \mathrm{diag}(\mathrm{e}^{2a_4}, \mathrm{e}^{2a_4}, \mathrm{e}^{2a_4}, -1)\)
  in the coordinates \(x_1, \dots, x_4\) (the induced metric at
  \(\sigma = 0\) without the extra times); \(\hat h^{\mu\nu}\) its
  inverse, \(\hat h\) its determinant,
  \(\sqrt{-\hat h} = \mathrm{e}^{3a_4}\); \(\hat h^{(1)}_{\mu\nu}\) its
  first-order cosmological perturbation. Unit: dimensionless. Status:
  DERIVED.
\item
  \(\rho_\Lambda\): The vacuum energy density of the fold above its
  static tension, the dark energy of UM-4:
  \(\rho_\Lambda = \rho_{\mathrm{ff}} + \rho_{\mathrm{vac}} - \lambda_{\mathrm{F}}V_*\)
  (four-dimensional normalisation). Unit: GeV\(^4\). Value:
  \(0.685\,\rho_{\mathrm{c}}^{(0)} = 2.5\times10^{-47}\) GeV\(^4\)
  (Planck 2018). Status: DERIVED.
\item
  \(G_{\mathrm{N}}\): Newton's constant of four-dimensional gravity (the
  gravity of our world, on the fold),
  \(G_{\mathrm{N}} = 1/M_{\mathrm{Pl}}^2\) with the Planck mass
  \(M_{\mathrm{Pl}} = 1.221\times10^{19}\) GeV. Unit: GeV\(^{-2}\).
  Value: \(6.708\times10^{-39}\) GeV\(^{-2}\). Status: STANDARD.
\item
  \(V_{\mathrm{t}}\): The coordinate volume of the (compact) extra
  times,
  \(V_{\mathrm{t}} = \int\mathrm{d}x_5\,\mathrm{d}x_6\,\mathrm{d}x_7\).
  Unit: length\(^3\). Status: ASSUMED.
\item
  \(\rho_{\mathrm{c}}\): The critical density, the energy density of a
  spatially flat fold at the expansion rate \(H_4\);
  \(\rho_{\mathrm{c}}^{(0)}\) its value today. Unit: GeV\(^4\). Value:
  \(\rho_{\mathrm{c}}^{(0)} = 1.0537\times10^{-5}\,h_{100}^2\)
  GeV/cm\(^3\). Status: STANDARD.
\item
  \(\rho_{\mathrm{ff}}\): The energy density of the fold fluid. Unit:
  GeV\(^4\). Value:
  \(\lambda_{\mathrm{F}}V_* + \rho_\Lambda - \rho_{\mathrm{vac}}\):
  together with the vacuum energy \(\rho_{\mathrm{vac}}\) of the fields
  of the fold it makes up the tension \(\lambda_{\mathrm{F}}V_*\) that
  the static block requires plus the dark energy \(\rho_\Lambda\).
  Status: ASSUMED (the fluid), DERIVED (the value).
\item
  \(\rho_{\mathrm{vac}}\): The vacuum energy of the fields of the fold:
  the minimum of their potential (for the Standard Model the minimum of
  the Higgs potential with its loop corrections), a tension of the
  seven-dimensional fold. Unit: GeV\(^4\). Status: STANDARD.
\item
  \(H_4\): The expansion rate (Hubble rate) of the fold, \(H_4 = a_4'\),
  as a function of the time or of the redshift; \(H_4^{(0)}\) its value
  today, the Hubble constant. Unit: GeV (1/time). Value:
  \(H_4^{(0)} = 100\,h_{100}\) km/s/Mpc with \(h_{100} = 0.674\) (Planck
  2018): 67.4 km/s/Mpc \(= 1.44\times10^{-42}\) GeV. Status: STANDARD.
\item
  \(M_{\mathrm{Pl}}\): The Planck mass,
  \(M_{\mathrm{Pl}} = G_{\mathrm{N}}^{-1/2}\). Unit: GeV. Value:
  \(1.221\times10^{19}\) GeV. Status: STANDARD.
\item
  \(h_{100}\): The present expansion rate of the fold in units of 100
  km/s/Mpc. Unit: dimensionless. Value: 0.674 (Planck 2018). Status:
  STANDARD.
\end{itemize}

(\(\hat G^a{}_b\) the Einstein tensor of the metric FD-2 of the fold):
the tension plus matter that obeys the Einstein equations of seven
dimensions, \(\hat G^a{}_b = \kappa_7\mathsf{t}^a{}_b\) with
\(\kappa_7 = 5\kappa H/2\) (\(\mathsf{t}^a{}_b\) the matter of the
fold), which for frozen extra times are the Friedmann equations of
general relativity. The fold is therefore not only a place but a
physical object: a thin layer that carries energy and momentum (a
brane). 4. \emph{The author's coordinate \(x_0\) degenerates on it.} The
metric component \(g_{x_0x_0} = \cot^2(6Hx_0)\) vanishes on the fold and
nowhere else in the period; near the fold the proper distance grows with
the square of the coordinate distance,

\[
\lvert\sigma\rvert = 3H\Big(x_0 - \frac{\pi}{12H}\Big)^2\Big[1 + 6H^2\Big(x_0 - \frac{\pi}{12H}\Big)^2 + \dots\Big] \tag{FD-6}
\] \textbf{Symbols in (FD-6).}

\begin{itemize}
\tightlist
\item
  \(\sigma\): The proper coordinate of the hidden direction:
  \(\mathrm{d}\sigma = \lvert\cot z\rvert\,\mathrm{d}x_0\),
  \(\sigma = \mathrm{sgn}(\pi/2 - z)\ln(\sin z)/(6H)\) on one period
  \(0 < z < \pi\); \(\sigma = 0\) is the fold \(z = \pi/2\),
  \(\sigma \to \mp\infty\) the tips \(z \to 0, \pi\); one period of
  \(x_0\) is the whole line \(-\infty < \sigma < \infty\). Unit: length.
  Status: DERIVED.
\item
  \(H\): The author's constant (an inverse length) that fixes the period
  \(\pi/(6H)\) of \(x_0\), the warp factor
  \(\mathrm{e}^{-2H\lvert\sigma\rvert}\) of the metric, the cosmological
  constant \(\Lambda = -21H^2\) and the tension
  \(\kappa\lambda_{\mathrm{F}} = 12H\) of the fold; in UM-4 a free
  parameter, bounded from below by the propagation of the gauge fields
  in the bulk. Unit: GeV (1/length). Value: \(H > 2534\) GeV, about 2.5
  TeV (from the invisible width of the Z boson). Status: AUTHOR.
\item
  \(x_0\): The hidden space-like coordinate (listed first); the metric
  depends on it through \(z = 6Hx_0\), periodically with the period
  \(\pi/(6H)\). Unit: length. Status: AUTHOR.
\item
  \(\pi\): The number pi, the ratio of the circumference of a circle to
  its diameter. Unit: dimensionless. Value: 3.141592654. Status:
  STANDARD.
\end{itemize}

so the two coordinate values on either side of \(x_0 = \pi/(12H)\) at
the same coordinate distance lie at the same proper distance from the
fold: the coordinate line of \(x_0\) is folded onto the line of proper
distances (a fold map, example 5).

\textbf{What lives on the fold.} In UM-4 the Higgs doublet (with the
author's ad hoc vacuum, a time-dependent choice of its vacuum value),
three generations of quarks and leptons, three right-handed neutrinos
and the fold fluid (its tension and the pressure along the extra times
that keeps their volume fixed) live on the fold: they are fields at
\(\sigma = 0\), functions of \(x_1, \dots, x_4\) (no field depends on
the extra times). The gauge fields of
\(SU(3)_C \times SU(2)_L \times U(1)_Y\) live in the bulk; their
massless modes are bound to the fold, so that the photon and the gluons
of the four-dimensional world are waves of the bulk that the fold holds,
and photons that leave the fold travel through the bulk at the speed of
light. dirac16complex (the author's spinor field of eight dimensions)
and gravity live in the bulk \(\sigma \ne 0\) as well as on the fold;
the states \(\chi\) of dirac16complex that the fold holds (in UM-4
long-lived resonances) are the part of it near the fold, and the portals
(couplings between dirac16complex and the fields of the fold) act at
\(\sigma = 0\). Integrated over the extra times, the fold is the
four-dimensional universe that we observe, with the metric
\(\mathrm{diag}(\mathrm{e}^{2a_4}, \mathrm{e}^{2a_4}, \mathrm{e}^{2a_4}, -1)\)
in \(x_1, \dots, x_4\). In UM-4 files ``the fold'' means this surface;
the examples below use the word for other surfaces with the same
properties. The word is not used in its other senses: a set of equal
energy levels is called degenerate, and a growth by the factor
\(\mathrm{e}\) is called an e-folding.

\textbf{Examples of folds that are meaningful for UM-4.}

\begin{enumerate}
\def\labelenumi{\arabic{enumi}.}
\tightlist
\item
  \emph{Our fold.} The fold \(z = \pi/2\) of the period that contains
  our world (FD-1): the seven-dimensional world-volume of the quarks,
  leptons and the Higgs field, and after the integration over the extra
  times the four-dimensional universe that we observe.
\item
  \emph{The copies of the fold.} Because the metric repeats with the
  period \(\pi/(6H)\), every surface
  \(x_0 = \pi/(12H) + (\text{an integer multiple of } \pi/(6H))\) is a
  copy of the fold, one per period. Neighbouring copies are separated by
  a tip (\(z\) a multiple of \(\pi\)), where the warp factor
  \((\sin^2z)^{1/6}\) vanishes and \(g_{x_0x_0} = \cot^2z\) diverges,
  and which lies at an infinite proper distance from both: in proper
  distance every copy has its own complete line
  \(-\infty < \sigma < \infty\). A compact direction of length \(L\)
  would give every field a tower of particles with the masses
  \(2\pi n/L\), \(n = 1, 2, \dots\) (a Kaluza-Klein tower); because the
  copies are infinitely far apart, the period \(\pi/(6H)\) of \(x_0\)
  produces no such tower with the spacing \(12H\): every field of the
  bulk has instead a continuum of modes that starts at zero mass, and
  the fold holds its massless mode.
\item
  \emph{A folded sheet of paper.} The crease of a folded sheet is a fold
  in this sense: the two halves are mirror images (property 1) joined
  along the crease, the sheet is continuous there and its direction
  changes abruptly (property 3). In UM-4 the direction across the crease
  is the hidden direction \(\sigma\), the seven directions along the
  crease are \(x_1, \dots, x_7\), and the crease is the fold.
\item
  \emph{A folded function.} The warp factor
  \(\mathrm{e}^{-2H\lvert\sigma\rvert}\) is the simplest folded
  function: continuous, mirror symmetric, largest at \(\sigma = 0\),
  with a jump of its derivative,
\end{enumerate}

\[
\frac{\mathrm{d}}{\mathrm{d}\sigma}\mathrm{e}^{-2H\lvert\sigma\rvert} = -2H\,\mathrm{sgn}(\sigma)\,\mathrm{e}^{-2H\lvert\sigma\rvert},\qquad
\frac{\mathrm{d}^2}{\mathrm{d}\sigma^2}\mathrm{e}^{-2H\lvert\sigma\rvert} = 4H^2\,\mathrm{e}^{-2H\lvert\sigma\rvert} - 4H\,\delta_{\mathrm{D}}(\sigma) \tag{FD-7}
\] \textbf{Symbols in (FD-7).}

\begin{itemize}
\tightlist
\item
  \(\mathrm{d}\): The differential: \(\mathrm{d}\) in front of a
  quantity is an infinitesimal change of it (as \(\mathrm{d}\sigma\), a
  small step of \(\sigma\)); the quotient of two such changes is a
  derivative; \(\mathrm{d}\) with an exponent, in front of several
  variables, is their volume element, the product of their
  differentials. Status: STANDARD.
\item
  \(\sigma\): The proper coordinate of the hidden direction:
  \(\mathrm{d}\sigma = \lvert\cot z\rvert\,\mathrm{d}x_0\),
  \(\sigma = \mathrm{sgn}(\pi/2 - z)\ln(\sin z)/(6H)\) on one period
  \(0 < z < \pi\); \(\sigma = 0\) is the fold \(z = \pi/2\),
  \(\sigma \to \mp\infty\) the tips \(z \to 0, \pi\); one period of
  \(x_0\) is the whole line \(-\infty < \sigma < \infty\). Unit: length.
  Status: DERIVED.
\item
  \(\mathrm{e}\): Euler's number, the base of the natural logarithm;
  \(\mathrm{e}^{X}\) is the exponential of \(X\). Unit: dimensionless.
  Value: 2.718281828. Status: STANDARD.
\item
  \(H\): The author's constant (an inverse length) that fixes the period
  \(\pi/(6H)\) of \(x_0\), the warp factor
  \(\mathrm{e}^{-2H\lvert\sigma\rvert}\) of the metric, the cosmological
  constant \(\Lambda = -21H^2\) and the tension
  \(\kappa\lambda_{\mathrm{F}} = 12H\) of the fold; in UM-4 a free
  parameter, bounded from below by the propagation of the gauge fields
  in the bulk. Unit: GeV (1/length). Value: \(H > 2534\) GeV, about 2.5
  TeV (from the invisible width of the Z boson). Status: AUTHOR.
\item
  \(\mathrm{sgn}\): The sign function: \(+1\) for a positive argument
  and \(-1\) for a negative one. Status: STANDARD.
\item
  \(\delta_{\mathrm{D}}\): The Dirac delta function of \(\sigma\): zero
  for \(\sigma \ne 0\), with the integral 1 over any range that contains
  \(\sigma = 0\); the integral of \(\delta_{\mathrm{D}}(\sigma)\) times
  a smooth function of \(\sigma\) is the value of that function at
  \(\sigma = 0\). Unit: 1/length (for \(\delta_{\mathrm{D}}(\sigma)\)).
  Status: STANDARD.
\end{itemize}

whose delta function is the origin of the thin source of property 3.
Every metric function of FD-4 depends on \(\sigma\) through
\(\lvert\sigma\rvert\) and is folded at \(\sigma = 0\) in this sense. 5.
\emph{A fold map.} In the mathematics of smooth maps the fold (Whitney's
fold singularity) is the map \(u \mapsto u^2\) of a line onto a
half-line: it folds the line at \(u = 0\), two points \(\pm u\) land on
the same point, and the derivative of the map vanishes on the fold. Near
the fold the map from the author's coordinate \(x_0\) to the distance
\(\lvert\sigma\rvert\) from the fold is of this kind (FD-6), and
\(g_{x_0x_0} = \cot^2z\), the square of the derivative
\(\mathrm{d}\sigma/\mathrm{d}x_0\), vanishes there. 6. \emph{The brane
of Randall and Sundrum.} The five-dimensional metric

\[
\mathrm{d}y_{\mathrm{RS}}^2 + \mathrm{e}^{-2k_{\mathrm{RS}}\lvert y_{\mathrm{RS}}\rvert}\big(\mathrm{d}x_1^2 + \mathrm{d}x_2^2 + \mathrm{d}x_3^2
 - \mathrm{d}x_4^2\big),\qquad \lambda_{\mathrm{RS}} = \frac{6k_{\mathrm{RS}}}{\kappa_5},\qquad \Lambda_{\mathrm{RS}} = -6k_{\mathrm{RS}}^2 \tag{FD-8}
\] \textbf{Symbols in (FD-8).}

\begin{itemize}
\tightlist
\item
  \(\mathrm{d}\): The differential: \(\mathrm{d}\) in front of a
  quantity is an infinitesimal change of it (as \(\mathrm{d}\sigma\), a
  small step of \(\sigma\)); the quotient of two such changes is a
  derivative; \(\mathrm{d}\) with an exponent, in front of several
  variables, is their volume element, the product of their
  differentials. Status: STANDARD.
\item
  \(y_{\mathrm{RS}}\): The proper coordinate across the brane of the
  five-dimensional model of Randall and Sundrum (FD-8, an example, not
  UM-4). Unit: length. Status: STANDARD.
\item
  \(\mathrm{e}\): Euler's number, the base of the natural logarithm;
  \(\mathrm{e}^{X}\) is the exponential of \(X\). Unit: dimensionless.
  Value: 2.718281828. Status: STANDARD.
\item
  \(k_{\mathrm{RS}}\): The inverse curvature length of the bulk of the
  model of Randall and Sundrum: the warp factor of the brane is
  \(\mathrm{e}^{-2k_{\mathrm{RS}}\lvert y_{\mathrm{RS}}\rvert}\) (FD-8,
  an example, not UM-4). Unit: GeV (1/length). Status: STANDARD.
\item
  \(x_1\): The first coordinate of ordinary 3-space. Unit: length.
  Status: AUTHOR.
\item
  \(x_2\): The second coordinate of ordinary 3-space. Unit: length.
  Status: AUTHOR.
\item
  \(x_3\): The third coordinate of ordinary 3-space. Unit: length.
  Status: AUTHOR.
\item
  \(x_4\): The time; a prime on a function of \(x_4\) alone is
  \(\mathrm{d}/\mathrm{d}x_4\). Unit: length (= time, \(c = 1\)).
  Status: AUTHOR.
\item
  \(\lambda_{\mathrm{RS}}\): The tension (energy per unit volume of
  ordinary space) of the brane of the model of Randall and Sundrum,
  \(\lambda_{\mathrm{RS}} = 6k_{\mathrm{RS}}/\kappa_5\) (FD-8, an
  example, not UM-4). Unit: GeV\(^4\). Status: STANDARD.
\item
  \(\kappa_5\): The gravitational coupling of five dimensions in the
  model of Randall and Sundrum, \(8\pi\) times the Newton constant of
  five dimensions (FD-8, an example, not UM-4). Unit: GeV\(^{-3}\).
  Status: STANDARD.
\item
  \(\Lambda_{\mathrm{RS}}\): The cosmological constant of the bulk of
  the model of Randall and Sundrum,
  \(\Lambda_{\mathrm{RS}} = -6k_{\mathrm{RS}}^2\), the vacuum energy
  density \(\Lambda_{\mathrm{RS}}/\kappa_5\) (FD-8, an example, not
  UM-4). Unit: GeV\(^2\). Status: STANDARD.
\end{itemize}

has a fold at \(y_{\mathrm{RS}} = 0\) with the properties 1 to 3 (its
coordinate \(y_{\mathrm{RS}}\) is already the proper distance, so
property 4 does not arise); Einstein's equations require there a brane
of the tension \(\lambda_{\mathrm{RS}}\) (Israel's junction condition)
and in the bulk the cosmological constant \(\Lambda_{\mathrm{RS}}\) (the
vacuum energy \(\Lambda_{\mathrm{RS}}/\kappa_5\)), and the bulk is
otherwise empty. The fold of UM-4 is of the same kind, with every
direction of the fold warped by the same factor as in the model of
Randall and Sundrum: its static block is a space of constant curvature
\(-H^2\) in all eight directions, a vacuum of Einstein gravity with
\(\Lambda = -21H^2\), with a fold of pure tension (FD-12), and the
graviton is bound to the fold as there. It differs in that the fold has
seven dimensions, three of them extra times, in that the space is not
cut at the brane (Randall and Sundrum identify the two sides of their
brane, one physical side; the fold of UM-4 has two distinct sides,
example 7), and in that the gauge fields live in the bulk as well. 7.
\emph{A particle held by a thin layer.} A particle held by an attractive
thin layer has a wave function that falls off exponentially with the
distance on both sides, and the layer makes its derivative jump (for the
first-order Dirac equation, the wave function itself). In UM-4 the
states \(\chi\) of dirac16complex near the fold are of this kind only
near the fold: at rest and without their self-interaction (the field is
\(\Psi = \mathrm{e}^{(7/2)H\lvert\sigma\rvert}\chi(\sigma)\mathrm{e}^{-\mathrm{i}\omega x_4}\))

\[
\begin{aligned}
&\chi(\sigma) = \mathrm{e}^{-q_{\mathrm{b}}\lvert\sigma\rvert}\chi(0^+)\ \ (0 < \sigma \ll \sigma_{\mathrm{t}}),\qquad
\chi(\sigma) = \mathrm{e}^{-q_{\mathrm{b}}\lvert\sigma\rvert}\chi(0^-)\ \ (0 < -\sigma \ll \sigma_{\mathrm{t}}),\\
&\chi(0^+) = \exp\Big(\gamma^{(0)}\big(\mu_{\mathrm{L}}P_{\mathrm{L}} + \mu_{\mathrm{R}}P_{\mathrm{R}}\big)\Big)\chi(0^-),\qquad
\omega_0 = \omega_{\mathrm{r}} - \frac{\mathrm{i}}{2}\,\Gamma_\chi,\qquad \omega_{\mathrm{r}} \to \omega_{\mathrm{b}} = \sqrt{m_\Psi^2 - q_{\mathrm{b}}^2}
\end{aligned} \tag{FD-9}
\] \textbf{Symbols in (FD-9).}

\begin{itemize}
\tightlist
\item
  \(\chi(\sigma)\): The profile of a mode of dirac16complex along the
  hidden direction,
  \(\Psi = \mathrm{e}^{\frac72H\lvert\sigma\rvert}\chi(\sigma)\mathrm{e}^{\mathrm{i}(\vec k\cdot\vec x - \omega x_4)}\)
  (a 16-component spinor function of \(\sigma\)); as a four-dimensional
  field, \(\chi\) denotes the particle of the level of dirac16complex at
  the fold (a long-lived resonance) and \(\bar\chi\) its antiparticle.
  Unit: GeV\(^{2}\) (with the \(V_*\) normalisation). Status: DERIVED.
\item
  \(\sigma\): The proper coordinate of the hidden direction:
  \(\mathrm{d}\sigma = \lvert\cot z\rvert\,\mathrm{d}x_0\),
  \(\sigma = \mathrm{sgn}(\pi/2 - z)\ln(\sin z)/(6H)\) on one period
  \(0 < z < \pi\); \(\sigma = 0\) is the fold \(z = \pi/2\),
  \(\sigma \to \mp\infty\) the tips \(z \to 0, \pi\); one period of
  \(x_0\) is the whole line \(-\infty < \sigma < \infty\). Unit: length.
  Status: DERIVED.
\item
  \(\mathrm{e}\): Euler's number, the base of the natural logarithm;
  \(\mathrm{e}^{X}\) is the exponential of \(X\). Unit: dimensionless.
  Value: 2.718281828. Status: STANDARD.
\item
  \(q_{\mathrm{b}}\): The inverse width of a state of dirac16complex at
  the fold: inside the turning point its profile falls as
  \(\mathrm{e}^{-q_{\mathrm{b}}\lvert\sigma\rvert}\). Unit: GeV. Status:
  DERIVED.
\item
  \(\sigma_{\mathrm{t}}\): The turning point of a state of
  dirac16complex of the frequency \(\omega\): the depth at which the
  warped frequency equals the bulk mass,
  \(\lvert\omega\rvert\mathrm{e}^{H\sigma_{\mathrm{t}}} = m_\Psi\);
  beyond it every solution oscillates. Unit: length. Status: DERIVED.
\item
  \(\gamma^M,\ \gamma^{(a)}\): The Dirac matrices of dirac16complex
  along the coordinates, \(\gamma^M = \mathbf{e}_a{}^MT^a\),
  \(\gamma_M = g_{MN}\gamma^N\); \(\gamma^{(a)} = T^a\) along the frame
  direction \((a)\) (\(\gamma^{(0)}\) along \(\sigma\), \(\gamma^{(4)}\)
  along the time). Unit: dimensionless. Status: AUTHOR.
\item
  \(\mu_X,\ \mu_{\mathrm{L}},\ \mu_{\mathrm{R}}\): The dimensionless
  strengths of the fold for the left-handed and right-handed halves of
  dirac16complex, set by the mean fields of the Standard Model at the
  fold: the jump of the profile across the fold is the exponential of
  \(\gamma^{(0)}(\mu_{\mathrm{L}}P_{\mathrm{L}} + \mu_{\mathrm{R}}P_{\mathrm{R}})\);
  negative values attract. Unit: dimensionless. Status: DERIVED.
\item
  \(P_{\mathrm{L}}\): The author's left-handed chiral projector,
  \(P_{\mathrm{L}} = (\mathbf{1} - \Gamma_8)/2\). Unit: dimensionless.
  Status: AUTHOR.
\item
  \(P_{\mathrm{R}}\): The author's right-handed chiral projector,
  \(P_{\mathrm{R}} = (\mathbf{1} + \Gamma_8)/2\). Unit: dimensionless.
  Status: AUTHOR.
\item
  \(\omega_0\): The energy at rest (zero momentum along ordinary space
  and along the extra times) of a state of dirac16complex at the fold;
  complex for a resonance,
  \(\omega_0 = \omega_{\mathrm{r}} - \frac{\mathrm{i}}{2}\Gamma_\chi\).
  Unit: GeV. Status: DERIVED.
\item
  \(\omega_{\mathrm{r}}\): The mass of the particle \(\chi\): the real
  part of the energy at rest of the resonance of dirac16complex at the
  fold,
  \(\omega_{\mathrm{r}} = \omega_{\mathrm{b}}(1 - H/(2q_{\mathrm{b}})) + \dots\).
  Unit: GeV. Status: DERIVED.
\item
  \(\mathrm{i}\): The imaginary unit, \(\mathrm{i}^2 = -1\). Unit:
  dimensionless. Status: STANDARD.
\item
  \(\Gamma_\chi\): The width of the particle \(\chi\) at rest: the rate
  of its decay into the waves of dirac16complex in the bulk. Unit: GeV.
  Value: below \(1/(11.6\,t_0) = 1.3\times10^{-43}\) GeV for dark matter
  (a lifetime above 160 Gyr). Status: DERIVED.
\item
  \(\omega_{\mathrm{b}}\): The energy at zero momentum of the state of
  dirac16complex at the fold in the limit \(H \to 0\) (the bound state
  of UM-3): the particle \(\chi\) of UM-4 has the mass
  \(\omega_{\mathrm{r}} = \omega_{\mathrm{b}}(1 - H/(2q_{\mathrm{b}})) + \dots\),
  equal to \(\omega_{\mathrm{b}}\) up to the relative correction
  \(H/(2q_{\mathrm{b}})\). Unit: GeV. Status: DERIVED.
\item
  \(m_\Psi\): The mass parameter of dirac16complex in the bulk. Unit:
  GeV. Status: RECORD.
\item
  \(\Gamma_8\): The chirality matrix of eight dimensions,
  \(\Gamma_8 = T_0T_1\cdots T_7\) (the author's T16\^{}A{[}8{]})
  \(= \mathrm{diag}(-\mathbf{1}_8, \mathbf{1}_8)\);
  \(\Gamma_8^2 = \mathbf{1}\); it anticommutes with every \(T_a\). Unit:
  dimensionless. Status: AUTHOR.
\item
  \(\Psi,\ \bar\Psi\): dirac16complex: the author's 16-component complex
  spinor field of eight dimensions;
  \(\bar\Psi = \Psi^\dagger\sigma_{16}\) is its Dirac adjoint; in a term
  of the fold, \(\Psi\) is its value at \(\sigma = 0\). Unit:
  GeV\(^{7/2}\). Status: AUTHOR.
\item
  \(\mathbf{1},\ \mathbf{1}_4,\ \mathbf{1}_8,\ \mathbf{1}_{16}\): The
  identity matrix of the size fixed by the context (\(16 \times 16\) on
  dirac16complex, \(4 \times 4\) on a four-dimensional spinor,
  \(3 \times 3\) in generation or flavour space); \(\mathbf{1}_4\),
  \(\mathbf{1}_8\) and \(\mathbf{1}_{16}\) are the identity matrices of
  4, 8 and 16 rows. Unit: dimensionless. Status: STANDARD.
\item
  \(\mathbf{e}_a{}^M,\ \mathbf{e}^a{}_M\): The vielbein (frame field) of
  the metric: \(g_{MN} = \eta_{ab}\,\mathbf{e}^a{}_M\mathbf{e}^b{}_N\);
  \(\mathbf{e}_a{}^M\) is its inverse;
  \(\mathbf{e}_{bN} = \eta_{bc}\mathbf{e}^c{}_N\); \(\mathbf{e}^{(a)}\)
  is the frame one-form \(\mathbf{e}^a{}_M\,\mathrm{d}X^M\). Unit:
  dimensionless. Status: STANDARD.
\item
  \(\omega\): A frequency: the energy of a mode
  \(\propto \mathrm{e}^{-\mathrm{i}\omega x_4}\) (\(\omega_o\) of the
  orbital \(o\)). Unit: GeV. Status: STANDARD.
\item
  \(\vec k\): The (comoving) momentum along ordinary space, \(\vec k^2\)
  its square. Unit: GeV. Status: STANDARD.
\item
  \(\vec x\): The position in ordinary space,
  \(\vec x = (x_1, x_2, x_3)\). Unit: length. Status: AUTHOR.
\item
  \(t_0\): The age of the universe. Unit: time. Value: 13.8 Gyr. Status:
  STANDARD.
\item
  \(\sigma_{16}\): The matrix of the Dirac adjoint,
  \(\sigma_{16} = T_0T_1T_2T_3\), real and symmetric; \(\sigma_{16}T_a\)
  is real and antisymmetric. Unit: dimensionless. Status: AUTHOR.
\item
  \({}^\dagger\): Hermitian conjugation: the complex conjugate of the
  transpose of a matrix or spinor; for an operator on the space of
  states, its adjoint with respect to the scalar product of that space
  (for the quantised dirac16complex the indefinite Krein scalar
  product). Status: STANDARD.
\item
  \(o\): The label of a single-particle state of dirac16complex: a
  Kohn-Sham orbital, or a mode of the quantised field; a sum over \(o\)
  runs over the orbitals or the modes. Status: STANDARD.
\end{itemize}

where \(\mu_{\mathrm{L}}\), \(\mu_{\mathrm{R}}\) are the strengths of
the fold for the left- and right-handed halves of dirac16complex (set by
the fields of the fold),
\(q_{\mathrm{b}} = -m_\Psi(\mu_{\mathrm{L}} + \mu_{\mathrm{R}})\tanh\varsigma/(2\varsigma)\),
\(\varsigma = \sqrt{\mu_{\mathrm{L}}\mu_{\mathrm{R}}}\)
(\(q_{\mathrm{b}} = -m_\Psi\mu_{\mathrm{L}}/2\) for
\(\mu_{\mathrm{R}} = 0\)), and \(\omega_{\mathrm{r}}\) tends to the
energy \(\omega_{\mathrm{b}}\) of the bound state of UM-3 when
\(H/q_{\mathrm{b}} \to 0\). Because the time is warped, the energy
\(\omega\) of the field enters its equation multiplied by
\(\mathrm{e}^{H\lvert\sigma\rvert}\): beyond the turning point
\(\sigma_{\mathrm{t}}\)
(\(\omega_{\mathrm{r}}\mathrm{e}^{H\sigma_{\mathrm{t}}} = m_\Psi\)) the
profile oscillates as a wave that leaves the fold, so the state is a
resonance of the width
\(\Gamma_\chi = (2/\mathrm{e})(q_{\mathrm{b}}^2/m_\Psi)\,\mathrm{e}^{-2F(\upsilon_{\mathrm{b}})m_\Psi/H}\)
(to the leading order in \(H/q_{\mathrm{b}}\)), with
\(F(\upsilon) = \ln((1 + \sqrt{1 - \upsilon^2})/\upsilon) - \sqrt{1 - \upsilon^2}\)
and \(\upsilon_{\mathrm{b}} = \omega_{\mathrm{b}}/m_\Psi\) (the
tunnelling through the region between the fold and
\(\sigma_{\mathrm{t}}\), where the profile decays at the local rate
\(\sqrt{m_\Psi^2 - \omega_{\mathrm{r}}^2\mathrm{e}^{2H\lvert\sigma\rvert}}\),
which falls to 0 at \(\sigma_{\mathrm{t}}\)). The profile is not a
function of \(\lvert\sigma\rvert\): it jumps at the fold,
\(\chi(0^+) \ne \chi(0^-)\). 8. \emph{The fold as a mirror for
dirac16complex.} The reflection across the fold acts on the profiles
\(\chi(\sigma)\) in two ways,

\[
(\mathsf{R}_1\chi)(\sigma) = T_0T_1T_2T_3\,\chi(-\sigma),\qquad (\mathsf{R}_2\chi)(\sigma) = T_0T_5T_6T_7\,\chi(-\sigma) \tag{FD-10}
\] \textbf{Symbols in (FD-10).}

\begin{itemize}
\tightlist
\item
  \(\mathsf{R}_1\): The reflection across the fold acting on a profile
  of dirac16complex together with the reflection of ordinary space,
  \((\mathsf{R}_1\chi)(\sigma) = T_0T_1T_2T_3\,\chi(-\sigma)\) (FD-10):
  a symmetry of dirac16complex, of gravity and of the mean field of the
  gas of the particles \(\chi\); the states at the fold of Krein charge
  \(-1\) are even, those of Krein charge \(+1\) odd. Unit:
  dimensionless. Status: DERIVED.
\item
  \(\mathsf{R}_2\): The reflection across the fold acting on a profile
  of dirac16complex together with the reflection of the extra times,
  \((\mathsf{R}_2\chi)(\sigma) = T_0T_5T_6T_7\,\chi(-\sigma)\) (FD-10):
  it maps the states at the fold of one Krein charge onto those of the
  other. Unit: dimensionless. Status: DERIVED.
\item
  \(\chi(\sigma)\): The profile of a mode of dirac16complex along the
  hidden direction,
  \(\Psi = \mathrm{e}^{\frac72H\lvert\sigma\rvert}\chi(\sigma)\mathrm{e}^{\mathrm{i}(\vec k\cdot\vec x - \omega x_4)}\)
  (a 16-component spinor function of \(\sigma\)); as a four-dimensional
  field, \(\chi\) denotes the particle of the level of dirac16complex at
  the fold (a long-lived resonance) and \(\bar\chi\) its antiparticle.
  Unit: GeV\(^{2}\) (with the \(V_*\) normalisation). Status: DERIVED.
\item
  \(\sigma\): The proper coordinate of the hidden direction:
  \(\mathrm{d}\sigma = \lvert\cot z\rvert\,\mathrm{d}x_0\),
  \(\sigma = \mathrm{sgn}(\pi/2 - z)\ln(\sin z)/(6H)\) on one period
  \(0 < z < \pi\); \(\sigma = 0\) is the fold \(z = \pi/2\),
  \(\sigma \to \mp\infty\) the tips \(z \to 0, \pi\); one period of
  \(x_0\) is the whole line \(-\infty < \sigma < \infty\). Unit: length.
  Status: DERIVED.
\item
  \(T_a,\ T^a\): The author's eight real \(16 \times 16\) Dirac matrices
  T16\^{}A{[}a{]} of 4+4.nb (\(a = 0, \dots, 7\); \(T_0, \dots, T_7\)),
  written out in FD-13: \(T_a\) has the block \(\bar\Xi_a\) in its upper
  right and \(\Xi_a\) in its lower left corner;
  \(\{T_a, T_b\} = 2\eta_{ab}\mathbf{1}_{16}\); \(T^a = \eta^{ab}T_b\).
  Unit: dimensionless. Status: AUTHOR.
\end{itemize}

(for states at rest \(\mathsf{R}_1\) is the reflection across the fold
combined with the reflection of ordinary space, \(\mathsf{R}_2\) the
reflection across the fold combined with the reflection of the extra
times). \(T_a\) (\(a = 0, \dots, 7\)) are the author's real
\(16 \times 16\) Dirac matrices \texttt{T16\^{}A{[}a{]}} of
\texttt{4+4.nb}, built there from six real \(4 \times 4\) matrices:

\[
\begin{aligned}
&\mathbf{s}_1 = \begin{pmatrix}0&0&0&1\\0&0&1&0\\0&-1&0&0\\-1&0&0&0\end{pmatrix},\quad
 \mathbf{s}_2 = \begin{pmatrix}0&0&-1&0\\0&0&0&1\\1&0&0&0\\0&-1&0&0\end{pmatrix},\quad
 \mathbf{s}_3 = \begin{pmatrix}0&1&0&0\\-1&0&0&0\\0&0&0&1\\0&0&-1&0\end{pmatrix},\\
&\mathbf{t}_1 = \begin{pmatrix}0&0&0&-1\\0&0&1&0\\0&-1&0&0\\1&0&0&0\end{pmatrix},\quad
 \mathbf{t}_2 = \begin{pmatrix}0&0&-1&0\\0&0&0&-1\\1&0&0&0\\0&1&0&0\end{pmatrix},\quad
 \mathbf{t}_3 = \begin{pmatrix}0&1&0&0\\-1&0&0&0\\0&0&0&-1\\0&0&1&0\end{pmatrix},\\
&\Xi_0 = \mathbf{1}_8,\quad
 \Xi_1 = \begin{pmatrix}0&\mathbf{s}_1\\ \mathbf{s}_1&0\end{pmatrix},\quad
 \Xi_2 = \begin{pmatrix}0&\mathbf{s}_2\\ \mathbf{s}_2&0\end{pmatrix},\quad
 \Xi_3 = \begin{pmatrix}0&\mathbf{s}_3\\ \mathbf{s}_3&0\end{pmatrix},\\
&\Xi_4 = \begin{pmatrix}0&\mathbf{t}_3\\ -\mathbf{t}_3&0\end{pmatrix},\quad
 \Xi_5 = \begin{pmatrix}0&\mathbf{t}_2\\ -\mathbf{t}_2&0\end{pmatrix},\quad
 \Xi_6 = \begin{pmatrix}0&\mathbf{t}_1\\ -\mathbf{t}_1&0\end{pmatrix},\quad
 \Xi_7 = \Xi_1\Xi_2\Xi_3\Xi_4\Xi_5\Xi_6 = \begin{pmatrix}-\mathbf{1}_4&0\\0&\mathbf{1}_4\end{pmatrix},\\
&\bar\Xi_0 = \mathbf{1}_8,\quad \bar\Xi_a = \sigma_8\,\Xi_a^T\,\sigma_8,\quad
 \sigma_8 = \begin{pmatrix}0&\mathbf{1}_4\\ \mathbf{1}_4&0\end{pmatrix},\qquad
 T_a = \begin{pmatrix}0&\bar\Xi_a\\ \Xi_a&0\end{pmatrix}
\end{aligned} \tag{FD-13}
\] \textbf{Symbols in (FD-13).}

\begin{itemize}
\tightlist
\item
  \(\mathbf{s}_1,\ \mathbf{s}_2,\ \mathbf{s}_3\): The author's three
  real antisymmetric \(4 \times 4\) matrices s4by4{[}1{]}, s4by4{[}2{]},
  s4by4{[}3{]} of 4+4.nb (written out in FD-13), the blocks of
  \(\Xi_1\), \(\Xi_2\), \(\Xi_3\). Unit: dimensionless. Status: AUTHOR.
\item
  \(\mathbf{t}_1,\ \mathbf{t}_2,\ \mathbf{t}_3\): The author's three
  real antisymmetric \(4 \times 4\) matrices t4by4{[}1{]}, t4by4{[}2{]},
  t4by4{[}3{]} of 4+4.nb (written out in FD-13), the blocks of
  \(\Xi_6\), \(\Xi_5\), \(\Xi_4\). Unit: dimensionless. Status: AUTHOR.
\item
  \(\Xi_a\): The author's eight real \(8 \times 8\) matrices tau{[}a{]}
  of 4+4.nb (\(a = 0, \dots, 7\); \(\Xi_0 = \mathbf{1}_8\),
  \(\Xi_7 = \Xi_1\Xi_2\cdots\Xi_6\)), the lower left blocks of the Dirac
  matrices \(T_a\) (FD-13). Unit: dimensionless. Status: AUTHOR.
\item
  \(\bar\Xi_a\): The author's eight real \(8 \times 8\) matrices
  taubar{[}a{]} of 4+4.nb, \(\bar\Xi_0 = \mathbf{1}_8\) and
  \(\bar\Xi_a = \sigma_8\Xi_a^T\sigma_8\) for \(a = 1, \dots, 7\), the
  upper right blocks of the Dirac matrices \(T_a\) (FD-13). Unit:
  dimensionless. Status: AUTHOR.
\item
  \(\mathbf{1},\ \mathbf{1}_4,\ \mathbf{1}_8,\ \mathbf{1}_{16}\): The
  identity matrix of the size fixed by the context (\(16 \times 16\) on
  dirac16complex, \(4 \times 4\) on a four-dimensional spinor,
  \(3 \times 3\) in generation or flavour space); \(\mathbf{1}_4\),
  \(\mathbf{1}_8\) and \(\mathbf{1}_{16}\) are the identity matrices of
  4, 8 and 16 rows. Unit: dimensionless. Status: STANDARD.
\item
  \(\sigma_8\): The author's real symmetric \(8 \times 8\) matrix
  sigma88 of 4+4.nb, with the blocks \(\mathbf{1}_4\) off the diagonal
  and zeros on it (FD-13); \(\sigma_8^2 = \mathbf{1}_8\) and
  \(\sigma_{16} = \mathrm{diag}(-\sigma_8, \sigma_8)\). Unit:
  dimensionless. Status: AUTHOR.
\item
  \({}^T\): The transpose of a matrix. Status: STANDARD.
\item
  \(T_a,\ T^a\): The author's eight real \(16 \times 16\) Dirac matrices
  T16\^{}A{[}a{]} of 4+4.nb (\(a = 0, \dots, 7\); \(T_0, \dots, T_7\)),
  written out in FD-13: \(T_a\) has the block \(\bar\Xi_a\) in its upper
  right and \(\Xi_a\) in its lower left corner;
  \(\{T_a, T_b\} = 2\eta_{ab}\mathbf{1}_{16}\); \(T^a = \eta^{ab}T_b\).
  Unit: dimensionless. Status: AUTHOR.
\item
  \(a\): A frame (orthonormal) index running over \(0, 1, \dots, 7\)
  (frame direction \((0)\) along \(\sigma\), \((1), (2), (3)\) along
  ordinary space, \((4)\) along the time, \((5), (6), (7)\) along the
  extra times); raised and lowered with the frame metric
  \(\eta_{ab} = \mathrm{diag}(1, 1, 1, 1, -1, -1, -1, -1)\); on the fold
  (\(a, b \in \{1, \dots, 7\}\)) the frame of the fold; in the tensors
  of the fold written with its coordinates \(x^a\) (its metric
  \(\hat g_{ab}\) and the tensors built from it) a coordinate index of
  the fold, \(a \in \{1, \dots, 7\}\); a repeated index is summed.
  Status: STANDARD.
\item
  \(\eta_{ab}\): The flat metric of the frame,
  \(\eta = \mathrm{diag}(1, 1, 1, 1, -1, -1, -1, -1)\) (eta4488 of
  4+4.nb); \(\eta^{ab}\) is its inverse (the same matrix). Unit:
  dimensionless. Status: AUTHOR.
\item
  \(\sigma_{16}\): The matrix of the Dirac adjoint,
  \(\sigma_{16} = T_0T_1T_2T_3\), real and symmetric; \(\sigma_{16}T_a\)
  is real and antisymmetric. Unit: dimensionless. Status: AUTHOR.
\item
  \(\mathrm{diag}\): A diagonal matrix with the listed entries on its
  diagonal (in the order of the coordinates \(x_0, x_1, \dots, x_7\) for
  a tensor of eight dimensions, \(x_1, \dots, x_7\) for a tensor of the
  fold, \(x_1, \dots, x_4\) for a four-dimensional tensor). Status:
  STANDARD.
\end{itemize}

(\(\bar\Xi_a = \sigma_8\Xi_a^T\sigma_8\) for \(a = 1, \dots, 7\), and
\(T_a\) for \(a = 0, \dots, 7\); a matrix of four blocks, each block
\(4 \times 4\) or \(8 \times 8\), 0 a block of zeros;
\(\mathbf{s}_1, \mathbf{s}_2, \mathbf{s}_3\) and
\(\mathbf{t}_1, \mathbf{t}_2, \mathbf{t}_3\) are
\texttt{s4by4{[}1..3{]}} and \texttt{t4by4{[}1..3{]}}, \(\Xi_a\) and
\(\bar\Xi_a\) are \texttt{tau{[}a{]}} and \texttt{taubar{[}a{]}},
\(\sigma_8\) is \texttt{sigma88} of the notebook). They obey
\(\{T_a, T_b\} = T_aT_b + T_bT_a = 2\eta_{ab}\mathbf{1}_{16}\) with
\(\eta = \mathrm{diag}(1, 1, 1, 1, -1, -1, -1, -1)\);
\(\gamma^{(a)} = \eta^{aa}T_a\) (no sum) are the Dirac matrices of the
frame, \(\gamma^{(0)}\) along \(\sigma\); the chirality matrix is
\(\Gamma_8 = T_0T_1\cdots T_7 =  \mathrm{diag}(-\mathbf{1}_8, \mathbf{1}_8)\),
the chiral projectors are
\(P_{\mathrm{L}} = (\mathbf{1}_{16} - \Gamma_8)/2\) and
\(P_{\mathrm{R}} = (\mathbf{1}_{16} + \Gamma_8)/2\), and
\(\sigma_{16} = T_0T_1T_2T_3 = \mathrm{diag}(-\sigma_8, \sigma_8)\)
defines the Dirac adjoint \(\bar\Psi = \Psi^\dagger\sigma_{16}\)
(\(\Psi^\dagger\) the complex conjugate transpose). The Krein charge of
a state is its share of the conserved charge of dirac16complex, which is
indefinite: a state can carry the charge \(-1\) or \(+1\).
\(\mathsf{R}_1\) is a symmetry of dirac16complex, of gravity and of a
mean field of a gas of the particles \(\chi\) that is the same on both
sides of the fold: the four resonances of Krein charge \(-1\) (the
particles \(\chi\)) are even under it and the four of Krein charge
\(+1\) are odd, so such a mean field cannot mix the two signs of the
indefinite charge. \(\mathsf{R}_2\) maps the states of one Krein charge
onto those of the other and makes the two sets degenerate without the
self-interaction. 9. \emph{The planes of Horava and Witten.} In M-theory
(the eleven-dimensional theory whose limits include the superstring
theories) on an interval (a circle whose points are identified with
their mirror images), the two ten-dimensional planes left in place by
the reflection are mirror surfaces (property 1); each carries one
\(E_8\) gauge group while gravity fills the bulk. As in UM-4, gravity
fills the bulk; unlike UM-4, the gauge fields live on the mirror planes
(in UM-4 they live in the bulk), and the two sides of each plane are
identified (one physical side).

\textbf{What is not a fold.}

\begin{itemize}
\tightlist
\item
  \emph{The tips} \(z = 0\) and \(z = \pi\) (mod \(\pi\)): the mirror
  exchanges them within a period, but there the warp factor vanishes,
  \(g_{x_0x_0}\) diverges, and they lie at an infinite proper distance;
  they carry no thin source.
\item
  \emph{A smooth maximum of a warp factor}, such as
  \(1/\cosh^2(H\sigma)\): mirror symmetric and largest at
  \(\sigma = 0\), but without a kink and without a thin source.
\end{itemize}

\hypertarget{how-to-read-this-document}{%
\section{How to read this document}\label{how-to-read-this-document}}

\hypertarget{what-this-document-contains}{%
\subsection{What this document
contains}\label{what-this-document-contains}}

This document states the complete set of coupled field equations of the
universe model UM-4, derives every one of them from one action, and
explains every term. It is self-contained: every equation it uses is
written in it, and every symbol of every equation is defined right after
the equation. The sections are:

\begin{itemize}
\tightlist
\item
  \textbf{Why UM-4} (this section): what was physically incorrect in
  UM-3, what the author ordered on 2026-10-25 at 13:15 PDT, and how UM-4
  meets the order.
\item
  \textbf{Where the fields live} (this section): the metric of UM-4 has
  a fold, the seven-dimensional surface \(z = \pi/2\), that carries
  energy; the Higgs doublet, the quarks and the leptons live on the
  fold, gravity, dirac16complex and the gauge fields in the bulk.
\item
  \textbf{The action} (IN-1 to IN-16): every term of the action, written
  out and explained; why the gravity of UM-4 is Einstein gravity (FE-68
  to FE-70); the gravitational waves of the fold (FE-63), the gravity
  that the fields of the fold induce (FE-64) and the fold fluid (FE-65).
\item
  \textbf{Gravity} (FE-1 to FE-11): the Einstein equations in the bulk,
  the junction condition at the fold, the energy-momentum tensor of
  every field and the conservation laws.
\item
  \textbf{The matter fields} (FE-12 to FE-26): the field equations of
  dirac16complex, of the gauge fields in the bulk, of the Higgs doublet
  and of the quarks and leptons, and the equations of the ad hoc Higgs
  vacuum.
\item
  \textbf{The background} (FE-27 to FE-35, FE-66, FE-71, FE-72): the
  metric of UM-4, its exact curvature, the static block of constant
  curvature, the exact bulk equations of the author's form of the
  metric, the bulk metric and the equations of an expanding fold, the
  volume of the extra times.
\item
  \textbf{Propagation in the bulk} (FE-73 to FE-84): the modes of every
  field of the bulk, the transverse spin-1 photon, the gravitational
  wave of spin 2 with its + and \(\times\) polarisations, the
  propagators from one point of the fold to another, Newton's and
  Coulomb's laws, the W and Z bosons, and the causal structure.
\item
  \textbf{dirac16complex: modes and Kohn-Sham equations} (FE-36 to
  FE-43, FE-67, FE-85, FE-86): the equation of the hidden direction, its
  exact solutions, the resonances at the fold, the Kohn-Sham equations
  of a gas of them, the Krein norm and its quantisation.
\item
  \textbf{The Standard Model on the fold} (FE-44 to FE-47): the
  four-dimensional couplings, the masses with the ad hoc vacuum, the
  mixing matrices and the four-dimensional couplings of the particles
  \(\chi\) of dirac16complex.
\item
  \textbf{The three neutrino flavours} (FE-48 to FE-53): the flavour
  evolution in vacuum, in matter, in dense neutrino gases (supernovae)
  and in the early Universe.
\item
  \textbf{The cosmological perturbations} (FE-54 to FE-58): the metric
  of the fold and the Einstein-Boltzmann equations of the cosmic
  microwave background (CMB), the cosmic neutrino background (CNB) and
  the baryon acoustic oscillations (BAO).
\item
  \textbf{The thermal history} (FE-59 to FE-62, FE-87, FE-88): the
  expansion of the fold, the degrees of freedom, the Boltzmann equations
  of the particles \(\chi\) and of the right-handed neutrinos, and the
  primordial era.
\item
  \textbf{The coupled system}: which field appears in which equation,
  and how the system is solved.
\item
  \textbf{What was verified}: the checks of the computations behind
  every derived statement.
\end{itemize}

\hypertarget{why-um-4}{%
\subsection{Why UM-4}\label{why-um-4}}

In UM-3 the time \(x_4\) was not warped: its metric was
\(\mathrm{d}\sigma^2 + \mathrm{e}^{-2H\lvert\sigma\rvert}\big(\mathrm{e}^{2a_4}\,\mathrm{d}\vec x^2 - \mathrm{e}^{-2a_4}b_4^2\,\mathrm{d}\vec x_5^2\big) - \mathrm{d}x_4^2\)
(with
\(\mathrm{d}\vec x^2 = \mathrm{d}x_1^2 + \mathrm{d}x_2^2 + \mathrm{d}x_3^2\)
and
\(\mathrm{d}\vec x_5^2 = \mathrm{d}x_5^2 + \mathrm{d}x_6^2 + \mathrm{d}x_7^2\)),
so that at the depth \(\sigma\) the light cone of the bulk was
\(\mathrm{e}^{H\lvert\sigma\rvert}\) times wider than that of the fold:
whatever propagated in the bulk outran light on the fold (bulk scalars
1.22 times, bulk photons 1.41 times the speed of light on the fold;
gravitons could be made luminal on the fold only by removing them from
the bulk). On 2026-10-25 at 13:15 PDT the author ordered:

\begin{quote}
``Since in the UM-3 model anything that propagates in the bulk outruns
light on the fold, and bulk scalars and photons faster than light, my
model is physically incorrect. You must find a new metric model in which
x4 is warped like the other directions, both photons and gravitons
actually propagate at c in the bulk, photon waves are transverse spin-1
phenomena, Gravitational waves (ripples in spacetime) are spin-2
phenomena (because they originate from the rank-2 stress-energy tensor)
and possess the following Polarization Modes: the spin-2 nature of
gravitational waves manifests physically as characteristic''plus'' (+)
and ``cross'' (×) polarization modes that stretch and squeeze space
perpendicularly. Find the best model and re-calculate and re-document;
however keep UM-3 and create a new UM-4, placed in new folders.''
\end{quote}

UM-4 meets the order as follows (UM-3 is kept unchanged in its folder
\texttt{universe\_model\_3}).

\begin{itemize}
\tightlist
\item
  \textbf{The metric.} The time is warped like the other six directions
  of the fold (FE-27, FD-4): the author's metric with the single change
  \(g_{x_4x_4} = -(\sin^2z)^{1/6}\) (UM-3: \(-1\)). Every slice
  \(\sigma = \text{const}\) then has the light cone of the fold (FE-30).
  This is the smallest change of the author's metric that meets the
  order, and it keeps everything else of the metric: the coordinates,
  the period, the fold, the functions \(a_4\) and \(b_4\).
\item
  \textbf{The gravity.} Among the author's Einstein-Lovelock theories,
  only Einstein gravity (the weights of the Gauss-Bonnet and of the
  cubic Lovelock terms zero) lets gravitons propagate at the speed of
  light on every background, not only around one (FE-68 to FE-70); with
  it the static block is a vacuum for \(\Lambda = -21H^2\). UM-4
  therefore uses Einstein gravity with \(\Lambda = -21H^2\) (IN-2). The
  author's original weights \((1, 1/32, 1/1152)\) would also give
  gravitons at the speed of light around the static block (with
  \(\Lambda = -161H^2/16\)), but not around an expanding fold.
\item
  \textbf{The fields of the bulk.} Gravity, dirac16complex and the gauge
  fields of \(SU(3)_C \times SU(2)_L \times U(1)_Y\) live in the bulk;
  their massless modes are bound to the fold and travel at the speed of
  light at every depth and every wavelength, their massive modes are
  slower (FE-73 to FE-78). The photon is a transverse wave of helicity
  \(\pm1\) (FE-75); the gravitational wave of the four-dimensional world
  is a transverse traceless wave of helicity \(\pm2\) with the two
  polarisations \(+\) and \(\times\), which stretch and squeeze a ring
  of test masses along perpendicular directions (FE-76); of the other
  polarisations of the graviton of the bulk the three vectors
  \(h_{\mu r}\) are ghosts, which the matter of the fold does not emit
  at the linear order. Nothing that crosses the static bulk at fixed
  extra times arrives earlier than light on the fold, and on an
  expanding fold the bulk is ahead by an unobservable amount (FE-83,
  FE-84); along the extra times the light cone has no such bound (four
  time-like directions).
\item
  \textbf{The consequences.} The fold is a seven-dimensional brane of
  pure tension with \(\kappa\lambda_{\mathrm{F}} = 12H\), and to the
  first order in its curvature it obeys the Einstein equations of seven
  dimensions (FE-71), which with frozen extra times are the Friedmann
  equations of general relativity. Newton's and Coulomb's laws get
  corrections of relative size \(0.34/(Hr)^5\) and \(0.42/(Hr)^3\), and
  the W and Z bosons a width into the bulk (FE-81, FE-82), which bound
  \(H\) from below by 2534 GeV. The states of dirac16complex at the
  fold, bound in UM-3, are resonances in UM-4 (FE-86) with the width
  \((2/\mathrm{e})(q_{\mathrm{b}}^2/m_\Psi)\exp(-2F(\upsilon_{\mathrm{b}})m_\Psi/H)\):
  they live longer than 160 Gyr (the bound for dark matter that decays
  into invisible particles) only if \(m_\Psi/H\) is large, about
  \(10^3\) for \(q_{\mathrm{b}} = m_\Psi/2\), so that the particles
  \(\chi\) of dark matter are heavier than \(2.5\times10^6\) GeV for
  that binding; for a strong binding the bound falls (\(1.2\times10^4\)
  GeV for \(q_{\mathrm{b}} = 0.99\,m_\Psi\)), and only particles bound
  more strongly than \(q_{\mathrm{b}} \approx 0.9\,m_\Psi\) can be a
  thermal relic (FE-61). The portals of IN-16 cannot provide such a
  binding: bound by the Higgs portal or by the gluon condensate the
  particles cannot live that long and escape direct detection at the
  same time, so the dark matter of UM-4 needs a binding of \(\chi\) to
  the fold that does not couple it to nucleons (OPEN).
\end{itemize}

\hypertarget{the-status-labels}{%
\subsection{The status labels}\label{the-status-labels}}

Every equation and every result carries one of the labels

\begin{itemize}
\tightlist
\item
  \textbf{AUTHOR}: from the author (his notebooks \texttt{4+4.nb} and
  \texttt{Dirac\_Einstein-Lovelock\_4+4\_Nash.nb}, and his orders of
  2026-10-24 and 2026-10-25: the coordinates, the metric with its change
  of 2026-10-25 13:15 (the time warped), the Einstein-Lovelock action,
  the Dirac matrices, the chiral projectors, the list of fields, the
  Higgs potential and the ad hoc vacuum);
\item
  \textbf{RECORD}: from the record of this repository (the Lagrangian of
  dirac16complex with its canonical spin connection and its contact
  self-interaction, the Kohn-Sham treatment);
\item
  \textbf{STANDARD}: textbook physics;
\item
  \textbf{ASSUMED}: a choice made in UM-4 without derivation;
\item
  \textbf{INVENTED}: a new interaction (the portals of dirac16complex);
\item
  \textbf{DERIVED}: computed and verified in UM-4 (the verification is
  stated);
\item
  \textbf{ESTIMATE}: an order-of-magnitude evaluation;
\item
  \textbf{OPEN}: not settled.
\end{itemize}

\hypertarget{every-equation-is-followed-by-its-symbols}{%
\subsection{Every equation is followed by its
symbols}\label{every-equation-is-followed-by-its-symbols}}

Every display equation is followed by the list ``Symbols in (FE-n)'' (or
``Symbols in this equation'' for a display without a tag) of EVERY
symbol that appears in it: every quantity, field, coupling, constant and
operator, every index letter with the values it takes, every mark
(prime, dagger, transpose, star), the imaginary unit, Euler's number and
the differential. Each item gives the meaning, the unit (natural units,
in which the reduced Planck constant \(\hbar\) and the speed of light
\(c\) are 1), the value where one exists, and the status. A symbol that
one of these meanings uses is itself listed in this document. A field
equation carries the tag FE-1 to FE-88 and a term of the action the tag
IN-1 to IN-16; the numbering is that of the inventory of UM-4 and is the
same in every UM-4 document. The equations that UM-4 shares with UM-3
(the matter fields on the fold, the neutrinos, the cosmological
perturbations, the thermal history) keep the numbers they have in UM-3;
the equations of gravity, of the background and of dirac16complex have
the numbers of the UM-3 equations on the same subjects, and FE-73 to
FE-88 are new. Displays without a tag are steps of a derivation.

\hypertarget{coordinates-signature-and-units}{%
\subsection{Coordinates, signature and
units}\label{coordinates-signature-and-units}}

The author's coordinates are \(X = \{x_0, x_1, \dots, x_7\}\), always
with the hidden direction first: \(x_0\) is the hidden space-like
coordinate, on which the metric depends periodically with the period
\(\pi/(6H)\) through the angle \(z = 6Hx_0\); \(x_1, x_2, x_3\) are
ordinary space, \(x_4\) is the time, and \(x_5, x_6, x_7\) are three
extra time-like coordinates. The signature is \((+,+,+,+,-,-,-,-)\) in
this order; the four observable coordinates \(x_1, \dots, x_4\) have the
signature \((+,+,+,-)\). \(H\) is the author's constant (an inverse
length); this document keeps it as a free parameter, bounded from below
by the propagation of the gauge fields in the bulk (FE-82). Units are
natural: the reduced Planck constant \(\hbar\) and the speed of light
\(c\) are 1; in eight dimensions the field dirac16complex has the
dimension GeV\(^{7/2}\) and the gravitational coupling
\(\kappa = 8\pi G_8\) the dimension GeV\(^{-6}\). A prime on a function
of the time \(x_4\) alone is its derivative \(\mathrm{d}/\mathrm{d}x_4\)
(in the cosmological perturbation equations, and only there, a dot is
the derivative with respect to the conformal time). The proper
coordinate of the hidden direction is \(\sigma\) (FD-3, FE-28):
\(\sigma = 0\) is the fold \(z = \pi/2\).

\hypertarget{the-fields-of-um-4-and-where-they-live}{%
\subsection{The fields of UM-4 and where they
live}\label{the-fields-of-um-4-and-where-they-live}}

The author orders a model universe with the following coupled fields
(his order of 2026-10-24, items 1 to 10, and his order of 2026-10-25
13:15 for the photons and gravitons); UM-4, as UM-3, adds the
right-handed quark singlets, the Yukawa couplings and three right-handed
neutrinos (accepted by the author on 2026-10-24 at 17:32), the fold
fluid and the induced gravity of the fold (ASSUMED and ESTIMATE,
explained below).

\begin{longtable}[]{@{}
  >{\raggedright\arraybackslash}p{(\columnwidth - 10\tabcolsep) * \real{0.1495}}
  >{\raggedright\arraybackslash}p{(\columnwidth - 10\tabcolsep) * \real{0.1308}}
  >{\raggedright\arraybackslash}p{(\columnwidth - 10\tabcolsep) * \real{0.2056}}
  >{\raggedright\arraybackslash}p{(\columnwidth - 10\tabcolsep) * \real{0.1028}}
  >{\raggedright\arraybackslash}p{(\columnwidth - 10\tabcolsep) * \real{0.2056}}
  >{\raggedright\arraybackslash}p{(\columnwidth - 10\tabcolsep) * \real{0.2056}}@{}}
\toprule\noalign{}
\begin{minipage}[b]{\linewidth}\raggedright
field
\end{minipage} & \begin{minipage}[b]{\linewidth}\raggedright
symbol
\end{minipage} & \begin{minipage}[b]{\linewidth}\raggedright
what it is
\end{minipage} & \begin{minipage}[b]{\linewidth}\raggedright
lives in
\end{minipage} & \begin{minipage}[b]{\linewidth}\raggedright
it couples to
\end{minipage} & \begin{minipage}[b]{\linewidth}\raggedright
field equations
\end{minipage} \\
\midrule\noalign{}
\endhead
\bottomrule\noalign{}
\endlastfoot
Einstein gravity & \(g_{MN}\) & the metric of eight dimensions & bulk
and fold & every field & FE-1 to FE-11, FE-27 to FE-35, FE-54, FE-59,
FE-63, FE-66, FE-68 to FE-72, FE-76, FE-79 to FE-81, FE-87 \\
dirac16complex & \(\Psi\) & the author's 16-component complex spinor
field; its left- and right-handed halves \(\Psi_{\mathrm{L}}\),
\(\Psi_{\mathrm{R}}\) & the bulk & gravity (spin connection), itself,
every field of the fold at the fold (IN-16) & FE-6, FE-12 to FE-14,
FE-36 to FE-43, FE-47, FE-55, FE-61, FE-85, FE-86 \\
gluons & \(G^A\) & the gauge field of \(SU(3)_C\) & the bulk; its
massless mode is bound to the fold & gravity, quarks,
\(\Psi_{\mathrm{L}}\), \(\Psi_{\mathrm{R}}\) (at the fold) & FE-8,
FE-16, FE-74, FE-79 \\
gauge fields of \(SU(2)_L\) and \(U(1)_Y\) & \(W^1\), \(W^2\), \(W^3\),
\(B\) & the electroweak gauge fields & the bulk; their massless modes
are bound to the fold & gravity, \(\Phi\), the fermions,
\(\Psi_{\mathrm{L}}\), \(\Psi_{\mathrm{R}}\) (at the fold) & FE-8,
FE-17, FE-18, FE-45, FE-74, FE-75, FE-82 \\
Higgs doublet & \(\Phi\) & the scalar doublet with
\(V = \mu^2\Phi^\dagger\Phi + \lambda(\Phi^\dagger\Phi)^2\), expanded
about the ad hoc vacuum \(\Phi = (\mathcal{A}, v + h)/\sqrt2\) & the
fold & gravity, \(W\), \(B\), the fermions (Yukawa),
\(\Psi_{\mathrm{L}}\) only & FE-7, FE-15, FE-26, FE-45, FE-49 \\
three left-handed lepton doublets & \(\ell_i\) &
\((\nu_{\mathrm{L}i}, e_{\mathrm{L}i})\), \(i = 1, 2, 3\) & the fold &
gravity, \(W\), \(B\), \(\Phi\), \(\Psi_{\mathrm{L}}\) and
\(\Psi_{\mathrm{R}}\) & FE-9, FE-22, FE-48 to FE-53 \\
right-handed electron, muon, tau & \(e_{\mathrm{R}i}\) & the
charged-lepton singlets & the fold & gravity, \(B\), \(\Phi\),
\(\Psi_{\mathrm{L}}\), \(\Psi_{\mathrm{R}}\) & FE-9, FE-23 \\
three right-handed neutrinos & \(\nu_{\mathrm{R}i}\) & gauge singlets;
Dirac neutrino masses & the fold & gravity, \(\Phi\),
\(\Psi_{\mathrm{L}}\), \(\Psi_{\mathrm{R}}\) & FE-9, FE-24, FE-62 \\
three quark generations & \(Q_i\), \(u_{\mathrm{R}i}\),
\(d_{\mathrm{R}i}\) & the quark doublets and singlets & the fold &
gravity, gluons, \(W\), \(B\), \(\Phi\), \(\Psi_{\mathrm{L}}\),
\(\Psi_{\mathrm{R}}\) & FE-9, FE-19 to FE-21, FE-46 \\
fold fluid & \(\mathcal{T}_{\mathrm{ff}}\) & the tension of the fold and
the pressure along the extra times that holds their volume fixed: the
part of the thin source at the fold not supplied by the other fields &
the fold & gravity & FE-4, FE-5, FE-30, FE-59, FE-65, FE-72 \\
induced gravity of the fold & \(\kappa_{\mathrm{ind}}\) & an
Einstein-Hilbert term of the fold induced by the quantum fluctuations of
its fields & the fold & gravity & FE-5, FE-59, FE-63, FE-64 \\
\end{longtable}

\textbf{Why the gauge fields live in the bulk (AUTHOR, by the order).}
The order requires that photons propagate in the bulk at the speed of
light; the photon is the massless mode of the gauge fields of
\(SU(2)_L \times U(1)_Y\), so UM-4 places all gauge fields of the
Standard Model in the bulk. In the metric of UM-4 every field of the
bulk has a massless mode bound to the fold that travels at the speed of
light at every depth (FE-73 to FE-78), so the four-dimensional photon
and gluons are these modes, and a photon that leaves the fold travels
through the bulk at the speed of light. The price is a continuum of
massive modes of every gauge field, which modify Coulomb's law at short
distances, add contact interactions between the currents of the fold and
give the W and Z bosons a width into the bulk (FE-80 to FE-82):
observation then requires \(H > 2534\) GeV (DERIVED, ESTIMATE). The
alternative, the gauge fields on the fold (as in UM-3), would keep the
Standard Model on the fold without these corrections, but its photon
could not propagate in the bulk at all: a gauge field confined to the
fold has no wave along \(\sigma\), so it does not meet the order that
photons propagate in the bulk at the speed of light, and UM-4 does not
take it.

\textbf{Why the Higgs doublet, the quarks and the leptons live on the
fold (ASSUMED).} In UM-4 the light cone of every slice of the bulk is
that of the fold, so the argument of UM-3 for placing the Standard Model
on the fold (fields with different profiles in the bulk would have
different limiting speeds) no longer applies. The fermions could be
fields of the bulk with massless modes bound to the fold (a kink mass
gives each a normalisable zero mode, FE-77), but then their Yukawa
couplings to the Higgs doublet would be set by the overlaps of their
profiles, and the flavour structure of three generations by the profiles
of the eighteen multiplets; UM-4 keeps the Higgs doublet, the quarks and
the leptons on the fold, where their couplings and their flavour physics
are those of the Standard Model with three Dirac neutrinos (the simplest
choice that reproduces the measured flavour physics). dirac16complex is
the author's field of eight dimensions and lives in the bulk; it couples
to the fields of the fold at the fold. Status: ASSUMED (the place of the
Higgs doublet and the fermions), AUTHOR (the place of gravity and of the
photon, by the order), DERIVED (the speeds of the fields of the bulk,
\texttt{um4\_bulk.py}, B1 to B8).

\textbf{Why the fold carries a source (DERIVED).} The warp factor
\(\mathrm{e}^{-2H\lvert\sigma\rvert}\) has a kink at the fold. The
Einstein equations, which are linear in the second derivatives of the
metric along \(\sigma\), then contain a delta function at
\(\sigma = 0\): the metric is a solution only if the fold carries a thin
source \(\mathcal{S}^a{}_b\) (FE-4). For the author's form of the metric
\(\kappa\mathcal{S}^a{}_b = -12H\delta^a{}_b\) in every direction of the
fold, a pure tension; on an expanding fold the tension plus matter that
obeys the Einstein equations of seven dimensions to the first order in
the curvature of the fold (FE-30, FE-71; DERIVED,
\texttt{um4\_geometry.py}, G7; \texttt{um4\_fold.py}, F2). This source
is the fold fluid together with the fields of the fold, the particles
\(\chi\) of dirac16complex and the portals.

\hypertarget{conventions}{%
\subsection{Conventions}\label{conventions}}

\begin{itemize}
\tightlist
\item
  Indices. \(M, N, P, Q\) (with numbered copies \(M_1, M_2, \dots\)):
  coordinate indices of eight dimensions; \(a, b, c, d\): frame indices
  of eight dimensions (the frame direction \((0)\) is the \(\sigma\)
  direction); \(\mu, \nu, \rho\): coordinate indices of the four
  observable directions \(x_1, \dots, x_4\); \(A, B, C = 1, \dots, 8\):
  colour (adjoint) indices; \(I, J, K = 1, 2, 3\): weak-isospin
  (adjoint) indices; \(i, j = 1, 2, 3\): generations (and the neutrino
  mass states); \(\alpha, \beta = e, \mu, \tau\): lepton flavours;
  \(p = 1, 2, 3\): the Lovelock order; \(l = 0, 1, 2, \dots\): the
  multipole index of the cosmological perturbations;
  \(X = \mathrm{L}, \mathrm{R}\): the two halves of dirac16complex;
  \(f \in \{Q, u, d, \ell, e, \nu\}\): the six kinds of Standard-Model
  multiplets. A repeated index is summed over its range.
\item
  Two kinds of Dirac matrices. The author's real \(16 \times 16\)
  matrices \(T_a\) (IN-5) act on dirac16complex; the four-dimensional
  complex \(4 \times 4\) matrices \(\hat\gamma^\mu\) (IN-14) act on the
  Standard-Model fermions. They are never multiplied with each other;
  the couplings between dirac16complex and the Standard Model are
  products of Lorentz tensors.
\item
  Normalisation of the fold. No field depends on the extra times
  \(x_5, x_6, x_7\) (ASSUMED: a dependence on them gives tachyonic
  four-dimensional modes, FE-38). Integrated over them, the extra times
  give their physical volume on the fold, \(V_*\); the fields of the
  fold are normalised as four-dimensional fields, the couplings are
  their measured four-dimensional values, and the fold action carries
  the factor \(1/V_*\) (IN-8). The gauge fields of the bulk are
  normalised in eight dimensions (IN-7); their zero modes have the
  measured four-dimensional couplings (FE-44).
\item
  A word or a particle name in a subscript is upright
  (\(m_{\mathrm{W}}\), \(\nu_{\mathrm{R}}\), \(\Psi_{\mathrm{L}}\)); an
  italic letter in a subscript is an index. \(\mathrm{e}\) is Euler's
  number, \(\mathrm{i}\) the imaginary unit, \(\mathrm{d}\) the
  differential; \(\delta\) is the Kronecker delta and
  \(\delta_{\mathrm{D}}\) the Dirac delta function.
\end{itemize}

\hypertarget{the-action}{%
\section{The action}\label{the-action}}

All field equations of UM-4 follow from one action by varying it with
respect to every field. This section writes out every term of the action
(IN-1 to IN-16) and says in plain words what it couples. The
Standard-Model part has exactly the gauge group, the fermion content
(three generations, each with a right-handed neutrino), the Higgs
potential and the gauge couplings of the Standard Model; its gauge
fields live in the bulk, its Higgs doublet and its fermions on the fold.
dirac16complex is a singlet of the Standard-Model gauge group; its
left-handed and right-handed halves couple to every Standard-Model field
at the fold through gauge-invariant portals with their own coupling
constants, and only its left-handed half couples to the Higgs doublet.

\hypertarget{the-sum-of-all-terms}{%
\subsection{The sum of all terms}\label{the-sum-of-all-terms}}

\[
\begin{aligned}
S_{\mathrm{tot}} &= S_{\mathrm{EL}} + S_\Psi + S_{\mathrm{gauge}} + S_{\mathrm{fold}},\\
S_{\mathrm{fold}} &= \int\mathrm{d}^8X\,\sqrt{\lvert g\rvert}\;\delta_{\mathrm{D}}(\sigma)\,\frac{1}{V_*}\,
\big(\mathcal{L}_{\mathrm{IG}} + \mathcal{L}_{\mathrm{ff}} + \mathcal{L}_{\mathrm{H}} + \mathcal{L}_{\mathrm{f}}
+ \mathcal{L}_{\mathrm{Y}} + \mathcal{L}_{\mathrm{P}}\big)
\end{aligned} \tag{IN-1}
\] \textbf{Symbols in (IN-1).}

\begin{itemize}
\tightlist
\item
  \(S_{\mathrm{tot}}\): The action of UM-4 (the integral of the
  Lagrangian density over spacetime); its variation with respect to
  every field gives the field equations. Unit: dimensionless (in units
  in which the reduced Planck constant is 1). Status: ASSUMED.
\item
  \(S_{\mathrm{EL}}\): The Einstein-Lovelock action of gravity in eight
  dimensions; in UM-4 the Einstein-Hilbert action with the cosmological
  constant. Unit: dimensionless. Status: AUTHOR.
\item
  \(S_\Psi\): The action of dirac16complex in the bulk (kinetic term
  with the canonical spin connection, mass term, contact
  self-interaction). Unit: dimensionless. Status: RECORD.
\item
  \(S_{\mathrm{gauge}}\): The action of the gauge fields in the bulk.
  Unit: dimensionless. Status: STANDARD.
\item
  \(S_{\mathrm{fold}}\): The action of the fold: the Standard Model, the
  portals of dirac16complex, the fold fluid and the Einstein-Hilbert
  term that the fields of the fold induce, all at \(\sigma = 0\). Unit:
  dimensionless. Status: ASSUMED.
\item
  \(\mathrm{d}\): The differential: \(\mathrm{d}\) in front of a
  quantity is an infinitesimal change of it (as \(\mathrm{d}\sigma\), a
  small step of \(\sigma\)); the quotient of two such changes is a
  derivative; \(\mathrm{d}\) with an exponent, in front of several
  variables, is their volume element, the product of their
  differentials. Status: STANDARD.
\item
  \(X = \{x_0, \dots, x_7\}\): The author's coordinates of the
  eight-dimensional spacetime, \(X = \{x_0, x_1, \dots, x_7\}\) with the
  hidden direction first (Dirac\_Einstein-Lovelock\_4+4\_Nash.nb,
  section DEFCHART); \(X^M\) is the coordinate with the index \(M\) and
  \(\mathrm{d}^8X = \mathrm{d}x_0\,\mathrm{d}x_1\cdots\mathrm{d}x_7\)
  the coordinate volume element. Unit: length (each coordinate). Status:
  AUTHOR.
\item
  \(g_{MN}\): The metric tensor of the eight-dimensional spacetime
  (signature \(+,+,+,+,-,-,-,-\) in the order \(x_0, \dots, x_7\)):
  \(g_{MN}\,\mathrm{d}X^M\mathrm{d}X^N\) is the squared interval,
  \(g^{MN}\) its inverse, \(\lvert g\rvert\) the absolute value of its
  determinant (\(\sqrt{\lvert g\rvert}\,\mathrm{d}^8X\) is the proper
  volume element). Unit: dimensionless. Status: AUTHOR.
\item
  \(\delta_{\mathrm{D}}\): The Dirac delta function of \(\sigma\): zero
  for \(\sigma \ne 0\), with the integral 1 over any range that contains
  \(\sigma = 0\); the integral of \(\delta_{\mathrm{D}}(\sigma)\) times
  a smooth function of \(\sigma\) is the value of that function at
  \(\sigma = 0\). Unit: 1/length (for \(\delta_{\mathrm{D}}(\sigma)\)).
  Status: STANDARD.
\item
  \(\sigma\): The proper coordinate of the hidden direction:
  \(\mathrm{d}\sigma = \lvert\cot z\rvert\,\mathrm{d}x_0\),
  \(\sigma = \mathrm{sgn}(\pi/2 - z)\ln(\sin z)/(6H)\) on one period
  \(0 < z < \pi\); \(\sigma = 0\) is the fold \(z = \pi/2\),
  \(\sigma \to \mp\infty\) the tips \(z \to 0, \pi\); one period of
  \(x_0\) is the whole line \(-\infty < \sigma < \infty\). Unit: length.
  Status: DERIVED.
\item
  \(V_*\): The physical volume of the extra times \(x_5, x_6, x_7\) on
  the fold in the late universe,
  \(V_* = V_{\mathrm{t}}(\mathrm{e}^{-a_4}b_4)^3\) (with
  \(V_{\mathrm{t}} = \int\mathrm{d}x_5\,\mathrm{d}x_6\,\mathrm{d}x_7\)
  the coordinate volume of the extra times) at a time when
  \(\mathrm{e}^{-a_4}b_4\) is constant; it converts the
  eight-dimensional normalisation into the four-dimensional one. Unit:
  length\(^3\) (GeV\(^{-3}\)). Status: ASSUMED.
\item
  \(\mathcal{L}_{\mathrm{IG}}\): The Lagrangian of the Einstein-Hilbert
  term of the four-dimensional metric of the fold that the quantum
  fluctuations of the fields of the fold induce; it adds to the
  Einstein-Hilbert term that the graviton of the bulk gives the fold and
  does not change the speed of gravitational waves. Unit: GeV\(^4\).
  Status: ESTIMATE.
\item
  \(\mathcal{L}_{\mathrm{ff}}\): The Lagrangian of the fold fluid (only
  its energy-momentum tensor \(\mathcal{T}_{\mathrm{ff}}\) enters the
  equations). Unit: GeV\(^4\). Status: ASSUMED.
\item
  \(\mathcal{L}_{\mathrm{H}}\): The Lagrangian of the Higgs doublet on
  the fold. Unit: GeV\(^4\). Status: AUTHOR.
\item
  \(\mathcal{L}_{\mathrm{f}}\): The kinetic Lagrangian of the quarks and
  leptons on the fold. Unit: GeV\(^4\). Status: STANDARD.
\item
  \(\mathcal{L}_{\mathrm{Y}}\): The Yukawa Lagrangian of the quarks and
  leptons with the Higgs doublet. Unit: GeV\(^4\). Status: STANDARD.
\item
  \(\mathcal{L}_{\mathrm{P}}\): The Lagrangian of the portals of
  dirac16complex at the fold. Unit: GeV\(^4\). Status: INVENTED.
\end{itemize}

The action has four parts: gravity in eight dimensions
(\(S_{\mathrm{EL}}\), IN-2, IN-3), the field dirac16complex in the bulk
(\(S_\Psi\), IN-4 to IN-6), the gauge fields in the bulk
(\(S_{\mathrm{gauge}}\), IN-7) and the fold (\(S_{\mathrm{fold}}\), IN-8
to IN-16). The Dirac delta \(\delta_{\mathrm{D}}(\sigma)\) of the proper
hidden coordinate \(\sigma\) confines the terms of the fold to
\(\sigma = 0\); the factor \(1/V_*\) turns the integral over the extra
times into the normalisation of four-dimensional fields (IN-8). On the
fold live: the gravity that its fields induce
(\(\mathcal{L}_{\mathrm{IG}}\)), the fold fluid
(\(\mathcal{L}_{\mathrm{ff}}\)), the Higgs doublet
(\(\mathcal{L}_{\mathrm{H}}\)), the quarks and leptons
(\(\mathcal{L}_{\mathrm{f}}\)), their Yukawa couplings
(\(\mathcal{L}_{\mathrm{Y}}\)) and the portals of dirac16complex
(\(\mathcal{L}_{\mathrm{P}}\)); the gauge fields enter the fold through
their values at \(\sigma = 0\) in the covariant derivatives of the
fields of the fold (IN-13). Varying \(S_{\mathrm{tot}}\) with respect to
the metric gives the gravitational equations (FE-1, FE-4); with respect
to \(\bar\Psi\) the equation of dirac16complex (FE-12, FE-13); with
respect to the gauge fields the Yang-Mills equations of the bulk with
the currents of the fold as their source (FE-16 to FE-18); with respect
to \(\Phi^\dagger\) the Higgs equation (FE-15); with respect to every
fermion its Dirac equation (FE-19 to FE-24). Status: ASSUMED (the choice
and the place of the terms; the content of each term is stated with its
own label).

\hypertarget{gravity}{%
\subsection{Gravity}\label{gravity}}

\[
\begin{aligned}
S_{\mathrm{EL}} &= \frac{1}{2\kappa}\int\mathrm{d}^8X\,\sqrt{\lvert g\rvert}\,\Big(\sum_{p=1}^{3}\alpha_{(p)}\,\mathcal{L}_{(p)}
 - 2\Lambda\Big) = \frac{1}{2\kappa}\int\mathrm{d}^8X\,\sqrt{\lvert g\rvert}\,\big(R - 2\Lambda\big),\qquad \kappa = 8\pi G_8,\\
\alpha_{(p)} &= w_{(p)}\,H^{2-2p},\qquad w_{(1)} = 1,\quad w_{(2)} = w_{(3)} = 0,\qquad \Lambda = -21H^2
\end{aligned} \tag{IN-2}
\] \textbf{Symbols in (IN-2).}

\begin{itemize}
\tightlist
\item
  \(S_{\mathrm{EL}}\): The Einstein-Lovelock action of gravity in eight
  dimensions; in UM-4 the Einstein-Hilbert action with the cosmological
  constant. Unit: dimensionless. Status: AUTHOR.
\item
  \(\kappa\): The eight-dimensional gravitational coupling,
  \(\kappa = 8\pi G_8\). Unit: GeV\(^{-6}\). Status: AUTHOR.
\item
  \(\mathrm{d}\): The differential: \(\mathrm{d}\) in front of a
  quantity is an infinitesimal change of it (as \(\mathrm{d}\sigma\), a
  small step of \(\sigma\)); the quotient of two such changes is a
  derivative; \(\mathrm{d}\) with an exponent, in front of several
  variables, is their volume element, the product of their
  differentials. Status: STANDARD.
\item
  \(X = \{x_0, \dots, x_7\}\): The author's coordinates of the
  eight-dimensional spacetime, \(X = \{x_0, x_1, \dots, x_7\}\) with the
  hidden direction first (Dirac\_Einstein-Lovelock\_4+4\_Nash.nb,
  section DEFCHART); \(X^M\) is the coordinate with the index \(M\) and
  \(\mathrm{d}^8X = \mathrm{d}x_0\,\mathrm{d}x_1\cdots\mathrm{d}x_7\)
  the coordinate volume element. Unit: length (each coordinate). Status:
  AUTHOR.
\item
  \(g_{MN}\): The metric tensor of the eight-dimensional spacetime
  (signature \(+,+,+,+,-,-,-,-\) in the order \(x_0, \dots, x_7\)):
  \(g_{MN}\,\mathrm{d}X^M\mathrm{d}X^N\) is the squared interval,
  \(g^{MN}\) its inverse, \(\lvert g\rvert\) the absolute value of its
  determinant (\(\sqrt{\lvert g\rvert}\,\mathrm{d}^8X\) is the proper
  volume element). Unit: dimensionless. Status: AUTHOR.
\item
  \(p\): The Lovelock order \(p = 1, 2, 3\) (Einstein, Gauss-Bonnet,
  cubic), as an index and as a number (UM-4 keeps only \(p = 1\)); in
  the Kohn-Sham equations also a second label of a Kohn-Sham orbital,
  with the Kronecker delta of the two labels. Status: STANDARD.
\item
  \(\alpha_{(p)}\): The Lovelock coupling of order \(p\),
  \(\alpha_{(p)} = w_{(p)}H^{2-2p}\): in UM-4 \(\alpha_{(1)} = 1\) and
  \(\alpha_{(2)} = \alpha_{(3)} = 0\) (Einstein gravity). Unit:
  GeV\(^{2-2p}\). Status: AUTHOR.
\item
  \(\mathcal{L}_{(p)}\): The Lovelock density of order \(p\),
  \(\mathcal{L}_{(p)} = 2^{-p}\,\delta^{P_1\cdots P_{2p}}_{Q_1\cdots Q_{2p}}\,R^{Q_1Q_2}{}_{P_1P_2}\cdots R^{Q_{2p-1}Q_{2p}}{}_{P_{2p-1}P_{2p}}\),
  a scalar built from \(p\) curvature tensors: \(\mathcal{L}_{(1)} = R\)
  (Einstein-Hilbert), \(\mathcal{L}_{(2)}\) the Gauss-Bonnet density,
  \(\mathcal{L}_{(3)}\) the cubic Lovelock density. Unit:
  1/length\(^{2p}\). Status: STANDARD.
\item
  \(\Lambda\): The eight-dimensional cosmological constant of UM-4,
  \(\Lambda = -21H^2\): the value for which the static block is a vacuum
  of Einstein gravity (its Einstein tensor is \(21H^2\) in every
  direction). The author writes the pure number \(\Lambda_{\mathrm{N}}\)
  in
  \(0 = -\Lambda_{\mathrm{N}} + \sum_p H^{-2p}w_{(p)}\mathcal{E}_{(p)}\),
  so that \(\Lambda = -\Lambda_{\mathrm{N}}H^2\). Unit: GeV\(^2\).
  Value: \(-21H^2\); the author's original \(-10H^2\); UM-3: \(-7H^2\).
  Status: AUTHOR (the term), DERIVED (the value).
\item
  \(R^{MN}{}_{PQ},\ R_{MN},\ R\): The curvature of the metric
  \(g_{MN}\): with four indices the Riemann tensor \(R^M{}_{NPQ}\) (in
  the convention in which the commutator of two covariant derivatives,
  acting on a vector, is the Riemann tensor, with its last two indices
  those of the two derivatives, contracted with the vector), with two
  indices the Ricci tensor \(R_{MN} = R^P{}_{MPN}\), without indices the
  scalar curvature \(R = g^{MN}R_{MN}\); with frame indices
  \(R^{(a)(b)}{}_{(c)(d)}\) its frame components,
  \(R^{(a)(b)}{}_{(a)(b)}\) (no sum) the sectional curvature of the
  frame plane \((a, b)\). Unit: 1/length\(^2\). Status: STANDARD.
\item
  \(\pi\): The number pi, the ratio of the circumference of a circle to
  its diameter. Unit: dimensionless. Value: 3.141592654. Status:
  STANDARD.
\item
  \(G_8\): The eight-dimensional Newton constant
  (\(\kappa = 8\pi G_8\)). Unit: GeV\(^{-6}\). Status: AUTHOR.
\item
  \(w_{(p)}\): The Lovelock weights of the author's Einstein-Lovelock
  action, \(w_{(1)} = 1\) and \(w_{(2)}\), \(w_{(3)}\) of the
  Gauss-Bonnet and the cubic densities. UM-4 sets
  \(w_{(2)} = w_{(3)} = 0\) (Einstein gravity): the only weights for
  which gravitons propagate on the light cone on every background, not
  only around the static block. The author's original weights
  \(w_{(2)} = 1/32\), \(w_{(3)} = 1/1152\)
  (Dirac\_Einstein-Lovelock\_4+4\_Nash.nb, In{[}134{]}) keep the
  gravitons on the light cone only around the static block (with
  \(\Lambda = -161H^2/16\)); UM-3 used \((1, 1/20, 1/360)\). Unit:
  dimensionless. Value: \((1, 0, 0)\) in UM-4; the author's original
  \((1, 1/32, 1/1152)\). Status: AUTHOR (the form), DERIVED (the values
  that the order of 2026-10-25 selects).
\item
  \(H\): The author's constant (an inverse length) that fixes the period
  \(\pi/(6H)\) of \(x_0\), the warp factor
  \(\mathrm{e}^{-2H\lvert\sigma\rvert}\) of the metric, the cosmological
  constant \(\Lambda = -21H^2\) and the tension
  \(\kappa\lambda_{\mathrm{F}} = 12H\) of the fold; in UM-4 a free
  parameter, bounded from below by the propagation of the gauge fields
  in the bulk. Unit: GeV (1/length). Value: \(H > 2534\) GeV, about 2.5
  TeV (from the invisible width of the Z boson). Status: AUTHOR.
\end{itemize}

Gravity is the author's Einstein-Lovelock theory in eight dimensions
with the weights of the Gauss-Bonnet density \(\mathcal{L}_{(2)}\) and
of the cubic Lovelock density \(\mathcal{L}_{(3)}\) set to zero:
Einstein gravity, whose density \(\mathcal{L}_{(1)} = R\) is the scalar
curvature, with the cosmological constant \(\Lambda = -21H^2\) that
makes the static block a vacuum (FE-70). The author writes the vacuum
equations as
\(0 = -\Lambda_{\mathrm{N}} + H^{-2}w_{(1)}\mathcal{E}_{(1)} + H^{-4}w_{(2)}\mathcal{E}_{(2)} + H^{-6}w_{(3)}\mathcal{E}_{(3)}\)
(\texttt{Dirac\_Einstein-Lovelock\_4+4\_Nash.nb}, In{[}134{]});
multiplied by \(H^2\) this is FE-1 without sources and with
\(\Lambda = -\Lambda_{\mathrm{N}}H^2\), here
\(\Lambda_{\mathrm{N}} = 21\). His original values were
\(w_{(2)} = 1/32\), \(w_{(3)} = 1/1152\) and
\(\Lambda_{\mathrm{N}} = 10\); UM-3 used \(w_{(2)} = 1/20\),
\(w_{(3)} = 1/360\), \(\Lambda_{\mathrm{N}} = 7\). UM-4 uses Einstein
gravity because it is the only choice of the weights for which gravitons
propagate at the speed of light in the bulk on every background (FE-69),
as the order of 2026-10-25 13:15 requires. Gravity couples to every
field through the metric that appears in every other term of the action
(in \(\sqrt{\lvert g\rvert}\), in the contractions of indices and in the
spin connection). Status: AUTHOR (the action), DERIVED (the weights and
\(\Lambda\) that the order selects, FE-68 to FE-70).

\[
\begin{aligned}
\mathcal{L}_{(p)} &= \frac{1}{2^{p}}\,\delta^{M_1N_1\cdots M_pN_p}_{P_1Q_1\cdots P_pQ_p}\,
 R^{P_1Q_1}{}_{M_1N_1}\cdots R^{P_pQ_p}{}_{M_pN_p},\\
\mathcal{L}_{(1)} &= R,\qquad \mathcal{L}_{(2)} = R^2 - 4R_{MN}R^{MN} + R_{MNPQ}R^{MNPQ}
\end{aligned} \tag{IN-3}
\] \textbf{Symbols in (IN-3).}

\begin{itemize}
\tightlist
\item
  \(\mathcal{L}_{(p)}\): The Lovelock density of order \(p\),
  \(\mathcal{L}_{(p)} = 2^{-p}\,\delta^{P_1\cdots P_{2p}}_{Q_1\cdots Q_{2p}}\,R^{Q_1Q_2}{}_{P_1P_2}\cdots R^{Q_{2p-1}Q_{2p}}{}_{P_{2p-1}P_{2p}}\),
  a scalar built from \(p\) curvature tensors: \(\mathcal{L}_{(1)} = R\)
  (Einstein-Hilbert), \(\mathcal{L}_{(2)}\) the Gauss-Bonnet density,
  \(\mathcal{L}_{(3)}\) the cubic Lovelock density. Unit:
  1/length\(^{2p}\). Status: STANDARD.
\item
  \(p\): The Lovelock order \(p = 1, 2, 3\) (Einstein, Gauss-Bonnet,
  cubic), as an index and as a number (UM-4 keeps only \(p = 1\)); in
  the Kohn-Sham equations also a second label of a Kohn-Sham orbital,
  with the Kronecker delta of the two labels. Status: STANDARD.
\item
  \(\delta\): The Kronecker delta: with two indices \(\delta^M_N = 1\)
  if \(M = N\) and 0 otherwise (likewise \(\delta_{X\mathrm{L}}\) is 1
  if the label \(X\) of a half of dirac16complex is L and 0 otherwise);
  with two index strings of equal length the generalised Kronecker delta
  \(\delta^{M_1M_2\cdots}_{N_1N_2\cdots}\), the determinant of the
  square matrix of the Kronecker deltas \(\delta^{M_a}_{N_b}\) (\(+1\)
  or \(-1\) if the upper indices are an even or odd permutation of the
  lower ones, 0 otherwise). Status: STANDARD.
\item
  \(M\): A coordinate index of eight dimensions (also with a numbered
  copy \(M_1, M_2, \dots\)): it runs over \(x_0, x_1, \dots, x_7\); a
  repeated index is summed. Status: STANDARD.
\item
  \(N\): A coordinate index of eight dimensions (also
  \(N_1, N_2, \dots\)), running over \(x_0, \dots, x_7\); a repeated
  index is summed. Status: STANDARD.
\item
  \(P\): A coordinate index of eight dimensions (also
  \(P_1, P_2, \dots\)), running over \(x_0, \dots, x_7\); a repeated
  index is summed. Status: STANDARD.
\item
  \(Q\): A coordinate index of eight dimensions (also
  \(Q_1, Q_2, \dots\)), running over \(x_0, \dots, x_7\); a repeated
  index is summed. Status: STANDARD.
\item
  \(R^{MN}{}_{PQ},\ R_{MN},\ R\): The curvature of the metric
  \(g_{MN}\): with four indices the Riemann tensor \(R^M{}_{NPQ}\) (in
  the convention in which the commutator of two covariant derivatives,
  acting on a vector, is the Riemann tensor, with its last two indices
  those of the two derivatives, contracted with the vector), with two
  indices the Ricci tensor \(R_{MN} = R^P{}_{MPN}\), without indices the
  scalar curvature \(R = g^{MN}R_{MN}\); with frame indices
  \(R^{(a)(b)}{}_{(c)(d)}\) its frame components,
  \(R^{(a)(b)}{}_{(a)(b)}\) (no sum) the sectional curvature of the
  frame plane \((a, b)\). Unit: 1/length\(^2\). Status: STANDARD.
\end{itemize}

The Lovelock densities are built with the generalised Kronecker delta
from \(p\) Riemann tensors. Their key property (Lovelock's theorem,
STANDARD) is that their metric variation contains no derivative of the
metric higher than the second, so that the field equations are
second-order differential equations, linear in the second derivatives
along any one direction (quasi-linearity): a thin source at the fold
produces a finite jump of first derivatives, the junction condition
FE-4. UM-4 uses only \(\mathcal{L}_{(1)}\); the others appear in FE-68
to FE-70, which show why. Status: STANDARD.

\hypertarget{why-einstein-gravity}{%
\subsection{Why Einstein gravity}\label{why-einstein-gravity}}

\[
\begin{aligned}
&\sum_{p=1}^{3}\alpha_{(p)}\,\mathcal{E}_{(p)}{}^a{}_b + \Lambda\,\delta^a{}_b = \mathsf{P}(\mathsf{y})\,H^2\,\delta^a{}_b,\qquad
 \mathsf{P}(\mathsf{y}) = 21\mathsf{y} - 420w_{(2)}\mathsf{y}^2 + 2520w_{(3)}\mathsf{y}^3 + \frac{\Lambda}{H^2},\\
&\text{Einstein } (w_{(2)} = w_{(3)} = 0,\ \Lambda = -21H^2):\quad \mathsf{P} = 21(\mathsf{y} - 1);\\
&\text{the author's original weights } \Big(\frac{1}{32}, \frac{1}{1152}\Big),\ \Lambda = -\frac{161}{16}H^2:\quad
 \mathsf{P} = \frac{7}{16}(\mathsf{y} - 1)\big(5\mathsf{y}^2 - 25\mathsf{y} + 23\big);\\
&\text{UM-3's weights } \Big(\frac{1}{20}, \frac{1}{360}\Big),\ \Lambda = -7H^2:\quad \mathsf{P} = 7(\mathsf{y} - 1)^3
\end{aligned} \tag{FE-70}
\] \textbf{Symbols in (FE-70).}

\begin{itemize}
\tightlist
\item
  \(p\): The Lovelock order \(p = 1, 2, 3\) (Einstein, Gauss-Bonnet,
  cubic), as an index and as a number (UM-4 keeps only \(p = 1\)); in
  the Kohn-Sham equations also a second label of a Kohn-Sham orbital,
  with the Kronecker delta of the two labels. Status: STANDARD.
\item
  \(\alpha_{(p)}\): The Lovelock coupling of order \(p\),
  \(\alpha_{(p)} = w_{(p)}H^{2-2p}\): in UM-4 \(\alpha_{(1)} = 1\) and
  \(\alpha_{(2)} = \alpha_{(3)} = 0\) (Einstein gravity). Unit:
  GeV\(^{2-2p}\). Status: AUTHOR.
\item
  \(\mathcal{E}_{(p)}{}^M{}_N\): The Lovelock tensor of order \(p\),
  \(\mathcal{E}_{(p)}{}^M{}_N = -2^{-(p+1)}\,\delta^{MP_1\cdots P_{2p}}_{NQ_1\cdots Q_{2p}}\,R^{Q_1Q_2}{}_{P_1P_2}\cdots R^{Q_{2p-1}Q_{2p}}{}_{P_{2p-1}P_{2p}}\)
  (with both indices down it is \(1/\sqrt{\lvert g\rvert}\) times the
  variation of \(\sqrt{\lvert g\rvert}\,\mathcal{L}_{(p)}\) with respect
  to \(g^{MN}\); \(\mathcal{E}_{(1)}\) the Einstein tensor,
  \(\mathcal{E}_{(2)}\) the Gauss-Bonnet tensor, \(\mathcal{E}_{(3)}\)
  the cubic Lovelock tensor); with frame indices its frame components.
  Unit: 1/length\(^{2p}\). Status: STANDARD.
\item
  \(a\): A frame (orthonormal) index running over \(0, 1, \dots, 7\)
  (frame direction \((0)\) along \(\sigma\), \((1), (2), (3)\) along
  ordinary space, \((4)\) along the time, \((5), (6), (7)\) along the
  extra times); raised and lowered with the frame metric
  \(\eta_{ab} = \mathrm{diag}(1, 1, 1, 1, -1, -1, -1, -1)\); on the fold
  (\(a, b \in \{1, \dots, 7\}\)) the frame of the fold; in the tensors
  of the fold written with its coordinates \(x^a\) (its metric
  \(\hat g_{ab}\) and the tensors built from it) a coordinate index of
  the fold, \(a \in \{1, \dots, 7\}\); a repeated index is summed.
  Status: STANDARD.
\item
  \(b\): A frame index (as \(a\)). Status: STANDARD.
\item
  \(\Lambda\): The eight-dimensional cosmological constant of UM-4,
  \(\Lambda = -21H^2\): the value for which the static block is a vacuum
  of Einstein gravity (its Einstein tensor is \(21H^2\) in every
  direction). The author writes the pure number \(\Lambda_{\mathrm{N}}\)
  in
  \(0 = -\Lambda_{\mathrm{N}} + \sum_p H^{-2p}w_{(p)}\mathcal{E}_{(p)}\),
  so that \(\Lambda = -\Lambda_{\mathrm{N}}H^2\). Unit: GeV\(^2\).
  Value: \(-21H^2\); the author's original \(-10H^2\); UM-3: \(-7H^2\).
  Status: AUTHOR (the term), DERIVED (the value).
\item
  \(\delta\): The Kronecker delta: with two indices \(\delta^M_N = 1\)
  if \(M = N\) and 0 otherwise (likewise \(\delta_{X\mathrm{L}}\) is 1
  if the label \(X\) of a half of dirac16complex is L and 0 otherwise);
  with two index strings of equal length the generalised Kronecker delta
  \(\delta^{M_1M_2\cdots}_{N_1N_2\cdots}\), the determinant of the
  square matrix of the Kronecker deltas \(\delta^{M_a}_{N_b}\) (\(+1\)
  or \(-1\) if the upper indices are an even or odd permutation of the
  lower ones, 0 otherwise). Status: STANDARD.
\item
  \(\mathsf{P}(\mathsf{y})\): The vacuum polynomial of the static blocks
  of curvature \(-\mathsf{y}H^2\): the left-hand side of the bulk
  equations divided by \(H^2\), the same in every direction,
  \(\mathsf{P}(\mathsf{y}) = 21\mathsf{y} - 420w_{(2)}\mathsf{y}^2 + 2520w_{(3)}\mathsf{y}^3 + \Lambda/H^2\);
  \(\mathsf{P}'(1)\) its derivative at \(\mathsf{y} = 1\), which sets
  the strength of the graviton around the static block relative to
  Einstein gravity (\(\mathsf{P}'(1)/21\)). Unit: dimensionless. Value:
  \(21(\mathsf{y} - 1)\) for Einstein gravity with \(\Lambda = -21H^2\).
  Status: DERIVED.
\item
  \(\mathsf{y}\): The curvature of a static block of constant curvature
  \(-\mathsf{y}H^2\) in all eight directions, in units of \(-H^2\) (its
  warp is \(\mathrm{e}^{-2\sqrt{\mathsf{y}}H\lvert\sigma\rvert}\) in all
  seven directions of the fold): \(\mathsf{y} = 1\) is the static block
  of UM-4. Unit: dimensionless. Status: DERIVED.
\item
  \(H\): The author's constant (an inverse length) that fixes the period
  \(\pi/(6H)\) of \(x_0\), the warp factor
  \(\mathrm{e}^{-2H\lvert\sigma\rvert}\) of the metric, the cosmological
  constant \(\Lambda = -21H^2\) and the tension
  \(\kappa\lambda_{\mathrm{F}} = 12H\) of the fold; in UM-4 a free
  parameter, bounded from below by the propagation of the gauge fields
  in the bulk. Unit: GeV (1/length). Value: \(H > 2534\) GeV, about 2.5
  TeV (from the invisible width of the Z boson). Status: AUTHOR.
\item
  \(w_{(p)}\): The Lovelock weights of the author's Einstein-Lovelock
  action, \(w_{(1)} = 1\) and \(w_{(2)}\), \(w_{(3)}\) of the
  Gauss-Bonnet and the cubic densities. UM-4 sets
  \(w_{(2)} = w_{(3)} = 0\) (Einstein gravity): the only weights for
  which gravitons propagate on the light cone on every background, not
  only around the static block. The author's original weights
  \(w_{(2)} = 1/32\), \(w_{(3)} = 1/1152\)
  (Dirac\_Einstein-Lovelock\_4+4\_Nash.nb, In{[}134{]}) keep the
  gravitons on the light cone only around the static block (with
  \(\Lambda = -161H^2/16\)); UM-3 used \((1, 1/20, 1/360)\). Unit:
  dimensionless. Value: \((1, 0, 0)\) in UM-4; the author's original
  \((1, 1/32, 1/1152)\). Status: AUTHOR (the form), DERIVED (the values
  that the order of 2026-10-25 selects).
\end{itemize}

The vacua of the static family: on the static blocks of curvature
\(-\mathsf{y}H^2\) in all eight directions (the warp
\(\mathrm{e}^{-2\sqrt{\mathsf{y}}H\lvert\sigma\rvert}\) in all seven
directions of the fold, FE-32 with \(H \to \sqrt{\mathsf{y}}H\)) the
left-hand side of the bulk equations is the same in every direction,
\(\mathsf{P}(\mathsf{y})H^2\); in a space of constant curvature \(-K\)
in \(n = 8\) dimensions the Lovelock tensor of order \(p\) is
\(\mathcal{E}_{(p)}{}^a{}_b = -\frac12(-K)^p\,\frac{(n - 1)!}{(n - 1 - 2p)!}\,\delta^a{}_b\),
which gives \(21K\), \(-420K^2\) and \(2520K^3\). The static block of
UM-4 (\(\mathsf{y} = 1\)) is a vacuum if \(\mathsf{P}(1) = 0\), and a
non-degenerate one if \(\mathsf{P}'(1) \ne 0\),
\(\mathsf{P}'(1) = 21 - 840w_{(2)} + 7560w_{(3)}\). For Einstein gravity
\(\Lambda = -21H^2\) and \(\mathsf{y} = 1\) is the only root, with
\(\mathsf{P}'(1) = 21\). With the author's original weights the
cosmological constant would have to be \(-161H^2/16\) (not his
\(-10H^2\)) and the polynomial has the roots \(\mathsf{y} = 1\),
\(1.2155\) and \(3.7845\), with \(\mathsf{P}'(1) = 21/16\). UM-3's
weights make \(\mathsf{y} = 1\) a triple root: a degenerate vacuum
without linear gravitons. Status: DERIVED (\texttt{um4\_geometry.py},
G4; \texttt{um4\_gravity.py}, L1).

\[
\frac{\mathsf{P}'(1)}{42}\Big[\mathrm{e}^{2H\lvert\sigma\rvert}\Big(\frac{\partial^2h_{\mathrm{gw}}}{\partial x_4^2}
 - \frac{\partial^2h_{\mathrm{gw}}}{\partial x_3^2}\Big) - \frac{\partial^2h_{\mathrm{gw}}}{\partial\sigma^2}
 + 7H\,\mathrm{sgn}(\sigma)\,\frac{\partial h_{\mathrm{gw}}}{\partial\sigma}\Big] = 0,\qquad
\mathsf{P}'(1) = 21 - 840w_{(2)} + 7560w_{(3)} \tag{FE-68}
\] \textbf{Symbols in (FE-68).}

\begin{itemize}
\tightlist
\item
  \(\mathsf{P}(\mathsf{y})\): The vacuum polynomial of the static blocks
  of curvature \(-\mathsf{y}H^2\): the left-hand side of the bulk
  equations divided by \(H^2\), the same in every direction,
  \(\mathsf{P}(\mathsf{y}) = 21\mathsf{y} - 420w_{(2)}\mathsf{y}^2 + 2520w_{(3)}\mathsf{y}^3 + \Lambda/H^2\);
  \(\mathsf{P}'(1)\) its derivative at \(\mathsf{y} = 1\), which sets
  the strength of the graviton around the static block relative to
  Einstein gravity (\(\mathsf{P}'(1)/21\)). Unit: dimensionless. Value:
  \(21(\mathsf{y} - 1)\) for Einstein gravity with \(\Lambda = -21H^2\).
  Status: DERIVED.
\item
  \({}'\): The derivative of a function of the time \(x_4\) alone with
  respect to \(x_4\) (\(a_4' = \mathrm{d}a_4/\mathrm{d}x_4\); two primes
  the second derivative); on a function of \(\sigma\) alone (a profile)
  the derivative with respect to \(\sigma\). Status: STANDARD.
\item
  \(\mathrm{e}\): Euler's number, the base of the natural logarithm;
  \(\mathrm{e}^{X}\) is the exponential of \(X\). Unit: dimensionless.
  Value: 2.718281828. Status: STANDARD.
\item
  \(H\): The author's constant (an inverse length) that fixes the period
  \(\pi/(6H)\) of \(x_0\), the warp factor
  \(\mathrm{e}^{-2H\lvert\sigma\rvert}\) of the metric, the cosmological
  constant \(\Lambda = -21H^2\) and the tension
  \(\kappa\lambda_{\mathrm{F}} = 12H\) of the fold; in UM-4 a free
  parameter, bounded from below by the propagation of the gauge fields
  in the bulk. Unit: GeV (1/length). Value: \(H > 2534\) GeV, about 2.5
  TeV (from the invisible width of the Z boson). Status: AUTHOR.
\item
  \(\sigma\): The proper coordinate of the hidden direction:
  \(\mathrm{d}\sigma = \lvert\cot z\rvert\,\mathrm{d}x_0\),
  \(\sigma = \mathrm{sgn}(\pi/2 - z)\ln(\sin z)/(6H)\) on one period
  \(0 < z < \pi\); \(\sigma = 0\) is the fold \(z = \pi/2\),
  \(\sigma \to \mp\infty\) the tips \(z \to 0, \pi\); one period of
  \(x_0\) is the whole line \(-\infty < \sigma < \infty\). Unit: length.
  Status: DERIVED.
\item
  \(\partial\): A partial derivative: the derivative with respect to one
  variable at fixed values of the others (\(\partial_\sigma\) with
  respect to \(\sigma\); \(\partial_M\) and \(\partial_\mu\) with
  respect to the coordinate that carries the same index, of the eight
  dimensions or of the fold). Status: STANDARD.
\item
  \(h_{\mathrm{gw}}\): The amplitude of a gravitational wave of the +
  polarisation travelling along \(x_3\): the metric components along
  \(x_1\) and \(x_2\) multiplied by \(\mathrm{e}^{h_{\mathrm{gw}}}\) and
  \(\mathrm{e}^{-h_{\mathrm{gw}}}\); a function of \(\sigma\), \(x_3\)
  and \(x_4\) in the bulk, and on the fold the gravitational wave of the
  four-dimensional world, to the first order its plus polarisation
  \(h_{+}\). Unit: dimensionless. Status: STANDARD.
\item
  \(x_4\): The time; a prime on a function of \(x_4\) alone is
  \(\mathrm{d}/\mathrm{d}x_4\). Unit: length (= time, \(c = 1\)).
  Status: AUTHOR.
\item
  \(x_3\): The third coordinate of ordinary 3-space. Unit: length.
  Status: AUTHOR.
\item
  \(\mathrm{sgn}\): The sign function: \(+1\) for a positive argument
  and \(-1\) for a negative one. Status: STANDARD.
\item
  \(w_{(p)}\): The Lovelock weights of the author's Einstein-Lovelock
  action, \(w_{(1)} = 1\) and \(w_{(2)}\), \(w_{(3)}\) of the
  Gauss-Bonnet and the cubic densities. UM-4 sets
  \(w_{(2)} = w_{(3)} = 0\) (Einstein gravity): the only weights for
  which gravitons propagate on the light cone on every background, not
  only around the static block. The author's original weights
  \(w_{(2)} = 1/32\), \(w_{(3)} = 1/1152\)
  (Dirac\_Einstein-Lovelock\_4+4\_Nash.nb, In{[}134{]}) keep the
  gravitons on the light cone only around the static block (with
  \(\Lambda = -161H^2/16\)); UM-3 used \((1, 1/20, 1/360)\). Unit:
  dimensionless. Value: \((1, 0, 0)\) in UM-4; the author's original
  \((1, 1/32, 1/1152)\). Status: AUTHOR (the form), DERIVED (the values
  that the order of 2026-10-25 selects).
\end{itemize}

The graviton around the static block of UM-4, for any weights: a
transverse traceless wave of the + polarisation travelling along \(x_3\)
multiplies the metric components along \(x_1\) and \(x_2\) by
\(\mathrm{e}^{h_{\mathrm{gw}}}\) and \(\mathrm{e}^{-h_{\mathrm{gw}}}\)
(\(h_{\mathrm{gw}}\) a function of \(\sigma\), \(x_3\) and \(x_4\));
FE-68 is the bulk equation (FE-1) to the first order in
\(h_{\mathrm{gw}}\), half the difference of its components along \(x_1\)
and \(x_2\), away from the fold. The time and ordinary space enter only
through
\(\mathrm{e}^{2H\lvert\sigma\rvert}(\partial^2/\partial x_4^2 - \partial^2/\partial x_3^2)\):
at every depth the wave has the light cone of the fold, for every choice
of the weights, with the strength \(\mathsf{P}'(1)/21\) relative to
Einstein gravity. The + polarisation along the extra times obeys the
same equation. A degenerate choice (\(\mathsf{P}'(1) = 0\), UM-3's
weights) gives no linear graviton, and \(\mathsf{P}'(1) < 0\) a ghost.
Around the static block, therefore, the order does not single out
Einstein gravity; FE-69 does. Status: DERIVED (\texttt{um4\_gravity.py},
L2).

\[
\begin{aligned}
&\mathrm{d}\sigma^2 + \mathrm{e}^{-2\mathfrak{p}H\sigma}\,\mathrm{d}\vec x^2 - \mathrm{e}^{-2\mathfrak{q}H\sigma}\,\mathrm{d}x_4^2
 - \mathrm{e}^{-2\mathfrak{r}H\sigma}\,\mathrm{d}\vec x_5^2\quad(\sigma > 0):\\
&c_{\mathrm{x}} - c_{\mathrm{t}} = -4\,(\mathfrak{p} - \mathfrak{q})\big(18\mathfrak{r}^2w_{(3)} - w_{(2)}\big)\big(\mathfrak{p} + \mathfrak{q} + 3\mathfrak{r}\big),\\
&c_\sigma - c_{\mathrm{t}} = -4\big[w_{(2)}\big(\mathfrak{p}\mathfrak{q} - \mathfrak{p}^2 + 3\mathfrak{q}\mathfrak{r} - 3\mathfrak{r}^2\big)
 + 18w_{(3)}\mathfrak{r}^2\big(\mathfrak{p}^2 - 3\mathfrak{p}\mathfrak{q} + 2\mathfrak{p}\mathfrak{r} - \mathfrak{q}\mathfrak{r} + \mathfrak{r}^2\big)\big]
\end{aligned} \tag{FE-69}
\] \textbf{Symbols in (FE-69).}

\begin{itemize}
\tightlist
\item
  \(\mathrm{d}\): The differential: \(\mathrm{d}\) in front of a
  quantity is an infinitesimal change of it (as \(\mathrm{d}\sigma\), a
  small step of \(\sigma\)); the quotient of two such changes is a
  derivative; \(\mathrm{d}\) with an exponent, in front of several
  variables, is their volume element, the product of their
  differentials. Status: STANDARD.
\item
  \(\sigma\): The proper coordinate of the hidden direction:
  \(\mathrm{d}\sigma = \lvert\cot z\rvert\,\mathrm{d}x_0\),
  \(\sigma = \mathrm{sgn}(\pi/2 - z)\ln(\sin z)/(6H)\) on one period
  \(0 < z < \pi\); \(\sigma = 0\) is the fold \(z = \pi/2\),
  \(\sigma \to \mp\infty\) the tips \(z \to 0, \pi\); one period of
  \(x_0\) is the whole line \(-\infty < \sigma < \infty\). Unit: length.
  Status: DERIVED.
\item
  \(\mathrm{e}\): Euler's number, the base of the natural logarithm;
  \(\mathrm{e}^{X}\) is the exponential of \(X\). Unit: dimensionless.
  Value: 2.718281828. Status: STANDARD.
\item
  \(\mathfrak{p},\ \mathfrak{q},\ \mathfrak{r}\): The rates of the warps
  of ordinary space (\(\mathfrak{p}\)), of the time (\(\mathfrak{q}\))
  and of the extra times (\(\mathfrak{r}\)) of a static warped
  background, in units of \(H\). Unit: dimensionless. Value:
  \(\mathfrak{p} = \mathfrak{q} = \mathfrak{r} = 1\) for the static
  block of UM-4; \(\mathfrak{p} = \mathfrak{r} = 1\),
  \(\mathfrak{q} = 0\) for that of UM-3. Status: DERIVED.
\item
  \(H\): The author's constant (an inverse length) that fixes the period
  \(\pi/(6H)\) of \(x_0\), the warp factor
  \(\mathrm{e}^{-2H\lvert\sigma\rvert}\) of the metric, the cosmological
  constant \(\Lambda = -21H^2\) and the tension
  \(\kappa\lambda_{\mathrm{F}} = 12H\) of the fold; in UM-4 a free
  parameter, bounded from below by the propagation of the gauge fields
  in the bulk. Unit: GeV (1/length). Value: \(H > 2534\) GeV, about 2.5
  TeV (from the invisible width of the Z boson). Status: AUTHOR.
\item
  \(\vec x\): The position in ordinary space,
  \(\vec x = (x_1, x_2, x_3)\). Unit: length. Status: AUTHOR.
\item
  \(x_4\): The time; a prime on a function of \(x_4\) alone is
  \(\mathrm{d}/\mathrm{d}x_4\). Unit: length (= time, \(c = 1\)).
  Status: AUTHOR.
\item
  \(\vec x_5\): The position along the extra times,
  \(\vec x_5 = (x_5, x_6, x_7)\);
  \(\mathrm{d}\vec x_5^2 = \mathrm{d}x_5^2 + \mathrm{d}x_6^2 + \mathrm{d}x_7^2\).
  Unit: length. Status: AUTHOR.
\item
  \(c_{\mathrm{x}}\): The coefficient of the second derivative along the
  proper length of \(x_3\) in the same equation; the wave moves along
  \(x_3\) with the local speed \(\sqrt{c_{\mathrm{x}}/c_{\mathrm{t}}}\)
  in units of the local speed of light. Unit: dimensionless. Value: 1
  for Einstein gravity on every background;
  \(1 - 24w_{(2)} + 72w_{(3)}\) at the static block of UM-3
  (\(c_{\mathrm{x}}/c_{\mathrm{t}} = 5\) with the author's original
  weights). Status: DERIVED.
\item
  \(c_{\mathrm{t}}\): The coefficient of the second derivative along the
  proper time in the linearised equation of a gravitational wave of the
  + polarisation on a static warped background, normalised so that
  Einstein gravity gives 1. Unit: dimensionless. Value: 1 for Einstein
  gravity on every background; \(1 - 40w_{(2)} + 360w_{(3)}\) at the
  static block of UM-3. Status: DERIVED.
\item
  \(c_\sigma\): The coefficient of the second derivative along
  \(\sigma\) in the linearised equation of a gravitational wave on a
  static warped background, normalised so that Einstein gravity gives 1;
  the wave moves along \(\sigma\) with the local speed
  \(\sqrt{c_\sigma/c_{\mathrm{t}}}\) in units of the local speed of
  light. Unit: dimensionless. Value: 1 for Einstein gravity on every
  background. Status: DERIVED.
\item
  \(w_{(p)}\): The Lovelock weights of the author's Einstein-Lovelock
  action, \(w_{(1)} = 1\) and \(w_{(2)}\), \(w_{(3)}\) of the
  Gauss-Bonnet and the cubic densities. UM-4 sets
  \(w_{(2)} = w_{(3)} = 0\) (Einstein gravity): the only weights for
  which gravitons propagate on the light cone on every background, not
  only around the static block. The author's original weights
  \(w_{(2)} = 1/32\), \(w_{(3)} = 1/1152\)
  (Dirac\_Einstein-Lovelock\_4+4\_Nash.nb, In{[}134{]}) keep the
  gravitons on the light cone only around the static block (with
  \(\Lambda = -161H^2/16\)); UM-3 used \((1, 1/20, 1/360)\). Unit:
  dimensionless. Value: \((1, 0, 0)\) in UM-4; the author's original
  \((1, 1/32, 1/1152)\). Status: AUTHOR (the form), DERIVED (the values
  that the order of 2026-10-25 selects).
\end{itemize}

The gravitons on backgrounds that are not of constant curvature: the
static warped metrics of the first line (\(\mathfrak{p}\),
\(\mathfrak{q}\), \(\mathfrak{r}\) constant;
\(\mathfrak{p} = \mathfrak{q} = \mathfrak{r} = 1\) is the static block
of UM-4, \(\mathfrak{p} = \mathfrak{r} = 1\), \(\mathfrak{q} = 0\) that
of UM-3), and the + polarisation along \(x_1\), \(x_2\) travelling along
\(x_3\). Its linearised equation contains the second derivatives in the
frame with the coefficients \(c_{\mathrm{t}}\) (of the second derivative
along the proper time), \(c_{\mathrm{x}}\) (along the proper length of
\(x_3\)) and \(c_\sigma\) (along \(\sigma\)), normalised so that
Einstein gravity gives 1; the wave moves along \(x_3\) with the local
speed \(\sqrt{c_{\mathrm{x}}/c_{\mathrm{t}}}\) and along \(\sigma\) with
\(\sqrt{c_\sigma/c_{\mathrm{t}}}\) in units of the local speed of light.
Einstein gravity gives
\(c_{\mathrm{t}} = c_{\mathrm{x}} = c_\sigma = 1\) on every such
background. The Gauss-Bonnet and cubic terms give the deviations of
FE-69, which vanish for all \(\mathfrak{p}\), \(\mathfrak{q}\),
\(\mathfrak{r}\) only for \(w_{(2)} = w_{(3)} = 0\). A curved fold (an
expanding universe, a star, a gravitational wave of finite amplitude)
makes the warps of the different directions differ, so with the
Gauss-Bonnet or cubic terms gravitons would leave the light cone there.
At the static block of UM-3 the same coefficients are
\(c_{\mathrm{t}} = 1 - 40w_{(2)} + 360w_{(3)}\) and
\(c_{\mathrm{x}} = 1 - 24w_{(2)} + 72w_{(3)}\), with the author's
original weights \(c_{\mathrm{x}}/c_{\mathrm{t}} = 5\): the bulk
gravitons of UM-3 that were \(\sqrt5\) times faster than light. Status:
DERIVED (\texttt{um4\_gravity.py}, L3 and L4).

The gravity of UM-4 is therefore Einstein gravity with
\(\Lambda = -21H^2\): the static block is its unique vacuum
(\(\mathsf{P}'(1) = 21 > 0\)), gravitons propagate on the light cone on
every background, and the fold is a pure tension
\(\kappa\lambda_{\mathrm{F}} = 12H\) (FE-30). The author's original
weights would also make the static block a vacuum (with
\(\Lambda = -161H^2/16\)) and give gravitons at the speed of light
around it with \(1/16\) of the Einstein strength, but on every
background that is not of constant curvature their gravitons would leave
the light cone; UM-4 does not use them. Status: DERIVED
(\texttt{um4\_gravity.py}, L4).

\hypertarget{the-field-dirac16complex}{%
\subsection{The field dirac16complex}\label{the-field-dirac16complex}}

\[
S_\Psi = \int\mathrm{d}^8X\,\sqrt{\lvert g\rvert}\,\Big[\frac12\big(\bar\Psi\gamma^MD_M\Psi - (D_M\bar\Psi)\gamma^M\Psi\big)
 - m_\Psi\,\bar\Psi\Psi - \frac{\lambda_\Psi}{2}\big(\bar\Psi\Psi\big)^2\Big] \tag{IN-4}
\] \textbf{Symbols in (IN-4).}

\begin{itemize}
\tightlist
\item
  \(S_\Psi\): The action of dirac16complex in the bulk (kinetic term
  with the canonical spin connection, mass term, contact
  self-interaction). Unit: dimensionless. Status: RECORD.
\item
  \(\mathrm{d}\): The differential: \(\mathrm{d}\) in front of a
  quantity is an infinitesimal change of it (as \(\mathrm{d}\sigma\), a
  small step of \(\sigma\)); the quotient of two such changes is a
  derivative; \(\mathrm{d}\) with an exponent, in front of several
  variables, is their volume element, the product of their
  differentials. Status: STANDARD.
\item
  \(X = \{x_0, \dots, x_7\}\): The author's coordinates of the
  eight-dimensional spacetime, \(X = \{x_0, x_1, \dots, x_7\}\) with the
  hidden direction first (Dirac\_Einstein-Lovelock\_4+4\_Nash.nb,
  section DEFCHART); \(X^M\) is the coordinate with the index \(M\) and
  \(\mathrm{d}^8X = \mathrm{d}x_0\,\mathrm{d}x_1\cdots\mathrm{d}x_7\)
  the coordinate volume element. Unit: length (each coordinate). Status:
  AUTHOR.
\item
  \(g_{MN}\): The metric tensor of the eight-dimensional spacetime
  (signature \(+,+,+,+,-,-,-,-\) in the order \(x_0, \dots, x_7\)):
  \(g_{MN}\,\mathrm{d}X^M\mathrm{d}X^N\) is the squared interval,
  \(g^{MN}\) its inverse, \(\lvert g\rvert\) the absolute value of its
  determinant (\(\sqrt{\lvert g\rvert}\,\mathrm{d}^8X\) is the proper
  volume element). Unit: dimensionless. Status: AUTHOR.
\item
  \(\Psi,\ \bar\Psi\): dirac16complex: the author's 16-component complex
  spinor field of eight dimensions;
  \(\bar\Psi = \Psi^\dagger\sigma_{16}\) is its Dirac adjoint; in a term
  of the fold, \(\Psi\) is its value at \(\sigma = 0\). Unit:
  GeV\(^{7/2}\). Status: AUTHOR.
\item
  \(\gamma^M,\ \gamma^{(a)}\): The Dirac matrices of dirac16complex
  along the coordinates, \(\gamma^M = \mathbf{e}_a{}^MT^a\),
  \(\gamma_M = g_{MN}\gamma^N\); \(\gamma^{(a)} = T^a\) along the frame
  direction \((a)\) (\(\gamma^{(0)}\) along \(\sigma\), \(\gamma^{(4)}\)
  along the time). Unit: dimensionless. Status: AUTHOR.
\item
  \(M\): A coordinate index of eight dimensions (also with a numbered
  copy \(M_1, M_2, \dots\)): it runs over \(x_0, x_1, \dots, x_7\); a
  repeated index is summed. Status: STANDARD.
\item
  \(D_M,\ D_\mu\): The covariant derivative: on dirac16complex
  \(D_M = \partial_M + \Omega_M\) (with
  \(D_M\bar\Psi = \partial_M\bar\Psi - \bar\Psi\Omega_M\)); on the
  fields of the fold the gauge and spin covariant derivative \(D_\mu\)
  (the partial derivative plus the spin connection of the fold and the
  gauge fields times their couplings and generators);
  \(D^\mu = \hat h^{\mu\nu}D_\nu\); on an adjoint field
  \((D_\mu X)^A = \partial_\mu X^A + g_3f^{ABC}G^B_\mu X^C\) (likewise
  for \(SU(2)\) with \(g_2\epsilon^{IJK}\)). Unit: 1/length. Status:
  STANDARD.
\item
  \(m_\Psi\): The mass parameter of dirac16complex in the bulk. Unit:
  GeV. Status: RECORD.
\item
  \(\lambda_\Psi\): The coupling of the contact self-interaction of
  dirac16complex. Unit: GeV\(^{-6}\). Status: RECORD.
\end{itemize}

dirac16complex \(\Psi\) is the author's complex spinor field of eight
dimensions with 16 components. Its Lagrangian is that of the record: the
kinetic term with the canonical spin connection (in the symmetric form,
which is real without a factor \(\mathrm{i}\) for the author's real
matrices, IN-5), the mass term and a contact self-interaction. It
couples to gravity through \(\sqrt{\lvert g\rvert}\), \(\gamma^M\) and
the spin connection \(\Omega_M\) inside \(D_M\) (IN-6), and to itself
through \(\lambda_\Psi\).
\(\bar\Psi\Psi = \bar\Psi_{\mathrm{L}}\Psi_{\mathrm{L}} + \bar\Psi_{\mathrm{R}}\Psi_{\mathrm{R}}\)
(the mass term does not mix the halves, IN-5), while the kinetic term
mixes them (it contains one Dirac matrix). dirac16complex is a singlet
of the gauge group: it has no gauge coupling in the bulk. Status:
RECORD.

\[
\begin{aligned}
&\{T_a, T_b\} = 2\eta_{ab}\mathbf{1},\qquad \eta = \mathrm{diag}(1, 1, 1, 1, -1, -1, -1, -1),\qquad T^a = \eta^{ab}T_b,
 \qquad \gamma^M = \mathbf{e}_a{}^M\,T^a,\\
&\bar\Psi = \Psi^\dagger\sigma_{16},\qquad \sigma_{16} = T_0T_1T_2T_3,\qquad \Gamma_8 = T_0T_1T_2T_3T_4T_5T_6T_7,\\
&P_{\mathrm{L}} = \frac{\mathbf{1} - \Gamma_8}{2},\qquad P_{\mathrm{R}} = \frac{\mathbf{1} + \Gamma_8}{2},\qquad
 \Psi_{\mathrm{L}} = P_{\mathrm{L}}\Psi,\qquad \Psi_{\mathrm{R}} = P_{\mathrm{R}}\Psi
\end{aligned} \tag{IN-5}
\] \textbf{Symbols in (IN-5).}

\begin{itemize}
\tightlist
\item
  \(T_a,\ T^a\): The author's eight real \(16 \times 16\) Dirac matrices
  T16\^{}A{[}a{]} of 4+4.nb (\(a = 0, \dots, 7\); \(T_0, \dots, T_7\)),
  written out in FD-13: \(T_a\) has the block \(\bar\Xi_a\) in its upper
  right and \(\Xi_a\) in its lower left corner;
  \(\{T_a, T_b\} = 2\eta_{ab}\mathbf{1}_{16}\); \(T^a = \eta^{ab}T_b\).
  Unit: dimensionless. Status: AUTHOR.
\item
  \(a\): A frame (orthonormal) index running over \(0, 1, \dots, 7\)
  (frame direction \((0)\) along \(\sigma\), \((1), (2), (3)\) along
  ordinary space, \((4)\) along the time, \((5), (6), (7)\) along the
  extra times); raised and lowered with the frame metric
  \(\eta_{ab} = \mathrm{diag}(1, 1, 1, 1, -1, -1, -1, -1)\); on the fold
  (\(a, b \in \{1, \dots, 7\}\)) the frame of the fold; in the tensors
  of the fold written with its coordinates \(x^a\) (its metric
  \(\hat g_{ab}\) and the tensors built from it) a coordinate index of
  the fold, \(a \in \{1, \dots, 7\}\); a repeated index is summed.
  Status: STANDARD.
\item
  \(b\): A frame index (as \(a\)). Status: STANDARD.
\item
  \(\eta_{ab}\): The flat metric of the frame,
  \(\eta = \mathrm{diag}(1, 1, 1, 1, -1, -1, -1, -1)\) (eta4488 of
  4+4.nb); \(\eta^{ab}\) is its inverse (the same matrix). Unit:
  dimensionless. Status: AUTHOR.
\item
  \(\mathbf{1},\ \mathbf{1}_4,\ \mathbf{1}_8,\ \mathbf{1}_{16}\): The
  identity matrix of the size fixed by the context (\(16 \times 16\) on
  dirac16complex, \(4 \times 4\) on a four-dimensional spinor,
  \(3 \times 3\) in generation or flavour space); \(\mathbf{1}_4\),
  \(\mathbf{1}_8\) and \(\mathbf{1}_{16}\) are the identity matrices of
  4, 8 and 16 rows. Unit: dimensionless. Status: STANDARD.
\item
  \(\mathrm{diag}\): A diagonal matrix with the listed entries on its
  diagonal (in the order of the coordinates \(x_0, x_1, \dots, x_7\) for
  a tensor of eight dimensions, \(x_1, \dots, x_7\) for a tensor of the
  fold, \(x_1, \dots, x_4\) for a four-dimensional tensor). Status:
  STANDARD.
\item
  \(\gamma^M,\ \gamma^{(a)}\): The Dirac matrices of dirac16complex
  along the coordinates, \(\gamma^M = \mathbf{e}_a{}^MT^a\),
  \(\gamma_M = g_{MN}\gamma^N\); \(\gamma^{(a)} = T^a\) along the frame
  direction \((a)\) (\(\gamma^{(0)}\) along \(\sigma\), \(\gamma^{(4)}\)
  along the time). Unit: dimensionless. Status: AUTHOR.
\item
  \(M\): A coordinate index of eight dimensions (also with a numbered
  copy \(M_1, M_2, \dots\)): it runs over \(x_0, x_1, \dots, x_7\); a
  repeated index is summed. Status: STANDARD.
\item
  \(\mathbf{e}_a{}^M,\ \mathbf{e}^a{}_M\): The vielbein (frame field) of
  the metric: \(g_{MN} = \eta_{ab}\,\mathbf{e}^a{}_M\mathbf{e}^b{}_N\);
  \(\mathbf{e}_a{}^M\) is its inverse;
  \(\mathbf{e}_{bN} = \eta_{bc}\mathbf{e}^c{}_N\); \(\mathbf{e}^{(a)}\)
  is the frame one-form \(\mathbf{e}^a{}_M\,\mathrm{d}X^M\). Unit:
  dimensionless. Status: STANDARD.
\item
  \(\Psi,\ \bar\Psi\): dirac16complex: the author's 16-component complex
  spinor field of eight dimensions;
  \(\bar\Psi = \Psi^\dagger\sigma_{16}\) is its Dirac adjoint; in a term
  of the fold, \(\Psi\) is its value at \(\sigma = 0\). Unit:
  GeV\(^{7/2}\). Status: AUTHOR.
\item
  \({}^\dagger\): Hermitian conjugation: the complex conjugate of the
  transpose of a matrix or spinor; for an operator on the space of
  states, its adjoint with respect to the scalar product of that space
  (for the quantised dirac16complex the indefinite Krein scalar
  product). Status: STANDARD.
\item
  \(\sigma_{16}\): The matrix of the Dirac adjoint,
  \(\sigma_{16} = T_0T_1T_2T_3\), real and symmetric; \(\sigma_{16}T_a\)
  is real and antisymmetric. Unit: dimensionless. Status: AUTHOR.
\item
  \(\Gamma_8\): The chirality matrix of eight dimensions,
  \(\Gamma_8 = T_0T_1\cdots T_7\) (the author's T16\^{}A{[}8{]})
  \(= \mathrm{diag}(-\mathbf{1}_8, \mathbf{1}_8)\);
  \(\Gamma_8^2 = \mathbf{1}\); it anticommutes with every \(T_a\). Unit:
  dimensionless. Status: AUTHOR.
\item
  \(P_{\mathrm{L}}\): The author's left-handed chiral projector,
  \(P_{\mathrm{L}} = (\mathbf{1} - \Gamma_8)/2\). Unit: dimensionless.
  Status: AUTHOR.
\item
  \(P_{\mathrm{R}}\): The author's right-handed chiral projector,
  \(P_{\mathrm{R}} = (\mathbf{1} + \Gamma_8)/2\). Unit: dimensionless.
  Status: AUTHOR.
\item
  \(\Psi_{\mathrm{L}}\): The left-handed half of dirac16complex,
  \(\Psi_{\mathrm{L}} = P_{\mathrm{L}}\Psi\) (the upper eight components
  in the basis of 4+4.nb). Unit: GeV\(^{7/2}\). Status: AUTHOR.
\item
  \(\Psi_{\mathrm{R}}\): The right-handed half of dirac16complex,
  \(\Psi_{\mathrm{R}} = P_{\mathrm{R}}\Psi\) (the lower eight
  components). Unit: GeV\(^{7/2}\). Status: AUTHOR.
\end{itemize}

\(T_a\) (\(a = 0, \dots, 7\)) are the author's eight real
\(16 \times 16\) Dirac matrices \texttt{T16\^{}A{[}a{]}} of
\texttt{4+4.nb}, written out entry by entry in FD-13 of the definition
of the fold at the beginning of this document (built from six real
\(4 \times 4\) matrices), and \(\Gamma_8\) is the author's chirality
matrix \texttt{T16\^{}A{[}8{]}}
\(= \mathrm{diag}(-\mathbf{1}_8, \mathbf{1}_8)\); \(P_{\mathrm{L}}\) and
\(P_{\mathrm{R}}\) are the author's chiral projectors (the Clifford
relations, the reality of the matrices and the projectors are checked
with the author's matrices in \texttt{um4\_dirac\_matrices.py}, 55
checks, and the matrices of FD-13 against the notebook in
\texttt{um4\_geometry.py}, G11; the facts used here are stated in this
paragraph). The Dirac adjoint uses \(\sigma_{16} = T_0T_1T_2T_3\):
\(\sigma_{16}\) is real and symmetric, \(\sigma_{16}T_a\) is real and
antisymmetric, and \(\sigma_{16}\) commutes with \(\Gamma_8\) (it is a
product of four Dirac matrices), so
\(\overline{P_X\Psi} = \bar\Psi P_X\) and every bilinear with an even
number of Dirac matrices joins a half with itself, every bilinear with
an odd number joins \(\Psi_{\mathrm{L}}\) with \(\Psi_{\mathrm{R}}\)
(DERIVED, \texttt{um4\_fields.py}, checks F1 and F2). Status: AUTHOR
(matrices, projectors), RECORD (adjoint).

\[
\begin{aligned}
&D_M\Psi = \partial_M\Psi + \Omega_M\Psi,\qquad D_M\bar\Psi = \partial_M\bar\Psi - \bar\Psi\,\Omega_M,\qquad
 \Omega_M = \frac12\,\varpi_{Mab}\,\Sigma^{ab},\qquad \Sigma^{ab} = \frac14\big[T^a, T^b\big],\\
&\varpi_{Mab} = \mathbf{e}_a{}^N\big(\partial_M\mathbf{e}_{bN} - \Gamma^P{}_{MN}\,\mathbf{e}_{bP}\big)
\end{aligned} \tag{IN-6}
\] \textbf{Symbols in (IN-6).}

\begin{itemize}
\tightlist
\item
  \(D_M,\ D_\mu\): The covariant derivative: on dirac16complex
  \(D_M = \partial_M + \Omega_M\) (with
  \(D_M\bar\Psi = \partial_M\bar\Psi - \bar\Psi\Omega_M\)); on the
  fields of the fold the gauge and spin covariant derivative \(D_\mu\)
  (the partial derivative plus the spin connection of the fold and the
  gauge fields times their couplings and generators);
  \(D^\mu = \hat h^{\mu\nu}D_\nu\); on an adjoint field
  \((D_\mu X)^A = \partial_\mu X^A + g_3f^{ABC}G^B_\mu X^C\) (likewise
  for \(SU(2)\) with \(g_2\epsilon^{IJK}\)). Unit: 1/length. Status:
  STANDARD.
\item
  \(M\): A coordinate index of eight dimensions (also with a numbered
  copy \(M_1, M_2, \dots\)): it runs over \(x_0, x_1, \dots, x_7\); a
  repeated index is summed. Status: STANDARD.
\item
  \(\partial\): A partial derivative: the derivative with respect to one
  variable at fixed values of the others (\(\partial_\sigma\) with
  respect to \(\sigma\); \(\partial_M\) and \(\partial_\mu\) with
  respect to the coordinate that carries the same index, of the eight
  dimensions or of the fold). Status: STANDARD.
\item
  \(\Psi,\ \bar\Psi\): dirac16complex: the author's 16-component complex
  spinor field of eight dimensions;
  \(\bar\Psi = \Psi^\dagger\sigma_{16}\) is its Dirac adjoint; in a term
  of the fold, \(\Psi\) is its value at \(\sigma = 0\). Unit:
  GeV\(^{7/2}\). Status: AUTHOR.
\item
  \(\Omega_M\): The spin connection acting on dirac16complex,
  \(\Omega_M = \frac12\varpi_{Mab}\Sigma^{ab}\). Unit: 1/length. Status:
  RECORD.
\item
  \(\varpi_{Mab}\): The canonical (Levi-Civita) spin connection of the
  vielbein, antisymmetric in \(a, b\). Unit: 1/length. Status: STANDARD.
\item
  \(a\): A frame (orthonormal) index running over \(0, 1, \dots, 7\)
  (frame direction \((0)\) along \(\sigma\), \((1), (2), (3)\) along
  ordinary space, \((4)\) along the time, \((5), (6), (7)\) along the
  extra times); raised and lowered with the frame metric
  \(\eta_{ab} = \mathrm{diag}(1, 1, 1, 1, -1, -1, -1, -1)\); on the fold
  (\(a, b \in \{1, \dots, 7\}\)) the frame of the fold; in the tensors
  of the fold written with its coordinates \(x^a\) (its metric
  \(\hat g_{ab}\) and the tensors built from it) a coordinate index of
  the fold, \(a \in \{1, \dots, 7\}\); a repeated index is summed.
  Status: STANDARD.
\item
  \(b\): A frame index (as \(a\)). Status: STANDARD.
\item
  \(\Sigma^{ab},\ \Sigma^{\mu\nu}\): The generators of Spin(4,4) on
  dirac16complex, \(\Sigma^{ab} = \frac14[T^a, T^b]\) (the scaled
  commutators of the author's matrices);
  \(\Sigma^{\mu\nu} = \mathbf{e}_a{}^\mu\mathbf{e}_b{}^\nu\Sigma^{ab}\)
  along the fold. Unit: dimensionless. Status: AUTHOR.
\item
  \(T_a,\ T^a\): The author's eight real \(16 \times 16\) Dirac matrices
  T16\^{}A{[}a{]} of 4+4.nb (\(a = 0, \dots, 7\); \(T_0, \dots, T_7\)),
  written out in FD-13: \(T_a\) has the block \(\bar\Xi_a\) in its upper
  right and \(\Xi_a\) in its lower left corner;
  \(\{T_a, T_b\} = 2\eta_{ab}\mathbf{1}_{16}\); \(T^a = \eta^{ab}T_b\).
  Unit: dimensionless. Status: AUTHOR.
\item
  \(\mathbf{e}_a{}^M,\ \mathbf{e}^a{}_M\): The vielbein (frame field) of
  the metric: \(g_{MN} = \eta_{ab}\,\mathbf{e}^a{}_M\mathbf{e}^b{}_N\);
  \(\mathbf{e}_a{}^M\) is its inverse;
  \(\mathbf{e}_{bN} = \eta_{bc}\mathbf{e}^c{}_N\); \(\mathbf{e}^{(a)}\)
  is the frame one-form \(\mathbf{e}^a{}_M\,\mathrm{d}X^M\). Unit:
  dimensionless. Status: STANDARD.
\item
  \(N\): A coordinate index of eight dimensions (also
  \(N_1, N_2, \dots\)), running over \(x_0, \dots, x_7\); a repeated
  index is summed. Status: STANDARD.
\item
  \(\Gamma^P{}_{MN}\): The Christoffel symbols of the metric \(g_{MN}\)
  (the Levi-Civita connection). Unit: 1/length. Status: STANDARD.
\item
  \(P\): A coordinate index of eight dimensions (also
  \(P_1, P_2, \dots\)), running over \(x_0, \dots, x_7\); a repeated
  index is summed. Status: STANDARD.
\end{itemize}

The canonical spin connection \(\varpi_{Mab}\) of the vielbein
(antisymmetric in \(a, b\)) acts on the spinor through the generators
\(\Sigma^{ab}\) of Spin(4,4). With it the Dirac matrices are covariantly
constant (\(D_M\gamma^N = 0\)), the kinetic term of IN-4 is a scalar and
the field equation FE-12 is covariant. For the metric of UM-4
\(\gamma^M\Omega_M\) is never zero (FE-14). Status: RECORD (the
canonical construction), DERIVED (the vielbein postulate for the metric
of UM-4, \texttt{um4\_geometry.py}, G9).

\hypertarget{the-gauge-fields-in-the-bulk}{%
\subsection{The gauge fields in the
bulk}\label{the-gauge-fields-in-the-bulk}}

\[
\begin{aligned}
S_{\mathrm{gauge}} &= \int\mathrm{d}^8X\,\sqrt{\lvert g\rvert}\;\mathcal{L}_{\mathrm{gauge}},\qquad
\mathcal{L}_{\mathrm{gauge}} = -\frac{1}{4\mathsf{g}_3^2}\,G^A_{MN}G^{AMN} - \frac{1}{4\mathsf{g}_2^2}\,W^I_{MN}W^{IMN}
 - \frac{1}{4\mathsf{g}_1^2}\,B_{MN}B^{MN},\\
G^A_{MN} &= \partial_MG^A_N - \partial_NG^A_M + f^{ABC}G^B_MG^C_N,\qquad
W^I_{MN} = \partial_MW^I_N - \partial_NW^I_M + \epsilon^{IJK}W^J_MW^K_N,\\
B_{MN} &= \partial_MB_N - \partial_NB_M
\end{aligned} \tag{IN-7}
\] \textbf{Symbols in (IN-7).}

\begin{itemize}
\tightlist
\item
  \(S_{\mathrm{gauge}}\): The action of the gauge fields in the bulk.
  Unit: dimensionless. Status: STANDARD.
\item
  \(\mathrm{d}\): The differential: \(\mathrm{d}\) in front of a
  quantity is an infinitesimal change of it (as \(\mathrm{d}\sigma\), a
  small step of \(\sigma\)); the quotient of two such changes is a
  derivative; \(\mathrm{d}\) with an exponent, in front of several
  variables, is their volume element, the product of their
  differentials. Status: STANDARD.
\item
  \(X = \{x_0, \dots, x_7\}\): The author's coordinates of the
  eight-dimensional spacetime, \(X = \{x_0, x_1, \dots, x_7\}\) with the
  hidden direction first (Dirac\_Einstein-Lovelock\_4+4\_Nash.nb,
  section DEFCHART); \(X^M\) is the coordinate with the index \(M\) and
  \(\mathrm{d}^8X = \mathrm{d}x_0\,\mathrm{d}x_1\cdots\mathrm{d}x_7\)
  the coordinate volume element. Unit: length (each coordinate). Status:
  AUTHOR.
\item
  \(g_{MN}\): The metric tensor of the eight-dimensional spacetime
  (signature \(+,+,+,+,-,-,-,-\) in the order \(x_0, \dots, x_7\)):
  \(g_{MN}\,\mathrm{d}X^M\mathrm{d}X^N\) is the squared interval,
  \(g^{MN}\) its inverse, \(\lvert g\rvert\) the absolute value of its
  determinant (\(\sqrt{\lvert g\rvert}\,\mathrm{d}^8X\) is the proper
  volume element). Unit: dimensionless. Status: AUTHOR.
\item
  \(\mathcal{L}_{\mathrm{gauge}}\): The Lagrangian of the gauge fields
  of \(SU(3)_C\times SU(2)_L\times U(1)_Y\) in the bulk. Unit:
  GeV\(^8\). Status: STANDARD.
\item
  \(\mathsf{g}_1,\ \mathsf{g}_2,\ \mathsf{g}_3\): The gauge couplings of
  eight dimensions of \(U(1)_Y\), \(SU(2)_L\) and \(SU(3)_C\): the
  coefficients of the kinetic terms of the gauge fields of the bulk in
  their geometric normalisation; the four-dimensional couplings are
  \(g_k^2 = 3H\mathsf{g}_k^2/(2V_*)\). Unit: GeV\(^{-2}\)
  (\(\mathsf{g}^2\): length\(^4\)). Status: ASSUMED.
\item
  \(G^A_M,\ G^A_\mu,\ G^A_{MN}\): The gluon field of \(SU(3)_C\)
  (\(A = 1, \dots, 8\)): in the bulk \(G^A_M\), a field of eight
  dimensions in its geometric normalisation (the covariant derivative
  \(\partial - \mathrm{i}G\)), and on the fold the canonically
  normalised four-dimensional field \(G^A_\mu(x)\) of its massless mode,
  \(g_3G^A_\mu(x) = G^A_\mu(\sigma = 0, x)\); \(G^A_{MN}\) and
  \(G^A_{\mu\nu}\) their field strengths. Unit: GeV (field), GeV\(^2\)
  (field strength). Status: AUTHOR.
\item
  \(A\): A colour (adjoint SU(3)) index running over \(1, \dots, 8\); a
  repeated index is summed. Status: STANDARD.
\item
  \(M\): A coordinate index of eight dimensions (also with a numbered
  copy \(M_1, M_2, \dots\)): it runs over \(x_0, x_1, \dots, x_7\); a
  repeated index is summed. Status: STANDARD.
\item
  \(N\): A coordinate index of eight dimensions (also
  \(N_1, N_2, \dots\)), running over \(x_0, \dots, x_7\); a repeated
  index is summed. Status: STANDARD.
\item
  \(W^I_M,\ W^I_\mu,\ W^I_{MN}\): The gauge field of \(SU(2)_L\)
  (\(I = 1, 2, 3\)): in the bulk \(W^I_M\) (geometric normalisation), on
  the fold the canonically normalised four-dimensional field
  \(W^I_\mu(x)\) of its massless mode,
  \(g_2W^I_\mu(x) = W^I_\mu(\sigma = 0, x)\); \(W^I_{MN}\),
  \(W^I_{\mu\nu}\) their field strengths;
  \(W^\pm_\mu = (W^1_\mu \mp \mathrm{i}W^2_\mu)/\sqrt2\). Unit: GeV
  (field), GeV\(^2\) (field strength). Status: AUTHOR.
\item
  \(I\): A weak-isospin (adjoint SU(2)) index running over \(1, 2, 3\);
  a repeated index is summed. Status: STANDARD.
\item
  \(B_M,\ B_\mu,\ B_{MN}\): The hypercharge gauge field of \(U(1)_Y\):
  in the bulk \(B_M\) (geometric normalisation), on the fold the
  canonically normalised four-dimensional field \(B_\mu(x)\) of its
  massless mode, \(g_1B_\mu(x) = B_\mu(\sigma = 0, x)\); \(B_{MN}\),
  \(B_{\mu\nu}\) their field strengths. Unit: GeV (field), GeV\(^2\)
  (field strength). Status: AUTHOR.
\item
  \(\partial\): A partial derivative: the derivative with respect to one
  variable at fixed values of the others (\(\partial_\sigma\) with
  respect to \(\sigma\); \(\partial_M\) and \(\partial_\mu\) with
  respect to the coordinate that carries the same index, of the eight
  dimensions or of the fold). Status: STANDARD.
\item
  \(f^{ABC}\): The structure constants of \(SU(3)\),
  \([t^A, t^B] = \mathrm{i}f^{ABC}t^C\). Unit: dimensionless. Status:
  STANDARD.
\item
  \(B\): A colour (adjoint SU(3)) index (as \(A\)). Status: STANDARD.
\item
  \(C\): A colour (adjoint SU(3)) index (as \(A\)). Status: STANDARD.
\item
  \(\epsilon^{IJK}\): The Levi-Civita symbol of three indices (the
  structure constants of \(SU(2)\)), \(\epsilon^{123} = 1\). Unit:
  dimensionless. Status: STANDARD.
\item
  \(J\): A weak-isospin (adjoint SU(2)) index (as \(I\)). Status:
  STANDARD.
\item
  \(K\): A weak-isospin (adjoint SU(2)) index (as \(I\)). Status:
  STANDARD.
\end{itemize}

The gauge fields of \(SU(3)_C \times SU(2)_L \times U(1)_Y\) in the bulk
of eight dimensions, with the couplings \(\mathsf{g}_3\),
\(\mathsf{g}_2\), \(\mathsf{g}_1\) of eight dimensions in front of the
kinetic terms (the fields normalised geometrically, so that the
covariant derivative is \(\partial - \mathrm{i}G\); their dimension is
that of an inverse length, and \(\mathsf{g}^2\) has the dimension of a
length to the fourth power). The order of 2026-10-25 13:15 places the
photon in the bulk; the photon is a combination of \(W^3\) and \(B\), so
the whole electroweak gauge field is a field of the bulk, and UM-4
places the gluons there as well (all gauge fields of the Standard Model
in one place). Their massless modes are bound to the fold with the
profile 1 along \(\sigma\) and have the four-dimensional couplings
\(g_k^2 = 3H\mathsf{g}_k^2/(2V_*)\) (FE-44, FE-74), the measured
couplings of the Standard Model. The gauge fields have no components
along the extra times and no dependence on them (ASSUMED: such
components would be scalars of wrong sign, FE-38), and no component
along \(\sigma\) in the gauge of FE-74 (a component \(B_\sigma\) has no
normalisable zero mode). They couple to the fields of the fold through
their values at \(\sigma = 0\) (IN-13) and to dirac16complex only
through the portals at the fold (IN-16). Status: AUTHOR (the photon in
the bulk, by the order), ASSUMED (all gauge fields in the bulk; no
components along the extra times), STANDARD (the Yang-Mills action).

\hypertarget{the-fold}{%
\subsection{The fold}\label{the-fold}}

\[
\begin{aligned}
&\int\mathrm{d}^8X\,\sqrt{\lvert g\rvert}\;\delta_{\mathrm{D}}(\sigma)\,\frac{1}{V_*}\,\mathbb{F}
 = \int\mathrm{d}^4x\,\sqrt{-\hat h}\;\mathcal{Z}\,\mathbb{F}\big\vert_{\sigma=0},\qquad \mathrm{d}^4x = \mathrm{d}x_1\,\mathrm{d}x_2\,\mathrm{d}x_3\,\mathrm{d}x_4,\\
&\hat h_{\mu\nu} = \mathrm{diag}\big(\mathrm{e}^{2a_4}, \mathrm{e}^{2a_4}, \mathrm{e}^{2a_4}, -1\big),\qquad
 \sqrt{-\hat h} = \mathrm{e}^{3a_4},\qquad
 \mathcal{Z} = \frac{V_{\mathrm{t}}\big(\mathrm{e}^{-a_4}b_4\big)^3}{V_*},\qquad V_{\mathrm{t}} = \int\mathrm{d}x_5\,\mathrm{d}x_6\,\mathrm{d}x_7
\end{aligned} \tag{IN-8}
\] \textbf{Symbols in (IN-8).}

\begin{itemize}
\tightlist
\item
  \(\mathrm{d}\): The differential: \(\mathrm{d}\) in front of a
  quantity is an infinitesimal change of it (as \(\mathrm{d}\sigma\), a
  small step of \(\sigma\)); the quotient of two such changes is a
  derivative; \(\mathrm{d}\) with an exponent, in front of several
  variables, is their volume element, the product of their
  differentials. Status: STANDARD.
\item
  \(X = \{x_0, \dots, x_7\}\): The author's coordinates of the
  eight-dimensional spacetime, \(X = \{x_0, x_1, \dots, x_7\}\) with the
  hidden direction first (Dirac\_Einstein-Lovelock\_4+4\_Nash.nb,
  section DEFCHART); \(X^M\) is the coordinate with the index \(M\) and
  \(\mathrm{d}^8X = \mathrm{d}x_0\,\mathrm{d}x_1\cdots\mathrm{d}x_7\)
  the coordinate volume element. Unit: length (each coordinate). Status:
  AUTHOR.
\item
  \(g_{MN}\): The metric tensor of the eight-dimensional spacetime
  (signature \(+,+,+,+,-,-,-,-\) in the order \(x_0, \dots, x_7\)):
  \(g_{MN}\,\mathrm{d}X^M\mathrm{d}X^N\) is the squared interval,
  \(g^{MN}\) its inverse, \(\lvert g\rvert\) the absolute value of its
  determinant (\(\sqrt{\lvert g\rvert}\,\mathrm{d}^8X\) is the proper
  volume element). Unit: dimensionless. Status: AUTHOR.
\item
  \(\delta_{\mathrm{D}}\): The Dirac delta function of \(\sigma\): zero
  for \(\sigma \ne 0\), with the integral 1 over any range that contains
  \(\sigma = 0\); the integral of \(\delta_{\mathrm{D}}(\sigma)\) times
  a smooth function of \(\sigma\) is the value of that function at
  \(\sigma = 0\). Unit: 1/length (for \(\delta_{\mathrm{D}}(\sigma)\)).
  Status: STANDARD.
\item
  \(\sigma\): The proper coordinate of the hidden direction:
  \(\mathrm{d}\sigma = \lvert\cot z\rvert\,\mathrm{d}x_0\),
  \(\sigma = \mathrm{sgn}(\pi/2 - z)\ln(\sin z)/(6H)\) on one period
  \(0 < z < \pi\); \(\sigma = 0\) is the fold \(z = \pi/2\),
  \(\sigma \to \mp\infty\) the tips \(z \to 0, \pi\); one period of
  \(x_0\) is the whole line \(-\infty < \sigma < \infty\). Unit: length.
  Status: DERIVED.
\item
  \(V_*\): The physical volume of the extra times \(x_5, x_6, x_7\) on
  the fold in the late universe,
  \(V_* = V_{\mathrm{t}}(\mathrm{e}^{-a_4}b_4)^3\) (with
  \(V_{\mathrm{t}} = \int\mathrm{d}x_5\,\mathrm{d}x_6\,\mathrm{d}x_7\)
  the coordinate volume of the extra times) at a time when
  \(\mathrm{e}^{-a_4}b_4\) is constant; it converts the
  eight-dimensional normalisation into the four-dimensional one. Unit:
  length\(^3\) (GeV\(^{-3}\)). Status: ASSUMED.
\item
  \(\mathbb{F}\): Any function of the fields at the fold (used to state
  a rule of integration). Status: STANDARD.
\item
  \(x\): The coordinates of the fold: \(x_1, x_2, x_3\) of ordinary
  space, \(x_4\) the time, \(x_5, x_6, x_7\) the extra times; a field of
  the fold is a function of the four observable coordinates
  \(x = (x_1, x_2, x_3, x_4)\), and
  \(\mathrm{d}^4x = \mathrm{d}x_1\,\mathrm{d}x_2\,\mathrm{d}x_3\,\mathrm{d}x_4\).
  Unit: length (each coordinate). Status: AUTHOR.
\item
  \(x_1\): The first coordinate of ordinary 3-space. Unit: length.
  Status: AUTHOR.
\item
  \(x_2\): The second coordinate of ordinary 3-space. Unit: length.
  Status: AUTHOR.
\item
  \(x_3\): The third coordinate of ordinary 3-space. Unit: length.
  Status: AUTHOR.
\item
  \(x_4\): The time; a prime on a function of \(x_4\) alone is
  \(\mathrm{d}/\mathrm{d}x_4\). Unit: length (= time, \(c = 1\)).
  Status: AUTHOR.
\item
  \(\hat h_{\mu\nu}\): The four-dimensional metric of the fold,
  \(\hat h_{\mu\nu} = \mathrm{diag}(\mathrm{e}^{2a_4}, \mathrm{e}^{2a_4}, \mathrm{e}^{2a_4}, -1)\)
  in the coordinates \(x_1, \dots, x_4\) (the induced metric at
  \(\sigma = 0\) without the extra times); \(\hat h^{\mu\nu}\) its
  inverse, \(\hat h\) its determinant,
  \(\sqrt{-\hat h} = \mathrm{e}^{3a_4}\); \(\hat h^{(1)}_{\mu\nu}\) its
  first-order cosmological perturbation. Unit: dimensionless. Status:
  DERIVED.
\item
  \(\mu\): A four-dimensional coordinate index running over the
  observable coordinates \(x_1, x_2, x_3, x_4\) of the fold (signature
  \(+,+,+,-\)); a repeated index is summed. Status: STANDARD.
\item
  \(\nu\): A four-dimensional coordinate index (as \(\mu\)). Status:
  STANDARD.
\item
  \(\mathrm{diag}\): A diagonal matrix with the listed entries on its
  diagonal (in the order of the coordinates \(x_0, x_1, \dots, x_7\) for
  a tensor of eight dimensions, \(x_1, \dots, x_7\) for a tensor of the
  fold, \(x_1, \dots, x_4\) for a four-dimensional tensor). Status:
  STANDARD.
\item
  \(\mathrm{e}\): Euler's number, the base of the natural logarithm;
  \(\mathrm{e}^{X}\) is the exponential of \(X\). Unit: dimensionless.
  Value: 2.718281828. Status: STANDARD.
\item
  \(a_4\): The author's metric function of the time \(x_4\) that scales
  ordinary 3-space: on the fold \(\mathrm{e}^{a_4}\) is the scale factor
  of cosmology and \(a_4'\) the Hubble rate. Unit: dimensionless.
  Status: AUTHOR.
\item
  \(\mathcal{Z}\): The fold factor,
  \(\mathcal{Z}(x_4) = V_{\mathrm{t}}(\mathrm{e}^{-a_4}b_4)^3/V_*\): the
  factor by which the size of the extra times multiplies every term of
  the fold and the kinetic terms of the massless modes of the fields of
  the bulk; \(\mathcal{Z} = 1\) when the extra times are frozen at their
  late value. Unit: dimensionless. Value: 1 in the late universe.
  Status: DERIVED.
\item
  \(V_{\mathrm{t}}\): The coordinate volume of the (compact) extra
  times,
  \(V_{\mathrm{t}} = \int\mathrm{d}x_5\,\mathrm{d}x_6\,\mathrm{d}x_7\).
  Unit: length\(^3\). Status: ASSUMED.
\item
  \(b_4\): The author's metric function of the time \(x_4\) in the
  extra-time components, \(b_4 > 0\), never identically 1 (the author's
  order of 2026-10-24); the scale factor (the factor of the lengths) of
  the extra times is \(\mathrm{e}^{-a_4}b_4\) on the fold and
  \(\mathrm{e}^{-a_4}b_4\,\mathrm{e}^{-H\lvert\sigma\rvert}\) at the
  distance \(\sigma\) from it (the square root of the warp factor times
  \(\mathrm{e}^{-a_4}b_4\)). Unit: dimensionless. Status: AUTHOR.
\item
  \(x_5\): The first of the three extra time-like coordinates. Unit:
  length. Status: AUTHOR.
\item
  \(x_6\): The second of the three extra time-like coordinates. Unit:
  length. Status: AUTHOR.
\item
  \(x_7\): The third of the three extra time-like coordinates. Unit:
  length. Status: AUTHOR.
\end{itemize}

The fold is the seven-dimensional surface \(\sigma = 0\) (\(z = \pi/2\))
with the coordinates \(x_1, \dots, x_7\); at \(\sigma = 0\) the metric
of UM-4 is that of the fold (FD-2), the same as in UM-3. No field
depends on the extra times (ASSUMED, FE-38), so the integral over them
gives their coordinate volume \(V_{\mathrm{t}}\) (the extra times are
compact) times their scale factors \((\mathrm{e}^{-a_4}b_4)^3\) at the
fold. \(V_*\) is the physical volume of the extra times on the fold in
the late universe, where \(\mathrm{e}^{-a_4}b_4\) is constant (FE-72);
with this normalisation \(\mathcal{Z} = 1\) in the whole late universe,
the fields of the fold are four-dimensional fields in the
four-dimensional metric \(\hat h_{\mu\nu}\) of the fold, and their
couplings are the measured four-dimensional ones. \(\mathcal{Z}(x_4)\)
is the factor by which the extra times rescale every term of the fold
when \(\mathrm{e}^{-a_4}b_4\) changes, and with it the couplings of the
fold (FE-44). \(\mathbb{F}\) stands for any function of the fields on
the fold. Status: ASSUMED (the place of the fields), DERIVED (the
reduction).

\[
\mathcal{L}_{\mathrm{IG}} = \frac{\hat R}{2\kappa_{\mathrm{ind}}} \tag{IN-9}
\] \textbf{Symbols in (IN-9).}

\begin{itemize}
\tightlist
\item
  \(\mathcal{L}_{\mathrm{IG}}\): The Lagrangian of the Einstein-Hilbert
  term of the four-dimensional metric of the fold that the quantum
  fluctuations of the fields of the fold induce; it adds to the
  Einstein-Hilbert term that the graviton of the bulk gives the fold and
  does not change the speed of gravitational waves. Unit: GeV\(^4\).
  Status: ESTIMATE.
\item
  \(\hat R\): The curvature of the metric \(\hat g_{ab}\) of the fold
  (FD-2): with four frame indices its Riemann tensor, with two its Ricci
  tensor \(\hat R_{ab}\), without indices its scalar curvature
  \(\hat R\); for frozen extra times the scalar curvature equals that of
  the four-dimensional metric \(\hat h_{\mu\nu}\). Unit: 1/length\(^2\).
  Status: STANDARD.
\item
  \(\kappa_{\mathrm{ind}}\): The coupling of the Einstein-Hilbert term
  \(\hat R/(2\kappa_{\mathrm{ind}})\) that the quantum fluctuations of
  the fields of the fold may induce,
  \(\kappa_{\mathrm{ind}} = 8\pi G_{\mathrm{ind}}\); it adds to the
  gravity of the fold,
  \(1/(8\pi G_{\mathrm{N}}) = V_*/\kappa_7 + 1/\kappa_{\mathrm{ind}}\).
  Unit: GeV\(^{-2}\). Value: not fixed by UM-4; negligible for a cutoff
  of the order of \(H\). Status: ESTIMATE.
\end{itemize}

The Einstein-Hilbert term of the four-dimensional metric
\(\hat h_{\mu\nu}\) of the fold (\(\hat R\) its scalar curvature) that
the quantum fluctuations of the fields of the fold induce (FE-64), with
the coupling \(\kappa_{\mathrm{ind}} = 8\pi G_{\mathrm{ind}}\). It adds
\(1/\kappa_{\mathrm{ind}}\) to the coefficient \(V_*/\kappa_7\) of the
Einstein-Hilbert term that the bulk gives the fold (FE-63, FE-71), so
that four-dimensional gravity has
\(1/(8\pi G_{\mathrm{N}}) = V_*/\kappa_7 + 1/\kappa_{\mathrm{ind}}\)
(FE-59), and it has the same light cone, so that it leaves the speed of
gravitational waves equal to the speed of light (FE-63). Its strength is
not fixed by UM-4 (FE-64); for a cutoff of the order of \(H\) it is
negligible. Status: ESTIMATE.

\[
\frac{1}{2}\Big(\frac{\partial^2h_{\mathrm{gw}}}{\partial x_4^2} + 3a_4'\,\frac{\partial h_{\mathrm{gw}}}{\partial x_4}
 - \mathrm{e}^{-2a_4}\,\frac{\partial^2h_{\mathrm{gw}}}{\partial x_3^2}\Big) = \kappa_7\,\Pi_{\mathrm{gw}},\qquad
\kappa_7 = \frac52\,\kappa H,\qquad c_{\mathrm{g}} = 1 \tag{FE-63}
\] \textbf{Symbols in (FE-63).}

\begin{itemize}
\tightlist
\item
  \(\partial\): A partial derivative: the derivative with respect to one
  variable at fixed values of the others (\(\partial_\sigma\) with
  respect to \(\sigma\); \(\partial_M\) and \(\partial_\mu\) with
  respect to the coordinate that carries the same index, of the eight
  dimensions or of the fold). Status: STANDARD.
\item
  \(h_{\mathrm{gw}}\): The amplitude of a gravitational wave of the +
  polarisation travelling along \(x_3\): the metric components along
  \(x_1\) and \(x_2\) multiplied by \(\mathrm{e}^{h_{\mathrm{gw}}}\) and
  \(\mathrm{e}^{-h_{\mathrm{gw}}}\); a function of \(\sigma\), \(x_3\)
  and \(x_4\) in the bulk, and on the fold the gravitational wave of the
  four-dimensional world, to the first order its plus polarisation
  \(h_{+}\). Unit: dimensionless. Status: STANDARD.
\item
  \(x_4\): The time; a prime on a function of \(x_4\) alone is
  \(\mathrm{d}/\mathrm{d}x_4\). Unit: length (= time, \(c = 1\)).
  Status: AUTHOR.
\item
  \(a_4\): The author's metric function of the time \(x_4\) that scales
  ordinary 3-space: on the fold \(\mathrm{e}^{a_4}\) is the scale factor
  of cosmology and \(a_4'\) the Hubble rate. Unit: dimensionless.
  Status: AUTHOR.
\item
  \({}'\): The derivative of a function of the time \(x_4\) alone with
  respect to \(x_4\) (\(a_4' = \mathrm{d}a_4/\mathrm{d}x_4\); two primes
  the second derivative); on a function of \(\sigma\) alone (a profile)
  the derivative with respect to \(\sigma\). Status: STANDARD.
\item
  \(\mathrm{e}\): Euler's number, the base of the natural logarithm;
  \(\mathrm{e}^{X}\) is the exponential of \(X\). Unit: dimensionless.
  Value: 2.718281828. Status: STANDARD.
\item
  \(x_3\): The third coordinate of ordinary 3-space. Unit: length.
  Status: AUTHOR.
\item
  \(\kappa_7\): The gravitational coupling of the fold: to the first
  order in its curvature over \(H^2\) the fold obeys the Einstein
  equations of seven dimensions
  \(\hat G^a{}_b = \kappa_7\mathsf{t}^a{}_b\) with
  \(\kappa_7 = 5\kappa H/2\), the coupling of the massless graviton of
  the bulk bound to the fold; four-dimensional gravity on the fold has
  \(8\pi G_{\mathrm{N}} = \kappa_7/V_*\) (plus the induced term, which
  is negligible). Unit: GeV\(^{-5}\). Value: \(5\kappa H/2\). Status:
  DERIVED.
\item
  \(\Pi_{\mathrm{gw}}\): The transverse traceless stress of the matter
  of the fold that drives the gravitational wave \(h_{\mathrm{gw}}\)
  (half the difference of its components along \(x_1\) and \(x_2\)),
  integrated across the fold (eight-dimensional normalisation). Unit:
  GeV\(^7\). Status: STANDARD.
\item
  \(\kappa\): The eight-dimensional gravitational coupling,
  \(\kappa = 8\pi G_8\). Unit: GeV\(^{-6}\). Status: AUTHOR.
\item
  \(H\): The author's constant (an inverse length) that fixes the period
  \(\pi/(6H)\) of \(x_0\), the warp factor
  \(\mathrm{e}^{-2H\lvert\sigma\rvert}\) of the metric, the cosmological
  constant \(\Lambda = -21H^2\) and the tension
  \(\kappa\lambda_{\mathrm{F}} = 12H\) of the fold; in UM-4 a free
  parameter, bounded from below by the propagation of the gauge fields
  in the bulk. Unit: GeV (1/length). Value: \(H > 2534\) GeV, about 2.5
  TeV (from the invisible width of the Z boson). Status: AUTHOR.
\item
  \(c_{\mathrm{g}}\): The speed of gravitational waves on the fold, in
  units of the speed of light. Unit: dimensionless. Value: 1 at every
  wavelength for the four-dimensional graviton on the static block and
  on every fold with frozen extra times, and for the massless graviton
  of the bulk at every depth; measured
  \(-3\times10^{-15} \le c_{\mathrm{g}} - 1 \le 7\times10^{-16}\)
  (GW170817). Status: DERIVED.
\end{itemize}

The gravitational wave of the fold: a wave of the + polarisation
travelling along \(x_3\) on a fold with frozen extra times (any
\(a_4(x_4)\), the static block included). The left-hand side is half the
difference of the components along \(x_1\) and \(x_2\) of the Einstein
tensor of the metric of the fold to the first order in
\(h_{\mathrm{gw}}\) (computed exactly), and the right-hand side is the
same difference of \(\kappa_7\) times the matter of the fold,
\(\Pi_{\mathrm{gw}}\) its transverse traceless stress: by FE-71 the fold
obeys the Einstein equations of seven dimensions with
\(\kappa_7 = 5\kappa H/2\), the same coupling as in its Friedmann
equations. FE-63 is the wave equation of four-dimensional general
relativity with the friction \(3a_4'\) of the expansion: the
gravitational waves of the four-dimensional world travel at the speed of
light at every wavelength, and their Planck mass equals that of the
expansion. In the bulk the same wave is the zero mode of FE-76:
\(h_{\mathrm{gw}}\) independent of \(\sigma\) solves the bulk equation
FE-68 (with \(\mathsf{P}'(1) = 21\)) and the junction condition, so the
wave of the fold is a wave of the bulk at the speed of light that the
fold holds. The gravitational waves of the neutron-star merger GW170817
arrived together with its light,
\(-3\times10^{-15} \le c_{\mathrm{g}} - 1 \le 7\times10^{-16}\) for the
distance 26 Mpc (B. P. Abbott et al., Astrophys. J. Lett. 848 (2017)
L13): UM-4 agrees with it for every value of \(H\), at every order of
the bulk (the light cone of every slice is that of the fold). The
induced term IN-9 adds to the coefficient of the same operator and keeps
\(c_{\mathrm{g}} = 1\). Status: DERIVED (\texttt{um4\_foldsources.py},
FS2; \texttt{um4\_bulk.py}, B4 and B5; \texttt{um4\_fold.py}, F2).

\[
\frac{1}{16\pi G_{\mathrm{ind}}} = \frac{N_{\mathrm{ind}}\,\Lambda_{\mathrm{c}}^2}{192\pi^2},\qquad
N_{\mathrm{ind}} = \sum_{\text{real scalars}}\big(1 - 6\xi\big) + N_{\mathrm{Weyl}} - 4N_{\mathrm{gauge}} = 52 - 24\xi_{\mathrm{H}} \tag{FE-64}
\] \textbf{Symbols in (FE-64).}

\begin{itemize}
\tightlist
\item
  \(\pi\): The number pi, the ratio of the circumference of a circle to
  its diameter. Unit: dimensionless. Value: 3.141592654. Status:
  STANDARD.
\item
  \(G_{\mathrm{ind}}\): The Newton constant of the induced
  Einstein-Hilbert term of the fold,
  \(\kappa_{\mathrm{ind}} = 8\pi G_{\mathrm{ind}}\). Unit: GeV\(^{-2}\).
  Status: ESTIMATE.
\item
  \(N_{\mathrm{ind}}\): The weighted number of the fields of the fold
  that induce the Einstein-Hilbert term: each real scalar counts
  \(1 - 6\xi\), each Weyl fermion 1, each gauge field \(-4\) (with its
  ghost fields). Unit: dimensionless. Value: 52 for the fields of the
  fold of UM-4 (the four real scalars of the Higgs doublet and 48 Weyl
  fermions; the gauge fields live in the bulk) with
  \(\xi_{\mathrm{H}} = 0\); \(52 - 24\xi_{\mathrm{H}}\) with the
  curvature coupling \(\xi_{\mathrm{H}}\) of an extension. Status:
  DERIVED.
\item
  \(\Lambda_{\mathrm{c}}\): The cutoff of the proper-time integral of
  the one-loop effective action of the fields of the fold: the highest
  energy of the quantum fluctuations that induce gravity. Unit: GeV.
  Value: taken of the order of \(H\) (2534 GeV); the induced term alone
  would give the measured Newton constant for
  \(\Lambda_{\mathrm{c}} = \pi M_{\mathrm{red}}\sqrt{96/52} = 1.0\times10^{19}\)
  GeV. Status: ESTIMATE.
\item
  \(\xi\): The coupling of a real scalar field to the scalar curvature
  \(\hat R\) of the fold: the term \(-\frac12\xi\hat R\) times the
  square of the field in its Lagrangian (\(\xi = 1/6\) is the conformal
  coupling). Unit: dimensionless. Status: STANDARD.
\item
  \(N_{\mathrm{Weyl}}\): The number of Weyl (two-component) fermion
  fields of the fold. Unit: dimensionless. Value: 48 (three generations
  of quarks and leptons with three right-handed neutrinos). Status:
  STANDARD.
\item
  \(N_{\mathrm{gauge}}\): The number of gauge fields of the fold. Unit:
  dimensionless. Value: 0 in UM-4 (the gauge fields live in the bulk);
  12 (eight gluons, three \(W\), one \(B\)) if they lived on the fold.
  Status: STANDARD.
\item
  \(\xi_{\mathrm{H}}\): The coupling of the Higgs doublet to the scalar
  curvature of the fold, the term
  \(-\xi_{\mathrm{H}}\hat R\,\Phi^\dagger\Phi\) in the Lagrangian of the
  fold: an extension of UM-4, whose Higgs Lagrangian has no such term
  (\(\xi_{\mathrm{H}} = 0\)). Unit: dimensionless. Value: 0 in UM-4;
  about \(1.6\times10^4\) if the Higgs field drove inflation (ESTIMATE,
  an extension of the action). Status: ESTIMATE (an extension of the
  action).
\end{itemize}

The Einstein-Hilbert term that the quantum fluctuations of the fields of
the fold induce (Sakharov's induced gravity, A. D. Sakharov, Dokl. Akad.
Nauk SSSR 177 (1967) 70), evaluated at one loop with the heat kernel and
the proper-time cutoff \(\Lambda_{\mathrm{c}}\) (the coefficient of the
scalar curvature in the heat kernel of a Laplace-type operator
\(-(\nabla^2 + E)\) is the trace of \(E + R/6\); D. V. Vassilevich,
Phys. Rep.~388 (2003) 279): each real scalar with the curvature coupling
\(\xi\) counts \(1 - 6\xi\), each Weyl fermion 1, each gauge field of
the fold \(-4\) with its ghost fields (S. N. Solodukhin,
arXiv:1502.03758). In UM-4 the gauge fields live in the bulk, so the
fold carries the four real scalars of the Higgs doublet and the 48 Weyl
fermions of three generations with three right-handed neutrinos, and no
gauge field (\(N_{\mathrm{gauge}} = 0\)), so
\(N_{\mathrm{ind}} = 4(1 - 6\xi_{\mathrm{H}}) + 48\). The Higgs
Lagrangian of UM-4 (IN-11) has no coupling to the curvature of the fold,
so \(\xi_{\mathrm{H}} = 0\) and \(N_{\mathrm{ind}} = 52\). For a cutoff
\(\Lambda_{\mathrm{c}}\) of the order of \(H\) (2534 GeV) the induced
term is a part \(5.9\times10^{-32}\) of \(M_{\mathrm{red}}^2\)
(\(M_{\mathrm{red}}\) the reduced Planck mass): four-dimensional gravity
is the gravity of the bulk (FE-59). The same loops give the fold a
vacuum energy of about
\(\lvert N_{\mathrm{B}} - N_{\mathrm{F}}\rvert\Lambda_{\mathrm{c}}^4/(64\pi^2)\)
(\(N_{\mathrm{B}} = 4\) bosonic and \(N_{\mathrm{F}} = 96\) fermionic
degrees of freedom; \(6.0\times10^{12}\) GeV\(^4\) for
\(\Lambda_{\mathrm{c}} = 2534\) GeV), a tension that adds to
\(\lambda_{\mathrm{F}}\) (FE-65). Status: ESTIMATE (one loop and
cutoff-dependent; \texttt{um4\_foldsources.py}, FS3).

\[
\mathcal{T}_{\mathrm{ff}}{}^a{}_b = \mathrm{diag}\big(p_{\mathrm{ff}}, p_{\mathrm{ff}}, p_{\mathrm{ff}}, -\rho_{\mathrm{ff}},
 q_{\mathrm{ff}}, q_{\mathrm{ff}}, q_{\mathrm{ff}}\big) \tag{IN-10}
\] \textbf{Symbols in (IN-10).}

\begin{itemize}
\tightlist
\item
  \(\mathcal{T}_{\mathrm{ff}}{}^a{}_b\): The energy-momentum tensor of
  the fold fluid (frame components of the fold, four-dimensional
  normalisation). Unit: GeV\(^4\). Status: ASSUMED.
\item
  \(a\): A frame (orthonormal) index running over \(0, 1, \dots, 7\)
  (frame direction \((0)\) along \(\sigma\), \((1), (2), (3)\) along
  ordinary space, \((4)\) along the time, \((5), (6), (7)\) along the
  extra times); raised and lowered with the frame metric
  \(\eta_{ab} = \mathrm{diag}(1, 1, 1, 1, -1, -1, -1, -1)\); on the fold
  (\(a, b \in \{1, \dots, 7\}\)) the frame of the fold; in the tensors
  of the fold written with its coordinates \(x^a\) (its metric
  \(\hat g_{ab}\) and the tensors built from it) a coordinate index of
  the fold, \(a \in \{1, \dots, 7\}\); a repeated index is summed.
  Status: STANDARD.
\item
  \(b\): A frame index (as \(a\)). Status: STANDARD.
\item
  \(\mathrm{diag}\): A diagonal matrix with the listed entries on its
  diagonal (in the order of the coordinates \(x_0, x_1, \dots, x_7\) for
  a tensor of eight dimensions, \(x_1, \dots, x_7\) for a tensor of the
  fold, \(x_1, \dots, x_4\) for a four-dimensional tensor). Status:
  STANDARD.
\item
  \(p_{\mathrm{ff}}\): The pressure of the fold fluid along ordinary
  space. Unit: GeV\(^4\). Value: \(-\rho_{\mathrm{ff}}\) (a tension).
  Status: ASSUMED (the fluid), DERIVED (the value).
\item
  \(\rho_{\mathrm{ff}}\): The energy density of the fold fluid. Unit:
  GeV\(^4\). Value:
  \(\lambda_{\mathrm{F}}V_* + \rho_\Lambda - \rho_{\mathrm{vac}}\):
  together with the vacuum energy \(\rho_{\mathrm{vac}}\) of the fields
  of the fold it makes up the tension \(\lambda_{\mathrm{F}}V_*\) that
  the static block requires plus the dark energy \(\rho_\Lambda\).
  Status: ASSUMED (the fluid), DERIVED (the value).
\item
  \(q_{\mathrm{ff}}\): The pressure of the fold fluid along the extra
  times: minus its energy density (its tension) plus the pressure
  \(q_{\mathrm{hold}}\) that holds the volume of the extra times fixed.
  Unit: GeV\(^4\). Value: \(-\rho_{\mathrm{ff}} + q_{\mathrm{hold}}\).
  Status: ASSUMED (the fluid), DERIVED (the value).
\end{itemize}

The fold fluid: a medium at rest on the fold (the frame indices \(a, b\)
run over the seven directions \(1, \dots, 7\) of the fold), with the
energy density \(\rho_{\mathrm{ff}}\), the pressure \(p_{\mathrm{ff}}\)
along ordinary space and the pressure \(q_{\mathrm{ff}}\) along the
extra times (four-dimensional normalisation, IN-8). The kink of the
metric at the fold requires a thin source (FE-4); the fold fluid is the
part of it that the fields of the fold, the particles \(\chi\) of
dirac16complex and the portals do not supply. It has two parts (FE-65):
a tension, a vacuum energy of the seven-dimensional fold like the
tension of the brane of Randall and Sundrum, which the static block
requires; and a pressure \(q_{\mathrm{hold}}\) along the extra times,
which holds their volume fixed. Without \(q_{\mathrm{hold}}\) the volume
of the extra times on the fold changes whenever the matter of the fold
is not radiation (FE-72), and with it Newton's constant, the squares of
the gauge and Yukawa couplings and the quartic coupling of the Higgs
doublet: in a matter era the extra times would expand as ordinary space
and Newton's constant would fall at three times the Hubble rate, which
lunar laser ranging excludes by a factor of about 8500. UM-4 therefore
assumes that the fold fluid supplies \(q_{\mathrm{hold}}\). Its
microscopic origin is OPEN: it acts as a potential for the volume of the
extra times, as the stabilisation of the sizes of compact extra
dimensions does in other theories. Status: ASSUMED (its existence),
DERIVED (its values, FE-65).

\[
\begin{aligned}
&p_{\mathrm{ff}} = -\rho_{\mathrm{ff}},\qquad q_{\mathrm{ff}} = -\rho_{\mathrm{ff}} + q_{\mathrm{hold}},\qquad
 \kappa\lambda_{\mathrm{F}} = 12H\ \ \big(\lambda_{\mathrm{F}}V_* = 30H^2M_{\mathrm{red}}^2\big),\\
&q_{\mathrm{hold}} = -\frac{\rho - 3\mathcal{P}}{2} - q_{\mathrm{SM}} - q_\chi - q_{\mathrm{P}}:\qquad
 \text{radiation } 0,\quad \text{dust } -\frac{\rho}{2},\quad \text{a vacuum energy } -\rho
\end{aligned} \tag{FE-65}
\] \textbf{Symbols in (FE-65).}

\begin{itemize}
\tightlist
\item
  \(p_{\mathrm{ff}}\): The pressure of the fold fluid along ordinary
  space. Unit: GeV\(^4\). Value: \(-\rho_{\mathrm{ff}}\) (a tension).
  Status: ASSUMED (the fluid), DERIVED (the value).
\item
  \(\rho_{\mathrm{ff}}\): The energy density of the fold fluid. Unit:
  GeV\(^4\). Value:
  \(\lambda_{\mathrm{F}}V_* + \rho_\Lambda - \rho_{\mathrm{vac}}\):
  together with the vacuum energy \(\rho_{\mathrm{vac}}\) of the fields
  of the fold it makes up the tension \(\lambda_{\mathrm{F}}V_*\) that
  the static block requires plus the dark energy \(\rho_\Lambda\).
  Status: ASSUMED (the fluid), DERIVED (the value).
\item
  \(q_{\mathrm{ff}}\): The pressure of the fold fluid along the extra
  times: minus its energy density (its tension) plus the pressure
  \(q_{\mathrm{hold}}\) that holds the volume of the extra times fixed.
  Unit: GeV\(^4\). Value: \(-\rho_{\mathrm{ff}} + q_{\mathrm{hold}}\).
  Status: ASSUMED (the fluid), DERIVED (the value).
\item
  \(q_{\mathrm{hold}}\): The pressure along the extra times that the
  fold fluid adds to its tension to hold the volume of the extra times
  fixed (four-dimensional normalisation). Unit: GeV\(^4\). Value:
  \(-(\rho - 3\mathcal{P})/2 - q_{\mathrm{SM}} - q_\chi - q_{\mathrm{P}}\):
  0 for radiation, \(-\rho/2\) for dust, \(-\rho\) for a vacuum energy;
  today \(-0.842\,\rho_{\mathrm{c}}^{(0)}\) if the baryons have no
  pressure along the extra times (ASSUMED). Status: DERIVED (the value),
  ASSUMED (that the fold fluid supplies it).
\item
  \(\kappa\): The eight-dimensional gravitational coupling,
  \(\kappa = 8\pi G_8\). Unit: GeV\(^{-6}\). Status: AUTHOR.
\item
  \(\lambda_{\mathrm{F}}\): The tension that the static block requires
  of the fold, \(\kappa\lambda_{\mathrm{F}} = 12H\): the total vacuum
  energy of the seven-dimensional fold, made of the tension of the fold
  fluid and the vacuum energy of the fields of the fold (both have the
  pressure minus their energy density along ordinary space and along the
  extra times), per unit volume of the six other directions of the fold
  (eight-dimensional normalisation); in the four-dimensional
  normalisation it is \(\lambda_{\mathrm{F}}V_*\), and the excess of the
  total over \(\lambda_{\mathrm{F}}V_*\) is the dark energy
  \(\rho_\Lambda\). Unit: GeV\(^7\). Value: \(12H/\kappa\);
  \(\lambda_{\mathrm{F}}V_* = 30H^2M_{\mathrm{red}}^2 = 1.1\times10^{45}\)
  GeV\(^4\) for \(H = 2534\) GeV. Status: DERIVED.
\item
  \(H\): The author's constant (an inverse length) that fixes the period
  \(\pi/(6H)\) of \(x_0\), the warp factor
  \(\mathrm{e}^{-2H\lvert\sigma\rvert}\) of the metric, the cosmological
  constant \(\Lambda = -21H^2\) and the tension
  \(\kappa\lambda_{\mathrm{F}} = 12H\) of the fold; in UM-4 a free
  parameter, bounded from below by the propagation of the gauge fields
  in the bulk. Unit: GeV (1/length). Value: \(H > 2534\) GeV, about 2.5
  TeV (from the invisible width of the Z boson). Status: AUTHOR.
\item
  \(V_*\): The physical volume of the extra times \(x_5, x_6, x_7\) on
  the fold in the late universe,
  \(V_* = V_{\mathrm{t}}(\mathrm{e}^{-a_4}b_4)^3\) (with
  \(V_{\mathrm{t}} = \int\mathrm{d}x_5\,\mathrm{d}x_6\,\mathrm{d}x_7\)
  the coordinate volume of the extra times) at a time when
  \(\mathrm{e}^{-a_4}b_4\) is constant; it converts the
  eight-dimensional normalisation into the four-dimensional one. Unit:
  length\(^3\) (GeV\(^{-3}\)). Status: ASSUMED.
\item
  \(M_{\mathrm{red}}\): The reduced Planck mass,
  \(M_{\mathrm{red}} = (8\pi G_{\mathrm{N}})^{-1/2}\). Unit: GeV. Value:
  \(2.435\times10^{18}\) GeV. Status: STANDARD.
\item
  \(\rho\): The total energy density of the matter of the fold (all
  species); \(\Delta\rho\) its perturbation. Unit: GeV\(^4\). Status:
  STANDARD.
\item
  \(\mathcal{P}\): The total pressure of the matter of the fold;
  \(\Delta\mathcal{P}\) its perturbation. Unit: GeV\(^4\). Status:
  STANDARD.
\item
  \(q_{\mathrm{SM}}\): The pressure of the Standard Model along the
  extra times at the fold: the mean on-shell Lagrangian density of the
  Standard-Model fields (the extra-time component of the energy-momentum
  tensor of every term of the fold is its Lagrangian, because the fields
  of the fold do not depend on the extra times). Unit: GeV\(^4\).
  Status: DERIVED (for radiation and free fields), ASSUMED (for the
  baryons).
\item
  \(q_\chi\): The pressure of the particles \(\chi\) of dirac16complex
  along the extra times, integrated across the fold (their on-shell
  Lagrangian, the self-interaction energy per volume; zero for a dilute
  gas). Unit: GeV\(^4\). Status: DERIVED.
\item
  \(q_{\mathrm{P}}\): The pressure of the portals along the extra times:
  their mean on-shell Lagrangian density. Unit: GeV\(^4\). Status:
  DERIVED.
\end{itemize}

The fold fluid of UM-4 with frozen extra times. Its energy density and
its pressure along ordinary space are those of a tension. The static
block requires the total tension \(\kappa\lambda_{\mathrm{F}} = 12H\)
(FD-12, FE-30), which in the four-dimensional normalisation is
\(\lambda_{\mathrm{F}}V_* = 30H^2M_{\mathrm{red}}^2\) (FE-71): the
tension of the fold fluid together with the vacuum energy of the fields
of the fold (their potential at the minimum), which is a tension of the
same kind, because the fields do not depend on the extra times and their
pressure along the extra times is their on-shell Lagrangian (FE-5). The
excess of the total tension over \(\lambda_{\mathrm{F}}\) is the vacuum
energy \(\rho_\Lambda\) with which the fold expands, the dark energy of
UM-4 (FE-59; its value needs the tension tuned to one part in
\(10^{92}\) for \(H = 2534\) GeV, the cosmological-constant problem of
the brane of Randall and Sundrum, OPEN). When the fold expands with
frozen extra times, the component of its equations along the extra times
requires of the matter above \(\lambda_{\mathrm{F}}\) the pressure
\(-(\rho - 3\mathcal{P})/2\) along the extra times (\(\rho\),
\(\mathcal{P}\) its energy density and its pressure along ordinary
space; FE-72). The pressure of a species of the fold along the extra
times is its mean on-shell Lagrangian density (FE-5): radiation has
none, and supplies what is required (\(\rho = 3\mathcal{P}\)); dust has
none and needs \(-\rho/2\) (dust of free particles: the mean on-shell
Lagrangian of a free field vanishes, for the Dirac field because its
Lagrangian vanishes on its equation of motion, for a free scalar because
the averages of its kinetic and mass terms are equal; for the baryons
this is ASSUMED, because the on-shell Lagrangian of the quark and gluon
fields inside a nucleon is not computed); a vacuum energy has \(-\rho\)
and needs \(-2\rho\). The particles \(\chi\) of dirac16complex have no
pressure along the extra times in a dilute gas (the on-shell Lagrangian
of a free Dirac field vanishes; the contact term gives
\(+(\lambda_\Psi/2)(\bar\Psi\Psi)^2\), negligible in a dilute gas).
\(q_{\mathrm{hold}}\) makes up the difference. Status: DERIVED
(\texttt{um4\_foldsources.py}, FS4 and FS5; \texttt{um4\_fold.py}, F4
and F5; \texttt{um4\_background.py}, BG2 and BG3), ASSUMED (the pressure
of the baryons along the extra times; that the fold fluid supplies
\(q_{\mathrm{hold}}\)), OPEN (what supplies it, and the tuning of the
tension).

\hypertarget{the-higgs-doublet-with-the-ad-hoc-vacuum}{%
\subsection{The Higgs doublet with the ad hoc
vacuum}\label{the-higgs-doublet-with-the-ad-hoc-vacuum}}

\[
\mathcal{L}_{\mathrm{H}} = -\hat h^{\mu\nu}\big(D_\mu\Phi\big)^\dagger D_\nu\Phi - V(\Phi),\qquad
V(\Phi) = \mu^2\,\Phi^\dagger\Phi + \lambda\big(\Phi^\dagger\Phi\big)^2,\qquad \lambda > 0,\quad \mu^2 < 0 \tag{IN-11}
\] \textbf{Symbols in (IN-11).}

\begin{itemize}
\tightlist
\item
  \(\mathcal{L}_{\mathrm{H}}\): The Lagrangian of the Higgs doublet on
  the fold. Unit: GeV\(^4\). Status: AUTHOR.
\item
  \(\hat h_{\mu\nu}\): The four-dimensional metric of the fold,
  \(\hat h_{\mu\nu} = \mathrm{diag}(\mathrm{e}^{2a_4}, \mathrm{e}^{2a_4}, \mathrm{e}^{2a_4}, -1)\)
  in the coordinates \(x_1, \dots, x_4\) (the induced metric at
  \(\sigma = 0\) without the extra times); \(\hat h^{\mu\nu}\) its
  inverse, \(\hat h\) its determinant,
  \(\sqrt{-\hat h} = \mathrm{e}^{3a_4}\); \(\hat h^{(1)}_{\mu\nu}\) its
  first-order cosmological perturbation. Unit: dimensionless. Status:
  DERIVED.
\item
  \(\mu\): A four-dimensional coordinate index running over the
  observable coordinates \(x_1, x_2, x_3, x_4\) of the fold (signature
  \(+,+,+,-\)); a repeated index is summed. Status: STANDARD.
\item
  \(\nu\): A four-dimensional coordinate index (as \(\mu\)). Status:
  STANDARD.
\item
  \(D_M,\ D_\mu\): The covariant derivative: on dirac16complex
  \(D_M = \partial_M + \Omega_M\) (with
  \(D_M\bar\Psi = \partial_M\bar\Psi - \bar\Psi\Omega_M\)); on the
  fields of the fold the gauge and spin covariant derivative \(D_\mu\)
  (the partial derivative plus the spin connection of the fold and the
  gauge fields times their couplings and generators);
  \(D^\mu = \hat h^{\mu\nu}D_\nu\); on an adjoint field
  \((D_\mu X)^A = \partial_\mu X^A + g_3f^{ABC}G^B_\mu X^C\) (likewise
  for \(SU(2)\) with \(g_2\epsilon^{IJK}\)). Unit: 1/length. Status:
  STANDARD.
\item
  \(\Phi\): The Higgs doublet (two complex components, hypercharge
  \(1/2\)), a field of the fold. Unit: GeV. Status: AUTHOR.
\item
  \({}^\dagger\): Hermitian conjugation: the complex conjugate of the
  transpose of a matrix or spinor; for an operator on the space of
  states, its adjoint with respect to the scalar product of that space
  (for the quantised dirac16complex the indefinite Krein scalar
  product). Status: STANDARD.
\item
  \(V(\Phi)\): The Higgs potential,
  \(V = \mu^2\Phi^\dagger\Phi + \lambda(\Phi^\dagger\Phi)^2\). Unit:
  GeV\(^4\). Status: AUTHOR.
\item
  \(\mu^2\): The mass parameter of the Higgs potential, \(\mu^2 < 0\)
  (it enters only as \(\mu^2\)). Unit: GeV\(^2\). Value:
  \(\mu^2 = -\lambda v^2 = -m_h^2/2 = -(88.6\ \mathrm{GeV})^2\). Status:
  AUTHOR.
\item
  \(\lambda\): The quartic coupling of the Higgs potential,
  \(\lambda > 0\). Unit: dimensionless. Value:
  \(\lambda = m_h^2/(2v^2) = 0.129\). Status: AUTHOR.
\end{itemize}

The Higgs doublet \(\Phi\) (two complex components, hypercharge \(1/2\))
with the author's potential. In the signature \((+,+,+,-)\) of the fold
the kinetic term carries a minus sign, so that the energy density is
positive. Status: AUTHOR (the potential), STANDARD (the rest).

\[
\Phi = \frac{1}{\sqrt2}\begin{pmatrix}\mathcal{A}(x_4)\\ v + h\end{pmatrix},\qquad
\mathcal{A}(x_4) = v_{\mathcal{A}}\,\mathrm{e}^{-2a_{40}x_4},\qquad a_{40} > 0,\quad x_4 \ge 0,\qquad v^2 = -\frac{\mu^2}{\lambda} \tag{IN-12}
\] \textbf{Symbols in (IN-12).}

\begin{itemize}
\tightlist
\item
  \(\Phi\): The Higgs doublet (two complex components, hypercharge
  \(1/2\)), a field of the fold. Unit: GeV. Status: AUTHOR.
\item
  \(\mathcal{A}\): The author's ad hoc upper component of the Higgs
  vacuum, almightyAdHoc,
  \(\mathcal{A}(x_4) = v_{\mathcal{A}}\mathrm{e}^{-2a_{40}x_4}\). Unit:
  GeV. Status: AUTHOR.
\item
  \(x_4\): The time; a prime on a function of \(x_4\) alone is
  \(\mathrm{d}/\mathrm{d}x_4\). Unit: length (= time, \(c = 1\)).
  Status: AUTHOR.
\item
  \(v\): The vacuum expectation value of the lower component of the
  Higgs doublet, \(v^2 = -\mu^2/\lambda\). Unit: GeV. Value: 246.22 GeV
  \(= (\sqrt2\,G_{\mathrm{F}})^{-1/2}\) with the Fermi constant
  \(G_{\mathrm{F}} = 1.1663788\times10^{-5}\) GeV\(^{-2}\) of the muon
  lifetime (Particle Data Group). Status: STANDARD.
\item
  \(h\): The Higgs field: the fluctuation of the lower component of the
  Higgs doublet about \(v\). Unit: GeV. Value: \(m_h = 125.25\) GeV.
  Status: STANDARD.
\item
  \(v_{\mathcal{A}}\): The amplitude of almightyAdHoc at \(x_4 = 0\);
  the author's almightyAdHoc, \(\mathrm{e}^{-2a_{40}x_4}\) (written
  Exp{[}-2 a40 x4{]} in his notebook), has \(v_{\mathcal{A}} = 1\) in
  the units of his notebook. Unit: GeV. Value: 1 in the units of the
  author's notebook; UM-4 keeps it free (the examples take
  \(v_{\mathcal{A}} = v\) and \(10v\)). Status: AUTHOR (the form),
  ASSUMED (a free amplitude).
\item
  \(\mathrm{e}\): Euler's number, the base of the natural logarithm;
  \(\mathrm{e}^{X}\) is the exponential of \(X\). Unit: dimensionless.
  Value: 2.718281828. Status: STANDARD.
\item
  \(a_{40}\): The decay constant of almightyAdHoc, \(a_{40} > 0\):
  \(\mathcal{A}\) falls by \(\mathrm{e}\) in the time \(1/(2a_{40})\).
  Unit: GeV (1/length). Status: AUTHOR.
\item
  \(\mu^2\): The mass parameter of the Higgs potential, \(\mu^2 < 0\)
  (it enters only as \(\mu^2\)). Unit: GeV\(^2\). Value:
  \(\mu^2 = -\lambda v^2 = -m_h^2/2 = -(88.6\ \mathrm{GeV})^2\). Status:
  AUTHOR.
\item
  \(\lambda\): The quartic coupling of the Higgs potential,
  \(\lambda > 0\). Unit: dimensionless. Value:
  \(\lambda = m_h^2/(2v^2) = 0.129\). Status: AUTHOR.
\end{itemize}

The author's ad hoc vacuum: the upper component of the Higgs doublet is
the prescribed function almightyAdHoc
\(\mathcal{A}(x_4) = \exp(-2a_{40}x_4)\) (the author writes it with the
amplitude \(v_{\mathcal{A}} = 1\) in the unit of the field; UM-4 keeps
\(v_{\mathcal{A}}\) to show where it enters), the lower component is the
vacuum expectation value \(v\) plus the Higgs field \(h\). The Goldstone
fields are absorbed by the gauge choice (unitary gauge). \(\mathcal{A}\)
decays with the rate \(2a_{40}\) from the beginning of the time
\(x_4 = 0\): it acts only in the primordial era. The ad hoc vacuum is
not an exact solution of the Higgs equation; the residuals are FE-26.
Status: AUTHOR.

\hypertarget{the-covariant-derivative-of-the-fields-of-the-fold}{%
\subsection{The covariant derivative of the fields of the
fold}\label{the-covariant-derivative-of-the-fields-of-the-fold}}

\[
\begin{aligned}
D_\mu &= \partial_\mu + \hat\omega_\mu - \mathrm{i}g_3\,G^A_\mu\,t^A - \mathrm{i}g_2\,W^I_\mu\,\frac{\tau^I}{2}
 - \mathrm{i}g_1\,Y_f\,B_\mu,\\
g_3\,G^A_\mu(x) &= G^A_\mu(\sigma = 0, x),\qquad g_2\,W^I_\mu(x) = W^I_\mu(\sigma = 0, x),\qquad g_1\,B_\mu(x) = B_\mu(\sigma = 0, x)
\end{aligned} \tag{IN-13}
\] \textbf{Symbols in (IN-13).}

\begin{itemize}
\tightlist
\item
  \(D_M,\ D_\mu\): The covariant derivative: on dirac16complex
  \(D_M = \partial_M + \Omega_M\) (with
  \(D_M\bar\Psi = \partial_M\bar\Psi - \bar\Psi\Omega_M\)); on the
  fields of the fold the gauge and spin covariant derivative \(D_\mu\)
  (the partial derivative plus the spin connection of the fold and the
  gauge fields times their couplings and generators);
  \(D^\mu = \hat h^{\mu\nu}D_\nu\); on an adjoint field
  \((D_\mu X)^A = \partial_\mu X^A + g_3f^{ABC}G^B_\mu X^C\) (likewise
  for \(SU(2)\) with \(g_2\epsilon^{IJK}\)). Unit: 1/length. Status:
  STANDARD.
\item
  \(\mu\): A four-dimensional coordinate index running over the
  observable coordinates \(x_1, x_2, x_3, x_4\) of the fold (signature
  \(+,+,+,-\)); a repeated index is summed. Status: STANDARD.
\item
  \(\partial\): A partial derivative: the derivative with respect to one
  variable at fixed values of the others (\(\partial_\sigma\) with
  respect to \(\sigma\); \(\partial_M\) and \(\partial_\mu\) with
  respect to the coordinate that carries the same index, of the eight
  dimensions or of the fold). Status: STANDARD.
\item
  \(\hat\omega_\mu\): The spin connection of the four-dimensional metric
  \(\hat h_{\mu\nu}\) of the fold acting on four-dimensional spinors.
  Unit: 1/length. Status: STANDARD.
\item
  \(\mathrm{i}\): The imaginary unit, \(\mathrm{i}^2 = -1\). Unit:
  dimensionless. Status: STANDARD.
\item
  \(g_3\): The \(SU(3)_C\) gauge coupling of the four-dimensional world
  (the strong coupling \(g_s\)): the coupling of the massless mode of
  the gluon field of the bulk, \(g_3^2 = 3H\mathsf{g}_3^2/(2V_*)\).
  Unit: dimensionless. Value: \(g_3(m_Z) = 1.2177\)
  (\(g_3^2/(4\pi) = 0.1180\)). Status: STANDARD.
\item
  \(G^A_M,\ G^A_\mu,\ G^A_{MN}\): The gluon field of \(SU(3)_C\)
  (\(A = 1, \dots, 8\)): in the bulk \(G^A_M\), a field of eight
  dimensions in its geometric normalisation (the covariant derivative
  \(\partial - \mathrm{i}G\)), and on the fold the canonically
  normalised four-dimensional field \(G^A_\mu(x)\) of its massless mode,
  \(g_3G^A_\mu(x) = G^A_\mu(\sigma = 0, x)\); \(G^A_{MN}\) and
  \(G^A_{\mu\nu}\) their field strengths. Unit: GeV (field), GeV\(^2\)
  (field strength). Status: AUTHOR.
\item
  \(A\): A colour (adjoint SU(3)) index running over \(1, \dots, 8\); a
  repeated index is summed. Status: STANDARD.
\item
  \(t^A\): The generators of \(SU(3)\) on colour triplets, \(t^A\): half
  the Gell-Mann matrices, with \(\mathrm{tr}(t^At^B) = \delta^{AB}/2\).
  Unit: dimensionless. Status: STANDARD.
\item
  \(g_2\): The \(SU(2)_L\) gauge coupling of the four-dimensional world
  (usually \(g\)): the coupling of the massless mode of the \(SU(2)_L\)
  field of the bulk, \(g_2^2 = 3H\mathsf{g}_2^2/(2V_*)\). Unit:
  dimensionless. Value: \(g_2(m_Z) = 0.6517\). Status: STANDARD.
\item
  \(W^I_M,\ W^I_\mu,\ W^I_{MN}\): The gauge field of \(SU(2)_L\)
  (\(I = 1, 2, 3\)): in the bulk \(W^I_M\) (geometric normalisation), on
  the fold the canonically normalised four-dimensional field
  \(W^I_\mu(x)\) of its massless mode,
  \(g_2W^I_\mu(x) = W^I_\mu(\sigma = 0, x)\); \(W^I_{MN}\),
  \(W^I_{\mu\nu}\) their field strengths;
  \(W^\pm_\mu = (W^1_\mu \mp \mathrm{i}W^2_\mu)/\sqrt2\). Unit: GeV
  (field), GeV\(^2\) (field strength). Status: AUTHOR.
\item
  \(I\): A weak-isospin (adjoint SU(2)) index running over \(1, 2, 3\);
  a repeated index is summed. Status: STANDARD.
\item
  \(\tau^I\): The Pauli matrices acting on weak doublets
  (\(\tau^1, \tau^2, \tau^3\); \(\tau^2\) is the second one, not a
  square). Unit: dimensionless. Status: STANDARD.
\item
  \(g_1\): The \(U(1)_Y\) gauge coupling of the four-dimensional world
  (usually \(g'\); not normalised for grand unification): the coupling
  of the massless mode of the hypercharge field of the bulk,
  \(g_1^2 = 3H\mathsf{g}_1^2/(2V_*)\). Unit: dimensionless. Value:
  \(g_1(m_Z) = 0.3574\). Status: STANDARD.
\item
  \(Y_f\): The hypercharge of a field, the charge of its coupling
  \(g_1Y_fB_\mu\) to the hypercharge field in the covariant derivative
  (\(1/6\), \(2/3\), \(-1/3\), \(-1/2\), \(-1\), \(0\) for
  \(Q, u_{\mathrm{R}}, d_{\mathrm{R}}, \ell, e_{\mathrm{R}}, \nu_{\mathrm{R}}\),
  \(1/2\) for \(\Phi\)). Unit: dimensionless. Status: STANDARD.
\item
  \(f\): The label of a kind of Standard-Model multiplet:
  \(f \in \{Q, u, d, \ell, e, \nu\}\) (quark doublets, up-quark
  singlets, down-quark singlets, lepton doublets, charged-lepton
  singlets, right-handed neutrinos); a sum over \(f\) runs over the six.
  Status: STANDARD.
\item
  \(B_M,\ B_\mu,\ B_{MN}\): The hypercharge gauge field of \(U(1)_Y\):
  in the bulk \(B_M\) (geometric normalisation), on the fold the
  canonically normalised four-dimensional field \(B_\mu(x)\) of its
  massless mode, \(g_1B_\mu(x) = B_\mu(\sigma = 0, x)\); \(B_{MN}\),
  \(B_{\mu\nu}\) their field strengths. Unit: GeV (field), GeV\(^2\)
  (field strength). Status: AUTHOR.
\item
  \(\sigma\): The proper coordinate of the hidden direction:
  \(\mathrm{d}\sigma = \lvert\cot z\rvert\,\mathrm{d}x_0\),
  \(\sigma = \mathrm{sgn}(\pi/2 - z)\ln(\sin z)/(6H)\) on one period
  \(0 < z < \pi\); \(\sigma = 0\) is the fold \(z = \pi/2\),
  \(\sigma \to \mp\infty\) the tips \(z \to 0, \pi\); one period of
  \(x_0\) is the whole line \(-\infty < \sigma < \infty\). Unit: length.
  Status: DERIVED.
\item
  \(x\): The coordinates of the fold: \(x_1, x_2, x_3\) of ordinary
  space, \(x_4\) the time, \(x_5, x_6, x_7\) the extra times; a field of
  the fold is a function of the four observable coordinates
  \(x = (x_1, x_2, x_3, x_4)\), and
  \(\mathrm{d}^4x = \mathrm{d}x_1\,\mathrm{d}x_2\,\mathrm{d}x_3\,\mathrm{d}x_4\).
  Unit: length (each coordinate). Status: AUTHOR.
\end{itemize}

The gauge fields act on the fields of the fold through their values at
the fold, \(\sigma = 0\) (the second line: the field of the bulk in its
geometric normalisation, IN-7, at \(\sigma = 0\) equals the coupling
times the canonically normalised field of four dimensions). For the
massless modes of the bulk, whose profile along \(\sigma\) is 1 (FE-74),
the canonical four-dimensional fields \(G^A_\mu\), \(W^I_\mu\),
\(B_\mu\) have the kinetic terms \(-\frac14 G^A_{\mu\nu}G^{A\mu\nu}\)
and so on, with the four-dimensional couplings
\(g_k^2 = 3H\mathsf{g}_k^2/(2V_*)\) (FE-44): the gauge theory of the
Standard Model on the fold. The massive modes of the bulk enter the same
covariant derivative with their values at the fold (FE-79). The
covariant derivative of a field acts with the part that fits the field:
\(t^A\) on colour triplets, \(\tau^I/2\) on weak doublets, the
hypercharge \(Y_f\) of the field, and the four-dimensional spin
connection \(\hat\omega_\mu\) on spinors. The hypercharges (STANDARD)
are \(Y = 1/6\) for \(Q\), \(2/3\) for \(u_{\mathrm{R}}\), \(-1/3\) for
\(d_{\mathrm{R}}\), \(-1/2\) for \(\ell\), \(-1\) for
\(e_{\mathrm{R}}\), \(0\) for \(\nu_{\mathrm{R}}\) and \(1/2\) for
\(\Phi\). Status: STANDARD (the covariant derivative), DERIVED (the
normalisation of the zero modes, \texttt{um4\_bulk.py}, B2).

\hypertarget{the-quarks-and-leptons}{%
\subsection{The quarks and leptons}\label{the-quarks-and-leptons}}

\[
\begin{aligned}
&\mathcal{L}_{\mathrm{f}} = \sum_{\psi}\mathrm{i}\,\bar\psi\,\hat\gamma^\mu D_\mu\psi,\qquad
 \psi \in \{Q, u_{\mathrm{R}}, d_{\mathrm{R}}, \ell, e_{\mathrm{R}}, \nu_{\mathrm{R}}\},\\
&\{\hat\gamma^\mu, \hat\gamma^\nu\} = -2\hat h^{\mu\nu}\mathbf{1},\qquad \bar\psi = \psi^\dagger\beta,\qquad
 \beta = \hat\gamma^{(4)},\qquad \hat\gamma_5 = \mathrm{i}\,\hat\gamma^{(4)}\hat\gamma^{(1)}\hat\gamma^{(2)}\hat\gamma^{(3)},\\
&\hat P_{\mathrm{L}} = \frac{\mathbf{1} - \hat\gamma_5}{2},\quad \hat P_{\mathrm{R}} = \frac{\mathbf{1} + \hat\gamma_5}{2},\qquad
 \hat P_{\mathrm{L}}Q = Q,\quad \hat P_{\mathrm{L}}\ell = \ell,\quad \hat P_{\mathrm{R}}u_{\mathrm{R}} = u_{\mathrm{R}},\quad
 \hat P_{\mathrm{R}}d_{\mathrm{R}} = d_{\mathrm{R}},\quad \hat P_{\mathrm{R}}e_{\mathrm{R}} = e_{\mathrm{R}},\quad
 \hat P_{\mathrm{R}}\nu_{\mathrm{R}} = \nu_{\mathrm{R}}
\end{aligned} \tag{IN-14}
\] \textbf{Symbols in (IN-14).}

\begin{itemize}
\tightlist
\item
  \(\mathcal{L}_{\mathrm{f}}\): The kinetic Lagrangian of the quarks and
  leptons on the fold. Unit: GeV\(^4\). Status: STANDARD.
\item
  \(\psi\): The label of a Standard-Model fermion field in a sum:
  \(\psi\) runs over
  \(Q, u_{\mathrm{R}}, d_{\mathrm{R}}, \ell, e_{\mathrm{R}}, \nu_{\mathrm{R}}\).
  Status: STANDARD.
\item
  \(\mathrm{i}\): The imaginary unit, \(\mathrm{i}^2 = -1\). Unit:
  dimensionless. Status: STANDARD.
\item
  \(\psi,\ \bar\psi\): A Standard-Model fermion field (a
  four-dimensional Weyl field of the fold with three generations,
  written as a four-component spinor of fixed chirality);
  \(\bar\psi = \psi^\dagger\beta\); \(\psi_f\) the fields of the kind
  \(f\). Unit: GeV\(^{3/2}\). Status: STANDARD.
\item
  \(\hat\gamma^\mu\): The four-dimensional complex \(4 \times 4\) Dirac
  matrices of the fold: \(\hat\gamma^{(a)}\) along the frame directions
  of the fold and \(\hat\gamma^\mu\) along its coordinates (the frame
  matrices contracted with the inverse frame field of the fold),
  \(\{\hat\gamma^\mu, \hat\gamma^\nu\} = -2\hat h^{\mu\nu}\mathbf{1}\);
  \(\hat\gamma_\mu = \hat h_{\mu\nu}\hat\gamma^\nu\). Unit:
  dimensionless. Status: STANDARD.
\item
  \(\mu\): A four-dimensional coordinate index running over the
  observable coordinates \(x_1, x_2, x_3, x_4\) of the fold (signature
  \(+,+,+,-\)); a repeated index is summed. Status: STANDARD.
\item
  \(D_M,\ D_\mu\): The covariant derivative: on dirac16complex
  \(D_M = \partial_M + \Omega_M\) (with
  \(D_M\bar\Psi = \partial_M\bar\Psi - \bar\Psi\Omega_M\)); on the
  fields of the fold the gauge and spin covariant derivative \(D_\mu\)
  (the partial derivative plus the spin connection of the fold and the
  gauge fields times their couplings and generators);
  \(D^\mu = \hat h^{\mu\nu}D_\nu\); on an adjoint field
  \((D_\mu X)^A = \partial_\mu X^A + g_3f^{ABC}G^B_\mu X^C\) (likewise
  for \(SU(2)\) with \(g_2\epsilon^{IJK}\)). Unit: 1/length. Status:
  STANDARD.
\item
  \(Q,\ \bar Q\): The left-handed quark doublets of the three
  generations, \(Q = (u_{\mathrm{L}}, d_{\mathrm{L}})\) (colour
  triplets, hypercharge \(1/6\)). Unit: GeV\(^{3/2}\). Status: AUTHOR.
\item
  \(u_{\mathrm{R}}\): The right-handed up-type quarks (up, charm and
  top; singlets of \(SU(2)_L\), hypercharge \(2/3\)). Unit:
  GeV\(^{3/2}\). Status: STANDARD.
\item
  \(d_{\mathrm{R}}\): The right-handed down-type quarks (down, strange
  and bottom; hypercharge \(-1/3\)). Unit: GeV\(^{3/2}\). Status:
  STANDARD.
\item
  \(\ell,\ \bar\ell\): The left-handed lepton doublets of the three
  generations, \(\ell = (\nu_{\mathrm{L}}, e_{\mathrm{L}})\)
  (hypercharge \(-1/2\)). Unit: GeV\(^{3/2}\). Status: AUTHOR.
\item
  \(e_{\mathrm{R}}\): The right-handed electron, muon and tau (singlets,
  hypercharge \(-1\)). Unit: GeV\(^{3/2}\). Status: AUTHOR.
\item
  \(\nu_{\mathrm{R}}\): The three right-handed neutrinos (gauge
  singlets). Unit: GeV\(^{3/2}\). Status: ASSUMED.
\item
  \(\nu\): A four-dimensional coordinate index (as \(\mu\)). Status:
  STANDARD.
\item
  \(\hat h_{\mu\nu}\): The four-dimensional metric of the fold,
  \(\hat h_{\mu\nu} = \mathrm{diag}(\mathrm{e}^{2a_4}, \mathrm{e}^{2a_4}, \mathrm{e}^{2a_4}, -1)\)
  in the coordinates \(x_1, \dots, x_4\) (the induced metric at
  \(\sigma = 0\) without the extra times); \(\hat h^{\mu\nu}\) its
  inverse, \(\hat h\) its determinant,
  \(\sqrt{-\hat h} = \mathrm{e}^{3a_4}\); \(\hat h^{(1)}_{\mu\nu}\) its
  first-order cosmological perturbation. Unit: dimensionless. Status:
  DERIVED.
\item
  \(\mathbf{1},\ \mathbf{1}_4,\ \mathbf{1}_8,\ \mathbf{1}_{16}\): The
  identity matrix of the size fixed by the context (\(16 \times 16\) on
  dirac16complex, \(4 \times 4\) on a four-dimensional spinor,
  \(3 \times 3\) in generation or flavour space); \(\mathbf{1}_4\),
  \(\mathbf{1}_8\) and \(\mathbf{1}_{16}\) are the identity matrices of
  4, 8 and 16 rows. Unit: dimensionless. Status: STANDARD.
\item
  \({}^\dagger\): Hermitian conjugation: the complex conjugate of the
  transpose of a matrix or spinor; for an operator on the space of
  states, its adjoint with respect to the scalar product of that space
  (for the quantised dirac16complex the indefinite Krein scalar
  product). Status: STANDARD.
\item
  \(\beta\): The matrix of the four-dimensional Dirac adjoint,
  \(\beta = \hat\gamma^{(4)}\) (Hermitian, \(\beta^2 = \mathbf{1}\)).
  Unit: dimensionless. Status: STANDARD.
\item
  \(\hat\gamma_5\): The four-dimensional chirality matrix,
  \(\hat\gamma_5 = \mathrm{i}\hat\gamma^{(4)}\hat\gamma^{(1)}\hat\gamma^{(2)}\hat\gamma^{(3)}\),
  \(\hat\gamma_5^2 = \mathbf{1}\). Unit: dimensionless. Status:
  STANDARD.
\item
  \(\hat P_{\mathrm{L}}\): The four-dimensional left-handed projector,
  \(\hat P_{\mathrm{L}} = (\mathbf{1} - \hat\gamma_5)/2\). Unit:
  dimensionless. Status: STANDARD.
\item
  \(\hat P_{\mathrm{R}}\): The four-dimensional right-handed projector,
  \(\hat P_{\mathrm{R}} = (\mathbf{1} + \hat\gamma_5)/2\). Unit:
  dimensionless. Status: STANDARD.
\end{itemize}

Every Standard-Model fermion is a four-dimensional Weyl field on the
fold, written as a four-component spinor with a fixed chirality: the
quark doublets \(Q\) and the lepton doublets \(\ell\) are left-handed,
the singlets \(u_{\mathrm{R}}\), \(d_{\mathrm{R}}\), \(e_{\mathrm{R}}\),
\(\nu_{\mathrm{R}}\) right-handed. Each symbol stands for the three
generations (a vector in generation space) and, for quarks, the three
colours. \(\hat\gamma^\mu = \hat{\mathbf{e}}_a{}^\mu\hat\gamma^{(a)}\)
(\(\hat{\mathbf{e}}_a{}^\mu\) the inverse frame field of the fold) are
complex \(4 \times 4\) Dirac matrices of the fold (with the frame
matrices \(\hat\gamma^{(a)}\), \(a = 1, \dots, 4\),
\((\hat\gamma^{(4)})^2 = \mathbf{1}\),
\((\hat\gamma^{(1)})^2 = (\hat\gamma^{(2)})^2 = (\hat\gamma^{(3)})^2 = -\mathbf{1}\));
\(\beta = \hat\gamma^{(4)}\) is Hermitian and
\((\hat\gamma^\mu)^\dagger = \beta\hat\gamma^\mu\beta\). With these
conventions the free Dirac equation is
\((-\mathrm{i}\hat\gamma^\mu\partial_\mu + m)\psi = 0\) and
\(\bar\psi\beta\psi\)-type energy densities are positive. Status:
STANDARD.

\[
\mathcal{L}_{\mathrm{Y}} = -\big(\bar Q\,\mathcal{Y}^{u}\,\tilde\Phi\,u_{\mathrm{R}} + \bar Q\,\mathcal{Y}^{d}\,\Phi\,d_{\mathrm{R}}
 + \bar\ell\,\mathcal{Y}^{e}\,\Phi\,e_{\mathrm{R}} + \bar\ell\,\mathcal{Y}^{\nu}\,\tilde\Phi\,\nu_{\mathrm{R}}\big)
 + \mathrm{h.c.},\qquad \tilde\Phi = \mathrm{i}\tau^2\Phi^* \tag{IN-15}
\] \textbf{Symbols in (IN-15).}

\begin{itemize}
\tightlist
\item
  \(\mathcal{L}_{\mathrm{Y}}\): The Yukawa Lagrangian of the quarks and
  leptons with the Higgs doublet. Unit: GeV\(^4\). Status: STANDARD.
\item
  \(Q,\ \bar Q\): The left-handed quark doublets of the three
  generations, \(Q = (u_{\mathrm{L}}, d_{\mathrm{L}})\) (colour
  triplets, hypercharge \(1/6\)). Unit: GeV\(^{3/2}\). Status: AUTHOR.
\item
  \(\mathcal{Y}^{u},\ \mathcal{Y}^{d},\ \mathcal{Y}^{e},\ \mathcal{Y}^{\nu}\):
  The Yukawa matrices (\(3 \times 3\) complex matrices in generation
  space) of the up quarks, the down quarks, the charged leptons and the
  neutrinos. Unit: dimensionless. Status: STANDARD.
\item
  \(\tilde\Phi\): The conjugate Higgs doublet,
  \(\tilde\Phi = \mathrm{i}\tau^2\Phi^*\) (hypercharge \(-1/2\)). Unit:
  GeV. Status: STANDARD.
\item
  \(u_{\mathrm{R}}\): The right-handed up-type quarks (up, charm and
  top; singlets of \(SU(2)_L\), hypercharge \(2/3\)). Unit:
  GeV\(^{3/2}\). Status: STANDARD.
\item
  \(\Phi\): The Higgs doublet (two complex components, hypercharge
  \(1/2\)), a field of the fold. Unit: GeV. Status: AUTHOR.
\item
  \(d_{\mathrm{R}}\): The right-handed down-type quarks (down, strange
  and bottom; hypercharge \(-1/3\)). Unit: GeV\(^{3/2}\). Status:
  STANDARD.
\item
  \(\ell,\ \bar\ell\): The left-handed lepton doublets of the three
  generations, \(\ell = (\nu_{\mathrm{L}}, e_{\mathrm{L}})\)
  (hypercharge \(-1/2\)). Unit: GeV\(^{3/2}\). Status: AUTHOR.
\item
  \(e_{\mathrm{R}}\): The right-handed electron, muon and tau (singlets,
  hypercharge \(-1\)). Unit: GeV\(^{3/2}\). Status: AUTHOR.
\item
  \(\nu_{\mathrm{R}}\): The three right-handed neutrinos (gauge
  singlets). Unit: GeV\(^{3/2}\). Status: ASSUMED.
\item
  \(\mathrm{h.c.}\): The Hermitian conjugate of the preceding terms.
  Status: STANDARD.
\item
  \(\mathrm{i}\): The imaginary unit, \(\mathrm{i}^2 = -1\). Unit:
  dimensionless. Status: STANDARD.
\item
  \(\tau^I\): The Pauli matrices acting on weak doublets
  (\(\tau^1, \tau^2, \tau^3\); \(\tau^2\) is the second one, not a
  square). Unit: dimensionless. Status: STANDARD.
\item
  \({}^*\): Complex conjugation of a number, or of each entry of a
  matrix or column. Status: STANDARD.
\end{itemize}

The Yukawa couplings, with the complex \(3 \times 3\) matrices
\(\mathcal{Y}^{u}\), \(\mathcal{Y}^{d}\), \(\mathcal{Y}^{e}\),
\(\mathcal{Y}^{\nu}\) in generation space; \(\tilde\Phi\) is the
conjugate doublet (hypercharge \(-1/2\)). In the vacuum they give the
quarks, the charged leptons and the neutrinos their Dirac masses
(FE-46); with three right-handed neutrinos the neutrino masses are Dirac
masses and the lepton mixing is that of the PMNS matrix. No Majorana
mass of \(\nu_{\mathrm{R}}\) is included (lepton number is conserved).
Status: STANDARD (with \(\nu_{\mathrm{R}}\): accepted by the author,
2026-10-24 17:32).

\hypertarget{the-portals-of-dirac16complex}{%
\subsection{The portals of
dirac16complex}\label{the-portals-of-dirac16complex}}

\[
\begin{aligned}
\mathcal{L}_{\mathrm{P}} = {}&-c_\Phi\,\big(\Phi^\dagger\Phi\big)\,\bar\Psi P_{\mathrm{L}}\Psi
 - \sum_{X}\Big[\frac14\big(c^{G}_X\,G^A_{\mu\nu}G^{A\mu\nu} + c^{W}_X\,W^I_{\mu\nu}W^{I\mu\nu}
 + c^{B}_X\,B_{\mu\nu}B^{\mu\nu}\big)\,\bar\Psi P_X\Psi\\
&\qquad + \frac{d_X}{2}\,B_{\mu\nu}\,\mathcal{D}_X^{\mu\nu} + \sum_{f} c^{f}_X\,J_f^\mu\,j_{X\mu}\Big],\\
\mathcal{D}_X^{\mu\nu} = {}&\mathrm{i}\,\bar\Psi\,\Sigma^{\mu\nu}P_X\Psi,\qquad
j_X^\mu = \mathrm{i}\big(\bar\Psi P_X D^\mu\Psi - (D^\mu\bar\Psi)P_X\Psi\big),\qquad
J_f^\mu = \bar\psi_f\,\hat\gamma^\mu\psi_f
\end{aligned} \tag{IN-16}
\] \textbf{Symbols in (IN-16).}

\begin{itemize}
\tightlist
\item
  \(\mathcal{L}_{\mathrm{P}}\): The Lagrangian of the portals of
  dirac16complex at the fold. Unit: GeV\(^4\). Status: INVENTED.
\item
  \(c_\Phi\): The Higgs-portal coupling of the left-handed half of
  dirac16complex (the right-handed half has none). Unit: GeV\(^{-5}\).
  Status: INVENTED.
\item
  \(\Phi\): The Higgs doublet (two complex components, hypercharge
  \(1/2\)), a field of the fold. Unit: GeV. Status: AUTHOR.
\item
  \({}^\dagger\): Hermitian conjugation: the complex conjugate of the
  transpose of a matrix or spinor; for an operator on the space of
  states, its adjoint with respect to the scalar product of that space
  (for the quantised dirac16complex the indefinite Krein scalar
  product). Status: STANDARD.
\item
  \(\Psi,\ \bar\Psi\): dirac16complex: the author's 16-component complex
  spinor field of eight dimensions;
  \(\bar\Psi = \Psi^\dagger\sigma_{16}\) is its Dirac adjoint; in a term
  of the fold, \(\Psi\) is its value at \(\sigma = 0\). Unit:
  GeV\(^{7/2}\). Status: AUTHOR.
\item
  \(P_{\mathrm{L}}\): The author's left-handed chiral projector,
  \(P_{\mathrm{L}} = (\mathbf{1} - \Gamma_8)/2\). Unit: dimensionless.
  Status: AUTHOR.
\item
  \(X\): The label of a half of dirac16complex:
  \(X \in \{\mathrm{L}, \mathrm{R}\}\) (left-handed or right-handed
  half); a sum over \(X\) runs over both. Status: STANDARD.
\item
  \(c^{G}_X\): The gluon-kinetic portal coupling of the half \(X\) of
  dirac16complex. Unit: GeV\(^{-7}\). Status: INVENTED.
\item
  \(G^A_M,\ G^A_\mu,\ G^A_{MN}\): The gluon field of \(SU(3)_C\)
  (\(A = 1, \dots, 8\)): in the bulk \(G^A_M\), a field of eight
  dimensions in its geometric normalisation (the covariant derivative
  \(\partial - \mathrm{i}G\)), and on the fold the canonically
  normalised four-dimensional field \(G^A_\mu(x)\) of its massless mode,
  \(g_3G^A_\mu(x) = G^A_\mu(\sigma = 0, x)\); \(G^A_{MN}\) and
  \(G^A_{\mu\nu}\) their field strengths. Unit: GeV (field), GeV\(^2\)
  (field strength). Status: AUTHOR.
\item
  \(A\): A colour (adjoint SU(3)) index running over \(1, \dots, 8\); a
  repeated index is summed. Status: STANDARD.
\item
  \(\mu\): A four-dimensional coordinate index running over the
  observable coordinates \(x_1, x_2, x_3, x_4\) of the fold (signature
  \(+,+,+,-\)); a repeated index is summed. Status: STANDARD.
\item
  \(\nu\): A four-dimensional coordinate index (as \(\mu\)). Status:
  STANDARD.
\item
  \(c^{W}_X\): The \(W\)-kinetic portal coupling of the half \(X\).
  Unit: GeV\(^{-7}\). Status: INVENTED.
\item
  \(W^I_M,\ W^I_\mu,\ W^I_{MN}\): The gauge field of \(SU(2)_L\)
  (\(I = 1, 2, 3\)): in the bulk \(W^I_M\) (geometric normalisation), on
  the fold the canonically normalised four-dimensional field
  \(W^I_\mu(x)\) of its massless mode,
  \(g_2W^I_\mu(x) = W^I_\mu(\sigma = 0, x)\); \(W^I_{MN}\),
  \(W^I_{\mu\nu}\) their field strengths;
  \(W^\pm_\mu = (W^1_\mu \mp \mathrm{i}W^2_\mu)/\sqrt2\). Unit: GeV
  (field), GeV\(^2\) (field strength). Status: AUTHOR.
\item
  \(I\): A weak-isospin (adjoint SU(2)) index running over \(1, 2, 3\);
  a repeated index is summed. Status: STANDARD.
\item
  \(c^{B}_X\): The \(B\)-kinetic portal coupling of the half \(X\).
  Unit: GeV\(^{-7}\). Status: INVENTED.
\item
  \(B_M,\ B_\mu,\ B_{MN}\): The hypercharge gauge field of \(U(1)_Y\):
  in the bulk \(B_M\) (geometric normalisation), on the fold the
  canonically normalised four-dimensional field \(B_\mu(x)\) of its
  massless mode, \(g_1B_\mu(x) = B_\mu(\sigma = 0, x)\); \(B_{MN}\),
  \(B_{\mu\nu}\) their field strengths. Unit: GeV (field), GeV\(^2\)
  (field strength). Status: AUTHOR.
\item
  \(P_X\): The chiral projector of the half \(X\):
  \(P_X = P_{\mathrm{L}}\) for \(X = \mathrm{L}\), \(P_{\mathrm{R}}\)
  for \(X = \mathrm{R}\). Unit: dimensionless. Status: AUTHOR.
\item
  \(d_X\): The dipole-portal coupling of the half \(X\) of
  dirac16complex with the hypercharge field strength. Unit:
  GeV\(^{-5}\). Status: INVENTED.
\item
  \(\mathcal{D}_X^{\mu\nu}\): The dipole (antisymmetric tensor) density
  of the half \(X\) of dirac16complex at the fold,
  \(\mathcal{D}_X^{\mu\nu} = \mathrm{i}\bar\Psi\Sigma^{\mu\nu}P_X\Psi\)
  (real, DERIVED). Unit: GeV\(^7\). Status: DERIVED.
\item
  \(f\): The label of a kind of Standard-Model multiplet:
  \(f \in \{Q, u, d, \ell, e, \nu\}\) (quark doublets, up-quark
  singlets, down-quark singlets, lepton doublets, charged-lepton
  singlets, right-handed neutrinos); a sum over \(f\) runs over the six.
  Status: STANDARD.
\item
  \(c^{f}_X\): The current-portal coupling of the half \(X\) of
  dirac16complex with the fermion-number current of the multiplets of
  kind \(f\) (flavour blind); \(c^{\ell}_X\) and \(c^{\nu}_X\) are those
  of the lepton doublets and of the right-handed neutrinos. Unit:
  GeV\(^{-7}\). Status: INVENTED.
\item
  \(J_f^\mu\): The fermion-number current of the Standard-Model
  multiplets of kind \(f\), \(J_f^\mu = \bar\psi_f\hat\gamma^\mu\psi_f\)
  summed over the three generations and the gauge components; with other
  labels the currents that are the sources of the gauge fields. Unit:
  GeV\(^3\). Status: STANDARD.
\item
  \(j_X^\mu\): The convective current of the half \(X\) of
  dirac16complex,
  \(j_X^\mu = \mathrm{i}(\bar\Psi P_XD^\mu\Psi - (D^\mu\bar\Psi)P_X\Psi)\)
  (real and joining the half with itself, DERIVED). Unit: GeV\(^8\).
  Status: DERIVED.
\item
  \(\mathrm{i}\): The imaginary unit, \(\mathrm{i}^2 = -1\). Unit:
  dimensionless. Status: STANDARD.
\item
  \(\Sigma^{ab},\ \Sigma^{\mu\nu}\): The generators of Spin(4,4) on
  dirac16complex, \(\Sigma^{ab} = \frac14[T^a, T^b]\) (the scaled
  commutators of the author's matrices);
  \(\Sigma^{\mu\nu} = \mathbf{e}_a{}^\mu\mathbf{e}_b{}^\nu\Sigma^{ab}\)
  along the fold. Unit: dimensionless. Status: AUTHOR.
\item
  \(D_M,\ D_\mu\): The covariant derivative: on dirac16complex
  \(D_M = \partial_M + \Omega_M\) (with
  \(D_M\bar\Psi = \partial_M\bar\Psi - \bar\Psi\Omega_M\)); on the
  fields of the fold the gauge and spin covariant derivative \(D_\mu\)
  (the partial derivative plus the spin connection of the fold and the
  gauge fields times their couplings and generators);
  \(D^\mu = \hat h^{\mu\nu}D_\nu\); on an adjoint field
  \((D_\mu X)^A = \partial_\mu X^A + g_3f^{ABC}G^B_\mu X^C\) (likewise
  for \(SU(2)\) with \(g_2\epsilon^{IJK}\)). Unit: 1/length. Status:
  STANDARD.
\item
  \(\psi,\ \bar\psi\): A Standard-Model fermion field (a
  four-dimensional Weyl field of the fold with three generations,
  written as a four-component spinor of fixed chirality);
  \(\bar\psi = \psi^\dagger\beta\); \(\psi_f\) the fields of the kind
  \(f\). Unit: GeV\(^{3/2}\). Status: STANDARD.
\item
  \(\hat\gamma^\mu\): The four-dimensional complex \(4 \times 4\) Dirac
  matrices of the fold: \(\hat\gamma^{(a)}\) along the frame directions
  of the fold and \(\hat\gamma^\mu\) along its coordinates (the frame
  matrices contracted with the inverse frame field of the fold),
  \(\{\hat\gamma^\mu, \hat\gamma^\nu\} = -2\hat h^{\mu\nu}\mathbf{1}\);
  \(\hat\gamma_\mu = \hat h_{\mu\nu}\hat\gamma^\nu\). Unit:
  dimensionless. Status: STANDARD.
\end{itemize}

The new interactions of UM-4 (INVENTED; the same as in UM-3), the order
of the author that the left-handed and the right-handed halves of
dirac16complex each couple to every Standard-Model field with its own
coupling, and that only the left-handed half couples to the Higgs
doublet. \(\Psi\) is the value of dirac16complex at the fold, and the
field strengths are those of the four-dimensional gauge fields at the
fold (IN-13). Every term is a product of a gauge-invariant
Standard-Model operator and a Lorentz-tensor bilinear of one half \(X\)
of dirac16complex, so every term is invariant under
\(SU(3)_C \times SU(2)_L \times U(1)_Y\) and under the Lorentz
transformations of the fold; none changes the flavour of a quark or a
lepton (DERIVED by construction):

\begin{itemize}
\tightlist
\item
  the Higgs portal: the scalar density \(\bar\Psi P_{\mathrm{L}}\Psi\)
  of the left-handed half times \(\Phi^\dagger\Phi\), with the coupling
  \(c_\Phi\); the right-handed half has no Higgs coupling;
\item
  the gauge-kinetic portals: the scalar density of each half times the
  gauge-invariant kinetic densities of the gluons, of the \(W\) fields
  and of the \(B\) field, with the six couplings \(c^{G}_X\),
  \(c^{W}_X\), \(c^{B}_X\): they make the gauge couplings at the fold
  depend on the density of dirac16complex at the fold and couple the
  particles \(\chi\) to gluons, photons and weak bosons;
\item
  the dipole portal: the antisymmetric tensor \(\mathcal{D}_X^{\mu\nu}\)
  of each half times the hypercharge field strength (the only
  Standard-Model field strength that is gauge invariant by itself), with
  the couplings \(d_X\): an electric and magnetic dipole moment of
  dirac16complex with respect to the photon and the Z boson;
\item
  the current portals: the convective current \(j_X^\mu\) of each half
  (the only real vector current that joins a half with itself:
  \(\bar\Psi\gamma^\mu P_X\Psi\) joins \(\Psi_{\mathrm{L}}\) with
  \(\Psi_{\mathrm{R}}\) and is not real, DERIVED,
  \texttt{um4\_fields.py}, F1) times the fermion-number current
  \(J_f^\mu\) of each kind \(f\) of Standard-Model multiplet, with the
  twelve couplings \(c^{f}_X\); \(J_f^\mu\) sums over the three
  generations and over the gauge components, so the portal is flavour
  blind.
\end{itemize}

The 21 couplings are \(c_\Phi\), \(c^{G}_{\mathrm{L}}\),
\(c^{G}_{\mathrm{R}}\), \(c^{W}_{\mathrm{L}}\), \(c^{W}_{\mathrm{R}}\),
\(c^{B}_{\mathrm{L}}\), \(c^{B}_{\mathrm{R}}\), \(d_{\mathrm{L}}\),
\(d_{\mathrm{R}}\) and \(c^{f}_{\mathrm{L}}\), \(c^{f}_{\mathrm{R}}\)
for \(f = Q, u, d, \ell, e, \nu\). Every term is bilinear in \(\Psi\):
the action is invariant under a change of the phase of \(\Psi\) by a
constant (the charge of dirac16complex is conserved), and baryon and
lepton numbers are conserved. Status: INVENTED.

\hypertarget{gravity-1}{%
\section{Gravity}\label{gravity-1}}

\hypertarget{the-einstein-equations-in-the-bulk}{%
\subsection{The Einstein equations in the
bulk}\label{the-einstein-equations-in-the-bulk}}

\[
\mathcal{E}_{(1)}{}^M{}_N + \Lambda\,\delta^M_N = \kappa\,\big(\mathcal{T}_\Psi{}^M{}_N + \mathcal{T}_{\mathrm{gauge}}{}^M{}_N\big)
 \qquad (\sigma \ne 0) \tag{FE-1}
\] \textbf{Symbols in (FE-1).}

\begin{itemize}
\tightlist
\item
  \(\mathcal{E}_{(p)}{}^M{}_N\): The Lovelock tensor of order \(p\),
  \(\mathcal{E}_{(p)}{}^M{}_N = -2^{-(p+1)}\,\delta^{MP_1\cdots P_{2p}}_{NQ_1\cdots Q_{2p}}\,R^{Q_1Q_2}{}_{P_1P_2}\cdots R^{Q_{2p-1}Q_{2p}}{}_{P_{2p-1}P_{2p}}\)
  (with both indices down it is \(1/\sqrt{\lvert g\rvert}\) times the
  variation of \(\sqrt{\lvert g\rvert}\,\mathcal{L}_{(p)}\) with respect
  to \(g^{MN}\); \(\mathcal{E}_{(1)}\) the Einstein tensor,
  \(\mathcal{E}_{(2)}\) the Gauss-Bonnet tensor, \(\mathcal{E}_{(3)}\)
  the cubic Lovelock tensor); with frame indices its frame components.
  Unit: 1/length\(^{2p}\). Status: STANDARD.
\item
  \(p\): The Lovelock order \(p = 1, 2, 3\) (Einstein, Gauss-Bonnet,
  cubic), as an index and as a number (UM-4 keeps only \(p = 1\)); in
  the Kohn-Sham equations also a second label of a Kohn-Sham orbital,
  with the Kronecker delta of the two labels. Status: STANDARD.
\item
  \(M\): A coordinate index of eight dimensions (also with a numbered
  copy \(M_1, M_2, \dots\)): it runs over \(x_0, x_1, \dots, x_7\); a
  repeated index is summed. Status: STANDARD.
\item
  \(N\): A coordinate index of eight dimensions (also
  \(N_1, N_2, \dots\)), running over \(x_0, \dots, x_7\); a repeated
  index is summed. Status: STANDARD.
\item
  \(\Lambda\): The eight-dimensional cosmological constant of UM-4,
  \(\Lambda = -21H^2\): the value for which the static block is a vacuum
  of Einstein gravity (its Einstein tensor is \(21H^2\) in every
  direction). The author writes the pure number \(\Lambda_{\mathrm{N}}\)
  in
  \(0 = -\Lambda_{\mathrm{N}} + \sum_p H^{-2p}w_{(p)}\mathcal{E}_{(p)}\),
  so that \(\Lambda = -\Lambda_{\mathrm{N}}H^2\). Unit: GeV\(^2\).
  Value: \(-21H^2\); the author's original \(-10H^2\); UM-3: \(-7H^2\).
  Status: AUTHOR (the term), DERIVED (the value).
\item
  \(\delta\): The Kronecker delta: with two indices \(\delta^M_N = 1\)
  if \(M = N\) and 0 otherwise (likewise \(\delta_{X\mathrm{L}}\) is 1
  if the label \(X\) of a half of dirac16complex is L and 0 otherwise);
  with two index strings of equal length the generalised Kronecker delta
  \(\delta^{M_1M_2\cdots}_{N_1N_2\cdots}\), the determinant of the
  square matrix of the Kronecker deltas \(\delta^{M_a}_{N_b}\) (\(+1\)
  or \(-1\) if the upper indices are an even or odd permutation of the
  lower ones, 0 otherwise). Status: STANDARD.
\item
  \(\kappa\): The eight-dimensional gravitational coupling,
  \(\kappa = 8\pi G_8\). Unit: GeV\(^{-6}\). Status: AUTHOR.
\item
  \(\mathcal{T}_\Psi\): The energy-momentum tensor of dirac16complex:
  \(-2/\sqrt{\lvert g\rvert}\) times the variation of the action
  \(S_\Psi\) of dirac16complex with respect to \(g^{MN}\). Unit:
  GeV\(^8\). Status: STANDARD.
\item
  \(\mathcal{T}_{\mathrm{gauge}}\): The energy-momentum tensor of the
  gauge fields of the bulk. Unit: GeV\(^8\). Status: STANDARD.
\item
  \(\sigma\): The proper coordinate of the hidden direction:
  \(\mathrm{d}\sigma = \lvert\cot z\rvert\,\mathrm{d}x_0\),
  \(\sigma = \mathrm{sgn}(\pi/2 - z)\ln(\sin z)/(6H)\) on one period
  \(0 < z < \pi\); \(\sigma = 0\) is the fold \(z = \pi/2\),
  \(\sigma \to \mp\infty\) the tips \(z \to 0, \pi\); one period of
  \(x_0\) is the whole line \(-\infty < \sigma < \infty\). Unit: length.
  Status: DERIVED.
\end{itemize}

The variation of \(S_{\mathrm{tot}}\) with respect to the metric, away
from the fold. The left-hand side is the Einstein tensor
\(\mathcal{E}_{(1)}{}^M{}_N\) (FE-2) plus the cosmological constant
\(\Lambda = -21H^2\) (IN-2); the right-hand side is \(\kappa\) times the
energy-momentum tensors of the fields of the bulk, dirac16complex (FE-6)
and the gauge fields (FE-8). With the energy-momentum tensor defined by
\(\mathcal{T}_{MN} = -(2/\sqrt{\lvert g\rvert})\,\delta S_{\mathrm{matter}}/\delta g^{MN}\)
(\(S_{\mathrm{matter}}\) the action of the field), the energy density
seen by an observer at rest is \(-\mathcal{T}^{x_4}{}_{x_4}\), which
equals the time-time component \(\mathcal{T}_{x_4x_4}\) on the fold
(where \(g_{x_4x_4} = -1\); at the depth \(\sigma\) of the bulk
\(\mathcal{T}_{x_4x_4} = -\mathrm{e}^{-2H\lvert\sigma\rvert}\mathcal{T}^{x_4}{}_{x_4}\)
for the static block). Status: AUTHOR (the theory, with the weights of
the order), STANDARD (the variation).

\[
\begin{aligned}
\mathcal{E}_{(p)}{}^M{}_N &= -\frac{1}{2^{p+1}}\,\delta^{MM_1N_1\cdots M_pN_p}_{NP_1Q_1\cdots P_pQ_p}\,
 R^{P_1Q_1}{}_{M_1N_1}\cdots R^{P_pQ_p}{}_{M_pN_p},\\
\mathcal{E}_{(1)}{}^M{}_N &= R^M{}_N - \frac12\,\delta^M_N\,R,\\
\mathcal{E}_{(2)}{}^M{}_N &= 2\big(R\,R^M{}_N - 2R^M{}_PR^P{}_N - 2R^{MP}{}_{NQ}R^Q{}_P + R^{MPQS}R_{NPQS}\big)
 - \frac12\,\delta^M_N\,\mathcal{L}_{(2)}
\end{aligned} \tag{FE-2}
\] \textbf{Symbols in (FE-2).}

\begin{itemize}
\tightlist
\item
  \(\mathcal{E}_{(p)}{}^M{}_N\): The Lovelock tensor of order \(p\),
  \(\mathcal{E}_{(p)}{}^M{}_N = -2^{-(p+1)}\,\delta^{MP_1\cdots P_{2p}}_{NQ_1\cdots Q_{2p}}\,R^{Q_1Q_2}{}_{P_1P_2}\cdots R^{Q_{2p-1}Q_{2p}}{}_{P_{2p-1}P_{2p}}\)
  (with both indices down it is \(1/\sqrt{\lvert g\rvert}\) times the
  variation of \(\sqrt{\lvert g\rvert}\,\mathcal{L}_{(p)}\) with respect
  to \(g^{MN}\); \(\mathcal{E}_{(1)}\) the Einstein tensor,
  \(\mathcal{E}_{(2)}\) the Gauss-Bonnet tensor, \(\mathcal{E}_{(3)}\)
  the cubic Lovelock tensor); with frame indices its frame components.
  Unit: 1/length\(^{2p}\). Status: STANDARD.
\item
  \(p\): The Lovelock order \(p = 1, 2, 3\) (Einstein, Gauss-Bonnet,
  cubic), as an index and as a number (UM-4 keeps only \(p = 1\)); in
  the Kohn-Sham equations also a second label of a Kohn-Sham orbital,
  with the Kronecker delta of the two labels. Status: STANDARD.
\item
  \(M\): A coordinate index of eight dimensions (also with a numbered
  copy \(M_1, M_2, \dots\)): it runs over \(x_0, x_1, \dots, x_7\); a
  repeated index is summed. Status: STANDARD.
\item
  \(N\): A coordinate index of eight dimensions (also
  \(N_1, N_2, \dots\)), running over \(x_0, \dots, x_7\); a repeated
  index is summed. Status: STANDARD.
\item
  \(\delta\): The Kronecker delta: with two indices \(\delta^M_N = 1\)
  if \(M = N\) and 0 otherwise (likewise \(\delta_{X\mathrm{L}}\) is 1
  if the label \(X\) of a half of dirac16complex is L and 0 otherwise);
  with two index strings of equal length the generalised Kronecker delta
  \(\delta^{M_1M_2\cdots}_{N_1N_2\cdots}\), the determinant of the
  square matrix of the Kronecker deltas \(\delta^{M_a}_{N_b}\) (\(+1\)
  or \(-1\) if the upper indices are an even or odd permutation of the
  lower ones, 0 otherwise). Status: STANDARD.
\item
  \(P\): A coordinate index of eight dimensions (also
  \(P_1, P_2, \dots\)), running over \(x_0, \dots, x_7\); a repeated
  index is summed. Status: STANDARD.
\item
  \(Q\): A coordinate index of eight dimensions (also
  \(Q_1, Q_2, \dots\)), running over \(x_0, \dots, x_7\); a repeated
  index is summed. Status: STANDARD.
\item
  \(R^{MN}{}_{PQ},\ R_{MN},\ R\): The curvature of the metric
  \(g_{MN}\): with four indices the Riemann tensor \(R^M{}_{NPQ}\) (in
  the convention in which the commutator of two covariant derivatives,
  acting on a vector, is the Riemann tensor, with its last two indices
  those of the two derivatives, contracted with the vector), with two
  indices the Ricci tensor \(R_{MN} = R^P{}_{MPN}\), without indices the
  scalar curvature \(R = g^{MN}R_{MN}\); with frame indices
  \(R^{(a)(b)}{}_{(c)(d)}\) its frame components,
  \(R^{(a)(b)}{}_{(a)(b)}\) (no sum) the sectional curvature of the
  frame plane \((a, b)\). Unit: 1/length\(^2\). Status: STANDARD.
\item
  \(S\): A coordinate index of eight dimensions, running over
  \(x_0, \dots, x_7\); a repeated index is summed. Status: STANDARD.
\item
  \(\mathcal{L}_{(p)}\): The Lovelock density of order \(p\),
  \(\mathcal{L}_{(p)} = 2^{-p}\,\delta^{P_1\cdots P_{2p}}_{Q_1\cdots Q_{2p}}\,R^{Q_1Q_2}{}_{P_1P_2}\cdots R^{Q_{2p-1}Q_{2p}}{}_{P_{2p-1}P_{2p}}\),
  a scalar built from \(p\) curvature tensors: \(\mathcal{L}_{(1)} = R\)
  (Einstein-Hilbert), \(\mathcal{L}_{(2)}\) the Gauss-Bonnet density,
  \(\mathcal{L}_{(3)}\) the cubic Lovelock density. Unit:
  1/length\(^{2p}\). Status: STANDARD.
\end{itemize}

The Lovelock tensor of order \(p\): the variation of
\(\sqrt{\lvert g\rvert}\,\mathcal{L}_{(p)}\) with respect to the metric.
For \(p = 1\) it is the Einstein tensor (the second line), which is the
left-hand side of FE-1; for \(p = 2\) the Gauss-Bonnet tensor (written
out in the third line); the cubic tensor \(\mathcal{E}_{(3)}\) is the
first line with \(p = 3\). The tensors of order 2 and 3 enter only the
comparison FE-68 to FE-70. UM-4 evaluates the first line exactly by
enumerating the non-zero frame components of the curvature
(\texttt{um4\_lovelock.py}); the second line is checked independently:
the Ricci tensor taken directly from the Christoffel symbols gives the
same Einstein tensor in all 64 components on both sides of the fold
(DERIVED, \texttt{um4\_geometry.py}, G5). Status: STANDARD.

\[
\mathcal{E}_{(p)}{}^M{}_M = -\frac{8 - 2p}{2}\,\mathcal{L}_{(p)},\qquad \nabla_M\,\mathcal{E}_{(p)}{}^M{}_N = 0 \tag{FE-3}
\] \textbf{Symbols in (FE-3).}

\begin{itemize}
\tightlist
\item
  \(\mathcal{E}_{(p)}{}^M{}_N\): The Lovelock tensor of order \(p\),
  \(\mathcal{E}_{(p)}{}^M{}_N = -2^{-(p+1)}\,\delta^{MP_1\cdots P_{2p}}_{NQ_1\cdots Q_{2p}}\,R^{Q_1Q_2}{}_{P_1P_2}\cdots R^{Q_{2p-1}Q_{2p}}{}_{P_{2p-1}P_{2p}}\)
  (with both indices down it is \(1/\sqrt{\lvert g\rvert}\) times the
  variation of \(\sqrt{\lvert g\rvert}\,\mathcal{L}_{(p)}\) with respect
  to \(g^{MN}\); \(\mathcal{E}_{(1)}\) the Einstein tensor,
  \(\mathcal{E}_{(2)}\) the Gauss-Bonnet tensor, \(\mathcal{E}_{(3)}\)
  the cubic Lovelock tensor); with frame indices its frame components.
  Unit: 1/length\(^{2p}\). Status: STANDARD.
\item
  \(p\): The Lovelock order \(p = 1, 2, 3\) (Einstein, Gauss-Bonnet,
  cubic), as an index and as a number (UM-4 keeps only \(p = 1\)); in
  the Kohn-Sham equations also a second label of a Kohn-Sham orbital,
  with the Kronecker delta of the two labels. Status: STANDARD.
\item
  \(M\): A coordinate index of eight dimensions (also with a numbered
  copy \(M_1, M_2, \dots\)): it runs over \(x_0, x_1, \dots, x_7\); a
  repeated index is summed. Status: STANDARD.
\item
  \(N\): A coordinate index of eight dimensions (also
  \(N_1, N_2, \dots\)), running over \(x_0, \dots, x_7\); a repeated
  index is summed. Status: STANDARD.
\item
  \(\mathcal{L}_{(p)}\): The Lovelock density of order \(p\),
  \(\mathcal{L}_{(p)} = 2^{-p}\,\delta^{P_1\cdots P_{2p}}_{Q_1\cdots Q_{2p}}\,R^{Q_1Q_2}{}_{P_1P_2}\cdots R^{Q_{2p-1}Q_{2p}}{}_{P_{2p-1}P_{2p}}\),
  a scalar built from \(p\) curvature tensors: \(\mathcal{L}_{(1)} = R\)
  (Einstein-Hilbert), \(\mathcal{L}_{(2)}\) the Gauss-Bonnet density,
  \(\mathcal{L}_{(3)}\) the cubic Lovelock density. Unit:
  1/length\(^{2p}\). Status: STANDARD.
\item
  \(\nabla\): The covariant derivative of a tensor with the Levi-Civita
  connection of the metric \(g_{MN}\) (Christoffel symbols
  \(\Gamma^P{}_{MN}\)). Status: STANDARD.
\end{itemize}

Two identities of every Lovelock tensor in eight dimensions: its trace
is a multiple of the Lovelock density, and its covariant divergence
vanishes (the contracted Bianchi identity). The second makes the
energy-momentum tensor of the sources conserved (FE-11). Status:
STANDARD.

\hypertarget{the-junction-condition-at-the-fold}{%
\subsection{The junction condition at the
fold}\label{the-junction-condition-at-the-fold}}

\[
\kappa\,\mathcal{S}^a{}_b = \lim_{\varepsilon\to0}\int_{-\varepsilon}^{\varepsilon}\mathrm{d}\sigma\,\mathcal{E}_{(1)}{}^a{}_b
 = -\Big(\big(K^a{}_b\big\vert_{0^+} - K^a{}_b\big\vert_{0^-}\big) - \delta^a{}_b\big(K\big\vert_{0^+} - K\big\vert_{0^-}\big)\Big),\qquad
 K = K^c{}_c \tag{FE-4}
\] \textbf{Symbols in (FE-4).}

\begin{itemize}
\tightlist
\item
  \(\kappa\): The eight-dimensional gravitational coupling,
  \(\kappa = 8\pi G_8\). Unit: GeV\(^{-6}\). Status: AUTHOR.
\item
  \(\mathcal{S}^a{}_b\): The thin source of the fold: the
  energy-momentum tensor concentrated on the fold, integrated across it,
  \(\kappa\mathcal{S}^a{}_b = \int_{0^-}^{0^+}\mathrm{d}\sigma\,\mathcal{E}_{(1)}{}^a{}_b\)
  (the junction condition of Israel); \(\mathcal{S}^4{}_4\) is minus its
  energy density, \(\mathcal{S}^1{}_1\) and \(\mathcal{S}^5{}_5\) its
  pressures along ordinary space and along the extra times (frame
  components along the fold, eight-dimensional normalisation: an energy
  per unit volume of the six other directions of the fold). Unit:
  GeV\(^7\). Value: \(\kappa\mathcal{S}^a{}_b = -12H\delta^a{}_b\) for
  every \(a_4\), \(b_4\) of the author's form. Status: DERIVED.
\item
  \(a\): A frame (orthonormal) index running over \(0, 1, \dots, 7\)
  (frame direction \((0)\) along \(\sigma\), \((1), (2), (3)\) along
  ordinary space, \((4)\) along the time, \((5), (6), (7)\) along the
  extra times); raised and lowered with the frame metric
  \(\eta_{ab} = \mathrm{diag}(1, 1, 1, 1, -1, -1, -1, -1)\); on the fold
  (\(a, b \in \{1, \dots, 7\}\)) the frame of the fold; in the tensors
  of the fold written with its coordinates \(x^a\) (its metric
  \(\hat g_{ab}\) and the tensors built from it) a coordinate index of
  the fold, \(a \in \{1, \dots, 7\}\); a repeated index is summed.
  Status: STANDARD.
\item
  \(b\): A frame index (as \(a\)). Status: STANDARD.
\item
  \(\varepsilon\): The half-thickness of a layer around the fold, taken
  to zero. Unit: length. Status: STANDARD.
\item
  \(\mathrm{d}\): The differential: \(\mathrm{d}\) in front of a
  quantity is an infinitesimal change of it (as \(\mathrm{d}\sigma\), a
  small step of \(\sigma\)); the quotient of two such changes is a
  derivative; \(\mathrm{d}\) with an exponent, in front of several
  variables, is their volume element, the product of their
  differentials. Status: STANDARD.
\item
  \(\sigma\): The proper coordinate of the hidden direction:
  \(\mathrm{d}\sigma = \lvert\cot z\rvert\,\mathrm{d}x_0\),
  \(\sigma = \mathrm{sgn}(\pi/2 - z)\ln(\sin z)/(6H)\) on one period
  \(0 < z < \pi\); \(\sigma = 0\) is the fold \(z = \pi/2\),
  \(\sigma \to \mp\infty\) the tips \(z \to 0, \pi\); one period of
  \(x_0\) is the whole line \(-\infty < \sigma < \infty\). Unit: length.
  Status: DERIVED.
\item
  \(\mathcal{E}_{(p)}{}^M{}_N\): The Lovelock tensor of order \(p\),
  \(\mathcal{E}_{(p)}{}^M{}_N = -2^{-(p+1)}\,\delta^{MP_1\cdots P_{2p}}_{NQ_1\cdots Q_{2p}}\,R^{Q_1Q_2}{}_{P_1P_2}\cdots R^{Q_{2p-1}Q_{2p}}{}_{P_{2p-1}P_{2p}}\)
  (with both indices down it is \(1/\sqrt{\lvert g\rvert}\) times the
  variation of \(\sqrt{\lvert g\rvert}\,\mathcal{L}_{(p)}\) with respect
  to \(g^{MN}\); \(\mathcal{E}_{(1)}\) the Einstein tensor,
  \(\mathcal{E}_{(2)}\) the Gauss-Bonnet tensor, \(\mathcal{E}_{(3)}\)
  the cubic Lovelock tensor); with frame indices its frame components.
  Unit: 1/length\(^{2p}\). Status: STANDARD.
\item
  \(p\): The Lovelock order \(p = 1, 2, 3\) (Einstein, Gauss-Bonnet,
  cubic), as an index and as a number (UM-4 keeps only \(p = 1\)); in
  the Kohn-Sham equations also a second label of a Kohn-Sham orbital,
  with the Kronecker delta of the two labels. Status: STANDARD.
\item
  \(K^a{}_b\): The extrinsic curvature of a surface of constant
  \(\sigma\) of the hidden direction,
  \(K^a{}_b = \frac12 g^{ac}\partial_\sigma g_{cb}\) for \(a, b\) along
  the surface (FD-5): \(-H\,\mathrm{sgn}(\sigma)\) along all seven
  directions of the fold for the author's form of the metric of UM-4;
  \(K = K^c{}_c\) its trace. Unit: GeV (1/length). Status: DERIVED.
\item
  \(\delta\): The Kronecker delta: with two indices \(\delta^M_N = 1\)
  if \(M = N\) and 0 otherwise (likewise \(\delta_{X\mathrm{L}}\) is 1
  if the label \(X\) of a half of dirac16complex is L and 0 otherwise);
  with two index strings of equal length the generalised Kronecker delta
  \(\delta^{M_1M_2\cdots}_{N_1N_2\cdots}\), the determinant of the
  square matrix of the Kronecker deltas \(\delta^{M_a}_{N_b}\) (\(+1\)
  or \(-1\) if the upper indices are an even or odd permutation of the
  lower ones, 0 otherwise). Status: STANDARD.
\item
  \(c\): A frame index (as \(a\)). Status: STANDARD.
\end{itemize}

The field equations at the fold: the junction condition of Israel (W.
Israel, Nuovo Cimento B 44 (1966) 1). The metric of UM-4 has a kink at
\(\sigma = 0\): its first \(\sigma\)-derivatives jump at the fold, so
its second derivatives contain \(\delta_{\mathrm{D}}(\sigma)\). The
Einstein tensor is linear in the second \(\sigma\)-derivatives, so FE-1
extended across the fold reads
\(\mathcal{E}_{(1)} + \Lambda = \kappa(\mathcal{T}_\Psi + \mathcal{T}_{\mathrm{gauge}} + \mathcal{S}\,\delta_{\mathrm{D}}(\sigma))\):
the fold carries the thin source \(\mathcal{S}^a{}_b\) (\(a, b\) along
the fold; an energy per unit volume of the six other directions of the
fold, \(\mathcal{S}^4{}_4\) minus its energy density,
\(\mathcal{S}^1{}_1\) and \(\mathcal{S}^5{}_5\) its pressures along
ordinary space and along the extra times), equal to the integral of the
Einstein tensor across the fold, which is the jump of the extrinsic
curvature \(K^a{}_b\) of the slices \(\sigma = \text{const}\) (FE-30).
The \((\sigma, \sigma)\) and \((\sigma, a)\) components contain no
second \(\sigma\)-derivative: they carry no thin part and hold on each
side separately (the constraints of the bulk). Status: STANDARD (DERIVED
for the metric of UM-4: \texttt{um4\_geometry.py}, G7;
\texttt{um4\_fold.py}, F2).

\[
\mathcal{S}^a{}_b = \frac{1}{V_*}\Big(\mathcal{T}_{\mathrm{H}}{}^a{}_b + \mathcal{T}_{\mathrm{f}}{}^a{}_b
 + \mathcal{T}_{\mathrm{P}}{}^a{}_b + \mathcal{T}_{\mathrm{ff}}{}^a{}_b + \mathcal{T}_{\mathrm{IG}}{}^a{}_b\Big),\qquad
\mathcal{T}_{\mathrm{IG}}{}^\mu{}_\nu = -\frac{1}{\kappa_{\mathrm{ind}}}\,\hat G^\mu{}_\nu,\qquad
\mathcal{T}^r{}_r = \mathcal{L}_{\mathrm{fold}} \tag{FE-5}
\] \textbf{Symbols in (FE-5).}

\begin{itemize}
\tightlist
\item
  \(\mathcal{S}^a{}_b\): The thin source of the fold: the
  energy-momentum tensor concentrated on the fold, integrated across it,
  \(\kappa\mathcal{S}^a{}_b = \int_{0^-}^{0^+}\mathrm{d}\sigma\,\mathcal{E}_{(1)}{}^a{}_b\)
  (the junction condition of Israel); \(\mathcal{S}^4{}_4\) is minus its
  energy density, \(\mathcal{S}^1{}_1\) and \(\mathcal{S}^5{}_5\) its
  pressures along ordinary space and along the extra times (frame
  components along the fold, eight-dimensional normalisation: an energy
  per unit volume of the six other directions of the fold). Unit:
  GeV\(^7\). Value: \(\kappa\mathcal{S}^a{}_b = -12H\delta^a{}_b\) for
  every \(a_4\), \(b_4\) of the author's form. Status: DERIVED.
\item
  \(a\): A frame (orthonormal) index running over \(0, 1, \dots, 7\)
  (frame direction \((0)\) along \(\sigma\), \((1), (2), (3)\) along
  ordinary space, \((4)\) along the time, \((5), (6), (7)\) along the
  extra times); raised and lowered with the frame metric
  \(\eta_{ab} = \mathrm{diag}(1, 1, 1, 1, -1, -1, -1, -1)\); on the fold
  (\(a, b \in \{1, \dots, 7\}\)) the frame of the fold; in the tensors
  of the fold written with its coordinates \(x^a\) (its metric
  \(\hat g_{ab}\) and the tensors built from it) a coordinate index of
  the fold, \(a \in \{1, \dots, 7\}\); a repeated index is summed.
  Status: STANDARD.
\item
  \(b\): A frame index (as \(a\)). Status: STANDARD.
\item
  \(V_*\): The physical volume of the extra times \(x_5, x_6, x_7\) on
  the fold in the late universe,
  \(V_* = V_{\mathrm{t}}(\mathrm{e}^{-a_4}b_4)^3\) (with
  \(V_{\mathrm{t}} = \int\mathrm{d}x_5\,\mathrm{d}x_6\,\mathrm{d}x_7\)
  the coordinate volume of the extra times) at a time when
  \(\mathrm{e}^{-a_4}b_4\) is constant; it converts the
  eight-dimensional normalisation into the four-dimensional one. Unit:
  length\(^3\) (GeV\(^{-3}\)). Status: ASSUMED.
\item
  \(\mathcal{T}_{\mathrm{H}}\): The energy-momentum tensor of the Higgs
  doublet (four-dimensional normalisation). Unit: GeV\(^4\). Status:
  STANDARD.
\item
  \(\mathcal{T}_{\mathrm{f}}\): The energy-momentum tensor of the quarks
  and leptons. Unit: GeV\(^4\). Status: STANDARD.
\item
  \(\mathcal{T}_{\mathrm{P}}\): The energy-momentum tensor of the
  portals of dirac16complex. Unit: GeV\(^4\). Status: DERIVED.
\item
  \(\mathcal{T}_{\mathrm{ff}}{}^a{}_b\): The energy-momentum tensor of
  the fold fluid (frame components of the fold, four-dimensional
  normalisation). Unit: GeV\(^4\). Status: ASSUMED.
\item
  \(\mathcal{T}_{\mathrm{IG}}\): The energy-momentum tensor of the
  induced Einstein-Hilbert term of the fold,
  \(\mathcal{T}_{\mathrm{IG}}{}^\mu{}_\nu = -\hat G^\mu{}_\nu/\kappa_{\mathrm{ind}}\).
  Unit: GeV\(^4\). Status: ESTIMATE.
\item
  \(\mu\): A four-dimensional coordinate index running over the
  observable coordinates \(x_1, x_2, x_3, x_4\) of the fold (signature
  \(+,+,+,-\)); a repeated index is summed. Status: STANDARD.
\item
  \(\nu\): A four-dimensional coordinate index (as \(\mu\)). Status:
  STANDARD.
\item
  \(\kappa_{\mathrm{ind}}\): The coupling of the Einstein-Hilbert term
  \(\hat R/(2\kappa_{\mathrm{ind}})\) that the quantum fluctuations of
  the fields of the fold may induce,
  \(\kappa_{\mathrm{ind}} = 8\pi G_{\mathrm{ind}}\); it adds to the
  gravity of the fold,
  \(1/(8\pi G_{\mathrm{N}}) = V_*/\kappa_7 + 1/\kappa_{\mathrm{ind}}\).
  Unit: GeV\(^{-2}\). Value: not fixed by UM-4; negligible for a cutoff
  of the order of \(H\). Status: ESTIMATE.
\item
  \(\hat G^\mu{}_\nu\): The Einstein tensor of the metric
  \(\hat g_{ab}\) of the fold (FD-2), with frame indices along its seven
  directions; with four-dimensional indices, \(\hat G^\mu{}_\nu\), its
  components along the four observable directions, which for frozen
  extra times are those of the Einstein tensor of the four-dimensional
  metric \(\hat h_{\mu\nu}\). Unit: 1/length\(^2\). Status: STANDARD.
\item
  \(r\): An extra-time direction of the fold: \(r \in \{5, 6, 7\}\)
  (\(x_5, x_6, x_7\)); not summed (the three components are equal).
  Status: STANDARD.
\item
  \(\mathcal{T}\): The energy-momentum tensor of any one term of the
  fold (in \(\mathcal{T}^r{}_r = \mathcal{L}_{\mathrm{fold}}\)). Unit:
  GeV\(^4\). Status: STANDARD.
\item
  \(\mathcal{L}_{\mathrm{fold}}\): The Lagrangian of any one term of the
  fold (four-dimensional normalisation). Unit: GeV\(^4\). Status:
  STANDARD.
\end{itemize}

The thin source of the fold is the energy-momentum tensor of everything
that lives on the fold, divided by \(V_*\) (IN-8): the Higgs doublet
(FE-7), the quarks and leptons (FE-9), the portals (FE-10), the fold
fluid (IN-10) and the induced gravity (IN-9), whose energy-momentum
tensor is \(-1/\kappa_{\mathrm{ind}}\) times the Einstein tensor
\(\hat G^\mu{}_\nu\) of the metric of the fold; the particles \(\chi\)
of dirac16complex enter through the jump of its profile at the fold. The
gauge fields live in the bulk: their energy-momentum tensor is finite at
the fold and has no thin part (their field strengths jump there, FE-16
to FE-18, but do not diverge). Because the fields of the fold do not
depend on the extra times and have no components along them, the extra
times enter only through \(\sqrt{\lvert g\rvert}\): for every term of
the fold the component along an extra time is the Lagrangian of the
term, \(\mathcal{T}^r{}_r = \mathcal{L}_{\mathrm{fold}}\) (DERIVED:
\(\delta\sqrt{\lvert g\rvert} = -\frac12\sqrt{\lvert g\rvert}g_{rr}\delta g^{rr}\)).
Combining FE-4 and FE-5 gives the equations of the fold (FE-71). Status:
DERIVED.

\hypertarget{the-energy-momentum-tensors}{%
\subsection{The energy-momentum
tensors}\label{the-energy-momentum-tensors}}

\[
{\mathcal{T}_\Psi}_{MN} = -\frac14\Big(\bar\Psi\gamma_MD_N\Psi + \bar\Psi\gamma_ND_M\Psi - (D_M\bar\Psi)\gamma_N\Psi
 - (D_N\bar\Psi)\gamma_M\Psi\Big) + g_{MN}\,\mathcal{L}_\Psi \tag{FE-6}
\] \textbf{Symbols in (FE-6).}

\begin{itemize}
\tightlist
\item
  \(\mathcal{T}_\Psi\): The energy-momentum tensor of dirac16complex:
  \(-2/\sqrt{\lvert g\rvert}\) times the variation of the action
  \(S_\Psi\) of dirac16complex with respect to \(g^{MN}\). Unit:
  GeV\(^8\). Status: STANDARD.
\item
  \(M\): A coordinate index of eight dimensions (also with a numbered
  copy \(M_1, M_2, \dots\)): it runs over \(x_0, x_1, \dots, x_7\); a
  repeated index is summed. Status: STANDARD.
\item
  \(N\): A coordinate index of eight dimensions (also
  \(N_1, N_2, \dots\)), running over \(x_0, \dots, x_7\); a repeated
  index is summed. Status: STANDARD.
\item
  \(\Psi,\ \bar\Psi\): dirac16complex: the author's 16-component complex
  spinor field of eight dimensions;
  \(\bar\Psi = \Psi^\dagger\sigma_{16}\) is its Dirac adjoint; in a term
  of the fold, \(\Psi\) is its value at \(\sigma = 0\). Unit:
  GeV\(^{7/2}\). Status: AUTHOR.
\item
  \(\gamma^M,\ \gamma^{(a)}\): The Dirac matrices of dirac16complex
  along the coordinates, \(\gamma^M = \mathbf{e}_a{}^MT^a\),
  \(\gamma_M = g_{MN}\gamma^N\); \(\gamma^{(a)} = T^a\) along the frame
  direction \((a)\) (\(\gamma^{(0)}\) along \(\sigma\), \(\gamma^{(4)}\)
  along the time). Unit: dimensionless. Status: AUTHOR.
\item
  \(D_M,\ D_\mu\): The covariant derivative: on dirac16complex
  \(D_M = \partial_M + \Omega_M\) (with
  \(D_M\bar\Psi = \partial_M\bar\Psi - \bar\Psi\Omega_M\)); on the
  fields of the fold the gauge and spin covariant derivative \(D_\mu\)
  (the partial derivative plus the spin connection of the fold and the
  gauge fields times their couplings and generators);
  \(D^\mu = \hat h^{\mu\nu}D_\nu\); on an adjoint field
  \((D_\mu X)^A = \partial_\mu X^A + g_3f^{ABC}G^B_\mu X^C\) (likewise
  for \(SU(2)\) with \(g_2\epsilon^{IJK}\)). Unit: 1/length. Status:
  STANDARD.
\item
  \(g_{MN}\): The metric tensor of the eight-dimensional spacetime
  (signature \(+,+,+,+,-,-,-,-\) in the order \(x_0, \dots, x_7\)):
  \(g_{MN}\,\mathrm{d}X^M\mathrm{d}X^N\) is the squared interval,
  \(g^{MN}\) its inverse, \(\lvert g\rvert\) the absolute value of its
  determinant (\(\sqrt{\lvert g\rvert}\,\mathrm{d}^8X\) is the proper
  volume element). Unit: dimensionless. Status: AUTHOR.
\item
  \(\mathcal{L}_\Psi\): The Lagrangian density of dirac16complex: the
  square bracket in its action \(S_\Psi\). Unit: GeV\(^8\). Status:
  RECORD.
\end{itemize}

The symmetric energy-momentum tensor of dirac16complex, from the
variation of the vielbein in IN-4 (\(\mathcal{L}_\Psi\) is the
Lagrangian in the square bracket of IN-4; the variation of the spin
connection adds only the divergence of the spin current, which drops out
of the symmetric part). On solutions of FE-12,
\(\mathcal{L}_\Psi = \frac{\lambda_\Psi}{2}(\bar\Psi\Psi)^2\). For a
stationary state \(\Psi \propto \mathrm{e}^{-\mathrm{i}\omega x_4}\) of
the free field (\(D_{x_4}\Psi = -\mathrm{i}\omega\Psi\),
\(D_{x_4}\bar\Psi = \mathrm{i}\omega\bar\Psi\),
\(\mathcal{L}_\Psi = 0\)) the time-time component is
\(\mathcal{T}_{\Psi\,x_4x_4} = \omega\,\mathrm{i}\bar\Psi\gamma_{x_4}\Psi = -\omega\,\lvert g_{x_4x_4}\rvert\,\mathrm{i}\bar\Psi\gamma^{x_4}\Psi\):
its sign is minus the product of the sign of the frequency and of the
sign of the charge density \(\mathrm{i}\bar\Psi\gamma^{x_4}\Psi\) of
dirac16complex, the Krein density (FE-43: the energy
\(E_o = -\omega_o\mathcal{Q}_o\) of an orbital of the frequency
\(\omega_o\) and the Krein charge \(\mathcal{Q}_o\)). Status: STANDARD
(DERIVED for the energy density).

\[
{\mathcal{T}_{\mathrm{H}}}_{\mu\nu} = (D_\mu\Phi)^\dagger D_\nu\Phi + (D_\nu\Phi)^\dagger D_\mu\Phi + \hat h_{\mu\nu}\,\mathcal{L}_{\mathrm{H}} \tag{FE-7}
\] \textbf{Symbols in (FE-7).}

\begin{itemize}
\tightlist
\item
  \(\mathcal{T}_{\mathrm{H}}\): The energy-momentum tensor of the Higgs
  doublet (four-dimensional normalisation). Unit: GeV\(^4\). Status:
  STANDARD.
\item
  \(\mu\): A four-dimensional coordinate index running over the
  observable coordinates \(x_1, x_2, x_3, x_4\) of the fold (signature
  \(+,+,+,-\)); a repeated index is summed. Status: STANDARD.
\item
  \(\nu\): A four-dimensional coordinate index (as \(\mu\)). Status:
  STANDARD.
\item
  \(D_M,\ D_\mu\): The covariant derivative: on dirac16complex
  \(D_M = \partial_M + \Omega_M\) (with
  \(D_M\bar\Psi = \partial_M\bar\Psi - \bar\Psi\Omega_M\)); on the
  fields of the fold the gauge and spin covariant derivative \(D_\mu\)
  (the partial derivative plus the spin connection of the fold and the
  gauge fields times their couplings and generators);
  \(D^\mu = \hat h^{\mu\nu}D_\nu\); on an adjoint field
  \((D_\mu X)^A = \partial_\mu X^A + g_3f^{ABC}G^B_\mu X^C\) (likewise
  for \(SU(2)\) with \(g_2\epsilon^{IJK}\)). Unit: 1/length. Status:
  STANDARD.
\item
  \(\Phi\): The Higgs doublet (two complex components, hypercharge
  \(1/2\)), a field of the fold. Unit: GeV. Status: AUTHOR.
\item
  \({}^\dagger\): Hermitian conjugation: the complex conjugate of the
  transpose of a matrix or spinor; for an operator on the space of
  states, its adjoint with respect to the scalar product of that space
  (for the quantised dirac16complex the indefinite Krein scalar
  product). Status: STANDARD.
\item
  \(\hat h_{\mu\nu}\): The four-dimensional metric of the fold,
  \(\hat h_{\mu\nu} = \mathrm{diag}(\mathrm{e}^{2a_4}, \mathrm{e}^{2a_4}, \mathrm{e}^{2a_4}, -1)\)
  in the coordinates \(x_1, \dots, x_4\) (the induced metric at
  \(\sigma = 0\) without the extra times); \(\hat h^{\mu\nu}\) its
  inverse, \(\hat h\) its determinant,
  \(\sqrt{-\hat h} = \mathrm{e}^{3a_4}\); \(\hat h^{(1)}_{\mu\nu}\) its
  first-order cosmological perturbation. Unit: dimensionless. Status:
  DERIVED.
\item
  \(\mathcal{L}_{\mathrm{H}}\): The Lagrangian of the Higgs doublet on
  the fold. Unit: GeV\(^4\). Status: AUTHOR.
\end{itemize}

The energy-momentum tensor of the Higgs doublet (four-dimensional
normalisation): its time-time component is
\(\lvert D_4\Phi\rvert^2 + \hat h^{ij}(D_i\Phi)^\dagger D_j\Phi + V \ge V\).
Status: STANDARD.

\[
{\mathcal{T}_{\mathrm{gauge}}}_{MN} = \frac{1}{\mathsf{g}_3^2}\,G^A_{MP}G^A{}_N{}^P + \frac{1}{\mathsf{g}_2^2}\,W^I_{MP}W^I{}_N{}^P
 + \frac{1}{\mathsf{g}_1^2}\,B_{MP}B_N{}^P + g_{MN}\,\mathcal{L}_{\mathrm{gauge}} \tag{FE-8}
\] \textbf{Symbols in (FE-8).}

\begin{itemize}
\tightlist
\item
  \(\mathcal{T}_{\mathrm{gauge}}\): The energy-momentum tensor of the
  gauge fields of the bulk. Unit: GeV\(^8\). Status: STANDARD.
\item
  \(M\): A coordinate index of eight dimensions (also with a numbered
  copy \(M_1, M_2, \dots\)): it runs over \(x_0, x_1, \dots, x_7\); a
  repeated index is summed. Status: STANDARD.
\item
  \(N\): A coordinate index of eight dimensions (also
  \(N_1, N_2, \dots\)), running over \(x_0, \dots, x_7\); a repeated
  index is summed. Status: STANDARD.
\item
  \(\mathsf{g}_1,\ \mathsf{g}_2,\ \mathsf{g}_3\): The gauge couplings of
  eight dimensions of \(U(1)_Y\), \(SU(2)_L\) and \(SU(3)_C\): the
  coefficients of the kinetic terms of the gauge fields of the bulk in
  their geometric normalisation; the four-dimensional couplings are
  \(g_k^2 = 3H\mathsf{g}_k^2/(2V_*)\). Unit: GeV\(^{-2}\)
  (\(\mathsf{g}^2\): length\(^4\)). Status: ASSUMED.
\item
  \(G^A_M,\ G^A_\mu,\ G^A_{MN}\): The gluon field of \(SU(3)_C\)
  (\(A = 1, \dots, 8\)): in the bulk \(G^A_M\), a field of eight
  dimensions in its geometric normalisation (the covariant derivative
  \(\partial - \mathrm{i}G\)), and on the fold the canonically
  normalised four-dimensional field \(G^A_\mu(x)\) of its massless mode,
  \(g_3G^A_\mu(x) = G^A_\mu(\sigma = 0, x)\); \(G^A_{MN}\) and
  \(G^A_{\mu\nu}\) their field strengths. Unit: GeV (field), GeV\(^2\)
  (field strength). Status: AUTHOR.
\item
  \(A\): A colour (adjoint SU(3)) index running over \(1, \dots, 8\); a
  repeated index is summed. Status: STANDARD.
\item
  \(P\): A coordinate index of eight dimensions (also
  \(P_1, P_2, \dots\)), running over \(x_0, \dots, x_7\); a repeated
  index is summed. Status: STANDARD.
\item
  \(W^I_M,\ W^I_\mu,\ W^I_{MN}\): The gauge field of \(SU(2)_L\)
  (\(I = 1, 2, 3\)): in the bulk \(W^I_M\) (geometric normalisation), on
  the fold the canonically normalised four-dimensional field
  \(W^I_\mu(x)\) of its massless mode,
  \(g_2W^I_\mu(x) = W^I_\mu(\sigma = 0, x)\); \(W^I_{MN}\),
  \(W^I_{\mu\nu}\) their field strengths;
  \(W^\pm_\mu = (W^1_\mu \mp \mathrm{i}W^2_\mu)/\sqrt2\). Unit: GeV
  (field), GeV\(^2\) (field strength). Status: AUTHOR.
\item
  \(I\): A weak-isospin (adjoint SU(2)) index running over \(1, 2, 3\);
  a repeated index is summed. Status: STANDARD.
\item
  \(B_M,\ B_\mu,\ B_{MN}\): The hypercharge gauge field of \(U(1)_Y\):
  in the bulk \(B_M\) (geometric normalisation), on the fold the
  canonically normalised four-dimensional field \(B_\mu(x)\) of its
  massless mode, \(g_1B_\mu(x) = B_\mu(\sigma = 0, x)\); \(B_{MN}\),
  \(B_{\mu\nu}\) their field strengths. Unit: GeV (field), GeV\(^2\)
  (field strength). Status: AUTHOR.
\item
  \(g_{MN}\): The metric tensor of the eight-dimensional spacetime
  (signature \(+,+,+,+,-,-,-,-\) in the order \(x_0, \dots, x_7\)):
  \(g_{MN}\,\mathrm{d}X^M\mathrm{d}X^N\) is the squared interval,
  \(g^{MN}\) its inverse, \(\lvert g\rvert\) the absolute value of its
  determinant (\(\sqrt{\lvert g\rvert}\,\mathrm{d}^8X\) is the proper
  volume element). Unit: dimensionless. Status: AUTHOR.
\item
  \(\mathcal{L}_{\mathrm{gauge}}\): The Lagrangian of the gauge fields
  of \(SU(3)_C\times SU(2)_L\times U(1)_Y\) in the bulk. Unit:
  GeV\(^8\). Status: STANDARD.
\end{itemize}

The energy-momentum tensor of the gauge fields in the bulk; for each
field its energy density is \((\vec E^2 + \vec B^2)/(2\mathsf{g}^2)\) in
terms of its electric and magnetic components \(\vec E\) and \(\vec B\)
in the frame (seven-dimensional vectors along the fold and along
\(\sigma\)). For the massless modes integrated over \(\sigma\) and over
the extra times it is the energy-momentum tensor of the gauge fields of
the Standard Model on the fold,
\(G^A_{\mu\rho}G^A{}_\nu{}^\rho + W^I_{\mu\rho}W^I{}_\nu{}^\rho + B_{\mu\rho}B_\nu{}^\rho + \hat h_{\mu\nu}\mathcal{L}_4\)
(the canonical fields of IN-13; \(\mathcal{L}_4\) their Lagrangian).
Status: STANDARD.

\[
{\mathcal{T}_{\mathrm{f}}}_{\mu\nu} = -\frac{\mathrm{i}}{4}\sum_\psi\Big(\bar\psi\hat\gamma_\mu D_\nu\psi + \bar\psi\hat\gamma_\nu D_\mu\psi
 - (D_\mu\bar\psi)\hat\gamma_\nu\psi - (D_\nu\bar\psi)\hat\gamma_\mu\psi\Big) + \hat h_{\mu\nu}\big(\mathcal{L}_{\mathrm{f}} + \mathcal{L}_{\mathrm{Y}}\big) \tag{FE-9}
\] \textbf{Symbols in (FE-9).}

\begin{itemize}
\tightlist
\item
  \(\mathcal{T}_{\mathrm{f}}\): The energy-momentum tensor of the quarks
  and leptons. Unit: GeV\(^4\). Status: STANDARD.
\item
  \(\mu\): A four-dimensional coordinate index running over the
  observable coordinates \(x_1, x_2, x_3, x_4\) of the fold (signature
  \(+,+,+,-\)); a repeated index is summed. Status: STANDARD.
\item
  \(\nu\): A four-dimensional coordinate index (as \(\mu\)). Status:
  STANDARD.
\item
  \(\mathrm{i}\): The imaginary unit, \(\mathrm{i}^2 = -1\). Unit:
  dimensionless. Status: STANDARD.
\item
  \(\psi\): The label of a Standard-Model fermion field in a sum:
  \(\psi\) runs over
  \(Q, u_{\mathrm{R}}, d_{\mathrm{R}}, \ell, e_{\mathrm{R}}, \nu_{\mathrm{R}}\).
  Status: STANDARD.
\item
  \(\psi,\ \bar\psi\): A Standard-Model fermion field (a
  four-dimensional Weyl field of the fold with three generations,
  written as a four-component spinor of fixed chirality);
  \(\bar\psi = \psi^\dagger\beta\); \(\psi_f\) the fields of the kind
  \(f\). Unit: GeV\(^{3/2}\). Status: STANDARD.
\item
  \(\hat\gamma^\mu\): The four-dimensional complex \(4 \times 4\) Dirac
  matrices of the fold: \(\hat\gamma^{(a)}\) along the frame directions
  of the fold and \(\hat\gamma^\mu\) along its coordinates (the frame
  matrices contracted with the inverse frame field of the fold),
  \(\{\hat\gamma^\mu, \hat\gamma^\nu\} = -2\hat h^{\mu\nu}\mathbf{1}\);
  \(\hat\gamma_\mu = \hat h_{\mu\nu}\hat\gamma^\nu\). Unit:
  dimensionless. Status: STANDARD.
\item
  \(D_M,\ D_\mu\): The covariant derivative: on dirac16complex
  \(D_M = \partial_M + \Omega_M\) (with
  \(D_M\bar\Psi = \partial_M\bar\Psi - \bar\Psi\Omega_M\)); on the
  fields of the fold the gauge and spin covariant derivative \(D_\mu\)
  (the partial derivative plus the spin connection of the fold and the
  gauge fields times their couplings and generators);
  \(D^\mu = \hat h^{\mu\nu}D_\nu\); on an adjoint field
  \((D_\mu X)^A = \partial_\mu X^A + g_3f^{ABC}G^B_\mu X^C\) (likewise
  for \(SU(2)\) with \(g_2\epsilon^{IJK}\)). Unit: 1/length. Status:
  STANDARD.
\item
  \(\hat h_{\mu\nu}\): The four-dimensional metric of the fold,
  \(\hat h_{\mu\nu} = \mathrm{diag}(\mathrm{e}^{2a_4}, \mathrm{e}^{2a_4}, \mathrm{e}^{2a_4}, -1)\)
  in the coordinates \(x_1, \dots, x_4\) (the induced metric at
  \(\sigma = 0\) without the extra times); \(\hat h^{\mu\nu}\) its
  inverse, \(\hat h\) its determinant,
  \(\sqrt{-\hat h} = \mathrm{e}^{3a_4}\); \(\hat h^{(1)}_{\mu\nu}\) its
  first-order cosmological perturbation. Unit: dimensionless. Status:
  DERIVED.
\item
  \(\mathcal{L}_{\mathrm{f}}\): The kinetic Lagrangian of the quarks and
  leptons on the fold. Unit: GeV\(^4\). Status: STANDARD.
\item
  \(\mathcal{L}_{\mathrm{Y}}\): The Yukawa Lagrangian of the quarks and
  leptons with the Higgs doublet. Unit: GeV\(^4\). Status: STANDARD.
\end{itemize}

The energy-momentum tensor of the quarks and leptons (with the kinetic
term in its symmetric form). For a free particle of energy \(E\) at rest
in the frame of the fold its time-time component is
\(E\,\psi^\dagger\psi > 0\). Status: STANDARD.

\[
\begin{aligned}
{\mathcal{T}_{\mathrm{P}}}_{\mu\nu} = {}&\sum_X\Big[\big(c^{G}_X\,G^A_{\mu\rho}G^A{}_\nu{}^\rho + c^{W}_X\,W^I_{\mu\rho}W^I{}_\nu{}^\rho
 + c^{B}_X\,B_{\mu\rho}B_\nu{}^\rho\big)\,\bar\Psi P_X\Psi\\
&+ \frac{d_X}{2}\big(B_{\mu\rho}\,\mathcal{D}_{X\nu}{}^\rho + B_{\nu\rho}\,\mathcal{D}_{X\mu}{}^\rho\big)
 + \sum_f\frac{c^{f}_X}{2}\big(J_{f\mu}\,j_{X\nu} + J_{f\nu}\,j_{X\mu}\big)\Big] + \hat h_{\mu\nu}\,\mathcal{L}_{\mathrm{P}}
\end{aligned} \tag{FE-10}
\] \textbf{Symbols in (FE-10).}

\begin{itemize}
\tightlist
\item
  \(\mathcal{T}_{\mathrm{P}}\): The energy-momentum tensor of the
  portals of dirac16complex. Unit: GeV\(^4\). Status: DERIVED.
\item
  \(\mu\): A four-dimensional coordinate index running over the
  observable coordinates \(x_1, x_2, x_3, x_4\) of the fold (signature
  \(+,+,+,-\)); a repeated index is summed. Status: STANDARD.
\item
  \(\nu\): A four-dimensional coordinate index (as \(\mu\)). Status:
  STANDARD.
\item
  \(X\): The label of a half of dirac16complex:
  \(X \in \{\mathrm{L}, \mathrm{R}\}\) (left-handed or right-handed
  half); a sum over \(X\) runs over both. Status: STANDARD.
\item
  \(c^{G}_X\): The gluon-kinetic portal coupling of the half \(X\) of
  dirac16complex. Unit: GeV\(^{-7}\). Status: INVENTED.
\item
  \(G^A_M,\ G^A_\mu,\ G^A_{MN}\): The gluon field of \(SU(3)_C\)
  (\(A = 1, \dots, 8\)): in the bulk \(G^A_M\), a field of eight
  dimensions in its geometric normalisation (the covariant derivative
  \(\partial - \mathrm{i}G\)), and on the fold the canonically
  normalised four-dimensional field \(G^A_\mu(x)\) of its massless mode,
  \(g_3G^A_\mu(x) = G^A_\mu(\sigma = 0, x)\); \(G^A_{MN}\) and
  \(G^A_{\mu\nu}\) their field strengths. Unit: GeV (field), GeV\(^2\)
  (field strength). Status: AUTHOR.
\item
  \(A\): A colour (adjoint SU(3)) index running over \(1, \dots, 8\); a
  repeated index is summed. Status: STANDARD.
\item
  \(\rho\): A four-dimensional coordinate index (as \(\mu\)). Status:
  STANDARD.
\item
  \(c^{W}_X\): The \(W\)-kinetic portal coupling of the half \(X\).
  Unit: GeV\(^{-7}\). Status: INVENTED.
\item
  \(W^I_M,\ W^I_\mu,\ W^I_{MN}\): The gauge field of \(SU(2)_L\)
  (\(I = 1, 2, 3\)): in the bulk \(W^I_M\) (geometric normalisation), on
  the fold the canonically normalised four-dimensional field
  \(W^I_\mu(x)\) of its massless mode,
  \(g_2W^I_\mu(x) = W^I_\mu(\sigma = 0, x)\); \(W^I_{MN}\),
  \(W^I_{\mu\nu}\) their field strengths;
  \(W^\pm_\mu = (W^1_\mu \mp \mathrm{i}W^2_\mu)/\sqrt2\). Unit: GeV
  (field), GeV\(^2\) (field strength). Status: AUTHOR.
\item
  \(I\): A weak-isospin (adjoint SU(2)) index running over \(1, 2, 3\);
  a repeated index is summed. Status: STANDARD.
\item
  \(c^{B}_X\): The \(B\)-kinetic portal coupling of the half \(X\).
  Unit: GeV\(^{-7}\). Status: INVENTED.
\item
  \(B_M,\ B_\mu,\ B_{MN}\): The hypercharge gauge field of \(U(1)_Y\):
  in the bulk \(B_M\) (geometric normalisation), on the fold the
  canonically normalised four-dimensional field \(B_\mu(x)\) of its
  massless mode, \(g_1B_\mu(x) = B_\mu(\sigma = 0, x)\); \(B_{MN}\),
  \(B_{\mu\nu}\) their field strengths. Unit: GeV (field), GeV\(^2\)
  (field strength). Status: AUTHOR.
\item
  \(\Psi,\ \bar\Psi\): dirac16complex: the author's 16-component complex
  spinor field of eight dimensions;
  \(\bar\Psi = \Psi^\dagger\sigma_{16}\) is its Dirac adjoint; in a term
  of the fold, \(\Psi\) is its value at \(\sigma = 0\). Unit:
  GeV\(^{7/2}\). Status: AUTHOR.
\item
  \(P_X\): The chiral projector of the half \(X\):
  \(P_X = P_{\mathrm{L}}\) for \(X = \mathrm{L}\), \(P_{\mathrm{R}}\)
  for \(X = \mathrm{R}\). Unit: dimensionless. Status: AUTHOR.
\item
  \(d_X\): The dipole-portal coupling of the half \(X\) of
  dirac16complex with the hypercharge field strength. Unit:
  GeV\(^{-5}\). Status: INVENTED.
\item
  \(\mathcal{D}_X^{\mu\nu}\): The dipole (antisymmetric tensor) density
  of the half \(X\) of dirac16complex at the fold,
  \(\mathcal{D}_X^{\mu\nu} = \mathrm{i}\bar\Psi\Sigma^{\mu\nu}P_X\Psi\)
  (real, DERIVED). Unit: GeV\(^7\). Status: DERIVED.
\item
  \(f\): The label of a kind of Standard-Model multiplet:
  \(f \in \{Q, u, d, \ell, e, \nu\}\) (quark doublets, up-quark
  singlets, down-quark singlets, lepton doublets, charged-lepton
  singlets, right-handed neutrinos); a sum over \(f\) runs over the six.
  Status: STANDARD.
\item
  \(c^{f}_X\): The current-portal coupling of the half \(X\) of
  dirac16complex with the fermion-number current of the multiplets of
  kind \(f\) (flavour blind); \(c^{\ell}_X\) and \(c^{\nu}_X\) are those
  of the lepton doublets and of the right-handed neutrinos. Unit:
  GeV\(^{-7}\). Status: INVENTED.
\item
  \(J_f^\mu\): The fermion-number current of the Standard-Model
  multiplets of kind \(f\), \(J_f^\mu = \bar\psi_f\hat\gamma^\mu\psi_f\)
  summed over the three generations and the gauge components; with other
  labels the currents that are the sources of the gauge fields. Unit:
  GeV\(^3\). Status: STANDARD.
\item
  \(j_X^\mu\): The convective current of the half \(X\) of
  dirac16complex,
  \(j_X^\mu = \mathrm{i}(\bar\Psi P_XD^\mu\Psi - (D^\mu\bar\Psi)P_X\Psi)\)
  (real and joining the half with itself, DERIVED). Unit: GeV\(^8\).
  Status: DERIVED.
\item
  \(\hat h_{\mu\nu}\): The four-dimensional metric of the fold,
  \(\hat h_{\mu\nu} = \mathrm{diag}(\mathrm{e}^{2a_4}, \mathrm{e}^{2a_4}, \mathrm{e}^{2a_4}, -1)\)
  in the coordinates \(x_1, \dots, x_4\) (the induced metric at
  \(\sigma = 0\) without the extra times); \(\hat h^{\mu\nu}\) its
  inverse, \(\hat h\) its determinant,
  \(\sqrt{-\hat h} = \mathrm{e}^{3a_4}\); \(\hat h^{(1)}_{\mu\nu}\) its
  first-order cosmological perturbation. Unit: dimensionless. Status:
  DERIVED.
\item
  \(\mathcal{L}_{\mathrm{P}}\): The Lagrangian of the portals of
  dirac16complex at the fold. Unit: GeV\(^4\). Status: INVENTED.
\end{itemize}

The energy-momentum tensor of the portals (from the variation of IN-16
with respect to the metric of the fold, which enters through the raised
indices, through \(\hat\gamma^\mu\) in \(J_f^\mu\) and through
\(\sqrt{-\hat h}\)); the Higgs portal enters only through
\(\hat h_{\mu\nu}\mathcal{L}_{\mathrm{P}}\); the field strengths are
those of the canonical four-dimensional fields at the fold (IN-13). It
is a source of the fold. Status: DERIVED.

\hypertarget{the-conservation-laws}{%
\subsection{The conservation laws}\label{the-conservation-laws}}

\[
\begin{aligned}
&\nabla_M\,\big(\mathcal{T}_\Psi{}^M{}_N + \mathcal{T}_{\mathrm{gauge}}{}^M{}_N\big) = 0\qquad (\sigma \ne 0),\\
&\hat\nabla_a\,\mathcal{S}^a{}_b = -\Big(\mathcal{T}_{\mathrm{bulk}}{}^{(0)}{}_b\big\vert_{\sigma=0^+}
 - \mathcal{T}_{\mathrm{bulk}}{}^{(0)}{}_b\big\vert_{\sigma=0^-}\Big),\qquad
 \mathcal{T}_{\mathrm{bulk}} = \mathcal{T}_\Psi + \mathcal{T}_{\mathrm{gauge}}
\end{aligned} \tag{FE-11}
\] \textbf{Symbols in (FE-11).}

\begin{itemize}
\tightlist
\item
  \(\nabla\): The covariant derivative of a tensor with the Levi-Civita
  connection of the metric \(g_{MN}\) (Christoffel symbols
  \(\Gamma^P{}_{MN}\)). Status: STANDARD.
\item
  \(\mathcal{T}_\Psi\): The energy-momentum tensor of dirac16complex:
  \(-2/\sqrt{\lvert g\rvert}\) times the variation of the action
  \(S_\Psi\) of dirac16complex with respect to \(g^{MN}\). Unit:
  GeV\(^8\). Status: STANDARD.
\item
  \(\mathcal{T}_{\mathrm{gauge}}\): The energy-momentum tensor of the
  gauge fields of the bulk. Unit: GeV\(^8\). Status: STANDARD.
\item
  \(M\): A coordinate index of eight dimensions (also with a numbered
  copy \(M_1, M_2, \dots\)): it runs over \(x_0, x_1, \dots, x_7\); a
  repeated index is summed. Status: STANDARD.
\item
  \(N\): A coordinate index of eight dimensions (also
  \(N_1, N_2, \dots\)), running over \(x_0, \dots, x_7\); a repeated
  index is summed. Status: STANDARD.
\item
  \(\sigma\): The proper coordinate of the hidden direction:
  \(\mathrm{d}\sigma = \lvert\cot z\rvert\,\mathrm{d}x_0\),
  \(\sigma = \mathrm{sgn}(\pi/2 - z)\ln(\sin z)/(6H)\) on one period
  \(0 < z < \pi\); \(\sigma = 0\) is the fold \(z = \pi/2\),
  \(\sigma \to \mp\infty\) the tips \(z \to 0, \pi\); one period of
  \(x_0\) is the whole line \(-\infty < \sigma < \infty\). Unit: length.
  Status: DERIVED.
\item
  \(\hat\nabla\): The covariant derivative of the fold: with the
  Levi-Civita connection of the metric of the fold (for four-dimensional
  tensors of its metric
  \(\hat h_{\mu\nu} = \mathrm{diag}(\mathrm{e}^{2a_4}, \mathrm{e}^{2a_4}, \mathrm{e}^{2a_4}, -1)\)
  in the observable coordinates \(x_1, \dots, x_4\); for
  \(\hat\nabla_a\mathcal{S}^a{}_b\) of the induced metric of the
  seven-dimensional fold). Status: STANDARD.
\item
  \(\mathcal{S}^a{}_b\): The thin source of the fold: the
  energy-momentum tensor concentrated on the fold, integrated across it,
  \(\kappa\mathcal{S}^a{}_b = \int_{0^-}^{0^+}\mathrm{d}\sigma\,\mathcal{E}_{(1)}{}^a{}_b\)
  (the junction condition of Israel); \(\mathcal{S}^4{}_4\) is minus its
  energy density, \(\mathcal{S}^1{}_1\) and \(\mathcal{S}^5{}_5\) its
  pressures along ordinary space and along the extra times (frame
  components along the fold, eight-dimensional normalisation: an energy
  per unit volume of the six other directions of the fold). Unit:
  GeV\(^7\). Value: \(\kappa\mathcal{S}^a{}_b = -12H\delta^a{}_b\) for
  every \(a_4\), \(b_4\) of the author's form. Status: DERIVED.
\item
  \(a\): A frame (orthonormal) index running over \(0, 1, \dots, 7\)
  (frame direction \((0)\) along \(\sigma\), \((1), (2), (3)\) along
  ordinary space, \((4)\) along the time, \((5), (6), (7)\) along the
  extra times); raised and lowered with the frame metric
  \(\eta_{ab} = \mathrm{diag}(1, 1, 1, 1, -1, -1, -1, -1)\); on the fold
  (\(a, b \in \{1, \dots, 7\}\)) the frame of the fold; in the tensors
  of the fold written with its coordinates \(x^a\) (its metric
  \(\hat g_{ab}\) and the tensors built from it) a coordinate index of
  the fold, \(a \in \{1, \dots, 7\}\); a repeated index is summed.
  Status: STANDARD.
\item
  \(b\): A frame index (as \(a\)). Status: STANDARD.
\item
  \(\mathcal{T}_{\mathrm{bulk}}\): The energy-momentum tensor of all
  sources of the bulk, dirac16complex and the gauge fields,
  \(\mathcal{T}_{\mathrm{bulk}} = \mathcal{T}_\Psi + \mathcal{T}_{\mathrm{gauge}}\);
  \(\mathcal{T}_{\mathrm{bulk}}{}^{(0)}{}_b\) is its flux along
  \(\sigma\) (frame components). Unit: GeV\(^8\). Status: STANDARD.
\end{itemize}

The energy and momentum of the bulk are conserved away from the fold
(FE-3 and FE-1). At the fold the thin source exchanges energy and
momentum with the bulk: its covariant divergence along the fold equals
minus the jump of the flux \(\mathcal{T}_{\mathrm{bulk}}{}^{(0)}{}_b\)
of the sources of the bulk along \(\sigma\) (the frame component
\((0), b\)). In UM-4 two kinds of flux leave the fold: the waves of
dirac16complex by which the particles \(\chi\) decay into the bulk
(FE-86), and the massive modes of the gauge fields that the currents of
the fold emit (the width of the W and Z bosons into the bulk, FE-82).
Without them (the static block, a fold with frozen extra times and no
such emission) the matter of the fold is conserved by itself. Status:
DERIVED.

\hypertarget{the-matter-fields}{%
\section{The matter fields}\label{the-matter-fields}}

\hypertarget{dirac16complex}{%
\subsection{dirac16complex}\label{dirac16complex}}

\[
\gamma^MD_M\Psi - m_\Psi\Psi - \lambda_\Psi\big(\bar\Psi\Psi\big)\Psi = 0\qquad (\sigma \ne 0) \tag{FE-12}
\] \textbf{Symbols in (FE-12).}

\begin{itemize}
\tightlist
\item
  \(\gamma^M,\ \gamma^{(a)}\): The Dirac matrices of dirac16complex
  along the coordinates, \(\gamma^M = \mathbf{e}_a{}^MT^a\),
  \(\gamma_M = g_{MN}\gamma^N\); \(\gamma^{(a)} = T^a\) along the frame
  direction \((a)\) (\(\gamma^{(0)}\) along \(\sigma\), \(\gamma^{(4)}\)
  along the time). Unit: dimensionless. Status: AUTHOR.
\item
  \(M\): A coordinate index of eight dimensions (also with a numbered
  copy \(M_1, M_2, \dots\)): it runs over \(x_0, x_1, \dots, x_7\); a
  repeated index is summed. Status: STANDARD.
\item
  \(D_M,\ D_\mu\): The covariant derivative: on dirac16complex
  \(D_M = \partial_M + \Omega_M\) (with
  \(D_M\bar\Psi = \partial_M\bar\Psi - \bar\Psi\Omega_M\)); on the
  fields of the fold the gauge and spin covariant derivative \(D_\mu\)
  (the partial derivative plus the spin connection of the fold and the
  gauge fields times their couplings and generators);
  \(D^\mu = \hat h^{\mu\nu}D_\nu\); on an adjoint field
  \((D_\mu X)^A = \partial_\mu X^A + g_3f^{ABC}G^B_\mu X^C\) (likewise
  for \(SU(2)\) with \(g_2\epsilon^{IJK}\)). Unit: 1/length. Status:
  STANDARD.
\item
  \(\Psi,\ \bar\Psi\): dirac16complex: the author's 16-component complex
  spinor field of eight dimensions;
  \(\bar\Psi = \Psi^\dagger\sigma_{16}\) is its Dirac adjoint; in a term
  of the fold, \(\Psi\) is its value at \(\sigma = 0\). Unit:
  GeV\(^{7/2}\). Status: AUTHOR.
\item
  \(m_\Psi\): The mass parameter of dirac16complex in the bulk. Unit:
  GeV. Status: RECORD.
\item
  \(\lambda_\Psi\): The coupling of the contact self-interaction of
  dirac16complex. Unit: GeV\(^{-6}\). Status: RECORD.
\item
  \(\sigma\): The proper coordinate of the hidden direction:
  \(\mathrm{d}\sigma = \lvert\cot z\rvert\,\mathrm{d}x_0\),
  \(\sigma = \mathrm{sgn}(\pi/2 - z)\ln(\sin z)/(6H)\) on one period
  \(0 < z < \pi\); \(\sigma = 0\) is the fold \(z = \pi/2\),
  \(\sigma \to \mp\infty\) the tips \(z \to 0, \pi\); one period of
  \(x_0\) is the whole line \(-\infty < \sigma < \infty\). Unit: length.
  Status: DERIVED.
\end{itemize}

The field equation of dirac16complex in the bulk: the variation of IN-4
with respect to \(\bar\Psi\) (the second half of the symmetric kinetic
term, integrated by parts with \(D_M\gamma^M = 0\), doubles the first).
With the author's real matrices the operator \(\gamma^MD_M - m_\Psi\) is
real: the field equation does not contain \(\mathrm{i}\), and a real
solution exists for every complex one. Status: RECORD.

\[
\begin{aligned}
&\Psi\big\vert_{\sigma=0^+} = \exp\Big(\frac{1}{V_*}\,\gamma^{(0)}\,\mathbb{M}\Big)\,\Psi\big\vert_{\sigma=0^-},\\
&\mathbb{M} = c_\Phi\big(\Phi^\dagger\Phi\big)P_{\mathrm{L}} + \sum_X\Big[\frac14\big(c^{G}_X\,G^A_{\mu\nu}G^{A\mu\nu}
 + c^{W}_X\,W^I_{\mu\nu}W^{I\mu\nu} + c^{B}_X\,B_{\mu\nu}B^{\mu\nu}\big)P_X + \frac{d_X}{2}\,B_{\mu\nu}\,\mathrm{i}\Sigma^{\mu\nu}P_X\\
&\qquad\qquad + \sum_f c^{f}_X\big(2\mathrm{i}\,J_f^\mu P_XD_\mu + \mathrm{i}(\hat\nabla_\mu J_f^\mu)P_X\big)\Big]
\end{aligned} \tag{FE-13}
\] \textbf{Symbols in (FE-13).}

\begin{itemize}
\tightlist
\item
  \(\Psi,\ \bar\Psi\): dirac16complex: the author's 16-component complex
  spinor field of eight dimensions;
  \(\bar\Psi = \Psi^\dagger\sigma_{16}\) is its Dirac adjoint; in a term
  of the fold, \(\Psi\) is its value at \(\sigma = 0\). Unit:
  GeV\(^{7/2}\). Status: AUTHOR.
\item
  \(\sigma\): The proper coordinate of the hidden direction:
  \(\mathrm{d}\sigma = \lvert\cot z\rvert\,\mathrm{d}x_0\),
  \(\sigma = \mathrm{sgn}(\pi/2 - z)\ln(\sin z)/(6H)\) on one period
  \(0 < z < \pi\); \(\sigma = 0\) is the fold \(z = \pi/2\),
  \(\sigma \to \mp\infty\) the tips \(z \to 0, \pi\); one period of
  \(x_0\) is the whole line \(-\infty < \sigma < \infty\). Unit: length.
  Status: DERIVED.
\item
  \(V_*\): The physical volume of the extra times \(x_5, x_6, x_7\) on
  the fold in the late universe,
  \(V_* = V_{\mathrm{t}}(\mathrm{e}^{-a_4}b_4)^3\) (with
  \(V_{\mathrm{t}} = \int\mathrm{d}x_5\,\mathrm{d}x_6\,\mathrm{d}x_7\)
  the coordinate volume of the extra times) at a time when
  \(\mathrm{e}^{-a_4}b_4\) is constant; it converts the
  eight-dimensional normalisation into the four-dimensional one. Unit:
  length\(^3\) (GeV\(^{-3}\)). Status: ASSUMED.
\item
  \(\gamma^M,\ \gamma^{(a)}\): The Dirac matrices of dirac16complex
  along the coordinates, \(\gamma^M = \mathbf{e}_a{}^MT^a\),
  \(\gamma_M = g_{MN}\gamma^N\); \(\gamma^{(a)} = T^a\) along the frame
  direction \((a)\) (\(\gamma^{(0)}\) along \(\sigma\), \(\gamma^{(4)}\)
  along the time). Unit: dimensionless. Status: AUTHOR.
\item
  \(\mathbb{M}\): The fold operator of dirac16complex: the
  \(16 \times 16\) matrix (with a derivative along the fold) built from
  the Standard-Model fields at the fold by varying the portals with
  respect to \(\bar\Psi\). Unit: GeV\(^{-3}\) (so that
  \(\mathbb{M}/V_*\) is dimensionless). Status: DERIVED.
\item
  \(c_\Phi\): The Higgs-portal coupling of the left-handed half of
  dirac16complex (the right-handed half has none). Unit: GeV\(^{-5}\).
  Status: INVENTED.
\item
  \(\Phi\): The Higgs doublet (two complex components, hypercharge
  \(1/2\)), a field of the fold. Unit: GeV. Status: AUTHOR.
\item
  \({}^\dagger\): Hermitian conjugation: the complex conjugate of the
  transpose of a matrix or spinor; for an operator on the space of
  states, its adjoint with respect to the scalar product of that space
  (for the quantised dirac16complex the indefinite Krein scalar
  product). Status: STANDARD.
\item
  \(P_{\mathrm{L}}\): The author's left-handed chiral projector,
  \(P_{\mathrm{L}} = (\mathbf{1} - \Gamma_8)/2\). Unit: dimensionless.
  Status: AUTHOR.
\item
  \(X\): The label of a half of dirac16complex:
  \(X \in \{\mathrm{L}, \mathrm{R}\}\) (left-handed or right-handed
  half); a sum over \(X\) runs over both. Status: STANDARD.
\item
  \(c^{G}_X\): The gluon-kinetic portal coupling of the half \(X\) of
  dirac16complex. Unit: GeV\(^{-7}\). Status: INVENTED.
\item
  \(G^A_M,\ G^A_\mu,\ G^A_{MN}\): The gluon field of \(SU(3)_C\)
  (\(A = 1, \dots, 8\)): in the bulk \(G^A_M\), a field of eight
  dimensions in its geometric normalisation (the covariant derivative
  \(\partial - \mathrm{i}G\)), and on the fold the canonically
  normalised four-dimensional field \(G^A_\mu(x)\) of its massless mode,
  \(g_3G^A_\mu(x) = G^A_\mu(\sigma = 0, x)\); \(G^A_{MN}\) and
  \(G^A_{\mu\nu}\) their field strengths. Unit: GeV (field), GeV\(^2\)
  (field strength). Status: AUTHOR.
\item
  \(A\): A colour (adjoint SU(3)) index running over \(1, \dots, 8\); a
  repeated index is summed. Status: STANDARD.
\item
  \(\mu\): A four-dimensional coordinate index running over the
  observable coordinates \(x_1, x_2, x_3, x_4\) of the fold (signature
  \(+,+,+,-\)); a repeated index is summed. Status: STANDARD.
\item
  \(\nu\): A four-dimensional coordinate index (as \(\mu\)). Status:
  STANDARD.
\item
  \(c^{W}_X\): The \(W\)-kinetic portal coupling of the half \(X\).
  Unit: GeV\(^{-7}\). Status: INVENTED.
\item
  \(W^I_M,\ W^I_\mu,\ W^I_{MN}\): The gauge field of \(SU(2)_L\)
  (\(I = 1, 2, 3\)): in the bulk \(W^I_M\) (geometric normalisation), on
  the fold the canonically normalised four-dimensional field
  \(W^I_\mu(x)\) of its massless mode,
  \(g_2W^I_\mu(x) = W^I_\mu(\sigma = 0, x)\); \(W^I_{MN}\),
  \(W^I_{\mu\nu}\) their field strengths;
  \(W^\pm_\mu = (W^1_\mu \mp \mathrm{i}W^2_\mu)/\sqrt2\). Unit: GeV
  (field), GeV\(^2\) (field strength). Status: AUTHOR.
\item
  \(I\): A weak-isospin (adjoint SU(2)) index running over \(1, 2, 3\);
  a repeated index is summed. Status: STANDARD.
\item
  \(c^{B}_X\): The \(B\)-kinetic portal coupling of the half \(X\).
  Unit: GeV\(^{-7}\). Status: INVENTED.
\item
  \(B_M,\ B_\mu,\ B_{MN}\): The hypercharge gauge field of \(U(1)_Y\):
  in the bulk \(B_M\) (geometric normalisation), on the fold the
  canonically normalised four-dimensional field \(B_\mu(x)\) of its
  massless mode, \(g_1B_\mu(x) = B_\mu(\sigma = 0, x)\); \(B_{MN}\),
  \(B_{\mu\nu}\) their field strengths. Unit: GeV (field), GeV\(^2\)
  (field strength). Status: AUTHOR.
\item
  \(P_X\): The chiral projector of the half \(X\):
  \(P_X = P_{\mathrm{L}}\) for \(X = \mathrm{L}\), \(P_{\mathrm{R}}\)
  for \(X = \mathrm{R}\). Unit: dimensionless. Status: AUTHOR.
\item
  \(d_X\): The dipole-portal coupling of the half \(X\) of
  dirac16complex with the hypercharge field strength. Unit:
  GeV\(^{-5}\). Status: INVENTED.
\item
  \(\mathrm{i}\): The imaginary unit, \(\mathrm{i}^2 = -1\). Unit:
  dimensionless. Status: STANDARD.
\item
  \(\Sigma^{ab},\ \Sigma^{\mu\nu}\): The generators of Spin(4,4) on
  dirac16complex, \(\Sigma^{ab} = \frac14[T^a, T^b]\) (the scaled
  commutators of the author's matrices);
  \(\Sigma^{\mu\nu} = \mathbf{e}_a{}^\mu\mathbf{e}_b{}^\nu\Sigma^{ab}\)
  along the fold. Unit: dimensionless. Status: AUTHOR.
\item
  \(f\): The label of a kind of Standard-Model multiplet:
  \(f \in \{Q, u, d, \ell, e, \nu\}\) (quark doublets, up-quark
  singlets, down-quark singlets, lepton doublets, charged-lepton
  singlets, right-handed neutrinos); a sum over \(f\) runs over the six.
  Status: STANDARD.
\item
  \(c^{f}_X\): The current-portal coupling of the half \(X\) of
  dirac16complex with the fermion-number current of the multiplets of
  kind \(f\) (flavour blind); \(c^{\ell}_X\) and \(c^{\nu}_X\) are those
  of the lepton doublets and of the right-handed neutrinos. Unit:
  GeV\(^{-7}\). Status: INVENTED.
\item
  \(J_f^\mu\): The fermion-number current of the Standard-Model
  multiplets of kind \(f\), \(J_f^\mu = \bar\psi_f\hat\gamma^\mu\psi_f\)
  summed over the three generations and the gauge components; with other
  labels the currents that are the sources of the gauge fields. Unit:
  GeV\(^3\). Status: STANDARD.
\item
  \(D_M,\ D_\mu\): The covariant derivative: on dirac16complex
  \(D_M = \partial_M + \Omega_M\) (with
  \(D_M\bar\Psi = \partial_M\bar\Psi - \bar\Psi\Omega_M\)); on the
  fields of the fold the gauge and spin covariant derivative \(D_\mu\)
  (the partial derivative plus the spin connection of the fold and the
  gauge fields times their couplings and generators);
  \(D^\mu = \hat h^{\mu\nu}D_\nu\); on an adjoint field
  \((D_\mu X)^A = \partial_\mu X^A + g_3f^{ABC}G^B_\mu X^C\) (likewise
  for \(SU(2)\) with \(g_2\epsilon^{IJK}\)). Unit: 1/length. Status:
  STANDARD.
\item
  \(\hat\nabla\): The covariant derivative of the fold: with the
  Levi-Civita connection of the metric of the fold (for four-dimensional
  tensors of its metric
  \(\hat h_{\mu\nu} = \mathrm{diag}(\mathrm{e}^{2a_4}, \mathrm{e}^{2a_4}, \mathrm{e}^{2a_4}, -1)\)
  in the observable coordinates \(x_1, \dots, x_4\); for
  \(\hat\nabla_a\mathcal{S}^a{}_b\) of the induced metric of the
  seven-dimensional fold). Status: STANDARD.
\end{itemize}

The equation of dirac16complex at the fold: the variation of the portals
IN-16 with respect to \(\bar\Psi\) adds
\(-\delta_{\mathrm{D}}(\sigma)\,\mathbb{M}\Psi/V_*\) to FE-12.
\(\mathbb{M}\) is the fold operator, a \(16 \times 16\) matrix (with a
derivative along the fold from the current portals) built from the
fields of the fold and the gauge fields at the fold. Integrating FE-12
with this term across a fold of vanishing thickness gives the matching
condition in the first line: the values of dirac16complex on the two
sides of the fold differ by the exponential of
\(\gamma^{(0)}\mathbb{M}/V_*\) (the exact solution of
\(\partial_\sigma\Psi = \gamma^{(0)}\mathbb{M}\,\delta_{\mathrm{D}}(\sigma)\Psi/V_*\)
in a layer in which the fold term dominates;
\((\gamma^{(0)})^2 = \mathbf{1}\)). The Higgs term acts on the
left-handed half only. Status: DERIVED (from INVENTED terms).

\[
\gamma^M\Omega_M = -\frac72\,H\,\mathrm{sgn}(\sigma)\,\gamma^{(0)} + \frac32\,\frac{b_4'}{b_4}\,\mathrm{e}^{H\lvert\sigma\rvert}\,\gamma^{(4)} \ne 0 \tag{FE-14}
\] \textbf{Symbols in (FE-14).}

\begin{itemize}
\tightlist
\item
  \(\gamma^M,\ \gamma^{(a)}\): The Dirac matrices of dirac16complex
  along the coordinates, \(\gamma^M = \mathbf{e}_a{}^MT^a\),
  \(\gamma_M = g_{MN}\gamma^N\); \(\gamma^{(a)} = T^a\) along the frame
  direction \((a)\) (\(\gamma^{(0)}\) along \(\sigma\), \(\gamma^{(4)}\)
  along the time). Unit: dimensionless. Status: AUTHOR.
\item
  \(M\): A coordinate index of eight dimensions (also with a numbered
  copy \(M_1, M_2, \dots\)): it runs over \(x_0, x_1, \dots, x_7\); a
  repeated index is summed. Status: STANDARD.
\item
  \(\Omega_M\): The spin connection acting on dirac16complex,
  \(\Omega_M = \frac12\varpi_{Mab}\Sigma^{ab}\). Unit: 1/length. Status:
  RECORD.
\item
  \(H\): The author's constant (an inverse length) that fixes the period
  \(\pi/(6H)\) of \(x_0\), the warp factor
  \(\mathrm{e}^{-2H\lvert\sigma\rvert}\) of the metric, the cosmological
  constant \(\Lambda = -21H^2\) and the tension
  \(\kappa\lambda_{\mathrm{F}} = 12H\) of the fold; in UM-4 a free
  parameter, bounded from below by the propagation of the gauge fields
  in the bulk. Unit: GeV (1/length). Value: \(H > 2534\) GeV, about 2.5
  TeV (from the invisible width of the Z boson). Status: AUTHOR.
\item
  \(\mathrm{sgn}\): The sign function: \(+1\) for a positive argument
  and \(-1\) for a negative one. Status: STANDARD.
\item
  \(\sigma\): The proper coordinate of the hidden direction:
  \(\mathrm{d}\sigma = \lvert\cot z\rvert\,\mathrm{d}x_0\),
  \(\sigma = \mathrm{sgn}(\pi/2 - z)\ln(\sin z)/(6H)\) on one period
  \(0 < z < \pi\); \(\sigma = 0\) is the fold \(z = \pi/2\),
  \(\sigma \to \mp\infty\) the tips \(z \to 0, \pi\); one period of
  \(x_0\) is the whole line \(-\infty < \sigma < \infty\). Unit: length.
  Status: DERIVED.
\item
  \(b_4\): The author's metric function of the time \(x_4\) in the
  extra-time components, \(b_4 > 0\), never identically 1 (the author's
  order of 2026-10-24); the scale factor (the factor of the lengths) of
  the extra times is \(\mathrm{e}^{-a_4}b_4\) on the fold and
  \(\mathrm{e}^{-a_4}b_4\,\mathrm{e}^{-H\lvert\sigma\rvert}\) at the
  distance \(\sigma\) from it (the square root of the warp factor times
  \(\mathrm{e}^{-a_4}b_4\)). Unit: dimensionless. Status: AUTHOR.
\item
  \({}'\): The derivative of a function of the time \(x_4\) alone with
  respect to \(x_4\) (\(a_4' = \mathrm{d}a_4/\mathrm{d}x_4\); two primes
  the second derivative); on a function of \(\sigma\) alone (a profile)
  the derivative with respect to \(\sigma\). Status: STANDARD.
\item
  \(\mathrm{e}\): Euler's number, the base of the natural logarithm;
  \(\mathrm{e}^{X}\) is the exponential of \(X\). Unit: dimensionless.
  Value: 2.718281828. Status: STANDARD.
\end{itemize}

The coupling of dirac16complex to gravity through the canonical spin
connection, for the author's form of the metric of UM-4: the contraction
\(\gamma^M\Omega_M\) that enters the Dirac operator. The first term
comes from the warp of the seven directions of the fold (each
contributes \(\frac12\partial_\sigma\) of the logarithm of its squared
scale factor, \(-H\,\mathrm{sgn}(\sigma)\); in UM-3 six directions gave
\(-3H\)), the second from the time dependence of their scale factors
(\(a_4\) cancels between ordinary space and the extra times; the factor
\(\mathrm{e}^{H\lvert\sigma\rvert}\) is the inverse of the warp of the
time). The first term never vanishes (\(H > 0\), and it only changes
sign at the fold); the second is non-zero whenever \(b_4\) depends on
\(x_4\). Hence the coupling of dirac16complex to gravity through the
spin connection is never identically zero (DERIVED,
\texttt{um4\_geometry.py}, G9, with the author's matrices; the vielbein
postulate \(D_M\gamma^N = 0\) holds, G9). The first term is removed
exactly by the factor \(\mathrm{e}^{(7/2)H\lvert\sigma\rvert}\) of
FE-36. Status: DERIVED.

\hypertarget{the-higgs-doublet}{%
\subsection{The Higgs doublet}\label{the-higgs-doublet}}

\[
\begin{aligned}
&\frac{1}{\sqrt{-\hat h}\,\mathcal{Z}}\,D_\mu\big(\sqrt{-\hat h}\,\mathcal{Z}\,\hat h^{\mu\nu}D_\nu\Phi\big) - \mu^2\Phi
 - 2\lambda\big(\Phi^\dagger\Phi\big)\Phi - c_\Phi\,\bar\Psi P_{\mathrm{L}}\Psi\,\Phi = \Theta,\\
&\Theta = \begin{pmatrix}\bar d_{\mathrm{R}}(\mathcal{Y}^{d})^\dagger u_{\mathrm{L}} + \bar e_{\mathrm{R}}(\mathcal{Y}^{e})^\dagger\nu_{\mathrm{L}}
 - \bar d_{\mathrm{L}}\mathcal{Y}^{u}u_{\mathrm{R}} - \bar e_{\mathrm{L}}\mathcal{Y}^{\nu}\nu_{\mathrm{R}}\\
 \bar d_{\mathrm{R}}(\mathcal{Y}^{d})^\dagger d_{\mathrm{L}} + \bar e_{\mathrm{R}}(\mathcal{Y}^{e})^\dagger e_{\mathrm{L}}
 + \bar u_{\mathrm{L}}\mathcal{Y}^{u}u_{\mathrm{R}} + \bar\nu_{\mathrm{L}}\mathcal{Y}^{\nu}\nu_{\mathrm{R}}\end{pmatrix}
\end{aligned} \tag{FE-15}
\] \textbf{Symbols in (FE-15).}

\begin{itemize}
\tightlist
\item
  \(\hat h_{\mu\nu}\): The four-dimensional metric of the fold,
  \(\hat h_{\mu\nu} = \mathrm{diag}(\mathrm{e}^{2a_4}, \mathrm{e}^{2a_4}, \mathrm{e}^{2a_4}, -1)\)
  in the coordinates \(x_1, \dots, x_4\) (the induced metric at
  \(\sigma = 0\) without the extra times); \(\hat h^{\mu\nu}\) its
  inverse, \(\hat h\) its determinant,
  \(\sqrt{-\hat h} = \mathrm{e}^{3a_4}\); \(\hat h^{(1)}_{\mu\nu}\) its
  first-order cosmological perturbation. Unit: dimensionless. Status:
  DERIVED.
\item
  \(\mathcal{Z}\): The fold factor,
  \(\mathcal{Z}(x_4) = V_{\mathrm{t}}(\mathrm{e}^{-a_4}b_4)^3/V_*\): the
  factor by which the size of the extra times multiplies every term of
  the fold and the kinetic terms of the massless modes of the fields of
  the bulk; \(\mathcal{Z} = 1\) when the extra times are frozen at their
  late value. Unit: dimensionless. Value: 1 in the late universe.
  Status: DERIVED.
\item
  \(D_M,\ D_\mu\): The covariant derivative: on dirac16complex
  \(D_M = \partial_M + \Omega_M\) (with
  \(D_M\bar\Psi = \partial_M\bar\Psi - \bar\Psi\Omega_M\)); on the
  fields of the fold the gauge and spin covariant derivative \(D_\mu\)
  (the partial derivative plus the spin connection of the fold and the
  gauge fields times their couplings and generators);
  \(D^\mu = \hat h^{\mu\nu}D_\nu\); on an adjoint field
  \((D_\mu X)^A = \partial_\mu X^A + g_3f^{ABC}G^B_\mu X^C\) (likewise
  for \(SU(2)\) with \(g_2\epsilon^{IJK}\)). Unit: 1/length. Status:
  STANDARD.
\item
  \(\mu\): A four-dimensional coordinate index running over the
  observable coordinates \(x_1, x_2, x_3, x_4\) of the fold (signature
  \(+,+,+,-\)); a repeated index is summed. Status: STANDARD.
\item
  \(\nu\): A four-dimensional coordinate index (as \(\mu\)). Status:
  STANDARD.
\item
  \(\Phi\): The Higgs doublet (two complex components, hypercharge
  \(1/2\)), a field of the fold. Unit: GeV. Status: AUTHOR.
\item
  \(\mu^2\): The mass parameter of the Higgs potential, \(\mu^2 < 0\)
  (it enters only as \(\mu^2\)). Unit: GeV\(^2\). Value:
  \(\mu^2 = -\lambda v^2 = -m_h^2/2 = -(88.6\ \mathrm{GeV})^2\). Status:
  AUTHOR.
\item
  \(\lambda\): The quartic coupling of the Higgs potential,
  \(\lambda > 0\). Unit: dimensionless. Value:
  \(\lambda = m_h^2/(2v^2) = 0.129\). Status: AUTHOR.
\item
  \({}^\dagger\): Hermitian conjugation: the complex conjugate of the
  transpose of a matrix or spinor; for an operator on the space of
  states, its adjoint with respect to the scalar product of that space
  (for the quantised dirac16complex the indefinite Krein scalar
  product). Status: STANDARD.
\item
  \(c_\Phi\): The Higgs-portal coupling of the left-handed half of
  dirac16complex (the right-handed half has none). Unit: GeV\(^{-5}\).
  Status: INVENTED.
\item
  \(\Psi,\ \bar\Psi\): dirac16complex: the author's 16-component complex
  spinor field of eight dimensions;
  \(\bar\Psi = \Psi^\dagger\sigma_{16}\) is its Dirac adjoint; in a term
  of the fold, \(\Psi\) is its value at \(\sigma = 0\). Unit:
  GeV\(^{7/2}\). Status: AUTHOR.
\item
  \(P_{\mathrm{L}}\): The author's left-handed chiral projector,
  \(P_{\mathrm{L}} = (\mathbf{1} - \Gamma_8)/2\). Unit: dimensionless.
  Status: AUTHOR.
\item
  \(\Theta\): The Yukawa source of the Higgs equation,
  \(\Theta = -\partial\mathcal{L}_{\mathrm{Y}}/\partial\Phi^\dagger\) (a
  doublet of fermion bilinears). Unit: GeV\(^3\). Status: STANDARD.
\item
  \(d_{\mathrm{R}}\): The right-handed down-type quarks (down, strange
  and bottom; hypercharge \(-1/3\)). Unit: GeV\(^{3/2}\). Status:
  STANDARD.
\item
  \(\mathcal{Y}^{u},\ \mathcal{Y}^{d},\ \mathcal{Y}^{e},\ \mathcal{Y}^{\nu}\):
  The Yukawa matrices (\(3 \times 3\) complex matrices in generation
  space) of the up quarks, the down quarks, the charged leptons and the
  neutrinos. Unit: dimensionless. Status: STANDARD.
\item
  \(u_{\mathrm{L}}\): The upper (up-type, left-handed) components of the
  quark doublets. Unit: GeV\(^{3/2}\). Status: STANDARD.
\item
  \(e_{\mathrm{R}}\): The right-handed electron, muon and tau (singlets,
  hypercharge \(-1\)). Unit: GeV\(^{3/2}\). Status: AUTHOR.
\item
  \(\nu_{\mathrm{L}}\): The upper (neutral, left-handed) components of
  the lepton doublets: the left-handed neutrinos. Unit: GeV\(^{3/2}\).
  Status: STANDARD.
\item
  \(d_{\mathrm{L}}\): The lower (down-type, left-handed) components of
  the quark doublets. Unit: GeV\(^{3/2}\). Status: STANDARD.
\item
  \(u_{\mathrm{R}}\): The right-handed up-type quarks (up, charm and
  top; singlets of \(SU(2)_L\), hypercharge \(2/3\)). Unit:
  GeV\(^{3/2}\). Status: STANDARD.
\item
  \(\nu_{\mathrm{R}}\): The three right-handed neutrinos (gauge
  singlets). Unit: GeV\(^{3/2}\). Status: ASSUMED.
\item
  \(e_{\mathrm{L}}\): The lower (charged, left-handed) components of the
  lepton doublets. Unit: GeV\(^{3/2}\). Status: STANDARD.
\end{itemize}

The field equation of the Higgs doublet on the fold: the variation of
IN-11, IN-15 and IN-16 with respect to \(\Phi^\dagger\), with the gauge
fields of the bulk at the fold in \(D_\mu\) (IN-13). The factor
\(\mathcal{Z}\) (IN-8) is 1 in the late universe. The Higgs portal adds
the term \(-c_\Phi\bar\Psi P_{\mathrm{L}}\Psi\,\Phi\): the scalar
density of the left-handed half of dirac16complex at the fold shifts the
mass parameter of the Higgs field. \(\Theta\) is the Yukawa source,
\(-\partial\mathcal{L}_{\mathrm{Y}}/\partial\Phi^\dagger\), written with
the components \(Q = (u_{\mathrm{L}}, d_{\mathrm{L}})\) and
\(\ell = (\nu_{\mathrm{L}}, e_{\mathrm{L}})\) of the doublets
(generation indices contracted with the Yukawa matrices). Status:
STANDARD (AUTHOR potential; INVENTED portal term).

\hypertarget{the-gauge-fields}{%
\subsection{The gauge fields}\label{the-gauge-fields}}

\[
\begin{aligned}
&\frac{1}{\sqrt{\lvert g\rvert}}\,D_M\big(\sqrt{\lvert g\rvert}\,G^{AMN}\big) = 0\qquad (\sigma \ne 0),\\
&\frac{1}{\mathsf{g}_3^2}\Big(G^{A\sigma\nu}\big\vert_{0^+} - G^{A\sigma\nu}\big\vert_{0^-}\Big)
 = -\frac{1}{V_*}\Big(J_{\mathrm{c}}^{A\nu} + \frac{1}{g_3}\,D_\mu\big[\big(\mathcal{K}_{G} - 1\big)G^{A\mu\nu}\big]\Big),\qquad
 \mathcal{K}_{G} = 1 + \sum_X c^{G}_X\,\bar\Psi P_X\Psi,\\
&J_{\mathrm{c}}^{A\nu} = \bar Q\,\hat\gamma^\nu t^AQ + \bar u_{\mathrm{R}}\,\hat\gamma^\nu t^Au_{\mathrm{R}}
 + \bar d_{\mathrm{R}}\,\hat\gamma^\nu t^Ad_{\mathrm{R}}
\end{aligned} \tag{FE-16}
\] \textbf{Symbols in (FE-16).}

\begin{itemize}
\tightlist
\item
  \(g_{MN}\): The metric tensor of the eight-dimensional spacetime
  (signature \(+,+,+,+,-,-,-,-\) in the order \(x_0, \dots, x_7\)):
  \(g_{MN}\,\mathrm{d}X^M\mathrm{d}X^N\) is the squared interval,
  \(g^{MN}\) its inverse, \(\lvert g\rvert\) the absolute value of its
  determinant (\(\sqrt{\lvert g\rvert}\,\mathrm{d}^8X\) is the proper
  volume element). Unit: dimensionless. Status: AUTHOR.
\item
  \(D_M,\ D_\mu\): The covariant derivative: on dirac16complex
  \(D_M = \partial_M + \Omega_M\) (with
  \(D_M\bar\Psi = \partial_M\bar\Psi - \bar\Psi\Omega_M\)); on the
  fields of the fold the gauge and spin covariant derivative \(D_\mu\)
  (the partial derivative plus the spin connection of the fold and the
  gauge fields times their couplings and generators);
  \(D^\mu = \hat h^{\mu\nu}D_\nu\); on an adjoint field
  \((D_\mu X)^A = \partial_\mu X^A + g_3f^{ABC}G^B_\mu X^C\) (likewise
  for \(SU(2)\) with \(g_2\epsilon^{IJK}\)). Unit: 1/length. Status:
  STANDARD.
\item
  \(M\): A coordinate index of eight dimensions (also with a numbered
  copy \(M_1, M_2, \dots\)): it runs over \(x_0, x_1, \dots, x_7\); a
  repeated index is summed. Status: STANDARD.
\item
  \(G^A_M,\ G^A_\mu,\ G^A_{MN}\): The gluon field of \(SU(3)_C\)
  (\(A = 1, \dots, 8\)): in the bulk \(G^A_M\), a field of eight
  dimensions in its geometric normalisation (the covariant derivative
  \(\partial - \mathrm{i}G\)), and on the fold the canonically
  normalised four-dimensional field \(G^A_\mu(x)\) of its massless mode,
  \(g_3G^A_\mu(x) = G^A_\mu(\sigma = 0, x)\); \(G^A_{MN}\) and
  \(G^A_{\mu\nu}\) their field strengths. Unit: GeV (field), GeV\(^2\)
  (field strength). Status: AUTHOR.
\item
  \(A\): A colour (adjoint SU(3)) index running over \(1, \dots, 8\); a
  repeated index is summed. Status: STANDARD.
\item
  \(N\): A coordinate index of eight dimensions (also
  \(N_1, N_2, \dots\)), running over \(x_0, \dots, x_7\); a repeated
  index is summed. Status: STANDARD.
\item
  \(\sigma\): The proper coordinate of the hidden direction:
  \(\mathrm{d}\sigma = \lvert\cot z\rvert\,\mathrm{d}x_0\),
  \(\sigma = \mathrm{sgn}(\pi/2 - z)\ln(\sin z)/(6H)\) on one period
  \(0 < z < \pi\); \(\sigma = 0\) is the fold \(z = \pi/2\),
  \(\sigma \to \mp\infty\) the tips \(z \to 0, \pi\); one period of
  \(x_0\) is the whole line \(-\infty < \sigma < \infty\). Unit: length.
  Status: DERIVED.
\item
  \(\mathsf{g}_1,\ \mathsf{g}_2,\ \mathsf{g}_3\): The gauge couplings of
  eight dimensions of \(U(1)_Y\), \(SU(2)_L\) and \(SU(3)_C\): the
  coefficients of the kinetic terms of the gauge fields of the bulk in
  their geometric normalisation; the four-dimensional couplings are
  \(g_k^2 = 3H\mathsf{g}_k^2/(2V_*)\). Unit: GeV\(^{-2}\)
  (\(\mathsf{g}^2\): length\(^4\)). Status: ASSUMED.
\item
  \(\sigma\): The label \(\sigma\) in \(\partial_\sigma\): the
  derivative with respect to the proper hidden coordinate \(\sigma\).
  Status: STANDARD.
\item
  \(\nu\): A four-dimensional coordinate index (as \(\mu\)). Status:
  STANDARD.
\item
  \(V_*\): The physical volume of the extra times \(x_5, x_6, x_7\) on
  the fold in the late universe,
  \(V_* = V_{\mathrm{t}}(\mathrm{e}^{-a_4}b_4)^3\) (with
  \(V_{\mathrm{t}} = \int\mathrm{d}x_5\,\mathrm{d}x_6\,\mathrm{d}x_7\)
  the coordinate volume of the extra times) at a time when
  \(\mathrm{e}^{-a_4}b_4\) is constant; it converts the
  eight-dimensional normalisation into the four-dimensional one. Unit:
  length\(^3\) (GeV\(^{-3}\)). Status: ASSUMED.
\item
  \(J_{\mathrm{c}}^{A\nu}\): The colour current of the quarks. Unit:
  GeV\(^3\). Status: STANDARD.
\item
  \(g_3\): The \(SU(3)_C\) gauge coupling of the four-dimensional world
  (the strong coupling \(g_s\)): the coupling of the massless mode of
  the gluon field of the bulk, \(g_3^2 = 3H\mathsf{g}_3^2/(2V_*)\).
  Unit: dimensionless. Value: \(g_3(m_Z) = 1.2177\)
  (\(g_3^2/(4\pi) = 0.1180\)). Status: STANDARD.
\item
  \(\mu\): A four-dimensional coordinate index running over the
  observable coordinates \(x_1, x_2, x_3, x_4\) of the fold (signature
  \(+,+,+,-\)); a repeated index is summed. Status: STANDARD.
\item
  \(\mathcal{K}_{G},\ \mathcal{K}_{W},\ \mathcal{K}_{B}\): The portal
  factors of the gauge kinetic terms,
  \(\mathcal{K}_{G} = 1 + \sum_Xc^{G}_X\bar\Psi P_X\Psi\) (likewise for
  \(W\) and \(B\)): the gauge couplings at the fold are divided by the
  square root of the portal factor. Unit: dimensionless. Value: 1 where
  dirac16complex vanishes. Status: DERIVED.
\item
  \(X\): The label of a half of dirac16complex:
  \(X \in \{\mathrm{L}, \mathrm{R}\}\) (left-handed or right-handed
  half); a sum over \(X\) runs over both. Status: STANDARD.
\item
  \(c^{G}_X\): The gluon-kinetic portal coupling of the half \(X\) of
  dirac16complex. Unit: GeV\(^{-7}\). Status: INVENTED.
\item
  \(\Psi,\ \bar\Psi\): dirac16complex: the author's 16-component complex
  spinor field of eight dimensions;
  \(\bar\Psi = \Psi^\dagger\sigma_{16}\) is its Dirac adjoint; in a term
  of the fold, \(\Psi\) is its value at \(\sigma = 0\). Unit:
  GeV\(^{7/2}\). Status: AUTHOR.
\item
  \(P_X\): The chiral projector of the half \(X\):
  \(P_X = P_{\mathrm{L}}\) for \(X = \mathrm{L}\), \(P_{\mathrm{R}}\)
  for \(X = \mathrm{R}\). Unit: dimensionless. Status: AUTHOR.
\item
  \(Q,\ \bar Q\): The left-handed quark doublets of the three
  generations, \(Q = (u_{\mathrm{L}}, d_{\mathrm{L}})\) (colour
  triplets, hypercharge \(1/6\)). Unit: GeV\(^{3/2}\). Status: AUTHOR.
\item
  \(\hat\gamma^\mu\): The four-dimensional complex \(4 \times 4\) Dirac
  matrices of the fold: \(\hat\gamma^{(a)}\) along the frame directions
  of the fold and \(\hat\gamma^\mu\) along its coordinates (the frame
  matrices contracted with the inverse frame field of the fold),
  \(\{\hat\gamma^\mu, \hat\gamma^\nu\} = -2\hat h^{\mu\nu}\mathbf{1}\);
  \(\hat\gamma_\mu = \hat h_{\mu\nu}\hat\gamma^\nu\). Unit:
  dimensionless. Status: STANDARD.
\item
  \(t^A\): The generators of \(SU(3)\) on colour triplets, \(t^A\): half
  the Gell-Mann matrices, with \(\mathrm{tr}(t^At^B) = \delta^{AB}/2\).
  Unit: dimensionless. Status: STANDARD.
\item
  \(u_{\mathrm{R}}\): The right-handed up-type quarks (up, charm and
  top; singlets of \(SU(2)_L\), hypercharge \(2/3\)). Unit:
  GeV\(^{3/2}\). Status: STANDARD.
\item
  \(d_{\mathrm{R}}\): The right-handed down-type quarks (down, strange
  and bottom; hypercharge \(-1/3\)). Unit: GeV\(^{3/2}\). Status:
  STANDARD.
\end{itemize}

The Yang-Mills equation of the gluons in the bulk (first line; \(D_M\)
with the gauge connection in the adjoint representation) and at the fold
(second line): the variation of IN-7, IN-14 and IN-16 with respect to
the gluon field of the bulk. The quarks on the fold couple to the value
of the field at \(\sigma = 0\) (IN-13), which puts their colour current
\(J_{\mathrm{c}}^{A\nu}\), times \(\delta_{\mathrm{D}}(\sigma)/V_*\), on
the right-hand side of the bulk equation; integrated across the fold it
makes the component \(G^{A\sigma\nu}\) of the field strength jump by
\(-\mathsf{g}_3^2J_{\mathrm{c}}^{A\nu}/V_*\) (convention of the
signature: \(\nabla_M F^{MN} = -J^N\) for a Lagrangian
\(-\frac14F^2 + A_MJ^M\)). The gluon-kinetic portals add a kinetic term
on the fold (\(\mathcal{K}_{G} - 1\), with the canonical
four-dimensional field strength \(G^{A\mu\nu}\) of IN-13 and the
covariant derivative of the fold). For the massless mode (the profile 1,
FE-74) the integral of the first line over \(\sigma\) with the weight of
the metric, together with the jump, gives the Yang-Mills equation of the
fold,
\(\hat\nabla_\mu(\mathcal{K}_{G}G^{A\mu\nu}) + g_3f^{ABC}G^B_\mu\mathcal{K}_{G}G^{C\mu\nu} = -g_3J_{\mathrm{c}}^{A\nu}\)
(in UM-3 the field equation of the gluons on the fold), with
\(g_3^2 = 3H\mathsf{g}_3^2/(2V_*)\); the massive modes add the
corrections of FE-79 and FE-80. Status: STANDARD (DERIVED for the
reduction to the fold, \texttt{um4\_bulk.py}, B2; INVENTED portal
factor).

\[
\begin{aligned}
&\frac{1}{\sqrt{\lvert g\rvert}}\,D_M\big(\sqrt{\lvert g\rvert}\,W^{IMN}\big) = 0\qquad (\sigma \ne 0),\\
&\frac{1}{\mathsf{g}_2^2}\Big(W^{I\sigma\nu}\big\vert_{0^+} - W^{I\sigma\nu}\big\vert_{0^-}\Big)
 = -\frac{1}{V_*}\Big(J_{\mathrm{w}}^{I\nu} + J_\Phi^{I\nu} + \frac{1}{g_2}\,D_\mu\big[\big(\mathcal{K}_{W} - 1\big)W^{I\mu\nu}\big]\Big),\\
&\mathcal{K}_{W} = 1 + \sum_X c^{W}_X\,\bar\Psi P_X\Psi,\qquad J_{\mathrm{w}}^{I\nu} = \bar Q\,\hat\gamma^\nu\frac{\tau^I}{2}Q + \bar\ell\,\hat\gamma^\nu\frac{\tau^I}{2}\ell,\qquad
 J_\Phi^{I\nu} = \mathrm{i}\Big[(D^\nu\Phi)^\dagger\frac{\tau^I}{2}\Phi - \Phi^\dagger\frac{\tau^I}{2}D^\nu\Phi\Big]
\end{aligned} \tag{FE-17}
\] \textbf{Symbols in (FE-17).}

\begin{itemize}
\tightlist
\item
  \(g_{MN}\): The metric tensor of the eight-dimensional spacetime
  (signature \(+,+,+,+,-,-,-,-\) in the order \(x_0, \dots, x_7\)):
  \(g_{MN}\,\mathrm{d}X^M\mathrm{d}X^N\) is the squared interval,
  \(g^{MN}\) its inverse, \(\lvert g\rvert\) the absolute value of its
  determinant (\(\sqrt{\lvert g\rvert}\,\mathrm{d}^8X\) is the proper
  volume element). Unit: dimensionless. Status: AUTHOR.
\item
  \(D_M,\ D_\mu\): The covariant derivative: on dirac16complex
  \(D_M = \partial_M + \Omega_M\) (with
  \(D_M\bar\Psi = \partial_M\bar\Psi - \bar\Psi\Omega_M\)); on the
  fields of the fold the gauge and spin covariant derivative \(D_\mu\)
  (the partial derivative plus the spin connection of the fold and the
  gauge fields times their couplings and generators);
  \(D^\mu = \hat h^{\mu\nu}D_\nu\); on an adjoint field
  \((D_\mu X)^A = \partial_\mu X^A + g_3f^{ABC}G^B_\mu X^C\) (likewise
  for \(SU(2)\) with \(g_2\epsilon^{IJK}\)). Unit: 1/length. Status:
  STANDARD.
\item
  \(M\): A coordinate index of eight dimensions (also with a numbered
  copy \(M_1, M_2, \dots\)): it runs over \(x_0, x_1, \dots, x_7\); a
  repeated index is summed. Status: STANDARD.
\item
  \(W^I_M,\ W^I_\mu,\ W^I_{MN}\): The gauge field of \(SU(2)_L\)
  (\(I = 1, 2, 3\)): in the bulk \(W^I_M\) (geometric normalisation), on
  the fold the canonically normalised four-dimensional field
  \(W^I_\mu(x)\) of its massless mode,
  \(g_2W^I_\mu(x) = W^I_\mu(\sigma = 0, x)\); \(W^I_{MN}\),
  \(W^I_{\mu\nu}\) their field strengths;
  \(W^\pm_\mu = (W^1_\mu \mp \mathrm{i}W^2_\mu)/\sqrt2\). Unit: GeV
  (field), GeV\(^2\) (field strength). Status: AUTHOR.
\item
  \(I\): A weak-isospin (adjoint SU(2)) index running over \(1, 2, 3\);
  a repeated index is summed. Status: STANDARD.
\item
  \(N\): A coordinate index of eight dimensions (also
  \(N_1, N_2, \dots\)), running over \(x_0, \dots, x_7\); a repeated
  index is summed. Status: STANDARD.
\item
  \(\sigma\): The proper coordinate of the hidden direction:
  \(\mathrm{d}\sigma = \lvert\cot z\rvert\,\mathrm{d}x_0\),
  \(\sigma = \mathrm{sgn}(\pi/2 - z)\ln(\sin z)/(6H)\) on one period
  \(0 < z < \pi\); \(\sigma = 0\) is the fold \(z = \pi/2\),
  \(\sigma \to \mp\infty\) the tips \(z \to 0, \pi\); one period of
  \(x_0\) is the whole line \(-\infty < \sigma < \infty\). Unit: length.
  Status: DERIVED.
\item
  \(\mathsf{g}_1,\ \mathsf{g}_2,\ \mathsf{g}_3\): The gauge couplings of
  eight dimensions of \(U(1)_Y\), \(SU(2)_L\) and \(SU(3)_C\): the
  coefficients of the kinetic terms of the gauge fields of the bulk in
  their geometric normalisation; the four-dimensional couplings are
  \(g_k^2 = 3H\mathsf{g}_k^2/(2V_*)\). Unit: GeV\(^{-2}\)
  (\(\mathsf{g}^2\): length\(^4\)). Status: ASSUMED.
\item
  \(\sigma\): The label \(\sigma\) in \(\partial_\sigma\): the
  derivative with respect to the proper hidden coordinate \(\sigma\).
  Status: STANDARD.
\item
  \(\nu\): A four-dimensional coordinate index (as \(\mu\)). Status:
  STANDARD.
\item
  \(V_*\): The physical volume of the extra times \(x_5, x_6, x_7\) on
  the fold in the late universe,
  \(V_* = V_{\mathrm{t}}(\mathrm{e}^{-a_4}b_4)^3\) (with
  \(V_{\mathrm{t}} = \int\mathrm{d}x_5\,\mathrm{d}x_6\,\mathrm{d}x_7\)
  the coordinate volume of the extra times) at a time when
  \(\mathrm{e}^{-a_4}b_4\) is constant; it converts the
  eight-dimensional normalisation into the four-dimensional one. Unit:
  length\(^3\) (GeV\(^{-3}\)). Status: ASSUMED.
\item
  \(J_{\mathrm{w}}^{I\nu}\): The weak-isospin current of the left-handed
  fermion doublets. Unit: GeV\(^3\). Status: STANDARD.
\item
  \(J_\Phi^{I\nu}\): The weak-isospin current of the Higgs doublet (it
  contains the mass terms of the \(W\) fields); \(J_\Phi^{2\,x_4}\) its
  isospin charge density with \(I = 2\). Unit: GeV\(^3\). Status:
  STANDARD.
\item
  \(g_2\): The \(SU(2)_L\) gauge coupling of the four-dimensional world
  (usually \(g\)): the coupling of the massless mode of the \(SU(2)_L\)
  field of the bulk, \(g_2^2 = 3H\mathsf{g}_2^2/(2V_*)\). Unit:
  dimensionless. Value: \(g_2(m_Z) = 0.6517\). Status: STANDARD.
\item
  \(\mu\): A four-dimensional coordinate index running over the
  observable coordinates \(x_1, x_2, x_3, x_4\) of the fold (signature
  \(+,+,+,-\)); a repeated index is summed. Status: STANDARD.
\item
  \(\mathcal{K}_{G},\ \mathcal{K}_{W},\ \mathcal{K}_{B}\): The portal
  factors of the gauge kinetic terms,
  \(\mathcal{K}_{G} = 1 + \sum_Xc^{G}_X\bar\Psi P_X\Psi\) (likewise for
  \(W\) and \(B\)): the gauge couplings at the fold are divided by the
  square root of the portal factor. Unit: dimensionless. Value: 1 where
  dirac16complex vanishes. Status: DERIVED.
\item
  \(X\): The label of a half of dirac16complex:
  \(X \in \{\mathrm{L}, \mathrm{R}\}\) (left-handed or right-handed
  half); a sum over \(X\) runs over both. Status: STANDARD.
\item
  \(c^{W}_X\): The \(W\)-kinetic portal coupling of the half \(X\).
  Unit: GeV\(^{-7}\). Status: INVENTED.
\item
  \(\Psi,\ \bar\Psi\): dirac16complex: the author's 16-component complex
  spinor field of eight dimensions;
  \(\bar\Psi = \Psi^\dagger\sigma_{16}\) is its Dirac adjoint; in a term
  of the fold, \(\Psi\) is its value at \(\sigma = 0\). Unit:
  GeV\(^{7/2}\). Status: AUTHOR.
\item
  \(P_X\): The chiral projector of the half \(X\):
  \(P_X = P_{\mathrm{L}}\) for \(X = \mathrm{L}\), \(P_{\mathrm{R}}\)
  for \(X = \mathrm{R}\). Unit: dimensionless. Status: AUTHOR.
\item
  \(Q,\ \bar Q\): The left-handed quark doublets of the three
  generations, \(Q = (u_{\mathrm{L}}, d_{\mathrm{L}})\) (colour
  triplets, hypercharge \(1/6\)). Unit: GeV\(^{3/2}\). Status: AUTHOR.
\item
  \(\hat\gamma^\mu\): The four-dimensional complex \(4 \times 4\) Dirac
  matrices of the fold: \(\hat\gamma^{(a)}\) along the frame directions
  of the fold and \(\hat\gamma^\mu\) along its coordinates (the frame
  matrices contracted with the inverse frame field of the fold),
  \(\{\hat\gamma^\mu, \hat\gamma^\nu\} = -2\hat h^{\mu\nu}\mathbf{1}\);
  \(\hat\gamma_\mu = \hat h_{\mu\nu}\hat\gamma^\nu\). Unit:
  dimensionless. Status: STANDARD.
\item
  \(\tau^I\): The Pauli matrices acting on weak doublets
  (\(\tau^1, \tau^2, \tau^3\); \(\tau^2\) is the second one, not a
  square). Unit: dimensionless. Status: STANDARD.
\item
  \(\ell,\ \bar\ell\): The left-handed lepton doublets of the three
  generations, \(\ell = (\nu_{\mathrm{L}}, e_{\mathrm{L}})\)
  (hypercharge \(-1/2\)). Unit: GeV\(^{3/2}\). Status: AUTHOR.
\item
  \(\mathrm{i}\): The imaginary unit, \(\mathrm{i}^2 = -1\). Unit:
  dimensionless. Status: STANDARD.
\item
  \(\Phi\): The Higgs doublet (two complex components, hypercharge
  \(1/2\)), a field of the fold. Unit: GeV. Status: AUTHOR.
\item
  \({}^\dagger\): Hermitian conjugation: the complex conjugate of the
  transpose of a matrix or spinor; for an operator on the space of
  states, its adjoint with respect to the scalar product of that space
  (for the quantised dirac16complex the indefinite Krein scalar
  product). Status: STANDARD.
\end{itemize}

The Yang-Mills equation of the \(SU(2)_L\) fields in the bulk and at the
fold. The sources on the fold are the weak-isospin current of the
left-handed doublets and the isospin current of the Higgs doublet, which
contains the mass terms of the \(W\) fields: in the vacuum
\(J_\Phi^{I\nu}\) contains \(-g_2(v^2 + \mathcal{A}^2)W^{I\nu}/4\), a
mass term of the field of the bulk concentrated on the fold, which makes
the W boson a resonance of the bulk (FE-82). For the massless mode the
reduction gives the field equation of UM-3 on the fold,
\(\hat\nabla_\mu(\mathcal{K}_{W}W^{I\mu\nu}) + \dots = -g_2(J_{\mathrm{w}}^{I\nu} + J_\Phi^{I\nu})\),
with \(g_2^2 = 3H\mathsf{g}_2^2/(2V_*)\). Status: STANDARD (INVENTED
portal factor).

\[
\begin{aligned}
&\frac{1}{\sqrt{\lvert g\rvert}}\,\partial_M\big(\sqrt{\lvert g\rvert}\,B^{MN}\big) = 0\qquad (\sigma \ne 0),\\
&\frac{1}{\mathsf{g}_1^2}\Big(B^{\sigma\nu}\big\vert_{0^+} - B^{\sigma\nu}\big\vert_{0^-}\Big)
 = -\frac{1}{V_*}\Big(J_Y^\nu + \frac{1}{g_1}\,\partial_\mu\Big[\big(\mathcal{K}_{B} - 1\big)B^{\mu\nu} + \sum_X d_X\,\mathcal{D}_X^{\mu\nu}\Big]\Big),\qquad
 \mathcal{K}_{B} = 1 + \sum_X c^{B}_X\,\bar\Psi P_X\Psi,\\
&J_Y^\nu = \sum_\psi Y_\psi\,\bar\psi\,\hat\gamma^\nu\psi + \frac{\mathrm{i}}{2}\big[(D^\nu\Phi)^\dagger\Phi - \Phi^\dagger D^\nu\Phi\big]
\end{aligned} \tag{FE-18}
\] \textbf{Symbols in (FE-18).}

\begin{itemize}
\tightlist
\item
  \(g_{MN}\): The metric tensor of the eight-dimensional spacetime
  (signature \(+,+,+,+,-,-,-,-\) in the order \(x_0, \dots, x_7\)):
  \(g_{MN}\,\mathrm{d}X^M\mathrm{d}X^N\) is the squared interval,
  \(g^{MN}\) its inverse, \(\lvert g\rvert\) the absolute value of its
  determinant (\(\sqrt{\lvert g\rvert}\,\mathrm{d}^8X\) is the proper
  volume element). Unit: dimensionless. Status: AUTHOR.
\item
  \(\partial\): A partial derivative: the derivative with respect to one
  variable at fixed values of the others (\(\partial_\sigma\) with
  respect to \(\sigma\); \(\partial_M\) and \(\partial_\mu\) with
  respect to the coordinate that carries the same index, of the eight
  dimensions or of the fold). Status: STANDARD.
\item
  \(M\): A coordinate index of eight dimensions (also with a numbered
  copy \(M_1, M_2, \dots\)): it runs over \(x_0, x_1, \dots, x_7\); a
  repeated index is summed. Status: STANDARD.
\item
  \(B_M,\ B_\mu,\ B_{MN}\): The hypercharge gauge field of \(U(1)_Y\):
  in the bulk \(B_M\) (geometric normalisation), on the fold the
  canonically normalised four-dimensional field \(B_\mu(x)\) of its
  massless mode, \(g_1B_\mu(x) = B_\mu(\sigma = 0, x)\); \(B_{MN}\),
  \(B_{\mu\nu}\) their field strengths. Unit: GeV (field), GeV\(^2\)
  (field strength). Status: AUTHOR.
\item
  \(N\): A coordinate index of eight dimensions (also
  \(N_1, N_2, \dots\)), running over \(x_0, \dots, x_7\); a repeated
  index is summed. Status: STANDARD.
\item
  \(\sigma\): The proper coordinate of the hidden direction:
  \(\mathrm{d}\sigma = \lvert\cot z\rvert\,\mathrm{d}x_0\),
  \(\sigma = \mathrm{sgn}(\pi/2 - z)\ln(\sin z)/(6H)\) on one period
  \(0 < z < \pi\); \(\sigma = 0\) is the fold \(z = \pi/2\),
  \(\sigma \to \mp\infty\) the tips \(z \to 0, \pi\); one period of
  \(x_0\) is the whole line \(-\infty < \sigma < \infty\). Unit: length.
  Status: DERIVED.
\item
  \(\mathsf{g}_1,\ \mathsf{g}_2,\ \mathsf{g}_3\): The gauge couplings of
  eight dimensions of \(U(1)_Y\), \(SU(2)_L\) and \(SU(3)_C\): the
  coefficients of the kinetic terms of the gauge fields of the bulk in
  their geometric normalisation; the four-dimensional couplings are
  \(g_k^2 = 3H\mathsf{g}_k^2/(2V_*)\). Unit: GeV\(^{-2}\)
  (\(\mathsf{g}^2\): length\(^4\)). Status: ASSUMED.
\item
  \(\sigma\): The label \(\sigma\) in \(\partial_\sigma\): the
  derivative with respect to the proper hidden coordinate \(\sigma\).
  Status: STANDARD.
\item
  \(\nu\): A four-dimensional coordinate index (as \(\mu\)). Status:
  STANDARD.
\item
  \(V_*\): The physical volume of the extra times \(x_5, x_6, x_7\) on
  the fold in the late universe,
  \(V_* = V_{\mathrm{t}}(\mathrm{e}^{-a_4}b_4)^3\) (with
  \(V_{\mathrm{t}} = \int\mathrm{d}x_5\,\mathrm{d}x_6\,\mathrm{d}x_7\)
  the coordinate volume of the extra times) at a time when
  \(\mathrm{e}^{-a_4}b_4\) is constant; it converts the
  eight-dimensional normalisation into the four-dimensional one. Unit:
  length\(^3\) (GeV\(^{-3}\)). Status: ASSUMED.
\item
  \(J_Y^\nu\): The hypercharge current of the fermions and of the Higgs
  doublet. Unit: GeV\(^3\). Status: STANDARD.
\item
  \(g_1\): The \(U(1)_Y\) gauge coupling of the four-dimensional world
  (usually \(g'\); not normalised for grand unification): the coupling
  of the massless mode of the hypercharge field of the bulk,
  \(g_1^2 = 3H\mathsf{g}_1^2/(2V_*)\). Unit: dimensionless. Value:
  \(g_1(m_Z) = 0.3574\). Status: STANDARD.
\item
  \(\mu\): A four-dimensional coordinate index running over the
  observable coordinates \(x_1, x_2, x_3, x_4\) of the fold (signature
  \(+,+,+,-\)); a repeated index is summed. Status: STANDARD.
\item
  \(\mathcal{K}_{G},\ \mathcal{K}_{W},\ \mathcal{K}_{B}\): The portal
  factors of the gauge kinetic terms,
  \(\mathcal{K}_{G} = 1 + \sum_Xc^{G}_X\bar\Psi P_X\Psi\) (likewise for
  \(W\) and \(B\)): the gauge couplings at the fold are divided by the
  square root of the portal factor. Unit: dimensionless. Value: 1 where
  dirac16complex vanishes. Status: DERIVED.
\item
  \(X\): The label of a half of dirac16complex:
  \(X \in \{\mathrm{L}, \mathrm{R}\}\) (left-handed or right-handed
  half); a sum over \(X\) runs over both. Status: STANDARD.
\item
  \(d_X\): The dipole-portal coupling of the half \(X\) of
  dirac16complex with the hypercharge field strength. Unit:
  GeV\(^{-5}\). Status: INVENTED.
\item
  \(\mathcal{D}_X^{\mu\nu}\): The dipole (antisymmetric tensor) density
  of the half \(X\) of dirac16complex at the fold,
  \(\mathcal{D}_X^{\mu\nu} = \mathrm{i}\bar\Psi\Sigma^{\mu\nu}P_X\Psi\)
  (real, DERIVED). Unit: GeV\(^7\). Status: DERIVED.
\item
  \(c^{B}_X\): The \(B\)-kinetic portal coupling of the half \(X\).
  Unit: GeV\(^{-7}\). Status: INVENTED.
\item
  \(\Psi,\ \bar\Psi\): dirac16complex: the author's 16-component complex
  spinor field of eight dimensions;
  \(\bar\Psi = \Psi^\dagger\sigma_{16}\) is its Dirac adjoint; in a term
  of the fold, \(\Psi\) is its value at \(\sigma = 0\). Unit:
  GeV\(^{7/2}\). Status: AUTHOR.
\item
  \(P_X\): The chiral projector of the half \(X\):
  \(P_X = P_{\mathrm{L}}\) for \(X = \mathrm{L}\), \(P_{\mathrm{R}}\)
  for \(X = \mathrm{R}\). Unit: dimensionless. Status: AUTHOR.
\item
  \(\psi\): The label of a Standard-Model fermion field in a sum:
  \(\psi\) runs over
  \(Q, u_{\mathrm{R}}, d_{\mathrm{R}}, \ell, e_{\mathrm{R}}, \nu_{\mathrm{R}}\).
  Status: STANDARD.
\item
  \(Y_f\): The hypercharge of a field, the charge of its coupling
  \(g_1Y_fB_\mu\) to the hypercharge field in the covariant derivative
  (\(1/6\), \(2/3\), \(-1/3\), \(-1/2\), \(-1\), \(0\) for
  \(Q, u_{\mathrm{R}}, d_{\mathrm{R}}, \ell, e_{\mathrm{R}}, \nu_{\mathrm{R}}\),
  \(1/2\) for \(\Phi\)). Unit: dimensionless. Status: STANDARD.
\item
  \(\psi,\ \bar\psi\): A Standard-Model fermion field (a
  four-dimensional Weyl field of the fold with three generations,
  written as a four-component spinor of fixed chirality);
  \(\bar\psi = \psi^\dagger\beta\); \(\psi_f\) the fields of the kind
  \(f\). Unit: GeV\(^{3/2}\). Status: STANDARD.
\item
  \(\hat\gamma^\mu\): The four-dimensional complex \(4 \times 4\) Dirac
  matrices of the fold: \(\hat\gamma^{(a)}\) along the frame directions
  of the fold and \(\hat\gamma^\mu\) along its coordinates (the frame
  matrices contracted with the inverse frame field of the fold),
  \(\{\hat\gamma^\mu, \hat\gamma^\nu\} = -2\hat h^{\mu\nu}\mathbf{1}\);
  \(\hat\gamma_\mu = \hat h_{\mu\nu}\hat\gamma^\nu\). Unit:
  dimensionless. Status: STANDARD.
\item
  \(\mathrm{i}\): The imaginary unit, \(\mathrm{i}^2 = -1\). Unit:
  dimensionless. Status: STANDARD.
\item
  \(D_M,\ D_\mu\): The covariant derivative: on dirac16complex
  \(D_M = \partial_M + \Omega_M\) (with
  \(D_M\bar\Psi = \partial_M\bar\Psi - \bar\Psi\Omega_M\)); on the
  fields of the fold the gauge and spin covariant derivative \(D_\mu\)
  (the partial derivative plus the spin connection of the fold and the
  gauge fields times their couplings and generators);
  \(D^\mu = \hat h^{\mu\nu}D_\nu\); on an adjoint field
  \((D_\mu X)^A = \partial_\mu X^A + g_3f^{ABC}G^B_\mu X^C\) (likewise
  for \(SU(2)\) with \(g_2\epsilon^{IJK}\)). Unit: 1/length. Status:
  STANDARD.
\item
  \(\Phi\): The Higgs doublet (two complex components, hypercharge
  \(1/2\)), a field of the fold. Unit: GeV. Status: AUTHOR.
\item
  \({}^\dagger\): Hermitian conjugation: the complex conjugate of the
  transpose of a matrix or spinor; for an operator on the space of
  states, its adjoint with respect to the scalar product of that space
  (for the quantised dirac16complex the indefinite Krein scalar
  product). Status: STANDARD.
\end{itemize}

The Maxwell equation of the hypercharge field in the bulk and at the
fold. The dipole portal adds the dipole density of dirac16complex at the
fold to the field strength, as a polarisation adds to the electric field
in a medium: dirac16complex at the fold carries electric and magnetic
dipole moments. The source on the fold is the hypercharge current of the
fermions and of the Higgs doublet. The photon is the combination of
\(W^3\) and \(B\) that the Higgs vacuum leaves massless (FE-45); in the
bulk it obeys the free Maxwell equations of the first lines of FE-17 and
FE-18, so a photon that leaves the fold travels through the bulk at the
speed of light (FE-74). Status: STANDARD (INVENTED portal terms).

\hypertarget{the-quarks-and-leptons-1}{%
\subsection{The quarks and leptons}\label{the-quarks-and-leptons-1}}

\[
\mathrm{i}\hat\gamma^\mu\big(D_\mu + \mathrm{i}\,\mathcal{V}^{Q}_\mu\big)Q = \mathcal{Y}^{u}\tilde\Phi\,u_{\mathrm{R}}
 + \mathcal{Y}^{d}\Phi\,d_{\mathrm{R}},\qquad \mathcal{V}^{f}_\mu = \sum_X c^{f}_X\,j_{X\mu} \tag{FE-19}
\] \textbf{Symbols in (FE-19).}

\begin{itemize}
\tightlist
\item
  \(\mathrm{i}\): The imaginary unit, \(\mathrm{i}^2 = -1\). Unit:
  dimensionless. Status: STANDARD.
\item
  \(\hat\gamma^\mu\): The four-dimensional complex \(4 \times 4\) Dirac
  matrices of the fold: \(\hat\gamma^{(a)}\) along the frame directions
  of the fold and \(\hat\gamma^\mu\) along its coordinates (the frame
  matrices contracted with the inverse frame field of the fold),
  \(\{\hat\gamma^\mu, \hat\gamma^\nu\} = -2\hat h^{\mu\nu}\mathbf{1}\);
  \(\hat\gamma_\mu = \hat h_{\mu\nu}\hat\gamma^\nu\). Unit:
  dimensionless. Status: STANDARD.
\item
  \(\mu\): A four-dimensional coordinate index running over the
  observable coordinates \(x_1, x_2, x_3, x_4\) of the fold (signature
  \(+,+,+,-\)); a repeated index is summed. Status: STANDARD.
\item
  \(D_M,\ D_\mu\): The covariant derivative: on dirac16complex
  \(D_M = \partial_M + \Omega_M\) (with
  \(D_M\bar\Psi = \partial_M\bar\Psi - \bar\Psi\Omega_M\)); on the
  fields of the fold the gauge and spin covariant derivative \(D_\mu\)
  (the partial derivative plus the spin connection of the fold and the
  gauge fields times their couplings and generators);
  \(D^\mu = \hat h^{\mu\nu}D_\nu\); on an adjoint field
  \((D_\mu X)^A = \partial_\mu X^A + g_3f^{ABC}G^B_\mu X^C\) (likewise
  for \(SU(2)\) with \(g_2\epsilon^{IJK}\)). Unit: 1/length. Status:
  STANDARD.
\item
  \(\mathcal{V}^{f}_\mu\): The vector potential that dirac16complex
  exerts on the Standard-Model fermions of kind \(f\) through the
  current portals, \(\mathcal{V}^{f}_\mu = \sum_Xc^{f}_Xj_{X\mu}\)
  (\(\mathcal{V}^{Q}, \mathcal{V}^{u}, \dots, \mathcal{V}^{\nu}\) for
  the six kinds); flavour blind. Unit: GeV. Status: DERIVED.
\item
  \(Q,\ \bar Q\): The left-handed quark doublets of the three
  generations, \(Q = (u_{\mathrm{L}}, d_{\mathrm{L}})\) (colour
  triplets, hypercharge \(1/6\)). Unit: GeV\(^{3/2}\). Status: AUTHOR.
\item
  \(\mathcal{Y}^{u},\ \mathcal{Y}^{d},\ \mathcal{Y}^{e},\ \mathcal{Y}^{\nu}\):
  The Yukawa matrices (\(3 \times 3\) complex matrices in generation
  space) of the up quarks, the down quarks, the charged leptons and the
  neutrinos. Unit: dimensionless. Status: STANDARD.
\item
  \(\tilde\Phi\): The conjugate Higgs doublet,
  \(\tilde\Phi = \mathrm{i}\tau^2\Phi^*\) (hypercharge \(-1/2\)). Unit:
  GeV. Status: STANDARD.
\item
  \(u_{\mathrm{R}}\): The right-handed up-type quarks (up, charm and
  top; singlets of \(SU(2)_L\), hypercharge \(2/3\)). Unit:
  GeV\(^{3/2}\). Status: STANDARD.
\item
  \(\Phi\): The Higgs doublet (two complex components, hypercharge
  \(1/2\)), a field of the fold. Unit: GeV. Status: AUTHOR.
\item
  \(d_{\mathrm{R}}\): The right-handed down-type quarks (down, strange
  and bottom; hypercharge \(-1/3\)). Unit: GeV\(^{3/2}\). Status:
  STANDARD.
\item
  \(X\): The label of a half of dirac16complex:
  \(X \in \{\mathrm{L}, \mathrm{R}\}\) (left-handed or right-handed
  half); a sum over \(X\) runs over both. Status: STANDARD.
\item
  \(c^{f}_X\): The current-portal coupling of the half \(X\) of
  dirac16complex with the fermion-number current of the multiplets of
  kind \(f\) (flavour blind); \(c^{\ell}_X\) and \(c^{\nu}_X\) are those
  of the lepton doublets and of the right-handed neutrinos. Unit:
  GeV\(^{-7}\). Status: INVENTED.
\item
  \(j_X^\mu\): The convective current of the half \(X\) of
  dirac16complex,
  \(j_X^\mu = \mathrm{i}(\bar\Psi P_XD^\mu\Psi - (D^\mu\bar\Psi)P_X\Psi)\)
  (real and joining the half with itself, DERIVED). Unit: GeV\(^8\).
  Status: DERIVED.
\item
  \(f\): The label of a kind of Standard-Model multiplet:
  \(f \in \{Q, u, d, \ell, e, \nu\}\) (quark doublets, up-quark
  singlets, down-quark singlets, lepton doublets, charged-lepton
  singlets, right-handed neutrinos); a sum over \(f\) runs over the six.
  Status: STANDARD.
\end{itemize}

The Dirac equation of the quark doublets: the variation of IN-14, IN-15
and IN-16 with respect to \(\bar Q\). The current portals act on every
Standard-Model fermion of kind \(f\) as a vector potential
\(\mathcal{V}^{f}_\mu\) built from the convective currents of the two
halves of dirac16complex at the fold; it is the same for the three
generations (flavour blind). Status: STANDARD (INVENTED potential).

\[
\mathrm{i}\hat\gamma^\mu\big(D_\mu + \mathrm{i}\,\mathcal{V}^{u}_\mu\big)u_{\mathrm{R}} = (\mathcal{Y}^{u})^\dagger\tilde\Phi^\dagger Q \tag{FE-20}
\] \textbf{Symbols in (FE-20).}

\begin{itemize}
\tightlist
\item
  \(\mathrm{i}\): The imaginary unit, \(\mathrm{i}^2 = -1\). Unit:
  dimensionless. Status: STANDARD.
\item
  \(\hat\gamma^\mu\): The four-dimensional complex \(4 \times 4\) Dirac
  matrices of the fold: \(\hat\gamma^{(a)}\) along the frame directions
  of the fold and \(\hat\gamma^\mu\) along its coordinates (the frame
  matrices contracted with the inverse frame field of the fold),
  \(\{\hat\gamma^\mu, \hat\gamma^\nu\} = -2\hat h^{\mu\nu}\mathbf{1}\);
  \(\hat\gamma_\mu = \hat h_{\mu\nu}\hat\gamma^\nu\). Unit:
  dimensionless. Status: STANDARD.
\item
  \(\mu\): A four-dimensional coordinate index running over the
  observable coordinates \(x_1, x_2, x_3, x_4\) of the fold (signature
  \(+,+,+,-\)); a repeated index is summed. Status: STANDARD.
\item
  \(D_M,\ D_\mu\): The covariant derivative: on dirac16complex
  \(D_M = \partial_M + \Omega_M\) (with
  \(D_M\bar\Psi = \partial_M\bar\Psi - \bar\Psi\Omega_M\)); on the
  fields of the fold the gauge and spin covariant derivative \(D_\mu\)
  (the partial derivative plus the spin connection of the fold and the
  gauge fields times their couplings and generators);
  \(D^\mu = \hat h^{\mu\nu}D_\nu\); on an adjoint field
  \((D_\mu X)^A = \partial_\mu X^A + g_3f^{ABC}G^B_\mu X^C\) (likewise
  for \(SU(2)\) with \(g_2\epsilon^{IJK}\)). Unit: 1/length. Status:
  STANDARD.
\item
  \(\mathcal{V}^{f}_\mu\): The vector potential that dirac16complex
  exerts on the Standard-Model fermions of kind \(f\) through the
  current portals, \(\mathcal{V}^{f}_\mu = \sum_Xc^{f}_Xj_{X\mu}\)
  (\(\mathcal{V}^{Q}, \mathcal{V}^{u}, \dots, \mathcal{V}^{\nu}\) for
  the six kinds); flavour blind. Unit: GeV. Status: DERIVED.
\item
  \(u_{\mathrm{R}}\): The right-handed up-type quarks (up, charm and
  top; singlets of \(SU(2)_L\), hypercharge \(2/3\)). Unit:
  GeV\(^{3/2}\). Status: STANDARD.
\item
  \(\mathcal{Y}^{u},\ \mathcal{Y}^{d},\ \mathcal{Y}^{e},\ \mathcal{Y}^{\nu}\):
  The Yukawa matrices (\(3 \times 3\) complex matrices in generation
  space) of the up quarks, the down quarks, the charged leptons and the
  neutrinos. Unit: dimensionless. Status: STANDARD.
\item
  \({}^\dagger\): Hermitian conjugation: the complex conjugate of the
  transpose of a matrix or spinor; for an operator on the space of
  states, its adjoint with respect to the scalar product of that space
  (for the quantised dirac16complex the indefinite Krein scalar
  product). Status: STANDARD.
\item
  \(\tilde\Phi\): The conjugate Higgs doublet,
  \(\tilde\Phi = \mathrm{i}\tau^2\Phi^*\) (hypercharge \(-1/2\)). Unit:
  GeV. Status: STANDARD.
\item
  \(Q,\ \bar Q\): The left-handed quark doublets of the three
  generations, \(Q = (u_{\mathrm{L}}, d_{\mathrm{L}})\) (colour
  triplets, hypercharge \(1/6\)). Unit: GeV\(^{3/2}\). Status: AUTHOR.
\end{itemize}

\[
\mathrm{i}\hat\gamma^\mu\big(D_\mu + \mathrm{i}\,\mathcal{V}^{d}_\mu\big)d_{\mathrm{R}} = (\mathcal{Y}^{d})^\dagger\Phi^\dagger Q \tag{FE-21}
\] \textbf{Symbols in (FE-21).}

\begin{itemize}
\tightlist
\item
  \(\mathrm{i}\): The imaginary unit, \(\mathrm{i}^2 = -1\). Unit:
  dimensionless. Status: STANDARD.
\item
  \(\hat\gamma^\mu\): The four-dimensional complex \(4 \times 4\) Dirac
  matrices of the fold: \(\hat\gamma^{(a)}\) along the frame directions
  of the fold and \(\hat\gamma^\mu\) along its coordinates (the frame
  matrices contracted with the inverse frame field of the fold),
  \(\{\hat\gamma^\mu, \hat\gamma^\nu\} = -2\hat h^{\mu\nu}\mathbf{1}\);
  \(\hat\gamma_\mu = \hat h_{\mu\nu}\hat\gamma^\nu\). Unit:
  dimensionless. Status: STANDARD.
\item
  \(\mu\): A four-dimensional coordinate index running over the
  observable coordinates \(x_1, x_2, x_3, x_4\) of the fold (signature
  \(+,+,+,-\)); a repeated index is summed. Status: STANDARD.
\item
  \(D_M,\ D_\mu\): The covariant derivative: on dirac16complex
  \(D_M = \partial_M + \Omega_M\) (with
  \(D_M\bar\Psi = \partial_M\bar\Psi - \bar\Psi\Omega_M\)); on the
  fields of the fold the gauge and spin covariant derivative \(D_\mu\)
  (the partial derivative plus the spin connection of the fold and the
  gauge fields times their couplings and generators);
  \(D^\mu = \hat h^{\mu\nu}D_\nu\); on an adjoint field
  \((D_\mu X)^A = \partial_\mu X^A + g_3f^{ABC}G^B_\mu X^C\) (likewise
  for \(SU(2)\) with \(g_2\epsilon^{IJK}\)). Unit: 1/length. Status:
  STANDARD.
\item
  \(\mathcal{V}^{f}_\mu\): The vector potential that dirac16complex
  exerts on the Standard-Model fermions of kind \(f\) through the
  current portals, \(\mathcal{V}^{f}_\mu = \sum_Xc^{f}_Xj_{X\mu}\)
  (\(\mathcal{V}^{Q}, \mathcal{V}^{u}, \dots, \mathcal{V}^{\nu}\) for
  the six kinds); flavour blind. Unit: GeV. Status: DERIVED.
\item
  \(d_{\mathrm{R}}\): The right-handed down-type quarks (down, strange
  and bottom; hypercharge \(-1/3\)). Unit: GeV\(^{3/2}\). Status:
  STANDARD.
\item
  \(\mathcal{Y}^{u},\ \mathcal{Y}^{d},\ \mathcal{Y}^{e},\ \mathcal{Y}^{\nu}\):
  The Yukawa matrices (\(3 \times 3\) complex matrices in generation
  space) of the up quarks, the down quarks, the charged leptons and the
  neutrinos. Unit: dimensionless. Status: STANDARD.
\item
  \({}^\dagger\): Hermitian conjugation: the complex conjugate of the
  transpose of a matrix or spinor; for an operator on the space of
  states, its adjoint with respect to the scalar product of that space
  (for the quantised dirac16complex the indefinite Krein scalar
  product). Status: STANDARD.
\item
  \(\Phi\): The Higgs doublet (two complex components, hypercharge
  \(1/2\)), a field of the fold. Unit: GeV. Status: AUTHOR.
\item
  \(Q,\ \bar Q\): The left-handed quark doublets of the three
  generations, \(Q = (u_{\mathrm{L}}, d_{\mathrm{L}})\) (colour
  triplets, hypercharge \(1/6\)). Unit: GeV\(^{3/2}\). Status: AUTHOR.
\end{itemize}

\[
\mathrm{i}\hat\gamma^\mu\big(D_\mu + \mathrm{i}\,\mathcal{V}^{\ell}_\mu\big)\ell = \mathcal{Y}^{e}\Phi\,e_{\mathrm{R}}
 + \mathcal{Y}^{\nu}\tilde\Phi\,\nu_{\mathrm{R}} \tag{FE-22}
\] \textbf{Symbols in (FE-22).}

\begin{itemize}
\tightlist
\item
  \(\mathrm{i}\): The imaginary unit, \(\mathrm{i}^2 = -1\). Unit:
  dimensionless. Status: STANDARD.
\item
  \(\hat\gamma^\mu\): The four-dimensional complex \(4 \times 4\) Dirac
  matrices of the fold: \(\hat\gamma^{(a)}\) along the frame directions
  of the fold and \(\hat\gamma^\mu\) along its coordinates (the frame
  matrices contracted with the inverse frame field of the fold),
  \(\{\hat\gamma^\mu, \hat\gamma^\nu\} = -2\hat h^{\mu\nu}\mathbf{1}\);
  \(\hat\gamma_\mu = \hat h_{\mu\nu}\hat\gamma^\nu\). Unit:
  dimensionless. Status: STANDARD.
\item
  \(\mu\): A four-dimensional coordinate index running over the
  observable coordinates \(x_1, x_2, x_3, x_4\) of the fold (signature
  \(+,+,+,-\)); a repeated index is summed. Status: STANDARD.
\item
  \(D_M,\ D_\mu\): The covariant derivative: on dirac16complex
  \(D_M = \partial_M + \Omega_M\) (with
  \(D_M\bar\Psi = \partial_M\bar\Psi - \bar\Psi\Omega_M\)); on the
  fields of the fold the gauge and spin covariant derivative \(D_\mu\)
  (the partial derivative plus the spin connection of the fold and the
  gauge fields times their couplings and generators);
  \(D^\mu = \hat h^{\mu\nu}D_\nu\); on an adjoint field
  \((D_\mu X)^A = \partial_\mu X^A + g_3f^{ABC}G^B_\mu X^C\) (likewise
  for \(SU(2)\) with \(g_2\epsilon^{IJK}\)). Unit: 1/length. Status:
  STANDARD.
\item
  \(\mathcal{V}^{f}_\mu\): The vector potential that dirac16complex
  exerts on the Standard-Model fermions of kind \(f\) through the
  current portals, \(\mathcal{V}^{f}_\mu = \sum_Xc^{f}_Xj_{X\mu}\)
  (\(\mathcal{V}^{Q}, \mathcal{V}^{u}, \dots, \mathcal{V}^{\nu}\) for
  the six kinds); flavour blind. Unit: GeV. Status: DERIVED.
\item
  \(\ell,\ \bar\ell\): The left-handed lepton doublets of the three
  generations, \(\ell = (\nu_{\mathrm{L}}, e_{\mathrm{L}})\)
  (hypercharge \(-1/2\)). Unit: GeV\(^{3/2}\). Status: AUTHOR.
\item
  \(\mathcal{Y}^{u},\ \mathcal{Y}^{d},\ \mathcal{Y}^{e},\ \mathcal{Y}^{\nu}\):
  The Yukawa matrices (\(3 \times 3\) complex matrices in generation
  space) of the up quarks, the down quarks, the charged leptons and the
  neutrinos. Unit: dimensionless. Status: STANDARD.
\item
  \(\Phi\): The Higgs doublet (two complex components, hypercharge
  \(1/2\)), a field of the fold. Unit: GeV. Status: AUTHOR.
\item
  \(e_{\mathrm{R}}\): The right-handed electron, muon and tau (singlets,
  hypercharge \(-1\)). Unit: GeV\(^{3/2}\). Status: AUTHOR.
\item
  \(\tilde\Phi\): The conjugate Higgs doublet,
  \(\tilde\Phi = \mathrm{i}\tau^2\Phi^*\) (hypercharge \(-1/2\)). Unit:
  GeV. Status: STANDARD.
\item
  \(\nu_{\mathrm{R}}\): The three right-handed neutrinos (gauge
  singlets). Unit: GeV\(^{3/2}\). Status: ASSUMED.
\end{itemize}

\[
\mathrm{i}\hat\gamma^\mu\big(D_\mu + \mathrm{i}\,\mathcal{V}^{e}_\mu\big)e_{\mathrm{R}} = (\mathcal{Y}^{e})^\dagger\Phi^\dagger\ell \tag{FE-23}
\] \textbf{Symbols in (FE-23).}

\begin{itemize}
\tightlist
\item
  \(\mathrm{i}\): The imaginary unit, \(\mathrm{i}^2 = -1\). Unit:
  dimensionless. Status: STANDARD.
\item
  \(\hat\gamma^\mu\): The four-dimensional complex \(4 \times 4\) Dirac
  matrices of the fold: \(\hat\gamma^{(a)}\) along the frame directions
  of the fold and \(\hat\gamma^\mu\) along its coordinates (the frame
  matrices contracted with the inverse frame field of the fold),
  \(\{\hat\gamma^\mu, \hat\gamma^\nu\} = -2\hat h^{\mu\nu}\mathbf{1}\);
  \(\hat\gamma_\mu = \hat h_{\mu\nu}\hat\gamma^\nu\). Unit:
  dimensionless. Status: STANDARD.
\item
  \(\mu\): A four-dimensional coordinate index running over the
  observable coordinates \(x_1, x_2, x_3, x_4\) of the fold (signature
  \(+,+,+,-\)); a repeated index is summed. Status: STANDARD.
\item
  \(D_M,\ D_\mu\): The covariant derivative: on dirac16complex
  \(D_M = \partial_M + \Omega_M\) (with
  \(D_M\bar\Psi = \partial_M\bar\Psi - \bar\Psi\Omega_M\)); on the
  fields of the fold the gauge and spin covariant derivative \(D_\mu\)
  (the partial derivative plus the spin connection of the fold and the
  gauge fields times their couplings and generators);
  \(D^\mu = \hat h^{\mu\nu}D_\nu\); on an adjoint field
  \((D_\mu X)^A = \partial_\mu X^A + g_3f^{ABC}G^B_\mu X^C\) (likewise
  for \(SU(2)\) with \(g_2\epsilon^{IJK}\)). Unit: 1/length. Status:
  STANDARD.
\item
  \(\mathcal{V}^{f}_\mu\): The vector potential that dirac16complex
  exerts on the Standard-Model fermions of kind \(f\) through the
  current portals, \(\mathcal{V}^{f}_\mu = \sum_Xc^{f}_Xj_{X\mu}\)
  (\(\mathcal{V}^{Q}, \mathcal{V}^{u}, \dots, \mathcal{V}^{\nu}\) for
  the six kinds); flavour blind. Unit: GeV. Status: DERIVED.
\item
  \(e_{\mathrm{R}}\): The right-handed electron, muon and tau (singlets,
  hypercharge \(-1\)). Unit: GeV\(^{3/2}\). Status: AUTHOR.
\item
  \(\mathcal{Y}^{u},\ \mathcal{Y}^{d},\ \mathcal{Y}^{e},\ \mathcal{Y}^{\nu}\):
  The Yukawa matrices (\(3 \times 3\) complex matrices in generation
  space) of the up quarks, the down quarks, the charged leptons and the
  neutrinos. Unit: dimensionless. Status: STANDARD.
\item
  \({}^\dagger\): Hermitian conjugation: the complex conjugate of the
  transpose of a matrix or spinor; for an operator on the space of
  states, its adjoint with respect to the scalar product of that space
  (for the quantised dirac16complex the indefinite Krein scalar
  product). Status: STANDARD.
\item
  \(\Phi\): The Higgs doublet (two complex components, hypercharge
  \(1/2\)), a field of the fold. Unit: GeV. Status: AUTHOR.
\item
  \(\ell,\ \bar\ell\): The left-handed lepton doublets of the three
  generations, \(\ell = (\nu_{\mathrm{L}}, e_{\mathrm{L}})\)
  (hypercharge \(-1/2\)). Unit: GeV\(^{3/2}\). Status: AUTHOR.
\end{itemize}

\[
\mathrm{i}\hat\gamma^\mu\big(D_\mu + \mathrm{i}\,\mathcal{V}^{\nu}_\mu\big)\nu_{\mathrm{R}} = (\mathcal{Y}^{\nu})^\dagger\tilde\Phi^\dagger\ell \tag{FE-24}
\] \textbf{Symbols in (FE-24).}

\begin{itemize}
\tightlist
\item
  \(\mathrm{i}\): The imaginary unit, \(\mathrm{i}^2 = -1\). Unit:
  dimensionless. Status: STANDARD.
\item
  \(\hat\gamma^\mu\): The four-dimensional complex \(4 \times 4\) Dirac
  matrices of the fold: \(\hat\gamma^{(a)}\) along the frame directions
  of the fold and \(\hat\gamma^\mu\) along its coordinates (the frame
  matrices contracted with the inverse frame field of the fold),
  \(\{\hat\gamma^\mu, \hat\gamma^\nu\} = -2\hat h^{\mu\nu}\mathbf{1}\);
  \(\hat\gamma_\mu = \hat h_{\mu\nu}\hat\gamma^\nu\). Unit:
  dimensionless. Status: STANDARD.
\item
  \(\mu\): A four-dimensional coordinate index running over the
  observable coordinates \(x_1, x_2, x_3, x_4\) of the fold (signature
  \(+,+,+,-\)); a repeated index is summed. Status: STANDARD.
\item
  \(D_M,\ D_\mu\): The covariant derivative: on dirac16complex
  \(D_M = \partial_M + \Omega_M\) (with
  \(D_M\bar\Psi = \partial_M\bar\Psi - \bar\Psi\Omega_M\)); on the
  fields of the fold the gauge and spin covariant derivative \(D_\mu\)
  (the partial derivative plus the spin connection of the fold and the
  gauge fields times their couplings and generators);
  \(D^\mu = \hat h^{\mu\nu}D_\nu\); on an adjoint field
  \((D_\mu X)^A = \partial_\mu X^A + g_3f^{ABC}G^B_\mu X^C\) (likewise
  for \(SU(2)\) with \(g_2\epsilon^{IJK}\)). Unit: 1/length. Status:
  STANDARD.
\item
  \(\mathcal{V}^{f}_\mu\): The vector potential that dirac16complex
  exerts on the Standard-Model fermions of kind \(f\) through the
  current portals, \(\mathcal{V}^{f}_\mu = \sum_Xc^{f}_Xj_{X\mu}\)
  (\(\mathcal{V}^{Q}, \mathcal{V}^{u}, \dots, \mathcal{V}^{\nu}\) for
  the six kinds); flavour blind. Unit: GeV. Status: DERIVED.
\item
  \(\nu_{\mathrm{R}}\): The three right-handed neutrinos (gauge
  singlets). Unit: GeV\(^{3/2}\). Status: ASSUMED.
\item
  \(\mathcal{Y}^{u},\ \mathcal{Y}^{d},\ \mathcal{Y}^{e},\ \mathcal{Y}^{\nu}\):
  The Yukawa matrices (\(3 \times 3\) complex matrices in generation
  space) of the up quarks, the down quarks, the charged leptons and the
  neutrinos. Unit: dimensionless. Status: STANDARD.
\item
  \({}^\dagger\): Hermitian conjugation: the complex conjugate of the
  transpose of a matrix or spinor; for an operator on the space of
  states, its adjoint with respect to the scalar product of that space
  (for the quantised dirac16complex the indefinite Krein scalar
  product). Status: STANDARD.
\item
  \(\tilde\Phi\): The conjugate Higgs doublet,
  \(\tilde\Phi = \mathrm{i}\tau^2\Phi^*\) (hypercharge \(-1/2\)). Unit:
  GeV. Status: STANDARD.
\item
  \(\ell,\ \bar\ell\): The left-handed lepton doublets of the three
  generations, \(\ell = (\nu_{\mathrm{L}}, e_{\mathrm{L}})\)
  (hypercharge \(-1/2\)). Unit: GeV\(^{3/2}\). Status: AUTHOR.
\end{itemize}

The Dirac equations of the up-quark singlets, the down-quark singlets,
the lepton doublets, the charged-lepton singlets and the right-handed
neutrinos (FE-20 to FE-24), each the variation of IN-14, IN-15 and IN-16
with respect to the conjugate field. Each fermion sees the gauge fields
of the bulk at the fold (in \(D_\mu\)), the Higgs doublet (through the
Yukawa matrices) and dirac16complex (through \(\mathcal{V}^{f}_\mu\)).
The right-handed neutrinos see no gauge field; they couple to the
Standard Model only through the neutrino Yukawa matrix, which gives the
neutrinos their Dirac masses. Status: STANDARD (INVENTED potentials).

\[
\begin{aligned}
\hat\nabla_\mu J_{\ell}^\mu &= -\hat\nabla_\mu\big(J_{e}^\mu + J_{\nu}^\mu\big)
 = -\mathrm{i}\big(\bar\ell\,\mathcal{Y}^{e}\Phi\,e_{\mathrm{R}} + \bar\ell\,\mathcal{Y}^{\nu}\tilde\Phi\,\nu_{\mathrm{R}}\big) + \mathrm{h.c.},\\
\hat\nabla_\mu J_{Q}^\mu &= -\hat\nabla_\mu\big(J_{u}^\mu + J_{d}^\mu\big)
 = -\mathrm{i}\big(\bar Q\,\mathcal{Y}^{u}\tilde\Phi\,u_{\mathrm{R}} + \bar Q\,\mathcal{Y}^{d}\Phi\,d_{\mathrm{R}}\big) + \mathrm{h.c.}
\end{aligned} \tag{FE-25}
\] \textbf{Symbols in (FE-25).}

\begin{itemize}
\tightlist
\item
  \(\hat\nabla\): The covariant derivative of the fold: with the
  Levi-Civita connection of the metric of the fold (for four-dimensional
  tensors of its metric
  \(\hat h_{\mu\nu} = \mathrm{diag}(\mathrm{e}^{2a_4}, \mathrm{e}^{2a_4}, \mathrm{e}^{2a_4}, -1)\)
  in the observable coordinates \(x_1, \dots, x_4\); for
  \(\hat\nabla_a\mathcal{S}^a{}_b\) of the induced metric of the
  seven-dimensional fold). Status: STANDARD.
\item
  \(\mu\): A four-dimensional coordinate index running over the
  observable coordinates \(x_1, x_2, x_3, x_4\) of the fold (signature
  \(+,+,+,-\)); a repeated index is summed. Status: STANDARD.
\item
  \(J_f^\mu\): The fermion-number current of the Standard-Model
  multiplets of kind \(f\), \(J_f^\mu = \bar\psi_f\hat\gamma^\mu\psi_f\)
  summed over the three generations and the gauge components; with other
  labels the currents that are the sources of the gauge fields. Unit:
  GeV\(^3\). Status: STANDARD.
\item
  \(\mathrm{i}\): The imaginary unit, \(\mathrm{i}^2 = -1\). Unit:
  dimensionless. Status: STANDARD.
\item
  \(\ell,\ \bar\ell\): The left-handed lepton doublets of the three
  generations, \(\ell = (\nu_{\mathrm{L}}, e_{\mathrm{L}})\)
  (hypercharge \(-1/2\)). Unit: GeV\(^{3/2}\). Status: AUTHOR.
\item
  \(\mathcal{Y}^{u},\ \mathcal{Y}^{d},\ \mathcal{Y}^{e},\ \mathcal{Y}^{\nu}\):
  The Yukawa matrices (\(3 \times 3\) complex matrices in generation
  space) of the up quarks, the down quarks, the charged leptons and the
  neutrinos. Unit: dimensionless. Status: STANDARD.
\item
  \(\Phi\): The Higgs doublet (two complex components, hypercharge
  \(1/2\)), a field of the fold. Unit: GeV. Status: AUTHOR.
\item
  \(e_{\mathrm{R}}\): The right-handed electron, muon and tau (singlets,
  hypercharge \(-1\)). Unit: GeV\(^{3/2}\). Status: AUTHOR.
\item
  \(\tilde\Phi\): The conjugate Higgs doublet,
  \(\tilde\Phi = \mathrm{i}\tau^2\Phi^*\) (hypercharge \(-1/2\)). Unit:
  GeV. Status: STANDARD.
\item
  \(\nu_{\mathrm{R}}\): The three right-handed neutrinos (gauge
  singlets). Unit: GeV\(^{3/2}\). Status: ASSUMED.
\item
  \(\mathrm{h.c.}\): The Hermitian conjugate of the preceding terms.
  Status: STANDARD.
\item
  \(Q,\ \bar Q\): The left-handed quark doublets of the three
  generations, \(Q = (u_{\mathrm{L}}, d_{\mathrm{L}})\) (colour
  triplets, hypercharge \(1/6\)). Unit: GeV\(^{3/2}\). Status: AUTHOR.
\item
  \(u_{\mathrm{R}}\): The right-handed up-type quarks (up, charm and
  top; singlets of \(SU(2)_L\), hypercharge \(2/3\)). Unit:
  GeV\(^{3/2}\). Status: STANDARD.
\item
  \(d_{\mathrm{R}}\): The right-handed down-type quarks (down, strange
  and bottom; hypercharge \(-1/3\)). Unit: GeV\(^{3/2}\). Status:
  STANDARD.
\end{itemize}

The divergences of the fermion-number currents of the six kinds of
multiplets (DERIVED from FE-19 to FE-24): the Yukawa couplings turn
doublets into singlets, so the current of each kind is not conserved by
itself, while the lepton number \(J_\ell + J_e + J_\nu\) and the baryon
number \((J_Q + J_u + J_d)/3\) are conserved. The divergences enter the
fold operator \(\mathbb{M}\) (FE-13). Status: DERIVED.

\hypertarget{the-ad-hoc-higgs-vacuum}{%
\subsection{The ad hoc Higgs vacuum}\label{the-ad-hoc-higgs-vacuum}}

\[
\begin{aligned}
&\mathcal{R}_{\mathcal{A}} = \mathcal{A}'' + 3\frac{b_4'}{b_4}\mathcal{A}' + \lambda\mathcal{A}^3 + c_\Phi\,\bar\Psi P_{\mathrm{L}}\Psi\,\mathcal{A}
 = \Big(4a_{40}^2 - 6a_{40}\frac{b_4'}{b_4}\Big)\mathcal{A} + \lambda\mathcal{A}^3 + c_\Phi\,\bar\Psi P_{\mathrm{L}}\Psi\,\mathcal{A},\\
&\mathcal{R}_v = \lambda\mathcal{A}^2v + c_\Phi\,\bar\Psi P_{\mathrm{L}}\Psi\,v,\qquad
 J_\Phi^{2\,x_4} = a_{40}\,v\,\mathcal{A},\qquad J_\Phi^{1\,x_4} = J_\Phi^{3\,x_4} = 0
\end{aligned} \tag{FE-26}
\] \textbf{Symbols in (FE-26).}

\begin{itemize}
\tightlist
\item
  \(\mathcal{R}_{\mathcal{A}},\ \mathcal{R}_v\): The residuals of the ad
  hoc Higgs vacuum in the upper (\(\mathcal{R}_{\mathcal{A}}\)) and
  lower (\(\mathcal{R}_v\)) components of the Higgs equation: what the
  ad hoc vacuum leaves unbalanced. Unit: GeV\(^3\). Status: DERIVED.
\item
  \(\mathcal{A}\): The author's ad hoc upper component of the Higgs
  vacuum, almightyAdHoc,
  \(\mathcal{A}(x_4) = v_{\mathcal{A}}\mathrm{e}^{-2a_{40}x_4}\). Unit:
  GeV. Status: AUTHOR.
\item
  \({}'\): The derivative of a function of the time \(x_4\) alone with
  respect to \(x_4\) (\(a_4' = \mathrm{d}a_4/\mathrm{d}x_4\); two primes
  the second derivative); on a function of \(\sigma\) alone (a profile)
  the derivative with respect to \(\sigma\). Status: STANDARD.
\item
  \(b_4\): The author's metric function of the time \(x_4\) in the
  extra-time components, \(b_4 > 0\), never identically 1 (the author's
  order of 2026-10-24); the scale factor (the factor of the lengths) of
  the extra times is \(\mathrm{e}^{-a_4}b_4\) on the fold and
  \(\mathrm{e}^{-a_4}b_4\,\mathrm{e}^{-H\lvert\sigma\rvert}\) at the
  distance \(\sigma\) from it (the square root of the warp factor times
  \(\mathrm{e}^{-a_4}b_4\)). Unit: dimensionless. Status: AUTHOR.
\item
  \(\lambda\): The quartic coupling of the Higgs potential,
  \(\lambda > 0\). Unit: dimensionless. Value:
  \(\lambda = m_h^2/(2v^2) = 0.129\). Status: AUTHOR.
\item
  \(c_\Phi\): The Higgs-portal coupling of the left-handed half of
  dirac16complex (the right-handed half has none). Unit: GeV\(^{-5}\).
  Status: INVENTED.
\item
  \(\Psi,\ \bar\Psi\): dirac16complex: the author's 16-component complex
  spinor field of eight dimensions;
  \(\bar\Psi = \Psi^\dagger\sigma_{16}\) is its Dirac adjoint; in a term
  of the fold, \(\Psi\) is its value at \(\sigma = 0\). Unit:
  GeV\(^{7/2}\). Status: AUTHOR.
\item
  \(P_{\mathrm{L}}\): The author's left-handed chiral projector,
  \(P_{\mathrm{L}} = (\mathbf{1} - \Gamma_8)/2\). Unit: dimensionless.
  Status: AUTHOR.
\item
  \(a_{40}\): The decay constant of almightyAdHoc, \(a_{40} > 0\):
  \(\mathcal{A}\) falls by \(\mathrm{e}\) in the time \(1/(2a_{40})\).
  Unit: GeV (1/length). Status: AUTHOR.
\item
  \(v\): The vacuum expectation value of the lower component of the
  Higgs doublet, \(v^2 = -\mu^2/\lambda\). Unit: GeV. Value: 246.22 GeV
  \(= (\sqrt2\,G_{\mathrm{F}})^{-1/2}\) with the Fermi constant
  \(G_{\mathrm{F}} = 1.1663788\times10^{-5}\) GeV\(^{-2}\) of the muon
  lifetime (Particle Data Group). Status: STANDARD.
\item
  \(J_\Phi^{I\nu}\): The weak-isospin current of the Higgs doublet (it
  contains the mass terms of the \(W\) fields); \(J_\Phi^{2\,x_4}\) its
  isospin charge density with \(I = 2\). Unit: GeV\(^3\). Status:
  STANDARD.
\item
  \(x_4\): The time; a prime on a function of \(x_4\) alone is
  \(\mathrm{d}/\mathrm{d}x_4\). Unit: length (= time, \(c = 1\)).
  Status: AUTHOR.
\end{itemize}

The author's ad hoc vacuum IN-12 inserted into the Higgs equation FE-15
(homogeneous fields, no gauge fields, the fold factor
\(\mathcal{Z} \propto b_4^3\mathrm{e}^{-3a_4}\) so that the time
derivative of \(\sqrt{-\hat h}\mathcal{Z}\) gives \(3b_4'/b_4\); the
metric of the fold is the same in UM-4 as in UM-3). The upper and lower
components of FE-15 are not satisfied: \(\mathcal{R}_{\mathcal{A}}\) and
\(\mathcal{R}_v\) are what remains (the residuals; FE-15 holds exactly
if the source \(\Theta\) is replaced by
\(\Theta - (\mathcal{R}_{\mathcal{A}}, \mathcal{R}_v)/\sqrt2\)). The ad
hoc vacuum is an exact solution only in the limit \(\mathcal{A} \to 0\).
It also carries a weak-isospin charge density
\(J_\Phi^{2\,x_4} = a_{40}v\mathcal{A}\) (the component along \(x_4\) of
the isospin current of FE-17 with \(I = 2\)), which the jump condition
of the \(W^2\) field (FE-17) must balance with a background field
\(W^2_{x_4}\); with the pure-gauge background \(g_2W^2_{x_4} = \theta'\)
(\(\tan(\theta/2) = \mathcal{A}/v\)) the charge density vanishes and the
Higgs equation reduces to the single radial residual
\(\mathcal{R}_V = V'' + 3(b_4'/b_4)V' + \lambda\mathcal{A}^2V + c_\Phi\bar\Psi P_{\mathrm{L}}\Psi\,V\)
of the modulus \(V = \sqrt{v^2 + \mathcal{A}^2}\) (DERIVED,
\texttt{um4\_fields.py}, F5; \texttt{um4\_background.py}, BG5 and BG6).
Status: DERIVED.

\hypertarget{the-background}{%
\section{The background}\label{the-background}}

The background is the metric of UM-4 with the two functions \(a_4(x_4)\)
and \(b_4(x_4)\) of the time, on which every other field moves. This
section writes out its geometry, the exact equations of the author's
form of the metric, the static block, and the bulk metric and the
equations of an expanding fold, to the first order in its curvature and
exactly.

\hypertarget{the-metric-of-um-4-and-the-proper-chart}{%
\subsection{The metric of UM-4 and the proper
chart}\label{the-metric-of-um-4-and-the-proper-chart}}

\[
\begin{aligned}
&g_{MN} = \mathrm{diag}\Big(\cot^2z,\ \mathrm{e}^{2a_4}(\sin^2z)^{1/6},\ \mathrm{e}^{2a_4}(\sin^2z)^{1/6},
 \ \mathrm{e}^{2a_4}(\sin^2z)^{1/6},\ -(\sin^2z)^{1/6},\\
&\qquad\qquad -\mathrm{e}^{-2a_4}b_4^2(\sin^2z)^{1/6},\ -\mathrm{e}^{-2a_4}b_4^2(\sin^2z)^{1/6},
 \ -\mathrm{e}^{-2a_4}b_4^2(\sin^2z)^{1/6}\Big),\qquad z = 6Hx_0,\\
&\sigma = \mathrm{sgn}\Big(\frac{\pi}{2} - z\Big)\frac{\ln\sin z}{6H},\qquad \mathrm{d}\sigma = \lvert\cot z\rvert\,\mathrm{d}x_0,\\
&g_{MN}\,\mathrm{d}X^M\mathrm{d}X^N = \mathrm{d}\sigma^2 + \mathrm{e}^{-2H\lvert\sigma\rvert}\Big[\mathrm{e}^{2a_4}\big(\mathrm{d}x_1^2
 + \mathrm{d}x_2^2 + \mathrm{d}x_3^2\big) - \mathrm{d}x_4^2 - \mathrm{e}^{-2a_4}b_4^2\big(\mathrm{d}x_5^2 + \mathrm{d}x_6^2 + \mathrm{d}x_7^2\big)\Big]
\end{aligned} \tag{FE-27}
\] \textbf{Symbols in (FE-27).}

\begin{itemize}
\tightlist
\item
  \(g_{MN}\): The metric tensor of the eight-dimensional spacetime
  (signature \(+,+,+,+,-,-,-,-\) in the order \(x_0, \dots, x_7\)):
  \(g_{MN}\,\mathrm{d}X^M\mathrm{d}X^N\) is the squared interval,
  \(g^{MN}\) its inverse, \(\lvert g\rvert\) the absolute value of its
  determinant (\(\sqrt{\lvert g\rvert}\,\mathrm{d}^8X\) is the proper
  volume element). Unit: dimensionless. Status: AUTHOR.
\item
  \(M\): A coordinate index of eight dimensions (also with a numbered
  copy \(M_1, M_2, \dots\)): it runs over \(x_0, x_1, \dots, x_7\); a
  repeated index is summed. Status: STANDARD.
\item
  \(N\): A coordinate index of eight dimensions (also
  \(N_1, N_2, \dots\)), running over \(x_0, \dots, x_7\); a repeated
  index is summed. Status: STANDARD.
\item
  \(\mathrm{diag}\): A diagonal matrix with the listed entries on its
  diagonal (in the order of the coordinates \(x_0, x_1, \dots, x_7\) for
  a tensor of eight dimensions, \(x_1, \dots, x_7\) for a tensor of the
  fold, \(x_1, \dots, x_4\) for a four-dimensional tensor). Status:
  STANDARD.
\item
  \(z\): The angle variable of the hidden direction, \(z = 6Hx_0\); one
  period of \(x_0\) is \(0 < z < \pi\), the fold is \(z = \pi/2\). Unit:
  dimensionless. Status: AUTHOR.
\item
  \(\mathrm{e}\): Euler's number, the base of the natural logarithm;
  \(\mathrm{e}^{X}\) is the exponential of \(X\). Unit: dimensionless.
  Value: 2.718281828. Status: STANDARD.
\item
  \(a_4\): The author's metric function of the time \(x_4\) that scales
  ordinary 3-space: on the fold \(\mathrm{e}^{a_4}\) is the scale factor
  of cosmology and \(a_4'\) the Hubble rate. Unit: dimensionless.
  Status: AUTHOR.
\item
  \(b_4\): The author's metric function of the time \(x_4\) in the
  extra-time components, \(b_4 > 0\), never identically 1 (the author's
  order of 2026-10-24); the scale factor (the factor of the lengths) of
  the extra times is \(\mathrm{e}^{-a_4}b_4\) on the fold and
  \(\mathrm{e}^{-a_4}b_4\,\mathrm{e}^{-H\lvert\sigma\rvert}\) at the
  distance \(\sigma\) from it (the square root of the warp factor times
  \(\mathrm{e}^{-a_4}b_4\)). Unit: dimensionless. Status: AUTHOR.
\item
  \(H\): The author's constant (an inverse length) that fixes the period
  \(\pi/(6H)\) of \(x_0\), the warp factor
  \(\mathrm{e}^{-2H\lvert\sigma\rvert}\) of the metric, the cosmological
  constant \(\Lambda = -21H^2\) and the tension
  \(\kappa\lambda_{\mathrm{F}} = 12H\) of the fold; in UM-4 a free
  parameter, bounded from below by the propagation of the gauge fields
  in the bulk. Unit: GeV (1/length). Value: \(H > 2534\) GeV, about 2.5
  TeV (from the invisible width of the Z boson). Status: AUTHOR.
\item
  \(x_0\): The hidden space-like coordinate (listed first); the metric
  depends on it through \(z = 6Hx_0\), periodically with the period
  \(\pi/(6H)\). Unit: length. Status: AUTHOR.
\item
  \(\sigma\): The proper coordinate of the hidden direction:
  \(\mathrm{d}\sigma = \lvert\cot z\rvert\,\mathrm{d}x_0\),
  \(\sigma = \mathrm{sgn}(\pi/2 - z)\ln(\sin z)/(6H)\) on one period
  \(0 < z < \pi\); \(\sigma = 0\) is the fold \(z = \pi/2\),
  \(\sigma \to \mp\infty\) the tips \(z \to 0, \pi\); one period of
  \(x_0\) is the whole line \(-\infty < \sigma < \infty\). Unit: length.
  Status: DERIVED.
\item
  \(\mathrm{sgn}\): The sign function: \(+1\) for a positive argument
  and \(-1\) for a negative one. Status: STANDARD.
\item
  \(\pi\): The number pi, the ratio of the circumference of a circle to
  its diameter. Unit: dimensionless. Value: 3.141592654. Status:
  STANDARD.
\item
  \(\mathrm{d}\): The differential: \(\mathrm{d}\) in front of a
  quantity is an infinitesimal change of it (as \(\mathrm{d}\sigma\), a
  small step of \(\sigma\)); the quotient of two such changes is a
  derivative; \(\mathrm{d}\) with an exponent, in front of several
  variables, is their volume element, the product of their
  differentials. Status: STANDARD.
\item
  \(X = \{x_0, \dots, x_7\}\): The author's coordinates of the
  eight-dimensional spacetime, \(X = \{x_0, x_1, \dots, x_7\}\) with the
  hidden direction first (Dirac\_Einstein-Lovelock\_4+4\_Nash.nb,
  section DEFCHART); \(X^M\) is the coordinate with the index \(M\) and
  \(\mathrm{d}^8X = \mathrm{d}x_0\,\mathrm{d}x_1\cdots\mathrm{d}x_7\)
  the coordinate volume element. Unit: length (each coordinate). Status:
  AUTHOR.
\item
  \(x_1\): The first coordinate of ordinary 3-space. Unit: length.
  Status: AUTHOR.
\item
  \(x_2\): The second coordinate of ordinary 3-space. Unit: length.
  Status: AUTHOR.
\item
  \(x_3\): The third coordinate of ordinary 3-space. Unit: length.
  Status: AUTHOR.
\item
  \(x_4\): The time; a prime on a function of \(x_4\) alone is
  \(\mathrm{d}/\mathrm{d}x_4\). Unit: length (= time, \(c = 1\)).
  Status: AUTHOR.
\item
  \(x_5\): The first of the three extra time-like coordinates. Unit:
  length. Status: AUTHOR.
\item
  \(x_6\): The second of the three extra time-like coordinates. Unit:
  length. Status: AUTHOR.
\item
  \(x_7\): The third of the three extra time-like coordinates. Unit:
  length. Status: AUTHOR.
\end{itemize}

The metric of UM-4: the author's metric of the primordial gravitational
field (his order of 2026-10-24: the metric of
\texttt{Dirac\_Einstein-Lovelock\_4+4\_Nash.nb} with \(b_4 \to b_4^2\);
\((\sin^2z)^{1/6}\) is the positive branch of
\texttt{Sin{[}6\ H\ x0{]}\^{}(1/3)}, the author's rule of 2026-10-16)
with the one change that his order of 2026-10-25 13:15 requires: the
time \(x_4\) is warped like the other directions,
\(g_{x_4x_4} = -(\sin^2z)^{1/6}\) instead of \(-1\). In the proper chart
\(\sigma\) of the hidden direction (second and third lines) the warp
factor \(\mathrm{e}^{-2H\lvert\sigma\rvert}\) multiplies the whole
metric of the fold (FD-2), time included. In the author's chart the
metric is smooth but degenerate at the fold \(z = \pi/2\)
(\(\cot z = 0\)); the proper distance from the fold is
\(\lvert\sigma\rvert\), and \(\sigma\) runs from \(-\infty\) (the tip
\(z \to 0\)) through \(0\) (the fold) to \(+\infty\) (the tip
\(z \to \pi\)): one period of \(x_0\) is the whole line, and the two
halves of a period are mirror images under \(z \to \pi - z\)
(\(\sigma \to -\sigma\)). \(a_4\) and \(b_4\) are functions of \(x_4\).
Status: AUTHOR (the metric and the order), ASSUMED (the way the order is
met: the time gets the warp factor of the other directions, the smallest
change of the metric that the order allows), DERIVED (the chart,
\texttt{um4\_geometry.py}, G1).

\[
\begin{aligned}
&\mathbf{e}^{(0)} = \mathrm{d}\sigma,\quad \mathbf{e}^{(i)} = \mathrm{e}^{a_4 - H\lvert\sigma\rvert}\mathrm{d}x_i,\quad
 \mathbf{e}^{(4)} = \mathrm{e}^{-H\lvert\sigma\rvert}\mathrm{d}x_4,\quad \mathbf{e}^{(r)} = b_4\,\mathrm{e}^{-a_4 - H\lvert\sigma\rvert}\mathrm{d}x_r,\qquad
 \sqrt{\lvert g\rvert} = b_4^3\,\mathrm{e}^{-7H\lvert\sigma\rvert},\\
&R^{(a)(b)}{}_{(c)(d)} = \mathrm{e}^{2H\lvert\sigma\rvert}\,\hat R^{(a)(b)}{}_{(c)(d)} - H^2\,\delta^{ab}_{cd}\quad(\text{all four along the fold}),\qquad
 R^{(0)(a)}{}_{(0)(b)} = -H^2\,\delta^a{}_b,\\
&\hat R^{(i)(j)}{}_{(i)(j)} = \mathsf{A}^2H^2,\quad \hat R^{(i)(4)}{}_{(i)(4)} = (\mathsf{A}^2 + \mathsf{A}_1)H^2,\quad
 \hat R^{(i)(r)}{}_{(i)(r)} = \mathsf{A}\mathsf{V}H^2,\\
&\hat R^{(r)(s)}{}_{(r)(s)} = \mathsf{V}^2H^2,\quad
 \hat R^{(4)(r)}{}_{(4)(r)} = (\mathsf{V}^2 + \mathsf{V}_1)H^2,\\
&\mathsf{A} = \frac{a_4'}{H},\qquad \mathsf{U} = \frac{b_4'}{Hb_4},\qquad \mathsf{V} = \mathsf{U} - \mathsf{A},\qquad
 \mathsf{A}_1 = \frac{\mathsf{A}'}{H},\qquad \mathsf{V}_1 = \frac{\mathsf{V}'}{H}
\end{aligned} \tag{FE-28}
\] \textbf{Symbols in (FE-28).}

\begin{itemize}
\tightlist
\item
  \(\mathbf{e}_a{}^M,\ \mathbf{e}^a{}_M\): The vielbein (frame field) of
  the metric: \(g_{MN} = \eta_{ab}\,\mathbf{e}^a{}_M\mathbf{e}^b{}_N\);
  \(\mathbf{e}_a{}^M\) is its inverse;
  \(\mathbf{e}_{bN} = \eta_{bc}\mathbf{e}^c{}_N\); \(\mathbf{e}^{(a)}\)
  is the frame one-form \(\mathbf{e}^a{}_M\,\mathrm{d}X^M\). Unit:
  dimensionless. Status: STANDARD.
\item
  \(\mathrm{d}\): The differential: \(\mathrm{d}\) in front of a
  quantity is an infinitesimal change of it (as \(\mathrm{d}\sigma\), a
  small step of \(\sigma\)); the quotient of two such changes is a
  derivative; \(\mathrm{d}\) with an exponent, in front of several
  variables, is their volume element, the product of their
  differentials. Status: STANDARD.
\item
  \(\sigma\): The proper coordinate of the hidden direction:
  \(\mathrm{d}\sigma = \lvert\cot z\rvert\,\mathrm{d}x_0\),
  \(\sigma = \mathrm{sgn}(\pi/2 - z)\ln(\sin z)/(6H)\) on one period
  \(0 < z < \pi\); \(\sigma = 0\) is the fold \(z = \pi/2\),
  \(\sigma \to \mp\infty\) the tips \(z \to 0, \pi\); one period of
  \(x_0\) is the whole line \(-\infty < \sigma < \infty\). Unit: length.
  Status: DERIVED.
\item
  \(\mathrm{e}\): Euler's number, the base of the natural logarithm;
  \(\mathrm{e}^{X}\) is the exponential of \(X\). Unit: dimensionless.
  Value: 2.718281828. Status: STANDARD.
\item
  \(a_4\): The author's metric function of the time \(x_4\) that scales
  ordinary 3-space: on the fold \(\mathrm{e}^{a_4}\) is the scale factor
  of cosmology and \(a_4'\) the Hubble rate. Unit: dimensionless.
  Status: AUTHOR.
\item
  \(H\): The author's constant (an inverse length) that fixes the period
  \(\pi/(6H)\) of \(x_0\), the warp factor
  \(\mathrm{e}^{-2H\lvert\sigma\rvert}\) of the metric, the cosmological
  constant \(\Lambda = -21H^2\) and the tension
  \(\kappa\lambda_{\mathrm{F}} = 12H\) of the fold; in UM-4 a free
  parameter, bounded from below by the propagation of the gauge fields
  in the bulk. Unit: GeV (1/length). Value: \(H > 2534\) GeV, about 2.5
  TeV (from the invisible width of the Z boson). Status: AUTHOR.
\item
  \(x\): The coordinates of the fold: \(x_1, x_2, x_3\) of ordinary
  space, \(x_4\) the time, \(x_5, x_6, x_7\) the extra times; a field of
  the fold is a function of the four observable coordinates
  \(x = (x_1, x_2, x_3, x_4)\), and
  \(\mathrm{d}^4x = \mathrm{d}x_1\,\mathrm{d}x_2\,\mathrm{d}x_3\,\mathrm{d}x_4\).
  Unit: length (each coordinate). Status: AUTHOR.
\item
  \(i\): A generation index (\(i = 1, 2, 3\)), which also labels the
  three neutrino mass states; in a frame component such as
  \(\mathcal{S}^{(i)}{}_{(i)}\) one of the three directions of ordinary
  space (the three components are equal), and in a tensor of ordinary
  space such as \(\mathsf{h}_{ij}\) one of the directions
  \(x_1, x_2, x_3\); a repeated index is not summed unless a sum sign is
  written. Status: STANDARD.
\item
  \(x_4\): The time; a prime on a function of \(x_4\) alone is
  \(\mathrm{d}/\mathrm{d}x_4\). Unit: length (= time, \(c = 1\)).
  Status: AUTHOR.
\item
  \(b_4\): The author's metric function of the time \(x_4\) in the
  extra-time components, \(b_4 > 0\), never identically 1 (the author's
  order of 2026-10-24); the scale factor (the factor of the lengths) of
  the extra times is \(\mathrm{e}^{-a_4}b_4\) on the fold and
  \(\mathrm{e}^{-a_4}b_4\,\mathrm{e}^{-H\lvert\sigma\rvert}\) at the
  distance \(\sigma\) from it (the square root of the warp factor times
  \(\mathrm{e}^{-a_4}b_4\)). Unit: dimensionless. Status: AUTHOR.
\item
  \(r\): An extra-time direction of the fold: \(r \in \{5, 6, 7\}\)
  (\(x_5, x_6, x_7\)); not summed (the three components are equal).
  Status: STANDARD.
\item
  \(g_{MN}\): The metric tensor of the eight-dimensional spacetime
  (signature \(+,+,+,+,-,-,-,-\) in the order \(x_0, \dots, x_7\)):
  \(g_{MN}\,\mathrm{d}X^M\mathrm{d}X^N\) is the squared interval,
  \(g^{MN}\) its inverse, \(\lvert g\rvert\) the absolute value of its
  determinant (\(\sqrt{\lvert g\rvert}\,\mathrm{d}^8X\) is the proper
  volume element). Unit: dimensionless. Status: AUTHOR.
\item
  \(R^{MN}{}_{PQ},\ R_{MN},\ R\): The curvature of the metric
  \(g_{MN}\): with four indices the Riemann tensor \(R^M{}_{NPQ}\) (in
  the convention in which the commutator of two covariant derivatives,
  acting on a vector, is the Riemann tensor, with its last two indices
  those of the two derivatives, contracted with the vector), with two
  indices the Ricci tensor \(R_{MN} = R^P{}_{MPN}\), without indices the
  scalar curvature \(R = g^{MN}R_{MN}\); with frame indices
  \(R^{(a)(b)}{}_{(c)(d)}\) its frame components,
  \(R^{(a)(b)}{}_{(a)(b)}\) (no sum) the sectional curvature of the
  frame plane \((a, b)\). Unit: 1/length\(^2\). Status: STANDARD.
\item
  \(a\): A frame (orthonormal) index running over \(0, 1, \dots, 7\)
  (frame direction \((0)\) along \(\sigma\), \((1), (2), (3)\) along
  ordinary space, \((4)\) along the time, \((5), (6), (7)\) along the
  extra times); raised and lowered with the frame metric
  \(\eta_{ab} = \mathrm{diag}(1, 1, 1, 1, -1, -1, -1, -1)\); on the fold
  (\(a, b \in \{1, \dots, 7\}\)) the frame of the fold; in the tensors
  of the fold written with its coordinates \(x^a\) (its metric
  \(\hat g_{ab}\) and the tensors built from it) a coordinate index of
  the fold, \(a \in \{1, \dots, 7\}\); a repeated index is summed.
  Status: STANDARD.
\item
  \(b\): A frame index (as \(a\)). Status: STANDARD.
\item
  \(c\): A frame index (as \(a\)). Status: STANDARD.
\item
  \(d\): A frame index (as \(a\)). Status: STANDARD.
\item
  \(\hat R\): The curvature of the metric \(\hat g_{ab}\) of the fold
  (FD-2): with four frame indices its Riemann tensor, with two its Ricci
  tensor \(\hat R_{ab}\), without indices its scalar curvature
  \(\hat R\); for frozen extra times the scalar curvature equals that of
  the four-dimensional metric \(\hat h_{\mu\nu}\). Unit: 1/length\(^2\).
  Status: STANDARD.
\item
  \(\delta\): The Kronecker delta: with two indices \(\delta^M_N = 1\)
  if \(M = N\) and 0 otherwise (likewise \(\delta_{X\mathrm{L}}\) is 1
  if the label \(X\) of a half of dirac16complex is L and 0 otherwise);
  with two index strings of equal length the generalised Kronecker delta
  \(\delta^{M_1M_2\cdots}_{N_1N_2\cdots}\), the determinant of the
  square matrix of the Kronecker deltas \(\delta^{M_a}_{N_b}\) (\(+1\)
  or \(-1\) if the upper indices are an even or odd permutation of the
  lower ones, 0 otherwise). Status: STANDARD.
\item
  \(j\): A second generation or mass-state index, or a second direction
  of ordinary space (as \(i\)). Status: STANDARD.
\item
  \(s\): A second extra-time direction of the fold (as \(r\), different
  from it). Status: STANDARD.
\item
  \(\mathsf{A}\): The expansion rate of ordinary space in units of
  \(H\), \(\mathsf{A} = a_4'/H\). Unit: dimensionless. Value: 0 at the
  static block; \(H_4^{(0)}/H\) today (\(H_4^{(0)}\) the Hubble
  constant); the first-order equations of the fold hold for
  \(\mathsf{A}^2 \ll 1\); 1 at the author's vacuum, which is not a
  solution of UM-4. Status: DERIVED.
\item
  \(\mathsf{A}_1\): The time derivative of \(\mathsf{A}\) in units of
  \(H\), \(\mathsf{A}_1 = \mathsf{A}'/H = a_4''/H^2\). Unit:
  dimensionless. Status: DERIVED.
\item
  \(\mathsf{V}\): The rate of change of the scale factor
  \(\mathrm{e}^{-a_4}b_4\) of the extra times on the fold in units of
  \(H\), \(\mathsf{V} = \mathsf{U} - \mathsf{A}\); zero with frozen
  extra times. Unit: dimensionless. Status: DERIVED.
\item
  \(\mathsf{V}_1\): The time derivative of \(\mathsf{V}\) in units of
  \(H\), \(\mathsf{V}_1 = \mathsf{V}'/H\). Unit: dimensionless. Status:
  DERIVED.
\item
  \(\mathsf{U}\): The rate of change of \(b_4\) in units of \(H\),
  \(\mathsf{U} = b_4'/(Hb_4)\). Unit: dimensionless. Value:
  \(\mathsf{U} = \mathsf{A}\) with frozen extra times. Status: DERIVED.
\item
  \({}'\): The derivative of a function of the time \(x_4\) alone with
  respect to \(x_4\) (\(a_4' = \mathrm{d}a_4/\mathrm{d}x_4\); two primes
  the second derivative); on a function of \(\sigma\) alone (a profile)
  the derivative with respect to \(\sigma\). Status: STANDARD.
\end{itemize}

The vielbein and the volume element (\(i, j\) any two different
directions of ordinary space, \(r, s\) any two different extra times;
\(x_i\) and \(x_r\) the coordinates of these directions), the exact
frame curvature of the metric of UM-4 for any \(a_4(x_4)\), \(b_4(x_4)\)
(second line), and the curvature \(\hat R\) of the metric of the fold
FD-2 in its 21 frame planes (third line), with the dimensionless rates
of the background (\(\mathsf{A}\) the expansion rate of ordinary space
in units of \(H\), \(\mathsf{U}\) the rate of change of \(b_4\),
\(\mathsf{V}\) the rate of the scale factor \(\mathrm{e}^{-a_4}b_4\) of
the extra times). In every plane of two directions of the fold the
curvature of the bulk is that of the fold, enlarged by
\(\mathrm{e}^{2H\lvert\sigma\rvert}\), minus \(H^2\); in every plane
that contains the hidden direction it is \(-H^2\); there is no other
component (UM-3 had mixed components, because its time was not warped).
The metric of UM-4 is a warped product: the bulk is a space of curvature
\(-H^2\) wherever the fold is flat (DERIVED, \texttt{um4\_geometry.py},
G2, exact, computed in the author's coordinates on both halves of the
period; \texttt{um4\_background.py}, BG1). Status: DERIVED.

\hypertarget{the-exact-bulk-equations-of-the-authors-form-of-the-metric}{%
\subsection{The exact bulk equations of the author's form of the
metric}\label{the-exact-bulk-equations-of-the-authors-form-of-the-metric}}

\[
\mathcal{E}_{(1)}{}^{(a)}{}_{(b)} + \Lambda\,\delta^a{}_b = \mathrm{e}^{2H\lvert\sigma\rvert}\,\hat G^a{}_b\quad(\text{both along the fold}),\qquad
\mathcal{E}_{(1)}{}^{(0)}{}_{(0)} + \Lambda = -\frac12\,\mathrm{e}^{2H\lvert\sigma\rvert}\,\hat R,\qquad
\mathcal{E}_{(1)}{}^{(0)}{}_{(a)} = 0 \tag{FE-29}
\] \textbf{Symbols in (FE-29).}

\begin{itemize}
\tightlist
\item
  \(\mathcal{E}_{(p)}{}^M{}_N\): The Lovelock tensor of order \(p\),
  \(\mathcal{E}_{(p)}{}^M{}_N = -2^{-(p+1)}\,\delta^{MP_1\cdots P_{2p}}_{NQ_1\cdots Q_{2p}}\,R^{Q_1Q_2}{}_{P_1P_2}\cdots R^{Q_{2p-1}Q_{2p}}{}_{P_{2p-1}P_{2p}}\)
  (with both indices down it is \(1/\sqrt{\lvert g\rvert}\) times the
  variation of \(\sqrt{\lvert g\rvert}\,\mathcal{L}_{(p)}\) with respect
  to \(g^{MN}\); \(\mathcal{E}_{(1)}\) the Einstein tensor,
  \(\mathcal{E}_{(2)}\) the Gauss-Bonnet tensor, \(\mathcal{E}_{(3)}\)
  the cubic Lovelock tensor); with frame indices its frame components.
  Unit: 1/length\(^{2p}\). Status: STANDARD.
\item
  \(p\): The Lovelock order \(p = 1, 2, 3\) (Einstein, Gauss-Bonnet,
  cubic), as an index and as a number (UM-4 keeps only \(p = 1\)); in
  the Kohn-Sham equations also a second label of a Kohn-Sham orbital,
  with the Kronecker delta of the two labels. Status: STANDARD.
\item
  \(a\): A frame (orthonormal) index running over \(0, 1, \dots, 7\)
  (frame direction \((0)\) along \(\sigma\), \((1), (2), (3)\) along
  ordinary space, \((4)\) along the time, \((5), (6), (7)\) along the
  extra times); raised and lowered with the frame metric
  \(\eta_{ab} = \mathrm{diag}(1, 1, 1, 1, -1, -1, -1, -1)\); on the fold
  (\(a, b \in \{1, \dots, 7\}\)) the frame of the fold; in the tensors
  of the fold written with its coordinates \(x^a\) (its metric
  \(\hat g_{ab}\) and the tensors built from it) a coordinate index of
  the fold, \(a \in \{1, \dots, 7\}\); a repeated index is summed.
  Status: STANDARD.
\item
  \(b\): A frame index (as \(a\)). Status: STANDARD.
\item
  \(\Lambda\): The eight-dimensional cosmological constant of UM-4,
  \(\Lambda = -21H^2\): the value for which the static block is a vacuum
  of Einstein gravity (its Einstein tensor is \(21H^2\) in every
  direction). The author writes the pure number \(\Lambda_{\mathrm{N}}\)
  in
  \(0 = -\Lambda_{\mathrm{N}} + \sum_p H^{-2p}w_{(p)}\mathcal{E}_{(p)}\),
  so that \(\Lambda = -\Lambda_{\mathrm{N}}H^2\). Unit: GeV\(^2\).
  Value: \(-21H^2\); the author's original \(-10H^2\); UM-3: \(-7H^2\).
  Status: AUTHOR (the term), DERIVED (the value).
\item
  \(\delta\): The Kronecker delta: with two indices \(\delta^M_N = 1\)
  if \(M = N\) and 0 otherwise (likewise \(\delta_{X\mathrm{L}}\) is 1
  if the label \(X\) of a half of dirac16complex is L and 0 otherwise);
  with two index strings of equal length the generalised Kronecker delta
  \(\delta^{M_1M_2\cdots}_{N_1N_2\cdots}\), the determinant of the
  square matrix of the Kronecker deltas \(\delta^{M_a}_{N_b}\) (\(+1\)
  or \(-1\) if the upper indices are an even or odd permutation of the
  lower ones, 0 otherwise). Status: STANDARD.
\item
  \(\mathrm{e}\): Euler's number, the base of the natural logarithm;
  \(\mathrm{e}^{X}\) is the exponential of \(X\). Unit: dimensionless.
  Value: 2.718281828. Status: STANDARD.
\item
  \(H\): The author's constant (an inverse length) that fixes the period
  \(\pi/(6H)\) of \(x_0\), the warp factor
  \(\mathrm{e}^{-2H\lvert\sigma\rvert}\) of the metric, the cosmological
  constant \(\Lambda = -21H^2\) and the tension
  \(\kappa\lambda_{\mathrm{F}} = 12H\) of the fold; in UM-4 a free
  parameter, bounded from below by the propagation of the gauge fields
  in the bulk. Unit: GeV (1/length). Value: \(H > 2534\) GeV, about 2.5
  TeV (from the invisible width of the Z boson). Status: AUTHOR.
\item
  \(\sigma\): The proper coordinate of the hidden direction:
  \(\mathrm{d}\sigma = \lvert\cot z\rvert\,\mathrm{d}x_0\),
  \(\sigma = \mathrm{sgn}(\pi/2 - z)\ln(\sin z)/(6H)\) on one period
  \(0 < z < \pi\); \(\sigma = 0\) is the fold \(z = \pi/2\),
  \(\sigma \to \mp\infty\) the tips \(z \to 0, \pi\); one period of
  \(x_0\) is the whole line \(-\infty < \sigma < \infty\). Unit: length.
  Status: DERIVED.
\item
  \(\hat G^\mu{}_\nu\): The Einstein tensor of the metric
  \(\hat g_{ab}\) of the fold (FD-2), with frame indices along its seven
  directions; with four-dimensional indices, \(\hat G^\mu{}_\nu\), its
  components along the four observable directions, which for frozen
  extra times are those of the Einstein tensor of the four-dimensional
  metric \(\hat h_{\mu\nu}\). Unit: 1/length\(^2\). Status: STANDARD.
\item
  \(\hat R\): The curvature of the metric \(\hat g_{ab}\) of the fold
  (FD-2): with four frame indices its Riemann tensor, with two its Ricci
  tensor \(\hat R_{ab}\), without indices its scalar curvature
  \(\hat R\); for frozen extra times the scalar curvature equals that of
  the four-dimensional metric \(\hat h_{\mu\nu}\). Unit: 1/length\(^2\).
  Status: STANDARD.
\end{itemize}

The bulk equations FE-1 for the metric of UM-4 with Einstein gravity and
\(\Lambda = -21H^2\), exactly, for any \(a_4(x_4)\), \(b_4(x_4)\)
(\(\hat G^a{}_b\) and \(\hat R\) the Einstein tensor and the scalar
curvature of the metric of the fold FD-2). They state what the bulk
sources must be: the author's form of the metric is a vacuum of the bulk
exactly when the metric of the fold is Ricci flat (\(\hat R_{ab} = 0\)),
and with sources only when these grow like
\(\mathrm{e}^{2H\lvert\sigma\rvert}\) away from the fold. An expanding
fold with matter therefore has a bulk metric that is not of the author's
form (FE-34). (In UM-3 the left-hand side was a polynomial in the rates,
with a vacuum branch for any \(a_4\) with frozen extra times; with
Einstein gravity and the time warped there is no such branch, FE-33.)
Status: DERIVED (\texttt{um4\_geometry.py}, G5).

\hypertarget{the-fold-of-the-authors-form-of-the-metric}{%
\subsection{The fold of the author's form of the
metric}\label{the-fold-of-the-authors-form-of-the-metric}}

\[
K^a{}_b = \frac12\,g^{ac}\,\partial_\sigma g_{cb} = -H\,\mathrm{sgn}(\sigma)\,\delta^a{}_b,\qquad
\kappa\,\mathcal{S}^a{}_b = -12H\,\delta^a{}_b\quad\big(\kappa\lambda_{\mathrm{F}} = 12H\big) \tag{FE-30}
\] \textbf{Symbols in (FE-30).}

\begin{itemize}
\tightlist
\item
  \(K^a{}_b\): The extrinsic curvature of a surface of constant
  \(\sigma\) of the hidden direction,
  \(K^a{}_b = \frac12 g^{ac}\partial_\sigma g_{cb}\) for \(a, b\) along
  the surface (FD-5): \(-H\,\mathrm{sgn}(\sigma)\) along all seven
  directions of the fold for the author's form of the metric of UM-4;
  \(K = K^c{}_c\) its trace. Unit: GeV (1/length). Status: DERIVED.
\item
  \(a\): A frame (orthonormal) index running over \(0, 1, \dots, 7\)
  (frame direction \((0)\) along \(\sigma\), \((1), (2), (3)\) along
  ordinary space, \((4)\) along the time, \((5), (6), (7)\) along the
  extra times); raised and lowered with the frame metric
  \(\eta_{ab} = \mathrm{diag}(1, 1, 1, 1, -1, -1, -1, -1)\); on the fold
  (\(a, b \in \{1, \dots, 7\}\)) the frame of the fold; in the tensors
  of the fold written with its coordinates \(x^a\) (its metric
  \(\hat g_{ab}\) and the tensors built from it) a coordinate index of
  the fold, \(a \in \{1, \dots, 7\}\); a repeated index is summed.
  Status: STANDARD.
\item
  \(b\): A frame index (as \(a\)). Status: STANDARD.
\item
  \(g_{MN}\): The metric tensor of the eight-dimensional spacetime
  (signature \(+,+,+,+,-,-,-,-\) in the order \(x_0, \dots, x_7\)):
  \(g_{MN}\,\mathrm{d}X^M\mathrm{d}X^N\) is the squared interval,
  \(g^{MN}\) its inverse, \(\lvert g\rvert\) the absolute value of its
  determinant (\(\sqrt{\lvert g\rvert}\,\mathrm{d}^8X\) is the proper
  volume element). Unit: dimensionless. Status: AUTHOR.
\item
  \(c\): A frame index (as \(a\)). Status: STANDARD.
\item
  \(\partial\): A partial derivative: the derivative with respect to one
  variable at fixed values of the others (\(\partial_\sigma\) with
  respect to \(\sigma\); \(\partial_M\) and \(\partial_\mu\) with
  respect to the coordinate that carries the same index, of the eight
  dimensions or of the fold). Status: STANDARD.
\item
  \(\sigma\): The label \(\sigma\) in \(\partial_\sigma\): the
  derivative with respect to the proper hidden coordinate \(\sigma\).
  Status: STANDARD.
\item
  \(H\): The author's constant (an inverse length) that fixes the period
  \(\pi/(6H)\) of \(x_0\), the warp factor
  \(\mathrm{e}^{-2H\lvert\sigma\rvert}\) of the metric, the cosmological
  constant \(\Lambda = -21H^2\) and the tension
  \(\kappa\lambda_{\mathrm{F}} = 12H\) of the fold; in UM-4 a free
  parameter, bounded from below by the propagation of the gauge fields
  in the bulk. Unit: GeV (1/length). Value: \(H > 2534\) GeV, about 2.5
  TeV (from the invisible width of the Z boson). Status: AUTHOR.
\item
  \(\mathrm{sgn}\): The sign function: \(+1\) for a positive argument
  and \(-1\) for a negative one. Status: STANDARD.
\item
  \(\sigma\): The proper coordinate of the hidden direction:
  \(\mathrm{d}\sigma = \lvert\cot z\rvert\,\mathrm{d}x_0\),
  \(\sigma = \mathrm{sgn}(\pi/2 - z)\ln(\sin z)/(6H)\) on one period
  \(0 < z < \pi\); \(\sigma = 0\) is the fold \(z = \pi/2\),
  \(\sigma \to \mp\infty\) the tips \(z \to 0, \pi\); one period of
  \(x_0\) is the whole line \(-\infty < \sigma < \infty\). Unit: length.
  Status: DERIVED.
\item
  \(\delta\): The Kronecker delta: with two indices \(\delta^M_N = 1\)
  if \(M = N\) and 0 otherwise (likewise \(\delta_{X\mathrm{L}}\) is 1
  if the label \(X\) of a half of dirac16complex is L and 0 otherwise);
  with two index strings of equal length the generalised Kronecker delta
  \(\delta^{M_1M_2\cdots}_{N_1N_2\cdots}\), the determinant of the
  square matrix of the Kronecker deltas \(\delta^{M_a}_{N_b}\) (\(+1\)
  or \(-1\) if the upper indices are an even or odd permutation of the
  lower ones, 0 otherwise). Status: STANDARD.
\item
  \(\kappa\): The eight-dimensional gravitational coupling,
  \(\kappa = 8\pi G_8\). Unit: GeV\(^{-6}\). Status: AUTHOR.
\item
  \(\mathcal{S}^a{}_b\): The thin source of the fold: the
  energy-momentum tensor concentrated on the fold, integrated across it,
  \(\kappa\mathcal{S}^a{}_b = \int_{0^-}^{0^+}\mathrm{d}\sigma\,\mathcal{E}_{(1)}{}^a{}_b\)
  (the junction condition of Israel); \(\mathcal{S}^4{}_4\) is minus its
  energy density, \(\mathcal{S}^1{}_1\) and \(\mathcal{S}^5{}_5\) its
  pressures along ordinary space and along the extra times (frame
  components along the fold, eight-dimensional normalisation: an energy
  per unit volume of the six other directions of the fold). Unit:
  GeV\(^7\). Value: \(\kappa\mathcal{S}^a{}_b = -12H\delta^a{}_b\) for
  every \(a_4\), \(b_4\) of the author's form. Status: DERIVED.
\item
  \(\lambda_{\mathrm{F}}\): The tension that the static block requires
  of the fold, \(\kappa\lambda_{\mathrm{F}} = 12H\): the total vacuum
  energy of the seven-dimensional fold, made of the tension of the fold
  fluid and the vacuum energy of the fields of the fold (both have the
  pressure minus their energy density along ordinary space and along the
  extra times), per unit volume of the six other directions of the fold
  (eight-dimensional normalisation); in the four-dimensional
  normalisation it is \(\lambda_{\mathrm{F}}V_*\), and the excess of the
  total over \(\lambda_{\mathrm{F}}V_*\) is the dark energy
  \(\rho_\Lambda\). Unit: GeV\(^7\). Value: \(12H/\kappa\);
  \(\lambda_{\mathrm{F}}V_* = 30H^2M_{\mathrm{red}}^2 = 1.1\times10^{45}\)
  GeV\(^4\) for \(H = 2534\) GeV. Status: DERIVED.
\end{itemize}

The extrinsic curvature of the slices \(\sigma = \text{const}\) in all
seven directions of the fold (UM-3: zero along the time) and the thin
source of the fold from Israel's condition FE-4: for every \(a_4(x_4)\),
\(b_4(x_4)\) of the author's form a pure tension
\(\lambda_{\mathrm{F}}\) with \(\kappa\lambda_{\mathrm{F}} = 12H\), the
same in every direction (\(p = q = -\rho\)). For a fold of \(n\)
dimensions in a bulk of \(n + 1\) dimensions with this warp the tension
is \(\kappa\lambda_{\mathrm{F}} = 2(n - 1)H\), the critical tension that
belongs to the cosmological constant \(\Lambda = -n(n - 1)H^2/2\) of the
bulk: for \(n = 4\) they are the tension
\(\lambda_{\mathrm{RS}} = 6k_{\mathrm{RS}}/\kappa_5\) and the
cosmological constant \(\Lambda_{\mathrm{RS}} = -6k_{\mathrm{RS}}^2\) of
the brane of Randall and Sundrum (FD-8); by Gauss's equation a fold of
pure tension in this bulk is flat only at this tension (a larger tension
curves it like de Sitter space, a smaller one like the space of negative
curvature). Status: DERIVED (\texttt{um4\_geometry.py}, G7: the jump of
the extrinsic curvature computed; \texttt{um4\_foldsources.py}, FS1: the
tension, the cosmological constant and the critical tension computed
from the metric for \(n = 7\) and \(n = 4\)).

\hypertarget{the-authors-vacuum}{%
\subsection{The author's vacuum}\label{the-authors-vacuum}}

\[
a_4 = Hx_4,\quad b_4 = \text{constant}\ (\mathsf{A} = 1,\ \mathsf{U} = 0):\qquad
\hat R_{ab} \ne 0,\qquad \kappa_7\rho_7 = -3H^2,\quad \kappa_7P_7 = \kappa_7q_7 = -3H^2 \tag{FE-31}
\] \textbf{Symbols in (FE-31).}

\begin{itemize}
\tightlist
\item
  \(a_4\): The author's metric function of the time \(x_4\) that scales
  ordinary 3-space: on the fold \(\mathrm{e}^{a_4}\) is the scale factor
  of cosmology and \(a_4'\) the Hubble rate. Unit: dimensionless.
  Status: AUTHOR.
\item
  \(H\): The author's constant (an inverse length) that fixes the period
  \(\pi/(6H)\) of \(x_0\), the warp factor
  \(\mathrm{e}^{-2H\lvert\sigma\rvert}\) of the metric, the cosmological
  constant \(\Lambda = -21H^2\) and the tension
  \(\kappa\lambda_{\mathrm{F}} = 12H\) of the fold; in UM-4 a free
  parameter, bounded from below by the propagation of the gauge fields
  in the bulk. Unit: GeV (1/length). Value: \(H > 2534\) GeV, about 2.5
  TeV (from the invisible width of the Z boson). Status: AUTHOR.
\item
  \(x_4\): The time; a prime on a function of \(x_4\) alone is
  \(\mathrm{d}/\mathrm{d}x_4\). Unit: length (= time, \(c = 1\)).
  Status: AUTHOR.
\item
  \(b_4\): The author's metric function of the time \(x_4\) in the
  extra-time components, \(b_4 > 0\), never identically 1 (the author's
  order of 2026-10-24); the scale factor (the factor of the lengths) of
  the extra times is \(\mathrm{e}^{-a_4}b_4\) on the fold and
  \(\mathrm{e}^{-a_4}b_4\,\mathrm{e}^{-H\lvert\sigma\rvert}\) at the
  distance \(\sigma\) from it (the square root of the warp factor times
  \(\mathrm{e}^{-a_4}b_4\)). Unit: dimensionless. Status: AUTHOR.
\item
  \(\mathsf{A}\): The expansion rate of ordinary space in units of
  \(H\), \(\mathsf{A} = a_4'/H\). Unit: dimensionless. Value: 0 at the
  static block; \(H_4^{(0)}/H\) today (\(H_4^{(0)}\) the Hubble
  constant); the first-order equations of the fold hold for
  \(\mathsf{A}^2 \ll 1\); 1 at the author's vacuum, which is not a
  solution of UM-4. Status: DERIVED.
\item
  \(\mathsf{U}\): The rate of change of \(b_4\) in units of \(H\),
  \(\mathsf{U} = b_4'/(Hb_4)\). Unit: dimensionless. Value:
  \(\mathsf{U} = \mathsf{A}\) with frozen extra times. Status: DERIVED.
\item
  \(\hat R\): The curvature of the metric \(\hat g_{ab}\) of the fold
  (FD-2): with four frame indices its Riemann tensor, with two its Ricci
  tensor \(\hat R_{ab}\), without indices its scalar curvature
  \(\hat R\); for frozen extra times the scalar curvature equals that of
  the four-dimensional metric \(\hat h_{\mu\nu}\). Unit: 1/length\(^2\).
  Status: STANDARD.
\item
  \(a\): A frame (orthonormal) index running over \(0, 1, \dots, 7\)
  (frame direction \((0)\) along \(\sigma\), \((1), (2), (3)\) along
  ordinary space, \((4)\) along the time, \((5), (6), (7)\) along the
  extra times); raised and lowered with the frame metric
  \(\eta_{ab} = \mathrm{diag}(1, 1, 1, 1, -1, -1, -1, -1)\); on the fold
  (\(a, b \in \{1, \dots, 7\}\)) the frame of the fold; in the tensors
  of the fold written with its coordinates \(x^a\) (its metric
  \(\hat g_{ab}\) and the tensors built from it) a coordinate index of
  the fold, \(a \in \{1, \dots, 7\}\); a repeated index is summed.
  Status: STANDARD.
\item
  \(b\): A frame index (as \(a\)). Status: STANDARD.
\item
  \(\kappa_7\): The gravitational coupling of the fold: to the first
  order in its curvature over \(H^2\) the fold obeys the Einstein
  equations of seven dimensions
  \(\hat G^a{}_b = \kappa_7\mathsf{t}^a{}_b\) with
  \(\kappa_7 = 5\kappa H/2\), the coupling of the massless graviton of
  the bulk bound to the fold; four-dimensional gravity on the fold has
  \(8\pi G_{\mathrm{N}} = \kappa_7/V_*\) (plus the induced term, which
  is negligible). Unit: GeV\(^{-5}\). Value: \(5\kappa H/2\). Status:
  DERIVED.
\item
  \(\rho_7\): The energy density of the matter of the fold above its
  tension in the seven-dimensional normalisation (minus the component
  \(\mathsf{t}^{(4)}{}_{(4)}\): an energy per unit volume of the six
  other directions of the fold); \(V_*\rho_7\) in the four-dimensional
  normalisation. Unit: GeV\(^7\). Status: DERIVED.
\item
  \(P_7\): The pressure along ordinary space of the matter of the fold
  above its tension in the seven-dimensional normalisation (the
  component \(\mathsf{t}^{(i)}{}_{(i)}\)). Unit: GeV\(^7\). Status:
  DERIVED.
\item
  \(q_7\): The pressure along the extra times of the matter of the fold
  above its tension in the seven-dimensional normalisation (the
  component \(\mathsf{t}^{(r)}{}_{(r)}\)). Unit: GeV\(^7\). Status:
  DERIVED.
\end{itemize}

The author's vacuum of the earlier version (a de Sitter expansion of
ordinary space at the rate \(H\) with constant \(b_4\), so that the
extra times contract) is not a solution of UM-4: the metric of its fold
is not Ricci flat, so the bulk is not a vacuum of the author's form
(FE-29), and as a fold with matter it would need, by the first-order
equations FE-71, the energy density \(\rho_7 = -3H^2/\kappa_7 < 0\) (a
violation of the weak energy condition) with the pressures
\(P_7 = q_7 = -3H^2/\kappa_7\) along ordinary space and the extra times;
its rate \(\mathsf{A} = 1\) is moreover outside the regime
\(\mathsf{A} \ll 1\) in which FE-71 holds. The primordial era of UM-4 is
the inflation of FE-87 and FE-88. Status: DERIVED
(\texttt{um4\_primordial.py}, PR0).

\hypertarget{the-static-block}{%
\subsection{The static block}\label{the-static-block}}

\[
a_4,\ b_4\ \text{constant}:\qquad R^{(a)(b)}{}_{(a)(b)} = -H^2\ \text{in all 28 frame planes},\qquad
\mathcal{E}_{(1)}{}^a{}_b = 21H^2\,\delta^a{}_b,\qquad \Lambda = -21H^2 \tag{FE-32}
\] \textbf{Symbols in (FE-32).}

\begin{itemize}
\tightlist
\item
  \(a_4\): The author's metric function of the time \(x_4\) that scales
  ordinary 3-space: on the fold \(\mathrm{e}^{a_4}\) is the scale factor
  of cosmology and \(a_4'\) the Hubble rate. Unit: dimensionless.
  Status: AUTHOR.
\item
  \(b_4\): The author's metric function of the time \(x_4\) in the
  extra-time components, \(b_4 > 0\), never identically 1 (the author's
  order of 2026-10-24); the scale factor (the factor of the lengths) of
  the extra times is \(\mathrm{e}^{-a_4}b_4\) on the fold and
  \(\mathrm{e}^{-a_4}b_4\,\mathrm{e}^{-H\lvert\sigma\rvert}\) at the
  distance \(\sigma\) from it (the square root of the warp factor times
  \(\mathrm{e}^{-a_4}b_4\)). Unit: dimensionless. Status: AUTHOR.
\item
  \(R^{MN}{}_{PQ},\ R_{MN},\ R\): The curvature of the metric
  \(g_{MN}\): with four indices the Riemann tensor \(R^M{}_{NPQ}\) (in
  the convention in which the commutator of two covariant derivatives,
  acting on a vector, is the Riemann tensor, with its last two indices
  those of the two derivatives, contracted with the vector), with two
  indices the Ricci tensor \(R_{MN} = R^P{}_{MPN}\), without indices the
  scalar curvature \(R = g^{MN}R_{MN}\); with frame indices
  \(R^{(a)(b)}{}_{(c)(d)}\) its frame components,
  \(R^{(a)(b)}{}_{(a)(b)}\) (no sum) the sectional curvature of the
  frame plane \((a, b)\). Unit: 1/length\(^2\). Status: STANDARD.
\item
  \(a\): A frame (orthonormal) index running over \(0, 1, \dots, 7\)
  (frame direction \((0)\) along \(\sigma\), \((1), (2), (3)\) along
  ordinary space, \((4)\) along the time, \((5), (6), (7)\) along the
  extra times); raised and lowered with the frame metric
  \(\eta_{ab} = \mathrm{diag}(1, 1, 1, 1, -1, -1, -1, -1)\); on the fold
  (\(a, b \in \{1, \dots, 7\}\)) the frame of the fold; in the tensors
  of the fold written with its coordinates \(x^a\) (its metric
  \(\hat g_{ab}\) and the tensors built from it) a coordinate index of
  the fold, \(a \in \{1, \dots, 7\}\); a repeated index is summed.
  Status: STANDARD.
\item
  \(b\): A frame index (as \(a\)). Status: STANDARD.
\item
  \(H\): The author's constant (an inverse length) that fixes the period
  \(\pi/(6H)\) of \(x_0\), the warp factor
  \(\mathrm{e}^{-2H\lvert\sigma\rvert}\) of the metric, the cosmological
  constant \(\Lambda = -21H^2\) and the tension
  \(\kappa\lambda_{\mathrm{F}} = 12H\) of the fold; in UM-4 a free
  parameter, bounded from below by the propagation of the gauge fields
  in the bulk. Unit: GeV (1/length). Value: \(H > 2534\) GeV, about 2.5
  TeV (from the invisible width of the Z boson). Status: AUTHOR.
\item
  \(\mathcal{E}_{(p)}{}^M{}_N\): The Lovelock tensor of order \(p\),
  \(\mathcal{E}_{(p)}{}^M{}_N = -2^{-(p+1)}\,\delta^{MP_1\cdots P_{2p}}_{NQ_1\cdots Q_{2p}}\,R^{Q_1Q_2}{}_{P_1P_2}\cdots R^{Q_{2p-1}Q_{2p}}{}_{P_{2p-1}P_{2p}}\)
  (with both indices down it is \(1/\sqrt{\lvert g\rvert}\) times the
  variation of \(\sqrt{\lvert g\rvert}\,\mathcal{L}_{(p)}\) with respect
  to \(g^{MN}\); \(\mathcal{E}_{(1)}\) the Einstein tensor,
  \(\mathcal{E}_{(2)}\) the Gauss-Bonnet tensor, \(\mathcal{E}_{(3)}\)
  the cubic Lovelock tensor); with frame indices its frame components.
  Unit: 1/length\(^{2p}\). Status: STANDARD.
\item
  \(p\): The Lovelock order \(p = 1, 2, 3\) (Einstein, Gauss-Bonnet,
  cubic), as an index and as a number (UM-4 keeps only \(p = 1\)); in
  the Kohn-Sham equations also a second label of a Kohn-Sham orbital,
  with the Kronecker delta of the two labels. Status: STANDARD.
\item
  \(\delta\): The Kronecker delta: with two indices \(\delta^M_N = 1\)
  if \(M = N\) and 0 otherwise (likewise \(\delta_{X\mathrm{L}}\) is 1
  if the label \(X\) of a half of dirac16complex is L and 0 otherwise);
  with two index strings of equal length the generalised Kronecker delta
  \(\delta^{M_1M_2\cdots}_{N_1N_2\cdots}\), the determinant of the
  square matrix of the Kronecker deltas \(\delta^{M_a}_{N_b}\) (\(+1\)
  or \(-1\) if the upper indices are an even or odd permutation of the
  lower ones, 0 otherwise). Status: STANDARD.
\item
  \(\Lambda\): The eight-dimensional cosmological constant of UM-4,
  \(\Lambda = -21H^2\): the value for which the static block is a vacuum
  of Einstein gravity (its Einstein tensor is \(21H^2\) in every
  direction). The author writes the pure number \(\Lambda_{\mathrm{N}}\)
  in
  \(0 = -\Lambda_{\mathrm{N}} + \sum_p H^{-2p}w_{(p)}\mathcal{E}_{(p)}\),
  so that \(\Lambda = -\Lambda_{\mathrm{N}}H^2\). Unit: GeV\(^2\).
  Value: \(-21H^2\); the author's original \(-10H^2\); UM-3: \(-7H^2\).
  Status: AUTHOR (the term), DERIVED (the value).
\end{itemize}

The static block (\(a_4\) and \(b_4\) constant; the units of
\(x_1, x_2, x_3\) and of \(x_5, x_6, x_7\) can make \(a_4 = 0\),
\(b_4 = 1\)): a space of constant curvature \(-H^2\) in all eight
directions (its Weyl tensor vanishes), on each side of the fold a region
of the eight-dimensional space of constant negative curvature of radius
\(1/H\) and signature (4, 4) (the analogue of anti-de Sitter space,
which has one time-like direction). Its Einstein tensor is \(21H^2\) in
every direction, so it is a vacuum of Einstein gravity with
\(\Lambda = -21H^2\) (in UM-3 the time was not warped, its planes with
the time were flat, and the static block was a vacuum only of the
Lovelock gravity with UM-3's weights). The fold of the static block is a
pure tension (FE-30). Status: DERIVED (\texttt{um4\_geometry.py}, G3 and
G4).

\[
\begin{aligned}
&\mathrm{e}^{a_4} = x_4^{\mathfrak{k}_1},\quad \mathrm{e}^{-a_4}b_4 = x_4^{\mathfrak{k}_2}:\quad
 3\mathfrak{k}_1 + 3\mathfrak{k}_2 = 1,\quad 3\mathfrak{k}_1^2 + 3\mathfrak{k}_2^2 = 1,\quad
 \{\mathfrak{k}_1, \mathfrak{k}_2\} = \Big\{\frac{1 + \sqrt5}{6}, \frac{1 - \sqrt5}{6}\Big\};\\
&\mathsf{V} = 0:\quad \hat G^{(4)}{}_{(4)} = -3a_4'^2
\end{aligned} \tag{FE-33}
\] \textbf{Symbols in (FE-33).}

\begin{itemize}
\tightlist
\item
  \(\mathrm{e}\): Euler's number, the base of the natural logarithm;
  \(\mathrm{e}^{X}\) is the exponential of \(X\). Unit: dimensionless.
  Value: 2.718281828. Status: STANDARD.
\item
  \(a_4\): The author's metric function of the time \(x_4\) that scales
  ordinary 3-space: on the fold \(\mathrm{e}^{a_4}\) is the scale factor
  of cosmology and \(a_4'\) the Hubble rate. Unit: dimensionless.
  Status: AUTHOR.
\item
  \(b_4\): The author's metric function of the time \(x_4\) in the
  extra-time components, \(b_4 > 0\), never identically 1 (the author's
  order of 2026-10-24); the scale factor (the factor of the lengths) of
  the extra times is \(\mathrm{e}^{-a_4}b_4\) on the fold and
  \(\mathrm{e}^{-a_4}b_4\,\mathrm{e}^{-H\lvert\sigma\rvert}\) at the
  distance \(\sigma\) from it (the square root of the warp factor times
  \(\mathrm{e}^{-a_4}b_4\)). Unit: dimensionless. Status: AUTHOR.
\item
  \(x_4\): The time; a prime on a function of \(x_4\) alone is
  \(\mathrm{d}/\mathrm{d}x_4\). Unit: length (= time, \(c = 1\)).
  Status: AUTHOR.
\item
  \(\mathfrak{k}_1,\ \mathfrak{k}_2\): The exponents of the Kasner folds
  (the expanding folds without matter of the author's form of the
  metric): \(\mathrm{e}^{a_4} = x_4^{\mathfrak{k}_1}\) and
  \(\mathrm{e}^{-a_4}b_4 = x_4^{\mathfrak{k}_2}\). Unit: dimensionless.
  Value: \((1 \pm \sqrt5)/6\). Status: DERIVED.
\item
  \(\mathsf{V}\): The rate of change of the scale factor
  \(\mathrm{e}^{-a_4}b_4\) of the extra times on the fold in units of
  \(H\), \(\mathsf{V} = \mathsf{U} - \mathsf{A}\); zero with frozen
  extra times. Unit: dimensionless. Status: DERIVED.
\item
  \(\hat G^\mu{}_\nu\): The Einstein tensor of the metric
  \(\hat g_{ab}\) of the fold (FD-2), with frame indices along its seven
  directions; with four-dimensional indices, \(\hat G^\mu{}_\nu\), its
  components along the four observable directions, which for frozen
  extra times are those of the Einstein tensor of the four-dimensional
  metric \(\hat h_{\mu\nu}\). Unit: 1/length\(^2\). Status: STANDARD.
\item
  \({}'\): The derivative of a function of the time \(x_4\) alone with
  respect to \(x_4\) (\(a_4' = \mathrm{d}a_4/\mathrm{d}x_4\); two primes
  the second derivative); on a function of \(\sigma\) alone (a profile)
  the derivative with respect to \(\sigma\). Status: STANDARD.
\end{itemize}

The vacua of the bulk of the author's form (FE-29: the Ricci-flat
folds): besides the static fold, the two Kasner folds in which ordinary
space and the extra times expand and contract with the powers
\(\mathfrak{k}_1\), \(\mathfrak{k}_2\) of the time (\(x_4 > 0\)). With
frozen extra times (\(\mathsf{V} = 0\), \(\mathrm{e}^{-a_4}b_4\)
constant) the energy component of the Einstein tensor of the fold is
\(-3a_4'^2\), so the only Ricci-flat fold with frozen extra times is
static: an expanding fold with frozen extra times needs matter, and its
bulk metric is FE-34. With constant rates the only vacuum of the
author's form is the static block. Status: DERIVED
(\texttt{um4\_geometry.py}, G6; \texttt{um4\_background.py}, BG1).

\[
\begin{aligned}
&\hat G^{(i)}{}_{(i)} = -\big(2a_4'' + 3\mathsf{c}'' + 3a_4'^2 + 6a_4'\mathsf{c}' + 6\mathsf{c}'^2\big),\qquad
 \hat G^{(4)}{}_{(4)} = -\big(3a_4'^2 + 9a_4'\mathsf{c}' + 3\mathsf{c}'^2\big),\\
&\hat G^{(r)}{}_{(r)} = -\big(3a_4'' + 2\mathsf{c}'' + 6a_4'^2 + 6a_4'\mathsf{c}' + 3\mathsf{c}'^2\big),\qquad
 \mathsf{c} = \ln\big(\mathrm{e}^{-a_4}b_4\big)
\end{aligned} \tag{FE-66}
\] \textbf{Symbols in (FE-66).}

\begin{itemize}
\tightlist
\item
  \(\hat G^\mu{}_\nu\): The Einstein tensor of the metric
  \(\hat g_{ab}\) of the fold (FD-2), with frame indices along its seven
  directions; with four-dimensional indices, \(\hat G^\mu{}_\nu\), its
  components along the four observable directions, which for frozen
  extra times are those of the Einstein tensor of the four-dimensional
  metric \(\hat h_{\mu\nu}\). Unit: 1/length\(^2\). Status: STANDARD.
\item
  \(i\): A generation index (\(i = 1, 2, 3\)), which also labels the
  three neutrino mass states; in a frame component such as
  \(\mathcal{S}^{(i)}{}_{(i)}\) one of the three directions of ordinary
  space (the three components are equal), and in a tensor of ordinary
  space such as \(\mathsf{h}_{ij}\) one of the directions
  \(x_1, x_2, x_3\); a repeated index is not summed unless a sum sign is
  written. Status: STANDARD.
\item
  \(r\): An extra-time direction of the fold: \(r \in \{5, 6, 7\}\)
  (\(x_5, x_6, x_7\)); not summed (the three components are equal).
  Status: STANDARD.
\item
  \(a_4\): The author's metric function of the time \(x_4\) that scales
  ordinary 3-space: on the fold \(\mathrm{e}^{a_4}\) is the scale factor
  of cosmology and \(a_4'\) the Hubble rate. Unit: dimensionless.
  Status: AUTHOR.
\item
  \({}'\): The derivative of a function of the time \(x_4\) alone with
  respect to \(x_4\) (\(a_4' = \mathrm{d}a_4/\mathrm{d}x_4\); two primes
  the second derivative); on a function of \(\sigma\) alone (a profile)
  the derivative with respect to \(\sigma\). Status: STANDARD.
\item
  \(\mathsf{c}\): The logarithm of the scale factor of the extra times
  on the fold, \(\mathsf{c} = \ln(\mathrm{e}^{-a_4}b_4)\);
  \(\mathsf{c}' = \mathsf{V}H\); frozen extra times have
  \(\mathsf{c}' = 0\). Unit: dimensionless. Status: DERIVED.
\item
  \(b_4\): The author's metric function of the time \(x_4\) in the
  extra-time components, \(b_4 > 0\), never identically 1 (the author's
  order of 2026-10-24); the scale factor (the factor of the lengths) of
  the extra times is \(\mathrm{e}^{-a_4}b_4\) on the fold and
  \(\mathrm{e}^{-a_4}b_4\,\mathrm{e}^{-H\lvert\sigma\rvert}\) at the
  distance \(\sigma\) from it (the square root of the warp factor times
  \(\mathrm{e}^{-a_4}b_4\)). Unit: dimensionless. Status: AUTHOR.
\item
  \(\mathrm{e}\): Euler's number, the base of the natural logarithm;
  \(\mathrm{e}^{X}\) is the exponential of \(X\). Unit: dimensionless.
  Value: 2.718281828. Status: STANDARD.
\end{itemize}

The Einstein tensor of the metric of the fold FD-2 for any \(a_4(x_4)\),
\(b_4(x_4)\), written with the logarithm \(\mathsf{c}\) of the scale
factor of the extra times (\(\mathsf{c}' = \mathsf{V}H\)); the
off-diagonal components vanish. With the equations of the fold FE-71 it
gives the Friedmann equations of the fold and the evolution of the
volume of the extra times (FE-72). Status: DERIVED
(\texttt{um4\_background.py}, BG2).

\hypertarget{the-bulk-metric-of-an-expanding-fold}{%
\subsection{The bulk metric of an expanding
fold}\label{the-bulk-metric-of-an-expanding-fold}}

\[
\begin{aligned}
&g_{MN}\,\mathrm{d}X^M\mathrm{d}X^N = \mathrm{d}\sigma^2 + \mathrm{e}^{-2H\lvert\sigma\rvert}\,\tilde g_{ab}\,\mathrm{d}x^a\mathrm{d}x^b,\\
&\tilde g_{ab} = \hat g_{ab} - \frac{\mathrm{e}^{2H\lvert\sigma\rvert} - 1}{5H^2}\Big(\hat R_{ab} - \frac{\hat R}{12}\,\hat g_{ab}\Big)
 + \mathrm{e}^{7H\lvert\sigma\rvert}\,\mathfrak{c}_{ab} + \dots
\end{aligned} \tag{FE-34}
\] \textbf{Symbols in (FE-34).}

\begin{itemize}
\tightlist
\item
  \(g_{MN}\): The metric tensor of the eight-dimensional spacetime
  (signature \(+,+,+,+,-,-,-,-\) in the order \(x_0, \dots, x_7\)):
  \(g_{MN}\,\mathrm{d}X^M\mathrm{d}X^N\) is the squared interval,
  \(g^{MN}\) its inverse, \(\lvert g\rvert\) the absolute value of its
  determinant (\(\sqrt{\lvert g\rvert}\,\mathrm{d}^8X\) is the proper
  volume element). Unit: dimensionless. Status: AUTHOR.
\item
  \(M\): A coordinate index of eight dimensions (also with a numbered
  copy \(M_1, M_2, \dots\)): it runs over \(x_0, x_1, \dots, x_7\); a
  repeated index is summed. Status: STANDARD.
\item
  \(N\): A coordinate index of eight dimensions (also
  \(N_1, N_2, \dots\)), running over \(x_0, \dots, x_7\); a repeated
  index is summed. Status: STANDARD.
\item
  \(\mathrm{d}\): The differential: \(\mathrm{d}\) in front of a
  quantity is an infinitesimal change of it (as \(\mathrm{d}\sigma\), a
  small step of \(\sigma\)); the quotient of two such changes is a
  derivative; \(\mathrm{d}\) with an exponent, in front of several
  variables, is their volume element, the product of their
  differentials. Status: STANDARD.
\item
  \(X = \{x_0, \dots, x_7\}\): The author's coordinates of the
  eight-dimensional spacetime, \(X = \{x_0, x_1, \dots, x_7\}\) with the
  hidden direction first (Dirac\_Einstein-Lovelock\_4+4\_Nash.nb,
  section DEFCHART); \(X^M\) is the coordinate with the index \(M\) and
  \(\mathrm{d}^8X = \mathrm{d}x_0\,\mathrm{d}x_1\cdots\mathrm{d}x_7\)
  the coordinate volume element. Unit: length (each coordinate). Status:
  AUTHOR.
\item
  \(\sigma\): The proper coordinate of the hidden direction:
  \(\mathrm{d}\sigma = \lvert\cot z\rvert\,\mathrm{d}x_0\),
  \(\sigma = \mathrm{sgn}(\pi/2 - z)\ln(\sin z)/(6H)\) on one period
  \(0 < z < \pi\); \(\sigma = 0\) is the fold \(z = \pi/2\),
  \(\sigma \to \mp\infty\) the tips \(z \to 0, \pi\); one period of
  \(x_0\) is the whole line \(-\infty < \sigma < \infty\). Unit: length.
  Status: DERIVED.
\item
  \(\mathrm{e}\): Euler's number, the base of the natural logarithm;
  \(\mathrm{e}^{X}\) is the exponential of \(X\). Unit: dimensionless.
  Value: 2.718281828. Status: STANDARD.
\item
  \(H\): The author's constant (an inverse length) that fixes the period
  \(\pi/(6H)\) of \(x_0\), the warp factor
  \(\mathrm{e}^{-2H\lvert\sigma\rvert}\) of the metric, the cosmological
  constant \(\Lambda = -21H^2\) and the tension
  \(\kappa\lambda_{\mathrm{F}} = 12H\) of the fold; in UM-4 a free
  parameter, bounded from below by the propagation of the gauge fields
  in the bulk. Unit: GeV (1/length). Value: \(H > 2534\) GeV, about 2.5
  TeV (from the invisible width of the Z boson). Status: AUTHOR.
\item
  \(\tilde g_{ab}\): The metric of the slice at the depth \(\sigma\)
  divided by the warp factor \(\mathrm{e}^{-2H\lvert\sigma\rvert}\):
  equal to the metric of the fold at \(\sigma = 0\). Unit:
  dimensionless. Status: DERIVED.
\item
  \(a\): A frame (orthonormal) index running over \(0, 1, \dots, 7\)
  (frame direction \((0)\) along \(\sigma\), \((1), (2), (3)\) along
  ordinary space, \((4)\) along the time, \((5), (6), (7)\) along the
  extra times); raised and lowered with the frame metric
  \(\eta_{ab} = \mathrm{diag}(1, 1, 1, 1, -1, -1, -1, -1)\); on the fold
  (\(a, b \in \{1, \dots, 7\}\)) the frame of the fold; in the tensors
  of the fold written with its coordinates \(x^a\) (its metric
  \(\hat g_{ab}\) and the tensors built from it) a coordinate index of
  the fold, \(a \in \{1, \dots, 7\}\); a repeated index is summed.
  Status: STANDARD.
\item
  \(b\): A frame index (as \(a\)). Status: STANDARD.
\item
  \(x\): The coordinates of the fold: \(x_1, x_2, x_3\) of ordinary
  space, \(x_4\) the time, \(x_5, x_6, x_7\) the extra times; a field of
  the fold is a function of the four observable coordinates
  \(x = (x_1, x_2, x_3, x_4)\), and
  \(\mathrm{d}^4x = \mathrm{d}x_1\,\mathrm{d}x_2\,\mathrm{d}x_3\,\mathrm{d}x_4\).
  Unit: length (each coordinate). Status: AUTHOR.
\item
  \(\hat g_{ab}\): The metric of the seven-dimensional fold (FD-2) in
  its coordinates \(x_1, \dots, x_7\) (or in its frame), \(\hat g^{ab}\)
  its inverse; \(\hat h_{\mu\nu}\) is its part along the four observable
  directions. Unit: dimensionless. Status: AUTHOR.
\item
  \(\hat R\): The curvature of the metric \(\hat g_{ab}\) of the fold
  (FD-2): with four frame indices its Riemann tensor, with two its Ricci
  tensor \(\hat R_{ab}\), without indices its scalar curvature
  \(\hat R\); for frozen extra times the scalar curvature equals that of
  the four-dimensional metric \(\hat h_{\mu\nu}\). Unit: 1/length\(^2\).
  Status: STANDARD.
\item
  \(\mathfrak{c}_{ab}\): A traceless conserved tensor of the fold that
  sets the homogeneous part of the bulk metric of an expanding fold: the
  Weyl tensor of the bulk that a field travelling through the bulk can
  leave (dark radiation); UM-4 takes \(\mathfrak{c}_{ab} = 0\). Unit:
  dimensionless. Value: 0. Status: ASSUMED.
\end{itemize}

The bulk metric of a fold with a metric \(\hat g_{ab}(x)\) of the fold
(\(x^a\) its seven coordinates), expanded in the number of derivatives
along the fold: each derivative counts once, the curvature of the fold
twice, so the expansion parameter is the curvature of the fold divided
by \(H^2\) (for a homogeneous fold \(\mathsf{A}^2 = (a_4'/H)^2\)). To
the first order the second term solves every component of the bulk
equations FE-1 (the diagonal components to the second order in the
derivatives, the mixed component \((0)(4)\) to the third) for the
homogeneous diagonal folds of the author's form, every \(a_4(x_4)\),
\(b_4(x_4)\) (the form of the second term is the general one of the
derivative expansion; for folds that depend on ordinary space it is not
checked here), and it vanishes at the fold, so that
\(\tilde g_{ab}(\sigma = 0) = \hat g_{ab}\). The third term is the
homogeneous solution: a traceless conserved tensor
\(\mathfrak{c}_{ab}(x)\) of the fold times
\(\mathrm{e}^{7H\lvert\sigma\rvert}\), the bulk Weyl tensor that a field
travelling through the bulk can leave (the dark radiation of the brane
of Randall and Sundrum); UM-4 takes \(\mathfrak{c}_{ab} = 0\) (ASSUMED:
no dark radiation). The dots are the terms of higher order. Status:
DERIVED (\texttt{um4\_fold.py}, F1), ASSUMED
(\(\mathfrak{c}_{ab} = 0\)).

\hypertarget{the-equations-of-the-fold}{%
\subsection{The equations of the fold}\label{the-equations-of-the-fold}}

\[
\kappa\,\mathcal{S}^a{}_b = -12H\,\delta^a{}_b + \frac{2}{5H}\,\hat G^a{}_b + \dots,\qquad
\mathcal{S}^a{}_b = -\lambda_{\mathrm{F}}\,\delta^a{}_b + \mathsf{t}^a{}_b:\qquad
\hat G^a{}_b = \kappa_7\,\mathsf{t}^a{}_b,\qquad \kappa_7 = \frac52\,\kappa H \tag{FE-71}
\] \textbf{Symbols in (FE-71).}

\begin{itemize}
\tightlist
\item
  \(\kappa\): The eight-dimensional gravitational coupling,
  \(\kappa = 8\pi G_8\). Unit: GeV\(^{-6}\). Status: AUTHOR.
\item
  \(\mathcal{S}^a{}_b\): The thin source of the fold: the
  energy-momentum tensor concentrated on the fold, integrated across it,
  \(\kappa\mathcal{S}^a{}_b = \int_{0^-}^{0^+}\mathrm{d}\sigma\,\mathcal{E}_{(1)}{}^a{}_b\)
  (the junction condition of Israel); \(\mathcal{S}^4{}_4\) is minus its
  energy density, \(\mathcal{S}^1{}_1\) and \(\mathcal{S}^5{}_5\) its
  pressures along ordinary space and along the extra times (frame
  components along the fold, eight-dimensional normalisation: an energy
  per unit volume of the six other directions of the fold). Unit:
  GeV\(^7\). Value: \(\kappa\mathcal{S}^a{}_b = -12H\delta^a{}_b\) for
  every \(a_4\), \(b_4\) of the author's form. Status: DERIVED.
\item
  \(a\): A frame (orthonormal) index running over \(0, 1, \dots, 7\)
  (frame direction \((0)\) along \(\sigma\), \((1), (2), (3)\) along
  ordinary space, \((4)\) along the time, \((5), (6), (7)\) along the
  extra times); raised and lowered with the frame metric
  \(\eta_{ab} = \mathrm{diag}(1, 1, 1, 1, -1, -1, -1, -1)\); on the fold
  (\(a, b \in \{1, \dots, 7\}\)) the frame of the fold; in the tensors
  of the fold written with its coordinates \(x^a\) (its metric
  \(\hat g_{ab}\) and the tensors built from it) a coordinate index of
  the fold, \(a \in \{1, \dots, 7\}\); a repeated index is summed.
  Status: STANDARD.
\item
  \(b\): A frame index (as \(a\)). Status: STANDARD.
\item
  \(H\): The author's constant (an inverse length) that fixes the period
  \(\pi/(6H)\) of \(x_0\), the warp factor
  \(\mathrm{e}^{-2H\lvert\sigma\rvert}\) of the metric, the cosmological
  constant \(\Lambda = -21H^2\) and the tension
  \(\kappa\lambda_{\mathrm{F}} = 12H\) of the fold; in UM-4 a free
  parameter, bounded from below by the propagation of the gauge fields
  in the bulk. Unit: GeV (1/length). Value: \(H > 2534\) GeV, about 2.5
  TeV (from the invisible width of the Z boson). Status: AUTHOR.
\item
  \(\delta\): The Kronecker delta: with two indices \(\delta^M_N = 1\)
  if \(M = N\) and 0 otherwise (likewise \(\delta_{X\mathrm{L}}\) is 1
  if the label \(X\) of a half of dirac16complex is L and 0 otherwise);
  with two index strings of equal length the generalised Kronecker delta
  \(\delta^{M_1M_2\cdots}_{N_1N_2\cdots}\), the determinant of the
  square matrix of the Kronecker deltas \(\delta^{M_a}_{N_b}\) (\(+1\)
  or \(-1\) if the upper indices are an even or odd permutation of the
  lower ones, 0 otherwise). Status: STANDARD.
\item
  \(\hat G^\mu{}_\nu\): The Einstein tensor of the metric
  \(\hat g_{ab}\) of the fold (FD-2), with frame indices along its seven
  directions; with four-dimensional indices, \(\hat G^\mu{}_\nu\), its
  components along the four observable directions, which for frozen
  extra times are those of the Einstein tensor of the four-dimensional
  metric \(\hat h_{\mu\nu}\). Unit: 1/length\(^2\). Status: STANDARD.
\item
  \(\lambda_{\mathrm{F}}\): The tension that the static block requires
  of the fold, \(\kappa\lambda_{\mathrm{F}} = 12H\): the total vacuum
  energy of the seven-dimensional fold, made of the tension of the fold
  fluid and the vacuum energy of the fields of the fold (both have the
  pressure minus their energy density along ordinary space and along the
  extra times), per unit volume of the six other directions of the fold
  (eight-dimensional normalisation); in the four-dimensional
  normalisation it is \(\lambda_{\mathrm{F}}V_*\), and the excess of the
  total over \(\lambda_{\mathrm{F}}V_*\) is the dark energy
  \(\rho_\Lambda\). Unit: GeV\(^7\). Value: \(12H/\kappa\);
  \(\lambda_{\mathrm{F}}V_* = 30H^2M_{\mathrm{red}}^2 = 1.1\times10^{45}\)
  GeV\(^4\) for \(H = 2534\) GeV. Status: DERIVED.
\item
  \(\mathsf{t}^a{}_b,\ \mathsf{t}\): The matter of the fold above its
  tension: the thin source minus the tension,
  \(\mathcal{S}^a{}_b = -\lambda_{\mathrm{F}}\delta^a{}_b + \mathsf{t}^a{}_b\)
  (eight-dimensional normalisation); \(\mathsf{t} = \mathsf{t}^c{}_c\)
  its trace; with frozen extra times
  \(\mathsf{t}^a{}_b = \mathrm{diag}(P_7, P_7, P_7, -\rho_7, q_7, q_7, q_7)\).
  Unit: GeV\(^7\). Status: DERIVED.
\item
  \(\kappa_7\): The gravitational coupling of the fold: to the first
  order in its curvature over \(H^2\) the fold obeys the Einstein
  equations of seven dimensions
  \(\hat G^a{}_b = \kappa_7\mathsf{t}^a{}_b\) with
  \(\kappa_7 = 5\kappa H/2\), the coupling of the massless graviton of
  the bulk bound to the fold; four-dimensional gravity on the fold has
  \(8\pi G_{\mathrm{N}} = \kappa_7/V_*\) (plus the induced term, which
  is negligible). Unit: GeV\(^{-5}\). Value: \(5\kappa H/2\). Status:
  DERIVED.
\end{itemize}

The junction condition FE-4 with the bulk metric FE-34: the thin source
is the tension of FE-30 plus \(2/(5H)\) times the Einstein tensor of the
fold, up to terms of the fourth order in the derivatives (computed for
the space, time and extra-time components). Written with the tension
\(\lambda_{\mathrm{F}}\) and the matter \(\mathsf{t}^a{}_b\) of the
fold, the fold obeys the Einstein equations of seven dimensions with the
coupling \(\kappa_7 = 5\kappa H/2\) and no cosmological constant of its
own (for the critical tension; a tension above it gives the fold the
cosmological constant \(\kappa_7(\lambda - \lambda_{\mathrm{F}})\),
FE-65). The fold is a seven-dimensional brane on which gravity is
Einstein's to the first order in its curvature over \(H^2\), as for the
brane of Randall and Sundrum. Integrated over the extra times (volume
\(V_*\), four-dimensional densities \(\rho = V_*\rho_7\)) Newton's
constant is \(8\pi G_{\mathrm{N}} = \kappa_7/V_*\), and the tension in
four-dimensional normalisation is
\(\lambda_{\mathrm{F}}V_* = 30H^2M_{\mathrm{red}}^2\). Status: DERIVED
(\texttt{um4\_fold.py}, F2 and F5).

\[
\begin{aligned}
&\hat G_{ab} = \kappa_7\,\mathsf{t}_{ab} + \kappa^2\,\Pi_{ab} - \mathfrak{E}_{ab},\\
&\Pi_{ab} = -\frac14\,\mathsf{t}_{ac}\mathsf{t}^c{}_b + \frac{1}{24}\,\mathsf{t}\,\mathsf{t}_{ab} + \frac18\,\hat g_{ab}\,\mathsf{t}_{cd}\mathsf{t}^{cd}
 - \frac{1}{48}\,\hat g_{ab}\,\mathsf{t}^2,\qquad \mathsf{t} = \mathsf{t}^c{}_c
\end{aligned} \tag{FE-35}
\] \textbf{Symbols in (FE-35).}

\begin{itemize}
\tightlist
\item
  \(\hat G^\mu{}_\nu\): The Einstein tensor of the metric
  \(\hat g_{ab}\) of the fold (FD-2), with frame indices along its seven
  directions; with four-dimensional indices, \(\hat G^\mu{}_\nu\), its
  components along the four observable directions, which for frozen
  extra times are those of the Einstein tensor of the four-dimensional
  metric \(\hat h_{\mu\nu}\). Unit: 1/length\(^2\). Status: STANDARD.
\item
  \(a\): A frame (orthonormal) index running over \(0, 1, \dots, 7\)
  (frame direction \((0)\) along \(\sigma\), \((1), (2), (3)\) along
  ordinary space, \((4)\) along the time, \((5), (6), (7)\) along the
  extra times); raised and lowered with the frame metric
  \(\eta_{ab} = \mathrm{diag}(1, 1, 1, 1, -1, -1, -1, -1)\); on the fold
  (\(a, b \in \{1, \dots, 7\}\)) the frame of the fold; in the tensors
  of the fold written with its coordinates \(x^a\) (its metric
  \(\hat g_{ab}\) and the tensors built from it) a coordinate index of
  the fold, \(a \in \{1, \dots, 7\}\); a repeated index is summed.
  Status: STANDARD.
\item
  \(b\): A frame index (as \(a\)). Status: STANDARD.
\item
  \(\kappa_7\): The gravitational coupling of the fold: to the first
  order in its curvature over \(H^2\) the fold obeys the Einstein
  equations of seven dimensions
  \(\hat G^a{}_b = \kappa_7\mathsf{t}^a{}_b\) with
  \(\kappa_7 = 5\kappa H/2\), the coupling of the massless graviton of
  the bulk bound to the fold; four-dimensional gravity on the fold has
  \(8\pi G_{\mathrm{N}} = \kappa_7/V_*\) (plus the induced term, which
  is negligible). Unit: GeV\(^{-5}\). Value: \(5\kappa H/2\). Status:
  DERIVED.
\item
  \(\mathsf{t}^a{}_b,\ \mathsf{t}\): The matter of the fold above its
  tension: the thin source minus the tension,
  \(\mathcal{S}^a{}_b = -\lambda_{\mathrm{F}}\delta^a{}_b + \mathsf{t}^a{}_b\)
  (eight-dimensional normalisation); \(\mathsf{t} = \mathsf{t}^c{}_c\)
  its trace; with frozen extra times
  \(\mathsf{t}^a{}_b = \mathrm{diag}(P_7, P_7, P_7, -\rho_7, q_7, q_7, q_7)\).
  Unit: GeV\(^7\). Status: DERIVED.
\item
  \(\kappa\): The eight-dimensional gravitational coupling,
  \(\kappa = 8\pi G_8\). Unit: GeV\(^{-6}\). Status: AUTHOR.
\item
  \(\Pi_{ab}\): The correction of high energies in the exact equations
  of the fold, quadratic in the matter \(\mathsf{t}_{ab}\) of the fold.
  Unit: GeV\(^{14}\). Status: DERIVED.
\item
  \(\mathfrak{E}_{ab}\): The projection of the Weyl tensor of the bulk
  onto the fold (its electric part), which enters the exact equations of
  the fold; zero for a bulk of constant curvature, and not fixed by the
  fold alone. Unit: 1/length\(^2\). Status: DERIVED.
\item
  \(c\): A frame index (as \(a\)). Status: STANDARD.
\item
  \(d\): A frame index (as \(a\)). Status: STANDARD.
\item
  \(\hat g_{ab}\): The metric of the seven-dimensional fold (FD-2) in
  its coordinates \(x_1, \dots, x_7\) (or in its frame), \(\hat g^{ab}\)
  its inverse; \(\hat h_{\mu\nu}\) is its part along the four observable
  directions. Unit: dimensionless. Status: AUTHOR.
\end{itemize}

The exact equations of the fold: Gauss's equation for the fold in the
bulk with Israel's condition FE-4 and the tension of FE-30 (the method
of T. Shiromizu, K. Maeda and M. Sasaki, Phys. Rev.~D 62 (2000) 024012,
here in eight dimensions). The quadratic term \(\Pi_{ab}\) is the
correction of high energies, and \(\mathfrak{E}_{ab}\) is the projection
of the Weyl tensor of the bulk onto the fold (zero for a bulk of
constant curvature; not fixed by the fold alone). Relative to
\(\kappa_7\mathsf{t}_{ab}\) the term \(\kappa^2\Pi_{ab}\) is of the
order \((24/5)\rho/(\lambda_{\mathrm{F}}V_*) = (12/25)\mathsf{A}^2\) for
a fold that expands at the rate \(\mathsf{A}H\): the equations of the
fold are those of general relativity in seven dimensions whenever the
expansion rate is small compared to \(H\). With frozen extra times the
conservation of \(\Pi_{ab}\) fails for matter other than vacuum energy,
radiation and stiff matter (its divergence is \(a_4'\rho^2\) times
\(-3(w - 1)(w + 1)(3w - 1)/32\) for a fluid with
\(\mathcal{P} = w\rho\)): for dust \(\mathfrak{E}_{ab}\) must be
non-zero at the order \(\rho^2/\lambda_{\mathrm{F}}\) (its value needs
the bulk solution: OPEN). Status: DERIVED (\texttt{um4\_fold.py}, F3, F5
and F6; the identity verified for a diagonal \(\mathsf{t}_{ab}\) with
independent entries).

\hypertarget{the-volume-of-the-extra-times}{%
\subsection{The volume of the extra
times}\label{the-volume-of-the-extra-times}}

\[
\begin{aligned}
&\mathsf{c}' = 0:\quad 3a_4'^2 = \kappa_7\rho_7,\quad -\big(2a_4'' + 3a_4'^2\big) = \kappa_7P_7,\quad
 -3\big(a_4'' + 2a_4'^2\big) = \kappa_7q_7\quad\Rightarrow\quad q_7 = -\frac{\rho_7 - 3P_7}{2};\\
&q_7 = P_7 = 0,\ \mathsf{c}' = 0:\ \mathsf{c}'' = \frac35\,a_4'^2
\end{aligned} \tag{FE-72}
\] \textbf{Symbols in (FE-72).}

\begin{itemize}
\tightlist
\item
  \(\mathsf{c}\): The logarithm of the scale factor of the extra times
  on the fold, \(\mathsf{c} = \ln(\mathrm{e}^{-a_4}b_4)\);
  \(\mathsf{c}' = \mathsf{V}H\); frozen extra times have
  \(\mathsf{c}' = 0\). Unit: dimensionless. Status: DERIVED.
\item
  \({}'\): The derivative of a function of the time \(x_4\) alone with
  respect to \(x_4\) (\(a_4' = \mathrm{d}a_4/\mathrm{d}x_4\); two primes
  the second derivative); on a function of \(\sigma\) alone (a profile)
  the derivative with respect to \(\sigma\). Status: STANDARD.
\item
  \(a_4\): The author's metric function of the time \(x_4\) that scales
  ordinary 3-space: on the fold \(\mathrm{e}^{a_4}\) is the scale factor
  of cosmology and \(a_4'\) the Hubble rate. Unit: dimensionless.
  Status: AUTHOR.
\item
  \(\kappa_7\): The gravitational coupling of the fold: to the first
  order in its curvature over \(H^2\) the fold obeys the Einstein
  equations of seven dimensions
  \(\hat G^a{}_b = \kappa_7\mathsf{t}^a{}_b\) with
  \(\kappa_7 = 5\kappa H/2\), the coupling of the massless graviton of
  the bulk bound to the fold; four-dimensional gravity on the fold has
  \(8\pi G_{\mathrm{N}} = \kappa_7/V_*\) (plus the induced term, which
  is negligible). Unit: GeV\(^{-5}\). Value: \(5\kappa H/2\). Status:
  DERIVED.
\item
  \(\rho_7\): The energy density of the matter of the fold above its
  tension in the seven-dimensional normalisation (minus the component
  \(\mathsf{t}^{(4)}{}_{(4)}\): an energy per unit volume of the six
  other directions of the fold); \(V_*\rho_7\) in the four-dimensional
  normalisation. Unit: GeV\(^7\). Status: DERIVED.
\item
  \(P_7\): The pressure along ordinary space of the matter of the fold
  above its tension in the seven-dimensional normalisation (the
  component \(\mathsf{t}^{(i)}{}_{(i)}\)). Unit: GeV\(^7\). Status:
  DERIVED.
\item
  \(q_7\): The pressure along the extra times of the matter of the fold
  above its tension in the seven-dimensional normalisation (the
  component \(\mathsf{t}^{(r)}{}_{(r)}\)). Unit: GeV\(^7\). Status:
  DERIVED.
\end{itemize}

The fold with frozen extra times (\(\mathsf{c}' = 0\), FE-66 with FE-71;
\(\mathsf{t}^a{}_b = \mathrm{diag}(P_7, P_7, P_7, -\rho_7, q_7, q_7, q_7)\)):
the first two equations are the Friedmann equations of four-dimensional
general relativity, and the component along the extra times requires the
pressure \(q_7 = -(\rho_7 - 3P_7)/2\) along them (the radion condition).
Radiation satisfies it by itself; dust (\(P_7 = q_7 = 0\)) does not:
without a pressure along the extra times the volume of the extra times
accelerates, \(\mathsf{c}'' = 3a_4'^2/5\), and through a matter era the
ratio \(\mathsf{c}'/a_4'\) grows from 0 to 0.75, 0.98 and 1.00 after
one, two and three e-folds (the seven-dimensional dust solution, in
which the extra times expand as ordinary space); a vacuum energy
(\(q_7 = P_7 = -\rho_7\)) drives it too, \(\mathsf{c}'' = 6a_4'^2/5\).
Newton's constant
\(8\pi G_{\mathrm{N}} = \kappa_7/(V_{\mathrm{t}}\mathrm{e}^{3\mathsf{c}})\)
would then change at the rate \(-3\mathsf{c}' = -3H_4\), which lunar
laser ranging (the relative rate of change of Newton's constant is
\((-5.0 \pm 9.6)\times10^{-15}\) per year; L. Biskupek, J. Müller and
J.-M. Torre, Universe 7 (2021) 34) excludes by a factor of about 8500.
The fold fluid supplies the missing pressure (IN-10, FE-65). Status:
DERIVED (\texttt{um4\_fold.py}, F4; \texttt{um4\_background.py}, BG2 to
BG4), ASSUMED (the fold fluid).

\hypertarget{propagation-in-the-bulk}{%
\section{Propagation in the bulk}\label{propagation-in-the-bulk}}

In the metric of UM-4 the warp factor
\(\mathrm{e}^{-2H\lvert\sigma\rvert}\) multiplies the whole metric of
the fold, the time included (FE-27). Every slice
\(\sigma = \text{const}\) of the static block is therefore a copy of the
flat fold, scaled by \(\mathrm{e}^{-H\lvert\sigma\rvert}\) in all seven
directions, and its light cone is the light cone of the fold: at every
depth a signal of light moves along ordinary space by one unit of
\(x_1, x_2, x_3\) per unit of the time \(x_4\), exactly as on the fold.
For every field of the bulk the dependence on \(\sigma\) then separates
from the dependence on the seven coordinates of the fold, and the field
of the fold that remains obeys the Lorentz-invariant equation of seven
dimensions with a mass \(m\) that its profile along \(\sigma\) fixes.
This section writes the equations of the profiles of every field of the
bulk on the static block (\(a_4 = 0\), \(b_4 = 1\) by the choice of the
units of \(x_1, \dots, x_7\)), the photon and the gravitational wave as
waves of four dimensions, the propagators of the bulk from one point of
the fold to another, Newton's and Coulomb's laws, the widths of the W
and Z bosons, and the causal structure of the bulk.

\hypertarget{the-modes-of-the-fields-of-the-bulk}{%
\subsection{The modes of the fields of the
bulk}\label{the-modes-of-the-fields-of-the-bulk}}

\[
\begin{aligned}
&\phi_{\mathrm{bulk}} = u(\sigma)\,\mathrm{e}^{\mathrm{i}(kx_3 - \omega x_4)}:\qquad
u'' - 7H\,\mathrm{sgn}(\sigma)\,u' + \mathrm{e}^{2H\lvert\sigma\rvert}\,m^2\,u = 0,\qquad m^2 = \omega^2 - k^2,\\
&u_0 = 1,\quad \int_{-\infty}^{\infty}\mathrm{d}\sigma\,\mathrm{e}^{-5H\lvert\sigma\rvert}\,u_0^2 = \frac{2}{5H}
\end{aligned} \tag{FE-73}
\] \textbf{Symbols in (FE-73).}

\begin{itemize}
\tightlist
\item
  \(\phi_{\mathrm{bulk}}\): A massless scalar field of the bulk,
  \(\nabla_M\nabla^M\phi_{\mathrm{bulk}} = 0\) (an illustration: each
  polarisation of the graviton obeys its equation). Unit: GeV\(^3\).
  Status: STANDARD.
\item
  \(u(\sigma),\ u_0\): The profile along the hidden direction of a mode
  of a field of the bulk (a scalar, a gauge field or a polarisation of
  the graviton): the field is \(u(\sigma)\) times a field of the fold;
  \(u_0 = 1\) is the profile of the massless mode (the zero mode); in a
  propagator the response of the field to a unit source on the fold.
  Unit: dimensionless (a relative amplitude); length in a propagator.
  Status: DERIVED.
\item
  \(\sigma\): The proper coordinate of the hidden direction:
  \(\mathrm{d}\sigma = \lvert\cot z\rvert\,\mathrm{d}x_0\),
  \(\sigma = \mathrm{sgn}(\pi/2 - z)\ln(\sin z)/(6H)\) on one period
  \(0 < z < \pi\); \(\sigma = 0\) is the fold \(z = \pi/2\),
  \(\sigma \to \mp\infty\) the tips \(z \to 0, \pi\); one period of
  \(x_0\) is the whole line \(-\infty < \sigma < \infty\). Unit: length.
  Status: DERIVED.
\item
  \(\mathrm{e}\): Euler's number, the base of the natural logarithm;
  \(\mathrm{e}^{X}\) is the exponential of \(X\). Unit: dimensionless.
  Value: 2.718281828. Status: STANDARD.
\item
  \(\mathrm{i}\): The imaginary unit, \(\mathrm{i}^2 = -1\). Unit:
  dimensionless. Status: STANDARD.
\item
  \(k\): A wave number: the momentum of a plane wave along \(x_3\),
  \(\mathrm{e}^{\mathrm{i}(kx_3 - \omega x_4)}\); in the cosmological
  perturbations the comoving wave number of a Fourier mode,
  \(k = \lvert\vec k\rvert\). Unit: GeV (1/length). Status: STANDARD.
\item
  \(x_3\): The third coordinate of ordinary 3-space. Unit: length.
  Status: AUTHOR.
\item
  \(\omega\): A frequency: the energy of a mode
  \(\propto \mathrm{e}^{-\mathrm{i}\omega x_4}\) (\(\omega_o\) of the
  orbital \(o\)). Unit: GeV. Status: STANDARD.
\item
  \(x_4\): The time; a prime on a function of \(x_4\) alone is
  \(\mathrm{d}/\mathrm{d}x_4\). Unit: length (= time, \(c = 1\)).
  Status: AUTHOR.
\item
  \({}'\): The derivative of a function of the time \(x_4\) alone with
  respect to \(x_4\) (\(a_4' = \mathrm{d}a_4/\mathrm{d}x_4\); two primes
  the second derivative); on a function of \(\sigma\) alone (a profile)
  the derivative with respect to \(\sigma\). Status: STANDARD.
\item
  \(H\): The author's constant (an inverse length) that fixes the period
  \(\pi/(6H)\) of \(x_0\), the warp factor
  \(\mathrm{e}^{-2H\lvert\sigma\rvert}\) of the metric, the cosmological
  constant \(\Lambda = -21H^2\) and the tension
  \(\kappa\lambda_{\mathrm{F}} = 12H\) of the fold; in UM-4 a free
  parameter, bounded from below by the propagation of the gauge fields
  in the bulk. Unit: GeV (1/length). Value: \(H > 2534\) GeV, about 2.5
  TeV (from the invisible width of the Z boson). Status: AUTHOR.
\item
  \(\mathrm{sgn}\): The sign function: \(+1\) for a positive argument
  and \(-1\) for a negative one. Status: STANDARD.
\item
  \(m\): A mass: \(m_i\) the mass of the neutrino mass state \(i\);
  \(\Delta m_{ij}^2 = m_i^2 - m_j^2\) the squared-mass difference of the
  states \(i\) and \(j\) (\(\Delta m_{21}^2\), \(\Delta m_{31}^2\),
  \(\Delta m_{32}^2\)); \(\Delta m^2\) without labels the difference of
  the two states of a resonance, \(m_2^2 - m_1^2\), multiplied by
  \(v_{\mathrm{loc}}^2/v^2\) where this differs from 1; \(m\) without a
  label the mass of a mode of a field of the bulk as the fold sees it,
  \(m^2 = \omega^2 - k^2\); for the W and Z bosons their complex pole
  mass, \(m^2 = m_0^2 - \mathrm{i}m_0\Gamma_{\mathrm{bulk}}\) to the
  leading order; in the relic abundance of the particles \(\chi\) their
  mass. Unit: GeV. Status: STANDARD.
\item
  \(\infty\): Infinity: as a limit of an integral or a sum it means that
  the range has no bound. Status: STANDARD.
\item
  \(\mathrm{d}\): The differential: \(\mathrm{d}\) in front of a
  quantity is an infinitesimal change of it (as \(\mathrm{d}\sigma\), a
  small step of \(\sigma\)); the quotient of two such changes is a
  derivative; \(\mathrm{d}\) with an exponent, in front of several
  variables, is their volume element, the product of their
  differentials. Status: STANDARD.
\end{itemize}

A massless scalar field of the bulk,
\(\nabla_M\nabla^M\phi_{\mathrm{bulk}} = 0\) (an illustration: it is the
equation of each polarisation of the graviton, FE-76), for a wave of the
frequency \(\omega\) and the wave number \(k\) along \(x_3\),
independent of the extra times. The profile \(u(\sigma)\) obeys the
equation in the middle (on both sides of the fold; the term
\(-7H\,\mathrm{sgn}(\sigma)u'\) comes from the volume element
\(\sqrt{\lvert g\rvert} \propto \mathrm{e}^{-7H\lvert\sigma\rvert}\)),
in which \(\omega\) and \(k\) appear only through
\(m^2 = \omega^2 - k^2\): the mode is a field of the fold of the mass
\(m\), with the Lorentz-invariant dispersion \(\omega^2 = k^2 + m^2\),
whatever its profile. The mode \(m = 0\) has the profile \(u_0 = 1\),
the same amplitude at every depth; it is normalisable, because its
kinetic term carries the weight
\(\sqrt{\lvert g\rvert}\,\lvert g^{x_4x_4}\rvert \propto \mathrm{e}^{-5H\lvert\sigma\rvert}\)
with the finite integral \(2/(5H)\): the zero mode is a massless field
of the seven-dimensional fold, bound to it, and it travels at the speed
of light at every wavelength. The other modes form a continuum of masses
from \(m = 0\) upwards (no gap) and move slower than light, with the
group velocity \(k/\omega\). (In UM-3, whose time was not warped, the
same equation had \(\omega^2\) without and \(k^2\) with the factor
\(\mathrm{e}^{2H\lvert\sigma\rvert}\), and the zero mode had
\(\omega^2 = \frac32k^2\): 1.22 times the speed of light.) Status:
DERIVED (\texttt{um4\_bulk.py}, B1: the profile equation exactly on both
halves; the lowest \(\omega^2\) of the operator in a box
\(\lvert\sigma\rvert < 6/H\) equals \(k^2\) to \(4\times10^{-12}\) for
\(k = 0.1H\) to \(2H\); B10: the massive modes, Bessel functions of the
warped distance from the fold, a continuum from zero mass).

\[
\begin{aligned}
&B_\sigma = 0,\quad B_a = u(\sigma)\,\hat B_a(x):\qquad
u'' - 5H\,\mathrm{sgn}(\sigma)\,u' + \mathrm{e}^{2H\lvert\sigma\rvert}\,m^2\,u = 0,\qquad
\eta^{ac}\,\partial_a\hat B_{cb} = m^2\,\hat B_b,\\
&u_0 = 1,\quad \int_{-\infty}^{\infty}\mathrm{d}\sigma\,\mathrm{e}^{-3H\lvert\sigma\rvert}\,u_0^2 = \frac{2}{3H}
\end{aligned} \tag{FE-74}
\] \textbf{Symbols in (FE-74).}

\begin{itemize}
\tightlist
\item
  \(B_M,\ B_\mu,\ B_{MN}\): The hypercharge gauge field of \(U(1)_Y\):
  in the bulk \(B_M\) (geometric normalisation), on the fold the
  canonically normalised four-dimensional field \(B_\mu(x)\) of its
  massless mode, \(g_1B_\mu(x) = B_\mu(\sigma = 0, x)\); \(B_{MN}\),
  \(B_{\mu\nu}\) their field strengths. Unit: GeV (field), GeV\(^2\)
  (field strength). Status: AUTHOR.
\item
  \(\sigma\): The label \(\sigma\) in \(\partial_\sigma\): the
  derivative with respect to the proper hidden coordinate \(\sigma\).
  Status: STANDARD.
\item
  \(a\): A frame (orthonormal) index running over \(0, 1, \dots, 7\)
  (frame direction \((0)\) along \(\sigma\), \((1), (2), (3)\) along
  ordinary space, \((4)\) along the time, \((5), (6), (7)\) along the
  extra times); raised and lowered with the frame metric
  \(\eta_{ab} = \mathrm{diag}(1, 1, 1, 1, -1, -1, -1, -1)\); on the fold
  (\(a, b \in \{1, \dots, 7\}\)) the frame of the fold; in the tensors
  of the fold written with its coordinates \(x^a\) (its metric
  \(\hat g_{ab}\) and the tensors built from it) a coordinate index of
  the fold, \(a \in \{1, \dots, 7\}\); a repeated index is summed.
  Status: STANDARD.
\item
  \(u(\sigma),\ u_0\): The profile along the hidden direction of a mode
  of a field of the bulk (a scalar, a gauge field or a polarisation of
  the graviton): the field is \(u(\sigma)\) times a field of the fold;
  \(u_0 = 1\) is the profile of the massless mode (the zero mode); in a
  propagator the response of the field to a unit source on the fold.
  Unit: dimensionless (a relative amplitude); length in a propagator.
  Status: DERIVED.
\item
  \(\sigma\): The proper coordinate of the hidden direction:
  \(\mathrm{d}\sigma = \lvert\cot z\rvert\,\mathrm{d}x_0\),
  \(\sigma = \mathrm{sgn}(\pi/2 - z)\ln(\sin z)/(6H)\) on one period
  \(0 < z < \pi\); \(\sigma = 0\) is the fold \(z = \pi/2\),
  \(\sigma \to \mp\infty\) the tips \(z \to 0, \pi\); one period of
  \(x_0\) is the whole line \(-\infty < \sigma < \infty\). Unit: length.
  Status: DERIVED.
\item
  \(\hat B_a,\ \hat B_{ab}\): The field of the fold of a mode of the
  hypercharge field of the bulk (of any gauge field of the bulk to the
  first order in the field), \(B_a = u(\sigma)\hat B_a(x)\);
  \(\hat B_{ab} = \partial_a\hat B_b - \partial_b\hat B_a\) its field
  strength. Unit: GeV (field), GeV\(^2\) (field strength). Status:
  STANDARD.
\item
  \(x\): The coordinates of the fold: \(x_1, x_2, x_3\) of ordinary
  space, \(x_4\) the time, \(x_5, x_6, x_7\) the extra times; a field of
  the fold is a function of the four observable coordinates
  \(x = (x_1, x_2, x_3, x_4)\), and
  \(\mathrm{d}^4x = \mathrm{d}x_1\,\mathrm{d}x_2\,\mathrm{d}x_3\,\mathrm{d}x_4\).
  Unit: length (each coordinate). Status: AUTHOR.
\item
  \({}'\): The derivative of a function of the time \(x_4\) alone with
  respect to \(x_4\) (\(a_4' = \mathrm{d}a_4/\mathrm{d}x_4\); two primes
  the second derivative); on a function of \(\sigma\) alone (a profile)
  the derivative with respect to \(\sigma\). Status: STANDARD.
\item
  \(H\): The author's constant (an inverse length) that fixes the period
  \(\pi/(6H)\) of \(x_0\), the warp factor
  \(\mathrm{e}^{-2H\lvert\sigma\rvert}\) of the metric, the cosmological
  constant \(\Lambda = -21H^2\) and the tension
  \(\kappa\lambda_{\mathrm{F}} = 12H\) of the fold; in UM-4 a free
  parameter, bounded from below by the propagation of the gauge fields
  in the bulk. Unit: GeV (1/length). Value: \(H > 2534\) GeV, about 2.5
  TeV (from the invisible width of the Z boson). Status: AUTHOR.
\item
  \(\mathrm{sgn}\): The sign function: \(+1\) for a positive argument
  and \(-1\) for a negative one. Status: STANDARD.
\item
  \(\mathrm{e}\): Euler's number, the base of the natural logarithm;
  \(\mathrm{e}^{X}\) is the exponential of \(X\). Unit: dimensionless.
  Value: 2.718281828. Status: STANDARD.
\item
  \(m\): A mass: \(m_i\) the mass of the neutrino mass state \(i\);
  \(\Delta m_{ij}^2 = m_i^2 - m_j^2\) the squared-mass difference of the
  states \(i\) and \(j\) (\(\Delta m_{21}^2\), \(\Delta m_{31}^2\),
  \(\Delta m_{32}^2\)); \(\Delta m^2\) without labels the difference of
  the two states of a resonance, \(m_2^2 - m_1^2\), multiplied by
  \(v_{\mathrm{loc}}^2/v^2\) where this differs from 1; \(m\) without a
  label the mass of a mode of a field of the bulk as the fold sees it,
  \(m^2 = \omega^2 - k^2\); for the W and Z bosons their complex pole
  mass, \(m^2 = m_0^2 - \mathrm{i}m_0\Gamma_{\mathrm{bulk}}\) to the
  leading order; in the relic abundance of the particles \(\chi\) their
  mass. Unit: GeV. Status: STANDARD.
\item
  \(\eta_{ab}\): The flat metric of the frame,
  \(\eta = \mathrm{diag}(1, 1, 1, 1, -1, -1, -1, -1)\) (eta4488 of
  4+4.nb); \(\eta^{ab}\) is its inverse (the same matrix). Unit:
  dimensionless. Status: AUTHOR.
\item
  \(c\): A frame index (as \(a\)). Status: STANDARD.
\item
  \(\partial\): A partial derivative: the derivative with respect to one
  variable at fixed values of the others (\(\partial_\sigma\) with
  respect to \(\sigma\); \(\partial_M\) and \(\partial_\mu\) with
  respect to the coordinate that carries the same index, of the eight
  dimensions or of the fold). Status: STANDARD.
\item
  \(b\): A frame index (as \(a\)). Status: STANDARD.
\item
  \(\infty\): Infinity: as a limit of an integral or a sum it means that
  the range has no bound. Status: STANDARD.
\item
  \(\mathrm{d}\): The differential: \(\mathrm{d}\) in front of a
  quantity is an infinitesimal change of it (as \(\mathrm{d}\sigma\), a
  small step of \(\sigma\)); the quotient of two such changes is a
  derivative; \(\mathrm{d}\) with an exponent, in front of several
  variables, is their volume element, the product of their
  differentials. Status: STANDARD.
\end{itemize}

The modes of a gauge field of the bulk, written for the hypercharge
field \(B_M\) (each gauge field of the bulk, to the first order in the
field, IN-7), in the gauge \(B_\sigma = 0\): the field of the fold
\(\hat B_a(x)\) (\(x\) the coordinates of the fold, \(a, b, c\) the
seven directions of the fold,
\(\hat B_{cb} = \partial_c\hat B_b - \partial_b\hat B_c\)) obeys the
Proca equation of the fold with the mass \(m\), and the profile the
equation in the middle (the factor \(-5H\) from
\(\sqrt{\lvert g\rvert}\,g^{\sigma\sigma}g^{bb} \propto \mathrm{e}^{-5H\lvert\sigma\rvert}\)).
The massless mode has the profile 1 and the weight
\(\sqrt{\lvert g\rvert}\,g^{aa}g^{bb} \propto \mathrm{e}^{-3H\lvert\sigma\rvert}\)
in its kinetic term, whose integral is \(2/(3H)\): a massless gauge
field of the fold, bound to it, with the coupling of seven dimensions
\(\mathsf{g}_1^2\cdot\frac32H\) and the four-dimensional coupling
\(g_1^2 = 3H\mathsf{g}_1^2/(2V_*)\) (FE-44). This is the photon, the
gluons and the massless electroweak gauge fields of the Standard Model
before the Higgs vacuum gives the W and Z bosons their masses (FE-82). A
component \(B_\sigma\) along the hidden direction has no normalisable
zero mode (its profile would be \(\mathrm{e}^{5H\lvert\sigma\rvert}\)),
so the four-dimensional world sees no scalar from it. The massive modes
(\(m > 0\), a continuum without gap) modify Coulomb's law at short
distances (FE-81). (In UM-3 the bulk photon had \(\omega^2 = 2k^2\):
1.41 times the speed of light.) Status: DERIVED (\texttt{um4\_bulk.py},
B2: exactly on both halves; the lowest \(\omega^2\) in the box equals
\(k^2\) to \(4\times10^{-12}\); B10: the continuum of massive modes from
zero mass).

\[
\begin{aligned}
&A_\mu = \mathfrak{e}_\mu\,\mathrm{e}^{\mathrm{i}k(x_3 - x_4)}:\qquad \mathfrak{e}_3 + \mathfrak{e}_4 = 0,\qquad
\mathfrak{e}_\mu \sim \mathfrak{e}_\mu + \beta_{\mathrm{g}}\,k\,(0, 0, 1, -1),\\
&\mathfrak{e}_\pm = \frac{1}{\sqrt2}\big(1, \pm\mathrm{i}, 0, 0\big) \to \mathrm{e}^{\mp\mathrm{i}\theta_{\mathrm{r}}}\,\mathfrak{e}_\pm,\qquad
\vec E\cdot\vec B = 0,\quad \vec E^2 = \vec B^2
\end{aligned} \tag{FE-75}
\] \textbf{Symbols in (FE-75).}

\begin{itemize}
\tightlist
\item
  \(A_\mu\): The photon field of the fold (canonically normalised): the
  massless combination of the massless modes of \(W^1, W^2, W^3\) and
  \(B\) that the Higgs vacuum leaves massless; a field of the bulk whose
  zero mode travels at the speed of light. Unit: GeV. Status: STANDARD.
\item
  \(\mu\): A four-dimensional coordinate index running over the
  observable coordinates \(x_1, x_2, x_3, x_4\) of the fold (signature
  \(+,+,+,-\)); a repeated index is summed. Status: STANDARD.
\item
  \(\mathfrak{e}_\mu,\ \mathfrak{e}_\pm\): The polarisation vector of a
  plane wave of the photon, with components along
  \(x_1, x_2, x_3, x_4\);
  \(\mathfrak{e}_\pm = (1, \pm\mathrm{i}, 0, 0)/\sqrt2\) the circular
  polarisations of a wave travelling along \(x_3\). Unit: dimensionless.
  Status: STANDARD.
\item
  \(\mathrm{e}\): Euler's number, the base of the natural logarithm;
  \(\mathrm{e}^{X}\) is the exponential of \(X\). Unit: dimensionless.
  Value: 2.718281828. Status: STANDARD.
\item
  \(\mathrm{i}\): The imaginary unit, \(\mathrm{i}^2 = -1\). Unit:
  dimensionless. Status: STANDARD.
\item
  \(k\): A wave number: the momentum of a plane wave along \(x_3\),
  \(\mathrm{e}^{\mathrm{i}(kx_3 - \omega x_4)}\); in the cosmological
  perturbations the comoving wave number of a Fourier mode,
  \(k = \lvert\vec k\rvert\). Unit: GeV (1/length). Status: STANDARD.
\item
  \(x_3\): The third coordinate of ordinary 3-space. Unit: length.
  Status: AUTHOR.
\item
  \(x_4\): The time; a prime on a function of \(x_4\) alone is
  \(\mathrm{d}/\mathrm{d}x_4\). Unit: length (= time, \(c = 1\)).
  Status: AUTHOR.
\item
  \(\beta_{\mathrm{g}}\): Any constant: the size of a gauge
  transformation, which adds a multiple of the wave vector to the
  polarisation vector of the photon. Unit: dimensionless. Status:
  STANDARD.
\item
  \(\theta_{\mathrm{r}}\): The angle of a rotation about the direction
  of motion of a wave. Unit: dimensionless (radian). Status: STANDARD.
\item
  \(\vec E\): The electric field of the photon on the fold: its
  component along the direction \(x_i\) of ordinary space is
  \(\partial A_4/\partial x_i - \partial A_i/\partial x_4\)
  (\(i = 1, 2, 3\); \(A_4\) and \(A_i\) the components of \(A_\mu\));
  \(\vec E^2\) its square. Unit: GeV\(^2\). Status: STANDARD.
\item
  \(\vec B\): The magnetic field of the photon on the fold, the curl of
  the components \(A_1, A_2, A_3\) of \(A_\mu\) along ordinary space
  (its component along \(x_1\) is
  \(\partial A_3/\partial x_2 - \partial A_2/\partial x_3\), and
  cyclically); \(\vec B^2\) its square. Unit: GeV\(^2\). Status:
  STANDARD.
\end{itemize}

The photon as a wave: the zero mode of FE-74 for the combination
\(A_\mu\) of \(W^3\) and \(B\) that the Higgs vacuum leaves massless
(FE-45), a plane wave of the wave number \(k\) along \(x_3\) with the
polarisation vector \(\mathfrak{e}_\mu\) (components along
\(x_1, \dots, x_4\); independent of the extra times). The wave satisfies
Maxwell's equations of the fold because its wave vector is null. The
Lorenz condition \(\partial^\mu A_\mu = 0\) requires
\(\mathfrak{e}_3 + \mathfrak{e}_4 = 0\), and adding the gradient of a
gauge function (any constant \(\beta_{\mathrm{g}}\) in the second
relation) removes the remaining component along the direction of motion:
the only physical polarisations are the two transverse ones, along
\(x_1\) and along \(x_2\); the longitudinal and the time-like
polarisations are pure gauge. Its electric field \(\vec E\) and its
magnetic field \(\vec B\) are orthogonal to the direction of motion
\(x_3\) and to each other and equal in size. Under a rotation by the
angle \(\theta_{\mathrm{r}}\) about the direction of motion the circular
polarisations \(\mathfrak{e}_\pm\) are multiplied by
\(\mathrm{e}^{\mp\mathrm{i}\theta_{\mathrm{r}}}\): the photon has the
helicities \(\pm1\); it is a transverse wave of spin 1, as the order of
2026-10-25 13:15 requires, and it travels at the speed of light. A
photon that leaves the fold travels through the bulk at the local speed
of light, which at every depth is that of the fold (FE-83). The massless
vector of seven dimensions has five polarisations; the three that point
along the extra times are the components \(A_r\), which UM-4 excludes
(IN-7: they have kinetic terms of the wrong sign, \texttt{um4\_bulk.py},
B11). Status: DERIVED (\texttt{um4\_bulk.py}, B3 and B11), ASSUMED (no
components along the extra times).

\[
\begin{aligned}
&\frac12\Big[\mathrm{e}^{2H\lvert\sigma\rvert}\Big(\frac{\partial^2h_{\mathrm{gw}}}{\partial x_4^2} - \frac{\partial^2h_{\mathrm{gw}}}{\partial x_3^2}\Big)
 - \frac{\partial^2h_{\mathrm{gw}}}{\partial\sigma^2} + 7H\,\mathrm{sgn}(\sigma)\,\frac{\partial h_{\mathrm{gw}}}{\partial\sigma}\Big] = 0:\qquad
 h_{\mathrm{gw}} = h_{+}(x_3 - x_4),\\
&\mathsf{h}_{ij} = \begin{pmatrix} h_{+} & h_{\times} & 0\\ h_{\times} & -h_{+} & 0\\ 0 & 0 & 0\end{pmatrix},\qquad
 \frac{\Delta L}{L} = \frac12\big(h_{+}\cos2\phi_{\mathrm{r}} + h_{\times}\sin2\phi_{\mathrm{r}}\big),\qquad
 h_{+} \pm \mathrm{i}h_{\times} \to \mathrm{e}^{\mp2\mathrm{i}\theta_{\mathrm{r}}}\big(h_{+} \pm \mathrm{i}h_{\times}\big)
\end{aligned} \tag{FE-76}
\] \textbf{Symbols in (FE-76).}

\begin{itemize}
\tightlist
\item
  \(\mathrm{e}\): Euler's number, the base of the natural logarithm;
  \(\mathrm{e}^{X}\) is the exponential of \(X\). Unit: dimensionless.
  Value: 2.718281828. Status: STANDARD.
\item
  \(H\): The author's constant (an inverse length) that fixes the period
  \(\pi/(6H)\) of \(x_0\), the warp factor
  \(\mathrm{e}^{-2H\lvert\sigma\rvert}\) of the metric, the cosmological
  constant \(\Lambda = -21H^2\) and the tension
  \(\kappa\lambda_{\mathrm{F}} = 12H\) of the fold; in UM-4 a free
  parameter, bounded from below by the propagation of the gauge fields
  in the bulk. Unit: GeV (1/length). Value: \(H > 2534\) GeV, about 2.5
  TeV (from the invisible width of the Z boson). Status: AUTHOR.
\item
  \(\sigma\): The proper coordinate of the hidden direction:
  \(\mathrm{d}\sigma = \lvert\cot z\rvert\,\mathrm{d}x_0\),
  \(\sigma = \mathrm{sgn}(\pi/2 - z)\ln(\sin z)/(6H)\) on one period
  \(0 < z < \pi\); \(\sigma = 0\) is the fold \(z = \pi/2\),
  \(\sigma \to \mp\infty\) the tips \(z \to 0, \pi\); one period of
  \(x_0\) is the whole line \(-\infty < \sigma < \infty\). Unit: length.
  Status: DERIVED.
\item
  \(\partial\): A partial derivative: the derivative with respect to one
  variable at fixed values of the others (\(\partial_\sigma\) with
  respect to \(\sigma\); \(\partial_M\) and \(\partial_\mu\) with
  respect to the coordinate that carries the same index, of the eight
  dimensions or of the fold). Status: STANDARD.
\item
  \(h_{\mathrm{gw}}\): The amplitude of a gravitational wave of the +
  polarisation travelling along \(x_3\): the metric components along
  \(x_1\) and \(x_2\) multiplied by \(\mathrm{e}^{h_{\mathrm{gw}}}\) and
  \(\mathrm{e}^{-h_{\mathrm{gw}}}\); a function of \(\sigma\), \(x_3\)
  and \(x_4\) in the bulk, and on the fold the gravitational wave of the
  four-dimensional world, to the first order its plus polarisation
  \(h_{+}\). Unit: dimensionless. Status: STANDARD.
\item
  \(x_4\): The time; a prime on a function of \(x_4\) alone is
  \(\mathrm{d}/\mathrm{d}x_4\). Unit: length (= time, \(c = 1\)).
  Status: AUTHOR.
\item
  \(x_3\): The third coordinate of ordinary 3-space. Unit: length.
  Status: AUTHOR.
\item
  \(\mathrm{sgn}\): The sign function: \(+1\) for a positive argument
  and \(-1\) for a negative one. Status: STANDARD.
\item
  \(h_{+}\): The plus polarisation of a gravitational wave travelling
  along \(x_3\): \(\mathsf{h}_{11} = -\mathsf{h}_{22} = h_{+}\), a
  function of \(x_3 - x_4\); to the first order the zero mode of
  \(h_{\mathrm{gw}}\). Unit: dimensionless. Status: STANDARD.
\item
  \(\mathsf{h}_{ij}\): The transverse traceless perturbation of the
  metric of ordinary space by a gravitational wave as the
  four-dimensional world sees it:
  \(\hat h_{ij} = \mathrm{e}^{2a_4}(\delta_{ij} + \mathsf{h}_{ij})\)
  (\(i, j = 1, 2, 3\)). Unit: dimensionless. Status: STANDARD.
\item
  \(i\): A generation index (\(i = 1, 2, 3\)), which also labels the
  three neutrino mass states; in a frame component such as
  \(\mathcal{S}^{(i)}{}_{(i)}\) one of the three directions of ordinary
  space (the three components are equal), and in a tensor of ordinary
  space such as \(\mathsf{h}_{ij}\) one of the directions
  \(x_1, x_2, x_3\); a repeated index is not summed unless a sum sign is
  written. Status: STANDARD.
\item
  \(j\): A second generation or mass-state index, or a second direction
  of ordinary space (as \(i\)). Status: STANDARD.
\item
  \(h_{\times}\): The cross polarisation of a gravitational wave
  travelling along \(x_3\):
  \(\mathsf{h}_{12} = \mathsf{h}_{21} = h_{\times}\), a function of
  \(x_3 - x_4\). Unit: dimensionless. Status: STANDARD.
\item
  \(\Delta\): A change, difference or small variation of the quantity
  that follows (for example an interval of the time, a perturbation of a
  density, a difference of two squared masses, or an addition to an
  effective number of species). Status: STANDARD.
\item
  \(L\): A distance between two points of the fold along ordinary space:
  for a gravitational wave the radius of a ring of free test masses
  (\(\Delta L\) the change of the distance of one of them from the
  centre), for a signal the distance between its emitter and its
  receiver. Unit: length. Status: STANDARD.
\item
  \(\phi_{\mathrm{r}}\): The angle at which a test mass sits on a ring
  of test masses in the plane of \(x_1\) and \(x_2\), measured from the
  \(x_1\) axis. Unit: dimensionless (radian). Status: STANDARD.
\item
  \(\mathrm{i}\): The imaginary unit, \(\mathrm{i}^2 = -1\). Unit:
  dimensionless. Status: STANDARD.
\item
  \(\theta_{\mathrm{r}}\): The angle of a rotation about the direction
  of motion of a wave. Unit: dimensionless (radian). Status: STANDARD.
\end{itemize}

The gravitational wave. First line: the + polarisation
\(h_{\mathrm{gw}}\) of a wave along \(x_3\) in the bulk of the static
block (the metric components along \(x_1\) and \(x_2\) multiplied by
\(\mathrm{e}^{h_{\mathrm{gw}}}\) and \(\mathrm{e}^{-h_{\mathrm{gw}}}\)),
to the first order: Einstein's equation FE-1 gives the equation of the
massless scalar FE-73, and the zero mode, independent of \(\sigma\), is
any function of \(x_3 - x_4\) (written \(h_{+}\)): it solves the bulk
equation on both sides and the junction condition at the fold (its
\(\sigma\)-derivative does not jump), so the graviton of the bulk has a
massless mode bound to the fold that travels at the speed of light at
every wavelength, with the coupling \(\kappa_7 = 5\kappa H/2\) of seven
dimensions (FE-63, FE-71). Second line: the wave as the four-dimensional
world sees it, the transverse traceless perturbation \(\mathsf{h}_{ij}\)
of the metric of ordinary space (\(i, j = 1, 2, 3\):
\(\hat h_{ij} = \mathrm{e}^{2a_4}(\delta_{ij} + \mathsf{h}_{ij})\)) of a
wave along \(x_3\). It has the two polarisations \(h_{+}\) (the ``plus''
mode, \(\mathsf{h}_{11} = -\mathsf{h}_{22}\)) and \(h_{\times}\) (the
``cross'' mode, \(\mathsf{h}_{12} = \mathsf{h}_{21}\)). A ring of free
test masses in the plane of \(x_1\) and \(x_2\), of the radius \(L\),
changes the distance \(L\) of the mass at the angle
\(\phi_{\mathrm{r}}\) (from the \(x_1\) axis) by \(\Delta L\): the plus
mode stretches the ring along \(x_1\) while it squeezes it along \(x_2\)
and then the reverse, the cross mode does the same along the two
diagonals, rotated by \(45^\circ\): they stretch and squeeze space
perpendicularly to each other and to the direction of motion. Under a
rotation by the angle \(\theta_{\mathrm{r}}\) about the direction of
motion \(h_{+}\) and \(h_{\times}\) rotate into each other by the angle
\(2\theta_{\mathrm{r}}\), so that the circular combinations
\(h_{+} \pm \mathrm{i}h_{\times}\) are multiplied by
\(\mathrm{e}^{\mp2\mathrm{i}\theta_{\mathrm{r}}}\): helicity \(\pm2\), a
wave of spin 2, as it must be for a wave sourced by the rank-2
energy-momentum tensor (FE-63: its source is the transverse traceless
stress \(\Pi_{\mathrm{gw}}\)). The massless graviton of seven dimensions
has \(7\cdot(7 - 3)/2 = 14\) polarisations (20 transverse traceless
components less 6 residual gauge modes), which the four-dimensional
world sees as the graviton (2), three vectors \(h_{\mu r}\) (6) and six
scalars \(h_{rs}\) (6). In the Einstein-Hilbert action to the second
order on the fold of signature (3, 4) the graviton and the six scalars
have kinetic terms of the right sign, but the three vectors
\(h_{\mu r}\), with one index along ordinary space or the time and one
along an extra time, have the wrong sign: they are ghosts, like the
components \(A_r\) of a gauge field that UM-4 excludes. Matter of the
fold that does not depend on the extra times and is the same along the
three of them has no stress \(\mathcal{T}_{\mu r}\) and only the trace
of \(\mathcal{T}_{rs}\), so at the linear order it emits only \(h_{+}\),
\(h_{\times}\) and the trace of \(h_{rs}\), the volume of the extra
times, which the fold fluid holds fixed (IN-10, FE-65), and no ghost
vectors: the gravitational waves that the matter of the fold emits have
only the plus and the cross polarisations, as the order requires and as
observed (no other polarisation has been detected). This holds for the
massless mode; the massive modes of the bulk carry other helicities, but
their emission is negligible at the energies of the fold, and the ghost
vectors remain a defect of the bulk gravity of signature (3, 4) at
higher orders (OPEN). Status: DERIVED (\texttt{um4\_bulk.py}, B4, B5 and
B11: the signs of the kinetic terms of all 14 polarisations from the
action; \texttt{um4\_foldsources.py}, FS2), ASSUMED (that the fold fluid
holds the volume of the extra times fixed for waves as well, a stiff
modulus), OPEN (the ghost vectors beyond the linear order).

\[
\begin{aligned}
&\Psi_{\mathrm{k}} = \mathrm{e}^{\frac72H\lvert\sigma\rvert}\,\mathrm{e}^{-M_{\mathrm{k}}\lvert\sigma\rvert}\,\chi_0(x),\qquad
\gamma^{(0)}\chi_0 = -\chi_0,\qquad \gamma^{(a)}\,\partial_a\chi_0 = 0\quad(\text{along the fold}),\\
&\sqrt{\lvert g\rvert}\,\bar\Psi_{\mathrm{k}}\gamma^{x_4}\Psi_{\mathrm{k}} \propto \mathrm{e}^{(H - 2M_{\mathrm{k}})\lvert\sigma\rvert}
\end{aligned} \tag{FE-77}
\] \textbf{Symbols in (FE-77).}

\begin{itemize}
\tightlist
\item
  \(\Psi_{\mathrm{k}},\ \bar\Psi_{\mathrm{k}}\): A fermion field of the
  bulk with the kink mass \(M_{\mathrm{k}}\,\mathrm{sgn}(\sigma)\) (an
  illustration, not a field of UM-4). Unit: GeV\(^{7/2}\). Status:
  ASSUMED.
\item
  \(\mathrm{e}\): Euler's number, the base of the natural logarithm;
  \(\mathrm{e}^{X}\) is the exponential of \(X\). Unit: dimensionless.
  Value: 2.718281828. Status: STANDARD.
\item
  \(H\): The author's constant (an inverse length) that fixes the period
  \(\pi/(6H)\) of \(x_0\), the warp factor
  \(\mathrm{e}^{-2H\lvert\sigma\rvert}\) of the metric, the cosmological
  constant \(\Lambda = -21H^2\) and the tension
  \(\kappa\lambda_{\mathrm{F}} = 12H\) of the fold; in UM-4 a free
  parameter, bounded from below by the propagation of the gauge fields
  in the bulk. Unit: GeV (1/length). Value: \(H > 2534\) GeV, about 2.5
  TeV (from the invisible width of the Z boson). Status: AUTHOR.
\item
  \(\sigma\): The proper coordinate of the hidden direction:
  \(\mathrm{d}\sigma = \lvert\cot z\rvert\,\mathrm{d}x_0\),
  \(\sigma = \mathrm{sgn}(\pi/2 - z)\ln(\sin z)/(6H)\) on one period
  \(0 < z < \pi\); \(\sigma = 0\) is the fold \(z = \pi/2\),
  \(\sigma \to \mp\infty\) the tips \(z \to 0, \pi\); one period of
  \(x_0\) is the whole line \(-\infty < \sigma < \infty\). Unit: length.
  Status: DERIVED.
\item
  \(M_{\mathrm{k}}\): The mass of a fermion of the bulk that changes
  sign at the fold (a kink mass
  \(M_{\mathrm{k}}\,\mathrm{sgn}(\sigma)\)), used to illustrate the
  speeds of bulk fermions. Unit: GeV. Status: ASSUMED.
\item
  \(\chi_0\): The spinor of the fold of the massless mode of a fermion
  of the bulk: a solution of the massless Dirac equation of the fold of
  one chirality along \(\sigma\). Unit: GeV\(^{7/2}\). Status: DERIVED.
\item
  \(x\): The coordinates of the fold: \(x_1, x_2, x_3\) of ordinary
  space, \(x_4\) the time, \(x_5, x_6, x_7\) the extra times; a field of
  the fold is a function of the four observable coordinates
  \(x = (x_1, x_2, x_3, x_4)\), and
  \(\mathrm{d}^4x = \mathrm{d}x_1\,\mathrm{d}x_2\,\mathrm{d}x_3\,\mathrm{d}x_4\).
  Unit: length (each coordinate). Status: AUTHOR.
\item
  \(\gamma^M,\ \gamma^{(a)}\): The Dirac matrices of dirac16complex
  along the coordinates, \(\gamma^M = \mathbf{e}_a{}^MT^a\),
  \(\gamma_M = g_{MN}\gamma^N\); \(\gamma^{(a)} = T^a\) along the frame
  direction \((a)\) (\(\gamma^{(0)}\) along \(\sigma\), \(\gamma^{(4)}\)
  along the time). Unit: dimensionless. Status: AUTHOR.
\item
  \(a\): A frame (orthonormal) index running over \(0, 1, \dots, 7\)
  (frame direction \((0)\) along \(\sigma\), \((1), (2), (3)\) along
  ordinary space, \((4)\) along the time, \((5), (6), (7)\) along the
  extra times); raised and lowered with the frame metric
  \(\eta_{ab} = \mathrm{diag}(1, 1, 1, 1, -1, -1, -1, -1)\); on the fold
  (\(a, b \in \{1, \dots, 7\}\)) the frame of the fold; in the tensors
  of the fold written with its coordinates \(x^a\) (its metric
  \(\hat g_{ab}\) and the tensors built from it) a coordinate index of
  the fold, \(a \in \{1, \dots, 7\}\); a repeated index is summed.
  Status: STANDARD.
\item
  \(\partial\): A partial derivative: the derivative with respect to one
  variable at fixed values of the others (\(\partial_\sigma\) with
  respect to \(\sigma\); \(\partial_M\) and \(\partial_\mu\) with
  respect to the coordinate that carries the same index, of the eight
  dimensions or of the fold). Status: STANDARD.
\item
  \(g_{MN}\): The metric tensor of the eight-dimensional spacetime
  (signature \(+,+,+,+,-,-,-,-\) in the order \(x_0, \dots, x_7\)):
  \(g_{MN}\,\mathrm{d}X^M\mathrm{d}X^N\) is the squared interval,
  \(g^{MN}\) its inverse, \(\lvert g\rvert\) the absolute value of its
  determinant (\(\sqrt{\lvert g\rvert}\,\mathrm{d}^8X\) is the proper
  volume element). Unit: dimensionless. Status: AUTHOR.
\item
  \(x_4\): The time; a prime on a function of \(x_4\) alone is
  \(\mathrm{d}/\mathrm{d}x_4\). Unit: length (= time, \(c = 1\)).
  Status: AUTHOR.
\end{itemize}

A fermion of the bulk with a mass that changes sign at the fold,
\(M_{\mathrm{k}}\,\mathrm{sgn}(\sigma)\) (an illustration: UM-4 has no
such field; it shows that the quarks and leptons could live in the bulk
without outrunning light). The factor
\(\mathrm{e}^{(7/2)H\lvert\sigma\rvert}\) removes the spin connection
(FE-14), and the zero mode is the profile
\(\mathrm{e}^{-M_{\mathrm{k}}\lvert\sigma\rvert}\) times a spinor
\(\chi_0(x)\) of the fold of one chirality along \(\sigma\) that obeys
the massless Dirac equation of the fold. Its charge density along
\(\sigma\) falls as
\(\mathrm{e}^{(H - 2M_{\mathrm{k}})\lvert\sigma\rvert}\), so it is
normalisable for \(M_{\mathrm{k}} > H/2\): a massless fermion of seven
dimensions, bound to the fold, at the speed of light (in UM-3 it moved
at \(2M_{\mathrm{k}}/(2M_{\mathrm{k}} - H)\) times the speed of light).
Status: DERIVED (\texttt{um4\_bulk.py}, B6).

\[
\mathsf{S} = \frac12\,\gamma^{(3)}\gamma^{(4)}:\qquad
\big[\mathsf{S}, \gamma^{(0)}\big] = \big[\mathsf{S}, \Gamma_8\big] = 0,\qquad
\big[\mathsf{S}, \gamma^{(3)}\big] = -\gamma^{(4)},\qquad \big[\mathsf{S}, \gamma^{(4)}\big] = -\gamma^{(3)}\qquad\Rightarrow\qquad
\omega^2 = \omega_0^2 + k^2 \tag{FE-78}
\] \textbf{Symbols in (FE-78).}

\begin{itemize}
\tightlist
\item
  \(\mathsf{S}\): The generator of the boosts along \(x_3\) on
  dirac16complex, \(\mathsf{S} = \frac12\gamma^{(3)}\gamma^{(4)}\).
  Unit: dimensionless. Status: DERIVED.
\item
  \(\gamma^M,\ \gamma^{(a)}\): The Dirac matrices of dirac16complex
  along the coordinates, \(\gamma^M = \mathbf{e}_a{}^MT^a\),
  \(\gamma_M = g_{MN}\gamma^N\); \(\gamma^{(a)} = T^a\) along the frame
  direction \((a)\) (\(\gamma^{(0)}\) along \(\sigma\), \(\gamma^{(4)}\)
  along the time). Unit: dimensionless. Status: AUTHOR.
\item
  \(\Gamma_8\): The chirality matrix of eight dimensions,
  \(\Gamma_8 = T_0T_1\cdots T_7\) (the author's T16\^{}A{[}8{]})
  \(= \mathrm{diag}(-\mathbf{1}_8, \mathbf{1}_8)\);
  \(\Gamma_8^2 = \mathbf{1}\); it anticommutes with every \(T_a\). Unit:
  dimensionless. Status: AUTHOR.
\item
  \(\omega\): A frequency: the energy of a mode
  \(\propto \mathrm{e}^{-\mathrm{i}\omega x_4}\) (\(\omega_o\) of the
  orbital \(o\)). Unit: GeV. Status: STANDARD.
\item
  \(\omega_0\): The energy at rest (zero momentum along ordinary space
  and along the extra times) of a state of dirac16complex at the fold;
  complex for a resonance,
  \(\omega_0 = \omega_{\mathrm{r}} - \frac{\mathrm{i}}{2}\Gamma_\chi\).
  Unit: GeV. Status: DERIVED.
\item
  \(k\): A wave number: the momentum of a plane wave along \(x_3\),
  \(\mathrm{e}^{\mathrm{i}(kx_3 - \omega x_4)}\); in the cosmological
  perturbations the comoving wave number of a Fourier mode,
  \(k = \lvert\vec k\rvert\). Unit: GeV (1/length). Status: STANDARD.
\end{itemize}

dirac16complex is invariant under the boosts of the fold. \(\mathsf{S}\)
is the generator of the boosts along \(x_3\) on dirac16complex, built
from the author's matrices (\(\gamma^{(a)} = T^a\),
\((\gamma^{(3)})^2 = \mathbf{1}\), \((\gamma^{(4)})^2 = -\mathbf{1}\)).
It commutes with \(\gamma^{(0)}\), with the chirality \(\Gamma_8\) and
with the mass term, and it rotates \(\gamma^{(3)}\) and \(\gamma^{(4)}\)
into each other hyperbolically; in the bulk equation (FE-36) the warp
multiplies the frequency and the momentum by the same factor
\(\mathrm{e}^{H\lvert\sigma\rvert}\), so the boost
\(\mathrm{e}^{\beta_{\mathrm{b}}\mathsf{S}}\) of the rapidity
\(\beta_{\mathrm{b}}\) maps a solution of frequency \(\omega\) and
momentum \(k\) along \(x_3\) onto one with
\((\omega\cosh\beta_{\mathrm{b}} - k\sinh\beta_{\mathrm{b}},\ k\cosh\beta_{\mathrm{b}} - \omega\sinh\beta_{\mathrm{b}})\).
The fold terms of the vacuum (the Higgs portal and the gluon condensate,
FE-39) are scalars of the fold and products of \(\gamma^{(0)}\) and the
projectors \(P_{\mathrm{L}}, P_{\mathrm{R}}\), so they are invariant
too. Every state of dirac16complex at the fold with the energy
\(\omega_0\) at rest therefore has \(\omega^2 = \omega_0^2 + k^2\)
exactly and moves slower than light (in UM-3 the time was not warped,
the boosts were not a symmetry, and the bound states had the speed
coefficient \(c_{\mathrm{b}}^2 = 1.19\) along ordinary space). Status:
DERIVED (\texttt{um4\_bulk.py}, B7, with the author's matrices;
\texttt{um4\_krein.py}, KR6: the moving states have
\(\sqrt{\omega_0^2 + k^2}\) to \(10^{-10}\)).

\hypertarget{the-propagators-from-one-point-of-the-fold-to-another}{%
\subsection{The propagators from one point of the fold to
another}\label{the-propagators-from-one-point-of-the-fold-to-another}}

\[
-\mathrm{e}^{nH\lvert\sigma\rvert}\,\partial_\sigma\big(\mathrm{e}^{-nH\lvert\sigma\rvert}\,\partial_\sigma u\big)
 + \mathrm{e}^{2H\lvert\sigma\rvert}\,\mathsf{p}^2\,u = \delta_{\mathrm{D}}(\sigma),\qquad
\mathbb{G}_n(\mathsf{p}) = u(0) = \frac{\mathrm{K}_{n/2}(\mathsf{p}/H)}{2\mathsf{p}\,\mathrm{K}_{n/2-1}(\mathsf{p}/H)} \tag{FE-79}
\] \textbf{Symbols in (FE-79).}

\begin{itemize}
\tightlist
\item
  \(\mathrm{e}\): Euler's number, the base of the natural logarithm;
  \(\mathrm{e}^{X}\) is the exponential of \(X\). Unit: dimensionless.
  Value: 2.718281828. Status: STANDARD.
\item
  \(n\): The label of a field of the bulk in its profile equation: the
  exponent of the weight \(\mathrm{e}^{-nH\lvert\sigma\rvert}\) of its
  \(\sigma\)-derivatives, \(n = 7\) for a scalar and for each
  polarisation of the graviton, \(n = 5\) for a gauge field. Status:
  STANDARD.
\item
  \(H\): The author's constant (an inverse length) that fixes the period
  \(\pi/(6H)\) of \(x_0\), the warp factor
  \(\mathrm{e}^{-2H\lvert\sigma\rvert}\) of the metric, the cosmological
  constant \(\Lambda = -21H^2\) and the tension
  \(\kappa\lambda_{\mathrm{F}} = 12H\) of the fold; in UM-4 a free
  parameter, bounded from below by the propagation of the gauge fields
  in the bulk. Unit: GeV (1/length). Value: \(H > 2534\) GeV, about 2.5
  TeV (from the invisible width of the Z boson). Status: AUTHOR.
\item
  \(\sigma\): The proper coordinate of the hidden direction:
  \(\mathrm{d}\sigma = \lvert\cot z\rvert\,\mathrm{d}x_0\),
  \(\sigma = \mathrm{sgn}(\pi/2 - z)\ln(\sin z)/(6H)\) on one period
  \(0 < z < \pi\); \(\sigma = 0\) is the fold \(z = \pi/2\),
  \(\sigma \to \mp\infty\) the tips \(z \to 0, \pi\); one period of
  \(x_0\) is the whole line \(-\infty < \sigma < \infty\). Unit: length.
  Status: DERIVED.
\item
  \(\partial\): A partial derivative: the derivative with respect to one
  variable at fixed values of the others (\(\partial_\sigma\) with
  respect to \(\sigma\); \(\partial_M\) and \(\partial_\mu\) with
  respect to the coordinate that carries the same index, of the eight
  dimensions or of the fold). Status: STANDARD.
\item
  \(\sigma\): The label \(\sigma\) in \(\partial_\sigma\): the
  derivative with respect to the proper hidden coordinate \(\sigma\).
  Status: STANDARD.
\item
  \(u(\sigma),\ u_0\): The profile along the hidden direction of a mode
  of a field of the bulk (a scalar, a gauge field or a polarisation of
  the graviton): the field is \(u(\sigma)\) times a field of the fold;
  \(u_0 = 1\) is the profile of the massless mode (the zero mode); in a
  propagator the response of the field to a unit source on the fold.
  Unit: dimensionless (a relative amplitude); length in a propagator.
  Status: DERIVED.
\item
  \(\mathsf{p}\): The Euclidean momentum along the fold of a source on
  the fold: \(\mathsf{p}^2\) the square of its seven-dimensional
  momentum (for a static source spread uniformly over the extra times
  the square of its momentum along ordinary space);
  \(\mathsf{p} = -\mathrm{i}m\) for a time-like momentum of the mass
  \(m\) with outgoing waves into the bulk. Unit: GeV. Status: STANDARD.
\item
  \(\delta_{\mathrm{D}}\): The Dirac delta function of \(\sigma\): zero
  for \(\sigma \ne 0\), with the integral 1 over any range that contains
  \(\sigma = 0\); the integral of \(\delta_{\mathrm{D}}(\sigma)\) times
  a smooth function of \(\sigma\) is the value of that function at
  \(\sigma = 0\). Unit: 1/length (for \(\delta_{\mathrm{D}}(\sigma)\)).
  Status: STANDARD.
\item
  \(\mathbb{G}_n,\ \mathbb{G}_7,\ \mathbb{G}_5\): The propagator of a
  massless field of the bulk from one point of the fold to another, in
  the momentum space of the fold: the value at the fold of its response
  to a unit source on the fold; \(\mathbb{G}_7\) for each polarisation
  of the graviton and for a scalar, \(\mathbb{G}_5\) for a gauge field.
  Unit: GeV\(^{-1}\). Status: DERIVED.
\item
  \(\mathrm{K}\): The modified Bessel function of the second kind (of
  the order written as its subscript): the solution of the modified
  Bessel equation that decays exponentially for a large argument; for a
  half-integer order it is an exponential times a polynomial in the
  inverse of the argument, which makes the propagators of the bulk
  rational functions. Unit: dimensionless. Status: STANDARD.
\end{itemize}

The response of a massless field of the bulk to a unit source on the
fold: the profile \(u(\sigma)\) of the field for the Euclidean momentum
\(\mathsf{p}\) along the fold (\(\mathsf{p}^2\) the square of the
seven-dimensional momentum; for a static source spread uniformly over
the extra times, the momentum along ordinary space), with the label
\(n = 7\) for the graviton (each polarisation) and a scalar, and
\(n = 5\) for a gauge field (FE-73, FE-74). With
\(\zeta_{\mathrm{w}} = \mathrm{e}^{H\lvert\sigma\rvert}/H\) the solution
that decays into the bulk is
\(\zeta_{\mathrm{w}}^{n/2}\mathrm{K}_{n/2}(\mathsf{p}\zeta_{\mathrm{w}})\),
and the jump of \(u'\) at the fold fixes its normalisation:
\(\mathbb{G}_n(\mathsf{p})\) is the propagator of the bulk from one
point of the fold to another, in the momentum space of the fold
(\(\mathrm{K}_\nu\) the modified Bessel function of the second kind).
For small momenta \(\mathbb{G}_n \to (n/2 - 1)H/\mathsf{p}^2\): the
massless propagator of the fold with the strength \(\frac52H\) for the
graviton (the coupling \(\kappa_7 = \frac52\kappa H\)) and \(\frac32H\)
for a gauge field (\(g_1^2 = \frac32H\mathsf{g}_1^2/V_*\)). Status:
DERIVED (\texttt{um4\_propagators.py}, P1: the solution and the jump
checked with sympy).

\[
\begin{aligned}
&\mathbb{G}_7 = \frac{H}{2\mathsf{p}^2}\,\frac{x_{\mathrm{p}}^3 + 6x_{\mathrm{p}}^2 + 15x_{\mathrm{p}} + 15}{x_{\mathrm{p}}^2 + 3x_{\mathrm{p}} + 3}
 = \frac{5H}{2\mathsf{p}^2}\Big(1 + \frac{x_{\mathrm{p}}^2}{15} - \frac{x_{\mathrm{p}}^4}{45} + \frac{x_{\mathrm{p}}^5}{45} + \dots\Big),\\
&\mathbb{G}_5 = \frac{H}{2\mathsf{p}^2}\,\frac{x_{\mathrm{p}}^2 + 3x_{\mathrm{p}} + 3}{x_{\mathrm{p}} + 1}
 = \frac{3H}{2\mathsf{p}^2}\Big(1 + \frac{x_{\mathrm{p}}^2}{3} - \frac{x_{\mathrm{p}}^3}{3} + \dots\Big),\qquad x_{\mathrm{p}} = \frac{\mathsf{p}}{H}
\end{aligned} \tag{FE-80}
\] \textbf{Symbols in (FE-80).}

\begin{itemize}
\tightlist
\item
  \(\mathbb{G}_n,\ \mathbb{G}_7,\ \mathbb{G}_5\): The propagator of a
  massless field of the bulk from one point of the fold to another, in
  the momentum space of the fold: the value at the fold of its response
  to a unit source on the fold; \(\mathbb{G}_7\) for each polarisation
  of the graviton and for a scalar, \(\mathbb{G}_5\) for a gauge field.
  Unit: GeV\(^{-1}\). Status: DERIVED.
\item
  \(H\): The author's constant (an inverse length) that fixes the period
  \(\pi/(6H)\) of \(x_0\), the warp factor
  \(\mathrm{e}^{-2H\lvert\sigma\rvert}\) of the metric, the cosmological
  constant \(\Lambda = -21H^2\) and the tension
  \(\kappa\lambda_{\mathrm{F}} = 12H\) of the fold; in UM-4 a free
  parameter, bounded from below by the propagation of the gauge fields
  in the bulk. Unit: GeV (1/length). Value: \(H > 2534\) GeV, about 2.5
  TeV (from the invisible width of the Z boson). Status: AUTHOR.
\item
  \(\mathsf{p}\): The Euclidean momentum along the fold of a source on
  the fold: \(\mathsf{p}^2\) the square of its seven-dimensional
  momentum (for a static source spread uniformly over the extra times
  the square of its momentum along ordinary space);
  \(\mathsf{p} = -\mathrm{i}m\) for a time-like momentum of the mass
  \(m\) with outgoing waves into the bulk. Unit: GeV. Status: STANDARD.
\item
  \(x_{\mathrm{p}}\): The Euclidean momentum along the fold in units of
  \(H\), \(x_{\mathrm{p}} = \mathsf{p}/H\). Unit: dimensionless. Status:
  DERIVED.
\end{itemize}

For the half-integer orders of the Bessel functions the propagators are
rational functions of \(x_{\mathrm{p}} = \mathsf{p}/H\). The expansions
for momenta small compared with \(H\) show what the massive modes add to
the zero mode: the even powers of \(x_{\mathrm{p}}\) are polynomials in
\(\mathsf{p}^2\), contact interactions between the sources on the fold
(local terms, which act only where the sources overlap); the first odd
power (\(x_{\mathrm{p}}^5\) for gravity, \(x_{\mathrm{p}}^3\) for a
gauge field) is not analytic in \(\mathsf{p}^2\) and gives a correction
of the potential of long range (FE-81). Status: DERIVED
(\texttt{um4\_propagators.py}, P2).

\[
\mathcal{V}_7(r) = \frac{5H}{8\pi r}\Big[1 + \frac{16}{15\pi}\,\frac{1}{(Hr)^5} + \dots\Big],\qquad
\mathcal{V}_5(r) = \frac{3H}{8\pi r}\Big[1 + \frac{4}{3\pi}\,\frac{1}{(Hr)^3} + \dots\Big] \tag{FE-81}
\] \textbf{Symbols in (FE-81).}

\begin{itemize}
\tightlist
\item
  \(\mathcal{V}_7(r),\ \mathcal{V}_5(r)\): The static potentials between
  two unit point sources on the fold at the distance \(r\) (spread
  uniformly over the extra times) from the exchange of the graviton
  (\(\mathcal{V}_7\)) and of a gauge field (\(\mathcal{V}_5\)) of the
  bulk: the three-dimensional Fourier transforms of \(\mathbb{G}_7\) and
  \(\mathbb{G}_5\). Unit: GeV\(^2\). Status: DERIVED.
\item
  \(r\): A distance from a centre: the distance from the centre of the
  star along the path of a neutrino; in Newton's and Coulomb's laws the
  distance between two point sources on the fold. Unit: length. Status:
  STANDARD.
\item
  \(H\): The author's constant (an inverse length) that fixes the period
  \(\pi/(6H)\) of \(x_0\), the warp factor
  \(\mathrm{e}^{-2H\lvert\sigma\rvert}\) of the metric, the cosmological
  constant \(\Lambda = -21H^2\) and the tension
  \(\kappa\lambda_{\mathrm{F}} = 12H\) of the fold; in UM-4 a free
  parameter, bounded from below by the propagation of the gauge fields
  in the bulk. Unit: GeV (1/length). Value: \(H > 2534\) GeV, about 2.5
  TeV (from the invisible width of the Z boson). Status: AUTHOR.
\item
  \(\pi\): The number pi, the ratio of the circumference of a circle to
  its diameter. Unit: dimensionless. Value: 3.141592654. Status:
  STANDARD.
\end{itemize}

The static potentials between two unit point sources on the fold at the
distance \(r\) (spread uniformly over the extra times): the
three-dimensional Fourier transforms of FE-80. The first terms are the
potentials of the zero modes, \(1/(4\pi r)\) times the strengths
\(\frac52H\) and \(\frac32H\); multiplied by the couplings they are
Newton's law (with \(8\pi G_{\mathrm{N}} = \kappa_7/V_*\), FE-59) and
Coulomb's law. The massive modes of the bulk add to Newton's law the
relative correction \(0.3395/(Hr)^5\) and to Coulomb's law
\(0.4244/(Hr)^3\) (\(16/(15\pi)\) and \(4/(3\pi)\)). For gravity this is
the scalar part of the exchange; the tensor structure of the massive
modes of spin 2 changes its coefficient by a factor of order one (as in
the brane of Randall and Sundrum, where the correction is \(1/(kr)^2\)
in five dimensions). Newton's law was tested by torsion balances at
separations from 52 micrometres to 3 mm (J. G. Lee, E. G. Adelberger, T.
S. Cook, S. M. Fleischer and B. R. Heckel, Phys. Rev.~Lett. 124 (2020)
101101, which excludes Yukawa forces of gravitational strength with a
range above 38.6 micrometres); taken to hold to one per cent at 52
micrometres, it requires \(H > 7.7\times10^{-3}\) eV. For the bound
\(H > 2534\) GeV of FE-82 the corrections are \(2.6\times10^{-75}\) at
52 micrometres (\(2.7\times10^{-75}\) for \(H = 2.5\) TeV) and
\(1.4\times10^{-27}\) for Coulomb's law at the Bohr radius:
unobservable. Status: DERIVED (\texttt{um4\_propagators.py}, P3: the
asymptotic coefficients exactly, and the numerical Fourier transforms at
\(Hr = 32\) within 1.2 per cent of them; P5: the corrections at the
bound), ESTIMATE (the factor of the tensor structure; the one per cent
at 52 micrometres; the bounds).

\[
\Big(1 + M_{\mathrm{loc}}^2\,\mathbb{G}_5(\mathsf{p})\Big)\Big\vert_{\mathsf{p} = -\mathrm{i}m} = 0:\qquad
m^2 = m_0^2\Big(1 - \frac{m_0^2}{3H^2}\Big) - \mathrm{i}\,\frac{m_0^5}{3H^3} + \dots,\qquad
m_0^2 = \frac32\,H M_{\mathrm{loc}}^2,\qquad \Gamma_{\mathrm{bulk}} = \frac{m_0^4}{3H^3} \tag{FE-82}
\] \textbf{Symbols in (FE-82).}

\begin{itemize}
\tightlist
\item
  \(M_{\mathrm{loc}}\): The mass term that the Higgs vacuum on the fold
  gives a gauge field of the bulk:
  \(M_{\mathrm{loc}}^2\delta_{\mathrm{D}}(\sigma)\) in its profile
  equation; for the W,
  \(M_{\mathrm{loc}}^2 = \mathsf{g}_2^2(v^2 + \mathcal{A}^2)/(4V_*)\).
  Unit: GeV\(^{1/2}\) (\(M_{\mathrm{loc}}^2\) in GeV). Status: DERIVED.
\item
  \(\mathbb{G}_n,\ \mathbb{G}_7,\ \mathbb{G}_5\): The propagator of a
  massless field of the bulk from one point of the fold to another, in
  the momentum space of the fold: the value at the fold of its response
  to a unit source on the fold; \(\mathbb{G}_7\) for each polarisation
  of the graviton and for a scalar, \(\mathbb{G}_5\) for a gauge field.
  Unit: GeV\(^{-1}\). Status: DERIVED.
\item
  \(\mathsf{p}\): The Euclidean momentum along the fold of a source on
  the fold: \(\mathsf{p}^2\) the square of its seven-dimensional
  momentum (for a static source spread uniformly over the extra times
  the square of its momentum along ordinary space);
  \(\mathsf{p} = -\mathrm{i}m\) for a time-like momentum of the mass
  \(m\) with outgoing waves into the bulk. Unit: GeV. Status: STANDARD.
\item
  \(\mathrm{i}\): The imaginary unit, \(\mathrm{i}^2 = -1\). Unit:
  dimensionless. Status: STANDARD.
\item
  \(m\): A mass: \(m_i\) the mass of the neutrino mass state \(i\);
  \(\Delta m_{ij}^2 = m_i^2 - m_j^2\) the squared-mass difference of the
  states \(i\) and \(j\) (\(\Delta m_{21}^2\), \(\Delta m_{31}^2\),
  \(\Delta m_{32}^2\)); \(\Delta m^2\) without labels the difference of
  the two states of a resonance, \(m_2^2 - m_1^2\), multiplied by
  \(v_{\mathrm{loc}}^2/v^2\) where this differs from 1; \(m\) without a
  label the mass of a mode of a field of the bulk as the fold sees it,
  \(m^2 = \omega^2 - k^2\); for the W and Z bosons their complex pole
  mass, \(m^2 = m_0^2 - \mathrm{i}m_0\Gamma_{\mathrm{bulk}}\) to the
  leading order; in the relic abundance of the particles \(\chi\) their
  mass. Unit: GeV. Status: STANDARD.
\item
  \(m_0\): The mass of the W or the Z boson without the massive modes of
  the bulk (the mass of the massless mode with the Higgs vacuum),
  \(m_0^2 = \frac32HM_{\mathrm{loc}}^2\);
  \(m_0^2 = g_2^2(v^2 + \mathcal{A}^2)/4\) for the W. Unit: GeV. Status:
  DERIVED.
\item
  \(H\): The author's constant (an inverse length) that fixes the period
  \(\pi/(6H)\) of \(x_0\), the warp factor
  \(\mathrm{e}^{-2H\lvert\sigma\rvert}\) of the metric, the cosmological
  constant \(\Lambda = -21H^2\) and the tension
  \(\kappa\lambda_{\mathrm{F}} = 12H\) of the fold; in UM-4 a free
  parameter, bounded from below by the propagation of the gauge fields
  in the bulk. Unit: GeV (1/length). Value: \(H > 2534\) GeV, about 2.5
  TeV (from the invisible width of the Z boson). Status: AUTHOR.
\item
  \(\Gamma_{\mathrm{bulk}}\): The width of the W or the Z boson into the
  bulk (its decay into the massive modes of the gauge fields of the
  bulk), \(\Gamma_{\mathrm{bulk}} = m_0(m_0/H)^3/3\). Unit: GeV. Value:
  1.42 MeV for the Z and 0.85 MeV for the W at \(H = 2534\) GeV. Status:
  DERIVED.
\end{itemize}

The W and Z bosons. The Higgs vacuum on the fold gives the gauge fields
of the bulk a mass term concentrated on the fold,
\(M_{\mathrm{loc}}^2\,\delta_{\mathrm{D}}(\sigma)\) in their profile
equation (from the isospin current \(J_\Phi\) of FE-17 and the
hypercharge current of FE-18:
\(M_{\mathrm{loc}}^2 = \mathsf{g}_2^2(v^2 + \mathcal{A}^2)/(4V_*)\) for
the W), and the propagator of the bulk becomes
\(\mathbb{G}_5/(1 + M_{\mathrm{loc}}^2\mathbb{G}_5)\). Its pole for a
time-like momentum (\(\mathsf{p} = -\mathrm{i}m\), the continuation with
outgoing waves into the bulk) lies at the complex \(m^2\) of the first
relation: \(m_0\) is the mass that the zero mode alone would have,
\(m_0^2 = g_2^2(v^2 + \mathcal{A}^2)/4\) for the W and
\((g_1^2 + g_2^2)(v^2 + \mathcal{A}^2)/4\) for the Z (FE-45), and the
massive modes shift the squared mass by \(-m_0^4/(3H^2)\) (the mass by
\(-m_0^3/(6H^2)\)) and give the W and the Z the width
\(\Gamma_{\mathrm{bulk}} = m_0(m_0/H)^3/3\) into the bulk
(\(m^2 = m_{\mathrm{W}}^2 - \mathrm{i}m_{\mathrm{W}}\Gamma_{\mathrm{bulk}}\)):
the W and the Z are resonances of the bulk. A Z boson that decays into
the bulk leaves the fold unseen and adds to its invisible width. The
number of light neutrinos measured at the Z,
\(N_\nu = 2.9963 \pm 0.0074\) (P. Janot and S. Jadach, Phys. Lett. B 803
(2020) 135319), with the width 167.2 MeV of the Z into one neutrino and
its antineutrino, allows at most 1.42 MeV more invisible width (95 per
cent, one-sided), so
\(H > m_{\mathrm{Z}}\big(m_{\mathrm{Z}}/(3\cdot1.42\ \mathrm{MeV})\big)^{1/3} = 2534\)
GeV (\(m_{\mathrm{Z}}\) the mass of the Z; about 2.5 TeV), the strongest
bound on \(H\). The contact term \(x_{\mathrm{p}}^2/3\) of
\(\mathbb{G}_5\) is a contact interaction between the currents of the
fold whose strength is the square of the four-dimensional gauge coupling
divided by \(3H^2\); the limit of ATLAS on contact interactions between
quarks and leptons in the spectra of electron and muon pairs, written as
\(4\pi/\Lambda_{\mathrm{CI}}^2\) (\(\Lambda_{\mathrm{CI}}\) the scale of
a contact interaction), \(\Lambda_{\mathrm{CI}} > 26.0\) TeV for the
destructive left-left structure (JHEP 11 (2020) 005), gives with the
photon alone and the charge \(2/3\) of the up quark \(H > 1.1\) TeV (1.3
TeV with unit charges; the exchange of the massive photons is a
vector-vector interaction, so this is an estimate); the shift of the
masses changes the ratio
\(m_{\mathrm{W}}^2/(\cos^2\theta_{\mathrm{W}}m_{\mathrm{Z}}^2)\) (the
parameter rho of the electroweak fits, \(\theta_{\mathrm{W}}\) the weak
mixing angle) by \(\sin^2\theta_{\mathrm{W}}m_{\mathrm{Z}}^2/(3H^2)\),
and its measured value \(1.00031 \pm 0.00019\) (Particle Data Group
2024) gives \(H > 0.97\) TeV. At \(H = 2534\) GeV the width of the W
into the bulk is 0.85 MeV (its measured total width is 2.1 GeV). Status:
DERIVED (\texttt{um4\_propagators.py}, P4: the exact pole of the
rational propagator differs from the expansion by the relative amounts
\(1.83(m_0/H)^2\) in the imaginary part and \(0.44(m_0/H)^4\) in the
real part, the next orders, for \(m_0/H = 10^{-2}\) to \(10^{-3}\)),
ESTIMATE (the bounds, P5: the measured values are inputs).

\hypertarget{the-causal-structure}{%
\subsection{The causal structure}\label{the-causal-structure}}

\[
\tilde\sigma = \mathrm{sgn}(\sigma)\,\frac{\mathrm{e}^{H\lvert\sigma\rvert} - 1}{H}:\qquad
g_{MN}\,\mathrm{d}X^M\mathrm{d}X^N = \frac{\mathrm{d}\tilde\sigma^2 + \mathrm{d}x_1^2 + \mathrm{d}x_2^2 + \mathrm{d}x_3^2 - \mathrm{d}x_4^2
 - \mathrm{d}x_5^2 - \mathrm{d}x_6^2 - \mathrm{d}x_7^2}{\big(1 + H\lvert\tilde\sigma\rvert\big)^2},\qquad \Delta x_4 \ge L \tag{FE-83}
\] \textbf{Symbols in (FE-83).}

\begin{itemize}
\tightlist
\item
  \(\tilde\sigma\): The conformal coordinate of the hidden direction on
  the static block,
  \(\tilde\sigma = \mathrm{sgn}(\sigma)(\mathrm{e}^{H\lvert\sigma\rvert} - 1)/H\),
  in which the static block is conformally flat. Unit: length. Status:
  DERIVED.
\item
  \(\mathrm{sgn}\): The sign function: \(+1\) for a positive argument
  and \(-1\) for a negative one. Status: STANDARD.
\item
  \(\sigma\): The proper coordinate of the hidden direction:
  \(\mathrm{d}\sigma = \lvert\cot z\rvert\,\mathrm{d}x_0\),
  \(\sigma = \mathrm{sgn}(\pi/2 - z)\ln(\sin z)/(6H)\) on one period
  \(0 < z < \pi\); \(\sigma = 0\) is the fold \(z = \pi/2\),
  \(\sigma \to \mp\infty\) the tips \(z \to 0, \pi\); one period of
  \(x_0\) is the whole line \(-\infty < \sigma < \infty\). Unit: length.
  Status: DERIVED.
\item
  \(\mathrm{e}\): Euler's number, the base of the natural logarithm;
  \(\mathrm{e}^{X}\) is the exponential of \(X\). Unit: dimensionless.
  Value: 2.718281828. Status: STANDARD.
\item
  \(H\): The author's constant (an inverse length) that fixes the period
  \(\pi/(6H)\) of \(x_0\), the warp factor
  \(\mathrm{e}^{-2H\lvert\sigma\rvert}\) of the metric, the cosmological
  constant \(\Lambda = -21H^2\) and the tension
  \(\kappa\lambda_{\mathrm{F}} = 12H\) of the fold; in UM-4 a free
  parameter, bounded from below by the propagation of the gauge fields
  in the bulk. Unit: GeV (1/length). Value: \(H > 2534\) GeV, about 2.5
  TeV (from the invisible width of the Z boson). Status: AUTHOR.
\item
  \(g_{MN}\): The metric tensor of the eight-dimensional spacetime
  (signature \(+,+,+,+,-,-,-,-\) in the order \(x_0, \dots, x_7\)):
  \(g_{MN}\,\mathrm{d}X^M\mathrm{d}X^N\) is the squared interval,
  \(g^{MN}\) its inverse, \(\lvert g\rvert\) the absolute value of its
  determinant (\(\sqrt{\lvert g\rvert}\,\mathrm{d}^8X\) is the proper
  volume element). Unit: dimensionless. Status: AUTHOR.
\item
  \(M\): A coordinate index of eight dimensions (also with a numbered
  copy \(M_1, M_2, \dots\)): it runs over \(x_0, x_1, \dots, x_7\); a
  repeated index is summed. Status: STANDARD.
\item
  \(N\): A coordinate index of eight dimensions (also
  \(N_1, N_2, \dots\)), running over \(x_0, \dots, x_7\); a repeated
  index is summed. Status: STANDARD.
\item
  \(\mathrm{d}\): The differential: \(\mathrm{d}\) in front of a
  quantity is an infinitesimal change of it (as \(\mathrm{d}\sigma\), a
  small step of \(\sigma\)); the quotient of two such changes is a
  derivative; \(\mathrm{d}\) with an exponent, in front of several
  variables, is their volume element, the product of their
  differentials. Status: STANDARD.
\item
  \(X = \{x_0, \dots, x_7\}\): The author's coordinates of the
  eight-dimensional spacetime, \(X = \{x_0, x_1, \dots, x_7\}\) with the
  hidden direction first (Dirac\_Einstein-Lovelock\_4+4\_Nash.nb,
  section DEFCHART); \(X^M\) is the coordinate with the index \(M\) and
  \(\mathrm{d}^8X = \mathrm{d}x_0\,\mathrm{d}x_1\cdots\mathrm{d}x_7\)
  the coordinate volume element. Unit: length (each coordinate). Status:
  AUTHOR.
\item
  \(x_1\): The first coordinate of ordinary 3-space. Unit: length.
  Status: AUTHOR.
\item
  \(x_2\): The second coordinate of ordinary 3-space. Unit: length.
  Status: AUTHOR.
\item
  \(x_3\): The third coordinate of ordinary 3-space. Unit: length.
  Status: AUTHOR.
\item
  \(x_4\): The time; a prime on a function of \(x_4\) alone is
  \(\mathrm{d}/\mathrm{d}x_4\). Unit: length (= time, \(c = 1\)).
  Status: AUTHOR.
\item
  \(x_5\): The first of the three extra time-like coordinates. Unit:
  length. Status: AUTHOR.
\item
  \(x_6\): The second of the three extra time-like coordinates. Unit:
  length. Status: AUTHOR.
\item
  \(x_7\): The third of the three extra time-like coordinates. Unit:
  length. Status: AUTHOR.
\item
  \(\Delta\): A change, difference or small variation of the quantity
  that follows (for example an interval of the time, a perturbation of a
  density, a difference of two squared masses, or an addition to an
  effective number of species). Status: STANDARD.
\item
  \(L\): A distance between two points of the fold along ordinary space:
  for a gravitational wave the radius of a ring of free test masses
  (\(\Delta L\) the change of the distance of one of them from the
  centre), for a signal the distance between its emitter and its
  receiver. Unit: length. Status: STANDARD.
\end{itemize}

The static block in the conformal coordinate \(\tilde\sigma\) of the
hidden direction: the metric is flat space of signature (4, 4) times the
factor \((1 + H\lvert\tilde\sigma\rvert)^{-2}\). A conformal factor does
not change which curves are null, and the null geodesics are the null
straight lines of flat space (each with its own parameter). A signal
that leaves a point of the fold, travels through the bulk at the local
speed of light at fixed extra times
(\(\mathrm{d}x_5 = \mathrm{d}x_6 = \mathrm{d}x_7 = 0\)) and arrives at
another point of the fold at the distance \(L\) along ordinary space
needs the time
\(\Delta x_4 = \int\sqrt{\mathrm{d}x_3^2 + \mathrm{d}\tilde\sigma^2} \ge L\),
with equality only for the path along the fold: nothing that stays at
fixed extra times crosses the static bulk faster than light on the fold.
The bound is a statement about the section of fixed extra times, the
spacetime of the fields that do not depend on the extra times (all
fields of the fold, and the modes of the bulk with no momentum along
them). In the full spacetime of signature (4, 4) it does not hold: a
null straight line that also moves by \(Y\) along an extra time reaches
the distance \(L\) in the time \(\Delta x_4 = \sqrt{L^2 - Y^2}\), down
to 0 for \(Y = L\) (the closed time-like curves of C4 are of the same
origin). (In UM-3 the light cone of the bulk at the depth \(\sigma\) was
\(\mathrm{e}^{H\lvert\sigma\rvert}\) times wider than that of the fold,
and signals through the bulk arrived first.) Status: DERIVED
(\texttt{um4\_causality.py}, C1: the conformal form exactly; the
geodesic equation of the conformal metric for a null straight line; at
fixed extra times the triangle inequality of the flat section, 200
sampled bulk paths all later than the path along the fold; the null line
along an extra time).

\[
c_{\mathrm{loc}}(\sigma) = 1 - \frac{a_4''}{5H^2}\big(\mathrm{e}^{2H\lvert\sigma\rvert} - 1\big),\qquad
\sigma(x_3) = \frac{\lvert a_4''\rvert}{5H}\,x_3\,(L - x_3),\qquad
\Delta x_4 = L - \frac{a_4''^2L^3}{150H^2} \tag{FE-84}
\] \textbf{Symbols in (FE-84).}

\begin{itemize}
\tightlist
\item
  \(c_{\mathrm{loc}}\): The coordinate speed of light along \(x_3\) at
  the depth \(\sigma\) (the rate of change of \(x_3\) with \(x_4\) along
  a null line at fixed \(\sigma\)), in units of that on the fold:
  \(\mathrm{e}^{H\lvert\sigma\rvert}\) on the static block of UM-3,
  whose time was not warped, 1 on the static block of UM-4, and on an
  expanding fold of UM-4 different from 1 in proportion to the second
  derivative \(a_4''\) of \(a_4\) (the rate of change of the expansion
  rate \(a_4'\)). Unit: dimensionless. Status: DERIVED.
\item
  \(\sigma\): The proper coordinate of the hidden direction:
  \(\mathrm{d}\sigma = \lvert\cot z\rvert\,\mathrm{d}x_0\),
  \(\sigma = \mathrm{sgn}(\pi/2 - z)\ln(\sin z)/(6H)\) on one period
  \(0 < z < \pi\); \(\sigma = 0\) is the fold \(z = \pi/2\),
  \(\sigma \to \mp\infty\) the tips \(z \to 0, \pi\); one period of
  \(x_0\) is the whole line \(-\infty < \sigma < \infty\). Unit: length.
  Status: DERIVED.
\item
  \(a_4\): The author's metric function of the time \(x_4\) that scales
  ordinary 3-space: on the fold \(\mathrm{e}^{a_4}\) is the scale factor
  of cosmology and \(a_4'\) the Hubble rate. Unit: dimensionless.
  Status: AUTHOR.
\item
  \({}'\): The derivative of a function of the time \(x_4\) alone with
  respect to \(x_4\) (\(a_4' = \mathrm{d}a_4/\mathrm{d}x_4\); two primes
  the second derivative); on a function of \(\sigma\) alone (a profile)
  the derivative with respect to \(\sigma\). Status: STANDARD.
\item
  \(H\): The author's constant (an inverse length) that fixes the period
  \(\pi/(6H)\) of \(x_0\), the warp factor
  \(\mathrm{e}^{-2H\lvert\sigma\rvert}\) of the metric, the cosmological
  constant \(\Lambda = -21H^2\) and the tension
  \(\kappa\lambda_{\mathrm{F}} = 12H\) of the fold; in UM-4 a free
  parameter, bounded from below by the propagation of the gauge fields
  in the bulk. Unit: GeV (1/length). Value: \(H > 2534\) GeV, about 2.5
  TeV (from the invisible width of the Z boson). Status: AUTHOR.
\item
  \(\mathrm{e}\): Euler's number, the base of the natural logarithm;
  \(\mathrm{e}^{X}\) is the exponential of \(X\). Unit: dimensionless.
  Value: 2.718281828. Status: STANDARD.
\item
  \(x_3\): The third coordinate of ordinary 3-space. Unit: length.
  Status: AUTHOR.
\item
  \(L\): A distance between two points of the fold along ordinary space:
  for a gravitational wave the radius of a ring of free test masses
  (\(\Delta L\) the change of the distance of one of them from the
  centre), for a signal the distance between its emitter and its
  receiver. Unit: length. Status: STANDARD.
\item
  \(\Delta\): A change, difference or small variation of the quantity
  that follows (for example an interval of the time, a perturbation of a
  density, a difference of two squared masses, or an addition to an
  effective number of species). Status: STANDARD.
\item
  \(x_4\): The time; a prime on a function of \(x_4\) alone is
  \(\mathrm{d}/\mathrm{d}x_4\). Unit: length (= time, \(c = 1\)).
  Status: AUTHOR.
\end{itemize}

An expanding fold with frozen extra times, with the bulk metric FE-34 to
the first order in its curvature: the coordinate speed of light
\(c_{\mathrm{loc}}\) along \(x_3\) at the depth \(\sigma\), in units of
the speed of light on the fold, differs from 1 in proportion to the
second derivative \(a_4'' = H_4' = -(1 + q_{\mathrm{dec}})H_4^2\) of the
logarithm of the scale factor (\(H_4 = a_4'\) the expansion rate,
\(q_{\mathrm{dec}} = -1 - a_4''/a_4'^2\) the deceleration parameter).
Wherever the expansion rate falls (\(a_4'' < 0\), i.e.
\(q_{\mathrm{dec}} > -1\): the eras of radiation and of matter, and also
today, \(q_{\mathrm{dec}} = -0.55\)) the light cone of the bulk is
slightly wider than that of the fold, and a signal can gain time by
dipping into the bulk; it is narrower only where the expansion rate
grows (\(q_{\mathrm{dec}} < -1\)), and then the path along the fold is
the fastest. Where \(a_4'' < 0\) the path that gains the most between
two points of the fold at the distance \(L\) (for distances small
compared with \(1/H_4\)) is the parabola of the second relation, and it
arrives earlier than light along the fold by \(a_4''^2L^3/(150H^2)\)
(the third relation; the second and third relations hold for
\(a_4'' < 0\)). Relative to the travel time this is
\((1 + q_{\mathrm{dec}})^2(H_4L)^2(H_4/H)^2/150\); extrapolated to the
Hubble distance today (\(L = 1/H_4\), \(q_{\mathrm{dec}} = -0.55\),
\(H_4 = 1.44\times10^{-42}\) GeV; an order of magnitude) it is
\(4.4\times10^{-94}\) for \(H = 2534\) GeV (\(4.5\times10^{-94}\) for
\(H = 2.5\) TeV, below \(3\times10^{-63}\) even for \(H = 10^{-3}\) eV):
the bulk is ahead of light on the fold, but by an unobservable amount;
on the static fold no signal at fixed extra times is ahead at all.
Status: DERIVED (\texttt{um4\_causality.py}, C2: the speed from the
metric of FE-34 and
\(\hat R^{(4)}{}_{(4)} - \hat R^{(1)}{}_{(1)} = 2a_4''\) computed from
the metric of the fold; C3: the advance by minimising the travel time
numerically, 0.997 of the formula, and no advance where \(a_4'' > 0\)).

\hypertarget{dirac16complex-modes-and-kohn-sham-equations}{%
\section{dirac16complex: modes and Kohn-Sham
equations}\label{dirac16complex-modes-and-kohn-sham-equations}}

dirac16complex lives in the bulk; its four-dimensional particles are its
states at the fold. This section writes the equation of those states on
the static block (the late universe; an expanding fold replaces \(a_4\),
\(b_4\) by their values at the time), its exact solutions, the
conditions at the fold, the resonances that the fold holds, and the
Kohn-Sham equations that make the description self-consistent with the
self-interaction, with gravity and with the mean fields of the Standard
Model.

\hypertarget{the-equation-of-the-hidden-direction}{%
\subsection{The equation of the hidden
direction}\label{the-equation-of-the-hidden-direction}}

\[
\Psi = \mathrm{e}^{\frac72H\lvert\sigma\rvert}\,\chi(\sigma)\,\mathrm{e}^{\mathrm{i}(\vec k\cdot\vec x - \omega x_4)},\qquad
\gamma^{(0)}\chi' + \mathrm{e}^{H\lvert\sigma\rvert}\big(-\mathrm{i}\omega\,\gamma^{(4)} + \mathrm{i}\,\vec k\cdot\vec\gamma\big)\chi
 - m_\Psi\chi = 0\qquad(\sigma \ne 0) \tag{FE-36}
\] \textbf{Symbols in (FE-36).}

\begin{itemize}
\tightlist
\item
  \(\Psi,\ \bar\Psi\): dirac16complex: the author's 16-component complex
  spinor field of eight dimensions;
  \(\bar\Psi = \Psi^\dagger\sigma_{16}\) is its Dirac adjoint; in a term
  of the fold, \(\Psi\) is its value at \(\sigma = 0\). Unit:
  GeV\(^{7/2}\). Status: AUTHOR.
\item
  \(\mathrm{e}\): Euler's number, the base of the natural logarithm;
  \(\mathrm{e}^{X}\) is the exponential of \(X\). Unit: dimensionless.
  Value: 2.718281828. Status: STANDARD.
\item
  \(H\): The author's constant (an inverse length) that fixes the period
  \(\pi/(6H)\) of \(x_0\), the warp factor
  \(\mathrm{e}^{-2H\lvert\sigma\rvert}\) of the metric, the cosmological
  constant \(\Lambda = -21H^2\) and the tension
  \(\kappa\lambda_{\mathrm{F}} = 12H\) of the fold; in UM-4 a free
  parameter, bounded from below by the propagation of the gauge fields
  in the bulk. Unit: GeV (1/length). Value: \(H > 2534\) GeV, about 2.5
  TeV (from the invisible width of the Z boson). Status: AUTHOR.
\item
  \(\sigma\): The proper coordinate of the hidden direction:
  \(\mathrm{d}\sigma = \lvert\cot z\rvert\,\mathrm{d}x_0\),
  \(\sigma = \mathrm{sgn}(\pi/2 - z)\ln(\sin z)/(6H)\) on one period
  \(0 < z < \pi\); \(\sigma = 0\) is the fold \(z = \pi/2\),
  \(\sigma \to \mp\infty\) the tips \(z \to 0, \pi\); one period of
  \(x_0\) is the whole line \(-\infty < \sigma < \infty\). Unit: length.
  Status: DERIVED.
\item
  \(\chi(\sigma)\): The profile of a mode of dirac16complex along the
  hidden direction,
  \(\Psi = \mathrm{e}^{\frac72H\lvert\sigma\rvert}\chi(\sigma)\mathrm{e}^{\mathrm{i}(\vec k\cdot\vec x - \omega x_4)}\)
  (a 16-component spinor function of \(\sigma\)); as a four-dimensional
  field, \(\chi\) denotes the particle of the level of dirac16complex at
  the fold (a long-lived resonance) and \(\bar\chi\) its antiparticle.
  Unit: GeV\(^{2}\) (with the \(V_*\) normalisation). Status: DERIVED.
\item
  \(\mathrm{i}\): The imaginary unit, \(\mathrm{i}^2 = -1\). Unit:
  dimensionless. Status: STANDARD.
\item
  \(\vec k\): The (comoving) momentum along ordinary space, \(\vec k^2\)
  its square. Unit: GeV. Status: STANDARD.
\item
  \(\vec x\): The position in ordinary space,
  \(\vec x = (x_1, x_2, x_3)\). Unit: length. Status: AUTHOR.
\item
  \(\omega\): A frequency: the energy of a mode
  \(\propto \mathrm{e}^{-\mathrm{i}\omega x_4}\) (\(\omega_o\) of the
  orbital \(o\)). Unit: GeV. Status: STANDARD.
\item
  \(x_4\): The time; a prime on a function of \(x_4\) alone is
  \(\mathrm{d}/\mathrm{d}x_4\). Unit: length (= time, \(c = 1\)).
  Status: AUTHOR.
\item
  \(\gamma^M,\ \gamma^{(a)}\): The Dirac matrices of dirac16complex
  along the coordinates, \(\gamma^M = \mathbf{e}_a{}^MT^a\),
  \(\gamma_M = g_{MN}\gamma^N\); \(\gamma^{(a)} = T^a\) along the frame
  direction \((a)\) (\(\gamma^{(0)}\) along \(\sigma\), \(\gamma^{(4)}\)
  along the time). Unit: dimensionless. Status: AUTHOR.
\item
  \({}'\): The derivative of a function of the time \(x_4\) alone with
  respect to \(x_4\) (\(a_4' = \mathrm{d}a_4/\mathrm{d}x_4\); two primes
  the second derivative); on a function of \(\sigma\) alone (a profile)
  the derivative with respect to \(\sigma\). Status: STANDARD.
\item
  \(\vec\gamma\): The three frame Dirac matrices of dirac16complex along
  ordinary space,
  \(\vec\gamma = (\gamma^{(1)}, \gamma^{(2)}, \gamma^{(3)})\). Unit:
  dimensionless. Status: AUTHOR.
\item
  \(m_\Psi\): The mass parameter of dirac16complex in the bulk. Unit:
  GeV. Status: RECORD.
\end{itemize}

The free field equation FE-12 on the static block (\(a_4\) and \(b_4\)
constant, ordinary space and the extra times normalised so that
\(a_4 = 0\) and \(b_4 = 1\)) for a mode of definite momentum \(\vec k\)
along ordinary space and frequency \(\omega\), independent of the extra
times. The factor \(\mathrm{e}^{(7/2)H\lvert\sigma\rvert}\) removes the
spin connection exactly
(\(\gamma^{(0)}\partial_\sigma \mathrm{e}^{(7/2)H\lvert\sigma\rvert} = \frac72H\,\mathrm{sgn}(\sigma)\gamma^{(0)}\mathrm{e}^{(7/2)H\lvert\sigma\rvert}\)
cancels the first term of FE-14), so that the profile \(\chi(\sigma)\)
obeys a Dirac equation of one variable in which the factor
\(\mathrm{e}^{H\lvert\sigma\rvert}\) (one over the square root of the
warp factor) multiplies the whole operator of the fold, the frequency as
well as the momentum: the time is warped like ordinary space (in UM-3
the factor was \(\mathrm{e}^{3H\lvert\sigma\rvert}\), and
\(\mathrm{e}^{H\lvert\sigma\rvert}\) multiplied only the momentum). The
density of the conserved charge along \(\sigma\) is
\(\mathrm{e}^{H\lvert\sigma\rvert}\chi^\dagger\mathfrak{B}\chi\) (times
\(b_4^3\); FE-43). Status: DERIVED (\texttt{um4\_geometry.py}, G9, with
the author's matrices).

\[
\begin{aligned}
&\chi' = \mathbb{K}\,\chi,\qquad \mathbb{K} = \gamma^{(0)}\big(m_\Psi + \mathrm{i}\omega\,\mathrm{e}^{H\lvert\sigma\rvert}\gamma^{(4)}\big),\qquad
\mathbb{K}^2 = \big(m_\Psi^2 - \omega^2\mathrm{e}^{2H\lvert\sigma\rvert}\big)\mathbf{1},\\
&q = \sqrt{m_\Psi^2 - \omega^2\mathrm{e}^{2H\lvert\sigma\rvert}},\qquad \lvert\omega\rvert\,\mathrm{e}^{H\sigma_{\mathrm{t}}} = m_\Psi
\end{aligned} \tag{FE-37}
\] \textbf{Symbols in (FE-37).}

\begin{itemize}
\tightlist
\item
  \(\chi(\sigma)\): The profile of a mode of dirac16complex along the
  hidden direction,
  \(\Psi = \mathrm{e}^{\frac72H\lvert\sigma\rvert}\chi(\sigma)\mathrm{e}^{\mathrm{i}(\vec k\cdot\vec x - \omega x_4)}\)
  (a 16-component spinor function of \(\sigma\)); as a four-dimensional
  field, \(\chi\) denotes the particle of the level of dirac16complex at
  the fold (a long-lived resonance) and \(\bar\chi\) its antiparticle.
  Unit: GeV\(^{2}\) (with the \(V_*\) normalisation). Status: DERIVED.
\item
  \({}'\): The derivative of a function of the time \(x_4\) alone with
  respect to \(x_4\) (\(a_4' = \mathrm{d}a_4/\mathrm{d}x_4\); two primes
  the second derivative); on a function of \(\sigma\) alone (a profile)
  the derivative with respect to \(\sigma\). Status: STANDARD.
\item
  \(\mathbb{K}\): The transfer matrix of the profile equation at zero
  momentum, \(\chi' = \mathbb{K}\chi\). Unit: GeV. Status: DERIVED.
\item
  \(\gamma^M,\ \gamma^{(a)}\): The Dirac matrices of dirac16complex
  along the coordinates, \(\gamma^M = \mathbf{e}_a{}^MT^a\),
  \(\gamma_M = g_{MN}\gamma^N\); \(\gamma^{(a)} = T^a\) along the frame
  direction \((a)\) (\(\gamma^{(0)}\) along \(\sigma\), \(\gamma^{(4)}\)
  along the time). Unit: dimensionless. Status: AUTHOR.
\item
  \(m_\Psi\): The mass parameter of dirac16complex in the bulk. Unit:
  GeV. Status: RECORD.
\item
  \(\mathrm{i}\): The imaginary unit, \(\mathrm{i}^2 = -1\). Unit:
  dimensionless. Status: STANDARD.
\item
  \(\omega\): A frequency: the energy of a mode
  \(\propto \mathrm{e}^{-\mathrm{i}\omega x_4}\) (\(\omega_o\) of the
  orbital \(o\)). Unit: GeV. Status: STANDARD.
\item
  \(\mathrm{e}\): Euler's number, the base of the natural logarithm;
  \(\mathrm{e}^{X}\) is the exponential of \(X\). Unit: dimensionless.
  Value: 2.718281828. Status: STANDARD.
\item
  \(H\): The author's constant (an inverse length) that fixes the period
  \(\pi/(6H)\) of \(x_0\), the warp factor
  \(\mathrm{e}^{-2H\lvert\sigma\rvert}\) of the metric, the cosmological
  constant \(\Lambda = -21H^2\) and the tension
  \(\kappa\lambda_{\mathrm{F}} = 12H\) of the fold; in UM-4 a free
  parameter, bounded from below by the propagation of the gauge fields
  in the bulk. Unit: GeV (1/length). Value: \(H > 2534\) GeV, about 2.5
  TeV (from the invisible width of the Z boson). Status: AUTHOR.
\item
  \(\sigma\): The proper coordinate of the hidden direction:
  \(\mathrm{d}\sigma = \lvert\cot z\rvert\,\mathrm{d}x_0\),
  \(\sigma = \mathrm{sgn}(\pi/2 - z)\ln(\sin z)/(6H)\) on one period
  \(0 < z < \pi\); \(\sigma = 0\) is the fold \(z = \pi/2\),
  \(\sigma \to \mp\infty\) the tips \(z \to 0, \pi\); one period of
  \(x_0\) is the whole line \(-\infty < \sigma < \infty\). Unit: length.
  Status: DERIVED.
\item
  \(\mathbf{1},\ \mathbf{1}_4,\ \mathbf{1}_8,\ \mathbf{1}_{16}\): The
  identity matrix of the size fixed by the context (\(16 \times 16\) on
  dirac16complex, \(4 \times 4\) on a four-dimensional spinor,
  \(3 \times 3\) in generation or flavour space); \(\mathbf{1}_4\),
  \(\mathbf{1}_8\) and \(\mathbf{1}_{16}\) are the identity matrices of
  4, 8 and 16 rows. Unit: dimensionless. Status: STANDARD.
\item
  \(q\): The local decay rate of a profile of dirac16complex along
  \(\sigma\) at the frequency \(\omega\),
  \(q = \sqrt{m_\Psi^2 - \omega^2\mathrm{e}^{2H\lvert\sigma\rvert}}\)
  (imaginary beyond the turning point, where the profile oscillates).
  Unit: GeV. Status: DERIVED.
\item
  \(\sigma_{\mathrm{t}}\): The turning point of a state of
  dirac16complex of the frequency \(\omega\): the depth at which the
  warped frequency equals the bulk mass,
  \(\lvert\omega\rvert\mathrm{e}^{H\sigma_{\mathrm{t}}} = m_\Psi\);
  beyond it every solution oscillates. Unit: length. Status: DERIVED.
\end{itemize}

At zero momentum along ordinary space (\(\vec k = 0\)) the profile
equation is a first-order system whose matrix \(\mathbb{K}\) depends on
\(\sigma\) through the warped frequency; its square is the local gap.
Where \(\lvert\omega\rvert\mathrm{e}^{H\lvert\sigma\rvert} < m_\Psi\)
the solutions grow or decay locally at the rate \(q(\sigma)\); at the
turning point \(\sigma_{\mathrm{t}}\) the local gap closes, and beyond
it every solution oscillates: the warped frequency exceeds the bulk mass
far from the fold for every \(\omega \ne 0\). There is therefore no gap
in the bulk, no threshold of the continuum other than \(\omega = 0\),
and no normalisable state of non-zero frequency: the states that UM-3
bound to the fold are resonances in UM-4 (FE-86). The 16 components
split into eight copies of two inequivalent problems of two components
(the element \(\gamma^{(0)}\gamma^{(4)}\Gamma_8\) squares to
\(\mathbf{1}\) and commutes with \(\gamma^{(0)}\), \(\gamma^{(4)}\) and
\(\Gamma_8\)), in each of which \(\gamma^{(0)}\),
\(\mathrm{i}\gamma^{(0)}\gamma^{(4)}\) and \(\Gamma_8\) act as the Pauli
matrices \(\mathsf{s}_3\), \(\mathrm{i}\mathsf{s}_1\) and
\(\pm\mathsf{s}_1\) (FE-85). With chirality-dependent masses
\(m_{\mathrm{L}}P_{\mathrm{L}} + m_{\mathrm{R}}P_{\mathrm{R}}\) in place
of \(m_\Psi\) the local gap would be
\(m_{\mathrm{L}}m_{\mathrm{R}} - \omega^2\mathrm{e}^{2H\lvert\sigma\rvert}\)
(tachyonic for \(m_{\mathrm{L}}m_{\mathrm{R}} < 0\)); UM-4 keeps the
record's single bulk mass and puts the difference of the halves into the
fold terms. Status: DERIVED (\texttt{um4\_spectra.py}, U1, with the
author's matrices).

\[
\mathrm{e}^{2H\lvert\sigma\rvert}\Big(\omega^2 - \mathrm{e}^{-2a_4}\,\vec k^2 + \mathrm{e}^{2a_4}\,b_4^{-2}\,\vec\kappa^2\Big) = m_\Psi^2 + p_\sigma^2,\qquad
\omega^2 = \omega_0^2 + \vec k^2 - \vec\kappa^2\quad\big(a_4 = 0,\ b_4 = 1\big) \tag{FE-38}
\] \textbf{Symbols in (FE-38).}

\begin{itemize}
\tightlist
\item
  \(\mathrm{e}\): Euler's number, the base of the natural logarithm;
  \(\mathrm{e}^{X}\) is the exponential of \(X\). Unit: dimensionless.
  Value: 2.718281828. Status: STANDARD.
\item
  \(H\): The author's constant (an inverse length) that fixes the period
  \(\pi/(6H)\) of \(x_0\), the warp factor
  \(\mathrm{e}^{-2H\lvert\sigma\rvert}\) of the metric, the cosmological
  constant \(\Lambda = -21H^2\) and the tension
  \(\kappa\lambda_{\mathrm{F}} = 12H\) of the fold; in UM-4 a free
  parameter, bounded from below by the propagation of the gauge fields
  in the bulk. Unit: GeV (1/length). Value: \(H > 2534\) GeV, about 2.5
  TeV (from the invisible width of the Z boson). Status: AUTHOR.
\item
  \(\sigma\): The proper coordinate of the hidden direction:
  \(\mathrm{d}\sigma = \lvert\cot z\rvert\,\mathrm{d}x_0\),
  \(\sigma = \mathrm{sgn}(\pi/2 - z)\ln(\sin z)/(6H)\) on one period
  \(0 < z < \pi\); \(\sigma = 0\) is the fold \(z = \pi/2\),
  \(\sigma \to \mp\infty\) the tips \(z \to 0, \pi\); one period of
  \(x_0\) is the whole line \(-\infty < \sigma < \infty\). Unit: length.
  Status: DERIVED.
\item
  \(\omega\): A frequency: the energy of a mode
  \(\propto \mathrm{e}^{-\mathrm{i}\omega x_4}\) (\(\omega_o\) of the
  orbital \(o\)). Unit: GeV. Status: STANDARD.
\item
  \(a_4\): The author's metric function of the time \(x_4\) that scales
  ordinary 3-space: on the fold \(\mathrm{e}^{a_4}\) is the scale factor
  of cosmology and \(a_4'\) the Hubble rate. Unit: dimensionless.
  Status: AUTHOR.
\item
  \(\vec k\): The (comoving) momentum along ordinary space, \(\vec k^2\)
  its square. Unit: GeV. Status: STANDARD.
\item
  \(b_4\): The author's metric function of the time \(x_4\) in the
  extra-time components, \(b_4 > 0\), never identically 1 (the author's
  order of 2026-10-24); the scale factor (the factor of the lengths) of
  the extra times is \(\mathrm{e}^{-a_4}b_4\) on the fold and
  \(\mathrm{e}^{-a_4}b_4\,\mathrm{e}^{-H\lvert\sigma\rvert}\) at the
  distance \(\sigma\) from it (the square root of the warp factor times
  \(\mathrm{e}^{-a_4}b_4\)). Unit: dimensionless. Status: AUTHOR.
\item
  \(\vec\kappa\): A momentum along the extra times (three components);
  the fields of UM-4 have \(\vec\kappa = 0\). Unit: GeV. Status:
  DERIVED.
\item
  \(m_\Psi\): The mass parameter of dirac16complex in the bulk. Unit:
  GeV. Status: RECORD.
\item
  \(p_\sigma\): The momentum along the hidden direction (local wave
  number of a profile along \(\sigma\)). Unit: GeV. Status: DERIVED.
\item
  \(\omega_0\): The energy at rest (zero momentum along ordinary space
  and along the extra times) of a state of dirac16complex at the fold;
  complex for a resonance,
  \(\omega_0 = \omega_{\mathrm{r}} - \frac{\mathrm{i}}{2}\Gamma_\chi\).
  Unit: GeV. Status: DERIVED.
\end{itemize}

The local dispersion relation of dirac16complex (WKB: for profiles that
vary slowly compared with the warp), with the momentum \(p_\sigma\)
along the hidden direction and a momentum \(\vec\kappa\) along the extra
times, and the exact dispersion of every state at the fold with the
energy \(\omega_0\) at rest. The second relation holds exactly because
the boosts along ordinary space (FE-78) and the rotations of the time
into an extra time (generated by \(\frac12\gamma^{(4)}\gamma^{(5)}\),
which commutes with \(\gamma^{(0)}\), with \(\Gamma_8\) and with the
mass term) are symmetries of the bulk equation and of the fold terms of
the vacuum: the warp multiplies the frequency and the momenta alike. A
momentum along a time-like direction enters with the sign of the
frequency. A state with \(\lvert\vec\kappa\rvert \le \omega_0\) is the
rotation of the state at rest in the plane of the time and an extra time
(by the angle whose sine is \(\lvert\vec\kappa\rvert/\omega_0\)) and has
its width; a state with \(\lvert\vec\kappa\rvert > \omega_0\) is reached
by no rotation: it has \(\omega^2 < 0\), the continuation of the same
dispersion, and grows exponentially in \(x_4\) (a tachyon). The same
holds for every field, in the bulk as on the fold (the sign comes from
\(g^{rr} < 0\)): a field of the mass \(m\) has growing modes for every
\(\vec\kappa^2 > m^2 + \vec k^2\) (a massless field for every
\(\lvert\vec\kappa\rvert > \lvert\vec k\rvert\)). UM-4 therefore keeps
only the fields that do not depend on the extra times (ASSUMED); this
restriction is consistent (the metric does not depend on
\(x_5, x_6, x_7\), so the momentum along the extra times is conserved
and a field without it stays without it classically) but not stable
(interactions create pairs of opposite momenta, an ESTIMATE: no rate is
computed). Compact extra times of the period \(L_{\mathrm{t}}\) do not
help: the first harmonic of a massless field stays below one e-folding
over the age of the universe only for \(L_{\mathrm{t}} \ge 26.6\) Gpc
(\(2\pi\) times the age) on the static fold, or about 13.5 Gpc with the
Hubble friction of the expanding fold held at its value today (an order
of magnitude), but the \(n\)-th harmonic grows \(n\) times faster, so
infinitely many modes grow for every \(L_{\mathrm{t}}\) (and a compact
time-like direction is itself a closed time-like curve). UM-4 needs a
gauge principle that removes the momentum along the extra times of every
field. (In UM-3 the time was not warped, the momentum along the extra
times carried the factor \(\mathrm{e}^{H\lvert\sigma\rvert}\) and the
frequency not, and a bound state with \(\vec\kappa \ne 0\) leaked into
the bulk beyond a finite depth.) Status: DERIVED (the dispersion and the
instability; \texttt{um4\_extratimes.py}, EX1 to EX5:
\(\omega^2 = \omega_0^2 - \vec\kappa^2\) and \(\omega_0^2 + \vec k^2\)
computed with the author's matrices for
\(\lvert\vec\kappa\rvert, \lvert\vec k\rvert = 0.1\), 0.3, 0.6
\(m_\Psi\) to ten digits), ASSUMED (the restriction), OPEN (the gauge
principle).

\[
\chi = (\omega\zeta_{\mathrm{w}})^{1/2}\begin{pmatrix}\mathfrak{z}_-(\omega\zeta_{\mathrm{w}})\\ \mathrm{i}\,\mathfrak{z}_+(\omega\zeta_{\mathrm{w}})\end{pmatrix}\quad(\sigma > 0),\qquad
\chi = (\omega\zeta_{\mathrm{w}})^{1/2}\begin{pmatrix}-\mathrm{i}\,\mathfrak{z}_+(\omega\zeta_{\mathrm{w}})\\ \mathfrak{z}_-(\omega\zeta_{\mathrm{w}})\end{pmatrix}\quad(\sigma < 0),\qquad
\zeta_{\mathrm{w}} = \frac{\mathrm{e}^{H\lvert\sigma\rvert}}{H} \tag{FE-85}
\] \textbf{Symbols in (FE-85).}

\begin{itemize}
\tightlist
\item
  \(\chi(\sigma)\): The profile of a mode of dirac16complex along the
  hidden direction,
  \(\Psi = \mathrm{e}^{\frac72H\lvert\sigma\rvert}\chi(\sigma)\mathrm{e}^{\mathrm{i}(\vec k\cdot\vec x - \omega x_4)}\)
  (a 16-component spinor function of \(\sigma\)); as a four-dimensional
  field, \(\chi\) denotes the particle of the level of dirac16complex at
  the fold (a long-lived resonance) and \(\bar\chi\) its antiparticle.
  Unit: GeV\(^{2}\) (with the \(V_*\) normalisation). Status: DERIVED.
\item
  \(\omega\): A frequency: the energy of a mode
  \(\propto \mathrm{e}^{-\mathrm{i}\omega x_4}\) (\(\omega_o\) of the
  orbital \(o\)). Unit: GeV. Status: STANDARD.
\item
  \(\zeta_{\mathrm{w}}\): The warped distance from the fold,
  \(\zeta_{\mathrm{w}} = \mathrm{e}^{H\lvert\sigma\rvert}/H\) (\(1/H\)
  at the fold): \(\omega\zeta_{\mathrm{w}}\) is the argument of the
  exact solutions of dirac16complex, and \(m\zeta_{\mathrm{w}}\) that of
  the massive modes of the mass \(m\) of the massless fields of the
  bulk. Unit: length. Status: DERIVED.
\item
  \(\mathfrak{z}_-,\ \mathfrak{z}_+\): Bessel functions of the orders
  \(m_\Psi/H - \frac12\) (\(\mathfrak{z}_-\)) and \(m_\Psi/H + \frac12\)
  (\(\mathfrak{z}_+\)), both of the same kind (of the first kind, of the
  second kind, or the Hankel functions of the first kind for waves that
  leave the fold), of the argument \(\omega\zeta_{\mathrm{w}}\). Unit:
  dimensionless. Status: STANDARD.
\item
  \(\mathrm{i}\): The imaginary unit, \(\mathrm{i}^2 = -1\). Unit:
  dimensionless. Status: STANDARD.
\item
  \(\sigma\): The proper coordinate of the hidden direction:
  \(\mathrm{d}\sigma = \lvert\cot z\rvert\,\mathrm{d}x_0\),
  \(\sigma = \mathrm{sgn}(\pi/2 - z)\ln(\sin z)/(6H)\) on one period
  \(0 < z < \pi\); \(\sigma = 0\) is the fold \(z = \pi/2\),
  \(\sigma \to \mp\infty\) the tips \(z \to 0, \pi\); one period of
  \(x_0\) is the whole line \(-\infty < \sigma < \infty\). Unit: length.
  Status: DERIVED.
\item
  \(\mathrm{e}\): Euler's number, the base of the natural logarithm;
  \(\mathrm{e}^{X}\) is the exponential of \(X\). Unit: dimensionless.
  Value: 2.718281828. Status: STANDARD.
\item
  \(H\): The author's constant (an inverse length) that fixes the period
  \(\pi/(6H)\) of \(x_0\), the warp factor
  \(\mathrm{e}^{-2H\lvert\sigma\rvert}\) of the metric, the cosmological
  constant \(\Lambda = -21H^2\) and the tension
  \(\kappa\lambda_{\mathrm{F}} = 12H\) of the fold; in UM-4 a free
  parameter, bounded from below by the propagation of the gauge fields
  in the bulk. Unit: GeV (1/length). Value: \(H > 2534\) GeV, about 2.5
  TeV (from the invisible width of the Z boson). Status: AUTHOR.
\end{itemize}

The exact solutions of FE-37 in one copy of two components
(\(\gamma^{(0)} = \mathsf{s}_3\),
\(\mathrm{i}\gamma^{(0)}\gamma^{(4)} = \mathrm{i}\mathsf{s}_1\)):
\(\mathfrak{z}_-\) and \(\mathfrak{z}_+\) are Bessel functions of the
orders \(m_\Psi/H - \frac12\) and \(m_\Psi/H + \frac12\) of the argument
\(\omega\zeta_{\mathrm{w}}\) (the same kind for both: of the first kind,
of the second kind, or a Hankel function), with \(\zeta_{\mathrm{w}}\)
the warped distance from the fold (\(\zeta_{\mathrm{w}} = 1/H\) at the
fold). For \(\omega\zeta_{\mathrm{w}} \ll m_\Psi/H\) (close to the fold)
the solutions grow or decay like
\(\mathrm{e}^{\pm q\lvert\sigma\rvert}\); for
\(\omega\zeta_{\mathrm{w}} \gg m_\Psi/H\) (beyond the turning point)
they are waves, and the outgoing ones, which carry the flux away from
the fold, are the Hankel functions of the first kind. Status: DERIVED
(\texttt{um4\_spectra.py}, U2: the solutions checked against the
differential equation at three depths).

\hypertarget{the-fold-and-the-resonances-that-it-holds}{%
\subsection{The fold and the resonances that it
holds}\label{the-fold-and-the-resonances-that-it-holds}}

\[
\begin{aligned}
&\chi(0^+) = \exp\Big(\gamma^{(0)}\big(\mu_{\mathrm{L}}P_{\mathrm{L}} + \mu_{\mathrm{R}}P_{\mathrm{R}}\big)\Big)\,\chi(0^-),\\
&\mu_X = \frac{1}{V_*}\Big[c_\Phi\,\langle\Phi^\dagger\Phi\rangle\,\delta_{X\mathrm{L}} + \frac14\big(c^{G}_X\langle G^2\rangle
 + c^{W}_X\langle W^2\rangle + c^{B}_X\langle B^2\rangle\big) + 2\omega\sum_f c^{f}_X\,n_f\Big],\qquad
 \langle\Phi^\dagger\Phi\rangle = \frac{v^2 + \mathcal{A}^2}{2}
\end{aligned} \tag{FE-39}
\] \textbf{Symbols in (FE-39).}

\begin{itemize}
\tightlist
\item
  \(\chi(\sigma)\): The profile of a mode of dirac16complex along the
  hidden direction,
  \(\Psi = \mathrm{e}^{\frac72H\lvert\sigma\rvert}\chi(\sigma)\mathrm{e}^{\mathrm{i}(\vec k\cdot\vec x - \omega x_4)}\)
  (a 16-component spinor function of \(\sigma\)); as a four-dimensional
  field, \(\chi\) denotes the particle of the level of dirac16complex at
  the fold (a long-lived resonance) and \(\bar\chi\) its antiparticle.
  Unit: GeV\(^{2}\) (with the \(V_*\) normalisation). Status: DERIVED.
\item
  \(\gamma^M,\ \gamma^{(a)}\): The Dirac matrices of dirac16complex
  along the coordinates, \(\gamma^M = \mathbf{e}_a{}^MT^a\),
  \(\gamma_M = g_{MN}\gamma^N\); \(\gamma^{(a)} = T^a\) along the frame
  direction \((a)\) (\(\gamma^{(0)}\) along \(\sigma\), \(\gamma^{(4)}\)
  along the time). Unit: dimensionless. Status: AUTHOR.
\item
  \(\mu_X,\ \mu_{\mathrm{L}},\ \mu_{\mathrm{R}}\): The dimensionless
  strengths of the fold for the left-handed and right-handed halves of
  dirac16complex, set by the mean fields of the Standard Model at the
  fold: the jump of the profile across the fold is the exponential of
  \(\gamma^{(0)}(\mu_{\mathrm{L}}P_{\mathrm{L}} + \mu_{\mathrm{R}}P_{\mathrm{R}})\);
  negative values attract. Unit: dimensionless. Status: DERIVED.
\item
  \(P_{\mathrm{L}}\): The author's left-handed chiral projector,
  \(P_{\mathrm{L}} = (\mathbf{1} - \Gamma_8)/2\). Unit: dimensionless.
  Status: AUTHOR.
\item
  \(P_{\mathrm{R}}\): The author's right-handed chiral projector,
  \(P_{\mathrm{R}} = (\mathbf{1} + \Gamma_8)/2\). Unit: dimensionless.
  Status: AUTHOR.
\item
  \(X\): The label of a half of dirac16complex:
  \(X \in \{\mathrm{L}, \mathrm{R}\}\) (left-handed or right-handed
  half); a sum over \(X\) runs over both. Status: STANDARD.
\item
  \(V_*\): The physical volume of the extra times \(x_5, x_6, x_7\) on
  the fold in the late universe,
  \(V_* = V_{\mathrm{t}}(\mathrm{e}^{-a_4}b_4)^3\) (with
  \(V_{\mathrm{t}} = \int\mathrm{d}x_5\,\mathrm{d}x_6\,\mathrm{d}x_7\)
  the coordinate volume of the extra times) at a time when
  \(\mathrm{e}^{-a_4}b_4\) is constant; it converts the
  eight-dimensional normalisation into the four-dimensional one. Unit:
  length\(^3\) (GeV\(^{-3}\)). Status: ASSUMED.
\item
  \(c_\Phi\): The Higgs-portal coupling of the left-handed half of
  dirac16complex (the right-handed half has none). Unit: GeV\(^{-5}\).
  Status: INVENTED.
\item
  \(\Phi\): The Higgs doublet (two complex components, hypercharge
  \(1/2\)), a field of the fold. Unit: GeV. Status: AUTHOR.
\item
  \({}^\dagger\): Hermitian conjugation: the complex conjugate of the
  transpose of a matrix or spinor; for an operator on the space of
  states, its adjoint with respect to the scalar product of that space
  (for the quantised dirac16complex the indefinite Krein scalar
  product). Status: STANDARD.
\item
  \(\langle\cdots\rangle\): The mean (expectation) value of an operator
  in the state of the medium (the vacuum, a plasma, the Kohn-Sham
  state). Status: STANDARD.
\item
  \(\delta\): The Kronecker delta: with two indices \(\delta^M_N = 1\)
  if \(M = N\) and 0 otherwise (likewise \(\delta_{X\mathrm{L}}\) is 1
  if the label \(X\) of a half of dirac16complex is L and 0 otherwise);
  with two index strings of equal length the generalised Kronecker delta
  \(\delta^{M_1M_2\cdots}_{N_1N_2\cdots}\), the determinant of the
  square matrix of the Kronecker deltas \(\delta^{M_a}_{N_b}\) (\(+1\)
  or \(-1\) if the upper indices are an even or odd permutation of the
  lower ones, 0 otherwise). Status: STANDARD.
\item
  \(c^{G}_X\): The gluon-kinetic portal coupling of the half \(X\) of
  dirac16complex. Unit: GeV\(^{-7}\). Status: INVENTED.
\item
  \(\langle G^2\rangle\): The mean square of the gluon field strength,
  \(\langle G^A_{\mu\nu}G^{A\mu\nu}\rangle\) (the gluon condensate in
  the QCD vacuum, a thermal value in the plasma). Unit: GeV\(^4\).
  Status: STANDARD.
\item
  \(G^A_M,\ G^A_\mu,\ G^A_{MN}\): The gluon field of \(SU(3)_C\)
  (\(A = 1, \dots, 8\)): in the bulk \(G^A_M\), a field of eight
  dimensions in its geometric normalisation (the covariant derivative
  \(\partial - \mathrm{i}G\)), and on the fold the canonically
  normalised four-dimensional field \(G^A_\mu(x)\) of its massless mode,
  \(g_3G^A_\mu(x) = G^A_\mu(\sigma = 0, x)\); \(G^A_{MN}\) and
  \(G^A_{\mu\nu}\) their field strengths. Unit: GeV (field), GeV\(^2\)
  (field strength). Status: AUTHOR.
\item
  \(c^{W}_X\): The \(W\)-kinetic portal coupling of the half \(X\).
  Unit: GeV\(^{-7}\). Status: INVENTED.
\item
  \(\langle W^2\rangle\): The mean square of the \(SU(2)\) field
  strength, \(\langle W^I_{\mu\nu}W^{I\mu\nu}\rangle\). Unit: GeV\(^4\).
  Status: STANDARD.
\item
  \(W^I_M,\ W^I_\mu,\ W^I_{MN}\): The gauge field of \(SU(2)_L\)
  (\(I = 1, 2, 3\)): in the bulk \(W^I_M\) (geometric normalisation), on
  the fold the canonically normalised four-dimensional field
  \(W^I_\mu(x)\) of its massless mode,
  \(g_2W^I_\mu(x) = W^I_\mu(\sigma = 0, x)\); \(W^I_{MN}\),
  \(W^I_{\mu\nu}\) their field strengths;
  \(W^\pm_\mu = (W^1_\mu \mp \mathrm{i}W^2_\mu)/\sqrt2\). Unit: GeV
  (field), GeV\(^2\) (field strength). Status: AUTHOR.
\item
  \(c^{B}_X\): The \(B\)-kinetic portal coupling of the half \(X\).
  Unit: GeV\(^{-7}\). Status: INVENTED.
\item
  \(\langle B^2\rangle\): The mean square of the hypercharge field
  strength, \(\langle B_{\mu\nu}B^{\mu\nu}\rangle\). Unit: GeV\(^4\).
  Status: STANDARD.
\item
  \(B_M,\ B_\mu,\ B_{MN}\): The hypercharge gauge field of \(U(1)_Y\):
  in the bulk \(B_M\) (geometric normalisation), on the fold the
  canonically normalised four-dimensional field \(B_\mu(x)\) of its
  massless mode, \(g_1B_\mu(x) = B_\mu(\sigma = 0, x)\); \(B_{MN}\),
  \(B_{\mu\nu}\) their field strengths. Unit: GeV (field), GeV\(^2\)
  (field strength). Status: AUTHOR.
\item
  \(\omega\): A frequency: the energy of a mode
  \(\propto \mathrm{e}^{-\mathrm{i}\omega x_4}\) (\(\omega_o\) of the
  orbital \(o\)). Unit: GeV. Status: STANDARD.
\item
  \(f\): The label of a kind of Standard-Model multiplet:
  \(f \in \{Q, u, d, \ell, e, \nu\}\) (quark doublets, up-quark
  singlets, down-quark singlets, lepton doublets, charged-lepton
  singlets, right-handed neutrinos); a sum over \(f\) runs over the six.
  Status: STANDARD.
\item
  \(c^{f}_X\): The current-portal coupling of the half \(X\) of
  dirac16complex with the fermion-number current of the multiplets of
  kind \(f\) (flavour blind); \(c^{\ell}_X\) and \(c^{\nu}_X\) are those
  of the lepton doublets and of the right-handed neutrinos. Unit:
  GeV\(^{-7}\). Status: INVENTED.
\item
  \(n_f\): The number density of the Standard-Model fermions of kind
  \(f\) (particles minus antiparticles, summed over generations and
  gauge components), the time component of the mean current
  \(\langle J_f^\mu\rangle\) in the rest frame of the medium. Unit:
  GeV\(^3\). Status: STANDARD.
\item
  \(v\): The vacuum expectation value of the lower component of the
  Higgs doublet, \(v^2 = -\mu^2/\lambda\). Unit: GeV. Value: 246.22 GeV
  \(= (\sqrt2\,G_{\mathrm{F}})^{-1/2}\) with the Fermi constant
  \(G_{\mathrm{F}} = 1.1663788\times10^{-5}\) GeV\(^{-2}\) of the muon
  lifetime (Particle Data Group). Status: STANDARD.
\item
  \(\mathcal{A}\): The author's ad hoc upper component of the Higgs
  vacuum, almightyAdHoc,
  \(\mathcal{A}(x_4) = v_{\mathcal{A}}\mathrm{e}^{-2a_{40}x_4}\). Unit:
  GeV. Status: AUTHOR.
\end{itemize}

The matching condition FE-13 for a mode, with the Standard-Model fields
replaced by their mean values in a homogeneous, isotropic medium at rest
on the fold (no mean field strength, no mean spatial current; the factor
\(\mathrm{e}^{(7/2)H\lvert\sigma\rvert}\) is 1 at the fold, so \(\chi\)
and \(\Psi\) have the same jump): the fold acts on dirac16complex as a
chirality-dependent mass concentrated at \(\sigma = 0\), with the
dimensionless strengths \(\mu_{\mathrm{L}}\) and \(\mu_{\mathrm{R}}\).
The Higgs portal acts on the left-handed half only
(\(\delta_{X\mathrm{L}} = 1\) for \(X = \mathrm{L}\), 0 for
\(X = \mathrm{R}\)); the gauge-kinetic portals act through the mean
squares \(\langle G^2\rangle = \langle G^A_{\mu\nu}G^{A\mu\nu}\rangle\)
(the gluon condensate of the QCD vacuum, or of the plasma),
\(\langle W^2\rangle\), \(\langle B^2\rangle\) of the canonical fields
of the fold (IN-13); the current portals act through the number
densities \(n_f\) of the Standard-Model fermions of kind \(f\) (they
give an energy-dependent strength). In the vacuum of the late universe
only the Higgs term and the gluon condensate contribute.
\((\gamma^{(0)}(\mu_{\mathrm{L}}P_{\mathrm{L}} + \mu_{\mathrm{R}}P_{\mathrm{R}}))^2 = \mu_{\mathrm{L}}\mu_{\mathrm{R}}\mathbf{1}\),
so the exponential is
\(\cosh\varsigma + (\sinh\varsigma/\varsigma)\gamma^{(0)}(\mu_{\mathrm{L}}P_{\mathrm{L}} + \mu_{\mathrm{R}}P_{\mathrm{R}})\)
with \(\varsigma^2 = \mu_{\mathrm{L}}\mu_{\mathrm{R}}\); for
\(\mu_{\mathrm{R}} = 0\) it is exactly
\(1 + \mu_{\mathrm{L}}\gamma^{(0)}P_{\mathrm{L}}\). Status: DERIVED (as
in UM-3; \texttt{um4\_spectra.py}, U3, uses it with the author's
matrices).

\[
\begin{aligned}
&q_{\mathrm{b}} = -m_\Psi\,\big(\mu_{\mathrm{L}} + \mu_{\mathrm{R}}\big)\,\frac{\tanh\varsigma}{2\varsigma},\qquad \varsigma = \sqrt{\mu_{\mathrm{L}}\mu_{\mathrm{R}}},\qquad
\omega_{\mathrm{b}} = \sqrt{m_\Psi^2 - q_{\mathrm{b}}^2},\\
&\omega_{\mathrm{r}} = \omega_{\mathrm{b}}\Big(1 - \frac{H}{2q_{\mathrm{b}}}\Big) + \dots,\qquad
\lvert\chi_{\mathrm{b}}(\sigma)\rvert \propto \mathrm{e}^{-q_{\mathrm{b}}\lvert\sigma\rvert}\quad\big(\lvert\sigma\rvert \ll \sigma_{\mathrm{t}}\big)
\end{aligned} \tag{FE-40}
\] \textbf{Symbols in (FE-40).}

\begin{itemize}
\tightlist
\item
  \(q_{\mathrm{b}}\): The inverse width of a state of dirac16complex at
  the fold: inside the turning point its profile falls as
  \(\mathrm{e}^{-q_{\mathrm{b}}\lvert\sigma\rvert}\). Unit: GeV. Status:
  DERIVED.
\item
  \(m_\Psi\): The mass parameter of dirac16complex in the bulk. Unit:
  GeV. Status: RECORD.
\item
  \(\mu_X,\ \mu_{\mathrm{L}},\ \mu_{\mathrm{R}}\): The dimensionless
  strengths of the fold for the left-handed and right-handed halves of
  dirac16complex, set by the mean fields of the Standard Model at the
  fold: the jump of the profile across the fold is the exponential of
  \(\gamma^{(0)}(\mu_{\mathrm{L}}P_{\mathrm{L}} + \mu_{\mathrm{R}}P_{\mathrm{R}})\);
  negative values attract. Unit: dimensionless. Status: DERIVED.
\item
  \(\varsigma\): The auxiliary variable of the bound-state condition,
  \(\varsigma = \sqrt{\mu_{\mathrm{L}}\mu_{\mathrm{R}}}\)
  (\(\varsigma_- = \sqrt{-\mu_{\mathrm{L}}\mu_{\mathrm{R}}}\) when
  \(\mu_{\mathrm{L}}\mu_{\mathrm{R}} < 0\)). Unit: dimensionless.
  Status: DERIVED.
\item
  \(\omega_{\mathrm{b}}\): The energy at zero momentum of the state of
  dirac16complex at the fold in the limit \(H \to 0\) (the bound state
  of UM-3): the particle \(\chi\) of UM-4 has the mass
  \(\omega_{\mathrm{r}} = \omega_{\mathrm{b}}(1 - H/(2q_{\mathrm{b}})) + \dots\),
  equal to \(\omega_{\mathrm{b}}\) up to the relative correction
  \(H/(2q_{\mathrm{b}})\). Unit: GeV. Status: DERIVED.
\item
  \(\omega_{\mathrm{r}}\): The mass of the particle \(\chi\): the real
  part of the energy at rest of the resonance of dirac16complex at the
  fold,
  \(\omega_{\mathrm{r}} = \omega_{\mathrm{b}}(1 - H/(2q_{\mathrm{b}})) + \dots\).
  Unit: GeV. Status: DERIVED.
\item
  \(H\): The author's constant (an inverse length) that fixes the period
  \(\pi/(6H)\) of \(x_0\), the warp factor
  \(\mathrm{e}^{-2H\lvert\sigma\rvert}\) of the metric, the cosmological
  constant \(\Lambda = -21H^2\) and the tension
  \(\kappa\lambda_{\mathrm{F}} = 12H\) of the fold; in UM-4 a free
  parameter, bounded from below by the propagation of the gauge fields
  in the bulk. Unit: GeV (1/length). Value: \(H > 2534\) GeV, about 2.5
  TeV (from the invisible width of the Z boson). Status: AUTHOR.
\item
  \(\chi_{\mathrm{b}}\): The profile of a state of dirac16complex at the
  fold (a resonance): it decays as
  \(\mathrm{e}^{-q_{\mathrm{b}}\lvert\sigma\rvert}\) out to the turning
  point and continues beyond it as an outgoing wave of exponentially
  small amplitude. Unit: GeV\(^{2}\). Status: DERIVED.
\item
  \(\mathrm{e}\): Euler's number, the base of the natural logarithm;
  \(\mathrm{e}^{X}\) is the exponential of \(X\). Unit: dimensionless.
  Value: 2.718281828. Status: STANDARD.
\item
  \(\sigma\): The proper coordinate of the hidden direction:
  \(\mathrm{d}\sigma = \lvert\cot z\rvert\,\mathrm{d}x_0\),
  \(\sigma = \mathrm{sgn}(\pi/2 - z)\ln(\sin z)/(6H)\) on one period
  \(0 < z < \pi\); \(\sigma = 0\) is the fold \(z = \pi/2\),
  \(\sigma \to \mp\infty\) the tips \(z \to 0, \pi\); one period of
  \(x_0\) is the whole line \(-\infty < \sigma < \infty\). Unit: length.
  Status: DERIVED.
\item
  \(\sigma_{\mathrm{t}}\): The turning point of a state of
  dirac16complex of the frequency \(\omega\): the depth at which the
  warped frequency equals the bulk mass,
  \(\lvert\omega\rvert\mathrm{e}^{H\sigma_{\mathrm{t}}} = m_\Psi\);
  beyond it every solution oscillates. Unit: length. Status: DERIVED.
\end{itemize}

The states of dirac16complex at the fold at zero momentum. Close to the
fold (\(\lvert\sigma\rvert \ll 1/H\)) the warped frequency is \(\omega\)
itself, the profile equation is that of UM-3 with constant coefficients,
and the decaying solutions on the two sides, joined by FE-39, are the
bound states of UM-3: they exist if the fold is attractive,
\(\mu_{\mathrm{L}} + \mu_{\mathrm{R}} < 0\), with the inverse width
\(0 < q_{\mathrm{b}} < m_\Psi\) and the energy \(\omega_{\mathrm{b}}\)
(for \(\mu_{\mathrm{L}}\mu_{\mathrm{R}} < 0\),
\(\tanh(\varsigma)/\varsigma\) means \(\tan(\varsigma_-)/\varsigma_-\)
with \(\varsigma_- = \sqrt{-\mu_{\mathrm{L}}\mu_{\mathrm{R}}}\); for the
Higgs portal alone, \(\mu_{\mathrm{R}} = 0\):
\(q_{\mathrm{b}} = -m_\Psi\mu_{\mathrm{L}}/2\),
\(\omega_{\mathrm{b}} = m_\Psi\sqrt{1 - \mu_{\mathrm{L}}^2/4}\)). In
UM-4 the profile \(\chi_{\mathrm{b}}\) decays at the local rate
\(q(\sigma)\) of FE-37, which equals \(q_{\mathrm{b}}\) close to the
fold and falls to 0 at the turning point \(\sigma_{\mathrm{t}}\), so
that its amplitude at the turning point is
\(\exp(-F(\upsilon_{\mathrm{b}})m_\Psi/H)\) times that at the fold
(FE-86; not \(\mathrm{e}^{-q_{\mathrm{b}}\sigma_{\mathrm{t}}}\), which
is smaller), and beyond it continues as an outgoing wave of
exponentially small amplitude: the state is a resonance, the particle
\(\chi\) of dirac16complex, with the four-dimensional mass
\(\omega_{\mathrm{r}}\) (its ``\(x_0\) mass'') and the width of FE-86.
The warped frequency lowers the energy: to the first order in \(H\) by
\(\omega_{\mathrm{b}}H\langle\lvert\sigma\rvert\rangle\), with
\(\langle\lvert\sigma\rvert\rangle = 1/(2q_{\mathrm{b}})\) the mean
distance from the fold of the density
\(\mathrm{e}^{-2q_{\mathrm{b}}\lvert\sigma\rvert}\). The resonances are
degenerate eight times at \(+\omega_{\mathrm{r}}\) (and eight times at
\(-\omega_{\mathrm{r}}\)). Status: DERIVED (\texttt{um4\_spectra.py},
U3: the zeros of the \(16 \times 16\) matching with outgoing waves on
both sides, for \((\mu_{\mathrm{L}}, \mu_{\mathrm{R}}) = (-1, 0)\) and
\((-0.6, -0.3)\) and \(m_\Psi/H = 10\) to 40, an eight-dimensional
kernel and \(\omega_{\mathrm{r}} \to \omega_{\mathrm{b}}\) as
\(H/m_\Psi \to 0\); \texttt{um4\_kohnsham.py}, K1:
\(\omega_{\mathrm{r}} = \omega_{\mathrm{b}}(1 - H/(2q_{\mathrm{b}}))\)
up to \(-27H^2/m_\Psi\) at \(\mu_{\mathrm{L}} = -0.6\),
\(H = 10^{-4}m_\Psi\), and the box method against the exact resonance to
\(2\times10^{-9}\); K9: the same for the weak binding
\(q_{\mathrm{b}} = 0.03\,m_\Psi\)).

\[
\begin{aligned}
&\omega_0 = \omega_{\mathrm{r}} - \frac{\mathrm{i}}{2}\,\Gamma_\chi,\qquad
\Gamma_\chi = p_{\mathrm{w}}\,m_\Psi\exp\Big(-2F(\upsilon_{\mathrm{b}})\,\frac{m_\Psi}{H}\Big),\qquad
p_{\mathrm{w}} \simeq \frac{2}{\mathrm{e}}\,\frac{q_{\mathrm{b}}^2}{m_\Psi^2},\\
&F(\upsilon) = \ln\frac{1 + \sqrt{1 - \upsilon^2}}{\upsilon} - \sqrt{1 - \upsilon^2},\qquad \upsilon_{\mathrm{b}} = \frac{\omega_{\mathrm{b}}}{m_\Psi}
\end{aligned} \tag{FE-86}
\] \textbf{Symbols in (FE-86).}

\begin{itemize}
\tightlist
\item
  \(\omega_0\): The energy at rest (zero momentum along ordinary space
  and along the extra times) of a state of dirac16complex at the fold;
  complex for a resonance,
  \(\omega_0 = \omega_{\mathrm{r}} - \frac{\mathrm{i}}{2}\Gamma_\chi\).
  Unit: GeV. Status: DERIVED.
\item
  \(\omega_{\mathrm{r}}\): The mass of the particle \(\chi\): the real
  part of the energy at rest of the resonance of dirac16complex at the
  fold,
  \(\omega_{\mathrm{r}} = \omega_{\mathrm{b}}(1 - H/(2q_{\mathrm{b}})) + \dots\).
  Unit: GeV. Status: DERIVED.
\item
  \(\mathrm{i}\): The imaginary unit, \(\mathrm{i}^2 = -1\). Unit:
  dimensionless. Status: STANDARD.
\item
  \(\Gamma_\chi\): The width of the particle \(\chi\) at rest: the rate
  of its decay into the waves of dirac16complex in the bulk. Unit: GeV.
  Value: below \(1/(11.6\,t_0) = 1.3\times10^{-43}\) GeV for dark matter
  (a lifetime above 160 Gyr). Status: DERIVED.
\item
  \(p_{\mathrm{w}}\): The prefactor of the width of the particles
  \(\chi\),
  \(\Gamma_\chi = p_{\mathrm{w}}m_\Psi\mathrm{e}^{-2F(\upsilon_{\mathrm{b}})m_\Psi/H}\):
  \(p_{\mathrm{w}} \simeq (2/\mathrm{e})(q_{\mathrm{b}}/m_\Psi)^2\) to
  the leading order in \(H/q_{\mathrm{b}}\) (FE-86; 0.184 for
  \(q_{\mathrm{b}} = m_\Psi/2\)). Unit: dimensionless. Status: DERIVED.
\item
  \(m_\Psi\): The mass parameter of dirac16complex in the bulk. Unit:
  GeV. Status: RECORD.
\item
  \(\mathrm{e}\): Euler's number, the base of the natural logarithm;
  \(\mathrm{e}^{X}\) is the exponential of \(X\). Unit: dimensionless.
  Value: 2.718281828. Status: STANDARD.
\item
  \(F(\upsilon)\): The tunnelling exponent of the width of the particles
  \(\chi\) per unit of \(m_\Psi/H\): the integral of the local decay
  rate \(q\) over the region inside the turning point, divided by
  \(m_\Psi/H\), as a function of the ratio \(\upsilon\) of the energy to
  the bulk mass. Unit: dimensionless. Value: 0.049 for
  \(\upsilon = 0.866\); \((q_{\mathrm{b}}/\omega_{\mathrm{b}})^3/3\) for
  weak binding. Status: DERIVED.
\item
  \(\upsilon,\ \upsilon_{\mathrm{b}},\ \upsilon_{\mathrm{r}}\): The
  argument of the tunnelling exponent \(F\), a ratio of an energy to the
  bulk mass \(m_\Psi\);
  \(\upsilon_{\mathrm{b}} = \omega_{\mathrm{b}}/m_\Psi\) that of the
  particle \(\chi\) in the limit \(H/q_{\mathrm{b}} \to 0\),
  \(\upsilon_{\mathrm{r}} = \omega_{\mathrm{r}}/m_\Psi\) that of the
  real part of its resonance. Unit: dimensionless. Status: DERIVED.
\item
  \(\omega_{\mathrm{b}}\): The energy at zero momentum of the state of
  dirac16complex at the fold in the limit \(H \to 0\) (the bound state
  of UM-3): the particle \(\chi\) of UM-4 has the mass
  \(\omega_{\mathrm{r}} = \omega_{\mathrm{b}}(1 - H/(2q_{\mathrm{b}})) + \dots\),
  equal to \(\omega_{\mathrm{b}}\) up to the relative correction
  \(H/(2q_{\mathrm{b}})\). Unit: GeV. Status: DERIVED.
\item
  \(H\): The author's constant (an inverse length) that fixes the period
  \(\pi/(6H)\) of \(x_0\), the warp factor
  \(\mathrm{e}^{-2H\lvert\sigma\rvert}\) of the metric, the cosmological
  constant \(\Lambda = -21H^2\) and the tension
  \(\kappa\lambda_{\mathrm{F}} = 12H\) of the fold; in UM-4 a free
  parameter, bounded from below by the propagation of the gauge fields
  in the bulk. Unit: GeV (1/length). Value: \(H > 2534\) GeV, about 2.5
  TeV (from the invisible width of the Z boson). Status: AUTHOR.
\item
  \(q_{\mathrm{b}}\): The inverse width of a state of dirac16complex at
  the fold: inside the turning point its profile falls as
  \(\mathrm{e}^{-q_{\mathrm{b}}\lvert\sigma\rvert}\). Unit: GeV. Status:
  DERIVED.
\end{itemize}

The resonances of dirac16complex at the fold: the solutions of FE-37
with outgoing waves on both sides of the fold (the Hankel functions of
FE-85) that satisfy the matching FE-39 exist only at the complex
energies \(\omega_0\), with the real part \(\omega_{\mathrm{r}}\) of
FE-40 and the width \(\Gamma_\chi\), the rate at which a particle
\(\chi\) at rest decays into the waves of the bulk. The width is the
tunnelling through the region
\(\lvert\sigma\rvert < \sigma_{\mathrm{t}}\) in which the warped
frequency is below the bulk mass:
\(2\int_0^{\sigma_{\mathrm{t}}}q\,\mathrm{d}\sigma = 2F(\upsilon_{\mathrm{b}})\,m_\Psi/H\)
(\(F(\upsilon) = \int_{\upsilon}^1\sqrt{1 - t^2}\,\mathrm{d}t/t\), with
\(t = \omega\mathrm{e}^{H\sigma}/m_\Psi\)), so the particles are
long-lived only if \(m_\Psi/H\) is large. The width with its prefactor
follows from the exact solutions FE-85 for large \(m_\Psi/H\): by the
uniform (Debye) asymptotic forms of the Hankel functions of the orders
\(m_\Psi/H \pm \frac12\), the ratio of the two outgoing solutions at the
fold is real up to an imaginary part of the relative size
\(\exp(-2F(\upsilon)m_\Psi/H)\,q_{\mathrm{b}}/\omega_{\mathrm{b}}\), and
the matching FE-39 turns this imaginary part into the width
\(\Gamma_\chi = 2(q_{\mathrm{b}}^2/m_\Psi)\exp(-2F(\upsilon_{\mathrm{r}})m_\Psi/H)\)
at the real part \(\upsilon_{\mathrm{r}} = \omega_{\mathrm{r}}/m_\Psi\)
of the resonance; the shift
\(\omega_{\mathrm{r}} = \omega_{\mathrm{b}}(1 - H/(2q_{\mathrm{b}}))\)
of FE-40 raises the exponent by exactly 1
(\(F(\upsilon_{\mathrm{r}}) = F(\upsilon_{\mathrm{b}}) + H/(2m_\Psi)\)),
the factor \(1/\mathrm{e}\) of the prefactor \(p_{\mathrm{w}}\), to the
leading order in \(H/q_{\mathrm{b}}\). The prefactor depends on the
binding: \(p_{\mathrm{w}} = 0.184\) for \(q_{\mathrm{b}} = m_\Psi/2\),
0.133 for \((\mu_{\mathrm{L}}, \mu_{\mathrm{R}}) = (-0.6, -0.3)\), 0.721
for \(q_{\mathrm{b}} = 0.99\,m_\Psi\). Particles \(\chi\) that make up
the dark matter today must live longer than 160 Gyr, 11.6 times the age
\(t_0 = 13.8\) Gyr of the universe (the bound on the lifetime of dark
matter that decays into invisible relativistic particles from the cosmic
microwave background and the large-scale structure: B. Audren, J.
Lesgourgues, G. Mangano, P. D. Serpico and T. Tram, JCAP 12 (2014) 028,
arXiv:1407.2418), \(\Gamma_\chi < 1.3\times10^{-43}\) GeV. With \(H\) at
its lower bound of 2534 GeV (FE-82) this requires \(m_\Psi/H\) and
\(\omega_{\mathrm{b}}\) above the values of the table, which fall with
the binding:

\begin{longtable}[]{@{}
  >{\raggedright\arraybackslash}p{(\columnwidth - 18\tabcolsep) * \real{0.1000}}
  >{\raggedright\arraybackslash}p{(\columnwidth - 18\tabcolsep) * \real{0.1000}}
  >{\raggedright\arraybackslash}p{(\columnwidth - 18\tabcolsep) * \real{0.1000}}
  >{\raggedright\arraybackslash}p{(\columnwidth - 18\tabcolsep) * \real{0.1000}}
  >{\raggedright\arraybackslash}p{(\columnwidth - 18\tabcolsep) * \real{0.1000}}
  >{\raggedright\arraybackslash}p{(\columnwidth - 18\tabcolsep) * \real{0.1000}}
  >{\raggedright\arraybackslash}p{(\columnwidth - 18\tabcolsep) * \real{0.1000}}
  >{\raggedright\arraybackslash}p{(\columnwidth - 18\tabcolsep) * \real{0.1000}}
  >{\raggedright\arraybackslash}p{(\columnwidth - 18\tabcolsep) * \real{0.1000}}
  >{\raggedright\arraybackslash}p{(\columnwidth - 18\tabcolsep) * \real{0.1000}}@{}}
\toprule\noalign{}
\begin{minipage}[b]{\linewidth}\raggedright
\(q_{\mathrm{b}}/m_\Psi\)
\end{minipage} & \begin{minipage}[b]{\linewidth}\raggedright
0.0995
\end{minipage} & \begin{minipage}[b]{\linewidth}\raggedright
0.3
\end{minipage} & \begin{minipage}[b]{\linewidth}\raggedright
0.5
\end{minipage} & \begin{minipage}[b]{\linewidth}\raggedright
0.7
\end{minipage} & \begin{minipage}[b]{\linewidth}\raggedright
0.8
\end{minipage} & \begin{minipage}[b]{\linewidth}\raggedright
0.9
\end{minipage} & \begin{minipage}[b]{\linewidth}\raggedright
0.95
\end{minipage} & \begin{minipage}[b]{\linewidth}\raggedright
0.99
\end{minipage} & \begin{minipage}[b]{\linewidth}\raggedright
0.999
\end{minipage} \\
\midrule\noalign{}
\endhead
\bottomrule\noalign{}
\endlastfoot
\(m_\Psi/H\) & \(1.72\times10^5\) & 5911 & 1135 & 333 & 186 & 96.7 &
62.5 & 33.1 & 19.5 \\
\(\omega_{\mathrm{b}}\) (GeV) & \(4.3\times10^8\) & \(1.4\times10^7\) &
\(2.5\times10^6\) & \(6.0\times10^5\) & \(2.8\times10^5\) &
\(1.1\times10^5\) & \(4.9\times10^4\) & \(1.2\times10^4\) &
\(2.2\times10^3\) \\
\end{longtable}

(\(q_{\mathrm{b}}/m_\Psi = 0.0995\) is the weak binding
\(q_{\mathrm{b}}/\omega_{\mathrm{b}} = 0.1\), where
\(F \simeq (q_{\mathrm{b}}/\omega_{\mathrm{b}})^3/3\); for the age
\(t_0\) instead of 160 Gyr, \(q_{\mathrm{b}} = m_\Psi/2\) needs
\(m_\Psi/H > 1110\) and \(\omega_{\mathrm{b}} > 2.44\times10^6\) GeV.)
The bound on \(\omega_{\mathrm{b}}\) vanishes as
\(q_{\mathrm{b}} \to m_\Psi\): strongly bound particles \(\chi\) can be
light. For them \(H/q_{\mathrm{b}}\) is about 0.03 (at
\(q_{\mathrm{b}} = 0.99\,m_\Psi\)), and the table uses the exact widths,
whose prefactor lies within 4 per cent of
\((2/\mathrm{e})(q_{\mathrm{b}}/m_\Psi)^2\) (within 1 per cent where
\(m_\Psi/H > 1000\)). Direct detection demands a far weaker binding by
the Higgs portal (its Yukawa coupling, FE-47): the largest
\(q_{\mathrm{b}}/\omega_{\mathrm{b}}\) that the limit of LZ allows falls
like \(\omega_{\mathrm{b}}^{-1/4}\), from 0.019 at 10 TeV to
\(6.2\times10^{-3}\) at \(10^6\) GeV and \(3\times10^{-6}\) at the
Planck mass, and with that binding the particles live longer than 160
Gyr only if \(H\) is below \(1.5\times10^{-3}\) GeV at \(10^6\) GeV and
below 2.5 GeV even at the Planck mass: never for \(H > 2534\) GeV. A
strong binding by the gluon condensate is excluded in the same way
(FE-61). A particle moving with the momentum \(\vec k\) has the energy
\(\sqrt{\omega_0^2 + \vec k^2}\) (FE-38) and the width
\(\Gamma_\chi\omega_{\mathrm{r}}/\sqrt{\omega_{\mathrm{r}}^2 + \vec k^2}\)
(the time dilation of the boost). Status: DERIVED (the width:
\texttt{um4\_spectra.py}, U3 and U4: the slope of \(\ln\Gamma_\chi\) in
\(m_\Psi/H\) between 30 and 40 agrees with
\(-2F(\upsilon_{\mathrm{b}})\) within 6 per cent, and the prefactor of
the exact widths approaches
\(p_{\mathrm{w}} = (2/\mathrm{e})(q_{\mathrm{b}}/m_\Psi)^2\) like
\(H/m_\Psi\): 0.18239 and 0.18356 at \(m_\Psi/H = 1000\) and 4000
against 0.18394; U6: the exact widths at the bounds of the table;
\texttt{um4\_kohnsham.py}, K8: the tunnelling exponent at the resonance
exceeds \(2F(\upsilon_{\mathrm{b}})m_\Psi/H\) by 1.002), ESTIMATE (the
bounds, U6, with the measured lifetime as input; the binding that direct
detection allows, \texttt{um4\_thermal.py}, T8 and T9).

\hypertarget{the-kohn-sham-equations}{%
\subsection{The Kohn-Sham equations}\label{the-kohn-sham-equations}}

\[
\begin{aligned}
&\big[\gamma^MD_M - m_\Psi - \lambda_\Psi\,n_{\mathrm{s}} + \lambda_\Psi\,\mathcal{N}\big]\varphi_o
 = \delta_{\mathrm{D}}(\sigma)\,\frac{1}{V_*}\,\langle\mathbb{M}\rangle\,\varphi_o,\\
&\mathcal{N} = \sum_o\mathsf{n}_o\,\varphi_o\bar\varphi_o,\qquad n_{\mathrm{s}} = \mathrm{tr}\,\mathcal{N} = \sum_o\mathsf{n}_o\,\bar\varphi_o\varphi_o,\qquad
 \mathsf{n}_o = \frac{1}{\mathrm{e}^{(E_o - \mu_\Psi)/\mathsf{T}_\Psi} + 1},\\
&\chi_o' = \gamma^{(0)}\Big[\mathrm{i}\omega\,\mathrm{e}^{H\lvert\sigma\rvert}\gamma^{(4)} + m_\Psi + \mathcal{M}(\sigma)\Big]\chi_o,\qquad
 \mathcal{M} = G_{\mathrm{KS}}\,\mathrm{e}^{7H\lvert\sigma\rvert}\big(\hat\rho\,\mathbf{1} - \hat{\mathcal{N}}\big),\qquad
 \hat\rho = \frac14\sum_o\bar\chi_o\chi_o,\qquad \hat{\mathcal{N}} = \frac14\sum_o\chi_o\bar\chi_o
\end{aligned} \tag{FE-41}
\] \textbf{Symbols in (FE-41).}

\begin{itemize}
\tightlist
\item
  \(\gamma^M,\ \gamma^{(a)}\): The Dirac matrices of dirac16complex
  along the coordinates, \(\gamma^M = \mathbf{e}_a{}^MT^a\),
  \(\gamma_M = g_{MN}\gamma^N\); \(\gamma^{(a)} = T^a\) along the frame
  direction \((a)\) (\(\gamma^{(0)}\) along \(\sigma\), \(\gamma^{(4)}\)
  along the time). Unit: dimensionless. Status: AUTHOR.
\item
  \(M\): A coordinate index of eight dimensions (also with a numbered
  copy \(M_1, M_2, \dots\)): it runs over \(x_0, x_1, \dots, x_7\); a
  repeated index is summed. Status: STANDARD.
\item
  \(D_M,\ D_\mu\): The covariant derivative: on dirac16complex
  \(D_M = \partial_M + \Omega_M\) (with
  \(D_M\bar\Psi = \partial_M\bar\Psi - \bar\Psi\Omega_M\)); on the
  fields of the fold the gauge and spin covariant derivative \(D_\mu\)
  (the partial derivative plus the spin connection of the fold and the
  gauge fields times their couplings and generators);
  \(D^\mu = \hat h^{\mu\nu}D_\nu\); on an adjoint field
  \((D_\mu X)^A = \partial_\mu X^A + g_3f^{ABC}G^B_\mu X^C\) (likewise
  for \(SU(2)\) with \(g_2\epsilon^{IJK}\)). Unit: 1/length. Status:
  STANDARD.
\item
  \(m_\Psi\): The mass parameter of dirac16complex in the bulk. Unit:
  GeV. Status: RECORD.
\item
  \(\lambda_\Psi\): The coupling of the contact self-interaction of
  dirac16complex. Unit: GeV\(^{-6}\). Status: RECORD.
\item
  \(n_{\mathrm{s}}\): The scalar density of dirac16complex in the
  Kohn-Sham state,
  \(n_{\mathrm{s}} = \langle\bar\Psi\Psi\rangle = \mathrm{tr}\,\mathcal{N}\).
  Unit: GeV\(^7\). Status: DERIVED.
\item
  \(\mathcal{N}\): The one-body density matrix of dirac16complex,
  \(\mathcal{N} = \sum_o\mathsf{n}_o\varphi_o\bar\varphi_o\) (a
  \(16 \times 16\) matrix at each point). Unit: GeV\(^7\). Status:
  DERIVED.
\item
  \(\varphi_o,\ \bar\varphi_o\): A Kohn-Sham orbital of dirac16complex
  (a single-particle solution of the Kohn-Sham equation: the Dirac
  equation of dirac16complex with the mean field of the gas added to its
  mass) and its Dirac adjoint. Unit: GeV\(^{7/2}\). Status: RECORD.
\item
  \(o\): The label of a single-particle state of dirac16complex: a
  Kohn-Sham orbital, or a mode of the quantised field; a sum over \(o\)
  runs over the orbitals or the modes. Status: STANDARD.
\item
  \(\delta_{\mathrm{D}}\): The Dirac delta function of \(\sigma\): zero
  for \(\sigma \ne 0\), with the integral 1 over any range that contains
  \(\sigma = 0\); the integral of \(\delta_{\mathrm{D}}(\sigma)\) times
  a smooth function of \(\sigma\) is the value of that function at
  \(\sigma = 0\). Unit: 1/length (for \(\delta_{\mathrm{D}}(\sigma)\)).
  Status: STANDARD.
\item
  \(\sigma\): The proper coordinate of the hidden direction:
  \(\mathrm{d}\sigma = \lvert\cot z\rvert\,\mathrm{d}x_0\),
  \(\sigma = \mathrm{sgn}(\pi/2 - z)\ln(\sin z)/(6H)\) on one period
  \(0 < z < \pi\); \(\sigma = 0\) is the fold \(z = \pi/2\),
  \(\sigma \to \mp\infty\) the tips \(z \to 0, \pi\); one period of
  \(x_0\) is the whole line \(-\infty < \sigma < \infty\). Unit: length.
  Status: DERIVED.
\item
  \(V_*\): The physical volume of the extra times \(x_5, x_6, x_7\) on
  the fold in the late universe,
  \(V_* = V_{\mathrm{t}}(\mathrm{e}^{-a_4}b_4)^3\) (with
  \(V_{\mathrm{t}} = \int\mathrm{d}x_5\,\mathrm{d}x_6\,\mathrm{d}x_7\)
  the coordinate volume of the extra times) at a time when
  \(\mathrm{e}^{-a_4}b_4\) is constant; it converts the
  eight-dimensional normalisation into the four-dimensional one. Unit:
  length\(^3\) (GeV\(^{-3}\)). Status: ASSUMED.
\item
  \(\langle\cdots\rangle\): The mean (expectation) value of an operator
  in the state of the medium (the vacuum, a plasma, the Kohn-Sham
  state). Status: STANDARD.
\item
  \(\mathbb{M}\): The fold operator of dirac16complex: the
  \(16 \times 16\) matrix (with a derivative along the fold) built from
  the Standard-Model fields at the fold by varying the portals with
  respect to \(\bar\Psi\). Unit: GeV\(^{-3}\) (so that
  \(\mathbb{M}/V_*\) is dimensionless). Status: DERIVED.
\item
  \(\mathsf{n}_o\): The occupation of the Kohn-Sham orbital \(o\)
  (between 0 and 1). Unit: dimensionless. Status: RECORD.
\item
  \(\mathrm{tr}\): The trace of a matrix: the sum of its diagonal
  entries. Status: STANDARD.
\item
  \(\mathrm{e}\): Euler's number, the base of the natural logarithm;
  \(\mathrm{e}^{X}\) is the exponential of \(X\). Unit: dimensionless.
  Value: 2.718281828. Status: STANDARD.
\item
  \(E\): An energy: \(E_o\) the energy of the Kohn-Sham orbital \(o\);
  in the neutrino sector the energy of a neutrino. Unit: GeV. Status:
  STANDARD.
\item
  \(\mu_\Psi\): The chemical potential of dirac16complex. Unit: GeV.
  Status: DERIVED.
\item
  \(\mathsf{T}_\Psi\): The temperature of dirac16complex. Unit: GeV.
  Status: DERIVED.
\item
  \(\chi_o(\sigma),\ \bar\chi_o\): The profile of the occupied Kohn-Sham
  orbital \(o\) (\(o = 1, \dots, 4\), the four positive-energy states of
  the level at the fold) of a homogeneous gas of the particles \(\chi\)
  at rest on the static block: a 16-component spinor function of
  \(\sigma\); the orbital of each particle is
  \(\mathrm{e}^{\frac72H\lvert\sigma\rvert}\chi_o(\sigma)\mathrm{e}^{-\mathrm{i}\omega x_4}\)
  divided by the square root of the volume of ordinary space and of
  \(V_*\); normalised to the Krein charge \(-1\) with the measure of
  UM-4,
  \(\int\mathrm{d}\sigma\,\mathrm{e}^{H\lvert\sigma\rvert}\chi_o^\dagger\mathfrak{B}\chi_o = -1\);
  \(\bar\chi_o = \chi_o^\dagger\sigma_{16}\) is its Dirac adjoint. Unit:
  GeV\(^{1/2}\). Status: DERIVED.
\item
  \({}'\): The derivative of a function of the time \(x_4\) alone with
  respect to \(x_4\) (\(a_4' = \mathrm{d}a_4/\mathrm{d}x_4\); two primes
  the second derivative); on a function of \(\sigma\) alone (a profile)
  the derivative with respect to \(\sigma\). Status: STANDARD.
\item
  \(\mathrm{i}\): The imaginary unit, \(\mathrm{i}^2 = -1\). Unit:
  dimensionless. Status: STANDARD.
\item
  \(\omega\): A frequency: the energy of a mode
  \(\propto \mathrm{e}^{-\mathrm{i}\omega x_4}\) (\(\omega_o\) of the
  orbital \(o\)). Unit: GeV. Status: STANDARD.
\item
  \(H\): The author's constant (an inverse length) that fixes the period
  \(\pi/(6H)\) of \(x_0\), the warp factor
  \(\mathrm{e}^{-2H\lvert\sigma\rvert}\) of the metric, the cosmological
  constant \(\Lambda = -21H^2\) and the tension
  \(\kappa\lambda_{\mathrm{F}} = 12H\) of the fold; in UM-4 a free
  parameter, bounded from below by the propagation of the gauge fields
  in the bulk. Unit: GeV (1/length). Value: \(H > 2534\) GeV, about 2.5
  TeV (from the invisible width of the Z boson). Status: AUTHOR.
\item
  \(\mathcal{M}(\sigma)\): The mean field of the self-interaction in the
  Kohn-Sham equation of the gas: the \(16 \times 16\) matrix
  \(\lambda_\Psi(n_{\mathrm{s}}\mathbf{1} - \mathcal{N})\) at each
  \(\sigma\), which adds to the mass \(m_\Psi\); positive (repulsive)
  for \(\lambda_\Psi > 0\). Unit: GeV. Status: DERIVED.
\item
  \(G_{\mathrm{KS}}\): The dimensionless strength of the
  self-interaction of the gas of particles \(\chi\) in the Kohn-Sham
  states: the self-coupling \(\lambda_\Psi\) times the number density
  \(n_\chi\) of the particles on the fold divided by the volume \(V_*\)
  of the extra times, \(G_{\mathrm{KS}} = \lambda_\Psi n_\chi/V_*\) (in
  units of \(m_\Psi\) when numbers are given). Unit: dimensionless.
  Value: \(10^{-6}\) to 0.5 in the solutions. Status: DERIVED.
\item
  \(\hat\rho(\sigma)\): The scalar density of the occupied orbitals,
  \(\hat\rho = \frac14\sum_o\bar\chi_o\chi_o\) (per particle); the
  scalar density of the gas is
  \(n_{\mathrm{s}} = (n_\chi/V_*)\,\mathrm{e}^{7H\lvert\sigma\rvert}\hat\rho\).
  Unit: GeV. Status: DERIVED.
\item
  \(\mathbf{1},\ \mathbf{1}_4,\ \mathbf{1}_8,\ \mathbf{1}_{16}\): The
  identity matrix of the size fixed by the context (\(16 \times 16\) on
  dirac16complex, \(4 \times 4\) on a four-dimensional spinor,
  \(3 \times 3\) in generation or flavour space); \(\mathbf{1}_4\),
  \(\mathbf{1}_8\) and \(\mathbf{1}_{16}\) are the identity matrices of
  4, 8 and 16 rows. Unit: dimensionless. Status: STANDARD.
\item
  \(\hat{\mathcal{N}}(\sigma)\): The one-body density matrix of the
  occupied orbitals,
  \(\hat{\mathcal{N}} = \frac14\sum_o\chi_o\bar\chi_o\) (a
  \(16 \times 16\) matrix at each \(\sigma\), per particle); the density
  matrix of the gas is
  \(\mathcal{N} = (n_\chi/V_*)\,\mathrm{e}^{7H\lvert\sigma\rvert}\hat{\mathcal{N}}\).
  Unit: GeV. Status: DERIVED.
\end{itemize}

The Kohn-Sham equations of dirac16complex coupled to itself, to gravity
and to every Standard-Model field: the quantised field is described by
single-particle orbitals \(\varphi_o\) with the occupations
\(\mathsf{n}_o\) (Fermi-Dirac with the temperature \(\mathsf{T}_\Psi\)
and the chemical potential \(\mu_\Psi\) of dirac16complex;
\(\mathsf{n}_o = 0\) or 1 in a degenerate state). Each orbital obeys the
field equation with the mean fields: gravity through \(D_M\) and
\(\gamma^M\) (the background metric and its spin connection), the
self-interaction through the Hartree term
\(\lambda_\Psi n_{\mathrm{s}}\) (the scalar density) and the exact
exchange term \(-\lambda_\Psi\mathcal{N}\) (the one-body density matrix,
a \(16 \times 16\) matrix at each point), and every Standard-Model field
through the fold operator with the mean fields of the Standard Model,
\(\langle\mathbb{M}\rangle\) (FE-13, FE-39). The mean-field
decomposition of \(-\frac{\lambda_\Psi}{2}(\bar\Psi\Psi)^2\) is
\(-\lambda_\Psi n_{\mathrm{s}}\bar\Psi\Psi + \lambda_\Psi\bar\Psi\mathcal{N}\Psi\)
(the second term with the sign of fermion exchange); correlation beyond
exchange is neglected and the negative-energy sea is not included
(no-sea approximation). The third line is the reduction to the hidden
direction for a homogeneous gas of the particles \(\chi\) at rest on the
static block: the four occupied orbitals \(\chi_o\) of the level of
FE-40 (Krein charge \(-1\) each, normalised with the measure of FE-43,
\(\int\mathrm{d}\sigma\,\mathrm{e}^{H\lvert\sigma\rvert}\chi_o^\dagger\mathfrak{B}\chi_o = -1\),
with \(\Psi = \mathrm{e}^{(7/2)H\lvert\sigma\rvert}\chi\) as in FE-36),
the mean field
\(\mathcal{M}(\sigma) = \lambda_\Psi(n_{\mathrm{s}} - \mathcal{N})\)
written with the scalar density \(\hat\rho\) and the density matrix
\(\hat{\mathcal{N}}\) per particle and the dimensionless strength
\(G_{\mathrm{KS}} = \lambda_\Psi n_\chi/V_*\) (\(n_\chi\) the number
density of the particles on the fold); the proper densities carry the
factor \(\mathrm{e}^{7H\lvert\sigma\rvert}\) of
\(\bar\Psi\Psi = \mathrm{e}^{7H\lvert\sigma\rvert}\bar\chi\chi\) (UM-3:
\(\mathrm{e}^{6H\lvert\sigma\rvert}\)), and the frequency carries the
warp \(\mathrm{e}^{H\lvert\sigma\rvert}\). The four occupied orbitals
are taken to form one multiplet of the rotations of ordinary space and
of the extra times (the symmetric occupation of the level): the
first-order energy is stationary there, but it is a saddle on the
32-dimensional manifold of the occupations of negative Krein norm and
has no lower bound on it (FE-67, \texttt{um4\_krein.py}, KR8), so the
symmetric occupation is ASSUMED, not a ground state. The solutions
(DERIVED, \texttt{um4\_kohnsham.py}, K1 to K10, for
\(\mu_{\mathrm{L}} = -0.6\), \(\mu_{\mathrm{R}} = 0\),
\(H = 10^{-4}m_\Psi\), solved in the box
\(\lvert\sigma\rvert \le 40/m_\Psi\) well inside the turning point at
\(473/m_\Psi\), with the author's matrices, reduced exactly by the
rotations of ordinary space and of the extra times and lifted back to
sixteen components for the checks): the self-consistent shift of the
level equals first-order perturbation theory in the dilute gas to
\(5\times10^{-4}\) for \(G_{\mathrm{KS}} \le 10^{-3}\),
\(\delta\omega \propto G_{\mathrm{KS}}\), and the occupied and the empty
multiplets split in proportion to \(G_{\mathrm{KS}}\) at every
\(G_{\mathrm{KS}}\); the residuals of FE-41 and of the matching are at
most \(10^{-4}\) for \(G_{\mathrm{KS}} = 10^{-6}\) to 0.5; within the
box the proper density decays away from the fold (also for an attractive
self-interaction; a property of the box, in which the weight
\(\mathrm{e}^{7H\lvert\sigma\rvert}\) differs from 1 by at most 0.028);
the Higgs-portal loop through FE-49 converges and lowers the level; the
pressure of the gas is positive and vanishes in the dilute limit (cold
dark matter); and the repulsive self-interaction raises the level
towards \(m_\Psi\), which shrinks the barrier of FE-86, so that at
\(m_\Psi/H = 10^4\) with \(H\) at its bound (FE-82) the gas lives longer
than 160 Gyr only for \(G_{\mathrm{KS}} < 0.139\). Status: RECORD (the
Kohn-Sham treatment), DERIVED (the exchange term, the reduction and the
solutions), ASSUMED (no correlation, no sea; the symmetric occupation).

\[
\langle{\mathcal{T}_\Psi}_{MN}\rangle = -\frac14\sum_o\mathsf{n}_o\Big(\bar\varphi_o\gamma_MD_N\varphi_o + \bar\varphi_o\gamma_ND_M\varphi_o
 - (D_M\bar\varphi_o)\gamma_N\varphi_o - (D_N\bar\varphi_o)\gamma_M\varphi_o\Big) + \frac{\lambda_\Psi}{2}\,g_{MN}\big(n_{\mathrm{s}}^2 - \mathrm{tr}\,\mathcal{N}^2\big) \tag{FE-42}
\] \textbf{Symbols in (FE-42).}

\begin{itemize}
\tightlist
\item
  \(\langle\cdots\rangle\): The mean (expectation) value of an operator
  in the state of the medium (the vacuum, a plasma, the Kohn-Sham
  state). Status: STANDARD.
\item
  \(\mathcal{T}_\Psi\): The energy-momentum tensor of dirac16complex:
  \(-2/\sqrt{\lvert g\rvert}\) times the variation of the action
  \(S_\Psi\) of dirac16complex with respect to \(g^{MN}\). Unit:
  GeV\(^8\). Status: STANDARD.
\item
  \(M\): A coordinate index of eight dimensions (also with a numbered
  copy \(M_1, M_2, \dots\)): it runs over \(x_0, x_1, \dots, x_7\); a
  repeated index is summed. Status: STANDARD.
\item
  \(N\): A coordinate index of eight dimensions (also
  \(N_1, N_2, \dots\)), running over \(x_0, \dots, x_7\); a repeated
  index is summed. Status: STANDARD.
\item
  \(o\): The label of a single-particle state of dirac16complex: a
  Kohn-Sham orbital, or a mode of the quantised field; a sum over \(o\)
  runs over the orbitals or the modes. Status: STANDARD.
\item
  \(\mathsf{n}_o\): The occupation of the Kohn-Sham orbital \(o\)
  (between 0 and 1). Unit: dimensionless. Status: RECORD.
\item
  \(\varphi_o,\ \bar\varphi_o\): A Kohn-Sham orbital of dirac16complex
  (a single-particle solution of the Kohn-Sham equation: the Dirac
  equation of dirac16complex with the mean field of the gas added to its
  mass) and its Dirac adjoint. Unit: GeV\(^{7/2}\). Status: RECORD.
\item
  \(\gamma^M,\ \gamma^{(a)}\): The Dirac matrices of dirac16complex
  along the coordinates, \(\gamma^M = \mathbf{e}_a{}^MT^a\),
  \(\gamma_M = g_{MN}\gamma^N\); \(\gamma^{(a)} = T^a\) along the frame
  direction \((a)\) (\(\gamma^{(0)}\) along \(\sigma\), \(\gamma^{(4)}\)
  along the time). Unit: dimensionless. Status: AUTHOR.
\item
  \(D_M,\ D_\mu\): The covariant derivative: on dirac16complex
  \(D_M = \partial_M + \Omega_M\) (with
  \(D_M\bar\Psi = \partial_M\bar\Psi - \bar\Psi\Omega_M\)); on the
  fields of the fold the gauge and spin covariant derivative \(D_\mu\)
  (the partial derivative plus the spin connection of the fold and the
  gauge fields times their couplings and generators);
  \(D^\mu = \hat h^{\mu\nu}D_\nu\); on an adjoint field
  \((D_\mu X)^A = \partial_\mu X^A + g_3f^{ABC}G^B_\mu X^C\) (likewise
  for \(SU(2)\) with \(g_2\epsilon^{IJK}\)). Unit: 1/length. Status:
  STANDARD.
\item
  \(\lambda_\Psi\): The coupling of the contact self-interaction of
  dirac16complex. Unit: GeV\(^{-6}\). Status: RECORD.
\item
  \(g_{MN}\): The metric tensor of the eight-dimensional spacetime
  (signature \(+,+,+,+,-,-,-,-\) in the order \(x_0, \dots, x_7\)):
  \(g_{MN}\,\mathrm{d}X^M\mathrm{d}X^N\) is the squared interval,
  \(g^{MN}\) its inverse, \(\lvert g\rvert\) the absolute value of its
  determinant (\(\sqrt{\lvert g\rvert}\,\mathrm{d}^8X\) is the proper
  volume element). Unit: dimensionless. Status: AUTHOR.
\item
  \(n_{\mathrm{s}}\): The scalar density of dirac16complex in the
  Kohn-Sham state,
  \(n_{\mathrm{s}} = \langle\bar\Psi\Psi\rangle = \mathrm{tr}\,\mathcal{N}\).
  Unit: GeV\(^7\). Status: DERIVED.
\item
  \(\mathrm{tr}\): The trace of a matrix: the sum of its diagonal
  entries. Status: STANDARD.
\item
  \(\mathcal{N}\): The one-body density matrix of dirac16complex,
  \(\mathcal{N} = \sum_o\mathsf{n}_o\varphi_o\bar\varphi_o\) (a
  \(16 \times 16\) matrix at each point). Unit: GeV\(^7\). Status:
  DERIVED.
\end{itemize}

The energy-momentum tensor of the Kohn-Sham state, the source of
dirac16complex in FE-1 and, through the portals, in FE-5: the sum over
the occupied orbitals of FE-6, and the mean value of the
self-interaction (on solutions of FE-41,
\(\langle\mathcal{L}_\Psi\rangle = \frac{\lambda_\Psi}{2}(n_{\mathrm{s}}^2 - \mathrm{tr}\,\mathcal{N}^2)\)).
It closes the loop: the orbitals determine the densities, the densities
the mean fields and the metric, the metric and the mean fields the
orbitals. For the gas of FE-41 the energy per particle is the orbital
energy minus the interaction energy per particle, and the pressure along
ordinary space, equal to that along the extra times, is the interaction
energy per particle times \(n_\chi\): the equation of state is
\(w_\chi = 5.4\times10^{-8}\) to \(5.0\times10^{-3}\) for
\(G_{\mathrm{KS}} = 10^{-6}\) to 0.1 (DERIVED,
\texttt{um4\_kohnsham.py}, K7). Status: DERIVED.

\[
\begin{aligned}
&Q_\Psi = \int\mathrm{d}\sigma\,\mathrm{d}^3x\,\sqrt{\lvert g\rvert}\;\mathrm{e}^{H\lvert\sigma\rvert}\,\Psi^\dagger\mathfrak{B}\Psi
 = b_4^3\int\mathrm{d}\sigma\,\mathrm{d}^3x\;\mathrm{e}^{H\lvert\sigma\rvert}\,\chi^\dagger\mathfrak{B}\chi,\\
&\mathfrak{B} = \mathrm{i}\,\sigma_{16}\,\gamma^{(4)},\qquad \mathfrak{B}^\dagger = \mathfrak{B},\qquad \mathfrak{B}^2 = \mathbf{1},\qquad
E_o = -\omega_o\,\mathcal{Q}_o
\end{aligned} \tag{FE-43}
\] \textbf{Symbols in (FE-43).}

\begin{itemize}
\tightlist
\item
  \(Q_\Psi\): The conserved charge of dirac16complex (the Noether charge
  of a change of the phase of \(\Psi\) by a constant), a Krein form
  (indefinite); as written it is the charge per unit coordinate volume
  of the extra times, on which no field depends. Unit: GeV\(^3\).
  Status: DERIVED.
\item
  \(\mathrm{d}\): The differential: \(\mathrm{d}\) in front of a
  quantity is an infinitesimal change of it (as \(\mathrm{d}\sigma\), a
  small step of \(\sigma\)); the quotient of two such changes is a
  derivative; \(\mathrm{d}\) with an exponent, in front of several
  variables, is their volume element, the product of their
  differentials. Status: STANDARD.
\item
  \(\sigma\): The proper coordinate of the hidden direction:
  \(\mathrm{d}\sigma = \lvert\cot z\rvert\,\mathrm{d}x_0\),
  \(\sigma = \mathrm{sgn}(\pi/2 - z)\ln(\sin z)/(6H)\) on one period
  \(0 < z < \pi\); \(\sigma = 0\) is the fold \(z = \pi/2\),
  \(\sigma \to \mp\infty\) the tips \(z \to 0, \pi\); one period of
  \(x_0\) is the whole line \(-\infty < \sigma < \infty\). Unit: length.
  Status: DERIVED.
\item
  \(x\): The coordinates of the fold: \(x_1, x_2, x_3\) of ordinary
  space, \(x_4\) the time, \(x_5, x_6, x_7\) the extra times; a field of
  the fold is a function of the four observable coordinates
  \(x = (x_1, x_2, x_3, x_4)\), and
  \(\mathrm{d}^4x = \mathrm{d}x_1\,\mathrm{d}x_2\,\mathrm{d}x_3\,\mathrm{d}x_4\).
  Unit: length (each coordinate). Status: AUTHOR.
\item
  \(g_{MN}\): The metric tensor of the eight-dimensional spacetime
  (signature \(+,+,+,+,-,-,-,-\) in the order \(x_0, \dots, x_7\)):
  \(g_{MN}\,\mathrm{d}X^M\mathrm{d}X^N\) is the squared interval,
  \(g^{MN}\) its inverse, \(\lvert g\rvert\) the absolute value of its
  determinant (\(\sqrt{\lvert g\rvert}\,\mathrm{d}^8X\) is the proper
  volume element). Unit: dimensionless. Status: AUTHOR.
\item
  \(\mathrm{e}\): Euler's number, the base of the natural logarithm;
  \(\mathrm{e}^{X}\) is the exponential of \(X\). Unit: dimensionless.
  Value: 2.718281828. Status: STANDARD.
\item
  \(H\): The author's constant (an inverse length) that fixes the period
  \(\pi/(6H)\) of \(x_0\), the warp factor
  \(\mathrm{e}^{-2H\lvert\sigma\rvert}\) of the metric, the cosmological
  constant \(\Lambda = -21H^2\) and the tension
  \(\kappa\lambda_{\mathrm{F}} = 12H\) of the fold; in UM-4 a free
  parameter, bounded from below by the propagation of the gauge fields
  in the bulk. Unit: GeV (1/length). Value: \(H > 2534\) GeV, about 2.5
  TeV (from the invisible width of the Z boson). Status: AUTHOR.
\item
  \(\Psi,\ \bar\Psi\): dirac16complex: the author's 16-component complex
  spinor field of eight dimensions;
  \(\bar\Psi = \Psi^\dagger\sigma_{16}\) is its Dirac adjoint; in a term
  of the fold, \(\Psi\) is its value at \(\sigma = 0\). Unit:
  GeV\(^{7/2}\). Status: AUTHOR.
\item
  \({}^\dagger\): Hermitian conjugation: the complex conjugate of the
  transpose of a matrix or spinor; for an operator on the space of
  states, its adjoint with respect to the scalar product of that space
  (for the quantised dirac16complex the indefinite Krein scalar
  product). Status: STANDARD.
\item
  \(\mathfrak{B}\): The Krein operator,
  \(\mathfrak{B} = \mathrm{i}\sigma_{16}\gamma^{(4)}\): a Hermitian
  involution with the signature (8, 8); \(\Psi^\dagger\mathfrak{B}\Psi\)
  is the charge density. Unit: dimensionless. Status: DERIVED.
\item
  \(b_4\): The author's metric function of the time \(x_4\) in the
  extra-time components, \(b_4 > 0\), never identically 1 (the author's
  order of 2026-10-24); the scale factor (the factor of the lengths) of
  the extra times is \(\mathrm{e}^{-a_4}b_4\) on the fold and
  \(\mathrm{e}^{-a_4}b_4\,\mathrm{e}^{-H\lvert\sigma\rvert}\) at the
  distance \(\sigma\) from it (the square root of the warp factor times
  \(\mathrm{e}^{-a_4}b_4\)). Unit: dimensionless. Status: AUTHOR.
\item
  \(\chi(\sigma)\): The profile of a mode of dirac16complex along the
  hidden direction,
  \(\Psi = \mathrm{e}^{\frac72H\lvert\sigma\rvert}\chi(\sigma)\mathrm{e}^{\mathrm{i}(\vec k\cdot\vec x - \omega x_4)}\)
  (a 16-component spinor function of \(\sigma\)); as a four-dimensional
  field, \(\chi\) denotes the particle of the level of dirac16complex at
  the fold (a long-lived resonance) and \(\bar\chi\) its antiparticle.
  Unit: GeV\(^{2}\) (with the \(V_*\) normalisation). Status: DERIVED.
\item
  \(\mathrm{i}\): The imaginary unit, \(\mathrm{i}^2 = -1\). Unit:
  dimensionless. Status: STANDARD.
\item
  \(\sigma_{16}\): The matrix of the Dirac adjoint,
  \(\sigma_{16} = T_0T_1T_2T_3\), real and symmetric; \(\sigma_{16}T_a\)
  is real and antisymmetric. Unit: dimensionless. Status: AUTHOR.
\item
  \(\gamma^M,\ \gamma^{(a)}\): The Dirac matrices of dirac16complex
  along the coordinates, \(\gamma^M = \mathbf{e}_a{}^MT^a\),
  \(\gamma_M = g_{MN}\gamma^N\); \(\gamma^{(a)} = T^a\) along the frame
  direction \((a)\) (\(\gamma^{(0)}\) along \(\sigma\), \(\gamma^{(4)}\)
  along the time). Unit: dimensionless. Status: AUTHOR.
\item
  \(\mathbf{1},\ \mathbf{1}_4,\ \mathbf{1}_8,\ \mathbf{1}_{16}\): The
  identity matrix of the size fixed by the context (\(16 \times 16\) on
  dirac16complex, \(4 \times 4\) on a four-dimensional spinor,
  \(3 \times 3\) in generation or flavour space); \(\mathbf{1}_4\),
  \(\mathbf{1}_8\) and \(\mathbf{1}_{16}\) are the identity matrices of
  4, 8 and 16 rows. Unit: dimensionless. Status: STANDARD.
\item
  \(E\): An energy: \(E_o\) the energy of the Kohn-Sham orbital \(o\);
  in the neutrino sector the energy of a neutrino. Unit: GeV. Status:
  STANDARD.
\item
  \(o\): The label of a single-particle state of dirac16complex: a
  Kohn-Sham orbital, or a mode of the quantised field; a sum over \(o\)
  runs over the orbitals or the modes. Status: STANDARD.
\item
  \(\omega\): A frequency: the energy of a mode
  \(\propto \mathrm{e}^{-\mathrm{i}\omega x_4}\) (\(\omega_o\) of the
  orbital \(o\)). Unit: GeV. Status: STANDARD.
\item
  \(\mathcal{Q}_o\): The charge \(Q_\Psi\) (the Krein charge) of the
  single-particle state \(o\) of dirac16complex, a Kohn-Sham orbital or
  a mode: \(+1\) or \(-1\) for a normalised state. Unit: dimensionless.
  Status: DERIVED.
\end{itemize}

The conserved charge of dirac16complex (the Noether charge of a constant
change of the phase of \(\Psi\), whose density is
\(\mathrm{i}\bar\Psi\gamma^{x_4}\Psi\), with
\(\gamma^{x_4} = \mathrm{e}^{H\lvert\sigma\rvert}\gamma^{(4)}\) in the
metric of UM-4; per unit coordinate volume of the extra times) is a
Krein form: \(\mathfrak{B}\) is a Hermitian involution with eight
eigenvalues \(+1\) and eight eigenvalues \(-1\), so the charge, and with
it the norm of the quantised field, is indefinite (the consequence of
the signature (4,4); DERIVED, \texttt{um4\_fields.py}, F3).
\(\mathfrak{B}\) commutes with \(\gamma^{(0)}, \dots, \gamma^{(4)}\) and
anticommutes with \(\gamma^{(5)}, \gamma^{(6)}, \gamma^{(7)}\) and with
\(\Gamma_8\): every term that treats the two halves differently (here:
the fold terms with \(\mu_{\mathrm{L}} \ne \mu_{\mathrm{R}}\)) mixes the
two Krein sectors. The energy of a stationary orbital of frequency
\(\omega_o\) and charge \(\mathcal{Q}_o\) is
\(E_o = -\omega_o\mathcal{Q}_o\) (from FE-6), positive when the
frequency and the charge have opposite signs. Of the eight resonances at
\(+\omega_{\mathrm{r}}\), four have positive and four negative charge
(integrated over the region of the fold inside the turning point; no
null states): four carry positive and four negative energy; the mirror
\(\chi(\sigma) \to T_0T_1T_2T_3\,\chi(-\sigma)\) of the fold (with that
of ordinary space) maps the eight onto themselves, the four of one Krein
sign even, the four of the other odd. The quantisation of this charge is
FE-67. Status: DERIVED (\texttt{um4\_spectra.py}, U7 and U8;
\texttt{um4\_krein.py}, KR1 and KR3).

\[
\big\{\mathsf{a}_o, \mathsf{a}_o^\dagger\big\} = -\mathcal{Q}_o,\qquad
\mathsf{H}_\Psi = \sum_o\big(-\mathcal{Q}_o\,\omega_o\big)\,\mathsf{a}_o^\dagger\mathsf{a}_o,\qquad
\big[\mathsf{H}_\Psi, \mathsf{a}_o^\dagger\big] = \omega_o\,\mathsf{a}_o^\dagger \tag{FE-67}
\] \textbf{Symbols in (FE-67).}

\begin{itemize}
\tightlist
\item
  \(\mathsf{a}_o,\ \mathsf{a}_o^\dagger\): The annihilation and the
  creation operator of the mode \(o\) of the quantised dirac16complex
  (frequency \(\omega_o\), Krein charge \(\mathcal{Q}_o\)); the dagger
  is the adjoint with respect to the indefinite (Krein) scalar product.
  Unit: dimensionless. Status: STANDARD.
\item
  \(o\): The label of a single-particle state of dirac16complex: a
  Kohn-Sham orbital, or a mode of the quantised field; a sum over \(o\)
  runs over the orbitals or the modes. Status: STANDARD.
\item
  \({}^\dagger\): Hermitian conjugation: the complex conjugate of the
  transpose of a matrix or spinor; for an operator on the space of
  states, its adjoint with respect to the scalar product of that space
  (for the quantised dirac16complex the indefinite Krein scalar
  product). Status: STANDARD.
\item
  \(\mathcal{Q}_o\): The charge \(Q_\Psi\) (the Krein charge) of the
  single-particle state \(o\) of dirac16complex, a Kohn-Sham orbital or
  a mode: \(+1\) or \(-1\) for a normalised state. Unit: dimensionless.
  Status: DERIVED.
\item
  \(\mathsf{H}_\Psi\): The Hamiltonian (the energy operator) of the
  quantised dirac16complex. Unit: GeV. Status: DERIVED.
\item
  \(\omega\): A frequency: the energy of a mode
  \(\propto \mathrm{e}^{-\mathrm{i}\omega x_4}\) (\(\omega_o\) of the
  orbital \(o\)). Unit: GeV. Status: STANDARD.
\end{itemize}

The quantised dirac16complex. The canonical anticommutator of the field
is \(\{\Psi, \Psi^\dagger\} = -\mathfrak{B}\) (times the delta function
of the positions), the overall sign chosen so that the occupied states
of charge \(-1\), the particles \(\chi\), have a positive norm. A mode
\(o\) of the frequency \(\omega_o\) and the Krein charge
\(\mathcal{Q}_o = \pm1\) then has the anticommutator \(-\mathcal{Q}_o\),
the field has the Hamiltonian \(\mathsf{H}_\Psi\) (with FE-43, the
energy of the mode is \(-\mathcal{Q}_o\omega_o\)), and every quantum
raises the energy by its frequency. In the positive-energy quantisation
(the vacuum is annihilated by \(\mathsf{a}_o\) for \(\omega_o > 0\) and
by \(\mathsf{a}_o^\dagger\) for \(\omega_o < 0\)) every particle and
every hole carries a positive energy and the norm \(-\mathcal{Q}_o\) of
its mode: the particles and the holes of the modes of charge \(-1\) are
physical, those of the modes of charge \(+1\) are ghosts (negative
norm). A fermion cannot trade its norm for its energy
(\(\{\mathsf{a}, \mathsf{a}^\dagger\} = -1\) has no representation with
a positive norm), so no choice of the vacuum removes the ghosts. The
consequences (DERIVED, \texttt{um4\_krein.py}, KR1 to KR9, on the static
block of UM-4 with the warped time): the vacuum cannot decay, because
every quantum has a positive energy and the energy is conserved (with
the indefinite metric the vacuum is not a state of lowest energy among
the states of positive norm: a superposition of the vacuum and a ghost
has a negative mean energy). Whether the gas of the particles \(\chi\)
is classically stable is not decided: the boosts between ordinary space
and the extra times, which mix the two charges, are exact flat
directions of the self-consistent gas; the symmetric occupation of the
level is a saddle of the first-order energy, which has no lower bound on
the occupations of negative norm (15 negative, 16 zero and one positive
eigenvalue of its Hessian at the multiplet; towards the null cone it
falls without bound along some directions); the linear response
truncated to the occupied level, which keeps exactly the tangent space
of these occupations, has frequencies far below the splitting of the two
multiplets (0.06 of it), two of them imaginary, and whether they are an
instability of the full response is OPEN; and in the self-consistent
mean field the partner level of charge \(+1\) moves away from the
occupied level as the self-interaction grows (the two have opposite
parity under the mirror of the fold) and stays real and at the fold up
to \(G_{\mathrm{KS}} = 0.3\) (in UM-3 it turned complex in the denser
gases). The couplings at the fold decide unitarity: the chirality-blind
fold operators (products of \(\gamma^{(1)}, \dots, \gamma^{(4)}\)) never
connect a physical state with a ghost at rest, the chiral ones (the same
times \(\Gamma_8\)) always do, and the left-handed Higgs portal of the
author creates a physical particle \(\chi\) with a ghost antiparticle
with exactly its elastic amplitude (all four singular values of the
amplitude matrix, which do not depend on the basis of the degenerate
multiplets, equal the elastic one, 0.157); for moving particles (whose
energies are \(\sqrt{\omega_0^2 + \vec k^2}\) exactly, FE-78) the
magnetic dipole of the fold connects the two sectors as well, in
proportion to the momentum. The convective current of the current
portals (IN-16) creates no pair at rest: the amplitude of a pair
\(\chi(\vec k)\bar\chi(-\vec k)\) grows like \(\vec k^2\) (KR9), and its
chirality-blind combination creates a physical particle with a ghost
antiparticle only at the next order in the momentum (the relative
amplitude of these ghost pairs is \(0.094\lvert\vec k\rvert/m_\Psi\),
their relative rate \(0.0089\,\vec k^2/m_\Psi^2\): small, but not zero).
With chiral couplings at the fold UM-4 is therefore not a unitary
quantum theory, and chirality-blind couplings are necessary for
unitarity but not shown to be sufficient. Status: DERIVED, OPEN (the
classical stability of the gas, a unitary set of couplings).

\hypertarget{the-standard-model-on-the-fold}{%
\section{The Standard Model on the
fold}\label{the-standard-model-on-the-fold}}

\hypertarget{the-four-dimensional-couplings}{%
\subsection{The four-dimensional
couplings}\label{the-four-dimensional-couplings}}

\[
\begin{aligned}
&g_1^2 = \frac{3H\mathsf{g}_1^2}{2V_*},\qquad g_2^2 = \frac{3H\mathsf{g}_2^2}{2V_*},\qquad g_3^2 = \frac{3H\mathsf{g}_3^2}{2V_*},\qquad
 8\pi G_{\mathrm{N}} = \frac{\kappa_7}{V_*},\\
&g_{1,\mathrm{eff}} = \frac{g_1}{\sqrt{\mathcal{Z}}},\quad g_{2,\mathrm{eff}} = \frac{g_2}{\sqrt{\mathcal{Z}}},\quad g_{3,\mathrm{eff}} = \frac{g_3}{\sqrt{\mathcal{Z}}},\qquad \lambda_{\mathrm{eff}} = \frac{\lambda}{\mathcal{Z}},\qquad
 \mathcal{Y}_{\mathrm{eff}} = \frac{\mathcal{Y}}{\sqrt{\mathcal{Z}}},\qquad G_{\mathrm{N,eff}} = \frac{G_{\mathrm{N}}}{\mathcal{Z}},\\
&\mathcal{Z} = \frac{V_{\mathrm{t}}\big(\mathrm{e}^{-a_4}b_4\big)^3}{V_*}
\end{aligned} \tag{FE-44}
\] \textbf{Symbols in (FE-44).}

\begin{itemize}
\tightlist
\item
  \(g_1\): The \(U(1)_Y\) gauge coupling of the four-dimensional world
  (usually \(g'\); not normalised for grand unification): the coupling
  of the massless mode of the hypercharge field of the bulk,
  \(g_1^2 = 3H\mathsf{g}_1^2/(2V_*)\). Unit: dimensionless. Value:
  \(g_1(m_Z) = 0.3574\). Status: STANDARD.
\item
  \(g_2\): The \(SU(2)_L\) gauge coupling of the four-dimensional world
  (usually \(g\)): the coupling of the massless mode of the \(SU(2)_L\)
  field of the bulk, \(g_2^2 = 3H\mathsf{g}_2^2/(2V_*)\). Unit:
  dimensionless. Value: \(g_2(m_Z) = 0.6517\). Status: STANDARD.
\item
  \(g_3\): The \(SU(3)_C\) gauge coupling of the four-dimensional world
  (the strong coupling \(g_s\)): the coupling of the massless mode of
  the gluon field of the bulk, \(g_3^2 = 3H\mathsf{g}_3^2/(2V_*)\).
  Unit: dimensionless. Value: \(g_3(m_Z) = 1.2177\)
  (\(g_3^2/(4\pi) = 0.1180\)). Status: STANDARD.
\item
  \(H\): The author's constant (an inverse length) that fixes the period
  \(\pi/(6H)\) of \(x_0\), the warp factor
  \(\mathrm{e}^{-2H\lvert\sigma\rvert}\) of the metric, the cosmological
  constant \(\Lambda = -21H^2\) and the tension
  \(\kappa\lambda_{\mathrm{F}} = 12H\) of the fold; in UM-4 a free
  parameter, bounded from below by the propagation of the gauge fields
  in the bulk. Unit: GeV (1/length). Value: \(H > 2534\) GeV, about 2.5
  TeV (from the invisible width of the Z boson). Status: AUTHOR.
\item
  \(\mathsf{g}_1,\ \mathsf{g}_2,\ \mathsf{g}_3\): The gauge couplings of
  eight dimensions of \(U(1)_Y\), \(SU(2)_L\) and \(SU(3)_C\): the
  coefficients of the kinetic terms of the gauge fields of the bulk in
  their geometric normalisation; the four-dimensional couplings are
  \(g_k^2 = 3H\mathsf{g}_k^2/(2V_*)\). Unit: GeV\(^{-2}\)
  (\(\mathsf{g}^2\): length\(^4\)). Status: ASSUMED.
\item
  \(V_*\): The physical volume of the extra times \(x_5, x_6, x_7\) on
  the fold in the late universe,
  \(V_* = V_{\mathrm{t}}(\mathrm{e}^{-a_4}b_4)^3\) (with
  \(V_{\mathrm{t}} = \int\mathrm{d}x_5\,\mathrm{d}x_6\,\mathrm{d}x_7\)
  the coordinate volume of the extra times) at a time when
  \(\mathrm{e}^{-a_4}b_4\) is constant; it converts the
  eight-dimensional normalisation into the four-dimensional one. Unit:
  length\(^3\) (GeV\(^{-3}\)). Status: ASSUMED.
\item
  \(\pi\): The number pi, the ratio of the circumference of a circle to
  its diameter. Unit: dimensionless. Value: 3.141592654. Status:
  STANDARD.
\item
  \(G_{\mathrm{N}}\): Newton's constant of four-dimensional gravity (the
  gravity of our world, on the fold),
  \(G_{\mathrm{N}} = 1/M_{\mathrm{Pl}}^2\) with the Planck mass
  \(M_{\mathrm{Pl}} = 1.221\times10^{19}\) GeV. Unit: GeV\(^{-2}\).
  Value: \(6.708\times10^{-39}\) GeV\(^{-2}\). Status: STANDARD.
\item
  \(\kappa_7\): The gravitational coupling of the fold: to the first
  order in its curvature over \(H^2\) the fold obeys the Einstein
  equations of seven dimensions
  \(\hat G^a{}_b = \kappa_7\mathsf{t}^a{}_b\) with
  \(\kappa_7 = 5\kappa H/2\), the coupling of the massless graviton of
  the bulk bound to the fold; four-dimensional gravity on the fold has
  \(8\pi G_{\mathrm{N}} = \kappa_7/V_*\) (plus the induced term, which
  is negligible). Unit: GeV\(^{-5}\). Value: \(5\kappa H/2\). Status:
  DERIVED.
\item
  \(g_{k,\mathrm{eff}}\): The effective (canonically normalised) gauge
  couplings of the fold when the fold factor is \(\mathcal{Z}\),
  \(g_{k,\mathrm{eff}} = g_k/\sqrt{\mathcal{Z}}\) (\(k = 1, 2, 3\)): the
  kinetic terms of the massless modes of the gauge fields of the bulk
  carry the volume of the extra times as the terms of the fold do. Unit:
  dimensionless. Status: DERIVED.
\item
  \(\mathcal{Z}\): The fold factor,
  \(\mathcal{Z}(x_4) = V_{\mathrm{t}}(\mathrm{e}^{-a_4}b_4)^3/V_*\): the
  factor by which the size of the extra times multiplies every term of
  the fold and the kinetic terms of the massless modes of the fields of
  the bulk; \(\mathcal{Z} = 1\) when the extra times are frozen at their
  late value. Unit: dimensionless. Value: 1 in the late universe.
  Status: DERIVED.
\item
  \(\lambda_{\mathrm{eff}}\): The effective quartic Higgs coupling,
  \(\lambda/\mathcal{Z}\). Unit: dimensionless. Status: DERIVED.
\item
  \(\lambda\): The quartic coupling of the Higgs potential,
  \(\lambda > 0\). Unit: dimensionless. Value:
  \(\lambda = m_h^2/(2v^2) = 0.129\). Status: AUTHOR.
\item
  \(\mathcal{Y}_{\mathrm{eff}}\): The effective Yukawa matrices,
  \(\mathcal{Y}/\sqrt{\mathcal{Z}}\) (for each of
  \(\mathcal{Y}^{u}, \mathcal{Y}^{d}, \mathcal{Y}^{e}, \mathcal{Y}^{\nu}\)).
  Unit: dimensionless. Status: DERIVED.
\item
  \(\mathcal{Y}^{u},\ \mathcal{Y}^{d},\ \mathcal{Y}^{e},\ \mathcal{Y}^{\nu}\):
  The Yukawa matrices (\(3 \times 3\) complex matrices in generation
  space) of the up quarks, the down quarks, the charged leptons and the
  neutrinos. Unit: dimensionless. Status: STANDARD.
\item
  \(G_{\mathrm{N,eff}}\): The effective Newton constant of the fold,
  \(G_{\mathrm{N}}/\mathcal{Z}\). Unit: GeV\(^{-2}\). Status: DERIVED.
\item
  \(V_{\mathrm{t}}\): The coordinate volume of the (compact) extra
  times,
  \(V_{\mathrm{t}} = \int\mathrm{d}x_5\,\mathrm{d}x_6\,\mathrm{d}x_7\).
  Unit: length\(^3\). Status: ASSUMED.
\item
  \(\mathrm{e}\): Euler's number, the base of the natural logarithm;
  \(\mathrm{e}^{X}\) is the exponential of \(X\). Unit: dimensionless.
  Value: 2.718281828. Status: STANDARD.
\item
  \(a_4\): The author's metric function of the time \(x_4\) that scales
  ordinary 3-space: on the fold \(\mathrm{e}^{a_4}\) is the scale factor
  of cosmology and \(a_4'\) the Hubble rate. Unit: dimensionless.
  Status: AUTHOR.
\item
  \(b_4\): The author's metric function of the time \(x_4\) in the
  extra-time components, \(b_4 > 0\), never identically 1 (the author's
  order of 2026-10-24); the scale factor (the factor of the lengths) of
  the extra times is \(\mathrm{e}^{-a_4}b_4\) on the fold and
  \(\mathrm{e}^{-a_4}b_4\,\mathrm{e}^{-H\lvert\sigma\rvert}\) at the
  distance \(\sigma\) from it (the square root of the warp factor times
  \(\mathrm{e}^{-a_4}b_4\)). Unit: dimensionless. Status: AUTHOR.
\end{itemize}

The gauge couplings of the Standard Model are those of the massless
modes of the gauge fields of the bulk (FE-74): the integral of the
kinetic weight \(\mathrm{e}^{-3H\lvert\sigma\rvert}\) over the hidden
direction, \(2/(3H)\), and the volume \(V_*\) of the extra times turn
the couplings \(\mathsf{g}_k\) of eight dimensions into the
four-dimensional couplings \(g_k\) (first line, \(k = 1, 2, 3\));
Newton's constant is that of the graviton of the fold (FE-71; the
induced term IN-9 adds \(1/\kappa_{\mathrm{ind}}\) to
\(1/(8\pi G_{\mathrm{N}})\), FE-59). When the extra times have the size
\(\mathrm{e}^{-a_4}b_4\), every term of the fold and the kinetic terms
of the zero modes of the bulk carry the factor \(\mathcal{Z}\) (IN-8);
rescaling every field by \(\sqrt{\mathcal{Z}}\) restores the canonical
normalisation and leaves the effective couplings of the second line,
while the masses (products of a coupling and a vacuum value) are
unchanged. With frozen extra times \(\mathcal{Z} = 1\) (FE-72), and the
couplings are constant; if the volume of the extra times changed, the
gauge and Yukawa couplings would change with \(\mathcal{Z}^{-1/2}\), the
quartic Higgs coupling and Newton's constant with \(\mathcal{Z}^{-1}\),
which observation excludes (FE-72): the fold fluid holds the volume
fixed (FE-65). Status: DERIVED (\texttt{um4\_bulk.py}, B2;
\texttt{um4\_background.py}, BG4).

\hypertarget{the-electroweak-vacuum-with-the-ad-hoc-component}{%
\subsection{The electroweak vacuum with the ad hoc
component}\label{the-electroweak-vacuum-with-the-ad-hoc-component}}

\[
\begin{aligned}
&m_{\mathrm{W}}^2 = \frac{g_2^2\,(v^2 + \mathcal{A}^2)}{4},\qquad m_{\mathrm{Z}}^2 = \frac{(g_1^2 + g_2^2)(v^2 + \mathcal{A}^2)}{4},\qquad
 \frac{G_{\mathrm{F}}}{\sqrt2} = \frac{g_2^2}{8m_{\mathrm{W}}^2} = \frac{1}{2(v^2 + \mathcal{A}^2)},\qquad m_{h}^2 = 2\lambda v^2,\\
&A_\mu = \frac{g_2\,B_\mu - g_1\,n^I W^I_\mu}{\sqrt{g_1^2 + g_2^2}},\qquad
 n^I = \frac{\big(2\mathcal{A}v,\ 0,\ \mathcal{A}^2 - v^2\big)}{\mathcal{A}^2 + v^2}
\end{aligned} \tag{FE-45}
\] \textbf{Symbols in (FE-45).}

\begin{itemize}
\tightlist
\item
  \(m_{\mathrm{W}}\): The mass of the \(W\) bosons. Unit: GeV. Value:
  80.37 GeV (\(\mathcal{A} = 0\)). Status: STANDARD.
\item
  \(g_2\): The \(SU(2)_L\) gauge coupling of the four-dimensional world
  (usually \(g\)): the coupling of the massless mode of the \(SU(2)_L\)
  field of the bulk, \(g_2^2 = 3H\mathsf{g}_2^2/(2V_*)\). Unit:
  dimensionless. Value: \(g_2(m_Z) = 0.6517\). Status: STANDARD.
\item
  \(v\): The vacuum expectation value of the lower component of the
  Higgs doublet, \(v^2 = -\mu^2/\lambda\). Unit: GeV. Value: 246.22 GeV
  \(= (\sqrt2\,G_{\mathrm{F}})^{-1/2}\) with the Fermi constant
  \(G_{\mathrm{F}} = 1.1663788\times10^{-5}\) GeV\(^{-2}\) of the muon
  lifetime (Particle Data Group). Status: STANDARD.
\item
  \(\mathcal{A}\): The author's ad hoc upper component of the Higgs
  vacuum, almightyAdHoc,
  \(\mathcal{A}(x_4) = v_{\mathcal{A}}\mathrm{e}^{-2a_{40}x_4}\). Unit:
  GeV. Status: AUTHOR.
\item
  \(m_{\mathrm{Z}}\): The mass of the Z boson. Unit: GeV. Value: 91.19
  GeV (\(\mathcal{A} = 0\)). Status: STANDARD.
\item
  \(g_1\): The \(U(1)_Y\) gauge coupling of the four-dimensional world
  (usually \(g'\); not normalised for grand unification): the coupling
  of the massless mode of the hypercharge field of the bulk,
  \(g_1^2 = 3H\mathsf{g}_1^2/(2V_*)\). Unit: dimensionless. Value:
  \(g_1(m_Z) = 0.3574\). Status: STANDARD.
\item
  \(G_{\mathrm{F}}\): The Fermi constant,
  \(G_{\mathrm{F}} = 1/(\sqrt2\,v^2)\) in the Standard-Model vacuum (its
  measured value); in the equations of the neutrino flavours it stands
  for the local Fermi constant of the ad hoc vacuum. Unit: GeV\(^{-2}\).
  Value: \(1.1664\times10^{-5}\) GeV\(^{-2}\). Status: STANDARD.
\item
  \(m_h\): The mass of the Higgs boson. Unit: GeV. Value: 125.25 GeV
  (Particle Data Group 2022). Status: STANDARD.
\item
  \(\lambda\): The quartic coupling of the Higgs potential,
  \(\lambda > 0\). Unit: dimensionless. Value:
  \(\lambda = m_h^2/(2v^2) = 0.129\). Status: AUTHOR.
\item
  \(A_\mu\): The photon field of the fold (canonically normalised): the
  massless combination of the massless modes of \(W^1, W^2, W^3\) and
  \(B\) that the Higgs vacuum leaves massless; a field of the bulk whose
  zero mode travels at the speed of light. Unit: GeV. Status: STANDARD.
\item
  \(\mu\): A four-dimensional coordinate index running over the
  observable coordinates \(x_1, x_2, x_3, x_4\) of the fold (signature
  \(+,+,+,-\)); a repeated index is summed. Status: STANDARD.
\item
  \(B_M,\ B_\mu,\ B_{MN}\): The hypercharge gauge field of \(U(1)_Y\):
  in the bulk \(B_M\) (geometric normalisation), on the fold the
  canonically normalised four-dimensional field \(B_\mu(x)\) of its
  massless mode, \(g_1B_\mu(x) = B_\mu(\sigma = 0, x)\); \(B_{MN}\),
  \(B_{\mu\nu}\) their field strengths. Unit: GeV (field), GeV\(^2\)
  (field strength). Status: AUTHOR.
\item
  \(n^I\): The unit isovector of the ad hoc vacuum,
  \(n = (2\mathcal{A}v, 0, \mathcal{A}^2 - v^2)/(\mathcal{A}^2 + v^2)\),
  which fixes which combination of \(W^1, W^2, W^3\) joins \(B\) in the
  photon. Unit: dimensionless. Value: \((0, 0, -1)\) for
  \(\mathcal{A} = 0\). Status: DERIVED.
\item
  \(I\): A weak-isospin (adjoint SU(2)) index running over \(1, 2, 3\);
  a repeated index is summed. Status: STANDARD.
\item
  \(W^I_M,\ W^I_\mu,\ W^I_{MN}\): The gauge field of \(SU(2)_L\)
  (\(I = 1, 2, 3\)): in the bulk \(W^I_M\) (geometric normalisation), on
  the fold the canonically normalised four-dimensional field
  \(W^I_\mu(x)\) of its massless mode,
  \(g_2W^I_\mu(x) = W^I_\mu(\sigma = 0, x)\); \(W^I_{MN}\),
  \(W^I_{\mu\nu}\) their field strengths;
  \(W^\pm_\mu = (W^1_\mu \mp \mathrm{i}W^2_\mu)/\sqrt2\). Unit: GeV
  (field), GeV\(^2\) (field strength). Status: AUTHOR.
\end{itemize}

The masses of the gauge bosons and the photon in the author's ad hoc
vacuum (the eigenvalues and eigenvectors of the mass matrix of the
canonical fields \(W^1, W^2, W^3, B\) of the fold from
\(\lvert D_\mu\Phi\rvert^2\) with \(\Phi = (\mathcal{A}, v)/\sqrt2\);
DERIVED, \texttt{um4\_fields.py}, F5): the vacuum value that breaks the
electroweak symmetry is \(\sqrt{v^2 + \mathcal{A}^2}\), and the unbroken
direction, the photon \(A_\mu\), is rotated from the usual combination
by the unit isovector \(n^I\) (for \(\mathcal{A} = 0\):
\(n = (0, 0, -1)\) and
\(A_\mu = (g_2B_\mu + g_1W^3_\mu)/\sqrt{g_1^2 + g_2^2}\)). The Fermi
constant measured in muon decay is that of the total vacuum value. As
\(\mathcal{A}\) decays (rate \(2a_{40}\)), the electroweak vacuum
becomes the Standard-Model vacuum. The Higgs mass is that of the
Standard-Model vacuum (\(\mathcal{A} \to 0\)). These are the masses of
the zero modes of the gauge fields of the bulk; their massive modes
lower the squares of the masses of the W and the Z by
\(m_{\mathrm{W}}^4/(3H^2)\) and \(m_{\mathrm{Z}}^4/(3H^2)\) and give
them the widths \(m(m/H)^3/3\) into the bulk (FE-82), while the photon
and the gluons stay massless. Status: DERIVED.

\hypertarget{the-masses-and-the-mixing-of-quarks-and-leptons}{%
\subsection{The masses and the mixing of quarks and
leptons}\label{the-masses-and-the-mixing-of-quarks-and-leptons}}

\[
\begin{aligned}
&\frac{v}{\sqrt2}\,\mathcal{Y}^{u} = U_{u\mathrm{L}}\,\mathrm{diag}(m_u, m_c, m_t)\,U_{u\mathrm{R}}^\dagger,\qquad
 \frac{v}{\sqrt2}\,\mathcal{Y}^{d} = U_{d\mathrm{L}}\,\mathrm{diag}(m_d, m_s, m_b)\,U_{d\mathrm{R}}^\dagger,\\
&\frac{v}{\sqrt2}\,\mathcal{Y}^{e} = U_{e\mathrm{L}}\,\mathrm{diag}(m_e, m_\mu, m_\tau)\,U_{e\mathrm{R}}^\dagger,\qquad
 \frac{v}{\sqrt2}\,\mathcal{Y}^{\nu} = U_{\nu\mathrm{L}}\,\mathrm{diag}(m_1, m_2, m_3)\,U_{\nu\mathrm{R}}^\dagger,\\
&V_{\mathrm{CKM}} = U_{u\mathrm{L}}^\dagger U_{d\mathrm{L}},\qquad U_{\mathrm{PMNS}} = U_{e\mathrm{L}}^\dagger U_{\nu\mathrm{L}},\qquad
 \nu_\alpha = \sum_{i=1}^{3}\big(U_{\mathrm{PMNS}}\big)_{\alpha i}\,\nu_i,\\
&\mathcal{L}_{\mathrm{CC}} = \frac{g_2}{\sqrt2}\,W^+_\mu\Big(\bar{\hat u}_{\mathrm{L}}\,\hat\gamma^\mu V_{\mathrm{CKM}}\,\hat d_{\mathrm{L}}
 + \sum_\alpha\bar\nu_\alpha\,\hat\gamma^\mu\hat e_{\mathrm{L}\alpha}\Big) + \mathrm{h.c.},\qquad W^\pm_\mu = \frac{W^1_\mu \mp \mathrm{i}W^2_\mu}{\sqrt2}
\end{aligned} \tag{FE-46}
\] \textbf{Symbols in (FE-46).}

\begin{itemize}
\tightlist
\item
  \(v\): The vacuum expectation value of the lower component of the
  Higgs doublet, \(v^2 = -\mu^2/\lambda\). Unit: GeV. Value: 246.22 GeV
  \(= (\sqrt2\,G_{\mathrm{F}})^{-1/2}\) with the Fermi constant
  \(G_{\mathrm{F}} = 1.1663788\times10^{-5}\) GeV\(^{-2}\) of the muon
  lifetime (Particle Data Group). Status: STANDARD.
\item
  \(\mathcal{Y}^{u},\ \mathcal{Y}^{d},\ \mathcal{Y}^{e},\ \mathcal{Y}^{\nu}\):
  The Yukawa matrices (\(3 \times 3\) complex matrices in generation
  space) of the up quarks, the down quarks, the charged leptons and the
  neutrinos. Unit: dimensionless. Status: STANDARD.
\item
  \(U_{f\mathrm{L}},\ U_{f\mathrm{R}}\): The unitary \(3 \times 3\)
  matrices that diagonalise the mass matrices of the up quarks, the down
  quarks, the charged leptons and the neutrinos (left-handed and
  right-handed rotations). Unit: dimensionless. Status: STANDARD.
\item
  \(\mathrm{diag}\): A diagonal matrix with the listed entries on its
  diagonal (in the order of the coordinates \(x_0, x_1, \dots, x_7\) for
  a tensor of eight dimensions, \(x_1, \dots, x_7\) for a tensor of the
  fold, \(x_1, \dots, x_4\) for a four-dimensional tensor). Status:
  STANDARD.
\item
  \(m_u, \dots, m_3\): The masses of the quarks (\(m_u, m_c, m_t\);
  \(m_d, m_s, m_b\)), of the charged leptons (\(m_e, m_\mu, m_\tau\))
  and of the three neutrino mass states (\(m_1, m_2, m_3\), Dirac
  masses). Unit: GeV. Value: measured values;
  \(\Delta m_{21}^2 = 7.41\times10^{-5}\) eV\(^2\),
  \(\lvert\Delta m_{31}^2\rvert = 2.511\times10^{-3}\) eV\(^2\) in the
  computations of UM-4 (the global fit NuFIT 5.2 of the oscillation
  data, November 2022), \(7.49\times10^{-5}\) and \(2.513\times10^{-3}\)
  eV\(^2\) in NuFIT 6.0 (2024); the lightest mass \(m_1 = 0\) in the
  computations, below 23 meV by DESI DR2 with the CMB and NuFIT 6.0.
  Status: STANDARD.
\item
  \({}^\dagger\): Hermitian conjugation: the complex conjugate of the
  transpose of a matrix or spinor; for an operator on the space of
  states, its adjoint with respect to the scalar product of that space
  (for the quantised dirac16complex the indefinite Krein scalar
  product). Status: STANDARD.
\item
  \(V_{\mathrm{CKM}}\): The quark mixing matrix
  (Cabibbo-Kobayashi-Maskawa). Unit: dimensionless. Status: STANDARD.
\item
  \(U_{\mathrm{PMNS}}\): The lepton mixing matrix
  (Pontecorvo-Maki-Nakagawa-Sakata) of the three Dirac neutrinos; its
  entries \((U_{\mathrm{PMNS}})_{\alpha i}\) connect the flavour state
  \(\nu_\alpha\) with the mass state \(\nu_i\). Unit: dimensionless.
  Value: \(\theta_{12} = 33.40^\circ\), \(\theta_{23} = 49.14^\circ\),
  \(\theta_{13} = 8.54^\circ\) (NuFIT 5.2, November 2022, normal
  ordering, the fit without the Super-Kamiokande atmospheric data); the
  Dirac phase is set to zero (ASSUMED; the fit prefers \(197^\circ\)).
  Status: STANDARD, ASSUMED (the phase).
\item
  \(\nu_\alpha,\ \nu_i\): A neutrino field: \(\nu_\alpha\) the flavour
  states (the flavour index runs over the electron, muon and tau
  flavours), \(\nu_i\) the mass states (\(i = 1, 2, 3\)). Unit:
  GeV\(^{3/2}\). Status: STANDARD.
\item
  \(\alpha\): A lepton-flavour index, \(\alpha \in \{e, \mu, \tau\}\).
  Status: STANDARD.
\item
  \(i\): A generation index (\(i = 1, 2, 3\)), which also labels the
  three neutrino mass states; in a frame component such as
  \(\mathcal{S}^{(i)}{}_{(i)}\) one of the three directions of ordinary
  space (the three components are equal), and in a tensor of ordinary
  space such as \(\mathsf{h}_{ij}\) one of the directions
  \(x_1, x_2, x_3\); a repeated index is not summed unless a sum sign is
  written. Status: STANDARD.
\item
  \(\mathcal{L}_{\mathrm{CC}}\): The charged-current interaction of the
  quarks and leptons with the \(W\) bosons. Unit: GeV\(^4\). Status:
  STANDARD.
\item
  \(g_2\): The \(SU(2)_L\) gauge coupling of the four-dimensional world
  (usually \(g\)): the coupling of the massless mode of the \(SU(2)_L\)
  field of the bulk, \(g_2^2 = 3H\mathsf{g}_2^2/(2V_*)\). Unit:
  dimensionless. Value: \(g_2(m_Z) = 0.6517\). Status: STANDARD.
\item
  \(W^I_M,\ W^I_\mu,\ W^I_{MN}\): The gauge field of \(SU(2)_L\)
  (\(I = 1, 2, 3\)): in the bulk \(W^I_M\) (geometric normalisation), on
  the fold the canonically normalised four-dimensional field
  \(W^I_\mu(x)\) of its massless mode,
  \(g_2W^I_\mu(x) = W^I_\mu(\sigma = 0, x)\); \(W^I_{MN}\),
  \(W^I_{\mu\nu}\) their field strengths;
  \(W^\pm_\mu = (W^1_\mu \mp \mathrm{i}W^2_\mu)/\sqrt2\). Unit: GeV
  (field), GeV\(^2\) (field strength). Status: AUTHOR.
\item
  \(\mu\): A four-dimensional coordinate index running over the
  observable coordinates \(x_1, x_2, x_3, x_4\) of the fold (signature
  \(+,+,+,-\)); a repeated index is summed. Status: STANDARD.
\item
  \(\hat u_{\mathrm{L}},\ \hat d_{\mathrm{L}},\ \hat e_{\mathrm{L}\alpha}\):
  The left-handed quarks and charged leptons in the mass basis
  (\(\hat e_{\mathrm{L}\alpha}\) the charged lepton of flavour
  \(\alpha\)). Unit: GeV\(^{3/2}\). Status: STANDARD.
\item
  \(\hat\gamma^\mu\): The four-dimensional complex \(4 \times 4\) Dirac
  matrices of the fold: \(\hat\gamma^{(a)}\) along the frame directions
  of the fold and \(\hat\gamma^\mu\) along its coordinates (the frame
  matrices contracted with the inverse frame field of the fold),
  \(\{\hat\gamma^\mu, \hat\gamma^\nu\} = -2\hat h^{\mu\nu}\mathbf{1}\);
  \(\hat\gamma_\mu = \hat h_{\mu\nu}\hat\gamma^\nu\). Unit:
  dimensionless. Status: STANDARD.
\item
  \(\mathrm{h.c.}\): The Hermitian conjugate of the preceding terms.
  Status: STANDARD.
\item
  \(\mathrm{i}\): The imaginary unit, \(\mathrm{i}^2 = -1\). Unit:
  dimensionless. Status: STANDARD.
\end{itemize}

The Dirac masses of the quarks, the charged leptons and the three
neutrinos (in the vacuum with \(\mathcal{A} = 0\); with
\(\mathcal{A} \ne 0\) the same with the rotated vacuum), the unitary
matrices that diagonalise the Yukawa matrices, the quark mixing matrix
\(V_{\mathrm{CKM}}\) and the lepton mixing matrix \(U_{\mathrm{PMNS}}\),
and the charged-current interaction in the mass basis (fields with a
hat; the neutrino flavour states \(\nu_\alpha\) are the partners of the
charged leptons of definite mass). This is the flavour structure of the
Standard Model with three Dirac neutrinos; the portals of dirac16complex
add nothing to it (they are flavour blind, IN-16). The values of the
masses and mixing angles are the measured ones (STANDARD inputs).
Status: STANDARD.

\hypertarget{the-four-dimensional-couplings-of-the-particles-of-dirac16complex}{%
\subsection{The four-dimensional couplings of the particles of
dirac16complex}\label{the-four-dimensional-couplings-of-the-particles-of-dirac16complex}}

\[
\begin{aligned}
&\mathfrak{g}_{\Phi} = \frac{\partial\omega_{\mathrm{b}}}{\partial\mu_{\mathrm{L}}}\,\frac{c_\Phi}{V_*},\qquad
 \mathfrak{g}_{G} = \sum_X\frac{\partial\omega_{\mathrm{b}}}{\partial\mu_X}\,\frac{c^{G}_X}{4V_*},\qquad
 \mathfrak{g}_{f} = \sum_X\frac{\partial\omega_{\mathrm{b}}}{\partial\mu_X}\,\frac{2\omega_{\mathrm{b}}\,c^{f}_X}{V_*},\qquad
 \frac{\partial\omega_{\mathrm{b}}}{\partial\mu_X} = -\frac{q_{\mathrm{b}}}{\omega_{\mathrm{b}}}\,\frac{\partial q_{\mathrm{b}}}{\partial\mu_X},\\
&\Delta\omega_{\mathrm{b}} = \mathfrak{g}_{\Phi}\,\Delta\langle\Phi^\dagger\Phi\rangle + \mathfrak{g}_{G}\,\Delta\langle G^2\rangle + \sum_f\mathfrak{g}_{f}\,\Delta n_f,\qquad
 \mathfrak{g}_{h\chi} = \mathfrak{g}_{\Phi}\,v = -\frac{m_\Psi^2\,\mu_{\mathrm{L}}^2}{2\,\omega_{\mathrm{b}}\,v}\quad(\mu_{\mathrm{R}} = 0)
\end{aligned} \tag{FE-47}
\] \textbf{Symbols in (FE-47).}

\begin{itemize}
\tightlist
\item
  \(\mathfrak{g}_{\Phi},\ \mathfrak{g}_{G},\ \mathfrak{g}_{f},\ \mathfrak{g}_{h\chi}\):
  The four-dimensional couplings of the particles \(\chi\) of
  dirac16complex to the Standard Model: to \(\Phi^\dagger\Phi\), to the
  gluon kinetic density, to the number densities of the fermions of kind
  \(f\), and the Yukawa coupling of the Higgs boson to \(\chi\). Unit:
  \(\mathfrak{g}_{\Phi}\): GeV\(^{-1}\); \(\mathfrak{g}_{G}\):
  GeV\(^{-3}\); \(\mathfrak{g}_{f}\): GeV\(^{-2}\);
  \(\mathfrak{g}_{h\chi}\): dimensionless. Status: DERIVED.
\item
  \(\partial\): A partial derivative: the derivative with respect to one
  variable at fixed values of the others (\(\partial_\sigma\) with
  respect to \(\sigma\); \(\partial_M\) and \(\partial_\mu\) with
  respect to the coordinate that carries the same index, of the eight
  dimensions or of the fold). Status: STANDARD.
\item
  \(\omega_{\mathrm{b}}\): The energy at zero momentum of the state of
  dirac16complex at the fold in the limit \(H \to 0\) (the bound state
  of UM-3): the particle \(\chi\) of UM-4 has the mass
  \(\omega_{\mathrm{r}} = \omega_{\mathrm{b}}(1 - H/(2q_{\mathrm{b}})) + \dots\),
  equal to \(\omega_{\mathrm{b}}\) up to the relative correction
  \(H/(2q_{\mathrm{b}})\). Unit: GeV. Status: DERIVED.
\item
  \(\mu_X,\ \mu_{\mathrm{L}},\ \mu_{\mathrm{R}}\): The dimensionless
  strengths of the fold for the left-handed and right-handed halves of
  dirac16complex, set by the mean fields of the Standard Model at the
  fold: the jump of the profile across the fold is the exponential of
  \(\gamma^{(0)}(\mu_{\mathrm{L}}P_{\mathrm{L}} + \mu_{\mathrm{R}}P_{\mathrm{R}})\);
  negative values attract. Unit: dimensionless. Status: DERIVED.
\item
  \(c_\Phi\): The Higgs-portal coupling of the left-handed half of
  dirac16complex (the right-handed half has none). Unit: GeV\(^{-5}\).
  Status: INVENTED.
\item
  \(V_*\): The physical volume of the extra times \(x_5, x_6, x_7\) on
  the fold in the late universe,
  \(V_* = V_{\mathrm{t}}(\mathrm{e}^{-a_4}b_4)^3\) (with
  \(V_{\mathrm{t}} = \int\mathrm{d}x_5\,\mathrm{d}x_6\,\mathrm{d}x_7\)
  the coordinate volume of the extra times) at a time when
  \(\mathrm{e}^{-a_4}b_4\) is constant; it converts the
  eight-dimensional normalisation into the four-dimensional one. Unit:
  length\(^3\) (GeV\(^{-3}\)). Status: ASSUMED.
\item
  \(c^{G}_X\): The gluon-kinetic portal coupling of the half \(X\) of
  dirac16complex. Unit: GeV\(^{-7}\). Status: INVENTED.
\item
  \(X\): The label of a half of dirac16complex:
  \(X \in \{\mathrm{L}, \mathrm{R}\}\) (left-handed or right-handed
  half); a sum over \(X\) runs over both. Status: STANDARD.
\item
  \(c^{f}_X\): The current-portal coupling of the half \(X\) of
  dirac16complex with the fermion-number current of the multiplets of
  kind \(f\) (flavour blind); \(c^{\ell}_X\) and \(c^{\nu}_X\) are those
  of the lepton doublets and of the right-handed neutrinos. Unit:
  GeV\(^{-7}\). Status: INVENTED.
\item
  \(f\): The label of a kind of Standard-Model multiplet:
  \(f \in \{Q, u, d, \ell, e, \nu\}\) (quark doublets, up-quark
  singlets, down-quark singlets, lepton doublets, charged-lepton
  singlets, right-handed neutrinos); a sum over \(f\) runs over the six.
  Status: STANDARD.
\item
  \(q_{\mathrm{b}}\): The inverse width of a state of dirac16complex at
  the fold: inside the turning point its profile falls as
  \(\mathrm{e}^{-q_{\mathrm{b}}\lvert\sigma\rvert}\). Unit: GeV. Status:
  DERIVED.
\item
  \(\Delta\): A change, difference or small variation of the quantity
  that follows (for example an interval of the time, a perturbation of a
  density, a difference of two squared masses, or an addition to an
  effective number of species). Status: STANDARD.
\item
  \(\Phi\): The Higgs doublet (two complex components, hypercharge
  \(1/2\)), a field of the fold. Unit: GeV. Status: AUTHOR.
\item
  \({}^\dagger\): Hermitian conjugation: the complex conjugate of the
  transpose of a matrix or spinor; for an operator on the space of
  states, its adjoint with respect to the scalar product of that space
  (for the quantised dirac16complex the indefinite Krein scalar
  product). Status: STANDARD.
\item
  \(\langle\cdots\rangle\): The mean (expectation) value of an operator
  in the state of the medium (the vacuum, a plasma, the Kohn-Sham
  state). Status: STANDARD.
\item
  \(\langle G^2\rangle\): The mean square of the gluon field strength,
  \(\langle G^A_{\mu\nu}G^{A\mu\nu}\rangle\) (the gluon condensate in
  the QCD vacuum, a thermal value in the plasma). Unit: GeV\(^4\).
  Status: STANDARD.
\item
  \(G^A_M,\ G^A_\mu,\ G^A_{MN}\): The gluon field of \(SU(3)_C\)
  (\(A = 1, \dots, 8\)): in the bulk \(G^A_M\), a field of eight
  dimensions in its geometric normalisation (the covariant derivative
  \(\partial - \mathrm{i}G\)), and on the fold the canonically
  normalised four-dimensional field \(G^A_\mu(x)\) of its massless mode,
  \(g_3G^A_\mu(x) = G^A_\mu(\sigma = 0, x)\); \(G^A_{MN}\) and
  \(G^A_{\mu\nu}\) their field strengths. Unit: GeV (field), GeV\(^2\)
  (field strength). Status: AUTHOR.
\item
  \(n_f\): The number density of the Standard-Model fermions of kind
  \(f\) (particles minus antiparticles, summed over generations and
  gauge components), the time component of the mean current
  \(\langle J_f^\mu\rangle\) in the rest frame of the medium. Unit:
  GeV\(^3\). Status: STANDARD.
\item
  \(v\): The vacuum expectation value of the lower component of the
  Higgs doublet, \(v^2 = -\mu^2/\lambda\). Unit: GeV. Value: 246.22 GeV
  \(= (\sqrt2\,G_{\mathrm{F}})^{-1/2}\) with the Fermi constant
  \(G_{\mathrm{F}} = 1.1663788\times10^{-5}\) GeV\(^{-2}\) of the muon
  lifetime (Particle Data Group). Status: STANDARD.
\item
  \(m_\Psi\): The mass parameter of dirac16complex in the bulk. Unit:
  GeV. Status: RECORD.
\end{itemize}

The couplings of the particles \(\chi\) (the resonances of FE-40 of the
mass \(\omega_{\mathrm{r}}\), equal to \(\omega_{\mathrm{b}}\) up to the
relative correction \(H/(2q_{\mathrm{b}})\)) to the Standard Model, by
the theorem of Feynman and Hellmann: a slowly varying Standard-Model
field shifts the strengths \(\mu_X\) of FE-39 and with them the energy
of the state, so the coupling of the particle to a Standard-Model
operator is the derivative of its energy with respect to its mean value
(valid for momentum transfers small compared with \(q_{\mathrm{b}}\)).
\(\mathfrak{g}_{h\chi}\) is the Yukawa coupling of the Higgs boson \(h\)
to the particle (\(\Delta\langle\Phi^\dagger\Phi\rangle = vh\)), here
for the Higgs portal alone; it is fixed by the mass of the particle and
the bulk mass,
\(\mathfrak{g}_{h\chi} = -2q_{\mathrm{b}}^2/(\omega_{\mathrm{b}}v)\),
and it grows with the mass at a fixed ratio
\(q_{\mathrm{b}}/\omega_{\mathrm{b}}\): the heavy particles \(\chi\)
that UM-4 requires (FE-86) are strongly coupled to the Higgs boson
unless they are very weakly bound (FE-61). The couplings to \(W\) and
\(B\) follow like \(\mathfrak{g}_{G}\); the dipole portal gives the
particles magnetic and electric dipole moments. These couplings govern
the annihilation (FE-61), the scattering on nuclei and the production of
the particles \(\chi\). Status: DERIVED (the theorem checked with the
resonance level of UM-4: the left-handed density at the fold equals
\(\partial\omega/\partial\mu_{\mathrm{L}}\) to \(3\times10^{-5}\),
\texttt{um4\_kohnsham.py}, K6).

\hypertarget{the-three-neutrino-flavours}{%
\section{The three neutrino
flavours}\label{the-three-neutrino-flavours}}

The neutrinos are the upper components \(\nu_{\mathrm{L}}\) of the three
lepton doublets on the fold; with the three right-handed neutrinos they
have Dirac masses and the lepton mixing of FE-46. Their flavour
conversion follows from FE-22 and FE-24 for neutrinos that propagate
through a medium of electrons, nucleons, other neutrinos and
dirac16complex, by the standard reduction of the Dirac equations of
ultrarelativistic neutrinos to equations for their flavour content (the
mean-field quantum kinetic equations of Sigl and Raffelt, 1993). The
author asks for correct flavour conversion because it shapes the fate of
core-collapse supernovae: Gogilashvili and Tamborra (``Neutrino Flavor
Conversion Shapes the Rate of Failed Core-collapse Supernovae'',
arXiv:2605.16504v3) simulate the collapse of 195 stars of 9 to 120 solar
masses and find that flavour equipartition below a critical baryon
density raises the fraction of failed explosions from 25.6\% (without
flavour conversion) to 50.8\%, 77.4\%, 93.3\% and 96.4\% for critical
densities of \(10^9\), \(10^{10}\), \(10^{11}\) and \(10^{13}\)
g/cm\(^3\) (from 27.0\% to between 47.6\% and 88.3\% when the stars are
weighted with the Salpeter initial mass function). This section writes
the equations of motion of the flavour content and shows exactly which
terms UM-4 adds to those of the Standard Model.

\hypertarget{the-equations-of-motion-of-the-flavour-content}{%
\subsection{The equations of motion of the flavour
content}\label{the-equations-of-motion-of-the-flavour-content}}

\[
\begin{aligned}
&\mathrm{i}\big(\partial_{x_4} + \hat v\cdot\vec\nabla\big)\varrho_{\vec p} = \big[\mathsf{H}_{\vec p},\ \varrho_{\vec p}\big] + \mathrm{i}\,\mathcal{C}_{\vec p},\qquad
 \mathrm{i}\big(\partial_{x_4} + \hat v\cdot\vec\nabla\big)\bar\varrho_{\vec p} = \big[\bar{\mathsf{H}}_{\vec p},\ \bar\varrho_{\vec p}\big] + \mathrm{i}\,\bar{\mathcal{C}}_{\vec p},\\
&\mathsf{H}_{\vec p} = \frac{v_{\mathrm{loc}}^2}{v^2}\,\frac{U_{\mathrm{PMNS}}\,\mathsf{M}^2\,U_{\mathrm{PMNS}}^\dagger}{2E} + \mathsf{V}_{\mathrm{m}} + \mathsf{V}_{\nu\nu} + V_\Psi\mathbf{1},\qquad
 \mathsf{M}^2 = \mathrm{diag}\big(m_1^2,\ m_2^2,\ m_3^2\big),\\
&\bar{\mathsf{H}}_{\vec p} = -\frac{v_{\mathrm{loc}}^2}{v^2}\,\frac{U_{\mathrm{PMNS}}\,\mathsf{M}^2\,U_{\mathrm{PMNS}}^\dagger}{2E} + \mathsf{V}_{\mathrm{m}} + \mathsf{V}_{\nu\nu} + V_\Psi\mathbf{1}
\end{aligned} \tag{FE-48}
\] \textbf{Symbols in (FE-48).}

\begin{itemize}
\tightlist
\item
  \(\mathrm{i}\): The imaginary unit, \(\mathrm{i}^2 = -1\). Unit:
  dimensionless. Status: STANDARD.
\item
  \(\partial\): A partial derivative: the derivative with respect to one
  variable at fixed values of the others (\(\partial_\sigma\) with
  respect to \(\sigma\); \(\partial_M\) and \(\partial_\mu\) with
  respect to the coordinate that carries the same index, of the eight
  dimensions or of the fold). Status: STANDARD.
\item
  \(x_4\): The time; a prime on a function of \(x_4\) alone is
  \(\mathrm{d}/\mathrm{d}x_4\). Unit: length (= time, \(c = 1\)).
  Status: AUTHOR.
\item
  \(\hat v,\ \hat v^\mu\): The direction of motion of a neutrino: the
  unit vector \(\hat v = \vec p/\lvert\vec p\rvert\)
  (\(\hat v_{\vec p}\), \(\hat v_{\vec q}\) for the momenta \(\vec p\),
  \(\vec q\)), and the four-vector \(\hat v^\mu = (1, \hat v)\) in the
  order \((x_4; x_1, x_2, x_3)\); \(\mathrm{d}^2\hat v\) is the element
  of solid angle of the directions. Unit: dimensionless. Status:
  STANDARD.
\item
  \(\vec\nabla\): The gradient along ordinary space,
  \(\vec\nabla = (\partial_{x_1}, \partial_{x_2}, \partial_{x_3})\) (in
  the local frame of the medium). Unit: 1/length. Status: STANDARD.
\item
  \(\varrho_{\vec p},\ \bar\varrho_{\vec p}\): The \(3 \times 3\)
  flavour density matrices of the neutrinos and of the antineutrinos of
  momentum \(\vec p\) in the flavour basis (electron, muon, tau): the
  diagonal entries are the occupation numbers of the three flavours, the
  off-diagonal entries the coherence between the flavours;
  \(\bar\varrho\) is defined with the flavour indices in the reversed
  order, so that it obeys an equation of the same form. Unit:
  dimensionless. Status: STANDARD.
\item
  \(\vec p\): The momentum of a neutrino (three components along
  ordinary space, in the rest frame of the medium);
  \(\lvert\vec p\rvert = E\) for a neutrino of negligible mass. Unit:
  GeV. Status: STANDARD.
\item
  \(\mathsf{H}_{\vec p},\ \bar{\mathsf{H}}_{\vec p}\): The flavour
  Hamiltonians (\(3 \times 3\) Hermitian matrices) of the neutrinos and
  of the antineutrinos of momentum \(\vec p\). Unit: GeV. Status:
  STANDARD.
\item
  \(\mathcal{C}_{\vec p},\ \bar{\mathcal{C}}_{\vec p}\): The collision
  terms (\(3 \times 3\) matrices) of the neutrinos and the antineutrinos
  of momentum \(\vec p\): the rates of change of \(\varrho_{\vec p}\)
  and \(\bar\varrho_{\vec p}\) by emission, absorption, scattering and
  pair processes with the matter. Unit: GeV. Status: STANDARD.
\item
  \(v_{\mathrm{loc}}\): The local value of the electroweak vacuum,
  \(v_{\mathrm{loc}}^2 = 2\langle\Phi^\dagger\Phi\rangle\): it sets the
  local masses of the quarks and leptons (each mass times
  \(v_{\mathrm{loc}}/v\)) and the local Fermi constant;
  \(v_{\mathrm{loc}} = v\) in the Standard-Model vacuum. Unit: GeV.
  Value: \(v\) in the late universe outside dense dirac16complex.
  Status: DERIVED.
\item
  \(v\): The vacuum expectation value of the lower component of the
  Higgs doublet, \(v^2 = -\mu^2/\lambda\). Unit: GeV. Value: 246.22 GeV
  \(= (\sqrt2\,G_{\mathrm{F}})^{-1/2}\) with the Fermi constant
  \(G_{\mathrm{F}} = 1.1663788\times10^{-5}\) GeV\(^{-2}\) of the muon
  lifetime (Particle Data Group). Status: STANDARD.
\item
  \(U_{\mathrm{PMNS}}\): The lepton mixing matrix
  (Pontecorvo-Maki-Nakagawa-Sakata) of the three Dirac neutrinos; its
  entries \((U_{\mathrm{PMNS}})_{\alpha i}\) connect the flavour state
  \(\nu_\alpha\) with the mass state \(\nu_i\). Unit: dimensionless.
  Value: \(\theta_{12} = 33.40^\circ\), \(\theta_{23} = 49.14^\circ\),
  \(\theta_{13} = 8.54^\circ\) (NuFIT 5.2, November 2022, normal
  ordering, the fit without the Super-Kamiokande atmospheric data); the
  Dirac phase is set to zero (ASSUMED; the fit prefers \(197^\circ\)).
  Status: STANDARD, ASSUMED (the phase).
\item
  \({}^\dagger\): Hermitian conjugation: the complex conjugate of the
  transpose of a matrix or spinor; for an operator on the space of
  states, its adjoint with respect to the scalar product of that space
  (for the quantised dirac16complex the indefinite Krein scalar
  product). Status: STANDARD.
\item
  \(\mathsf{M}^2\): The diagonal matrix of the squared masses of the
  three neutrino mass states,
  \(\mathsf{M}^2 = \mathrm{diag}(m_1^2, m_2^2, m_3^2)\) (\(\mathsf{M}\)
  appears only squared). Unit: GeV\(^2\) (of \(\mathsf{M}^2\)). Status:
  STANDARD.
\item
  \(E\): An energy: \(E_o\) the energy of the Kohn-Sham orbital \(o\);
  in the neutrino sector the energy of a neutrino. Unit: GeV. Status:
  STANDARD.
\item
  \(\mathsf{V}_{\mathrm{m}}\): The matter potential of the neutrinos (a
  \(3 \times 3\) matrix in flavour space) from their forward scattering
  on the electrons, positrons and neutrons of the medium. Unit: GeV.
  Status: STANDARD.
\item
  \(\mathsf{V}_{\nu\nu}\): The neutrino-neutrino potential (a
  \(3 \times 3\) matrix in flavour space) from the forward scattering of
  a neutrino on the other neutrinos and antineutrinos of the medium; it
  depends on their flavour content and their directions. Unit: GeV.
  Status: STANDARD.
\item
  \(V_\Psi\): The potential energy that dirac16complex gives every
  neutrino through the current portal of the lepton doublets,
  \(V_\Psi = \mathcal{V}^{\ell}_{x_4}\) (the same for the three
  flavours). Unit: GeV. Status: DERIVED.
\item
  \(\mathbf{1},\ \mathbf{1}_4,\ \mathbf{1}_8,\ \mathbf{1}_{16}\): The
  identity matrix of the size fixed by the context (\(16 \times 16\) on
  dirac16complex, \(4 \times 4\) on a four-dimensional spinor,
  \(3 \times 3\) in generation or flavour space); \(\mathbf{1}_4\),
  \(\mathbf{1}_8\) and \(\mathbf{1}_{16}\) are the identity matrices of
  4, 8 and 16 rows. Unit: dimensionless. Status: STANDARD.
\item
  \(\mathrm{diag}\): A diagonal matrix with the listed entries on its
  diagonal (in the order of the coordinates \(x_0, x_1, \dots, x_7\) for
  a tensor of eight dimensions, \(x_1, \dots, x_7\) for a tensor of the
  fold, \(x_1, \dots, x_4\) for a four-dimensional tensor). Status:
  STANDARD.
\item
  \(m_u, \dots, m_3\): The masses of the quarks (\(m_u, m_c, m_t\);
  \(m_d, m_s, m_b\)), of the charged leptons (\(m_e, m_\mu, m_\tau\))
  and of the three neutrino mass states (\(m_1, m_2, m_3\), Dirac
  masses). Unit: GeV. Value: measured values;
  \(\Delta m_{21}^2 = 7.41\times10^{-5}\) eV\(^2\),
  \(\lvert\Delta m_{31}^2\rvert = 2.511\times10^{-3}\) eV\(^2\) in the
  computations of UM-4 (the global fit NuFIT 5.2 of the oscillation
  data, November 2022), \(7.49\times10^{-5}\) and \(2.513\times10^{-3}\)
  eV\(^2\) in NuFIT 6.0 (2024); the lightest mass \(m_1 = 0\) in the
  computations, below 23 meV by DESI DR2 with the CMB and NuFIT 6.0.
  Status: STANDARD.
\end{itemize}

The equations of motion of the \(3 \times 3\) flavour density matrices
of the neutrinos (\(\varrho_{\vec p}\)) and of the antineutrinos
(\(\bar\varrho_{\vec p}\), defined with the flavour indices in the
reversed order) of momentum \(\vec p\), energy
\(E = \lvert\vec p\rvert\) and direction \(\hat v\); the diagonal
entries are the occupation numbers of the three flavours, the
off-diagonal entries the coherence between them. Free streaming stands
on the left, the commutator with the flavour Hamiltonian and the
collisions \(\mathcal{C}\) (emission, absorption, scattering and pair
processes) on the right. The Hamiltonian contains the vacuum term of the
three Dirac masses (with the opposite sign for antineutrinos in this
convention), multiplied by \(v_{\mathrm{loc}}^2/v^2\) because the masses
are the Yukawa couplings times the local value of the Higgs field
(FE-49); the matter potential; the neutrino-neutrino potential; and the
potential \(V_\Psi\) of dirac16complex. Status: STANDARD (UM-4 adds
\(V_\Psi\) and the factor \(v_{\mathrm{loc}}^2/v^2\)).

\[
\begin{aligned}
&\mathsf{V}_{\mathrm{m}} = \sqrt2\,G_{\mathrm{F,loc}}\,\mathrm{diag}\big(n_e - n_{\bar e},\ 0,\ 0\big) - \frac{G_{\mathrm{F,loc}}}{\sqrt2}\,n_n\mathbf{1},\qquad
 \mathsf{V}_{\nu\nu} = \sqrt2\,G_{\mathrm{F,loc}}\int\frac{\mathrm{d}^3\vec q}{(2\pi)^3}\big(1 - \hat v_{\vec p}\cdot\hat v_{\vec q}\big)\big(\varrho_{\vec q} - \bar\varrho_{\vec q}\big),\\
&V_\Psi = \mathcal{V}^{\ell}_{x_4} = \sum_X c^{\ell}_X\,j_{X\,x_4},\qquad G_{\mathrm{F,loc}} = \frac{1}{\sqrt2\,v_{\mathrm{loc}}^2},\\
&v_{\mathrm{loc}}^2 = 2\langle\Phi^\dagger\Phi\rangle,\qquad v_{\mathrm{loc}}^2 = v^2 + \mathcal{A}^2\ \ \text{(the ad hoc vacuum)},\\
&v_{\mathrm{loc}}^2 = v^2 - \frac{c_\Phi}{\lambda}\,\langle\bar\Psi P_{\mathrm{L}}\Psi\rangle\ \ \text{(the relaxed vacuum with dirac16complex)}
\end{aligned} \tag{FE-49}
\] \textbf{Symbols in (FE-49).}

\begin{itemize}
\tightlist
\item
  \(\mathsf{V}_{\mathrm{m}}\): The matter potential of the neutrinos (a
  \(3 \times 3\) matrix in flavour space) from their forward scattering
  on the electrons, positrons and neutrons of the medium. Unit: GeV.
  Status: STANDARD.
\item
  \(G_{\mathrm{F,loc}}\): The local Fermi constant,
  \(G_{\mathrm{F,loc}} = 1/(\sqrt2\,v_{\mathrm{loc}}^2)\)
  (\(= G_{\mathrm{F}}\) where \(v_{\mathrm{loc}} = v\)). Unit:
  GeV\(^{-2}\). Status: DERIVED.
\item
  \(\mathrm{diag}\): A diagonal matrix with the listed entries on its
  diagonal (in the order of the coordinates \(x_0, x_1, \dots, x_7\) for
  a tensor of eight dimensions, \(x_1, \dots, x_7\) for a tensor of the
  fold, \(x_1, \dots, x_4\) for a four-dimensional tensor). Status:
  STANDARD.
\item
  \(n_e\): The number density of electrons (bound and free; in a
  supernova all are free, in the recombining universe
  \(x_{\mathrm{e}}n_e\) is the density of the free electrons). Unit:
  GeV\(^3\) (cm\(^{-3}\)). Status: STANDARD.
\item
  \(n_{\bar e}\): The number density of positrons. Unit: GeV\(^3\).
  Status: STANDARD.
\item
  \(n_n\): The number density of neutrons. Unit: GeV\(^3\). Status:
  STANDARD.
\item
  \(\mathbf{1},\ \mathbf{1}_4,\ \mathbf{1}_8,\ \mathbf{1}_{16}\): The
  identity matrix of the size fixed by the context (\(16 \times 16\) on
  dirac16complex, \(4 \times 4\) on a four-dimensional spinor,
  \(3 \times 3\) in generation or flavour space); \(\mathbf{1}_4\),
  \(\mathbf{1}_8\) and \(\mathbf{1}_{16}\) are the identity matrices of
  4, 8 and 16 rows. Unit: dimensionless. Status: STANDARD.
\item
  \(\mathsf{V}_{\nu\nu}\): The neutrino-neutrino potential (a
  \(3 \times 3\) matrix in flavour space) from the forward scattering of
  a neutrino on the other neutrinos and antineutrinos of the medium; it
  depends on their flavour content and their directions. Unit: GeV.
  Status: STANDARD.
\item
  \(\mathrm{d}\): The differential: \(\mathrm{d}\) in front of a
  quantity is an infinitesimal change of it (as \(\mathrm{d}\sigma\), a
  small step of \(\sigma\)); the quotient of two such changes is a
  derivative; \(\mathrm{d}\) with an exponent, in front of several
  variables, is their volume element, the product of their
  differentials. Status: STANDARD.
\item
  \(\vec q\): The momentum of another neutrino of the medium (the
  integration variable of the neutrino-neutrino potential);
  \(\mathrm{d}^3\vec q\) its volume element. Unit: GeV. Status:
  STANDARD.
\item
  \(\pi\): The number pi, the ratio of the circumference of a circle to
  its diameter. Unit: dimensionless. Value: 3.141592654. Status:
  STANDARD.
\item
  \(\hat v,\ \hat v^\mu\): The direction of motion of a neutrino: the
  unit vector \(\hat v = \vec p/\lvert\vec p\rvert\)
  (\(\hat v_{\vec p}\), \(\hat v_{\vec q}\) for the momenta \(\vec p\),
  \(\vec q\)), and the four-vector \(\hat v^\mu = (1, \hat v)\) in the
  order \((x_4; x_1, x_2, x_3)\); \(\mathrm{d}^2\hat v\) is the element
  of solid angle of the directions. Unit: dimensionless. Status:
  STANDARD.
\item
  \(\vec p\): The momentum of a neutrino (three components along
  ordinary space, in the rest frame of the medium);
  \(\lvert\vec p\rvert = E\) for a neutrino of negligible mass. Unit:
  GeV. Status: STANDARD.
\item
  \(\varrho_{\vec p},\ \bar\varrho_{\vec p}\): The \(3 \times 3\)
  flavour density matrices of the neutrinos and of the antineutrinos of
  momentum \(\vec p\) in the flavour basis (electron, muon, tau): the
  diagonal entries are the occupation numbers of the three flavours, the
  off-diagonal entries the coherence between the flavours;
  \(\bar\varrho\) is defined with the flavour indices in the reversed
  order, so that it obeys an equation of the same form. Unit:
  dimensionless. Status: STANDARD.
\item
  \(V_\Psi\): The potential energy that dirac16complex gives every
  neutrino through the current portal of the lepton doublets,
  \(V_\Psi = \mathcal{V}^{\ell}_{x_4}\) (the same for the three
  flavours). Unit: GeV. Status: DERIVED.
\item
  \(\mathcal{V}^{f}_\mu\): The vector potential that dirac16complex
  exerts on the Standard-Model fermions of kind \(f\) through the
  current portals, \(\mathcal{V}^{f}_\mu = \sum_Xc^{f}_Xj_{X\mu}\)
  (\(\mathcal{V}^{Q}, \mathcal{V}^{u}, \dots, \mathcal{V}^{\nu}\) for
  the six kinds); flavour blind. Unit: GeV. Status: DERIVED.
\item
  \(x_4\): The time; a prime on a function of \(x_4\) alone is
  \(\mathrm{d}/\mathrm{d}x_4\). Unit: length (= time, \(c = 1\)).
  Status: AUTHOR.
\item
  \(X\): The label of a half of dirac16complex:
  \(X \in \{\mathrm{L}, \mathrm{R}\}\) (left-handed or right-handed
  half); a sum over \(X\) runs over both. Status: STANDARD.
\item
  \(c^{f}_X\): The current-portal coupling of the half \(X\) of
  dirac16complex with the fermion-number current of the multiplets of
  kind \(f\) (flavour blind); \(c^{\ell}_X\) and \(c^{\nu}_X\) are those
  of the lepton doublets and of the right-handed neutrinos. Unit:
  GeV\(^{-7}\). Status: INVENTED.
\item
  \(j_X^\mu\): The convective current of the half \(X\) of
  dirac16complex,
  \(j_X^\mu = \mathrm{i}(\bar\Psi P_XD^\mu\Psi - (D^\mu\bar\Psi)P_X\Psi)\)
  (real and joining the half with itself, DERIVED). Unit: GeV\(^8\).
  Status: DERIVED.
\item
  \(v_{\mathrm{loc}}\): The local value of the electroweak vacuum,
  \(v_{\mathrm{loc}}^2 = 2\langle\Phi^\dagger\Phi\rangle\): it sets the
  local masses of the quarks and leptons (each mass times
  \(v_{\mathrm{loc}}/v\)) and the local Fermi constant;
  \(v_{\mathrm{loc}} = v\) in the Standard-Model vacuum. Unit: GeV.
  Value: \(v\) in the late universe outside dense dirac16complex.
  Status: DERIVED.
\item
  \(\Phi\): The Higgs doublet (two complex components, hypercharge
  \(1/2\)), a field of the fold. Unit: GeV. Status: AUTHOR.
\item
  \({}^\dagger\): Hermitian conjugation: the complex conjugate of the
  transpose of a matrix or spinor; for an operator on the space of
  states, its adjoint with respect to the scalar product of that space
  (for the quantised dirac16complex the indefinite Krein scalar
  product). Status: STANDARD.
\item
  \(\langle\cdots\rangle\): The mean (expectation) value of an operator
  in the state of the medium (the vacuum, a plasma, the Kohn-Sham
  state). Status: STANDARD.
\item
  \(v\): The vacuum expectation value of the lower component of the
  Higgs doublet, \(v^2 = -\mu^2/\lambda\). Unit: GeV. Value: 246.22 GeV
  \(= (\sqrt2\,G_{\mathrm{F}})^{-1/2}\) with the Fermi constant
  \(G_{\mathrm{F}} = 1.1663788\times10^{-5}\) GeV\(^{-2}\) of the muon
  lifetime (Particle Data Group). Status: STANDARD.
\item
  \(\mathcal{A}\): The author's ad hoc upper component of the Higgs
  vacuum, almightyAdHoc,
  \(\mathcal{A}(x_4) = v_{\mathcal{A}}\mathrm{e}^{-2a_{40}x_4}\). Unit:
  GeV. Status: AUTHOR.
\item
  \(c_\Phi\): The Higgs-portal coupling of the left-handed half of
  dirac16complex (the right-handed half has none). Unit: GeV\(^{-5}\).
  Status: INVENTED.
\item
  \(\lambda\): The quartic coupling of the Higgs potential,
  \(\lambda > 0\). Unit: dimensionless. Value:
  \(\lambda = m_h^2/(2v^2) = 0.129\). Status: AUTHOR.
\item
  \(\Psi,\ \bar\Psi\): dirac16complex: the author's 16-component complex
  spinor field of eight dimensions;
  \(\bar\Psi = \Psi^\dagger\sigma_{16}\) is its Dirac adjoint; in a term
  of the fold, \(\Psi\) is its value at \(\sigma = 0\). Unit:
  GeV\(^{7/2}\). Status: AUTHOR.
\item
  \(P_{\mathrm{L}}\): The author's left-handed chiral projector,
  \(P_{\mathrm{L}} = (\mathbf{1} - \Gamma_8)/2\). Unit: dimensionless.
  Status: AUTHOR.
\end{itemize}

The potentials and the local electroweak vacuum. Forward scattering on
electrons and positrons through \(W\) exchange gives the electron
flavour the potential \(\sqrt2G_{\mathrm{F,loc}}(n_e - n_{\bar e})\);
neutral currents give all flavours the same potential, here from the
neutrons (the neutral-current potentials of electrons and protons cancel
in neutral matter for any weak mixing angle); forward scattering on the
other neutrinos gives the neutrino-neutrino potential, which depends on
the flavour content and the directions of all other neutrinos and makes
the equations non-linear. dirac16complex acts on the neutrinos in
exactly two ways, both fixed by the action (IN-16):

\begin{itemize}
\tightlist
\item
  through the current portal of the lepton doublets: \(V_\Psi\) is the
  time component of the vector potential \(\mathcal{V}^{\ell}_\mu\) of
  FE-22, the same for the three flavours because the portals are flavour
  blind. A multiple of the unit matrix drops out of every commutator:
  \(V_\Psi\) shifts the energies of the neutrinos but does not change
  their flavour;
\item
  through the Higgs portal: the scalar density of the left-handed half
  shifts the mass parameter of the Higgs field (FE-15). Where the Higgs
  field relaxes to the minimum of its potential (a homogeneous, static
  vacuum with \(\mathcal{A} = 0\) in FE-15), its local value is
  \(v_{\mathrm{loc}}^2 = v^2 - c_\Phi\langle\bar\Psi P_{\mathrm{L}}\Psi\rangle/\lambda\);
  in the author's ad hoc vacuum
  \(v_{\mathrm{loc}}^2 = v^2 + \mathcal{A}^2\) (FE-45). All masses scale
  with \(v_{\mathrm{loc}}/v\) and the Fermi constant with
  \(v^2/v_{\mathrm{loc}}^2\) (at tree level
  \(G_{\mathrm{F}} = 1/(\sqrt2 v^2)\) does not depend on the gauge
  couplings, so the gauge-kinetic portals do not change it), while the
  Yukawa matrices, and with them \(U_{\mathrm{PMNS}}\), are unchanged.
\end{itemize}

So dirac16complex changes neither the mixing angles nor the flavour
pattern of the conversion: it rescales the vacuum oscillation
frequencies by \(v_{\mathrm{loc}}^2/v^2\) and the potentials by
\(v^2/v_{\mathrm{loc}}^2\) (the resonance density of FE-50 by
\(v_{\mathrm{loc}}^4/v^4\)). In the late universe \(\mathcal{A} = 0\),
and the dark matter is far too dilute to move \(v_{\mathrm{loc}}\): with
FE-47 the shift is
\(\Delta v^2/v^2 = -\mathfrak{g}_{h\chi}\rho_\chi/(\lambda v^3\omega_{\mathrm{b}})\),
about \(10^{-50}\,\mathfrak{g}_{h\chi}\) for the local dark-matter
density \(0.3\) GeV/cm\(^3\) and \(\omega_{\mathrm{b}} = 100\) GeV, and
less for the heavier particles \(\chi\) that UM-4 requires (FE-86). The
flavour conversion of UM-4 is therefore that of the Standard Model with
three Dirac neutrinos wherever \(v_{\mathrm{loc}} = v\), which includes
every supernova and the Sun (DERIVED). The Dirac neutrinos of UM-4 have
the magnetic moments of the Standard Model, about \(3.2\times10^{-19}\)
Bohr magnetons times the mass in eV, too small to turn left-handed into
right-handed neutrinos in the magnetic fields of supernovae; the dipole
portal does not act on neutrinos directly. Status: STANDARD (the
potentials), DERIVED (\(V_\Psi\), \(v_{\mathrm{loc}}\) and their
consequences; \texttt{um4\_neutrino.py}, N1).

\hypertarget{matter-enhanced-conversion}{%
\subsection{Matter-enhanced
conversion}\label{matter-enhanced-conversion}}

\[
\begin{aligned}
&\tan2\theta_{\mathrm{m}} = \frac{\sin2\theta}{\cos2\theta - 2EV_{\mathrm{CC}}/\Delta m^2},\qquad
 V_{\mathrm{CC}} = \sqrt2\,G_{\mathrm{F,loc}}\,n_e,\qquad
 n_{e,\mathrm{res}} = \frac{\Delta m^2\cos2\theta}{2\sqrt2\,G_{\mathrm{F,loc}}\,E},\\
&P_{\mathrm{LZ}} = \exp\Big(-\frac{\pi}{2}\,\frac{\Delta m^2\sin^22\theta}{2E\cos2\theta}\,\Big\lvert\frac{\mathrm{d}\ln n_e}{\mathrm{d}r}\Big\rvert^{-1}_{\mathrm{res}}\Big)
\end{aligned} \tag{FE-50}
\] \textbf{Symbols in (FE-50).}

\begin{itemize}
\tightlist
\item
  \(\theta_{\mathrm{m}}\): The mixing angle of two neutrino flavours in
  matter. Unit: dimensionless. Status: STANDARD.
\item
  \(\theta\): The vacuum mixing angle of two neutrino flavours (in the
  two-flavour description of a resonance, one of the three angles
  \(\theta_{12}, \theta_{13}, \theta_{23}\) of \(U_{\mathrm{PMNS}}\)).
  Unit: dimensionless. Status: STANDARD.
\item
  \(E\): An energy: \(E_o\) the energy of the Kohn-Sham orbital \(o\);
  in the neutrino sector the energy of a neutrino. Unit: GeV. Status:
  STANDARD.
\item
  \(V_{\mathrm{CC}}\): The charged-current matter potential of the
  electron neutrinos (forward scattering on electrons through \(W\)
  exchange). Unit: GeV. Status: STANDARD.
\item
  \(\Delta\): A change, difference or small variation of the quantity
  that follows (for example an interval of the time, a perturbation of a
  density, a difference of two squared masses, or an addition to an
  effective number of species). Status: STANDARD.
\item
  \(m\): A mass: \(m_i\) the mass of the neutrino mass state \(i\);
  \(\Delta m_{ij}^2 = m_i^2 - m_j^2\) the squared-mass difference of the
  states \(i\) and \(j\) (\(\Delta m_{21}^2\), \(\Delta m_{31}^2\),
  \(\Delta m_{32}^2\)); \(\Delta m^2\) without labels the difference of
  the two states of a resonance, \(m_2^2 - m_1^2\), multiplied by
  \(v_{\mathrm{loc}}^2/v^2\) where this differs from 1; \(m\) without a
  label the mass of a mode of a field of the bulk as the fold sees it,
  \(m^2 = \omega^2 - k^2\); for the W and Z bosons their complex pole
  mass, \(m^2 = m_0^2 - \mathrm{i}m_0\Gamma_{\mathrm{bulk}}\) to the
  leading order; in the relic abundance of the particles \(\chi\) their
  mass. Unit: GeV. Status: STANDARD.
\item
  \(G_{\mathrm{F,loc}}\): The local Fermi constant,
  \(G_{\mathrm{F,loc}} = 1/(\sqrt2\,v_{\mathrm{loc}}^2)\)
  (\(= G_{\mathrm{F}}\) where \(v_{\mathrm{loc}} = v\)). Unit:
  GeV\(^{-2}\). Status: DERIVED.
\item
  \(n_e\): The number density of electrons (bound and free; in a
  supernova all are free, in the recombining universe
  \(x_{\mathrm{e}}n_e\) is the density of the free electrons). Unit:
  GeV\(^3\) (cm\(^{-3}\)). Status: STANDARD.
\item
  \(n_{e,\mathrm{res}}\): The electron density at which the flavour
  conversion in matter is resonant (the Mikheyev-Smirnov-Wolfenstein
  resonance). Unit: GeV\(^3\). Status: STANDARD.
\item
  \(P_{\mathrm{LZ}}\): The probability that a neutrino jumps between the
  eigenstates of matter at the resonance (Landau-Zener);
  \(P_{\mathrm{LZ}} = 0\) is an adiabatic resonance. Unit:
  dimensionless. Status: STANDARD.
\item
  \(\pi\): The number pi, the ratio of the circumference of a circle to
  its diameter. Unit: dimensionless. Value: 3.141592654. Status:
  STANDARD.
\item
  \(\mathrm{d}\): The differential: \(\mathrm{d}\) in front of a
  quantity is an infinitesimal change of it (as \(\mathrm{d}\sigma\), a
  small step of \(\sigma\)); the quotient of two such changes is a
  derivative; \(\mathrm{d}\) with an exponent, in front of several
  variables, is their volume element, the product of their
  differentials. Status: STANDARD.
\item
  \(r\): A distance from a centre: the distance from the centre of the
  star along the path of a neutrino; in Newton's and Coulomb's laws the
  distance between two point sources on the fold. Unit: length. Status:
  STANDARD.
\item
  \(\mathrm{res}\): The label of a quantity evaluated at the resonance,
  where \(n_e = n_{e,\mathrm{res}}\). Status: STANDARD.
\end{itemize}

The resonant flavour conversion in matter (Mikheyev-Smirnov-Wolfenstein)
for two flavours with the vacuum mixing angle \(\theta\) and the
squared-mass difference \(\Delta m^2\) (which includes the factor
\(v_{\mathrm{loc}}^2/v^2\) of FE-48): the mixing angle in matter
\(\theta_{\mathrm{m}}\) (FE-48 with the diagonal potential
\(V_{\mathrm{CC}}\) in the electron flavour), the electron density of
the resonance, where \(\theta_{\mathrm{m}} = 45^\circ\), and the
probability \(P_{\mathrm{LZ}}\) that the neutrino jumps between the
eigenstates of matter at the resonance for the density profile
\(n_e(r)\) along its path (the Landau-Zener formula, exact for a linear
profile). With the measured masses and mixing angles this gives the
survival probability of solar neutrinos and the conversion in the outer
layers of supernovae (the resonances of \(\Delta m^2_{31}\) and
\(\Delta m^2_{21}\)). Status: STANDARD (\texttt{um4\_neutrino.py}, N2).

\hypertarget{the-fast-flavour-instability-of-dense-neutrino-gases}{%
\subsection{The fast flavour instability of dense neutrino
gases}\label{the-fast-flavour-instability-of-dense-neutrino-gases}}

\[
\begin{aligned}
&\mathcal{G}_{\mathrm{ELN}}(\hat v) = \sqrt2\,G_{\mathrm{F,loc}}\int_0^\infty\frac{E^2\,\mathrm{d}E}{(2\pi)^3}\Big[\mathcal{F}_{e} - \bar{\mathcal{F}}_{e}
 - \big(\mathcal{F}_{x} - \bar{\mathcal{F}}_{x}\big)\Big],\\
&\det\Big[\eta_{\mathrm{M}}^{\mu\nu} + \int\mathrm{d}^2\hat v\;\mathcal{G}_{\mathrm{ELN}}(\hat v)\,
 \frac{\hat v^\mu\hat v^\nu}{\Omega_{\mathrm{fw}} - \vec K\cdot\hat v}\Big] = 0,\qquad \hat v^\mu = (1,\ \hat v)
\end{aligned} \tag{FE-51}
\] \textbf{Symbols in (FE-51).}

\begin{itemize}
\tightlist
\item
  \(\mathcal{G}_{\mathrm{ELN}}(\hat v)\): The angular distribution of
  the electron lepton number carried by the neutrinos minus that carried
  by the heavy-lepton neutrinos, per unit solid angle of the direction
  \(\hat v\), multiplied by \(\sqrt2G_{\mathrm{F,loc}}\). Unit: GeV.
  Status: STANDARD.
\item
  \(\hat v,\ \hat v^\mu\): The direction of motion of a neutrino: the
  unit vector \(\hat v = \vec p/\lvert\vec p\rvert\)
  (\(\hat v_{\vec p}\), \(\hat v_{\vec q}\) for the momenta \(\vec p\),
  \(\vec q\)), and the four-vector \(\hat v^\mu = (1, \hat v)\) in the
  order \((x_4; x_1, x_2, x_3)\); \(\mathrm{d}^2\hat v\) is the element
  of solid angle of the directions. Unit: dimensionless. Status:
  STANDARD.
\item
  \(G_{\mathrm{F,loc}}\): The local Fermi constant,
  \(G_{\mathrm{F,loc}} = 1/(\sqrt2\,v_{\mathrm{loc}}^2)\)
  (\(= G_{\mathrm{F}}\) where \(v_{\mathrm{loc}} = v\)). Unit:
  GeV\(^{-2}\). Status: DERIVED.
\item
  \(\infty\): Infinity: as a limit of an integral or a sum it means that
  the range has no bound. Status: STANDARD.
\item
  \(E\): An energy: \(E_o\) the energy of the Kohn-Sham orbital \(o\);
  in the neutrino sector the energy of a neutrino. Unit: GeV. Status:
  STANDARD.
\item
  \(\mathrm{d}\): The differential: \(\mathrm{d}\) in front of a
  quantity is an infinitesimal change of it (as \(\mathrm{d}\sigma\), a
  small step of \(\sigma\)); the quotient of two such changes is a
  derivative; \(\mathrm{d}\) with an exponent, in front of several
  variables, is their volume element, the product of their
  differentials. Status: STANDARD.
\item
  \(\pi\): The number pi, the ratio of the circumference of a circle to
  its diameter. Unit: dimensionless. Value: 3.141592654. Status:
  STANDARD.
\item
  \(\mathcal{F}_{e},\ \bar{\mathcal{F}}_{e},\ \mathcal{F}_{x},\ \bar{\mathcal{F}}_{x}\):
  The occupation numbers of the electron neutrinos, of the electron
  antineutrinos, and of one heavy-lepton neutrino and antineutrino
  (\(x = \mu\) or \(\tau\)), as functions of the energy \(E\) and the
  direction \(\hat v\) (the diagonal entries of \(\varrho_{\vec p}\) and
  \(\bar\varrho_{\vec p}\)); in a transport code, the same for the
  number of particles in an energy bin. Unit: dimensionless. Status:
  STANDARD.
\item
  \(\eta_{\mathrm{M}}^{\mu\nu}\): The Minkowski metric of the
  four-vectors \(\hat v^\mu\),
  \(\eta_{\mathrm{M}} = \mathrm{diag}(1, -1, -1, -1)\) in the order
  \((x_4; x_1, x_2, x_3)\). Unit: dimensionless. Status: STANDARD.
\item
  \(\mu\): A four-dimensional coordinate index running over the
  observable coordinates \(x_1, x_2, x_3, x_4\) of the fold (signature
  \(+,+,+,-\)); a repeated index is summed. Status: STANDARD.
\item
  \(\nu\): A four-dimensional coordinate index (as \(\mu\)). Status:
  STANDARD.
\item
  \(\Omega_{\mathrm{fw}}\): The (complex) frequency of a flavour wave of
  the neutrino gas, measured from the shift caused by the potentials; an
  imaginary part greater than zero is an instability. Unit: GeV. Status:
  STANDARD.
\item
  \(\vec K\): The (real) wave vector of a flavour wave of the neutrino
  gas, measured from the shift caused by the potentials. Unit: GeV.
  Status: STANDARD.
\end{itemize}

The fast flavour instability (STANDARD: FE-48 linearised in the flavour
coherence, without the vacuum term, which is negligible in the dense
gas; Izaguirre, Raffelt and Tamborra 2017).
\(\mathcal{G}_{\mathrm{ELN}}\) is the angular distribution of the
electron lepton number carried by the neutrinos minus that of the
heavy-lepton flavours (\(\mathcal{F}\) the occupation numbers of
\(\nu_e\), \(\bar\nu_e\) and of one heavy flavour \(\nu_x\),
\(\bar\nu_x\), as functions of the energy and the direction;
\(\mathrm{d}^2\hat v\) the element of solid angle). A small flavour
coherence
\(\propto \mathrm{e}^{-\mathrm{i}(\Omega_{\mathrm{fw}}x_4 - \vec K\cdot\vec x)}\)
grows exponentially if the dispersion relation (the second line, written
with the Minkowski metric \(\eta_{\mathrm{M}}\) of the four-vectors
\(\hat v^\mu\)) has a solution with an imaginary part of
\(\Omega_{\mathrm{fw}}\) greater than zero for a real wave vector
\(\vec K\) (the frequency and the wave vector are measured from the
shifts that the matter and the neutrino-neutrino potentials cause, which
drop out). This happens if and only if \(\mathcal{G}_{\mathrm{ELN}}\)
changes sign as a function of the direction, an angular crossing
(Morinaga 2022). The instability then grows on lengths of the order of
\(1/(\sqrt2G_{\mathrm{F}}n_\nu)\), centimetres in a supernova core, much
shorter than the scales of the explosion, and drives the flavours toward
equipartition; transport codes therefore model its outcome by a
prescription such as FE-52. Status: STANDARD (\texttt{um4\_neutrino.py},
N3).

\[
\begin{aligned}
&\mathcal{F}_{e} \to \mathcal{F}_{e} - \delta_{\mathrm{FC}},\qquad \bar{\mathcal{F}}_{e} \to \bar{\mathcal{F}}_{e} - \delta_{\mathrm{FC}},\qquad
 \mathcal{F}_{x} = \bar{\mathcal{F}}_{x} \to \mathcal{F}_{x} + \frac{\delta_{\mathrm{FC}}}{2},\\
&\delta_{\mathrm{FC}} = \max\Big(\frac23\Big(\min\big(\mathcal{F}_{e},\ \bar{\mathcal{F}}_{e}\big) - \mathcal{F}_{x}\Big),\ \max\big(\mathcal{F}_{e},\ \bar{\mathcal{F}}_{e}\big) - 1\Big),\\
&\text{applied where}\ \ \rho_{\mathrm{b}} < \rho_{\mathrm{FC}}\ \ \text{and}\ \ x_4 - x_{4,\mathrm{b}} > t_{\mathrm{FC}}
\end{aligned} \tag{FE-52}
\] \textbf{Symbols in (FE-52).}

\begin{itemize}
\tightlist
\item
  \(\mathcal{F}_{e},\ \bar{\mathcal{F}}_{e},\ \mathcal{F}_{x},\ \bar{\mathcal{F}}_{x}\):
  The occupation numbers of the electron neutrinos, of the electron
  antineutrinos, and of one heavy-lepton neutrino and antineutrino
  (\(x = \mu\) or \(\tau\)), as functions of the energy \(E\) and the
  direction \(\hat v\) (the diagonal entries of \(\varrho_{\vec p}\) and
  \(\bar\varrho_{\vec p}\)); in a transport code, the same for the
  number of particles in an energy bin. Unit: dimensionless. Status:
  STANDARD.
\item
  \(\delta_{\mathrm{FC}}\): The occupation (or number per energy bin) of
  electron neutrino-antineutrino pairs converted into heavy-lepton pairs
  by the flavour equipartition (the conversion that equalises the
  occupations of the three flavours where the fast flavour instability
  acts); a negative value converts heavy-lepton pairs into electron
  pairs. Unit: dimensionless. Status: ASSUMED.
\item
  \(\rho_{\mathrm{b}}\): The (rest-mass) energy density of the baryons:
  in a supernova the local density of the matter (quoted in g/cm\(^3\)),
  in cosmology the mean density of the baryons
  (\(\rho_{\mathrm{b}}^{(0)}\) today). Unit: GeV\(^4\). Status:
  STANDARD.
\item
  \(\rho_{\mathrm{FC}}\): The baryon density below which the flavour
  equipartition is applied (the critical density \(\rho_{\mathrm{c}}\)
  of Gogilashvili and Tamborra). Unit: GeV\(^4\) (g/cm\(^3\)). Value:
  \(10^9\), \(10^{10}\), \(10^{11}\) or \(10^{13}\) g/cm\(^3\). Status:
  STANDARD.
\item
  \(x_4\): The time; a prime on a function of \(x_4\) alone is
  \(\mathrm{d}/\mathrm{d}x_4\). Unit: length (= time, \(c = 1\)).
  Status: AUTHOR.
\item
  \(x_{4,\mathrm{b}}\): The time of the bounce of the collapsing stellar
  core. Unit: length (= time). Status: STANDARD.
\item
  \(t_{\mathrm{FC}}\): The time after the bounce from which the flavour
  equipartition is applied. Unit: time. Value: 0.02 s. Status: STANDARD.
\end{itemize}

The flavour equipartition that models the outcome of the fast
instability in the transport of a core-collapse supernova, written in
the form that satisfies the constraints stated by Gogilashvili and
Tamborra for their scheme: in every energy bin (and in the same way for
the flux of every energy bin) electron neutrino-antineutrino pairs are
converted into heavy-lepton pairs, which conserves the electron lepton
number \(\mathcal{F}_{e} - \bar{\mathcal{F}}_{e}\), the number of
neutrinos \(\mathcal{F}_{e} + 2\mathcal{F}_{x}\), the number of
antineutrinos \(\bar{\mathcal{F}}_{e} + 2\mathcal{F}_{x}\) and the
momentum, with the four heavy-lepton species equal as in the transport
(\(\mathcal{F}_{x} = \bar{\mathcal{F}}_{x}\)). The amount
\(\delta_{\mathrm{FC}}\) brings the less abundant of \(\nu_e\) and
\(\bar\nu_e\) into equipartition with the heavy-lepton flavours
(complete equipartition of all six species is impossible when the
electron lepton number is not zero, because flavour conversion conserves
it); the second argument of the maximum keeps every occupation number at
most 1 (the Pauli principle) when heavy-lepton pairs are converted into
electron pairs, and every occupation number stays between 0 and 1. The
map is applied where the baryon density is below \(\rho_{\mathrm{FC}}\)
(their \(\rho_{\mathrm{c}}\), varied from \(10^9\) to \(10^{13}\)
g/cm\(^3\)) and from \(t_{\mathrm{FC}} = 0.02\) s after the bounce of
the core at \(x_{4,\mathrm{b}}\). In UM-4, FE-48 and FE-49 are those of
the Standard Model wherever \(v_{\mathrm{loc}} = v\), so the
prescription and its consequences carry over unchanged: the supernova
physics of UM-4 is that of the Standard Model with three Dirac
neutrinos, and the failed-explosion fractions of the paper apply
(DERIVED from FE-49). Status: STANDARD (the constraints), ASSUMED (this
explicit form), DERIVED (its validity in UM-4;
\texttt{um4\_neutrino.py}, N4).

\hypertarget{neutrinos-in-the-expanding-universe}{%
\subsection{Neutrinos in the expanding
universe}\label{neutrinos-in-the-expanding-universe}}

\[
\begin{aligned}
&\Big(\partial_{x_4} - a_4'\,\lvert\vec p\rvert\,\frac{\partial}{\partial\lvert\vec p\rvert}\Big)\varrho_{\vec p}
 = -\mathrm{i}\big[\mathsf{H}_{\vec p},\ \varrho_{\vec p}\big] + \mathcal{C}_{\vec p},\\
&N_{\mathrm{eff}} = \frac87\Big(\frac{11}{4}\Big)^{4/3}\frac{\rho_\nu}{\rho_\gamma},\qquad
 \mathsf{T}_\nu^{(0)} = \Big(\frac{4}{11}\Big)^{1/3}\mathsf{T}_\gamma^{(0)},\qquad
 n_\nu^{(0)} = \frac{3\zeta(3)}{2\pi^2}\big(\mathsf{T}_\nu^{(0)}\big)^3
\end{aligned} \tag{FE-53}
\] \textbf{Symbols in (FE-53).}

\begin{itemize}
\tightlist
\item
  \(\partial\): A partial derivative: the derivative with respect to one
  variable at fixed values of the others (\(\partial_\sigma\) with
  respect to \(\sigma\); \(\partial_M\) and \(\partial_\mu\) with
  respect to the coordinate that carries the same index, of the eight
  dimensions or of the fold). Status: STANDARD.
\item
  \(x_4\): The time; a prime on a function of \(x_4\) alone is
  \(\mathrm{d}/\mathrm{d}x_4\). Unit: length (= time, \(c = 1\)).
  Status: AUTHOR.
\item
  \(a_4\): The author's metric function of the time \(x_4\) that scales
  ordinary 3-space: on the fold \(\mathrm{e}^{a_4}\) is the scale factor
  of cosmology and \(a_4'\) the Hubble rate. Unit: dimensionless.
  Status: AUTHOR.
\item
  \({}'\): The derivative of a function of the time \(x_4\) alone with
  respect to \(x_4\) (\(a_4' = \mathrm{d}a_4/\mathrm{d}x_4\); two primes
  the second derivative); on a function of \(\sigma\) alone (a profile)
  the derivative with respect to \(\sigma\). Status: STANDARD.
\item
  \(\vec p\): The momentum of a neutrino (three components along
  ordinary space, in the rest frame of the medium);
  \(\lvert\vec p\rvert = E\) for a neutrino of negligible mass. Unit:
  GeV. Status: STANDARD.
\item
  \(\varrho_{\vec p},\ \bar\varrho_{\vec p}\): The \(3 \times 3\)
  flavour density matrices of the neutrinos and of the antineutrinos of
  momentum \(\vec p\) in the flavour basis (electron, muon, tau): the
  diagonal entries are the occupation numbers of the three flavours, the
  off-diagonal entries the coherence between the flavours;
  \(\bar\varrho\) is defined with the flavour indices in the reversed
  order, so that it obeys an equation of the same form. Unit:
  dimensionless. Status: STANDARD.
\item
  \(\mathrm{i}\): The imaginary unit, \(\mathrm{i}^2 = -1\). Unit:
  dimensionless. Status: STANDARD.
\item
  \(\mathsf{H}_{\vec p},\ \bar{\mathsf{H}}_{\vec p}\): The flavour
  Hamiltonians (\(3 \times 3\) Hermitian matrices) of the neutrinos and
  of the antineutrinos of momentum \(\vec p\). Unit: GeV. Status:
  STANDARD.
\item
  \(\mathcal{C}_{\vec p},\ \bar{\mathcal{C}}_{\vec p}\): The collision
  terms (\(3 \times 3\) matrices) of the neutrinos and the antineutrinos
  of momentum \(\vec p\): the rates of change of \(\varrho_{\vec p}\)
  and \(\bar\varrho_{\vec p}\) by emission, absorption, scattering and
  pair processes with the matter. Unit: GeV. Status: STANDARD.
\item
  \(N_{\mathrm{eff}}\): The effective number of neutrino species: the
  energy density of the relativistic particles other than photons in
  units of that of one neutrino species whose temperature is
  \((4/11)^{1/3}\) times that of the photons. Unit: dimensionless.
  Value: 3.044 (Standard Model); measured \(2.99 \pm 0.17\) (Planck
  2018). Status: STANDARD.
\item
  \(\rho_\nu\): The energy density of the neutrinos and antineutrinos of
  all flavours. Unit: GeV\(^4\). Status: STANDARD.
\item
  \(\rho_\gamma\): The energy density of the photons. Unit: GeV\(^4\).
  Status: STANDARD.
\item
  \(\mathsf{T}_\nu\): The temperature of the cosmic neutrino background;
  \(\mathsf{T}_\nu^{(0)}\) its value today. Unit: GeV (or K). Value:
  1.945 K today. Status: STANDARD.
\item
  \(\mathsf{T}_\gamma\): The temperature of the photons (of the CMB);
  \(\mathsf{T}_\gamma^{(0)}\) its value today. Unit: GeV (or K). Value:
  2.7255 K today. Status: STANDARD.
\item
  \(n_\nu\): The number density of the cosmic neutrino background per
  flavour (neutrinos plus antineutrinos, one helicity each);
  \(n_\nu^{(0)}\) its value today. Unit: GeV\(^3\). Value: 112
  cm\(^{-3}\) today. Status: STANDARD.
\item
  \(\zeta\): The Riemann zeta function; \(\zeta(3) = 1.2020569\). Unit:
  dimensionless. Status: STANDARD.
\item
  \(\pi\): The number pi, the ratio of the circumference of a circle to
  its diameter. Unit: dimensionless. Value: 3.141592654. Status:
  STANDARD.
\end{itemize}

The neutrinos in the homogeneous, isotropic, expanding fold (the
momentum redshifts with the expansion rate \(a_4'\)) and the cosmic
neutrino background (CNB); the superscript \((0)\) marks a value today.
Until the temperature falls to about 1 MeV the collisions keep the
neutrinos in equilibrium with the plasma; they then decouple, the
annihilation of electrons and positrons heats the photons, and the
effective number of neutrino species is \(N_{\mathrm{eff}} = 3.044\) in
the Standard Model (with the small heating of the neutrinos and the
flavour conversion during decoupling included). Today the CNB has the
temperature \(\mathsf{T}_\nu^{(0)} = 1.945\) K and \(n_\nu^{(0)} = 112\)
neutrinos and antineutrinos per cm\(^3\) and flavour (one helicity each:
the right-handed states of the Dirac neutrinos are not populated by the
Standard-Model collisions). The right-handed neutrinos add to
\(N_{\mathrm{eff}}\) only if something else brings them into equilibrium
(FE-62). Status: STANDARD (\texttt{um4\_neutrino.py}, N5).

\hypertarget{the-cosmological-perturbations}{%
\section{The cosmological
perturbations}\label{the-cosmological-perturbations}}

The CMB, the CNB and the BAO are computed on the fold, where
four-dimensional gravity is general relativity with \(G_{\mathrm{N}}\)
(FE-59: the fold obeys the Einstein equations of seven dimensions to the
first order in its curvature over \(H^2\), FE-71, and with the volume of
the extra times held fixed these are the equations of four-dimensional
general relativity), the matter is that of the Standard Model, the dark
matter is the particles \(\chi\) of dirac16complex (the long-lived
resonances of FE-40 and FE-86), and the dark energy is the vacuum energy
\(\rho_\Lambda\) of the fold (FE-59). The metric of the fold is the
background plus the first-order cosmological perturbation (of the order
\(\epsilon_{\mathrm{c}}\), the size of the perturbations); the equations
are those of Ma and Bertschinger (1995) in the conformal Newtonian
gauge, written out here.

\[
\begin{aligned}
&\hat h_{\mu\nu}\,\mathrm{d}x^\mu\mathrm{d}x^\nu = a^2\Big[-(1 + 2\psi_{\mathrm{B}})\,\mathrm{d}\tau_{\mathrm{c}}^2
 + (1 - 2\phi_{\mathrm{B}})\big(\mathrm{d}x_1^2 + \mathrm{d}x_2^2 + \mathrm{d}x_3^2\big)\Big],\\
&a = \mathrm{e}^{a_4},\qquad \mathrm{d}\tau_{\mathrm{c}} = \frac{\mathrm{d}x_4}{a},\qquad \mathcal{H} = \frac{\dot a}{a},\\
&k^2\phi_{\mathrm{B}} + 3\mathcal{H}\big(\dot\phi_{\mathrm{B}} + \mathcal{H}\psi_{\mathrm{B}}\big) = -4\pi G_{\mathrm{N}}a^2\,\Delta\rho,\qquad
 k^2\big(\dot\phi_{\mathrm{B}} + \mathcal{H}\psi_{\mathrm{B}}\big) = 4\pi G_{\mathrm{N}}a^2(\rho + \mathcal{P})\,\vartheta,\\
&\ddot\phi_{\mathrm{B}} + \mathcal{H}\big(\dot\psi_{\mathrm{B}} + 2\dot\phi_{\mathrm{B}}\big) + \Big(2\frac{\ddot a}{a} - \mathcal{H}^2\Big)\psi_{\mathrm{B}}
 + \frac{k^2}{3}\big(\phi_{\mathrm{B}} - \psi_{\mathrm{B}}\big) = 4\pi G_{\mathrm{N}}a^2\,\Delta\mathcal{P},\\
&k^2\big(\phi_{\mathrm{B}} - \psi_{\mathrm{B}}\big) = 12\pi G_{\mathrm{N}}a^2(\rho + \mathcal{P})\,\Sigma_{\mathrm{s}}
\end{aligned} \tag{FE-54}
\] \textbf{Symbols in (FE-54).}

\begin{itemize}
\tightlist
\item
  \(\hat h_{\mu\nu}\): The four-dimensional metric of the fold,
  \(\hat h_{\mu\nu} = \mathrm{diag}(\mathrm{e}^{2a_4}, \mathrm{e}^{2a_4}, \mathrm{e}^{2a_4}, -1)\)
  in the coordinates \(x_1, \dots, x_4\) (the induced metric at
  \(\sigma = 0\) without the extra times); \(\hat h^{\mu\nu}\) its
  inverse, \(\hat h\) its determinant,
  \(\sqrt{-\hat h} = \mathrm{e}^{3a_4}\); \(\hat h^{(1)}_{\mu\nu}\) its
  first-order cosmological perturbation. Unit: dimensionless. Status:
  DERIVED.
\item
  \(\mu\): A four-dimensional coordinate index running over the
  observable coordinates \(x_1, x_2, x_3, x_4\) of the fold (signature
  \(+,+,+,-\)); a repeated index is summed. Status: STANDARD.
\item
  \(\nu\): A four-dimensional coordinate index (as \(\mu\)). Status:
  STANDARD.
\item
  \(\mathrm{d}\): The differential: \(\mathrm{d}\) in front of a
  quantity is an infinitesimal change of it (as \(\mathrm{d}\sigma\), a
  small step of \(\sigma\)); the quotient of two such changes is a
  derivative; \(\mathrm{d}\) with an exponent, in front of several
  variables, is their volume element, the product of their
  differentials. Status: STANDARD.
\item
  \(x\): The coordinates of the fold: \(x_1, x_2, x_3\) of ordinary
  space, \(x_4\) the time, \(x_5, x_6, x_7\) the extra times; a field of
  the fold is a function of the four observable coordinates
  \(x = (x_1, x_2, x_3, x_4)\), and
  \(\mathrm{d}^4x = \mathrm{d}x_1\,\mathrm{d}x_2\,\mathrm{d}x_3\,\mathrm{d}x_4\).
  Unit: length (each coordinate). Status: AUTHOR.
\item
  \(a,\ \dot a,\ \ddot a\): The scale factor of the fold,
  \(a = \mathrm{e}^{a_4}\) (normalised to 1 today by the choice of the
  comoving coordinates \(x_1, x_2, x_3\)), and its first and second
  derivatives with respect to the conformal time \(\tau_{\mathrm{c}}\).
  Unit: dimensionless. Status: STANDARD.
\item
  \(\psi_{\mathrm{B}},\ \dot\psi_{\mathrm{B}}\): The Bardeen potential
  of the time-time part of the perturbed metric of the fold (the
  Newtonian potential) in the conformal Newtonian gauge, a function of
  \(\tau_{\mathrm{c}}\) and of the wave number \(k\), and its derivative
  with respect to \(\tau_{\mathrm{c}}\). Unit: dimensionless. Status:
  STANDARD.
\item
  \(\tau_{\mathrm{c}}\): The conformal time of the fold,
  \(\mathrm{d}\tau_{\mathrm{c}} = \mathrm{e}^{-a_4}\mathrm{d}x_4\) (the
  time divided by the scale factor \(\mathrm{e}^{a_4}\) of the fold); a
  dot on a quantity of the cosmological perturbations is
  \(\mathrm{d}/\mathrm{d}\tau_{\mathrm{c}}\). Unit: length (= time).
  Status: STANDARD.
\item
  \(\phi_{\mathrm{B}},\ \dot\phi_{\mathrm{B}},\ \ddot\phi_{\mathrm{B}}\):
  The Bardeen potential of the space-space part of the perturbed metric
  of the fold (the curvature potential) in the conformal Newtonian
  gauge, and its first and second derivatives with respect to
  \(\tau_{\mathrm{c}}\). Unit: dimensionless. Status: STANDARD.
\item
  \(x_1\): The first coordinate of ordinary 3-space. Unit: length.
  Status: AUTHOR.
\item
  \(x_2\): The second coordinate of ordinary 3-space. Unit: length.
  Status: AUTHOR.
\item
  \(x_3\): The third coordinate of ordinary 3-space. Unit: length.
  Status: AUTHOR.
\item
  \(\mathrm{e}\): Euler's number, the base of the natural logarithm;
  \(\mathrm{e}^{X}\) is the exponential of \(X\). Unit: dimensionless.
  Value: 2.718281828. Status: STANDARD.
\item
  \(a_4\): The author's metric function of the time \(x_4\) that scales
  ordinary 3-space: on the fold \(\mathrm{e}^{a_4}\) is the scale factor
  of cosmology and \(a_4'\) the Hubble rate. Unit: dimensionless.
  Status: AUTHOR.
\item
  \(x_4\): The time; a prime on a function of \(x_4\) alone is
  \(\mathrm{d}/\mathrm{d}x_4\). Unit: length (= time, \(c = 1\)).
  Status: AUTHOR.
\item
  \(\mathcal{H}\): The conformal expansion rate,
  \(\mathcal{H} = \dot a/a = a\,a_4'\). Unit: 1/length. Status:
  STANDARD.
\item
  \(k\): A wave number: the momentum of a plane wave along \(x_3\),
  \(\mathrm{e}^{\mathrm{i}(kx_3 - \omega x_4)}\); in the cosmological
  perturbations the comoving wave number of a Fourier mode,
  \(k = \lvert\vec k\rvert\). Unit: GeV (1/length). Status: STANDARD.
\item
  \(\pi\): The number pi, the ratio of the circumference of a circle to
  its diameter. Unit: dimensionless. Value: 3.141592654. Status:
  STANDARD.
\item
  \(G_{\mathrm{N}}\): Newton's constant of four-dimensional gravity (the
  gravity of our world, on the fold),
  \(G_{\mathrm{N}} = 1/M_{\mathrm{Pl}}^2\) with the Planck mass
  \(M_{\mathrm{Pl}} = 1.221\times10^{19}\) GeV. Unit: GeV\(^{-2}\).
  Value: \(6.708\times10^{-39}\) GeV\(^{-2}\). Status: STANDARD.
\item
  \(\Delta\): A change, difference or small variation of the quantity
  that follows (for example an interval of the time, a perturbation of a
  density, a difference of two squared masses, or an addition to an
  effective number of species). Status: STANDARD.
\item
  \(\rho\): The total energy density of the matter of the fold (all
  species); \(\Delta\rho\) its perturbation. Unit: GeV\(^4\). Status:
  STANDARD.
\item
  \(\mathcal{P}\): The total pressure of the matter of the fold;
  \(\Delta\mathcal{P}\) its perturbation. Unit: GeV\(^4\). Status:
  STANDARD.
\item
  \(\vartheta\): The divergence of the velocity of the total matter:
  \((\rho + \mathcal{P})\vartheta\) is the sum over the species of their
  energy density plus pressure times the divergence of their velocity.
  Unit: 1/length. Status: STANDARD.
\item
  \(\Sigma_{\mathrm{s}}\): The anisotropic stress (shear) of the total
  matter: \((\rho + \mathcal{P})\Sigma_{\mathrm{s}}\) is the sum over
  the species (photons and neutrinos) of their energy density plus
  pressure times their shear. Unit: dimensionless. Status: STANDARD.
\end{itemize}

The metric of the fold with the scalar perturbations in the conformal
Newtonian gauge (\(\psi_{\mathrm{B}}\) and \(\phi_{\mathrm{B}}\) the two
Bardeen potentials, functions of the conformal time
\(\tau_{\mathrm{c}}\) and of the comoving wave number \(k\) of a Fourier
mode; a dot is \(\mathrm{d}/\mathrm{d}\tau_{\mathrm{c}}\); the comoving
coordinates are chosen so that \(a = 1\) today), and the Einstein
equations of the fold (FE-59) to first order: the energy constraint, the
momentum constraint, the trace of the space-space equations and their
traceless part. \(\Delta\rho\) and \(\Delta\mathcal{P}\) are the
perturbations of the total energy density and pressure;
\((\rho + \mathcal{P})\vartheta\) and
\((\rho + \mathcal{P})\Sigma_{\mathrm{s}}\) are the sums over the
species of (energy density plus pressure) times the divergence of their
velocity and times their anisotropic stress. By FE-34 and FE-71 the bulk
metric of any metric of the fold, perturbed or not, is fixed to the
first order in the curvature of the fold over \(H^2\), and the fold then
obeys the Einstein equations of seven dimensions: with the volume of the
extra times held fixed by the fold fluid (for the perturbations as well,
ASSUMED) the first-order perturbations of the fold are those of general
relativity, with corrections of the relative orders \((H_4/H)^2\) and
\((k/(aH))^2\) (\(k/a\) the physical wave number) from the massive modes
of the bulk, below \(10^{-78}\) at the scales of the CMB for
\(H > 2534\) GeV (ESTIMATE), and its gravitational waves are those of
general relativity (FE-63). The Weyl term \(\mathfrak{c}_{ab}\) of the
bulk in FE-34 (dark radiation) is taken to be zero (ASSUMED). Status:
STANDARD (DERIVED: the equations of the fold are those of general
relativity to the first order in its curvature over \(H^2\), FE-71;
ASSUMED: the volume held fixed for the perturbations, no dark
radiation).

\[
\begin{aligned}
&\dot\delta_\chi = -\vartheta_\chi + 3\dot\phi_{\mathrm{B}},\qquad
 \dot\vartheta_\chi = -\mathcal{H}\vartheta_\chi + k^2\psi_{\mathrm{B}} + \Gamma_{\chi\mathrm{b}}\big(\vartheta_{\mathrm{b}} - \vartheta_\chi\big),\\
&\dot\delta_{\mathrm{b}} = -\vartheta_{\mathrm{b}} + 3\dot\phi_{\mathrm{B}},\qquad
 \dot\vartheta_{\mathrm{b}} = -\mathcal{H}\vartheta_{\mathrm{b}} + c_{\mathrm{s}}^2k^2\delta_{\mathrm{b}} + k^2\psi_{\mathrm{B}}
 + \frac{a\,x_{\mathrm{e}}n_e\sigma_{\mathrm{T}}}{\mathsf{R}}\big(\vartheta_\gamma - \vartheta_{\mathrm{b}}\big)
 + \frac{\rho_\chi}{\rho_{\mathrm{b}}}\,\Gamma_{\chi\mathrm{b}}\big(\vartheta_\chi - \vartheta_{\mathrm{b}}\big),\qquad
 \mathsf{R} = \frac{3\rho_{\mathrm{b}}}{4\rho_\gamma},\\
&\dot\delta_\gamma = -\frac43\vartheta_\gamma + 4\dot\phi_{\mathrm{B}},\qquad
 \dot\vartheta_\gamma = k^2\Big(\frac{\delta_\gamma}{4} - \Sigma_\gamma\Big) + k^2\psi_{\mathrm{B}} + a\,x_{\mathrm{e}}n_e\sigma_{\mathrm{T}}\big(\vartheta_{\mathrm{b}} - \vartheta_\gamma\big),\\
&\dot{\mathsf{F}}_l = \frac{k}{2l + 1}\big(l\,\mathsf{F}_{l-1} - (l + 1)\,\mathsf{F}_{l+1}\big)
 - a\,x_{\mathrm{e}}n_e\sigma_{\mathrm{T}}\Big(\mathsf{F}_l - \frac{\mathsf{G}_0 + \mathsf{F}_2 + \mathsf{G}_2}{10}\,\delta_{l2}\Big)\quad(l \ge 2),\\
&\dot{\mathsf{G}}_l = \frac{k}{2l + 1}\big(l\,\mathsf{G}_{l-1} - (l + 1)\,\mathsf{G}_{l+1}\big)
 - a\,x_{\mathrm{e}}n_e\sigma_{\mathrm{T}}\Big(\mathsf{G}_l - \frac{\mathsf{G}_0 + \mathsf{F}_2 + \mathsf{G}_2}{2}\Big(\delta_{l0} + \frac{\delta_{l2}}{5}\Big)\Big)\quad(l \ge 0),\\
&\mathsf{F}_0 = \delta_\gamma,\qquad \mathsf{F}_1 = \frac{4\vartheta_\gamma}{3k},\qquad \mathsf{F}_2 = 2\Sigma_\gamma,\\
&\dot\delta_\nu = -\frac43\vartheta_\nu + 4\dot\phi_{\mathrm{B}},\qquad \dot\vartheta_\nu = k^2\Big(\frac{\delta_\nu}{4} - \Sigma_\nu\Big) + k^2\psi_{\mathrm{B}},\qquad
 \dot{\mathsf{N}}_l = \frac{k}{2l + 1}\big(l\,\mathsf{N}_{l-1} - (l + 1)\,\mathsf{N}_{l+1}\big)\quad(l \ge 2),\\
&\mathsf{N}_1 = \frac{4\vartheta_\nu}{3k},\qquad \mathsf{N}_2 = 2\Sigma_\nu,\\
&\dot\Upsilon_0 = -\frac{\mathsf{q}k}{\mathsf{E}}\Upsilon_1 - \dot\phi_{\mathrm{B}}\frac{\mathrm{d}\ln\mathcal{F}_0}{\mathrm{d}\ln\mathsf{q}},\qquad
 \dot\Upsilon_1 = \frac{\mathsf{q}k}{3\mathsf{E}}\big(\Upsilon_0 - 2\Upsilon_2\big) - \frac{\mathsf{E}k}{3\mathsf{q}}\,\psi_{\mathrm{B}}\frac{\mathrm{d}\ln\mathcal{F}_0}{\mathrm{d}\ln\mathsf{q}},\\
&\dot\Upsilon_l = \frac{\mathsf{q}k}{(2l + 1)\mathsf{E}}\big(l\,\Upsilon_{l-1} - (l + 1)\,\Upsilon_{l+1}\big)\quad(l \ge 2),\qquad
 \mathsf{E} = \sqrt{\mathsf{q}^2 + a^2m_i^2},\qquad \mathcal{F}_0 = \frac{1}{\mathrm{e}^{\mathsf{q}/\mathsf{T}_\nu^{(0)}} + 1}
\end{aligned} \tag{FE-55}
\] \textbf{Symbols in (FE-55).}

\begin{itemize}
\tightlist
\item
  \(\delta_\chi,\ \dot\delta_\chi\): The density contrast of the dark
  matter (the particles \(\chi\) of dirac16complex), the perturbation of
  its energy density divided by its mean, and its derivative with
  respect to \(\tau_{\mathrm{c}}\). Unit: dimensionless. Status:
  STANDARD.
\item
  \(\vartheta_\chi,\ \dot\vartheta_\chi\): The divergence of the
  velocity of the dark matter (for a Fourier mode, \(\mathrm{i}k\) times
  its velocity along the wave vector), and its derivative. Unit:
  1/length. Status: STANDARD.
\item
  \(\phi_{\mathrm{B}},\ \dot\phi_{\mathrm{B}},\ \ddot\phi_{\mathrm{B}}\):
  The Bardeen potential of the space-space part of the perturbed metric
  of the fold (the curvature potential) in the conformal Newtonian
  gauge, and its first and second derivatives with respect to
  \(\tau_{\mathrm{c}}\). Unit: dimensionless. Status: STANDARD.
\item
  \(\mathcal{H}\): The conformal expansion rate,
  \(\mathcal{H} = \dot a/a = a\,a_4'\). Unit: 1/length. Status:
  STANDARD.
\item
  \(k\): A wave number: the momentum of a plane wave along \(x_3\),
  \(\mathrm{e}^{\mathrm{i}(kx_3 - \omega x_4)}\); in the cosmological
  perturbations the comoving wave number of a Fourier mode,
  \(k = \lvert\vec k\rvert\). Unit: GeV (1/length). Status: STANDARD.
\item
  \(\psi_{\mathrm{B}},\ \dot\psi_{\mathrm{B}}\): The Bardeen potential
  of the time-time part of the perturbed metric of the fold (the
  Newtonian potential) in the conformal Newtonian gauge, a function of
  \(\tau_{\mathrm{c}}\) and of the wave number \(k\), and its derivative
  with respect to \(\tau_{\mathrm{c}}\). Unit: dimensionless. Status:
  STANDARD.
\item
  \(\Gamma_{\chi\mathrm{b}}\): The rate (per conformal time) at which
  the dark matter exchanges momentum with the baryons through the portal
  couplings (elastic scattering of \(\chi\) on nuclei and electrons).
  Unit: 1/length. Value: negligible for the couplings that direct
  detection allows. Status: DERIVED.
\item
  \(\vartheta_{\mathrm{b}},\ \dot\vartheta_{\mathrm{b}}\): The
  divergence of the velocity of the baryons and its derivative. Unit:
  1/length. Status: STANDARD.
\item
  \(\delta_{\mathrm{b}},\ \dot\delta_{\mathrm{b}}\): The density
  contrast of the baryons and its derivative. Unit: dimensionless.
  Status: STANDARD.
\item
  \(c_{\mathrm{s}}\): The sound speed of the baryon gas: its square is
  the temperature of the gas divided by the mean mass per particle,
  times one minus a third of the logarithmic derivative of that
  temperature with respect to \(a\). Unit: dimensionless. Status:
  STANDARD.
\item
  \(a,\ \dot a,\ \ddot a\): The scale factor of the fold,
  \(a = \mathrm{e}^{a_4}\) (normalised to 1 today by the choice of the
  comoving coordinates \(x_1, x_2, x_3\)), and its first and second
  derivatives with respect to the conformal time \(\tau_{\mathrm{c}}\).
  Unit: dimensionless. Status: STANDARD.
\item
  \(x_{\mathrm{e}}\): The ionisation fraction: the fraction of the
  electrons that are free. Unit: dimensionless. Status: STANDARD.
\item
  \(n_e\): The number density of electrons (bound and free; in a
  supernova all are free, in the recombining universe
  \(x_{\mathrm{e}}n_e\) is the density of the free electrons). Unit:
  GeV\(^3\) (cm\(^{-3}\)). Status: STANDARD.
\item
  \(\sigma_{\mathrm{T}}\): The Thomson cross-section of the scattering
  of photons on free electrons. Unit: GeV\(^{-2}\). Value:
  \(6.6524587\times10^{-25}\) cm\(^2\). Status: STANDARD.
\item
  \(\mathsf{R}\): The ratio of the inertia of the baryons to that of the
  photons, \(\mathsf{R} = 3\rho_{\mathrm{b}}/(4\rho_\gamma)\). Unit:
  dimensionless. Status: STANDARD.
\item
  \(\vartheta_\gamma,\ \dot\vartheta_\gamma\): The divergence of the
  velocity of the photons and its derivative. Unit: 1/length. Status:
  STANDARD.
\item
  \(\rho_\chi\): The energy density of the dark matter, the particles
  \(\chi\) of dirac16complex; \(\rho_\chi^{(0)}\) its value today. Unit:
  GeV\(^4\). Status: DERIVED.
\item
  \(\rho_{\mathrm{b}}\): The (rest-mass) energy density of the baryons:
  in a supernova the local density of the matter (quoted in g/cm\(^3\)),
  in cosmology the mean density of the baryons
  (\(\rho_{\mathrm{b}}^{(0)}\) today). Unit: GeV\(^4\). Status:
  STANDARD.
\item
  \(\rho_\gamma\): The energy density of the photons. Unit: GeV\(^4\).
  Status: STANDARD.
\item
  \(\delta_\gamma,\ \dot\delta_\gamma\): The density contrast of the
  photons and its derivative. Unit: dimensionless. Status: STANDARD.
\item
  \(\Sigma_\gamma\): The anisotropic stress (shear) of the photons.
  Unit: dimensionless. Status: STANDARD.
\item
  \(\mathsf{F}_l,\ \dot{\mathsf{F}}_l\): The multipoles of the
  temperature perturbation of the photons
  (\(\mathsf{F}_0 = \delta_\gamma\),
  \(\mathsf{F}_1 = 4\vartheta_\gamma/(3k)\),
  \(\mathsf{F}_2 = 2\Sigma_\gamma\); the photon multipoles of Ma and
  Bertschinger) and their derivatives with respect to
  \(\tau_{\mathrm{c}}\). Unit: dimensionless. Status: STANDARD.
\item
  \(l\): The multipole index and number (\(l = 0, 1, 2, \dots\)) of an
  expansion in Legendre polynomials of the direction of a photon or a
  neutrino relative to the wave vector, or of the CMB sky in spherical
  harmonics. Status: STANDARD.
\item
  \(\mathsf{G}_l,\ \dot{\mathsf{G}}_l\): The multipoles of the
  polarisation perturbation of the photons (the \(G_{\gamma l}\) of Ma
  and Bertschinger) and their derivatives. Unit: dimensionless. Status:
  STANDARD.
\item
  \(\delta\): The Kronecker delta: with two indices \(\delta^M_N = 1\)
  if \(M = N\) and 0 otherwise (likewise \(\delta_{X\mathrm{L}}\) is 1
  if the label \(X\) of a half of dirac16complex is L and 0 otherwise);
  with two index strings of equal length the generalised Kronecker delta
  \(\delta^{M_1M_2\cdots}_{N_1N_2\cdots}\), the determinant of the
  square matrix of the Kronecker deltas \(\delta^{M_a}_{N_b}\) (\(+1\)
  or \(-1\) if the upper indices are an even or odd permutation of the
  lower ones, 0 otherwise). Status: STANDARD.
\item
  \(\delta_\nu,\ \dot\delta_\nu\): The density contrast of the
  relativistic neutrinos and its derivative. Unit: dimensionless.
  Status: STANDARD.
\item
  \(\vartheta_\nu,\ \dot\vartheta_\nu\): The divergence of the velocity
  of the relativistic neutrinos and its derivative. Unit: 1/length.
  Status: STANDARD.
\item
  \(\Sigma_\nu\): The anisotropic stress (shear) of the relativistic
  neutrinos. Unit: dimensionless. Status: STANDARD.
\item
  \(\mathsf{N}_l,\ \dot{\mathsf{N}}_l\): The multipoles of the
  perturbation of the relativistic neutrinos
  (\(\mathsf{N}_0 = \delta_\nu\),
  \(\mathsf{N}_1 = 4\vartheta_\nu/(3k)\),
  \(\mathsf{N}_2 = 2\Sigma_\nu\)) and their derivatives. Unit:
  dimensionless. Status: STANDARD.
\item
  \(\Upsilon_l(\mathsf{q}),\ \dot\Upsilon_l\): The multipoles of the
  relative perturbation of the distribution of a massive neutrino state
  at the comoving momentum \(\mathsf{q}\) (the \(\Psi_l\) of Ma and
  Bertschinger) and their derivatives. Unit: dimensionless. Status:
  STANDARD.
\item
  \(\mathsf{q}\): The comoving momentum of a neutrino,
  \(a\lvert\vec p\rvert\). Unit: GeV. Status: STANDARD.
\item
  \(\mathsf{E}\): The comoving energy of a massive neutrino (\(a\) times
  its energy). Unit: GeV. Status: STANDARD.
\item
  \(\mathrm{d}\): The differential: \(\mathrm{d}\) in front of a
  quantity is an infinitesimal change of it (as \(\mathrm{d}\sigma\), a
  small step of \(\sigma\)); the quotient of two such changes is a
  derivative; \(\mathrm{d}\) with an exponent, in front of several
  variables, is their volume element, the product of their
  differentials. Status: STANDARD.
\item
  \(\mathcal{F}_0(\mathsf{q})\): The unperturbed (Fermi-Dirac)
  occupation number of the cosmic neutrino background at the comoving
  momentum \(\mathsf{q}\). Unit: dimensionless. Status: STANDARD.
\item
  \(m\): A mass: \(m_i\) the mass of the neutrino mass state \(i\);
  \(\Delta m_{ij}^2 = m_i^2 - m_j^2\) the squared-mass difference of the
  states \(i\) and \(j\) (\(\Delta m_{21}^2\), \(\Delta m_{31}^2\),
  \(\Delta m_{32}^2\)); \(\Delta m^2\) without labels the difference of
  the two states of a resonance, \(m_2^2 - m_1^2\), multiplied by
  \(v_{\mathrm{loc}}^2/v^2\) where this differs from 1; \(m\) without a
  label the mass of a mode of a field of the bulk as the fold sees it,
  \(m^2 = \omega^2 - k^2\); for the W and Z bosons their complex pole
  mass, \(m^2 = m_0^2 - \mathrm{i}m_0\Gamma_{\mathrm{bulk}}\) to the
  leading order; in the relic abundance of the particles \(\chi\) their
  mass. Unit: GeV. Status: STANDARD.
\item
  \(i\): A generation index (\(i = 1, 2, 3\)), which also labels the
  three neutrino mass states; in a frame component such as
  \(\mathcal{S}^{(i)}{}_{(i)}\) one of the three directions of ordinary
  space (the three components are equal), and in a tensor of ordinary
  space such as \(\mathsf{h}_{ij}\) one of the directions
  \(x_1, x_2, x_3\); a repeated index is not summed unless a sum sign is
  written. Status: STANDARD.
\item
  \(\mathrm{e}\): Euler's number, the base of the natural logarithm;
  \(\mathrm{e}^{X}\) is the exponential of \(X\). Unit: dimensionless.
  Value: 2.718281828. Status: STANDARD.
\item
  \(\mathsf{T}_\nu\): The temperature of the cosmic neutrino background;
  \(\mathsf{T}_\nu^{(0)}\) its value today. Unit: GeV (or K). Value:
  1.945 K today. Status: STANDARD.
\end{itemize}

The Boltzmann equations of the species in the same gauge (Ma and
Bertschinger 1995, their equations 42, 66, 63, 49 and 57): the dark
matter \(\chi\) (the long-lived resonances of dirac16complex at the
fold, FE-86: pressureless, FE-42, with the drag
\(\Gamma_{\chi\mathrm{b}}\) on the baryons from its portal couplings,
FE-47, which is negligible for the couplings that direct detection
allows); the baryons (with the sound speed \(c_{\mathrm{s}}\) of the gas
and the Thomson coupling to the photons through the free electrons,
\(x_{\mathrm{e}}\) the ionisation fraction); the photons (the multipoles
\(\mathsf{F}_l\) of the temperature and \(\mathsf{G}_l\) of the
polarisation, \(l\) the multipole index; \(\delta_{l0}\),
\(\delta_{l2}\) are Kronecker deltas, and the term
\(l\,\mathsf{G}_{l-1}\) vanishes for \(l = 0\)); the neutrinos while
they are relativistic (the multipoles \(\mathsf{N}_l\)); and each
massive neutrino state \(i\) of mass \(m_i\) (the multipoles
\(\Upsilon_l\) of the relative perturbation of its distribution at the
comoving momentum \(\mathsf{q}\) and energy \(\mathsf{E}\), with the
Fermi-Dirac occupation \(\mathcal{F}_0\) of the CNB). The hierarchies
are truncated at a large \(l\) in a numerical solution. Status: STANDARD
(UM-4: the dark matter is \(\chi\), with the drag term).

\[
\begin{aligned}
&r_{\mathrm{s}} = \int_0^{\tau_{\mathrm{d}}}c_{\gamma\mathrm{b}}\,\mathrm{d}\tau_{\mathrm{c}},\qquad c_{\gamma\mathrm{b}} = \frac{1}{\sqrt{3(1 + \mathsf{R})}},\\
&D_{\mathrm{M}}(\mathsf{z}) = \int_0^{\mathsf{z}}\frac{\mathrm{d}\mathsf{z}'}{H_4(\mathsf{z}')},\qquad
 D_{\mathrm{H}}(\mathsf{z}) = \frac{1}{H_4(\mathsf{z})},\qquad
 D_{\mathrm{V}}(\mathsf{z}) = \big(\mathsf{z}\,D_{\mathrm{M}}^2(\mathsf{z})\,D_{\mathrm{H}}(\mathsf{z})\big)^{1/3}
\end{aligned} \tag{FE-56}
\] \textbf{Symbols in (FE-56).}

\begin{itemize}
\tightlist
\item
  \(r_{\mathrm{s}}\): The comoving sound horizon at the drag epoch: the
  comoving distance that a sound wave of the photon-baryon fluid travels
  until the baryons are released from the photons. Unit: length. Value:
  about 147 Mpc (Planck 2018). Status: STANDARD.
\item
  \(\tau_{\mathrm{d}}\): The conformal time of the drag epoch, when the
  baryons are released from the photons (redshift about 1060). Unit:
  length (= time). Status: STANDARD.
\item
  \(c_{\gamma\mathrm{b}}\): The sound speed of the tightly coupled
  photon-baryon fluid. Unit: dimensionless. Status: STANDARD.
\item
  \(\mathrm{d}\): The differential: \(\mathrm{d}\) in front of a
  quantity is an infinitesimal change of it (as \(\mathrm{d}\sigma\), a
  small step of \(\sigma\)); the quotient of two such changes is a
  derivative; \(\mathrm{d}\) with an exponent, in front of several
  variables, is their volume element, the product of their
  differentials. Status: STANDARD.
\item
  \(\tau_{\mathrm{c}}\): The conformal time of the fold,
  \(\mathrm{d}\tau_{\mathrm{c}} = \mathrm{e}^{-a_4}\mathrm{d}x_4\) (the
  time divided by the scale factor \(\mathrm{e}^{a_4}\) of the fold); a
  dot on a quantity of the cosmological perturbations is
  \(\mathrm{d}/\mathrm{d}\tau_{\mathrm{c}}\). Unit: length (= time).
  Status: STANDARD.
\item
  \(\mathsf{R}\): The ratio of the inertia of the baryons to that of the
  photons, \(\mathsf{R} = 3\rho_{\mathrm{b}}/(4\rho_\gamma)\). Unit:
  dimensionless. Status: STANDARD.
\item
  \(D_{\mathrm{M}}(\mathsf{z})\): The transverse comoving distance
  (comoving angular-diameter distance) to the redshift \(\mathsf{z}\),
  for a spatially flat fold. Unit: length. Status: STANDARD.
\item
  \(\mathsf{z}\): The redshift, \(1 + \mathsf{z} = 1/a\);
  \(\mathsf{z}'\) an integration variable. Unit: dimensionless. Status:
  STANDARD.
\item
  \({}'\): The derivative of a function of the time \(x_4\) alone with
  respect to \(x_4\) (\(a_4' = \mathrm{d}a_4/\mathrm{d}x_4\); two primes
  the second derivative); on a function of \(\sigma\) alone (a profile)
  the derivative with respect to \(\sigma\). Status: STANDARD.
\item
  \(H_4\): The expansion rate (Hubble rate) of the fold, \(H_4 = a_4'\),
  as a function of the time or of the redshift; \(H_4^{(0)}\) its value
  today, the Hubble constant. Unit: GeV (1/time). Value:
  \(H_4^{(0)} = 100\,h_{100}\) km/s/Mpc with \(h_{100} = 0.674\) (Planck
  2018): 67.4 km/s/Mpc \(= 1.44\times10^{-42}\) GeV. Status: STANDARD.
\item
  \(D_{\mathrm{H}}(\mathsf{z})\): The Hubble distance at the redshift
  \(\mathsf{z}\). Unit: length. Status: STANDARD.
\item
  \(D_{\mathrm{V}}(\mathsf{z})\): The volume-averaged distance to the
  redshift \(\mathsf{z}\). Unit: length. Status: STANDARD.
\end{itemize}

The baryon acoustic oscillations: the sound horizon \(r_{\mathrm{s}}\)
at the drag epoch \(\tau_{\mathrm{d}}\) (the comoving distance that a
sound wave of the photon-baryon fluid, of speed
\(c_{\gamma\mathrm{b}}\), travels before the baryons are released from
the photons), and the distances to the redshift \(\mathsf{z}\) at which
galaxy surveys measure it: the transverse comoving distance
\(D_{\mathrm{M}}\) (for a spatially flat fold), the Hubble distance
\(D_{\mathrm{H}}\) and their volume average \(D_{\mathrm{V}}\)
(\(H_4 = a_4'\) the expansion rate of the fold as a function of
redshift). The surveys measure \(D_{\mathrm{M}}/r_{\mathrm{s}}\),
\(D_{\mathrm{H}}/r_{\mathrm{s}}\) and \(D_{\mathrm{V}}/r_{\mathrm{s}}\).
Status: STANDARD.

\[
\begin{aligned}
&C_l^{\mathrm{TT}} = 4\pi\int\frac{\mathrm{d}k}{k}\,\mathcal{P}_{\mathcal{R}}(k)\,\big(\mathcal{I}_l^{\mathrm{T}}(k)\big)^2,\qquad
 C_l^{\mathrm{TE}} = 4\pi\int\frac{\mathrm{d}k}{k}\,\mathcal{P}_{\mathcal{R}}(k)\,\mathcal{I}_l^{\mathrm{T}}(k)\,\mathcal{I}_l^{\mathrm{E}}(k),\\
&C_l^{\mathrm{EE}} = 4\pi\int\frac{\mathrm{d}k}{k}\,\mathcal{P}_{\mathcal{R}}(k)\,\big(\mathcal{I}_l^{\mathrm{E}}(k)\big)^2,\qquad
 \mathcal{P}_{\mathcal{R}}(k) = A_{\mathrm{s}}\Big(\frac{k}{k_*}\Big)^{n_{\mathrm{S}} - 1}
\end{aligned} \tag{FE-57}
\] \textbf{Symbols in (FE-57).}

\begin{itemize}
\tightlist
\item
  \(C_l^{\mathrm{TT}},\ C_l^{\mathrm{TE}},\ C_l^{\mathrm{EE}}\): The
  angular power spectra of the CMB at the multipole \(l\): temperature
  with temperature, temperature with E-mode polarisation, E-mode
  polarisation with itself. Unit: dimensionless. Status: STANDARD.
\item
  \(l\): The multipole index and number (\(l = 0, 1, 2, \dots\)) of an
  expansion in Legendre polynomials of the direction of a photon or a
  neutrino relative to the wave vector, or of the CMB sky in spherical
  harmonics. Status: STANDARD.
\item
  \(\mathrm{T},\ \mathrm{E}\): The labels of the CMB observables:
  \(\mathrm{T}\) the temperature anisotropy, \(\mathrm{E}\) the E-mode
  (gradient) polarisation; \(\mathrm{TT}\), \(\mathrm{TE}\),
  \(\mathrm{EE}\) label the spectra of their products. Status: STANDARD.
\item
  \(\pi\): The number pi, the ratio of the circumference of a circle to
  its diameter. Unit: dimensionless. Value: 3.141592654. Status:
  STANDARD.
\item
  \(\mathrm{d}\): The differential: \(\mathrm{d}\) in front of a
  quantity is an infinitesimal change of it (as \(\mathrm{d}\sigma\), a
  small step of \(\sigma\)); the quotient of two such changes is a
  derivative; \(\mathrm{d}\) with an exponent, in front of several
  variables, is their volume element, the product of their
  differentials. Status: STANDARD.
\item
  \(k\): A wave number: the momentum of a plane wave along \(x_3\),
  \(\mathrm{e}^{\mathrm{i}(kx_3 - \omega x_4)}\); in the cosmological
  perturbations the comoving wave number of a Fourier mode,
  \(k = \lvert\vec k\rvert\). Unit: GeV (1/length). Status: STANDARD.
\item
  \(\mathcal{P}_{\mathcal{R}}(k)\): The primordial power spectrum of the
  comoving curvature perturbation, per logarithmic interval of \(k\).
  Unit: dimensionless. Status: ASSUMED.
\item
  \(\mathcal{I}_l^{\mathrm{T}}(k),\ \mathcal{I}_l^{\mathrm{E}}(k)\): The
  transfer functions of the CMB: the multipole \(l\) of the temperature
  or E-mode anisotropy today produced by a Fourier mode \(k\) of unit
  primordial curvature perturbation, computed by the line-of-sight
  integral. Unit: dimensionless. Status: STANDARD.
\item
  \(A_{\mathrm{s}}\): The amplitude of the primordial curvature spectrum
  at the pivot \(k_*\). Unit: dimensionless. Value:
  \(2.10\times10^{-9}\) (Planck 2018). Status: ASSUMED.
\item
  \(k_*\): The pivot wave number of the primordial spectrum. Unit: GeV
  (1/length). Value: 0.05 Mpc\(^{-1}\). Status: STANDARD.
\item
  \(n_{\mathrm{S}}\): The spectral index of the primordial curvature
  spectrum. Unit: dimensionless. Value: 0.965 (Planck 2018). Status:
  ASSUMED.
\end{itemize}

The angular power spectra of the CMB temperature (T) and E-mode
polarisation (E) at the multipole \(l\): the transfer functions
\(\mathcal{I}_l^{\mathrm{T}}\), \(\mathcal{I}_l^{\mathrm{E}}\) follow
from FE-54 and FE-55, integrated from the early universe to today (the
line-of-sight integral), and are weighted with the primordial spectrum
\(\mathcal{P}_{\mathcal{R}}\) of the comoving curvature perturbation
(amplitude \(A_{\mathrm{s}}\), spectral index \(n_{\mathrm{S}}\), pivot
\(k_*\)). The primordial spectrum is an input of UM-4 (ASSUMED): the
primordial era can be an inflation of the fold driven by a vacuum energy
of the fold with the volume of the extra times held fixed (FE-88: with
the holding pressure of the fold fluid, whose origin is OPEN, and a Weyl
tensor of the bulk that the exact equations of the fold require, OPEN),
in the regime of general relativity only if its rate is below \(H/10\),
and UM-4 does not fix its spectrum (the Higgs field could drive it only
with a large coupling to the curvature of the fold, which the action of
UM-4 does not contain, IN-11, and then needs \(H > 1.6\times10^{14}\)
GeV, FE-88). Status: STANDARD (ASSUMED: the primordial spectrum).

\[
\begin{aligned}
&\Omega_\chi = \frac{\rho_\chi^{(0)}}{\rho_{\mathrm{c}}^{(0)}},\qquad \Omega_{\mathrm{b}} = \frac{\rho_{\mathrm{b}}^{(0)}}{\rho_{\mathrm{c}}^{(0)}},\qquad
 \Omega_\Lambda = \frac{\rho_\Lambda}{\rho_{\mathrm{c}}^{(0)}},\qquad \Omega_\nu h_{100}^2 = \frac{\sum_im_i}{93.14\ \text{eV}},\\
&\rho_{\mathrm{c}}^{(0)} = \frac{3\big(H_4^{(0)}\big)^2}{8\pi G_{\mathrm{N}}},\qquad H_4^{(0)} = 100\,h_{100}\ \text{km/s/Mpc}
\end{aligned} \tag{FE-58}
\] \textbf{Symbols in (FE-58).}

\begin{itemize}
\tightlist
\item
  \(\Omega_\chi\): The density parameter of the dark matter today. Unit:
  dimensionless. Value: \(\Omega_\chi h_{100}^2 = 0.120\) (Planck 2018).
  Status: DERIVED.
\item
  \(\rho_\chi\): The energy density of the dark matter, the particles
  \(\chi\) of dirac16complex; \(\rho_\chi^{(0)}\) its value today. Unit:
  GeV\(^4\). Status: DERIVED.
\item
  \(\rho_{\mathrm{c}}\): The critical density, the energy density of a
  spatially flat fold at the expansion rate \(H_4\);
  \(\rho_{\mathrm{c}}^{(0)}\) its value today. Unit: GeV\(^4\). Value:
  \(\rho_{\mathrm{c}}^{(0)} = 1.0537\times10^{-5}\,h_{100}^2\)
  GeV/cm\(^3\). Status: STANDARD.
\item
  \(\Omega_{\mathrm{b}}\): The density parameter of the baryons today.
  Unit: dimensionless. Value: \(\Omega_{\mathrm{b}}h_{100}^2 = 0.0224\)
  (Planck 2018). Status: STANDARD.
\item
  \(\rho_{\mathrm{b}}\): The (rest-mass) energy density of the baryons:
  in a supernova the local density of the matter (quoted in g/cm\(^3\)),
  in cosmology the mean density of the baryons
  (\(\rho_{\mathrm{b}}^{(0)}\) today). Unit: GeV\(^4\). Status:
  STANDARD.
\item
  \(\Omega_\Lambda\): The density parameter of the dark energy (the
  vacuum energy of the fold) today. Unit: dimensionless. Value: 0.685
  (Planck 2018). Status: DERIVED.
\item
  \(\rho_\Lambda\): The vacuum energy density of the fold above its
  static tension, the dark energy of UM-4:
  \(\rho_\Lambda = \rho_{\mathrm{ff}} + \rho_{\mathrm{vac}} - \lambda_{\mathrm{F}}V_*\)
  (four-dimensional normalisation). Unit: GeV\(^4\). Value:
  \(0.685\,\rho_{\mathrm{c}}^{(0)} = 2.5\times10^{-47}\) GeV\(^4\)
  (Planck 2018). Status: DERIVED.
\item
  \(\Omega_\nu,\ \Omega_\nu^{(i)}\): The density parameter of the
  neutrinos of the CNB: \(\Omega_\nu\) that of the massive states today;
  \(\Omega_\nu^{(i)}(\mathsf{z})\) the energy density of the mass state
  \(i\) at the redshift \(\mathsf{z}\) in units of the critical density
  today. Unit: dimensionless. Status: STANDARD.
\item
  \(h_{100}\): The present expansion rate of the fold in units of 100
  km/s/Mpc. Unit: dimensionless. Value: 0.674 (Planck 2018). Status:
  STANDARD.
\item
  \(i\): A generation index (\(i = 1, 2, 3\)), which also labels the
  three neutrino mass states; in a frame component such as
  \(\mathcal{S}^{(i)}{}_{(i)}\) one of the three directions of ordinary
  space (the three components are equal), and in a tensor of ordinary
  space such as \(\mathsf{h}_{ij}\) one of the directions
  \(x_1, x_2, x_3\); a repeated index is not summed unless a sum sign is
  written. Status: STANDARD.
\item
  \(m\): A mass: \(m_i\) the mass of the neutrino mass state \(i\);
  \(\Delta m_{ij}^2 = m_i^2 - m_j^2\) the squared-mass difference of the
  states \(i\) and \(j\) (\(\Delta m_{21}^2\), \(\Delta m_{31}^2\),
  \(\Delta m_{32}^2\)); \(\Delta m^2\) without labels the difference of
  the two states of a resonance, \(m_2^2 - m_1^2\), multiplied by
  \(v_{\mathrm{loc}}^2/v^2\) where this differs from 1; \(m\) without a
  label the mass of a mode of a field of the bulk as the fold sees it,
  \(m^2 = \omega^2 - k^2\); for the W and Z bosons their complex pole
  mass, \(m^2 = m_0^2 - \mathrm{i}m_0\Gamma_{\mathrm{bulk}}\) to the
  leading order; in the relic abundance of the particles \(\chi\) their
  mass. Unit: GeV. Status: STANDARD.
\item
  \(H_4\): The expansion rate (Hubble rate) of the fold, \(H_4 = a_4'\),
  as a function of the time or of the redshift; \(H_4^{(0)}\) its value
  today, the Hubble constant. Unit: GeV (1/time). Value:
  \(H_4^{(0)} = 100\,h_{100}\) km/s/Mpc with \(h_{100} = 0.674\) (Planck
  2018): 67.4 km/s/Mpc \(= 1.44\times10^{-42}\) GeV. Status: STANDARD.
\item
  \(\pi\): The number pi, the ratio of the circumference of a circle to
  its diameter. Unit: dimensionless. Value: 3.141592654. Status:
  STANDARD.
\item
  \(G_{\mathrm{N}}\): Newton's constant of four-dimensional gravity (the
  gravity of our world, on the fold),
  \(G_{\mathrm{N}} = 1/M_{\mathrm{Pl}}^2\) with the Planck mass
  \(M_{\mathrm{Pl}} = 1.221\times10^{19}\) GeV. Unit: GeV\(^{-2}\).
  Value: \(6.708\times10^{-39}\) GeV\(^{-2}\). Status: STANDARD.
\end{itemize}

The density parameters today: the dark matter (the particles \(\chi\),
whose origin is OPEN, FE-61), the baryons, the dark energy (the vacuum
energy of the fold, FE-59) and the massive neutrinos of the CNB (for the
three left-handed states with the temperature and the small non-thermal
distortion of the Standard Model; the right-handed states are not
populated). Together with \(A_{\mathrm{s}}\), \(n_{\mathrm{S}}\) and the
optical depth of reionisation they fix the CMB, the CNB and the BAO.
Status: STANDARD.

\hypertarget{the-thermal-history}{%
\section{The thermal history}\label{the-thermal-history}}

\[
\begin{aligned}
&H_4^2 = \frac{8\pi G_{\mathrm{N}}}{3}\big(\rho_{\mathrm{SM}} + \rho_\chi + \rho_\Lambda\big),\qquad
 2H_4' + 3H_4^2 = -8\pi G_{\mathrm{N}}\,\mathcal{P},\qquad
 \frac{1}{8\pi G_{\mathrm{N}}} = \frac{V_*}{\kappa_7} + \frac{1}{\kappa_{\mathrm{ind}}},\qquad \kappa_7 = \frac52\,\kappa H,\\
&\rho_\Lambda = \rho_{\mathrm{ff}} + \rho_{\mathrm{vac}} - \lambda_{\mathrm{F}}V_*,\qquad
 q_{\mathrm{SM}} + q_\chi + q_{\mathrm{P}} + q_{\mathrm{ff}} + \lambda_{\mathrm{F}}V_* + \frac{\rho - 3\mathcal{P}}{2} = 0,\\
&q_{\mathrm{SM}} = \big\langle\mathcal{L}_{\mathrm{H}} + \mathcal{L}_{\mathrm{f}} + \mathcal{L}_{\mathrm{Y}}\big\rangle,\qquad
 q_{\mathrm{P}} = \langle\mathcal{L}_{\mathrm{P}}\rangle
\end{aligned} \tag{FE-59}
\] \textbf{Symbols in (FE-59).}

\begin{itemize}
\tightlist
\item
  \(H_4\): The expansion rate (Hubble rate) of the fold, \(H_4 = a_4'\),
  as a function of the time or of the redshift; \(H_4^{(0)}\) its value
  today, the Hubble constant. Unit: GeV (1/time). Value:
  \(H_4^{(0)} = 100\,h_{100}\) km/s/Mpc with \(h_{100} = 0.674\) (Planck
  2018): 67.4 km/s/Mpc \(= 1.44\times10^{-42}\) GeV. Status: STANDARD.
\item
  \({}'\): The derivative of a function of the time \(x_4\) alone with
  respect to \(x_4\) (\(a_4' = \mathrm{d}a_4/\mathrm{d}x_4\); two primes
  the second derivative); on a function of \(\sigma\) alone (a profile)
  the derivative with respect to \(\sigma\). Status: STANDARD.
\item
  \(\pi\): The number pi, the ratio of the circumference of a circle to
  its diameter. Unit: dimensionless. Value: 3.141592654. Status:
  STANDARD.
\item
  \(G_{\mathrm{N}}\): Newton's constant of four-dimensional gravity (the
  gravity of our world, on the fold),
  \(G_{\mathrm{N}} = 1/M_{\mathrm{Pl}}^2\) with the Planck mass
  \(M_{\mathrm{Pl}} = 1.221\times10^{19}\) GeV. Unit: GeV\(^{-2}\).
  Value: \(6.708\times10^{-39}\) GeV\(^{-2}\). Status: STANDARD.
\item
  \(\rho_{\mathrm{SM}}\): The energy density of the Standard-Model
  particles. Unit: GeV\(^4\). Status: STANDARD.
\item
  \(\rho_\chi\): The energy density of the dark matter, the particles
  \(\chi\) of dirac16complex; \(\rho_\chi^{(0)}\) its value today. Unit:
  GeV\(^4\). Status: DERIVED.
\item
  \(\rho_\Lambda\): The vacuum energy density of the fold above its
  static tension, the dark energy of UM-4:
  \(\rho_\Lambda = \rho_{\mathrm{ff}} + \rho_{\mathrm{vac}} - \lambda_{\mathrm{F}}V_*\)
  (four-dimensional normalisation). Unit: GeV\(^4\). Value:
  \(0.685\,\rho_{\mathrm{c}}^{(0)} = 2.5\times10^{-47}\) GeV\(^4\)
  (Planck 2018). Status: DERIVED.
\item
  \(\mathcal{P}\): The total pressure of the matter of the fold;
  \(\Delta\mathcal{P}\) its perturbation. Unit: GeV\(^4\). Status:
  STANDARD.
\item
  \(V_*\): The physical volume of the extra times \(x_5, x_6, x_7\) on
  the fold in the late universe,
  \(V_* = V_{\mathrm{t}}(\mathrm{e}^{-a_4}b_4)^3\) (with
  \(V_{\mathrm{t}} = \int\mathrm{d}x_5\,\mathrm{d}x_6\,\mathrm{d}x_7\)
  the coordinate volume of the extra times) at a time when
  \(\mathrm{e}^{-a_4}b_4\) is constant; it converts the
  eight-dimensional normalisation into the four-dimensional one. Unit:
  length\(^3\) (GeV\(^{-3}\)). Status: ASSUMED.
\item
  \(\kappa_7\): The gravitational coupling of the fold: to the first
  order in its curvature over \(H^2\) the fold obeys the Einstein
  equations of seven dimensions
  \(\hat G^a{}_b = \kappa_7\mathsf{t}^a{}_b\) with
  \(\kappa_7 = 5\kappa H/2\), the coupling of the massless graviton of
  the bulk bound to the fold; four-dimensional gravity on the fold has
  \(8\pi G_{\mathrm{N}} = \kappa_7/V_*\) (plus the induced term, which
  is negligible). Unit: GeV\(^{-5}\). Value: \(5\kappa H/2\). Status:
  DERIVED.
\item
  \(\kappa_{\mathrm{ind}}\): The coupling of the Einstein-Hilbert term
  \(\hat R/(2\kappa_{\mathrm{ind}})\) that the quantum fluctuations of
  the fields of the fold may induce,
  \(\kappa_{\mathrm{ind}} = 8\pi G_{\mathrm{ind}}\); it adds to the
  gravity of the fold,
  \(1/(8\pi G_{\mathrm{N}}) = V_*/\kappa_7 + 1/\kappa_{\mathrm{ind}}\).
  Unit: GeV\(^{-2}\). Value: not fixed by UM-4; negligible for a cutoff
  of the order of \(H\). Status: ESTIMATE.
\item
  \(\kappa\): The eight-dimensional gravitational coupling,
  \(\kappa = 8\pi G_8\). Unit: GeV\(^{-6}\). Status: AUTHOR.
\item
  \(H\): The author's constant (an inverse length) that fixes the period
  \(\pi/(6H)\) of \(x_0\), the warp factor
  \(\mathrm{e}^{-2H\lvert\sigma\rvert}\) of the metric, the cosmological
  constant \(\Lambda = -21H^2\) and the tension
  \(\kappa\lambda_{\mathrm{F}} = 12H\) of the fold; in UM-4 a free
  parameter, bounded from below by the propagation of the gauge fields
  in the bulk. Unit: GeV (1/length). Value: \(H > 2534\) GeV, about 2.5
  TeV (from the invisible width of the Z boson). Status: AUTHOR.
\item
  \(\rho_{\mathrm{ff}}\): The energy density of the fold fluid. Unit:
  GeV\(^4\). Value:
  \(\lambda_{\mathrm{F}}V_* + \rho_\Lambda - \rho_{\mathrm{vac}}\):
  together with the vacuum energy \(\rho_{\mathrm{vac}}\) of the fields
  of the fold it makes up the tension \(\lambda_{\mathrm{F}}V_*\) that
  the static block requires plus the dark energy \(\rho_\Lambda\).
  Status: ASSUMED (the fluid), DERIVED (the value).
\item
  \(\rho_{\mathrm{vac}}\): The vacuum energy of the fields of the fold:
  the minimum of their potential (for the Standard Model the minimum of
  the Higgs potential with its loop corrections), a tension of the
  seven-dimensional fold. Unit: GeV\(^4\). Status: STANDARD.
\item
  \(\lambda_{\mathrm{F}}\): The tension that the static block requires
  of the fold, \(\kappa\lambda_{\mathrm{F}} = 12H\): the total vacuum
  energy of the seven-dimensional fold, made of the tension of the fold
  fluid and the vacuum energy of the fields of the fold (both have the
  pressure minus their energy density along ordinary space and along the
  extra times), per unit volume of the six other directions of the fold
  (eight-dimensional normalisation); in the four-dimensional
  normalisation it is \(\lambda_{\mathrm{F}}V_*\), and the excess of the
  total over \(\lambda_{\mathrm{F}}V_*\) is the dark energy
  \(\rho_\Lambda\). Unit: GeV\(^7\). Value: \(12H/\kappa\);
  \(\lambda_{\mathrm{F}}V_* = 30H^2M_{\mathrm{red}}^2 = 1.1\times10^{45}\)
  GeV\(^4\) for \(H = 2534\) GeV. Status: DERIVED.
\item
  \(q_{\mathrm{SM}}\): The pressure of the Standard Model along the
  extra times at the fold: the mean on-shell Lagrangian density of the
  Standard-Model fields (the extra-time component of the energy-momentum
  tensor of every term of the fold is its Lagrangian, because the fields
  of the fold do not depend on the extra times). Unit: GeV\(^4\).
  Status: DERIVED (for radiation and free fields), ASSUMED (for the
  baryons).
\item
  \(q_\chi\): The pressure of the particles \(\chi\) of dirac16complex
  along the extra times, integrated across the fold (their on-shell
  Lagrangian, the self-interaction energy per volume; zero for a dilute
  gas). Unit: GeV\(^4\). Status: DERIVED.
\item
  \(q_{\mathrm{P}}\): The pressure of the portals along the extra times:
  their mean on-shell Lagrangian density. Unit: GeV\(^4\). Status:
  DERIVED.
\item
  \(q_{\mathrm{ff}}\): The pressure of the fold fluid along the extra
  times: minus its energy density (its tension) plus the pressure
  \(q_{\mathrm{hold}}\) that holds the volume of the extra times fixed.
  Unit: GeV\(^4\). Value: \(-\rho_{\mathrm{ff}} + q_{\mathrm{hold}}\).
  Status: ASSUMED (the fluid), DERIVED (the value).
\item
  \(\rho\): The total energy density of the matter of the fold (all
  species); \(\Delta\rho\) its perturbation. Unit: GeV\(^4\). Status:
  STANDARD.
\item
  \(\langle\cdots\rangle\): The mean (expectation) value of an operator
  in the state of the medium (the vacuum, a plasma, the Kohn-Sham
  state). Status: STANDARD.
\item
  \(\mathcal{L}_{\mathrm{H}}\): The Lagrangian of the Higgs doublet on
  the fold. Unit: GeV\(^4\). Status: AUTHOR.
\item
  \(\mathcal{L}_{\mathrm{f}}\): The kinetic Lagrangian of the quarks and
  leptons on the fold. Unit: GeV\(^4\). Status: STANDARD.
\item
  \(\mathcal{L}_{\mathrm{Y}}\): The Yukawa Lagrangian of the quarks and
  leptons with the Higgs doublet. Unit: GeV\(^4\). Status: STANDARD.
\item
  \(\mathcal{L}_{\mathrm{P}}\): The Lagrangian of the portals of
  dirac16complex at the fold. Unit: GeV\(^4\). Status: INVENTED.
\end{itemize}

The Einstein equations of the fold with frozen extra times: FE-71 and
FE-72 with the sources of FE-5, written in the four-dimensional
normalisation (IN-8; \(H_4 = a_4'\) the expansion rate of the fold). The
first line is general relativity: the Friedmann equation and the
acceleration equation of a flat universe with Newton's constant
\(8\pi G_{\mathrm{N}} = \kappa_7/V_*\) of the graviton of the fold
(FE-63, FE-71), plus the induced term IN-9, which is negligible (FE-64),
with corrections of the relative order \((12/25)(H_4/H)^2\) from the
high-energy term of FE-35, below \(10^{-50}\) after nucleosynthesis for
\(H > 2534\) GeV. \(\rho_{\mathrm{SM}}\) is the energy density of the
Standard-Model particles (the gauge fields included: their zero modes
are the radiation of the fold, and their massive modes are not excited
at temperatures far below \(H\)), \(\rho_\chi\) that of the particles
\(\chi\) of dirac16complex, \(\mathcal{P}\) the total pressure along
ordinary space and
\(\rho = \rho_{\mathrm{SM}} + \rho_\chi + \rho_\Lambda\) the total
energy density above the tension of the static block. \(\rho_\Lambda\)
is the vacuum energy of the fold above that tension, the dark energy of
UM-4 (\(w = -1\) exactly): the tension of the fold fluid plus the vacuum
energy \(\rho_{\mathrm{vac}}\) of the fields of the fold, minus the
tension \(\lambda_{\mathrm{F}}V_* = 30H^2/(8\pi G_{\mathrm{N}})\) that
the static block requires (FE-65). UM-4 does not explain why the two
nearly cancel (for \(H = 2534\) GeV,
\(\lambda_{\mathrm{F}}V_* = 1.1\times10^{45}\) GeV\(^4\) against
\(\rho_\Lambda = 2.5\times10^{-47}\) GeV\(^4\), one part in
\(10^{92}\)): the cosmological-constant problem in the form it takes for
a brane of Randall and Sundrum (OPEN). The second line, last equation,
is the component along the extra times (the pressures
\(q_{\mathrm{SM}}\), \(q_\chi\), \(q_{\mathrm{P}}\) and
\(q_{\mathrm{ff}}\) along the extra times; for every term of the fold
its on-shell Lagrangian density by FE-5, the mean \(\langle\ \rangle\)
over the state of the fold; the gauge fields live in the bulk and their
pressure along the extra times enters through the bulk, where their zero
modes carry, like radiation, none on average): it is the radion
condition FE-72, the condition that the volume of the extra times stays
fixed, and with the fold fluid holding it fixed it determines
\(q_{\mathrm{hold}}\) in
\(q_{\mathrm{ff}} = -\rho_{\mathrm{ff}} + q_{\mathrm{hold}}\) (FE-65).
Status: DERIVED (\texttt{um4\_fold.py}, F2, F4 and F5;
\texttt{um4\_background.py}, BG2 to BG4), ASSUMED (that the fold fluid
holds the volume of the extra times fixed), OPEN (the tuning of the
tension).

\[
\rho_{\mathrm{SM}} = \frac{\pi^2}{30}\,g_\rho(\mathsf{T})\,\mathsf{T}^4,\qquad s = \frac{2\pi^2}{45}\,g_{\mathrm{s}}(\mathsf{T})\,\mathsf{T}^3,\qquad
H_4 = \sqrt{\frac{8\pi^3g_\rho}{90}}\,\frac{\mathsf{T}^2}{M_{\mathrm{Pl}}},\qquad M_{\mathrm{Pl}} = G_{\mathrm{N}}^{-1/2} \tag{FE-60}
\] \textbf{Symbols in (FE-60).}

\begin{itemize}
\tightlist
\item
  \(\rho_{\mathrm{SM}}\): The energy density of the Standard-Model
  particles. Unit: GeV\(^4\). Status: STANDARD.
\item
  \(\pi\): The number pi, the ratio of the circumference of a circle to
  its diameter. Unit: dimensionless. Value: 3.141592654. Status:
  STANDARD.
\item
  \(g_\rho\): The effective number of relativistic degrees of freedom of
  the energy density. Unit: dimensionless. Value: 106.75 in the Standard
  Model above the electroweak scale. Status: STANDARD.
\item
  \(\mathsf{T}\): The temperature of the Standard-Model plasma (of the
  photons after the annihilation of electrons and positrons). Unit: GeV.
  Status: STANDARD.
\item
  \(s\): The entropy density of the Standard-Model plasma. Unit:
  GeV\(^3\). Status: STANDARD.
\item
  \(g_{\mathrm{s}}\): The effective number of relativistic degrees of
  freedom of the entropy density. Unit: dimensionless. Value: 3.909
  today. Status: STANDARD.
\item
  \(H_4\): The expansion rate (Hubble rate) of the fold, \(H_4 = a_4'\),
  as a function of the time or of the redshift; \(H_4^{(0)}\) its value
  today, the Hubble constant. Unit: GeV (1/time). Value:
  \(H_4^{(0)} = 100\,h_{100}\) km/s/Mpc with \(h_{100} = 0.674\) (Planck
  2018): 67.4 km/s/Mpc \(= 1.44\times10^{-42}\) GeV. Status: STANDARD.
\item
  \(M_{\mathrm{Pl}}\): The Planck mass,
  \(M_{\mathrm{Pl}} = G_{\mathrm{N}}^{-1/2}\). Unit: GeV. Value:
  \(1.221\times10^{19}\) GeV. Status: STANDARD.
\item
  \(G_{\mathrm{N}}\): Newton's constant of four-dimensional gravity (the
  gravity of our world, on the fold),
  \(G_{\mathrm{N}} = 1/M_{\mathrm{Pl}}^2\) with the Planck mass
  \(M_{\mathrm{Pl}} = 1.221\times10^{19}\) GeV. Unit: GeV\(^{-2}\).
  Value: \(6.708\times10^{-39}\) GeV\(^{-2}\). Status: STANDARD.
\end{itemize}

The radiation of the Standard Model in thermal equilibrium at the
temperature \(\mathsf{T}\): the energy density and the entropy density
with the effective numbers of relativistic degrees of freedom \(g_\rho\)
and \(g_{\mathrm{s}}\) (106.75 above the electroweak scale in the
Standard Model; three right-handed neutrinos in equilibrium would add
\(\frac78\times6 = 5.25\); the particles \(\chi\) add theirs while they
are relativistic), and the expansion rate of the radiation era. The
gauge fields count with the two polarisations of their zero modes for
\(\mathsf{T} \ll H\). Status: STANDARD (\texttt{um4\_thermal.py}, T1).

\[
\frac{\mathrm{d}n_\chi}{\mathrm{d}x_4} + 3H_4\,n_\chi = -\langle\sigma_{\mathrm{ann}}\mathsf{v}\rangle\big(n_\chi^2 - n_{\chi,\mathrm{eq}}^2\big),\qquad
n_{\chi,\mathrm{eq}} = g_\chi\Big(\frac{\omega_{\mathrm{b}}\mathsf{T}}{2\pi}\Big)^{3/2}\mathrm{e}^{-\omega_{\mathrm{b}}/\mathsf{T}},\qquad
\Omega_\chi h_{100}^2 \simeq \frac{1.07\times10^9\ \text{GeV}^{-1}\,x_{\mathrm{f}}}{\sqrt{g_*}\,M_{\mathrm{Pl}}\,\langle\sigma_{\mathrm{ann}}\mathsf{v}\rangle} \tag{FE-61}
\] \textbf{Symbols in (FE-61).}

\begin{itemize}
\tightlist
\item
  \(\mathrm{d}\): The differential: \(\mathrm{d}\) in front of a
  quantity is an infinitesimal change of it (as \(\mathrm{d}\sigma\), a
  small step of \(\sigma\)); the quotient of two such changes is a
  derivative; \(\mathrm{d}\) with an exponent, in front of several
  variables, is their volume element, the product of their
  differentials. Status: STANDARD.
\item
  \(n_\chi\): The number density of the dark-matter particles \(\chi\)
  (particles plus antiparticles). Unit: GeV\(^3\). Status: DERIVED.
\item
  \(x_4\): The time; a prime on a function of \(x_4\) alone is
  \(\mathrm{d}/\mathrm{d}x_4\). Unit: length (= time, \(c = 1\)).
  Status: AUTHOR.
\item
  \(H_4\): The expansion rate (Hubble rate) of the fold, \(H_4 = a_4'\),
  as a function of the time or of the redshift; \(H_4^{(0)}\) its value
  today, the Hubble constant. Unit: GeV (1/time). Value:
  \(H_4^{(0)} = 100\,h_{100}\) km/s/Mpc with \(h_{100} = 0.674\) (Planck
  2018): 67.4 km/s/Mpc \(= 1.44\times10^{-42}\) GeV. Status: STANDARD.
\item
  \(\sigma_{\mathrm{ann}}\): The annihilation cross-section of the
  particles \(\chi\) into Standard-Model particles (through the portals)
  in the normalisation of the relic abundance, where \(n_\chi\) counts
  the four particle and the four antiparticle states of the level at the
  fold: \(\langle\sigma_{\mathrm{ann}}\mathsf{v}\rangle\) is a quarter
  of the annihilation rate \((\sigma\mathsf{v})_{\chi\bar\chi}\) of one
  particle \(\chi_a\) with one antiparticle \(\bar\chi_a\). Unit:
  GeV\(^{-2}\). Status: DERIVED.
\item
  \(\mathsf{v}\): The relative velocity of two annihilating particles;
  \(\langle\sigma_{\mathrm{ann}}\mathsf{v}\rangle\) is the thermal
  average of the product. Unit: dimensionless. Status: STANDARD.
\item
  \(\langle\cdots\rangle\): The mean (expectation) value of an operator
  in the state of the medium (the vacuum, a plasma, the Kohn-Sham
  state). Status: STANDARD.
\item
  \(n_{\chi,\mathrm{eq}}\): The number density of \(\chi\) in thermal
  equilibrium with the plasma (non-relativistic). Unit: GeV\(^3\).
  Status: STANDARD.
\item
  \(g_\chi\): The number of internal states of the level of \(\chi\) at
  the fold that are in thermal equilibrium. Unit: dimensionless. Value:
  8 (four particles and four antiparticles). Status: DERIVED.
\item
  \(\omega_{\mathrm{b}}\): The energy at zero momentum of the state of
  dirac16complex at the fold in the limit \(H \to 0\) (the bound state
  of UM-3): the particle \(\chi\) of UM-4 has the mass
  \(\omega_{\mathrm{r}} = \omega_{\mathrm{b}}(1 - H/(2q_{\mathrm{b}})) + \dots\),
  equal to \(\omega_{\mathrm{b}}\) up to the relative correction
  \(H/(2q_{\mathrm{b}})\). Unit: GeV. Status: DERIVED.
\item
  \(\mathsf{T}\): The temperature of the Standard-Model plasma (of the
  photons after the annihilation of electrons and positrons). Unit: GeV.
  Status: STANDARD.
\item
  \(\pi\): The number pi, the ratio of the circumference of a circle to
  its diameter. Unit: dimensionless. Value: 3.141592654. Status:
  STANDARD.
\item
  \(\mathrm{e}\): Euler's number, the base of the natural logarithm;
  \(\mathrm{e}^{X}\) is the exponential of \(X\). Unit: dimensionless.
  Value: 2.718281828. Status: STANDARD.
\item
  \(\Omega_\chi\): The density parameter of the dark matter today. Unit:
  dimensionless. Value: \(\Omega_\chi h_{100}^2 = 0.120\) (Planck 2018).
  Status: DERIVED.
\item
  \(h_{100}\): The present expansion rate of the fold in units of 100
  km/s/Mpc. Unit: dimensionless. Value: 0.674 (Planck 2018). Status:
  STANDARD.
\item
  \(x_{\mathrm{f}}\): The ratio of the mass of \(\chi\) to the
  temperature at which its annihilations stop (freeze-out). Unit:
  dimensionless. Value: 23 to 30 (for \(\omega_{\mathrm{b}}\) = 10 GeV
  to 10 TeV). Status: STANDARD.
\item
  \(g_*\): The effective number of relativistic degrees of freedom at
  the freeze-out (\(g_\rho\) there). Unit: dimensionless. Status:
  STANDARD.
\item
  \(M_{\mathrm{Pl}}\): The Planck mass,
  \(M_{\mathrm{Pl}} = G_{\mathrm{N}}^{-1/2}\). Unit: GeV. Value:
  \(1.221\times10^{19}\) GeV. Status: STANDARD.
\end{itemize}

The relic density of the particles \(\chi\) of dirac16complex if they
were once in thermal equilibrium: their number density \(n_\chi\)
(particles plus antiparticles) follows the Boltzmann equation with the
thermally averaged annihilation cross-section
\(\langle\sigma_{\mathrm{ann}}\mathsf{v}\rangle\) into Standard-Model
particles through the portal couplings of FE-47 (Higgs exchange, pairs
of gluons and photons, pairs of fermions); \(g_\chi\) is the number of
states of the level that are in equilibrium (8: four particles and four
antiparticles). The freeze-out at
\(x_{\mathrm{f}} = \omega_{\mathrm{b}}/\mathsf{T}_{\mathrm{f}} \approx 20\)
to 30 (\(\mathsf{T}_{\mathrm{f}}\) the temperature at which the
annihilations stop) gives the standard estimate of the relic density for
an annihilation that does not depend on the velocity (\(g_*\) the
degrees of freedom at freeze-out); the numerical solution (integrated to
\(x = 10^5\), where it has converged to \(10^{-4}\)) gives
\(\Omega_\chi h_{100}^2 = 0.12\) for
\(\langle\sigma_{\mathrm{ann}}\mathsf{v}\rangle = 2.14\) to
\(2.29\times10^{-26}\) cm\(^3\)/s from 10 GeV to 10 TeV, and the
estimate agrees with it within 10 per cent. The largest annihilation
that partial-wave unitarity allows,
\((\sigma\mathsf{v})_{\chi\bar\chi} = 4\pi/(m^2\mathsf{v})\) for the
s-wave, gives \(\Omega_\chi h_{100}^2 = 0.12\) at
\(\omega_{\mathrm{b}} = 9.7\times10^4\) GeV (\(6.9\times10^4\) GeV if
only the singlet of the extra-time SU(2) annihilates, \(1.37\times10^5\)
GeV if every channel is open), and heavier particles would be
overproduced. The lifetime of the particles \(\chi\) (FE-86) puts
\(\omega_{\mathrm{b}}\) above this mass unless they are strongly bound,
\(q_{\mathrm{b}} > 0.906\,m_\Psi\) (0.875 to 0.929 over the range of the
bookkeeping): only strongly bound particles \(\chi\) can be a thermal
relic. The portals of IN-16 cannot bind them so strongly. If the fold
binds them through the Higgs portal (FE-47), the limit of LZ on their
scattering on nuclei (Phys. Rev.~Lett. 135 (2025) 011802, extrapolated
in proportion to the mass above 40 GeV) allows at most
\(q_{\mathrm{b}}/\omega_{\mathrm{b}} = 0.019\) at 10 TeV, falling like
\(\omega_{\mathrm{b}}^{-1/4}\); with that binding the width of FE-86
(\(F \simeq (q_{\mathrm{b}}/\omega_{\mathrm{b}})^3/3\)) stays below that
of a lifetime of 160 Gyr only for \(H\) below \(1.5\times10^{-3}\) GeV
at \(10^6\) GeV and below 2.5 GeV at the Planck mass, never for
\(H > 2534\) GeV. A binding by the gauge-kinetic portal of the gluons
acts through the gluon condensate
\(\langle(\alpha_3/\pi)G^2\rangle = 0.012\) GeV\(^4\) of the vacuum
(\(\alpha_3 = g_3^2/(4\pi)\) the strong coupling), and by the theorem of
Feynman and Hellmann (FE-47) the same coupling makes \(\chi\) scatter on
nucleons, because the matrix element of \((\alpha_3/\pi)G^2\) in a
nucleon at rest is \(-\frac89m_{\mathrm{N}}f_{\mathrm{TG}}\) (the trace
anomaly; \(m_{\mathrm{N}}\) the nucleon mass, \(f_{\mathrm{TG}} = 0.92\)
its gluonic fraction): the nucleon coupling is
\(f_{\mathrm{N}} = 8m_{\mathrm{N}}f_{\mathrm{TG}}q_{\mathrm{b}}^2/(9\langle(\alpha_3/\pi)G^2\rangle\omega_{\mathrm{b}}) = 64\,q_{\mathrm{b}}^2/\omega_{\mathrm{b}}\)
GeV\(^{-3}\) (\(\alpha_3\) cancels), and the scattering cross-section
\(\sigma_{\mathrm{SI}} = \mu_{\mathrm{red}}^2f_{\mathrm{N}}^2/\pi\) is
\(9\times10^{-14}\) to \(8\times10^{-12}\) cm\(^2\) at the lifetime
bound of every binding, 29 to 33 orders of magnitude above the limit of
LZ. The other condensates of the portals vanish in the vacuum
(\(\langle W^2\rangle\), \(\langle B^2\rangle\), the currents and the
field strength \(B_{\mu\nu}\)). The dark matter of UM-4 therefore needs
a binding of \(\chi\) to the fold that does not couple it to nucleons,
and, unless that binding is strong, a non-thermal origin (both OPEN).
The annihilation through the current portals is velocity-suppressed:
their convective current creates a pair \(\chi\bar\chi\) only with an
amplitude of the second order in the momentum (FE-67), so a pair
annihilates into a massless chiral fermion state with
\((\sigma\mathsf{v})_{\chi\bar\chi} = \mathfrak{g}_{f}^2\,\omega_{\mathrm{b}}^2\,\mathsf{v}^4/(96\pi)\)
(\(\mathfrak{g}_{f}\) the coupling of FE-47, \(\mathsf{v}\) the relative
velocity; the thermal average of \(\mathsf{v}^4\) is
\(60\mathsf{T}^2/\omega_{\mathrm{b}}^2\)). For the masses 10 GeV to
\(10^5\) GeV a relic through the quark current portals is excluded by
direct detection by factors \(4\times10^{10}\) (100 GeV) to 60 (\(10^5\)
GeV); one through the lepton portals scatters on nuclei only at one loop
and is allowed by LZ only above about \(2.3\times10^4\) GeV, where the
strength \(\mathfrak{g}_{\mathrm{l}}\omega_{\mathrm{b}}^2\) of the
contact coupling that the relic needs is about 80, beyond the
perturbative range (it exceeds \(4\pi\) above about \(4\times10^3\)
GeV): the current portals cannot set the relic density of a
perturbatively coupled \(\chi\). The Higgs portal requires
\(q_{\mathrm{b}}/\omega_{\mathrm{b}} < 0.062\) to 0.019 from 100 GeV to
10 TeV. Status: STANDARD (the equation; UM-4: the cross-section from the
portals), DERIVED (\texttt{um4\_thermal.py}, T2 and T7; the velocity
dependence, \texttt{um4\_krein.py}, KR9), ESTIMATE (T3, T4, T8 and T9:
the limit of LZ and its extrapolation, the loop cut-off of the lepton
portals, the condensate and the gluonic fraction of the nucleon), OPEN
(the origin of the particles \(\chi\) and a binding that escapes direct
detection).

\[
\frac{\mathrm{d}\rho_{\nu_{\mathrm{R}}}}{\mathrm{d}x_4} + 4H_4\,\rho_{\nu_{\mathrm{R}}} = \mathcal{C}_{\nu_{\mathrm{R}}},\qquad
\Delta N_{\mathrm{eff}} = 3\Big(\frac{10.75}{g_{\mathrm{s}}(\mathsf{T}_{\mathrm{dec}})}\Big)^{4/3}\ \ \text{(three right-handed neutrinos decoupling at } \mathsf{T}_{\mathrm{dec}}) \tag{FE-62}
\] \textbf{Symbols in (FE-62).}

\begin{itemize}
\tightlist
\item
  \(\mathrm{d}\): The differential: \(\mathrm{d}\) in front of a
  quantity is an infinitesimal change of it (as \(\mathrm{d}\sigma\), a
  small step of \(\sigma\)); the quotient of two such changes is a
  derivative; \(\mathrm{d}\) with an exponent, in front of several
  variables, is their volume element, the product of their
  differentials. Status: STANDARD.
\item
  \(\rho_{\nu_{\mathrm{R}}}\): The energy density of the three
  right-handed neutrinos. Unit: GeV\(^4\). Status: STANDARD.
\item
  \(x_4\): The time; a prime on a function of \(x_4\) alone is
  \(\mathrm{d}/\mathrm{d}x_4\). Unit: length (= time, \(c = 1\)).
  Status: AUTHOR.
\item
  \(H_4\): The expansion rate (Hubble rate) of the fold, \(H_4 = a_4'\),
  as a function of the time or of the redshift; \(H_4^{(0)}\) its value
  today, the Hubble constant. Unit: GeV (1/time). Value:
  \(H_4^{(0)} = 100\,h_{100}\) km/s/Mpc with \(h_{100} = 0.674\) (Planck
  2018): 67.4 km/s/Mpc \(= 1.44\times10^{-42}\) GeV. Status: STANDARD.
\item
  \(\mathcal{C}_{\nu_{\mathrm{R}}}\): The collision term of the
  right-handed neutrinos: the rate at which collisions produce (or
  remove) their energy density. Unit: GeV\(^5\). Status: STANDARD.
\item
  \(\Delta\): A change, difference or small variation of the quantity
  that follows (for example an interval of the time, a perturbation of a
  density, a difference of two squared masses, or an addition to an
  effective number of species). Status: STANDARD.
\item
  \(N_{\mathrm{eff}}\): The effective number of neutrino species: the
  energy density of the relativistic particles other than photons in
  units of that of one neutrino species whose temperature is
  \((4/11)^{1/3}\) times that of the photons. Unit: dimensionless.
  Value: 3.044 (Standard Model); measured \(2.99 \pm 0.17\) (Planck
  2018). Status: STANDARD.
\item
  \(g_{\mathrm{s}}\): The effective number of relativistic degrees of
  freedom of the entropy density. Unit: dimensionless. Value: 3.909
  today. Status: STANDARD.
\item
  \(\mathsf{T}_{\mathrm{dec}}\): The temperature at which the
  right-handed neutrinos decouple from the plasma. Unit: GeV. Status:
  STANDARD.
\end{itemize}

The right-handed neutrinos in the early universe: their energy density
grows by the collision term \(\mathcal{C}_{\nu_{\mathrm{R}}}\)
(production by the neutrino Yukawa couplings, negligible for Dirac
masses below 0.1 eV because it is suppressed by the square of the mass
over the temperature, and by the current portal \(c^{\nu}_X\) with
dirac16complex), and if they were in equilibrium until the temperature
\(\mathsf{T}_{\mathrm{dec}}\) they add \(\Delta N_{\mathrm{eff}}\) to
the effective number of neutrinos (from the conservation of entropy
after their decoupling). The measured
\(N_{\mathrm{eff}} = 2.99 \pm 0.17\) (Planck 2018) and \(2.89 \pm 0.11\)
(ACT DR6, the sixth data release of the Atacama Cosmology Telescope,
with Planck, lensing, BAO and nucleosynthesis, 2025) bound
\(\Delta N_{\mathrm{eff}}\) and with it the portal couplings
\(c^{\nu}_X\): if the portal kept the right-handed neutrinos in
equilibrium with \(\chi\) (with the annihilation of FE-61) they would
decouple at \(\mathsf{T}_{\mathrm{dec}} \approx \omega_{\mathrm{b}}/15\)
to \(\omega_{\mathrm{b}}/18\) and add \(\Delta N_{\mathrm{eff}} = 0.14\)
to \(0.24\) (for \(\omega_{\mathrm{b}} = 10\) TeV to 10 GeV; at the
masses of UM-4 the decoupling lies far above the electroweak scale,
\(\Delta N_{\mathrm{eff}} = 0.141\)), 2.7 to 3.6 standard deviations
above the ACT value; the right-handed neutrinos of UM-4 must therefore
never thermalise (\(c^{\nu}_X\) much smaller than the lepton-doublet
couplings), and then \(N_{\mathrm{eff}} = 3.044\). Status: STANDARD (the
equations), DERIVED (the bound; \texttt{um4\_thermal.py}, T5).

\hypertarget{the-primordial-era}{%
\subsection{The primordial era}\label{the-primordial-era}}

\[
\begin{aligned}
&g_{MN}\,\mathrm{d}X^M\mathrm{d}X^N = \mathrm{d}\sigma^2 + f_{\mathrm{dS}}^2\Big[\mathrm{e}^{2H_{\mathrm{i}}x_4}\big(\mathrm{d}\vec x^2 - \mathrm{d}\vec x_5^2\big) - \mathrm{d}x_4^2\Big],\qquad
 f_{\mathrm{dS}} = \frac{\sinh\big(H(\sigma_0 - \lvert\sigma\rvert)\big)}{\sinh(H\sigma_0)},\qquad H_{\mathrm{i}} = \frac{H}{\sinh(H\sigma_0)},\\
&\kappa\big(\lambda_{\mathrm{F}} + \rho_{\mathrm{i}}\big) = 12H\coth(H\sigma_0),\qquad
 H_{\mathrm{i}}^2 = \Big(\frac{\kappa(\lambda_{\mathrm{F}} + \rho_{\mathrm{i}})}{12}\Big)^2 - H^2 = \frac{\kappa_7\rho_{\mathrm{i}}}{15}\Big(1 + \frac{\rho_{\mathrm{i}}}{2\lambda_{\mathrm{F}}}\Big)
\end{aligned} \tag{FE-87}
\] \textbf{Symbols in (FE-87).}

\begin{itemize}
\tightlist
\item
  \(g_{MN}\): The metric tensor of the eight-dimensional spacetime
  (signature \(+,+,+,+,-,-,-,-\) in the order \(x_0, \dots, x_7\)):
  \(g_{MN}\,\mathrm{d}X^M\mathrm{d}X^N\) is the squared interval,
  \(g^{MN}\) its inverse, \(\lvert g\rvert\) the absolute value of its
  determinant (\(\sqrt{\lvert g\rvert}\,\mathrm{d}^8X\) is the proper
  volume element). Unit: dimensionless. Status: AUTHOR.
\item
  \(M\): A coordinate index of eight dimensions (also with a numbered
  copy \(M_1, M_2, \dots\)): it runs over \(x_0, x_1, \dots, x_7\); a
  repeated index is summed. Status: STANDARD.
\item
  \(N\): A coordinate index of eight dimensions (also
  \(N_1, N_2, \dots\)), running over \(x_0, \dots, x_7\); a repeated
  index is summed. Status: STANDARD.
\item
  \(\mathrm{d}\): The differential: \(\mathrm{d}\) in front of a
  quantity is an infinitesimal change of it (as \(\mathrm{d}\sigma\), a
  small step of \(\sigma\)); the quotient of two such changes is a
  derivative; \(\mathrm{d}\) with an exponent, in front of several
  variables, is their volume element, the product of their
  differentials. Status: STANDARD.
\item
  \(X = \{x_0, \dots, x_7\}\): The author's coordinates of the
  eight-dimensional spacetime, \(X = \{x_0, x_1, \dots, x_7\}\) with the
  hidden direction first (Dirac\_Einstein-Lovelock\_4+4\_Nash.nb,
  section DEFCHART); \(X^M\) is the coordinate with the index \(M\) and
  \(\mathrm{d}^8X = \mathrm{d}x_0\,\mathrm{d}x_1\cdots\mathrm{d}x_7\)
  the coordinate volume element. Unit: length (each coordinate). Status:
  AUTHOR.
\item
  \(\sigma\): The proper coordinate of the hidden direction:
  \(\mathrm{d}\sigma = \lvert\cot z\rvert\,\mathrm{d}x_0\),
  \(\sigma = \mathrm{sgn}(\pi/2 - z)\ln(\sin z)/(6H)\) on one period
  \(0 < z < \pi\); \(\sigma = 0\) is the fold \(z = \pi/2\),
  \(\sigma \to \mp\infty\) the tips \(z \to 0, \pi\); one period of
  \(x_0\) is the whole line \(-\infty < \sigma < \infty\). Unit: length.
  Status: DERIVED.
\item
  \(f_{\mathrm{dS}}\): The warp of the de Sitter fold: the factor of the
  lengths of the fold at the depth \(\sigma\). Unit: dimensionless.
  Status: DERIVED.
\item
  \(\mathrm{e}\): Euler's number, the base of the natural logarithm;
  \(\mathrm{e}^{X}\) is the exponential of \(X\). Unit: dimensionless.
  Value: 2.718281828. Status: STANDARD.
\item
  \(H_{\mathrm{i}}\): A constant expansion rate of ordinary space: of
  the de Sitter fold (whose extra times expand at the same rate) or of
  an inflation of the fold with frozen extra times. Unit: GeV. Status:
  DERIVED (FE-87), ASSUMED (FE-88).
\item
  \(x_4\): The time; a prime on a function of \(x_4\) alone is
  \(\mathrm{d}/\mathrm{d}x_4\). Unit: length (= time, \(c = 1\)).
  Status: AUTHOR.
\item
  \(\vec x\): The position in ordinary space,
  \(\vec x = (x_1, x_2, x_3)\). Unit: length. Status: AUTHOR.
\item
  \(\vec x_5\): The position along the extra times,
  \(\vec x_5 = (x_5, x_6, x_7)\);
  \(\mathrm{d}\vec x_5^2 = \mathrm{d}x_5^2 + \mathrm{d}x_6^2 + \mathrm{d}x_7^2\).
  Unit: length. Status: AUTHOR.
\item
  \(\sigma_0\): The depth at which the warp of the de Sitter fold
  vanishes (a horizon of the bulk),
  \(\sinh(H\sigma_0) = H/H_{\mathrm{i}}\). Unit: length. Status:
  DERIVED.
\item
  \(H\): The author's constant (an inverse length) that fixes the period
  \(\pi/(6H)\) of \(x_0\), the warp factor
  \(\mathrm{e}^{-2H\lvert\sigma\rvert}\) of the metric, the cosmological
  constant \(\Lambda = -21H^2\) and the tension
  \(\kappa\lambda_{\mathrm{F}} = 12H\) of the fold; in UM-4 a free
  parameter, bounded from below by the propagation of the gauge fields
  in the bulk. Unit: GeV (1/length). Value: \(H > 2534\) GeV, about 2.5
  TeV (from the invisible width of the Z boson). Status: AUTHOR.
\item
  \(\kappa\): The eight-dimensional gravitational coupling,
  \(\kappa = 8\pi G_8\). Unit: GeV\(^{-6}\). Status: AUTHOR.
\item
  \(\lambda_{\mathrm{F}}\): The tension that the static block requires
  of the fold, \(\kappa\lambda_{\mathrm{F}} = 12H\): the total vacuum
  energy of the seven-dimensional fold, made of the tension of the fold
  fluid and the vacuum energy of the fields of the fold (both have the
  pressure minus their energy density along ordinary space and along the
  extra times), per unit volume of the six other directions of the fold
  (eight-dimensional normalisation); in the four-dimensional
  normalisation it is \(\lambda_{\mathrm{F}}V_*\), and the excess of the
  total over \(\lambda_{\mathrm{F}}V_*\) is the dark energy
  \(\rho_\Lambda\). Unit: GeV\(^7\). Value: \(12H/\kappa\);
  \(\lambda_{\mathrm{F}}V_* = 30H^2M_{\mathrm{red}}^2 = 1.1\times10^{45}\)
  GeV\(^4\) for \(H = 2534\) GeV. Status: DERIVED.
\item
  \(\rho_{\mathrm{i}}\): The vacuum energy of the fold above the
  critical tension \(\lambda_{\mathrm{F}}\) that expands the de Sitter
  fold (eight-dimensional normalisation). Unit: GeV\(^7\). Status:
  DERIVED.
\item
  \(\kappa_7\): The gravitational coupling of the fold: to the first
  order in its curvature over \(H^2\) the fold obeys the Einstein
  equations of seven dimensions
  \(\hat G^a{}_b = \kappa_7\mathsf{t}^a{}_b\) with
  \(\kappa_7 = 5\kappa H/2\), the coupling of the massless graviton of
  the bulk bound to the fold; four-dimensional gravity on the fold has
  \(8\pi G_{\mathrm{N}} = \kappa_7/V_*\) (plus the induced term, which
  is negligible). Unit: GeV\(^{-5}\). Value: \(5\kappa H/2\). Status:
  DERIVED.
\end{itemize}

An exact solution of UM-4 with an inflating fold: a fold with the
tension \(\lambda_{\mathrm{F}} + \rho_{\mathrm{i}}\) above the critical
one (the vacuum energy \(\rho_{\mathrm{i}}\) of the fold,
seven-dimensional normalisation) in a bulk of the constant curvature
\(-H^2\) (the static block's bulk, sliced differently). The metric of
the fold is a de Sitter space of seven dimensions with the expansion
rate \(H_{\mathrm{i}}\) in all its space-like and time-like directions
(\(\mathrm{d}\vec x^2\) and \(\mathrm{d}\vec x_5^2\) the squared steps
along ordinary space and along the extra times), and the warp factor
\(f_{\mathrm{dS}}^2\) vanishes at the depth \(\sigma_0\) (a horizon of
the bulk). The bulk equations hold exactly
(\(G^M{}_N = 21H^2\delta^M_N\) in all eight directions), Israel's
condition gives the tension of the second line, and the expansion rate
obeys the Friedmann equation of seven dimensions with the correction of
high energies \(\rho_{\mathrm{i}}/(2\lambda_{\mathrm{F}})\), as for the
brane of Randall and Sundrum (DERIVED, \texttt{um4\_primordial.py}, PR1,
exactly). It is not the primordial era of our universe: the extra times
inflate together with ordinary space, so the factor \(\mathcal{Z}\) of
the couplings of the fold grows like \(\mathrm{e}^{3N}\) after \(N\)
e-folds and every coupling of the fold falls like \(\mathrm{e}^{-3N}\),
by \(7\times10^{-79}\) after 60 e-folds (PR5). Status: DERIVED.

\[
\mathsf{c}' = 0,\quad a_4' = H_{\mathrm{i}}:\qquad \kappa_7\rho_7 = 3H_{\mathrm{i}}^2,\qquad P_7 = -\rho_7,\qquad q_7 = -2\rho_7,\qquad
\frac{H_{\mathrm{i}}}{H} < \frac{1}{10} \tag{FE-88}
\] \textbf{Symbols in (FE-88).}

\begin{itemize}
\tightlist
\item
  \(\mathsf{c}\): The logarithm of the scale factor of the extra times
  on the fold, \(\mathsf{c} = \ln(\mathrm{e}^{-a_4}b_4)\);
  \(\mathsf{c}' = \mathsf{V}H\); frozen extra times have
  \(\mathsf{c}' = 0\). Unit: dimensionless. Status: DERIVED.
\item
  \({}'\): The derivative of a function of the time \(x_4\) alone with
  respect to \(x_4\) (\(a_4' = \mathrm{d}a_4/\mathrm{d}x_4\); two primes
  the second derivative); on a function of \(\sigma\) alone (a profile)
  the derivative with respect to \(\sigma\). Status: STANDARD.
\item
  \(a_4\): The author's metric function of the time \(x_4\) that scales
  ordinary 3-space: on the fold \(\mathrm{e}^{a_4}\) is the scale factor
  of cosmology and \(a_4'\) the Hubble rate. Unit: dimensionless.
  Status: AUTHOR.
\item
  \(H_{\mathrm{i}}\): A constant expansion rate of ordinary space: of
  the de Sitter fold (whose extra times expand at the same rate) or of
  an inflation of the fold with frozen extra times. Unit: GeV. Status:
  DERIVED (FE-87), ASSUMED (FE-88).
\item
  \(\kappa_7\): The gravitational coupling of the fold: to the first
  order in its curvature over \(H^2\) the fold obeys the Einstein
  equations of seven dimensions
  \(\hat G^a{}_b = \kappa_7\mathsf{t}^a{}_b\) with
  \(\kappa_7 = 5\kappa H/2\), the coupling of the massless graviton of
  the bulk bound to the fold; four-dimensional gravity on the fold has
  \(8\pi G_{\mathrm{N}} = \kappa_7/V_*\) (plus the induced term, which
  is negligible). Unit: GeV\(^{-5}\). Value: \(5\kappa H/2\). Status:
  DERIVED.
\item
  \(\rho_7\): The energy density of the matter of the fold above its
  tension in the seven-dimensional normalisation (minus the component
  \(\mathsf{t}^{(4)}{}_{(4)}\): an energy per unit volume of the six
  other directions of the fold); \(V_*\rho_7\) in the four-dimensional
  normalisation. Unit: GeV\(^7\). Status: DERIVED.
\item
  \(P_7\): The pressure along ordinary space of the matter of the fold
  above its tension in the seven-dimensional normalisation (the
  component \(\mathsf{t}^{(i)}{}_{(i)}\)). Unit: GeV\(^7\). Status:
  DERIVED.
\item
  \(q_7\): The pressure along the extra times of the matter of the fold
  above its tension in the seven-dimensional normalisation (the
  component \(\mathsf{t}^{(r)}{}_{(r)}\)). Unit: GeV\(^7\). Status:
  DERIVED.
\item
  \(H\): The author's constant (an inverse length) that fixes the period
  \(\pi/(6H)\) of \(x_0\), the warp factor
  \(\mathrm{e}^{-2H\lvert\sigma\rvert}\) of the metric, the cosmological
  constant \(\Lambda = -21H^2\) and the tension
  \(\kappa\lambda_{\mathrm{F}} = 12H\) of the fold; in UM-4 a free
  parameter, bounded from below by the propagation of the gauge fields
  in the bulk. Unit: GeV (1/length). Value: \(H > 2534\) GeV, about 2.5
  TeV (from the invisible width of the Z boson). Status: AUTHOR.
\end{itemize}

The inflation of the fold with frozen extra times (the primordial era
that UM-4 assumes): a vacuum energy \(\rho_7\) of the fold expands
ordinary space at the constant rate \(H_{\mathrm{i}}\) by the Friedmann
equation of FE-72, and the radion condition requires the pressure
\(q_7 = -2\rho_7\) along the extra times, twice that of a vacuum energy:
the fold fluid adds \(q_{\mathrm{hold}} = -\rho_7\) (FE-65). These are
the equations of the first order; the exact equations FE-35 of a fold
with this extrinsic curvature have no solution in a bulk of constant
curvature for \(H_{\mathrm{i}} \ne 0\), so the inflation with frozen
extra times needs a Weyl tensor of the bulk, whose solution is not
computed (OPEN). The regime of general relativity,
\((12/25)(H_{\mathrm{i}}/H)^2 \ll 1\), needs \(H_{\mathrm{i}} < H/10\)
(fourth relation). The bound \(r_{\mathrm{T}} < 0.032\) on the ratio of
tensor to scalar perturbations limits any inflation to
\(H_{\mathrm{i}} < 4.4\times10^{13}\) GeV, so an inflation at that bound
needs \(H > 4.4\times10^{14}\) GeV; with \(H\) at its lower bound of
2534 GeV only an inflation with \(H_{\mathrm{i}} < 253\) GeV (an energy
scale \(\rho^{1/4} < 3.3\times10^{10}\) GeV) is in the regime of general
relativity. The Higgs field could drive the inflation (Higgs inflation,
with the coupling \(\xi_{\mathrm{H}} = 1.6\times10^4\) of the Higgs
doublet to the curvature of the fold, which the action of UM-4 does not
contain: \(n_{\mathrm{S}} = 0.964\), \(r_{\mathrm{T}} = 0.004\), and
\(H_{\mathrm{i}} = 1.6\times10^{13}\) GeV, so UM-4 would need
\(H > 1.6\times10^{14}\) GeV). Status: DERIVED
(\texttt{um4\_primordial.py}, PR2 and PR3), ESTIMATE (PR4, Higgs
inflation: an extension outside the action), OPEN (the Weyl tensor of
the bulk, the origin of the holding pressure).

\hypertarget{the-coupled-system}{%
\section{The coupled system}\label{the-coupled-system}}

\hypertarget{which-field-appears-in-which-equation}{%
\subsection{Which field appears in which
equation}\label{which-field-appears-in-which-equation}}

\begin{longtable}[]{@{}
  >{\raggedright\arraybackslash}p{(\columnwidth - 4\tabcolsep) * \real{0.2333}}
  >{\raggedright\arraybackslash}p{(\columnwidth - 4\tabcolsep) * \real{0.3000}}
  >{\raggedright\arraybackslash}p{(\columnwidth - 4\tabcolsep) * \real{0.4667}}@{}}
\toprule\noalign{}
\begin{minipage}[b]{\linewidth}\raggedright
equation
\end{minipage} & \begin{minipage}[b]{\linewidth}\raggedright
the fields in it
\end{minipage} & \begin{minipage}[b]{\linewidth}\raggedright
how they enter
\end{minipage} \\
\midrule\noalign{}
\endhead
\bottomrule\noalign{}
\endlastfoot
bulk gravity FE-1, FE-29, FE-32, FE-34 & metric, dirac16complex, gauge
fields & the Einstein tensor of the metric; \(\mathcal{T}_\Psi\) and
\(\mathcal{T}_{\mathrm{gauge}}\) \\
fold junction FE-4, FE-5, FE-30, FE-35, FE-59, FE-71, FE-72 & metric,
every field of the fold, dirac16complex at the fold & the jump of the
extrinsic curvature (to the first order the Einstein tensor of the fold
divided by \(\kappa_7\)); \(\mathcal{T}_{\mathrm{H}}\),
\(\mathcal{T}_{\mathrm{f}}\), \(\mathcal{T}_{\mathrm{P}}\),
\(\mathcal{T}_{\mathrm{ff}}\); the induced gravity
\(\hat G/\kappa_{\mathrm{ind}}\) \\
dirac16complex FE-12, FE-13, FE-36 to FE-43, FE-85, FE-86 & metric,
itself, every Standard-Model field & the spin connection and the warped
frequency; the self-interaction; the fold operator \(\mathbb{M}\) of the
portals (the Higgs portal for \(\Psi_{\mathrm{L}}\) only) \\
gauge fields FE-16 to FE-18, FE-74, FE-79 to FE-82 & metric,
dirac16complex, Higgs, quarks and leptons & the metric of the bulk; at
the fold the currents of the fold, the portal factors \(\mathcal{K}_G\),
\(\mathcal{K}_W\), \(\mathcal{K}_B\), the dipole density
\(\mathcal{D}_X\) and the mass term of the Higgs vacuum \\
Higgs FE-15 & metric, dirac16complex, the gauge fields at the fold,
quarks and leptons & \(\hat h\) and \(\mathcal{Z}\);
\(c_\Phi\bar\Psi P_{\mathrm{L}}\Psi\); \(D_\mu\); the Yukawa source
\(\Theta\) \\
quarks and leptons FE-19 to FE-24 & metric, the gauge fields at the
fold, Higgs, dirac16complex & \(\hat h\); \(D_\mu\); the Yukawa
matrices; the potentials \(\mathcal{V}^{f}\) \\
neutrino flavours FE-48 to FE-53 & expansion, Higgs, dirac16complex,
matter, other neutrinos & \(a_4'\); \(v_{\mathrm{loc}}\) and
\(G_{\mathrm{F,loc}}\); \(V_\Psi\); \(\mathsf{V}_{\mathrm{m}}\) and
\(\mathsf{V}_{\nu\nu}\) \\
cosmology FE-54 to FE-58 & metric perturbations, dark matter \(\chi\),
baryons, photons, neutrinos & \(\phi_{\mathrm{B}}\) and
\(\psi_{\mathrm{B}}\); \(\Gamma_{\chi\mathrm{b}}\); Thomson scattering;
\(\rho_\Lambda\) \\
thermal history FE-59 to FE-62, FE-87, FE-88 & metric, plasma, dark
matter \(\chi\), right-handed neutrinos, fold fluid & \(H_4\),
\(G_{\mathrm{N}}\) and \(\rho_\Lambda\);
\(\langle\sigma_{\mathrm{ann}}\mathsf{v}\rangle\);
\(\mathcal{C}_{\nu_{\mathrm{R}}}\); \(q_{\mathrm{hold}}\) \\
\end{longtable}

\hypertarget{how-the-system-is-solved}{%
\subsection{How the system is solved}\label{how-the-system-is-solved}}

The equations are coupled in a chain that UM-4 solves in steps, each
with the results of the earlier ones:

\begin{enumerate}
\def\labelenumi{\arabic{enumi}.}
\tightlist
\item
  The background: the static block (FE-32, exact: a space of constant
  curvature \(-H^2\) on each side of a fold of pure tension) and the
  expanding fold with frozen extra times to the first order in its
  curvature over \(H^2\) (FE-34, FE-71, FE-72: the Friedmann equations
  of general relativity), from an inflation of the fold (FE-88) to the
  late universe; the fold fluid holds the volume of the extra times
  fixed (FE-65; ASSUMED, its origin OPEN). The exact de Sitter fold
  FE-87 checks the first-order equations and shows why the volume must
  be held.
\item
  The propagation in the bulk (FE-73 to FE-84): the zero modes of
  gravity and of the gauge fields at the speed of light, the photon and
  the gravitational wave as waves of spin 1 and 2, the propagators from
  one point of the fold to another, and the bounds on \(H\) that the
  massive modes give (FE-81, FE-82). At energies far below \(H\) the
  zero modes of the graviton and of the gauge fields form a
  four-dimensional effective theory on the fold, in which the radiation
  of the gauge fields enters the Friedmann equations as four-dimensional
  radiation (the low-energy limit of FE-74 and FE-79; the source term of
  the bulk in the exact equations FE-35 is not computed).
\item
  The states of dirac16complex at the fold on the background: the
  resonances (FE-36 to FE-40, FE-85, FE-86) with the mean fields of the
  Standard-Model vacuum (FE-39), and their lifetime.
\item
  The Kohn-Sham states (FE-41, FE-42): the level of the resonances
  filled to a given density, with the self-interaction and the mean
  fields, to self-consistency; their energy-momentum tensor enters the
  bulk equations and the fold (their density integrated across the fold
  counts as matter of the fold, and their decay into the bulk as the
  flux of FE-11).
\item
  The Standard Model on the fold (FE-44 to FE-47): the couplings of the
  zero modes, the masses with the ad hoc vacuum, the residuals FE-26,
  the widths of the W and Z into the bulk (FE-82), the couplings of the
  particles \(\chi\).
\item
  The neutrino flavours (FE-48 to FE-53), the thermal history (FE-59 to
  FE-62) and the abundance of \(\chi\) (whose origin is OPEN).
\item
  The cosmological perturbations (FE-54 to FE-58): the CMB, the CNB and
  the BAO.
\end{enumerate}

\hypertarget{what-was-verified}{%
\section{What was verified}\label{what-was-verified}}

Every statement labelled DERIVED in this document rests on a computation
in the folder \texttt{universe\_model\_4/code/} of the repository, which
checks itself and writes a report to
\texttt{universe\_model\_4/reports/} (the equations and results
themselves are all written out above). The 189 checks of the 20 modules,
all passed:

\begin{itemize}
\tightlist
\item
  \texttt{um4\_dirac\_matrices.py} (55 checks): the author's matrices,
  their reality, the Clifford relations, the chiral projectors, Pin(4,4)
  (the matrices themselves are written out in FD-13 of the definition of
  the fold above).
\item
  \texttt{um4\_geometry.py} (11 checks): the proper chart with the
  warped time and the determinant of the metric (G1), the exact
  curvature of the metric of UM-4 for any \(a_4\), \(b_4\) on both
  halves of the period (G2), the static block of constant curvature with
  a vanishing Weyl tensor (G3), its Lovelock tensors (G4), the exact
  bulk equations FE-29, with the Einstein tensor computed in two ways
  (G5), the Ricci-flat folds FE-33 (G6), the junction FE-30 of a pure
  tension from the computed jump of the extrinsic curvature (G7), the
  light cone of every slice (G8), the spin connection FE-14 and the
  vielbein postulate on both halves with the author's matrices (G9), the
  static warps of the author's form, whose only vacua of Einstein
  gravity have equal rates in all seven directions of the fold (G10),
  and the matrices of FD-13 against the author's notebook (G11).
\item
  \texttt{um4\_gravity.py} (4 checks): the vacua of the static blocks
  for any weights FE-70 (L1), the graviton around the static block for
  any weights FE-68 (L2), the light cone of the gravitons on every
  warped background FE-69 (L3), the choice of Einstein gravity (L4).
\item
  \texttt{um4\_fold.py} (6 checks): the bulk metric of an expanding fold
  to the first order FE-34 (F1), the fold as a seven-dimensional
  Einstein brane FE-71 (F2), the exact equations of the fold FE-35 (F3),
  the Friedmann equations and the radion condition FE-72 (F4), the size
  of the corrections (F5), the Weyl term that dust requires (F6).
\item
  \texttt{um4\_foldsources.py} (5 checks): the tension of a fold in
  \(n\) dimensions FE-30 (FS1), the gravitational wave of the fold FE-63
  (FS2), the induced gravity FE-64 (FS3), the tuning of the tension
  (FS4), the fold fluid FE-65 (FS5).
\item
  \texttt{um4\_background.py} (6 checks): the vacua of the author's form
  (BG1), the expanding fold FE-66 and FE-72 (BG2), the free volume of
  the extra times excluded by lunar laser ranging (BG3), the factor
  \(\mathcal{Z}\) of the couplings (BG4), the gauge equivalence of the
  ad hoc vacuum and its radial residual (BG5), with numbers for three
  values of \(a_{40}\) (BG6).
\item
  \texttt{um4\_bulk.py} (10 checks): the modes of a scalar FE-73 (B1)
  and of a gauge field FE-74 (B2) of the bulk, the photon as a
  transverse wave of spin 1 FE-75 (B3), the graviton of the bulk (B4)
  and the gravitational wave with its plus and cross polarisations FE-76
  (B5), the bulk fermion FE-77 (B6), the boost invariance of
  dirac16complex FE-78 (B7), the speeds of UM-3 recomputed (B9), the
  massive modes of FE-73 and FE-74, a continuum from zero mass whose
  spectral density at the fold reproduces the propagators FE-79 (B10),
  and the signs of the kinetic terms of the 14 polarisations of the
  graviton and of the components of a gauge field, with the ghost
  vectors \(h_{\mu r}\) (B11).
\item
  \texttt{um4\_propagators.py} (5 checks): the propagators FE-79 (P1),
  their rational forms and expansions FE-80 (P2), Newton's and Coulomb's
  laws FE-81 (P3), the W and Z resonances FE-82 (P4), the bounds on
  \(H\) (P5).
\item
  \texttt{um4\_causality.py} (4 checks): the static block conformally
  flat FE-83, no shortcut at fixed extra times, and the null lines that
  move along an extra time (C1), the light cone of an expanding fold and
  the largest advance FE-84 (C2, C3), the closed time-like curves of the
  four time-like directions (C4).
\item
  \texttt{um4\_fields.py} (16 checks): the reality and chirality of the
  bilinears of dirac16complex, the Krein operator, the four-dimensional
  chirality, the electroweak masses and the photon of the ad hoc vacuum,
  its isospin charge and residual, the Yukawa invariance, the anomalies
  with three right-handed neutrinos, the electric charges.
\item
  \texttt{um4\_spectra.py} (8 checks): the reduction to eight copies
  (U1), the exact solutions FE-85 (U2), the resonances FE-40 with the
  \(16 \times 16\) matching (U3), their widths FE-86 (U4), their
  dispersion (U5), the lifetime condition (U6), their Krein form (U7)
  and their mirror (U8).
\item
  \texttt{um4\_kohnsham.py} (10 checks): the Kohn-Sham equations FE-41
  for a homogeneous gas of the particles \(\chi\), solved
  self-consistently with the author's matrices (reduced exactly by the
  rotations of ordinary space and of the extra times, and lifted back to
  sixteen components for the checks): the free level FE-40 and the box
  method against the exact resonance (K1), first-order perturbation
  theory in the dilute limit and the splitting of the multiplets (K2),
  the residuals of FE-41 and of the matching (K3), the exponent of the
  dilute limit in the box (K4), the gas at the fold within the box and
  the attractive case (K5), the Higgs-portal loop through FE-49 (K6),
  the equation of state from FE-42 (K7), the width that grows with the
  self-interaction and the factor \(1/\mathrm{e}\) of FE-86 (K8), a
  weakly bound point (K9), and how dilute the gas of the universe is
  (K10).
\item
  \texttt{um4\_krein.py} (9 checks): the Krein form and its quantisation
  FE-67 (KR1), the symmetries at rest and in motion (KR2), the mirrors
  of the fold (KR3), the flat directions of the dilute gas and its
  truncated response (KR4), the partner level of the gas (KR5), the
  selection rules at the fold with the invariant amplitudes (KR6), the
  consequences for stability and unitarity (KR7), the occupation of the
  level as a saddle of the first-order energy (KR8), the pairs that the
  convective current creates (KR9).
\item
  \texttt{um4\_extratimes.py} (5 checks): the dispersion with a momentum
  along the extra times on the fold (EX1), the states of dirac16complex
  with such a momentum FE-38 (EX2), the growth with frozen extra times
  (EX3), compact extra times in the late universe: no size helps (EX4),
  the truncation to fields without that momentum (EX5).
\item
  \texttt{um4\_primordial.py} (6 checks): the author's vacuum no longer
  a solution FE-31 (PR0), the exact de Sitter fold FE-87 (PR1), the
  inflation with frozen extra times FE-88 (PR2), its scales (PR3), Higgs
  inflation, an extension outside the action (PR4), the factor
  \(\mathcal{Z}\) during inflation (PR5).
\item
  \texttt{um4\_neutrino.py} (5 checks): the flavour blindness of
  \(V_\Psi\) and the invariance of the mixing under \(v_{\mathrm{loc}}\)
  (N1), the resonance formulas FE-50 against an exact integration (N2),
  the crossing criterion of FE-51 (N3), the conservation laws of FE-52
  (N4), the CNB numbers of FE-53 (N5).
\item
  \texttt{um4\_flavour.py} (4 checks): the survival probability of solar
  neutrinos by exact integration of FE-48 and FE-49, against the
  adiabatic formula of FE-50 and the measurements of Borexino; the
  three-flavour conversion in the envelope of a supernova for both mass
  orderings; the fast instability near a neutrinosphere and the
  dispersion relation FE-51 at its solutions; the equipartition FE-52
  (its conservation laws, its fixed point, the bound
  \(0 \le \mathcal{F} \le 1\)) and its effect on the heating behind the
  shock.
\item
  \texttt{um4\_numass.py} (5 checks): the neutrino masses of both
  orderings (NM1), the bounds of DESI DR2 and KATRIN (NM2), the Dirac
  phase in the results of UM-4 (NM3), what the lightest mass changes
  (NM4), the record against NuFIT 6.0 (NM5).
\item
  \texttt{um4\_cosmology.py} (6 checks): the expansion rate of FE-59 and
  the ruler FE-56 computed independently and by the Boltzmann code CAMB;
  the derived quantities of the CMB against Planck 2018; the CMB spectra
  (FE-57) against the Planck 2018 high-\(l\) data, with a refit for
  thermal right-handed neutrinos; the BAO distances against DESI DR2
  (the second data release of the Dark Energy Spectroscopic Instrument);
  the CNB numbers of FE-53 and FE-58.
\item
  \texttt{um4\_thermal.py} (9 checks): \(g_\rho\) and \(g_{\mathrm{s}}\)
  of FE-60 (T1), the relic density FE-61 by numerical solution,
  converged in the end point (T2), the current portals with the
  convective current against the limits of direct detection (T3), the
  Higgs-portal coupling of FE-47 and the binding of \(\chi\) (T4), the
  right-handed neutrinos and \(N_{\mathrm{eff}}\) FE-62 (T5), the ad hoc
  vacuum before nucleosynthesis: the helium abundance from the
  neutron-proton freeze-out and the energy of the ad hoc vacuum (T6),
  the thermal relic and the lifetime: only strongly bound particles
  \(\chi\) (T7), the binding that direct detection allows and the
  lifetime that it gives (T8), the binding by the gluon condensate (T9).
\end{itemize}

The helper modules \texttt{um4\_common.py} (the report),
\texttt{um4\_gammas.py} (the author's matrices) and
\texttt{um4\_lovelock.py} (the Lovelock tensors from the frame
curvature) have no checks of their own. The Kohn-Sham module runs in
parts (\texttt{um4\_kohnsham.py\ part\ NAME} for NAME = K1, I, K6, K9,
then \texttt{um4\_kohnsham.py\ assemble}), so that no single calculation
takes longer than 423 seconds; the Krein module reads the snapshots of
the parts, the modules of the sources of the fold, of the causal
structure and of the spectra the bounds of the propagators, the assembly
of the Kohn-Sham module the lifetimes of the spectra, and the thermal
module the lifetimes of the spectra and the pair amplitudes of the Krein
module.

\hypertarget{further-symbols}{%
\section{Further symbols}\label{further-symbols}}

\textbf{Further symbols.} These symbols appear in the definitions of the
symbols listed after the equations of this document.

\begin{itemize}
\tightlist
\item
  \(M_{\mathrm{red}}\): The reduced Planck mass,
  \(M_{\mathrm{red}} = (8\pi G_{\mathrm{N}})^{-1/2}\). Unit: GeV. Value:
  \(2.435\times10^{18}\) GeV. Status: STANDARD.
\item
  \(\theta_{12},\ \theta_{13},\ \theta_{23}\): The three mixing angles
  of \(U_{\mathrm{PMNS}}\) (its fourth parameter is the Dirac phase,
  which does not enter the electron-flavour results of UM-4). Unit:
  dimensionless. Value: \(33.40^\circ\), \(8.54^\circ\), \(49.14^\circ\)
  in the computations of UM-4 (\(\sin^2\theta_{12} = 0.303\),
  \(\sin^2\theta_{13} = 0.02203\), \(\sin^2\theta_{23} = 0.572\): NuFIT
  5.2, November 2022, normal ordering, the fit without the
  Super-Kamiokande atmospheric data); \(33.71^\circ\), \(8.56^\circ\),
  \(43.28^\circ\) in NuFIT 6.0 (2024, with the Super-Kamiokande
  atmospheric data; \(\sin^2\theta_{12} = 0.308\),
  \(\sin^2\theta_{13} = 0.02215\), \(\sin^2\theta_{23} = 0.470\)).
  Status: STANDARD.
\item
  \(L_{\mathrm{t}}\): The physical period of the extra times, if they
  are compact. Unit: length. Status: ASSUMED.
\item
  \(\mathcal{W}\): The warp factor: in a warped metric, the function of
  the position along one direction that multiplies the squared steps of
  other directions (its square root multiplies their lengths); in UM-4
  the function of \(\sigma\) that multiplies all seven directions of the
  fold, ordinary space, the time and the extra times,
  \(\mathcal{W} = \mathrm{e}^{-2H\lvert\sigma\rvert} = (\sin^2z)^{1/6}\),
  equal to 1 on the fold. Unit: dimensionless. Status: DERIVED.
\item
  \(\alpha_3\): The strong coupling constant,
  \(\alpha_3 = g_3^2/(4\pi)\); \(\langle(\alpha_3/\pi)G^2\rangle\) is
  the gluon condensate of the vacuum, 0.012 GeV\(^4\) (Shifman,
  Vainshtein and Zakharov). Unit: dimensionless. Status: STANDARD.
\item
  \(c_{\mathrm{b}}\): The speed coefficient of the particles \(\chi\) of
  dirac16complex along ordinary space, in units of the speed of light on
  the fold: at small momentum
  \(\omega^2 = \omega_0^2 + c_{\mathrm{b}}^2\vec k^2\). In UM-3, whose
  time was not warped, \(c_{\mathrm{b}}^2 = 1.19\) for
  \(H = 0.05\,m_\Psi\) and \(q_{\mathrm{b}} = 0.3\,m_\Psi\) (faster than
  light); in UM-4 the boosts of the fold are a symmetry and
  \(c_{\mathrm{b}} = 1\) exactly. Unit: dimensionless. Value: 1 in UM-4;
  1.09 in UM-3. Status: DERIVED.
\item
  \(\epsilon_{\mathrm{c}}\): The bookkeeping parameter of the
  cosmological perturbations of the metric of the fold (the size of the
  CMB anisotropies). Unit: dimensionless. Value: \(\sim 10^{-5}\).
  Status: STANDARD.
\item
  \(f_{\mathrm{N}}\): The effective coupling of \(\chi\) to a nucleon in
  the coherent scattering (a vector or scalar contact interaction).
  Unit: GeV\(^{-2}\). Status: ESTIMATE.
\item
  \(f_{\mathrm{TG}}\): The fraction of the nucleon mass carried by the
  gluons through the trace anomaly (one minus the fractions carried by
  the masses of the u, d and s quarks): the matrix element of
  \((\alpha_3/\pi)G^2\) in a nucleon at rest is
  \(-\frac89m_{\mathrm{N}}f_{\mathrm{TG}}\). Unit: dimensionless. Value:
  0.92. Status: STANDARD.
\item
  \(\mathfrak{g}_{\mathrm{q}},\ \mathfrak{g}_{\mathrm{l}},\ \mathfrak{g}_{\mathrm{p}}\):
  The current-portal couplings \(\mathfrak{g}_f\) taken equal for all
  chiral quark states (\(\mathfrak{g}_{\mathrm{q}}\)) or for the nine
  chiral lepton states that take part, the six of the lepton doublets
  and the three right-handed charged leptons
  (\(\mathfrak{g}_{\mathrm{l}}\)), and the coupling to the proton that
  the lepton loop induces (\(\mathfrak{g}_{\mathrm{p}}\)). Unit:
  GeV\(^{-2}\). Status: ESTIMATE.
\item
  \(m_{\mathrm{N}}\): The mass of a nucleon. Unit: GeV. Value: 0.939
  GeV. Status: STANDARD.
\item
  \(\mu_{\mathrm{red}}\): The reduced mass of \(\chi\) and a nucleon,
  \(\omega_{\mathrm{b}}m_{\mathrm{N}}/(\omega_{\mathrm{b}} + m_{\mathrm{N}})\).
  Unit: GeV. Status: STANDARD.
\item
  \(r_{\mathrm{T}}\): The ratio of the amplitudes of the primordial
  spectra of gravitational waves (tensor) and of curvature (scalar) at
  the pivot scale. Unit: dimensionless. Value: below 0.032 (Planck PR4
  with BICEP/Keck 2018 and BAO, M. Tristram et al.~2022, 95 per cent).
  Status: STANDARD.
\item
  \(\sigma_{\mathrm{SI}}\): The spin-independent cross-section of the
  scattering of \(\chi\) on a nucleon (normalised per nucleon, as the
  direct-detection limits are quoted). Unit: GeV\(^{-2}\) (cm\(^2\)).
  Status: ESTIMATE.
\item
  \(w_\chi\): The equation of state of the gas of the particles
  \(\chi\): its pressure divided by its energy density,
  \(w_\chi = P_{\mathrm{KS}}/(n_\chi\epsilon_{\mathrm{p}})\) (0 for cold
  dark matter, \(-1\) for a vacuum energy). Unit: dimensionless. Status:
  DERIVED.
\item
  \(P_{\mathrm{KS}}\): The pressure of the gas of \(\chi\) along
  ordinary space in the Kohn-Sham state: the particles are at rest, so
  it is the mean self-interaction,
  \(P_{\mathrm{KS}} = n_\chi\mathsf{E}_{\mathrm{int}}\). Unit:
  GeV\(^4\). Status: DERIVED.
\item
  \(\epsilon_{\mathrm{p}}\): The energy per particle of the gas of
  \(\chi\): the orbital energy minus the interaction energy per
  particle, which the orbital energy counts twice. Unit: GeV. Status:
  DERIVED.
\item
  \(\mathsf{E}_{\mathrm{int}}\): The interaction energy per particle of
  the gas of \(\chi\): the mean self-interaction energy
  \(\frac{\lambda_\Psi}{2}(n_{\mathrm{s}}^2 - \mathrm{tr}\,\mathcal{N}^2)\)
  integrated across the hidden direction, per particle. Unit: GeV.
  Status: DERIVED.
\end{itemize}

\end{document}
